% !TEX program = xelatex
\documentclass[11pt,oneside]{book}

% ==============================================================================
% 1. CÁC THƯ VIỆN & CẤU HÌNH GÓI (PACKAGES)
% ==============================================================================
\usepackage{silence}
\WarningFilter{tcolorbox}{Using nobreak failed}
\WarningFilter{babel}{No hyphenation patterns}
\WarningFilter{microtype}{Unable to apply patch}
\WarningFilter{rerunfilecheck}{File `main.out'}
\WarningFilter{latex}{Overfull}
\WarningFilter{latex}{Underfull}

\usepackage{fontspec}
\usepackage[vietnamese]{babel}
\shorthandoff{"}

\usepackage[
    paper=a4paper,
    top=2.6cm,
    bottom=2.6cm,
    left=2.5cm,
    right=2.5cm,
    headheight=26pt,
    headsep=20pt,
    footskip=30pt
]{geometry}

\usepackage{xcolor}
\usepackage{colortbl}
\definecolor{primaryBlue}{HTML}{0F4C81}     % Classic Navy/Royal Blue
\definecolor{secondaryTeal}{HTML}{008080}   % Teal Accent
\definecolor{darkText}{HTML}{1A1A1A}        % Màu chữ đen dịu
\definecolor{accentOrange}{HTML}{E65100}    % Cam điểm nhấn
\definecolor{accentPurple}{HTML}{5E35B1}    % Tím sang trọng
\definecolor{codeBg}{HTML}{F8FAFC}          % Nền code Slate nhạt
\definecolor{codeBorder}{HTML}{CBD5E1}      % Viền code Slate
\definecolor{codeComment}{HTML}{64748B}     % Màu chú thích Slate
\definecolor{codeKeyword}{HTML}{D97706}     % Từ khóa Cam/Hổ phách đậm
\definecolor{codeString}{HTML}{059669}      % Chuỗi ký tự Emerald
\definecolor{codeKey}{HTML}{0284C7}         % Tên thuộc tính Cyan
\definecolor{codeNumber}{HTML}{7C3AED}      % Giá trị số Indigo

\definecolor{noteBg}{HTML}{F0F7FF}
\definecolor{noteBorder}{HTML}{0F4C81}
\definecolor{tipBg}{HTML}{F0FDF4}
\definecolor{tipBorder}{HTML}{16A34A}
\definecolor{warnBg}{HTML}{FFFBEB}
\definecolor{warnBorder}{HTML}{D97706}
\definecolor{conceptBg}{HTML}{F5F3FF}
\definecolor{conceptBorder}{HTML}{7C3AED}

\usepackage{graphicx}
\usepackage{float}
\usepackage{caption}
\usepackage{subcaption}
\usepackage{tikz}
\usetikzlibrary{shapes, shadows, arrows.meta, positioning, calc, decorations.pathreplacing, fit}

\usepackage{tabularx}
\usepackage{booktabs}
\usepackage{multirow}
\usepackage{array}
\usepackage{ragged2e}

\usepackage[most]{tcolorbox}
\usepackage{listings}
\usepackage{fancyhdr}
\usepackage{titlesec}
\usepackage{titletoc}

\usepackage{hyperref}
\hypersetup{
    colorlinks=true,
    linkcolor=primaryBlue,
    citecolor=secondaryTeal,
    urlcolor=accentOrange,
    pdfauthor={Đội Ngũ Lập Trình Web},
    pdftitle={Giáo Trình CSS3 - Thiết Kế Giao Diện Web},
    pdfkeywords={CSS3, Flexbox, Grid, Responsive Design, CSS Variables, BEM, Animations, XeLaTeX}
}

\usepackage{indentfirst}
\setlength{\parindent}{18pt}
\setlength{\parskip}{2pt}

\usepackage{enumitem}
\usepackage[normalem]{ulem}
\setlist[itemize]{leftmargin=1.8em, itemsep=2pt, parsep=0pt, topsep=3pt, partopsep=0pt}
\setlist[enumerate]{leftmargin=1.8em, itemsep=2pt, parsep=0pt, topsep=3pt, partopsep=0pt}

\usepackage{setspace}
\setstretch{1.15}
\usepackage{needspace}
\usepackage{seqsplit}

% ==============================================================================
% 2. ĐỊNH NGHĨA KHUNG HỘP (CALLOUT BOXES) & NÂNG CAO (COMMANDS)
% ==============================================================================
\hfuzz=500pt
\vfuzz=500pt
\hbadness=10000
\vbadness=10000
\tcbset{lines before break=1}

\newtcolorbox{noteBox}[1][Ghi Chú Quan Trọng]{
    enhanced,
    before=\par\vspace{6pt}\noindent,
    after=\par\vspace{6pt},
    width=\linewidth,
    colback=noteBg,
    colframe=noteBorder,
    boxrule=0.5pt,
    arc=5pt,
    title={\small\sffamily\bfseries\color{noteBorder} [Ghi Chú]\quad #1},
    coltitle=noteBorder,
    fonttitle=\bfseries\sffamily,
    fontupper=\RaggedRight,
    attach title to upper,
    after title={\par\vspace{3pt}{\color{noteBorder!20}\hrule height 0.5pt}\par\vspace{5pt}},
    top=7pt, bottom=7pt, left=12pt, right=12pt
}

\newtcolorbox{tipBox}[1][Mẹo Thực Hành]{
    enhanced,
    before=\par\vspace{6pt}\noindent,
    after=\par\vspace{6pt},
    width=\linewidth,
    colback=tipBg,
    colframe=tipBorder,
    boxrule=0.5pt,
    arc=5pt,
    title={\small\sffamily\bfseries\color{tipBorder} [Mẹo]\quad #1},
    coltitle=tipBorder,
    fonttitle=\bfseries\sffamily,
    fontupper=\RaggedRight,
    attach title to upper,
    after title={\par\vspace{3pt}{\color{tipBorder!20}\hrule height 0.5pt}\par\vspace{5pt}},
    top=7pt, bottom=7pt, left=12pt, right=12pt
}

\newtcolorbox{warnBox}[1][Cảnh Báo \& Lỗi Thường Gặp]{
    enhanced,
    before=\par\vspace{6pt}\noindent,
    after=\par\vspace{6pt},
    width=\linewidth,
    colback=warnBg,
    colframe=warnBorder,
    boxrule=0.5pt,
    arc=5pt,
    title={\small\sffamily\bfseries\color{warnBorder} [Cảnh Báo]\quad #1},
    coltitle=warnBorder,
    fonttitle=\bfseries\sffamily,
    fontupper=\RaggedRight,
    attach title to upper,
    after title={\par\vspace{3pt}{\color{warnBorder!20}\hrule height 0.5pt}\par\vspace{5pt}},
    top=7pt, bottom=7pt, left=12pt, right=12pt
}

\newtcolorbox{conceptBox}[1][Khái Niệm Cốt Lõi]{
    enhanced,
    before=\par\vspace{6pt}\noindent,
    after=\par\vspace{6pt},
    width=\linewidth,
    colback=conceptBg,
    colframe=conceptBorder,
    boxrule=0.5pt,
    arc=5pt,
    title={\small\sffamily\bfseries\color{conceptBorder} [Khái Niệm]\quad #1},
    coltitle=conceptBorder,
    fonttitle=\bfseries\sffamily,
    fontupper=\RaggedRight,
    attach title to upper,
    after title={\par\vspace{3pt}{\color{conceptBorder!20}\hrule height 0.5pt}\par\vspace{5pt}},
    top=7pt, bottom=7pt, left=12pt, right=12pt
}

\newcommand{\macDots}{%
  \tikz[baseline=-0.5ex]{
    \fill[color=red!75!black] (0,0) circle (3pt);
    \fill[color=yellow!85!black] (8pt,0) circle (3pt);
    \fill[color=green!75!black] (16pt,0) circle (3pt);
  }\hspace{8pt}%
}

\newtcolorbox{codeWindow}[1][Cửa Sổ Mã Nguồn]{
    enhanced,
    before=\par\vspace{8pt}\noindent,
    after=\par\vspace{8pt},
    width=\linewidth,
    colback=codeBg,
    colframe=primaryBlue!30,
    colbacktitle=primaryBlue!8,
    boxrule=0.8pt,
    arc=7pt,
    title={\macDots \small\sffamily\bfseries\color{primaryBlue} #1},
    coltitle=primaryBlue,
    fonttitle=\bfseries\sffamily,
    top=6pt, bottom=4pt, left=8pt, right=8pt,
    toptitle=5pt, bottomtitle=5pt, lefttitle=10pt, righttitle=10pt
}

\newtcolorbox{codeWindowBreakable}[1][Cửa Sổ Mã Nguồn]{
    enhanced,
    breakable,
    before=\par\vspace{8pt}\noindent,
    after=\par\vspace{8pt},
    width=\linewidth,
    colback=codeBg,
    colframe=primaryBlue!30,
    colbacktitle=primaryBlue!8,
    boxrule=0.8pt,
    arc=7pt,
    title={\macDots \small\sffamily\bfseries\color{primaryBlue} #1},
    coltitle=primaryBlue,
    fonttitle=\bfseries\sffamily,
    top=6pt, bottom=2pt, left=8pt, right=8pt,
    toptitle=5pt, bottomtitle=5pt, lefttitle=10pt, righttitle=10pt
}

\newtcolorbox{browserWindow}[1][Mô Phỏng Trình Duyệt Web]{
    enhanced,
    before=\par\vspace{10pt}\noindent,
    after=\par\vspace{10pt},
    width=\linewidth,
    colback=white,
    colframe=primaryBlue!35,
    colbacktitle=gray!10,
    boxrule=0.8pt,
    arc=8pt,
    title={%
        \macDots \;\; {\small\sffamily\bfseries\color{primaryBlue} BROWSER PREVIEW} \hfill
        {\normalfont\small\sffamily\bfseries\color{darkText} #1}%
    },
    coltitle=darkText,
    fonttitle=\bfseries\sffamily,
    top=10pt, bottom=10pt, left=14pt, right=14pt,
    toptitle=6pt, bottomtitle=6pt, lefttitle=10pt, righttitle=10pt
}

\newtcolorbox{prettyTableBox}[1][Bảng Thống Kê]{
    enhanced,
    before=\par\vspace{8pt}\needspace{4.5\baselineskip}\noindent,
    after=\par\vspace{8pt},
    width=\linewidth,
    colback=white,
    colframe=primaryBlue!30,
    colbacktitle=primaryBlue!8,
    boxrule=0.8pt,
    arc=7pt,
    title={\small\sffamily\bfseries\color{primaryBlue} [Bảng]\quad #1},
    coltitle=primaryBlue,
    fonttitle=\bfseries\sffamily,
    top=0pt, bottom=0pt, left=0pt, right=0pt,
    toptitle=5pt, bottomtitle=5pt, lefttitle=10pt, righttitle=10pt
}

\renewcommand{\arraystretch}{1.2}
\arrayrulecolor{primaryBlue!20}

\newcolumntype{C}[1]{>{\centering\arraybackslash}m{#1}}
\newcolumntype{L}[1]{>{\raggedright\arraybackslash}m{#1}}
\newcolumntype{Y}{>{\centering\arraybackslash}X}
\newcolumntype{Z}{>{\raggedright\arraybackslash}X}
\newcolumntype{K}[1]{>{\ttfamily\color{accentOrange}\bfseries\raggedright\arraybackslash}m{#1}}

\newcommand{\thd}[1]{\textcolor{white}{\textbf{\strut #1}}}
\newcommand{\codeinline}[1]{\textcolor{accentOrange}{\texttt{#1}}}
\newcommand{\term}[1]{\textbf{\textcolor{primaryBlue}{#1}}}
\newcommand{\badge}[1]{\tcbox[on line, arc=3pt, colback=primaryBlue!10, colframe=primaryBlue!30, boxsep=1pt, left=3pt, right=3pt, top=1pt, bottom=1pt]{\small\sffamily\color{primaryBlue}#1}}

% ==============================================================================
% 3. CẤU HÌNH MÃ NGUỒN (CODE LISTINGS & UNICODE)
% ==============================================================================
\lstdefinelanguage{JavaScript}{
    keywords={break, case, catch, class, const, continue, debugger, default, delete, do, else, export, extends, false, finally, for, function, if, import, in, instanceof, new, null, return, super, switch, this, throw, true, try, typeof, var, void, while, with, yield, async, await, let},
    keywordstyle=\color{codeKeyword}\bfseries,
    ndkeywords={class, export, default, boolean, throw, implements, import, this, fetch, console, log, then},
    ndkeywordstyle=\color{codeKey}\bfseries,
    identifierstyle=\color{darkText},
    sensitive=true,
    comment=[l]{//},
    morecomment=[s]{/*}{*/},
    commentstyle=\color{codeComment}\upshape,
    stringstyle=\color{codeString},
    morestring=[b]',
    morestring=[b]",
    morestring=[b]`
}

\lstdefinelanguage{HTML5}{
    language=html,
    sensitive=true,
    alsoletter={<>=-/},
    morekeywords={
        <!DOCTYPE, html, head, title, meta, link, script, body,
        header, nav, main, section, article, aside, footer,
        h1, h2, h3, h4, h5, h6, p, a, img, ul, ol, li, div, span,
        form, input, button, select, option, textarea, label, table, tr, td, th,
        figure, figcaption, details, summary, dialog, meter, progress, picture, source, video, audio, track, iframe, svg, canvas, math
    },
    keywordstyle=\color{primaryBlue}\bfseries,
    tag=[s]{<}{>},
    tagstyle=\color{primaryBlue}\bfseries,
    stringstyle=\color{codeString},
    morestring=[b]",
    comment=[s]{<!--}{-->},
    commentstyle=\color{codeComment}\upshape
}

\lstdefinelanguage{CSS3}{
    keywords={color, background, background-color, border, border-radius, padding, margin, margin-bottom, font-weight, font-size, display, flex, grid, width, height, appearance},
    keywordstyle=\color{primaryBlue}\bfseries,
    sensitive=false,
    comment=[s]{/*}{*/},
    commentstyle=\color{codeComment}\upshape,
    stringstyle=\color{codeString},
    morestring=[b]",
    morestring=[b]'
}
\lstdefinelanguage{CSS}{language=CSS3}

\lstdefinelanguage{PHP}{
    keywords={<?php, ?>, echo, print, date, if, else, foreach, while, return, function, class, array},
    keywordstyle=\color{codeKeyword}\bfseries,
    ndkeywords={echo, print, date},
    ndkeywordstyle=\color{codeKey}\bfseries,
    stringstyle=\color{codeString},
    morestring=[b]",
    morestring=[b]',
    comment=[l]{//},
    commentstyle=\color{codeComment}\upshape
}

\lstdefinelanguage{HTTP}{
    keywords={GET, POST, PUT, DELETE, PATCH, HEAD, OPTIONS, HTTP, HTTPS, OK},
    keywordstyle=\color{codeKeyword}\bfseries,
    ndkeywords={Host, Content-Type, Content-Length, Cache-Control, Authorization, User-Agent, Accept, Connection, Server, Date, Status, Bearer, status, data, id, name, role},
    ndkeywordstyle=\color{codeKey}\bfseries,
    identifierstyle=\color{darkText},
    sensitive=false,
    stringstyle=\color{codeString},
    morestring=[b]",
    morestring=[b]',
    comment=[l]{//},
    commentstyle=\color{codeComment}\upshape
}

\lstdefinestyle{webcode}{
    backgroundcolor=\color{codeBg},
    basicstyle=\ttfamily\small,
    columns=fullflexible,
    keepspaces=true,
    breakatwhitespace=false,
    breaklines=true,
    captionpos=b,
    deletekeywords={...},
    escapeinside={(*@}{@*)},
    extendedchars=false,
    inputencoding=utf8,
    firstnumber=1,
    frame=none,
    rulecolor=\color{codeBorder},
    language=JavaScript,
    numbersep=8pt,
    numbers=left,
    numberstyle=\ttfamily\scriptsize\color{codeComment},
    showspaces=false,
    showstringspaces=false,
    showtabs=false,
    stepnumber=1,
    tabsize=2,
    aboveskip=3pt,
    belowskip=1pt,
    xleftmargin=16pt,
    xrightmargin=4pt
}

\lstset{style=webcode}

% ==============================================================================
% 4. PHÔNG CHỮ & ĐỊNH DẠNG TIÊU ĐỀ (TITLES & FONTS)
% ==============================================================================
\setmainfont{Times New Roman}[
    BoldFont = Times New Roman Bold,
    ItalicFont = Times New Roman Italic,
    BoldItalicFont = Times New Roman Bold Italic
]
\setsansfont{Arial}[
    BoldFont = Arial Bold,
    ItalicFont = Arial Italic,
    BoldItalicFont = Arial Bold Italic
]
\setmonofont{Consolas}[Scale=0.92]

\titleformat{\chapter}[display]
  {\normalfont\sffamily}
  {
    \begin{tikzpicture}
      \node[fill=primaryBlue, text=white, font=\sffamily\bfseries\small, inner xsep=10pt, inner ysep=4pt, rounded corners=3pt] {CHƯƠNG \thechapter};
    \end{tikzpicture}
  }
  {6pt}
  {\Huge\bfseries\color{darkText}\raggedright}
  [\vspace{8pt}{\color{primaryBlue}\hrule height 1.5pt}\vspace{15pt}]

\titlespacing*{\chapter}{0pt}{-10pt}{20pt}

\titleformat{name=\chapter,numberless}[block]
  {\normalfont\Large\sffamily\bfseries\color{primaryBlue}}
  {}
  {0pt}
  {}
  [\vspace{4pt}{\color{accentOrange}\rule{5cm}{1.2pt}}\vspace{12pt}]

\titlespacing*{name=\chapter,numberless}{0pt}{0pt}{15pt}

\titleformat{\section}
  {\normalfont\Large\sffamily\bfseries\color{primaryBlue}}
  {\thesection}
  {10pt}
  {}

\titlespacing*{\section}{0pt}{12pt}{6pt}

\titleformat{\subsection}
  {\normalfont\large\sffamily\bfseries\color{secondaryTeal}}
  {\thesubsection}
  {8pt}
  {}

\titlespacing*{\subsection}{0pt}{9pt}{4pt}

\titleformat{\paragraph}[hang]
  {\normalfont\sffamily\bfseries\color{darkText}}
  {\theparagraph}
  {0pt}
  {}

\titlespacing*{\paragraph}{0pt}{6pt}{3pt}

\titleformat{\part}[display]
  {\normalfont\sffamily\centering}
  {
    \begin{tikzpicture}
      \node[fill=primaryBlue, text=white, font=\sffamily\bfseries\Large, inner xsep=18pt, inner ysep=6pt, rounded corners=4pt] {PHẦN \thepart};
    \end{tikzpicture}
  }
  {12pt}
  {\Huge\bfseries\color{darkText}}
  [\vspace{15pt}\centerline{\color{primaryBlue!40}\rule{8cm}{1pt}}]

\titlespacing*{\part}{0pt}{100pt}{0pt}

\pagestyle{fancy}
\fancyhf{}
\fancyhead[R]{\sffamily\small\color{darkText}\rightmark}
\fancyfoot[R]{\sffamily\bfseries\color{primaryBlue}\thepage}
\fancyfoot[C]{\sffamily\small\color{codeComment}Giáo Trình Chuyên Sâu CSS3 \& Thiết Kế Giao Diện Web}

\renewcommand{\headrulewidth}{0.6pt}
\renewcommand{\headrule}{{\color{primaryBlue!40}\hrule width\headwidth height\headrulewidth}}
\renewcommand{\footrulewidth}{0.4pt}
\renewcommand{\footrule}{{\color{codeBorder}\hrule width\headwidth height\footrulewidth}}

\fancypagestyle{plain}{
  \fancyhf{}
  \fancyfoot[R]{\sffamily\bfseries\color{primaryBlue}\thepage}
  \renewcommand{\headrulewidth}{0pt}
  \renewcommand{\footrulewidth}{0.4pt}
}

% ==============================================================================
% BẮT ĐẦU TÀI LIỆU
% ==============================================================================
\begin{document}
\shorthandoff{"}

% ------------------------------------------------------------------------------
% TRANG BÌA CHÍNH (CSS MASTERCLASS COVER PAGE)
% ------------------------------------------------------------------------------
\begin{titlepage}
    \begin{tikzpicture}[remember picture, overlay]
        \fill[primaryBlue] (current page.north west) -- (current page.north east) -- ($(current page.north east)+(0,-10.5cm)$) -- ($(current page.north west)+(0,-8.5cm)$) -- cycle;
        \fill[accentOrange] ($(current page.north west)+(0,-8.5cm)$) -- ($(current page.north east)+(0,-10.5cm)$) -- ($(current page.north east)+(0,-10.9cm)$) -- ($(current page.north west)+(0,-8.9cm)$) -- cycle;
        \fill[white!15] ($(current page.north east)+(-6cm,0)$) -- (current page.north east) -- ($(current page.north east)+(0,-4.5cm)$) -- cycle;

        \fill[primaryBlue] ($(current page.south west)+(0,0)$) rectangle ($(current page.south east)+(0,1.3cm)$);
        \fill[accentOrange] ($(current page.south west)+(0,1.3cm)$) rectangle ($(current page.south east)+(0,1.48cm)$);

        \node[anchor=north west] at ($(current page.north west)+(2.2cm,-1.5cm)$) {
            {\sffamily \bfseries \small \color{white!85} GIÁO TRÌNH ĐÀO TẠO KỸ SƯ FRONTEND CHUYÊN SÂU}
        };

        \node[anchor=north west, text width=16cm] at ($(current page.north west)+(2.2cm,-2.8cm)$) {
            {\fontsize{34pt}{40pt}\selectfont \sffamily \bfseries \color{white} GIÁO TRÌNH CSS3} \\[0.4cm]
            {\fontsize{24pt}{30pt}\selectfont \sffamily \bfseries \color{accentOrange} THIẾT KẾ GIAO DIỆN WEB} \\[0.5cm]
            {\large \sffamily \color{white!95} Từ Nền Tảng Layout đến Flexbox, Grid, CSS Variables \& Animations}
        };

        \node[anchor=north west] at ($(current page.north west)+(2.2cm,-9.4cm)$) {
            \begin{tcolorbox}[
                on line,
                colback=white,
                colframe=primaryBlue!40,
                boxrule=1.2pt,
                arc=16pt,
                left=16pt, right=16pt, top=8pt, bottom=8pt,
                drop shadow=black!12!white
            ]
                {\sffamily \bfseries \small \color{primaryBlue}
                \textcolor{accentOrange}{\rule{5pt}{5pt}}\; CSS Foundations \quad$\bullet$\quad 
                \textcolor{accentOrange}{\rule{5pt}{5pt}}\; Flexbox \& Grid \quad$\bullet$\quad 
                \textcolor{accentOrange}{\rule{5pt}{5pt}}\; Responsive \quad$\bullet$\quad 
                \textcolor{accentOrange}{\rule{5pt}{5pt}}\; BEM \& Variables}
            \end{tcolorbox}
        };

        \node[anchor=center, text width=15.5cm, align=center] at ($(current page.center)+(0,-2.5cm)$) {
            \begin{tcolorbox}[
                colback=white,
                colframe=primaryBlue!25,
                boxrule=1pt,
                arc=10pt,
                top=16pt, bottom=16pt, left=20pt, right=20pt,
                drop shadow=black!6!white
            ]
                \centering
                {\Large \sffamily \bfseries \color{primaryBlue} TÀI LIỆU CHUYÊN SÂU TẬP TRUNG CSS3} \\[0.3cm]
                {\color{accentOrange}\rule{6cm}{1.2pt}} \\[0.3cm]
                {\small \sffamily \color{darkText} 
                Tệp LaTeX độc lập tích hợp toàn bộ cấu hình \& gói thư viện $\bullet$ Sẵn sàng di chuyển sang mọi trình biên dịch}
            \end{tcolorbox}
        };

        \node[anchor=south west] at ($(current page.south west)+(2.2cm,1.9cm)$) {
            {\sffamily \bfseries \small \color{primaryBlue} ĐỘI NGŨ BIÊN SOẠN WEB} \\
            {\sffamily \scriptsize \color{codeComment} Tài liệu lưu hành nội bộ}
        };

        \node[anchor=south east] at ($(current page.south east)+(-2.2cm,1.9cm)$) {
            {\sffamily \bfseries \small \color{primaryBlue} NĂM 2026} \\
            {\sffamily \scriptsize \color{codeComment} Phiên bản Standalone}
        };

        \node[anchor=south, color=white, font=\sffamily\scriptsize\bfseries] at ($(current page.south)+(0,0.5cm)$) {
            CHINH PHỤC NGHỆ THUẬT THIẾT KẾ GIAO DIỆN WEB HIỆN ĐẠI
        };
    \end{tikzpicture}
\end{titlepage}

% ------------------------------------------------------------------------------
% TRANG BÌA PHỤ (SECONDARY COVER PAGE)
% ------------------------------------------------------------------------------
\clearpage
\thispagestyle{empty}
\begin{tikzpicture}[remember picture, overlay]
    \draw[primaryBlue, line width=1.4pt] 
        ($(current page.north west)+(2.0cm,-2.0cm)$) rectangle ($(current page.south east)+(-2.0cm,2.0cm)$);
    \draw[accentOrange, line width=0.7pt] 
        ($(current page.north west)+(2.15cm,-2.15cm)$) rectangle ($(current page.south east)+(-2.15cm,2.15cm)$);

    \node[anchor=north, align=center, text width=15cm] at ($(current page.north)+(0,-3.5cm)$) {
        {\small \sffamily \bfseries \color{primaryBlue} ĐỘI NGŨ BIÊN SOẠN WEB} \\[0.2cm]
        {\scriptsize \sffamily \color{codeComment} TÀI LIỆU LƯU HÀNH NỘI BỘ -- CHUYÊN ĐỀ CSS3}
    };

    \node[anchor=center, align=center, text width=15cm] at ($(current page.center)+(0,2.2cm)$) {
        {\fontsize{28pt}{34pt}\selectfont \sffamily \bfseries \color{primaryBlue} GIÁO TRÌNH CSS3} \\[0.4cm]
        {\fontsize{20pt}{26pt}\selectfont \sffamily \bfseries \color{accentOrange} THIẾT KẾ GIAO DIỆN WEB} \\[0.6cm]
        {\large \sffamily \color{darkText} Từ Nền Tảng Layout đến Flexbox, Grid, CSS Variables \& Animations}
    };

    \draw[primaryBlue!35, line width=0.8pt] ($(current page.center)+(-4cm,-1.0cm)$) -- ($(current page.center)+(4cm,-1.0cm)$);

    \node[anchor=center, align=center, text width=15cm] at ($(current page.center)+(0,-2.8cm)$) {
        {\large \sffamily \bfseries \color{primaryBlue} TẬP TRUNG BIÊN SOẠN \& THỰC HÀNH CSS3} \\[0.3cm]
        {\small \sffamily \color{codeComment} Biên soạn chi tiết $\bullet$ Minh họa trực quan $\bullet$ Tối ưu tốc độ biên dịch}
    };

    \node[anchor=south, align=center, text width=15cm] at ($(current page.south)+(0,3.2cm)$) {
        {\sffamily \bfseries \small \color{primaryBlue} NĂM 2026} \\[0.1cm]
        {\sffamily \scriptsize \color{darkText} Phiên bản XeLaTeX Unicode}
    };
\end{tikzpicture}
\clearpage

% ==============================================================================
% MỤC LỤC TÀI LIỆU (TABLE OF CONTENTS)
% ==============================================================================
\frontmatter
\renewcommand{\contentsname}{MỤC LỤC TÀI LIỆU CSS3}
\setcounter{tocdepth}{2}

\tableofcontents
\clearpage

% ==============================================================================
% NỘI DUNG GIÁO TRÌNH CSS3 CHUYÊN SÂU
% ==============================================================================
\mainmatter

\part{CSS3 -- Thiết Kế Giao Diện \& Bố Cục Trang Web}

% ==============================================================================
% CHƯƠNG 6: NỀN TẢNG CSS (CASCADE, SPECIFICITY, BOX MODEL \& OUTLINE)
% ==============================================================================
\chapter{Nền Tảng CSS: Cascading Style Sheets, Box Model \& Thuộc Tính Outline}
\label{chap:css_foundations}

% ==============================================================================
% MỤC 1: SƠ LƯỢC \& KHÁI NIỆM CSS
% ==============================================================================
\section{Sơ Lược \& Khái Niệm Về CSS}

CSS (\term{Cascading Style Sheets}) giúp bạn thiết kế và định dạng giao diện trang web. Nó quyết định mọi thứ về hình thức hiển thị, từ màu sắc, kích thước chữ đến bố cục và hiệu ứng.

\begin{conceptBox}[Khái Niệm CSS Là Gì?]
\textbf{CSS (Cascading Style Sheets)} là ngôn ngữ dùng để mô tả cách trình bày của một tài liệu HTML hoặc XML. Nó giúp định dạng màu sắc, kiểu chữ, bố cục và hiệu ứng cho trang web. CSS giúp \textbf{tách biệt phần nội dung (HTML) với phần hiển thị}, giúp trang web dễ quản lý, tăng tính bảo trì và tối ưu hiệu suất tải trang.
\end{conceptBox}

\subsection{Mục Đích Cốt Lõi Của Ngôn Ngữ CSS}
\begin{enumerate}[leftmargin=1.5em, itemsep=4pt, parsep=0pt]
    \item \textbf{Tách biệt nội dung và trình bày}: Thẻ HTML tập trung vào nội dung ngữ nghĩa, CSS lo phần thiết kế giao diện thị giác.
    \item \textbf{Giúp trang web đẹp và chuyên nghiệp}: Mang lại giao diện hiện đại, nâng cao trải nghiệm thị giác.
    \item \textbf{Dễ bảo trì và tái sử dụng}: Một tập tin CSS duy nhất có thể dùng và định kiểu đồng bộ cho nhiều trang web.
    \item \textbf{Cải thiện trải nghiệm người dùng và tối ưu tốc độ tải trang}: Giảm kích thước file HTML, tối ưu hóa bộ nhớ đệm trình duyệt (\term{browser cache}).
\end{enumerate}

% ==============================================================================
% MỤC 2: CÚ PHÁP CSS (CSS SYNTAX)
% ==============================================================================
\section{Cú Pháp CSS (CSS Syntax)}

\subsection{Cấu Trúc Cơ Bản Của Cú Pháp CSS}
Một quy tắc CSS bao gồm hai thành phần chính:
\begin{itemize}[leftmargin=1.5em, itemsep=2pt, parsep=0pt]
    \item \textbf{Bộ chọn (Selector)} -- Xác định phần tử HTML cần áp dụng CSS.
    \item \textbf{Khối khai báo (Declaration Block)} -- Chứa các thuộc tính và giá trị đặt trong cặp dấu ngoặc nhọn \codeinline{\{\}}.
\end{itemize}

\begin{codeWindow}[Cú Pháp Định Dạng CSS Chuẩn]
\begin{lstlisting}[language=CSS3]
selector {
    thuoc_tinh: gia_tri;
    thuoc_tinh: gia_tri;
}
\end{lstlisting}
\end{codeWindow}

\begin{browserWindow}[Kết Quả: Sơ Đồ Cấu Trúc Khai Báo CSS]
\begin{tcolorbox}[arc=4pt, colback=white, colframe=primaryBlue!30, boxrule=0.8pt, top=10pt, bottom=10pt, left=12pt, right=12pt]
  {\sffamily\bfseries\color{primaryBlue} Khái niệm cốt lõi:}\par\vspace{4pt}
  {\sffamily Trình duyệt sẽ tìm tất cả các phần tử khớp với \codeinline{selector} trên trang HTML và áp dụng lần lượt các thuộc tính định kiểu bên trong khối khai báo \codeinline{\{\}}.}
\end{tcolorbox}
\end{browserWindow}

\subsection{Ví Dụ Cụ Thể \& Giải Thích Cú Pháp}
\begin{codeWindow}[Khai Báo Style Cho Thẻ h1]
\begin{lstlisting}[language=CSS3]
h1 {
    color: red;
    font-size: 24px;
}
\end{lstlisting}
\end{codeWindow}

\begin{browserWindow}[Kết Quả Hiển Thị: Thẻ h1 Chữ Đỏ, Cỡ Chữ 24px]
\sffamily\small
{\color{red}\bfseries\Large Tiêu Đề Văn Bản h1} \par\vspace{4pt}
\textcolor{gray!70}{\small\textit{($\rightarrow$ Kết quả: Thẻ <h1> có chữ màu đỏ và kích thước 24px)}}
\end{browserWindow}

\begin{noteBox}[Giải Thích Chi Tiết Cú Pháp]
\begin{itemize}[leftmargin=1.5em, itemsep=2pt, parsep=0pt]
    \item \codeinline{h1} $\rightarrow$ Chọn tất cả thẻ \codeinline{<h1>} trên trang.
    \item \codeinline{color: red;} $\rightarrow$ Thiết lập màu chữ đỏ.
    \item \codeinline{font-size: 24px;} $\rightarrow$ Thiết lập cỡ chữ \codeinline{24px}.
\end{itemize}
\end{noteBox}

% ==============================================================================
% MỤC 3: CÁC CÁCH THÊM STYLE VÀO TRANG WEB
% ==============================================================================
\section{Các Cách Thêm Style Vào Trang Web}

CSS có thể được áp dụng vào HTML theo 3 cách:

\begin{prettyTableBox}[Bảng So Sánh 3 Cách Thêm Style Vào Trang Web]
\rowcolors{2}{white}{primaryBlue!4}
\sffamily\small
\begin{tabularx}{\linewidth}{|>{\raggedright\arraybackslash \bfseries\color{primaryBlue}}K{2.8cm}|Z|>{\sffamily\small}K{3.8cm}|>{\sffamily\small}K{3.8cm}|}
\hline
\rowcolor{primaryBlue} \thd{Cách viết} & \thd{Mô tả} & \thd{Ưu điểm} & \thd{Nhược điểm} \\
\hline
\textbf{Inline CSS} & Viết trực tiếp vào thẻ HTML bằng thuộc tính \codeinline{style}. & Dễ dùng, nhanh chóng. & Khó bảo trì, không tái sử dụng được. \\
\hline
\textbf{Internal CSS} & Viết trong thẻ \codeinline{<style>} trong file HTML. & Dễ chỉnh sửa toàn bộ trang, không cần file ngoài. & Không tách biệt nội dung và kiểu dáng. \\
\hline
\textbf{External CSS} & Link tới file \codeinline{.css} bên ngoài qua thẻ \codeinline{<link>}. & Tái sử dụng cao, dễ bảo trì. & Phải tải thêm file, trang load lâu hơn chút. \\
\hline
\end{tabularx}
\end{prettyTableBox}

\subsection{Inline CSS (CSS Trong Thẻ HTML)}
Viết trực tiếp vào thẻ HTML bằng thuộc tính \codeinline{style}.
\begin{codeWindow}[Mã Nguồn Inline CSS]
\begin{lstlisting}[language=HTML5]
<h1 style="color: red; font-size: 30px;">Header Title</h1>
\end{lstlisting}
\end{codeWindow}

\begin{browserWindow}[Kết Quả Hiển Thị Inline CSS]
\sffamily\small
\textcolor{red}{\Huge Tiêu đề đỏ}
\end{browserWindow}
\textcolor{warnBorder}{\textbf{$\rightarrow$ Hạn chế}}: \textcolor{accentOrange}{\textbf{Không khuyến khích}} vì làm code khó bảo trì.

\subsection{Internal CSS (CSS Nội Bộ)}
Viết trong thẻ \codeinline{<style>} trong file HTML (bên trong \codeinline{<head>}). Dùng khi chỉ áp dụng CSS cho một trang web duy nhất.
\begin{codeWindow}[Mã Nguồn Internal CSS]
\begin{lstlisting}[language=HTML5]
<head>
  <style>
    h1 {
      color: blue;
      font-size: 24px;
    }
  </style>
</head>
\end{lstlisting}
\end{codeWindow}

\begin{browserWindow}[Kết Quả Hiển Thị Internal CSS]
\sffamily\small
\textcolor{primaryBlue}{\Large Tiêu Đề Xanh Dương}
\end{browserWindow}

\subsection{External CSS (CSS Bên Ngoài)}
Viết CSS trong file \codeinline{.css} riêng biệt, sau đó liên kết vào HTML bằng thẻ \codeinline{<link>}.
\begin{codeWindow}[Liên Kết File CSS Qua Thẻ Link]
\begin{lstlisting}[language=HTML5]
<head>
    <link rel="stylesheet" href="style.css">
</head>
\end{lstlisting}
\end{codeWindow}

\begin{codeWindow}[Nội Dung Trong File style.css]
\begin{lstlisting}[language=CSS3]
/* Sample Text */
p {
    color: green;
    font-size: 16px;
}
\end{lstlisting}
\end{codeWindow}

\begin{browserWindow}[Kết Quả Hiển Thị External CSS]
\sffamily\small
\textcolor{tipBorder}{Đoạn văn bản màu xanh lá được định kiểu từ file style.css bên ngoài.}
\end{browserWindow}
\textcolor{tipBorder}{\textbf{$\rightarrow$ Lời khuyên}}: \textcolor{primaryBlue}{\textbf{Khuyến khích sử dụng}} vì giúp code gọn gàng và dễ bảo trì.

% ==============================================================================
% MỤC 4: MỘT SỐ THUỘC TÍNH CSS CƠ BẢN \& ĐƠN VỊ ĐO KÍCH THƯỚC
% ==============================================================================
\section{Một Số Thuộc Tính CSS Cơ Bản \& Đơn Vị Đo Kích Thước}

\subsection{Bảng Tra Cứu Một Số Thuộc Tính CSS Cơ Bản}

\begin{prettyTableBox}[Bảng Tra Cứu Các Thuộc Tính CSS Cơ Bản Phổ Biến]
\rowcolors{2}{white}{primaryBlue!4}
\sffamily\small
\begin{tabularx}{\linewidth}{|>{\raggedright\arraybackslash \ttfamily\color{primaryBlue}\bfseries}K{2.8cm}|Z|>{\ttfamily\small}K{4.2cm}|}
\hline
\rowcolor{primaryBlue} \thd{\sffamily Thuộc tính} & \thd{\sffamily Chức năng} & \thd{\sffamily Ví dụ} \\
\hline
color & Đặt màu chữ. & color: red; \\
\hline
font-size & Đặt kích thước chữ. & font-size: 20px; \\
\hline
background & Đặt màu nền hoặc ảnh nền. & background: \#f0f0f0; \\
\hline
border & Tạo viền cho phần tử. & border: 2px solid black; \\
\hline
padding & Khoảng cách bên trong phần tử. & padding: 10px; \\
\hline
margin & Khoảng cách bên ngoài phần tử. & margin: 20px; \\
\hline
text-align & Căn chỉnh văn bản. & text-align: center; \\
\hline
display & Quyết định kiểu hiển thị của phần tử. & display: flex; \\
\hline
\end{tabularx}
\end{prettyTableBox}

\paragraph{Ví Dụ Kết Hợp Thuộc Tính Định Kiểu Cho Thẻ p:}
\begin{codeWindow}[Ví Dụ Định Kiểu Thẻ p Tổng Hợp]
\begin{lstlisting}[language=CSS3]
p {
    color: green;
    font-size: 18px;
    background: #e0f7fa;
    padding: 10px;
    border: 1px solid #00796b;
}
\end{lstlisting}
\end{codeWindow}

\begin{browserWindow}[Kết Quả Hiển Thị Của Thẻ p Định Kiểu]
\sffamily\small
\begin{tcolorbox}[arc=3pt, colback=cyan!8, colframe=teal!80!black, boxrule=1pt, boxsep=6pt]
\textcolor{teal!70!black}{\textbf{Đoạn văn có chữ màu xanh lá, nền xanh nhạt (\#e0f7fa), viền xanh đậm (\#00796b), khoảng cách nội dung padding 10px.}}
\end{tcolorbox} \par\vspace{4pt}
\textcolor{primaryBlue}{\textbf{$\rightarrow$ Kết quả}}: Đoạn văn sẽ có chữ màu xanh lá, nền xanh nhạt, viền xanh đậm, khoảng cách nội dung 10px.
\end{browserWindow}

% ==============================================================================
% MỤC 4.1: CÁC ĐƠN VỊ TRONG CSS (CSS UNITS MASTERCLASS TỪ A ĐẾN Z)
% ==============================================================================
\subsection{Chuyên Đề Các Đơn Vị Trong CSS (CSS Units Masterclass Từ A đến Z)}
\label{sec:css_units_masterclass}

Trong CSS, các đơn vị đo (\term{CSS Units}) quyết định kích thước của văn bản, khoảng cách cách lề (\codeinline{margin}, \codeinline{padding}), chiều rộng (\codeinline{width}), chiều cao (\codeinline{height}) và toàn bộ bố cục của trang web. Việc lựa chọn đúng loại đơn vị đo là yếu tố cốt lõi giúp giao diện web hiển thị chuẩn xác, đáp ứng linh hoạt trên mọi kích thước màn hình (\term{Responsive Web Design}).

\paragraph{I. Phân Loại Các Đơn Vị Trong CSS (CSS Units Classification)}

CSS Units được chia làm 2 nhóm chính:

\paragraph{1. Absolute Units (Đơn Vị Tuyệt Đối -- Không Thay Đổi Theo Bối Cảnh):}

Đơn vị tuyệt đối là các đơn vị cố định, có kích thước không bao giờ thay đổi bất kể kích thước màn hình, độ phân giải thiết bị hay bối cảnh phần tử cha.

\begin{prettyTableBox}[Bảng Chi Tiết Các Đơn Vị Tuyệt Đối (Absolute Units)]
\rowcolors{2}{white}{primaryBlue!4}
\sffamily\small
\begin{tabularx}{\linewidth}{|>{\raggedright\arraybackslash \ttfamily\color{primaryBlue}\bfseries}K{2.2cm}|Z|>{\sffamily\small}Z|}
\hline
\rowcolor{primaryBlue} \thd{Đơn Vị} & \thd{Ý Nghĩa} & \thd{Ghi Chú \& Đặc Điểm} \\
\hline
px & Pixels & Phổ biến nhất, kích thước cố định trên màn hình digital. \\
\hline
pt & Point ($1\text{pt} = 1/72\text{ inch}$) & Dành cho in ấn (ít dùng trên giao diện web). \\
\hline
cm & Centimeter & Dùng cho thiết kế in ấn, không phù hợp màn hình. \\
\hline
mm & Millimeter & Tương tự cm, dùng cho in ấn. \\
\hline
in & Inch & $1\text{in} = 2.54\text{cm} = 96\text{px}$. \\
\hline
pc & Pica ($1\text{pc} = 12\text{pt}$) & Hiếm dùng trên web. \\
\hline
\end{tabularx}
\end{prettyTableBox}

% ==============================================================================
% CHUYÊN ĐỀ SÂU VỀ ĐƠN VỊ PIXEL (PX MASTERCLASS TRONG CSS)
% ==============================================================================
\paragraph{Chuyên Đề Chi Tiết Về Đơn Vị Pixel (px Masterclass Trong CSS)}
\label{sec:css_pixel_masterclass}

Trong CSS, \codeinline{px} (viết tắt của \term{Pixel}) là đơn vị độ dài tuyệt đối vô cùng phổ biến và là một trong những đơn vị cơ bản nhất được sử dụng để xác định kích thước cho các thuộc tính như: \codeinline{width}, \codeinline{height}, \codeinline{margin}, \codeinline{padding}, \codeinline{font-size}, \codeinline{border-width}, \codeinline{border-radius}, \codeinline{box-shadow}, \codeinline{text-shadow}, \codeinline{min-width}, \codeinline{max-width}, \codeinline{min-height}, \codeinline{max-height}...

\paragraph{1. Pixel (px) Trong Ngữ Cảnh Màn Hình Là Gì?}

\begin{conceptBox}[Khái Niệm Pixel Và Pixel Tham Chiếu Trong CSS]
\begin{itemize}[leftmargin=1.5em, itemsep=3pt, parsep=0pt]
    \item \textbf{Pixel trên màn hình (Picture Element)}: Trong thế giới kỹ thuật số, một pixel (viết tắt của \term{Picture Element}) là điểm ảnh nhỏ nhất trên màn hình phần cứng có thể điều khiển riêng lẻ để hiển thị một màu sắc. Màn hình của bạn được tạo thành từ hàng triệu pixel xếp cạnh nhau.
    \item \textbf{Pixel tham chiếu trong CSS (\codeinline{px})}: Khi bạn sử dụng đơn vị \codeinline{px} trong CSS, bạn đang tham chiếu đến một \term{Pixel tham chiếu (Reference Pixel)}. Khác với pixel vật lý, pixel tham chiếu trong CSS là một đơn vị tuyệt đối -- giá trị của nó cố định và không thay đổi dựa trên kích thước phần tử cha hay kích thước trình duyệt (\term{Viewport}).
    \item \textbf{Tính cố định tuyệt đối}: Nếu bạn đặt \codeinline{width: 200px;} cho một phần tử, nó sẽ luôn có chiều rộng đúng $200$ pixel tham chiếu. Giá trị này không bị ảnh hưởng bất kể phần tử nằm trong một thẻ \codeinline{<div>} nhỏ hay một \codeinline{<div>} lớn hơn.
\end{itemize}
\end{conceptBox}

\paragraph{2. Sự Khác Biệt Giữa Pixel CSS Và Pixel Vật Lý (HiDPI / Retina Display):}
Đây là điểm quan trọng và thường gây nhầm lẫn nhất trong thiết kế giao diện web:
\begin{itemize}[leftmargin=1.5em, itemsep=3pt, parsep=0pt]
    \item \textbf{Pixel vật lý (Physical Pixel)}: Là chấm sáng nhỏ nhất trên màn hình phần cứng của thiết bị.
    \item \textbf{Pixel CSS (CSS Reference Pixel)}: Là đơn vị trừu tượng mà trình duyệt sử dụng để tính toán kích thước.
    \item \textbf{Cách trình duyệt xử lý trên các màn hình}:
    \begin{itemize}[leftmargin=1.2em, itemsep=2pt, parsep=0pt]
        \item \textbf{Màn hình tiêu chuẩn (Standard DPI)}: $1\text{px}$ CSS thường tương đương với $1$ pixel vật lý.
        \item \textbf{Màn hình mật độ cao (HiDPI / Retina)}: Để nội dung hiển thị sắc nét hơn và không bị ``rỗ'', trình duyệt sẽ tự động ``phóng to'' pixel CSS. Trên màn hình Retina, $1\text{px}$ CSS có thể được hiển thị bằng $2 \times 2$ (4) hoặc thậm chí $3 \times 3$ (9) pixel vật lý.
    \end{itemize}
    \item \textbf{Ví dụ thực tế}: Một đoạn văn bản có \codeinline{font-size: 16px;} sẽ có kích thước vật lý tương tự nhau trên cả màn hình tiêu chuẩn và màn hình Retina. Tuy nhiên, trên màn hình Retina, nó được vẽ bằng nhiều pixel vật lý hơn, khiến chữ hiển thị mịn màng và sắc nét hơn.
    \item \textbf{Mục đích cốt lõi}: Đảm bảo các phần tử trên trang web có kích thước ổn định và dễ đọc về mặt vật lý trên các thiết bị khác nhau, bất kể mật độ điểm ảnh.
\end{itemize}

\paragraph{3. Đơn Vị Tuyệt Đối Là Gì?}
Đơn vị \codeinline{px} là một đơn vị tuyệt đối, có nghĩa là giá trị của nó cố định và không thay đổi dựa trên kích thước của phần tử cha hoặc kích thước khung nhìn (\term{Viewport}). Nếu đặt \codeinline{width: 200px;} cho một thẻ \codeinline{<div>}, thẻ đó sẽ luôn có chiều rộng đúng $200$ pixel tham chiếu, bất kể kích thước màn hình hay kích thước phần tử chứa nó.

\paragraph{4. Ví Dụ Đơn Giản Khai Báo px Trong CSS:}

\begin{codeWindow}[Khai Báo Kích Thước Khối .box Bằng Đơn Vị px]
\begin{lstlisting}[language=HTML5]
<style>
.box {
  width: 200px;
  height: 100px;
  font-size: 16px;
  padding: 10px;
  border: 2px solid black;
}
</style>

<div class="box">Khoi Box dung px</div>
\end{lstlisting}
\end{codeWindow}

\begin{browserWindow}[Mô Phỏng Trực Quan Khối .box Đơn Vị px]
\sffamily\small
\centering
\begin{tcolorbox}[enhanced, colback=white, colframe=primaryBlue!30, width=200pt, height=100pt, arc=4pt, boxrule=1.5pt, boxsep=4pt, top=6pt, bottom=6pt, left=8pt, right=8pt]
  \textcolor{darkText}{\bfseries Khối Box ($200\text{px} \times 100\text{px}$)}\par
  \textcolor{codeComment}{\scriptsize (Font $16\text{px}$, Padding $10\text{px}$, Border $2\text{px}$ solid).}
\end{tcolorbox}
\end{browserWindow}

\paragraph{5. Đặc Điểm Của Đơn Vị px (Ưu Điểm \& Nhược Điểm):}

\begin{prettyTableBox}[Bảng Tổng Hợp Ưu Điểm Và Nhược Điểm Của Đơn Vị px]
\rowcolors{2}{white}{primaryBlue!4}
\sffamily\small
\begin{tabularx}{\linewidth}{|Z|Z|}
\hline
\rowcolor{primaryBlue} \thd{Ưu Điểm Của Pixel (px)} & \thd{Nhược Điểm \& Cân Nhắc} \\
\hline
Dễ kiểm soát kích thước chính xác tuyệt đối. & Không linh hoạt với màn hình quá lớn hoặc quá nhỏ. \\
\hline
Rất trực quan, dễ hình dung con số đại diện. & Khó tạo Responsive Web Design (cần viết nhiều Media Queries). \\
\hline
Phù hợp khi thiết kế chuẩn từng Pixel (\term{Pixel-Perfect}). & Không tự động co giãn theo thiết bị hoặc cài đặt Zoom phông chữ của người dùng. \\
\hline
Nhất quán trên mọi trình duyệt và hệ điều hành. & Gây cản trở khả năng truy cập (\term{Accessibility}) nếu dùng cho \codeinline{font-size} chính. \\
\hline
Tuyệt vời cho phần tử cố định (logo, icon, đường viền). & Có thể bị mờ hình ảnh raster nếu không dùng ảnh HiDPI/$2\times$. \\
\hline
\end{tabularx}
\end{prettyTableBox}

\paragraph{6. So Sánh px Với Các Đơn Vị Khác (em, rem, \%, vw, vh):}

\begin{prettyTableBox}[Bảng So Sánh Yếu Tố Phụ Thuộc Của Các Đơn Vị CSS]
\rowcolors{2}{white}{primaryBlue!4}
\sffamily\small
\begin{tabularx}{\linewidth}{|>{\raggedright\arraybackslash \ttfamily\color{primaryBlue}\bfseries}K{2.5cm}|Z|Z|}
\hline
\rowcolor{primaryBlue} \thd{Đơn Vị} & \thd{Ý Nghĩa} & \thd{Phụ Thuộc Vào Ai?} \\
\hline
px & Pixel tuyệt đối & Không phụ thuộc ai (Cố định tuyệt đối) \\
\hline
em & Tương đối theo phông chữ gần nhất & Dựa vào \codeinline{font-size} của phần tử cha \\
\hline
rem & Tương đối theo phông chữ gốc & Dựa vào \codeinline{font-size} của thẻ \codeinline{<html>} \\
\hline
\% & Tỷ lệ phần trăm & Dựa vào kích thước của phần tử cha \\
\hline
vw / vh & Phần trăm khung hình & Dựa vào kích thước màn hình (\term{Viewport}) \\
\hline
\end{tabularx}
\end{prettyTableBox}

\begin{prettyTableBox}[Bảng Ma Trận Sự Khác Biệt Tính Chất Giữa px Và Các Đơn Vị Tương Đối]
\rowcolors{2}{white}{primaryBlue!4}
\sffamily\small
\begin{tabularx}{\linewidth}{|>{\raggedright\arraybackslash \bfseries\color{primaryBlue}}K{2.6cm}|Z|Z|Z|Z|Z|}
\hline
\rowcolor{primaryBlue} \thd{Thuộc Tính} & \thd{px} & \thd{em} & \thd{rem} & \thd{\%} & \thd{vw / vh} \\
\hline
\textbf{Cố định} & \textcolor{green!60!black}{\textbf{[Có]}} & \textcolor{red!80!black}{\textbf{[Không]}} & \textcolor{red!80!black}{\textbf{[Không]}} & \textcolor{red!80!black}{\textbf{[Không]}} & \textcolor{red!80!black}{\textbf{[Không]}} \\
\hline
\textbf{Responsive} & \textcolor{red!80!black}{\textbf{[Không]}} & \textcolor{green!60!black}{\textbf{[Có]}} (tùy cha) & \textcolor{green!60!black}{\textbf{[Có]}} (toàn trang) & \textcolor{green!60!black}{\textbf{[Có]}} (tùy cha) & \textcolor{green!60!black}{\textbf{[Có]}} (theo Viewport) \\
\hline
\textbf{Dễ kiểm soát} & \textcolor{green!60!black}{\textbf{[Có]}} (nhỏ lẻ) & \textcolor{red!80!black}{\textbf{[Khó]}} (dễ bị lồng) & \textcolor{green!60!black}{\textbf{[Có]}} (tùy ngữ cảnh) & \textcolor{green!60!black}{\textbf{[Có]}} (layout lớn) & \textcolor{green!60!black}{\textbf{[Có]}} (layout tràn màn) \\
\hline
\end{tabularx}
\end{prettyTableBox}

\paragraph{7. Các Trường Hợp Sử Dụng Phổ Biến Của Pixel (px):}

\begin{prettyTableBox}[Bảng Tổng Hợp 6 Trường Hợp Sử Dụng Phổ Biến Của Đơn Vị Pixel]
\rowcolors{2}{white}{primaryBlue!4}
\sffamily\small
\begin{tabularx}{\linewidth}{|>{\raggedright\arraybackslash \bfseries\color{primaryBlue}}K{3.8cm}|>{\ttfamily\color{accentOrange}\bfseries}K{3.6cm}|Z|}
\hline
\rowcolor{primaryBlue} \thd{Trường Hợp Sử Dụng} & \thd{Thuộc Tính CSS} & \thd{Giải Thích \& Đánh Giá} \\
\hline
Độ dày đường viền & border-width & Thích hợp nhất cho viền cố định $1\text{px}$ -- $3\text{px}$, không bị biến dạng. \\
\hline
Khoảng cách lề cố định & padding, margin & Giữ khoảng cách nhỏ chuẩn xác (ví dụ $10\text{px}$, $15\text{px}$, $20\text{px}$). \\
\hline
Kích thước Icon / Logo & width, height & Đảm bảo biểu tượng nhỏ không bị lệch tỷ lệ hiển thị. \\
\hline
Giới hạn kích thước tối đa & max-width, min-width & Giới hạn chiều rộng khung nội dung (ví dụ $960\text{px}$, $1200\text{px}$). \\
\hline
Hiệu ứng viền \& bóng đổ & border-radius, box-shadow & Tạo độ bo góc và độ nhòe bóng đổ chuẩn xác từng pixel. \\
\hline
Cỡ chữ cố định (Font Size) & font-size & Kiểm soát cỡ chữ chuẩn tuyệt đối (nhược điểm: khó zoom riêng chữ). \\
\hline
\end{tabularx}
\end{prettyTableBox}

\begin{codeWindow}[Mã Nguồn Mẫu Tổng Hợp Các Trường Hợp Sử Dụng px Thực Tế]
\begin{lstlisting}[language=CSS3]
/* 1. Do day duong vien co dinh: */
.element { 
  border: 1px solid #ccc; 
}

/* 2. Khoang cach padding va margin nho: */
.button { 
  padding: 10px 20px; 
  margin-top: 15px; 
}

/* 3. Kich thuoc Icon va Logo nho: */
.icon { 
  width: 24px; 
  height: 24px; 
}

/* 4. Gioi han chieu rong khung noi dung: */
.content { 
  max-width: 960px; /* Noi dung khong bao gio rong hon 960px */ 
}

/* 5. Bo goc va tao bong do: */
.card {
  border-radius: 8px;
  box-shadow: 0 4px 12px rgba(0, 0, 0, 0.1);
}

/* 6. Co chu co dinh: */
h1 { 
  font-size: 32px; 
}
\end{lstlisting}
\end{codeWindow}

\begin{browserWindow}[Mô Phỏng Trực Quan Kết Quả Hiển Thị Các Trường Hợp Sử Dụng px]
\sffamily\small
\begin{tcolorbox}[enhanced, colback=white, colframe=primaryBlue!25, arc=6pt, boxrule=0.8pt, top=10pt, bottom=10pt, left=12pt, right=12pt]
  \begin{tabularx}{\linewidth}{Y Y}
    \tcbox[on line, arc=3pt, colback=gray!5, colframe=gray!40, boxsep=4pt]{%
      \small\textcolor{darkText}{\textbf{Viền cố định 1px solid \#ccc}}%
    } &
    \tcbox[on line, arc=4pt, colback=primaryBlue, colframe=primaryBlue, boxsep=4pt, left=10pt, right=10pt]{%
      \color{white}\bfseries\small Nút bấm (Padding 10px 20px)%
    } \\
    \vspace{2pt}\textcolor{primaryBlue}{\scriptsize (1. border: 1px solid \#ccc)} &
    \vspace{2pt}\textcolor{primaryBlue}{\scriptsize (2. padding: 10px 20px, margin-top: 15px)} \\[8pt]

    \tcbox[on line, arc=2pt, colback=orange!15, colframe=orange!80!black, boxsep=3pt]{%
      \color{orange!90!black}\bfseries\small Icon 24px $\times$ 24px%
    } &
    \tcbox[on line, arc=8pt, colback=white, colframe=gray!30, boxsep=6pt, drop shadow=black!10!white]{%
      \textcolor{primaryBlue}{\bfseries Thẻ Card (Radius 8px, Shadow)}%
    } \\
    \vspace{2pt}\textcolor{primaryBlue}{\scriptsize (3. width: 24px, height: 24px)} &
    \vspace{2pt}\textcolor{primaryBlue}{\scriptsize (5. border-radius: 8px, box-shadow)} \\
  \end{tabularx}
\end{tcolorbox}
\end{browserWindow}

\paragraph{8. Khi Nào Nên Dùng px Và Khi Nào Nên Dùng Đơn Vị Tương Đối?}
\begin{itemize}[leftmargin=1.5em, itemsep=3pt, parsep=0pt]
    \item \textbf{Nên sử dụng \codeinline{px} khi}:
    \begin{enumerate}[leftmargin=1.2em, itemsep=2pt, parsep=0pt]
        \item Cần độ chính xác cao tuyệt đối và kích thước cố định (đường viền, icon, khoảng cách nhỏ cụ thể).
        \item Đảm bảo một phần tử không thay đổi kích thước bất kể kích thước màn hình hay phần tử cha.
        \item Đối với các thuộc tính hiệu ứng kỹ thuật như \codeinline{border-radius}, \codeinline{box-shadow}, \codeinline{text-shadow}.
        \item Làm thiết kế tĩnh (\term{Fixed Layout}) hoặc UI nội bộ không yêu cầu responsive quá linh hoạt.
    \end{enumerate}
    \item \textbf{Nên sử dụng đơn vị tương đối (\%, \codeinline{em}, \codeinline{rem}, \codeinline{vw}, \codeinline{vh}) khi}:
    \begin{enumerate}[leftmargin=1.2em, itemsep=2pt, parsep=0pt]
        \item Xây dựng bố cục đáp ứng (\term{Responsive Layout}) tự động co giãn theo kích thước màn hình.
        \item Đảm bảo khả năng truy cập (\term{Accessibility}), đặc biệt là cho \codeinline{font-size} toàn trang.
        \item Duy trì tỷ lệ tương quan giữa các phần tử với nhau hoặc với khung hình Viewport.
    \end{enumerate}
\end{itemize}

\begin{noteBox}[Kết Luận Thực Tế Về Sử Dụng Đơn Vị Pixel]
Trong thực tế phát triển web hiện đại, một trang web chuyên nghiệp **luôn kết hợp cả đơn vị tuyệt đối (\codeinline{px}) và đơn vị tương đối (\codeinline{rem}, \codeinline{em}, \codeinline{\%})** để đạt được sự cân bằng hoàn hảo giữa kiểm soát chính xác và tính linh hoạt đáp ứng.
\end{noteBox}

\begin{warnBox}[Lưu Ý Về Đơn Vị Tuyệt Đối (Absolute Units)]
Các đơn vị tuyệt đối (\term{Absolute Units}) \textbf{không phản ứng với kích thước màn hình hoặc cài đặt phông chữ của trình duyệt}, do đó chúng \textbf{không có tính responsive}. Hạn chế dùng \codeinline{px} cho bố cục lớn hoặc kích thước phông chữ toàn trang.
\end{warnBox}

\begin{prettyTableBox}[Bảng Chi Tiết Các Đơn Vị Tương Đối (Relative Units)]
\rowcolors{2}{white}{primaryBlue!4}
\sffamily\small
\begin{tabularx}{\linewidth}{|>{\raggedright\arraybackslash \ttfamily\color{primaryBlue}\bfseries}K{2.2cm}|Z|>{\sffamily\small}Z|}
\hline
\rowcolor{primaryBlue} \thd{Đơn Vị} & \thd{Ý Nghĩa} & \thd{Dựa Vào Đâu / Ghi Chú} \\
\hline
\% & Tỷ lệ phần trăm & So với phần tử cha (thường là \codeinline{width}). \\
\hline
em & Relative to \codeinline{font-size} của chính nó & Kế thừa từ phần tử cha nếu không tự thiết lập. \\
\hline
rem & Relative to \codeinline{font-size} của \codeinline{<html>} & Dễ kiểm soát hơn \codeinline{em}, làm chuẩn cho toàn trang. \\
\hline
vw & Viewport width -- theo chiều ngang & $1\text{vw} = 1\%$ chiều rộng cửa sổ trình duyệt. \\
\hline
vh & Viewport height -- theo chiều cao & $1\text{vh} = 1\%$ chiều cao cửa sổ trình duyệt. \\
\hline
vmin & Nhỏ hơn giữa \codeinline{vw} và \codeinline{vh} & Phản ứng linh hoạt theo chiều nhỏ nhất của màn hình. \\
\hline
vmax & Lớn hơn giữa \codeinline{vw} và \codeinline{vh} & Hiếm dùng, theo chiều lớn nhất của màn hình. \\
\hline
ch & Chiều rộng ký tự số ``0'' của phông hiện tại & Thường dùng để giới hạn độ rộng cho vùng nhập văn bản (\codeinline{input text}). \\
\hline
ex & Chiều cao ký tự ``x'' của phông hiện tại & Ít dùng, không chính xác trên mọi phông chữ. \\
\hline
lh & Line height hiện tại & Tương đối với thuộc tính \codeinline{line-height} của phần tử. \\
\hline
\end{tabularx}
\end{prettyTableBox}

% ==============================================================================
% CHUYÊN ĐỀ SÂU VỀ ĐƠN VỊ PHẦN TRĂM (%) - PERCENTAGE MASTERCLASS TRONG CSS
% ==============================================================================
\paragraph{Chuyên Đề Chi Tiết Về Đơn Vị Phần Trăm (\%) -- Percentage Masterclass Trong CSS}
\label{sec:css_percentage_masterclass}

Đơn vị phần trăm (\codeinline{\%}) trong CSS là một đơn vị tương đối cực kỳ mạnh mẽ và linh hoạt, thường được sử dụng làm nòng cốt để tạo ra các bố cục trang web tự động co giãn và đáp ứng (\term{Responsive Web Design}).

\paragraph{1. Đơn Vị Tương Đối (\%) Là Gì?}

\begin{conceptBox}[Khái Niệm Đơn Vị Tương Đối Phần Trăm (\%)]
Trước hết, cần hiểu \codeinline{\%} là một \textbf{đơn vị tương đối (\term{Relative Unit})}. Điều này có nghĩa là giá trị thực tế của nó không cố định mà sẽ thay đổi tùy thuộc vào một \textbf{yếu tố tham chiếu} khác. Yếu tố tham chiếu này thường là kích thước của phần tử cha (\term{Parent Element}) hoặc một thuộc tính cụ thể của chính phần tử đó.
\par\vspace{4pt}
Đơn vị \codeinline{\%} được áp dụng rộng rãi cho nhiều thuộc tính CSS quan trọng như: \codeinline{width}, \codeinline{height}, \codeinline{margin}, \codeinline{padding}, \codeinline{font-size}, \codeinline{line-height}, \codeinline{transform}, \codeinline{background-size}, \codeinline{top}, \codeinline{left}, \codeinline{right}, \codeinline{bottom}...
\end{conceptBox}

\paragraph{2. Nguyên Tắc Hoạt Động Chung \& Công Thức Tính Đơn Vị \%:}

Khi bạn gán một giá trị \codeinline{\%} cho một thuộc tính CSS, trình duyệt sẽ tính toán giá trị độ dài thực tế (\term{Computed Value} tính bằng pixel) bằng cách lấy tỷ lệ phần trăm đó nhân với giá trị cơ sở (\term{Base Value}):

\begin{conceptBox}[Công Thức Chuyển Đổi Đơn Vị Phần Trăm Sang Pixel Thực Tế]
\[
\text{Giá Trị Thực Tế (px)} = \left( \frac{\text{Giá Trị Phần Trăm}}{100} \right) \times \text{Giá Trị Cơ Sở (px)}
\]
\begin{itemize}[leftmargin=1.5em, itemsep=2pt, parsep=0pt]
    \item \textbf{Giá trị phần trăm}: Là con số bạn chỉ định trong CSS (ví dụ: \codeinline{50\%}, \codeinline{20\%}, \codeinline{150\%}...).
    \item \textbf{Giá trị cơ sở}: Là giá trị tham chiếu mà phần trăm dựa vào (ví dụ: \codeinline{width} của phần tử cha, \codeinline{font-size} của phần tử cha...).
    \item \textbf{Giá trị thực tế}: Là số pixel thực tế được trình duyệt tính ra và áp dụng lên màn hình.
\end{itemize}
\end{conceptBox}

\begin{prettyTableBox}[Bảng Minh Họa Các Phép Tính Chuyển Đổi \% Cơ Bản Trong CSS]
\rowcolors{2}{white}{primaryBlue!4}
\sffamily\small
\begin{tabularx}{\linewidth}{|>{\ttfamily\color{primaryBlue}\bfseries}K{2.8cm}|Z|>{\raggedright\arraybackslash \ttfamily\color{codeNumber}\bfseries}K{4.2cm}|Z|}
\hline
\rowcolor{primaryBlue} \thd{Thuộc Tính} & \thd{Khai Báo CSS} & \thd{Công Thức Tính Toán} & \thd{Giá Trị Thực Tế} \\
\hline
width & Cha $500\text{px}$, con \codeinline{width: 60\%} & (60 / 100) * 500\text{px} & $300\text{px}$ \\
\hline
font-size & \codeinline{body} $20\text{px}$, \codeinline{p} \codeinline{font-size: 150\%} & (150 / 100) * 20\text{px} & $30\text{px}$ \\
\hline
padding & Khung $400\text{px}$, \codeinline{padding: 10\%} & (10 / 100) * 400\text{px} & $40\text{px}$ (cả 4 phía) \\
\hline
\end{tabularx}
\end{prettyTableBox}

\begin{prettyTableBox}[Bảng Tóm Tắt Nguyên Lý Hoạt Động Của Đơn Vị \%]
\rowcolors{2}{white}{primaryBlue!4}
\sffamily\small
\begin{tabularx}{\linewidth}{|>{\bfseries\color{primaryBlue}}K{3.2cm}|Z|}
\hline
\rowcolor{primaryBlue} \thd{Thành Phần} & \thd{Ý Nghĩa / Quy Tắc} \\
\hline
\textbf{Bản chất đơn vị} & Đơn vị tương đối (\term{Relative Unit}). \\
\hline
\textbf{Điểm tựa tính toán} & Phụ thuộc vào thuộc tính tương ứng của phần tử cha hoặc chính nó. \\
\hline
\textbf{Công thức cốt lõi} & $\text{Giá trị thực tế} = (\text{Phần trăm} / 100) \times \text{Giá trị cơ sở}$. \\
\hline
\end{tabularx}
\end{prettyTableBox}

\paragraph{3. Điểm Tham Chiếu Của Đơn Vị \% Theo Từng Thuộc Tính CSS:}

Đây là điểm quan trọng nhất cần nắm rõ khi sử dụng \codeinline{\%}, vì điểm tham chiếu thay đổi tùy thuộc vào thuộc tính mà bạn áp dụng nó:

\begin{prettyTableBox}[Bảng Tra Cứu Điểm Tham Chiếu Của Đơn Vị \% Theo Thuộc Tính CSS]
\rowcolors{2}{white}{primaryBlue!4}
\sffamily\small
\begin{tabularx}{\linewidth}{|>{\ttfamily\color{primaryBlue}\bfseries}K{3.2cm}|>{\bfseries\color{darkText}}K{3.8cm}|Z|}
\hline
\rowcolor{primaryBlue} \thd{Thuộc Tính CSS} & \thd{Phụ Thuộc Vào...} & \thd{Giải Thích Chi Tiết \& Điểm Tham Chiếu} \\
\hline
width \& height & \codeinline{width} hoặc \codeinline{height} của phần tử cha & Phần tử con có kích thước bằng tỷ lệ phần trăm của phần tử cha. \\
\hline
padding \& margin & \textbf{\codeinline{width} của phần tử cha} & \textbf{Lưu ý cốt lõi}: Cả 4 hướng (\codeinline{top}, \codeinline{bottom}, \codeinline{left}, \codeinline{right}) đều tính theo \codeinline{width} phần tử cha! \\
\hline
font-size & \codeinline{font-size} của phần tử cha & Kích thước chữ con tính theo tỷ lệ phần trăm cỡ chữ phần tử cha gần nhất. \\
\hline
line-height & \codeinline{font-size} của chính phần tử đó & Tính dựa trên kích thước phông chữ thực tế của phần tử hiện tại. \\
\hline
transform & Kích thước của chính phần tử đó & Tính dựa trên chiều rộng hoặc chiều cao của chính phần tử đang biến đổi. \\
\hline
background-size & Khung chứa ảnh nền phần tử & Tính dựa trên kích thước vùng bao chứa ảnh nền. \\
\hline
\end{tabularx}
\end{prettyTableBox}

\paragraph{a) Với width (chiều rộng) và height (chiều cao):}
\begin{itemize}[leftmargin=1.5em, itemsep=2pt, parsep=0pt]
    \item \textbf{Điểm tham chiếu}: Giá trị \codeinline{width} hoặc \codeinline{height} của phần tử cha (\term{Parent Element}).
    \item \textbf{Cách hoạt động}: Phần tử con sẽ có chiều rộng/chiều cao bằng một tỷ lệ phần trăm nhất định của phần tử cha.
\end{itemize}

\begin{codeWindow}[Ví Dụ Tính width và height Của Khối Con Theo \% Khối Cha]
\begin{lstlisting}[language=HTML5]
<div class="parent" style="width: 400px; height: 200px; background: #e0e0e0;">
    <div class="child" style="width: 50%; height: 75%; background: #00bcd4; color: white;">
        Child Block
    </div>
</div>
\end{lstlisting}
\end{codeWindow}

\begin{conceptBox}[Giải Thích Chi Tiết Phép Tính Kích Thước Khối Con]
\begin{itemize}[leftmargin=1.5em, itemsep=2pt, parsep=0pt]
    \item \textbf{Chiều rộng thực tế}: \codeinline{width: 50\%} $\rightarrow (50 / 100) \times 400\text{px} = \mathbf{200\text{px}}$.
    \item \textbf{Chiều cao thực tế}: \codeinline{height: 75\%} $\rightarrow (75 / 100) \times 200\text{px} = \mathbf{150\text{px}}$.
\end{itemize}
\end{conceptBox}

\begin{browserWindow}[Mô Phỏng Trực Quan Kết Quả Hiển Thị Của Khối .child (\%)]
\sffamily\small
\centering
\begin{tcolorbox}[enhanced, colback=gray!20, colframe=gray!50, width=240pt, height=120pt, arc=6pt, boxrule=1pt, top=8pt, bottom=8pt, left=8pt, right=8pt]
  \textcolor{darkText}{\textbf{Khối Cha (.parent): $400\text{px} \times 200\text{px}$}}\par\vspace{4pt}
  \begin{tcolorbox}[colback=cyan!80!black, colframe=cyan!90!black, width=120pt, height=75pt, arc=4pt, boxrule=1pt, top=4pt, bottom=4pt, left=6pt, right=6pt]
    \color{white}\bfseries\small .child ($200\text{px} \times 150\text{px}$)\par
    \tiny (Width 50\% cha = 200px, Height 75\% cha = 150px)
  \end{tcolorbox}
\end{tcolorbox}
\end{browserWindow}

\paragraph{b) Với padding (khoảng đệm) và margin (khoảng lề):}
\begin{itemize}[leftmargin=1.5em, itemsep=2pt, parsep=0pt]
    \item \textbf{Điểm tham chiếu}: Giá trị \textbf{\codeinline{width} (chiều rộng) của phần tử cha}.
    \item \textbf{Cách hoạt động (Điểm thường gây nhầm lẫn nhất)}:
    \begin{itemize}[leftmargin=1.2em, itemsep=2pt, parsep=0pt]
        \item Cả \codeinline{padding-top}, \codeinline{padding-bottom}, \codeinline{margin-top}, \codeinline{margin-bottom} (\textbf{chiều dọc}) \textbf{đều được tính dựa trên chiều rộng của phần tử cha}, chứ KHÔNG phải chiều cao.
        \item \codeinline{padding-left}, \codeinline{padding-right}, \codeinline{margin-left}, \codeinline{margin-right} (\textbf{chiều ngang}) cũng dựa trên chiều rộng của phần tử cha.
    \end{itemize}
    \item \textbf{Lý do thiết kế}: Nhằm duy trì tỷ lệ khung hình (\term{Aspect Ratio}) và tính linh hoạt trong bố cục, đặc biệt là khi tạo các khối nội dung co giãn có tỷ lệ cố định.
\end{itemize}

\begin{codeWindow}[Mã Nguồn Minh Họa padding và margin Theo Chiều Rộng Phần Tử Cha]
\begin{lstlisting}[language=HTML5]
<div class="parent" style="width: 500px;">
    <div class="child" style="padding-top: 10%; margin-left: 5%; background: #fff3e0; border: 1px solid #ff9800;">
        Noi dung khoi con
    </div>
</div>
\end{lstlisting}
\end{codeWindow}

\begin{conceptBox}[Giải Thích Phép Tính padding \& margin Theo \%]
Trong ví dụ trên với phần tử cha có \codeinline{width: 500px}:
\begin{itemize}[leftmargin=1.5em, itemsep=2pt, parsep=0pt]
    \item \textbf{\codeinline{padding-top: 10\%}}: Khoảng đệm trên thực tế là $(10 / 100) \times 500\text{px} = \mathbf{50\text{px}}$.
    \item \textbf{\codeinline{margin-left: 5\%}}: Khoảng lề trái thực tế là $(5 / 100) \times 500\text{px} = \mathbf{25\text{px}}$.
\end{itemize}
\end{conceptBox}

\paragraph{c) Với font-size (kích thước chữ):}
\begin{itemize}[leftmargin=1.5em, itemsep=2pt, parsep=0pt]
    \item \textbf{Điểm tham chiếu}: Giá trị \codeinline{font-size} của phần tử cha gần nhất.
    \item \textbf{Cách hoạt động}: Kích thước chữ của phần tử con sẽ là một tỷ lệ phần trăm của kích thước chữ của phần tử cha.
\end{itemize}

\begin{codeWindow}[Ví Dụ Kế Thừa font-size Đa Cấp Theo Tỷ Lệ \%]
\begin{lstlisting}[language=HTML5]
<body style="font-size: 16px;">
    <div style="font-size: 125%;">
        <p style="font-size: 80%;">Day la doan van ban da cap.</p>
    </div>
</body>
\end{lstlisting}
\end{codeWindow}

\begin{conceptBox}[Phân Tích Kích Thước Chữ Tính Toán Từng Cấp]
\begin{itemize}[leftmargin=1.5em, itemsep=2pt, parsep=0pt]
    \item Thẻ \codeinline{<body>} có cỡ chữ gốc: $16\text{px}$.
    \item Thẻ \codeinline{<div>} có \codeinline{font-size: 125\%} $\rightarrow (125 / 100) \times 16\text{px} = \mathbf{20\text{px}}$.
    \item Thẻ \codeinline{<p>} có \codeinline{font-size: 80\%} (dựa trên thẻ \codeinline{<div>} cha có $20\text{px}$) $\rightarrow (80 / 100) \times 20\text{px} = \mathbf{16\text{px}}$.
\end{itemize}
\end{conceptBox}

\paragraph{d) Với line-height (chiều cao dòng):}
\begin{itemize}[leftmargin=1.5em, itemsep=2pt, parsep=0pt]
    \item \textbf{Điểm tham chiếu}: Giá trị \codeinline{font-size} của chính phần tử đó.
    \item \textbf{Cách hoạt động}: Chiều cao dòng bằng tỷ lệ phần trăm của kích thước chữ của chính nó.
\end{itemize}

\begin{codeWindow}[Khai Báo line-height Theo \%]
\begin{lstlisting}[language=CSS3]
p {
    font-size: 18px;
    line-height: 150%; /* 150% cua 18px = (150 / 100) * 18px = 27px */
}
\end{lstlisting}
\end{codeWindow}

\paragraph{e) Với transform (biến đổi hình học):}
\begin{itemize}[leftmargin=1.5em, itemsep=2pt, parsep=0pt]
    \item \textbf{Điểm tham chiếu}: Kích thước của chính phần tử đang được biến đổi.
    \item \textbf{Cách hoạt động}: Giá trị dịch chuyển (\codeinline{translateX()}, \codeinline{translateY()}), phóng to... sẽ dựa trên kích thước của bản thân phần tử đó.
\end{itemize}

\begin{codeWindow}[Khai Báo transform: translateX(50\%) Dựa Trên Kích Thước Bản Thân]
\begin{lstlisting}[language=CSS3]
.box {
    width: 100px;
    height: 100px;
    transform: translateX(50%); /* Dich sang phai 50% cua width (100px) = 50px */
}
\end{lstlisting}
\end{codeWindow}

\paragraph{4. Ưu Điểm Nổi Bật Của Đơn Vị Phần Trăm (\%):}

\begin{prettyTableBox}[Bảng Tổng Hợp Ưu Điểm Của Đơn Vị \% Trong Thiết Kế Web]
\rowcolors{2}{white}{primaryBlue!4}
\sffamily\small
\begin{tabularx}{\linewidth}{|>{\bfseries\color{primaryBlue}}K{3.8cm}|Z|}
\hline
\rowcolor{primaryBlue} \thd{Ưu Điểm Nổi Bật} & \thd{Phân Tích Chi Tiết \& Lợi Ích Thực Tế} \\
\hline
\textbf{Responsive Design (Thiết kế đáp ứng)} & Ưu điểm lớn nhất. Khi kích thước trình duyệt hoặc phần tử cha thay đổi, phần tử tự động co giãn theo, thích ứng mượt mà trên desktop, tablet và mobile. \\
\hline
\textbf{Fluid Layouts (Bố cục linh hoạt)} & Tạo ra các bố cục dạng ``dòng chảy'' mềm mại, tự lấp đầy không gian có sẵn, tận dụng tối đa không gian màn hình. \\
\hline
\textbf{Giảm thiểu Media Queries} & Giảm đáng kể số lượng câu lệnh \codeinline{@media} phải viết, giúp mã nguồn CSS gọn gàng và dễ bảo trì hơn. \\
\hline
\end{tabularx}
\end{prettyTableBox}

\paragraph{5. Nhược Điểm Và Những Điều Cần Lưu Ý Khi Sử Dụng \%:}

\begin{prettyTableBox}[Bảng Nhược Điểm Và Cân Nhắc Kỹ Thuật Khi Dùng Đơn Vị \%]
\rowcolors{2}{white}{primaryBlue!4}
\sffamily\small
\begin{tabularx}{\linewidth}{|>{\bfseries\color{accentOrange}}K{3.8cm}|Z|}
\hline
\rowcolor{primaryBlue} \thd{Nhược Điểm / Cân Nhắc} & \thd{Phân Tích Chi Tiết \& Rủi Ro Cần Tránh} \\
\hline
\textbf{Tính kế thừa phức tạp (Cascade)} & Việc lồng ghép nhiều phần tử con liên tiếp dùng \codeinline{\%} dễ gây hiện tượng ``phình to'' hoặc ``teo nhỏ'' không mong muốn nếu không tính toán kỹ. \\
\hline
\textbf{Khó kiểm soát chính xác} & Khi cần kích thước cố định tuyệt đối (như đường viền, icon nhỏ), \codeinline{\%} không phải là lựa chọn tốt (nên dùng \codeinline{px}). \\
\hline
\textbf{Sự phụ thuộc vào phần tử cha} & Phần tử cha phải có kích thước xác định. Nếu cha không set \codeinline{height} rõ ràng, \codeinline{height: 100\%} ở con sẽ \textbf{hoàn toàn không có tác dụng}. \\
\hline
\textbf{Vấn đề với typography} & \codeinline{rem} và \codeinline{em} thường được ưu tiên hơn \codeinline{\%} cho kích thước chữ vì mang lại khả năng quản lý tỷ lệ phông chữ nhất quán hơn. \\
\hline
\end{tabularx}
\end{prettyTableBox}

\paragraph{6. Khi Nào Nên Sử Dụng Đơn Vị \% Trong Thực Tế?}

\begin{prettyTableBox}[Bảng Tổng Hợp Các Trường Hợp Ứng Dụng Thực Tế Của Đơn Vị \%]
\rowcolors{2}{white}{primaryBlue!4}
\sffamily\small
\begin{tabularx}{\linewidth}{|>{\raggedright\arraybackslash \bfseries\color{primaryBlue}}K{3.8cm}|>{\ttfamily\color{accentOrange}\bfseries}K{3.4cm}|Z|}
\hline
\rowcolor{primaryBlue} \thd{Trường Hợp Thực Tế} & \thd{Thuộc Tính CSS} & \thd{Mục Đích Áp Dụng} \\
\hline
Chia cột trong hệ thống lưới & width: 33.33\%, width: 49\% & Chia $3$ cột đều nhau ($33.33\%$) hoặc chia $2$ cột ($49\%$) kèm khoảng cách lề. \\
\hline
Khung bố cục chính (Container) & width: 80\%, margin: 0 auto & Tạo khung trung tâm chiếm $80\%$ độ rộng màn hình và tự động căn giữa. \\
\hline
Hình ảnh đáp ứng (Responsive Img) & max-width: 100\%, height: auto & Đảm bảo ảnh không bao giờ tràn ra ngoài khung chứa và giữ đúng tỷ lệ khung hình. \\
\hline
Duy trì tỷ lệ khoảng cách & padding: 5\% & Khoảng đệm tự co giãn tỷ lệ thuận với kích thước tổng thể của giao diện. \\
\hline
\end{tabularx}
\end{prettyTableBox}

\begin{codeWindow}[Mã Nguồn Mẫu Tổng Hợp Các Tình Huống Dùng Đơn Vị \% Chuẩn Mực]
\begin{lstlisting}[language=HTML5]
<style>
/* 1. Thiet lap khung container trung tam: */
.container {
    width: 80%;       /* Chiem 80% chieu rong cua phan tu cha */
    margin: 0 auto;   /* Can giua khung */
}

/* 2. Chia cot co gian trong Grid/Layout: */
.column-3 {
    width: 33.33%;    /* 3 cot deu nhau */
    float: left;
}
.column-2 {
    width: 49%;       /* 2 cot co khoang cach */
    float: left;
    margin-right: 2%;
}

/* 3. Padding duy tri ty le: */
.card {
    padding: 5%;      /* Padding 5% dua tren width phan tu cha */
}

/* 4. Hinh anh Responsive linh hoat: */
img.responsive-img {
    max-width: 100%;  /* Anh khong bao gio vuot qua chieu rong khung cha */
    height: auto;     /* Tu dong giu dung ty le khung hinh */
}
</style>
\end{lstlisting}
\end{codeWindow}

\begin{browserWindow}[Mô Phỏng Trực Quan Kết Quả Layout Chia Cột Và Ảnh Responsive (\%)]
\sffamily\small
\begin{tcolorbox}[enhanced, colback=white, colframe=primaryBlue!25, arc=6pt, boxrule=0.8pt, top=10pt, bottom=10pt, left=12pt, right=12pt]
  \textcolor{primaryBlue}{\bfseries Khung Container (Width 80\% màn hình, căn giữa margin: 0 auto)}
  \par\vspace{6pt}
  \begin{tcolorbox}[colback=blue!8, colframe=blue!60!black, arc=4pt, boxrule=1pt, top=6pt, bottom=6pt, left=8pt, right=8pt]
    \textcolor{primaryBlue}{\bfseries Cột 1 (Width 49\%)}\par
    \textcolor{darkText}{\small Thẻ Card Padding 5\% co giãn theo độ rộng cha.}
  \end{tcolorbox}
  \vspace{4pt}
  \begin{tcolorbox}[colback=blue!8, colframe=blue!60!black, arc=4pt, boxrule=1pt, top=6pt, bottom=6pt, left=8pt, right=8pt]
    \textcolor{primaryBlue}{\bfseries Cột 2 (Width 49\%)}\par
    \textcolor{darkText}{\small Ảnh \codeinline{max-width: 100\%} không tràn khung.}
  \end{tcolorbox}
\end{tcolorbox}
\end{browserWindow}

\begin{noteBox}[Tóm Tắt Cốt Lõi Về Đơn Vị Phần Trăm (\%)]
Tóm lại, đơn vị \codeinline{\%} là một công cụ không thể thiếu trong CSS hiện đại, giúp bạn tạo ra các thiết kế linh hoạt và đáp ứng. Hiểu rõ \textbf{điểm tham chiếu của nó cho từng thuộc tính CSS} là chìa khóa để sử dụng nó một cách hiệu quả và tránh những lỗi bố cục không mong muốn.
\end{noteBox}

\paragraph{II. Khi Nào Dùng Đơn Vị Nào? (Best Practices Guidance)}

\begin{prettyTableBox}[Bảng Hướng Dẫn Lựa Lựa Chọn Đơn Vị CSS Theo Tình Huống]
\rowcolors{2}{white}{primaryBlue!4}
\sffamily\small
\begin{tabularx}{\linewidth}{|>{\raggedright\arraybackslash \bfseries\color{primaryBlue}}K{3.5cm}|>{\ttfamily\color{accentOrange}\bfseries}K{3.0cm}|Z|}
\hline
\rowcolor{primaryBlue} \thd{Tình Huống Sử Dụng} & \thd{Nên Dùng Đơn Vị Gì?} & \thd{Vì Sao? (Lý Do Kỹ Thuật)} \\
\hline
Thiết kế Responsive Layout & \%, vw, vh, rem & Thích nghi tốt với mọi kích thước màn hình. \\
\hline
Phông chữ dễ kiểm soát & rem & Dễ điều chỉnh tập trung từ thẻ gốc (\codeinline{<html>}). \\
\hline
Khoảng cách Padding / Margin & em, rem, \% & Tùy thuộc vào tính kế thừa và ngữ cảnh thành phần. \\
\hline
Kích thước ảnh, khung cố định & px & Chính xác tuyệt đối, dễ canh chỉnh các chi tiết nhỏ. \\
\hline
\end{tabularx}
\end{prettyTableBox}

\paragraph{III. So Sánh Nhanh: em vs rem (Comparative Matrix)}

\begin{prettyTableBox}[Bảng So Sánh Chi Tiết Giữa Hai Đơn Vị em Và rem]
\rowcolors{2}{white}{primaryBlue!4}
\sffamily\small
\begin{tabularx}{\linewidth}{|>{\raggedright\arraybackslash \bfseries\color{primaryBlue}}K{3.2cm}|Z|Z|}
\hline
\rowcolor{primaryBlue} \thd{Thuộc Tính So Sánh} & \thd{Đơn Vị em} & \thd{Đơn Vị rem} \\
\hline
\textbf{Dựa vào đâu?} & \codeinline{font-size} của phần tử cha gần nhất & \codeinline{font-size} của thẻ gốc \codeinline{<html>} \\
\hline
\textbf{Dễ kiểm soát?} & \textcolor{red!80!black}{\textbf{[Khó]}} Có thể bị lồng nhau quá mức (phóng đại kích thước) & \textcolor{green!60!black}{\textbf{[Dễ]}} Dễ kiểm soát, đồng nhất toàn bộ trang web \\
\hline
\textbf{Dùng khi nào?} & Muốn kích thước co giãn theo bối cảnh phần tử gần nhất & Muốn toàn trang đồng bộ và dễ điều chỉnh tỷ lệ \\
\hline
\end{tabularx}
\end{prettyTableBox}

\paragraph{IV. Ví Dụ Minh Họa Khai Báo Các Đơn Vị CSS}

\begin{codeWindow}[Ví Dụ Sử Dụng rem, em, \%, vw và vh Trong CSS]
\begin{lstlisting}[language=HTML5]
<style>
/* 1. Dat font-size goc cho toan trang: 1rem = 16px */
html {
  font-size: 16px; 
}

/* 2. Su dung phan tram, em va rem: */
.box1 {
  width: 50%;        /* 50% chieu rong cua phan tu cha */
  padding: 2em;      /* 2 * font-size cua chinh no (32px neu font 16px) */
  font-size: 1.2rem; /* 1.2 * font-size html = 19.2px */
  background: #e3f2fd;
  border: 2px solid #2196f3;
}

/* 3. Su dung don vi Viewport (vw, vh): */
.box2 {
  width: 80vw;       /* 80% chieu ngang khung hinh trinh duyet */
  height: 30vh;      /* 30% chieu cao khung hinh trinh duyet */
  background: #f1f8e9;
  border: 2px solid #8bc34a;
}
</style>

<div class="box1">Box 1: Dung phan tram, em va rem</div>
<div class="box2">Box 2: Dung 80vw x 30vh Viewport</div>
\end{lstlisting}
\end{codeWindow}

\begin{browserWindow}[Mô Phỏng Trực Quan Kết Quả Hiển Thị Của Đơn Vị CSS (.box1 \& .box2)]
\sffamily\small
\begin{tcolorbox}[enhanced, colback=white, colframe=primaryBlue!25, arc=6pt, boxrule=0.8pt, top=10pt, bottom=10pt, left=12pt, right=12pt]
  % Box 1 simulation
  \begin{tcolorbox}[colback=blue!8, colframe=blue!60!black, arc=4pt, boxrule=1pt, width=0.6\linewidth, top=8pt, bottom=8pt, left=10pt, right=10pt]
    \textcolor{primaryBlue}{\bfseries\large Box 1: Width 50\%, Padding 2em, Font 1.2rem}\par
    \textcolor{darkText}{\small (Kích thước chữ $19.2\text{px}$, Padding $38.4\text{px}$ co giãn theo ngữ cảnh thẻ gốc \& thẻ cha).}
  \end{tcolorbox}
  \vspace{8pt}
  % Box 2 simulation
  \begin{tcolorbox}[colback=green!8, colframe=green!60!black, arc=4pt, boxrule=1pt, width=0.9\linewidth, top=8pt, bottom=8pt, left=10pt, right=10pt]
    \textcolor{green!40!black}{\bfseries\large Box 2: Width 80vw, Height 30vh Viewport}\par
    \textcolor{darkText}{\small (Kích thước luôn chiếm đúng $80\%$ chiều rộng màn hình trình duyệt hiện tại).}
  \end{tcolorbox}
\end{tcolorbox}
\end{browserWindow}

\paragraph{V. Mẹo \& Lưu Ý Quan Trọng (Tips \& Best Practices)}

\begin{tipBox}[Mẹo Thực Chiến \& Lưu Ý Cốt Lõi Khi Dùng Đơn Vị CSS]
\begin{itemize}[leftmargin=1.5em, itemsep=4pt, parsep=0pt]
    \item $\rightarrow$ \textbf{\codeinline{rem} nên dùng cho kích thước chữ (\codeinline{font-size})}: Giúp thay đổi quy mô chữ toàn bộ trang dễ dàng chỉ qua 1 dòng CSS thay đổi trên thẻ gốc \codeinline{html}.
    \item $\rightarrow$ \textbf{\codeinline{px} hạn chế dùng trong layout}: \codeinline{px} rất chính xác nhưng không responsive, nên hạn chế dùng cho độ rộng khung hoặc phông chữ chính.
    \item $\rightarrow$ \textbf{\codeinline{em} chú ý hiệu ứng nhân dồn (\term{Em Compounding Effect})}: \codeinline{em} có thể bị lồng nhau gây phóng đại kích thước chữ ngoài ý muốn -- cần sử dụng cẩn thận.
    \item $\rightarrow$ \textbf{\codeinline{vw} / \codeinline{vh} tối ưu cho Fullscreen Layout}: Phù hợp nhất khi thiết kế các khối Hero Section, Banner tràn màn hình hoặc giao diện App $100\%$ chiều cao khung hình.
\end{itemize}
\end{tipBox}

\subsection{Mục "Đọc Thêm": Mở Rộng Kiến Thức Về Các Đơn Vị Trong CSS}

\begin{conceptBox}[Đọc Thêm: Các Kỹ Thuật Đơn Vị Hiện Đại Trong CSS Modern]
Mục này giới thiệu các đơn vị Viewport thế hệ mới, Container Query Units và kỹ thuật Fluid Typography giúp tối ưu Responsive Web Design chuẩn hiện đại.
\end{conceptBox}

\paragraph{1. Các Đơn Vị Viewport Thế Hệ Mới (\codeinline{dvh}, \codeinline{svh}, \codeinline{lvh}):}
Trên các trình duyệt di động (như Safari trên iOS hay Chrome trên Android), thanh địa chỉ (\term{Address Bar}) có thể tự động ẩn hoặc hiện khi người dùng cuộn trang. Việc này làm giá trị \codeinline{100vh} truyền thống bị co giật giao diện. CSS Viewport Level 4 đã bổ sung các đơn vị giải quyết triệt để vấn đề này:
\begin{itemize}[leftmargin=1.5em, itemsep=3pt, parsep=0pt]
    \item \textbf{\codeinline{svh} (Small Viewport Height)}: Chiều cao Viewport khi thanh địa chỉ \textbf{hiển thị đầy đủ} (kích thước nhỏ nhất).
    \item \textbf{\codeinline{lvh} (Large Viewport Height)}: Chiều cao Viewport khi thanh địa chỉ \textbf{thu gọn hoàn toàn} (kích thước lớn nhất).
    \item \textbf{\codeinline{dvh} (Dynamic Viewport Height)}: Chiều cao Viewport \textbf{tự động cập nhật linh hoạt} theo trạng thái ẩn/hiện của thanh địa chỉ trình duyệt thực tế.
\end{itemize}

\paragraph{2. Đơn Vị Container Query Units (\codeinline{cqw}, \codeinline{cqh}, \codeinline{cqi}):}
Khi sử dụng tính năng Modern CSS Container Queries, thay vì phụ thuộc vào kích thước toàn màn hình (Viewport), bạn có thể định cỡ các thành phần con dựa trên kích thước của chính khung chứa (\term{Container}) bao ngoài nó:
\begin{itemize}[leftmargin=1.5em, itemsep=3pt, parsep=0pt]
    \item \textbf{\codeinline{cqw} / \codeinline{cqh}}: $1\%$ chiều rộng / chiều cao của Container bao ngoài.
    \item \textbf{\codeinline{cqi} / \codeinline{cqb}}: $1\%$ kích thước theo hướng dòng (\term{Inline}) / hướng khối (\term{Block}) của Container.
\end{itemize}

\paragraph{3. Kỹ Thuật Fluid Typography Không Cần Media Query Với Hàm \codeinline{clamp()}:}
Kết hợp đơn vị \codeinline{rem} và \codeinline{vw} với hàm \codeinline{clamp(min, val, max)} tạo ra cỡ chữ co giãn mượt mà theo độ rộng màn hình mà không cần viết hàng loạt đoạn \codeinline{@media}:

\begin{codeWindow}[Kỹ Thuật Fluid Typography Với clamp()]
\begin{lstlisting}[language=CSS3]
/* Co chu co gian tu 16px (1rem) den 32px (2rem), linh hoat theo 2.5vw + 0.5rem */
h1 {
    font-size: clamp(1rem, 2.5vw + 0.5rem, 2rem);
}
\end{lstlisting}
\end{codeWindow}

\begin{browserWindow}[Mô Phỏng Trực Quan Kích Thước Chữ Fluid Typography Theo Độ Rộng Màn Hình]
\sffamily\small
\begin{tcolorbox}[enhanced, colback=white, colframe=primaryBlue!25, arc=6pt, boxrule=0.8pt, top=10pt, bottom=10pt, left=12pt, right=12pt]
  \begin{tabularx}{\linewidth}{Y Y}
    \begin{tcolorbox}[colback=orange!8, colframe=orange!60!black, arc=4pt, boxrule=1pt, top=6pt, bottom=6pt, left=8pt, right=8pt]
      \textcolor{orange!90!black}{\bfseries\small Màn Hình Mobile ($375\text{px}$)}\par\vspace{4pt}
      \textcolor{darkText}{\textbf{\small Tiêu đề H1 ($16\text{px} / 1\text{rem}$)}}
    \end{tcolorbox} &
    \begin{tcolorbox}[colback=teal!8, colframe=teal!60!black, arc=4pt, boxrule=1pt, top=6pt, bottom=6pt, left=8pt, right=8pt]
      \textcolor{teal!90!black}{\bfseries\small Màn Hình Desktop ($1440\text{px}$)}\par\vspace{4pt}
      \textcolor{darkText}{\bfseries\Large Tiêu đề H1 ($32\text{px} / 2\text{rem}$)}
    \end{tcolorbox} \\
    \vspace{2pt}\textcolor{codeComment}{\scriptsize (1rem min-bound: Cỡ chữ nhỏ nhất)} &
    \vspace{2pt}\textcolor{codeComment}{\scriptsize (2rem max-bound: Cỡ chữ lớn nhất)} \\
  \end{tabularx}
\end{tcolorbox}
\end{browserWindow}

\paragraph{4. Tip Phỏng Vấn Senior Frontend: Kỹ Thuật Reset 62.5\% Cho Thẻ Root \codeinline{html}:}
\begin{tipBox}[Câu Hỏi Phỏng Vấn: Tại Sao Nhiều Dự Án Đặt html \{ font-size: 62.5\%; \}?]
\textbf{Câu hỏi:} Tại sao trong một số thư viện CSS Reset người ta thường viết \codeinline{html \{ font-size: 62.5\%; \}}?
\par\vspace{4pt}
\textbf{Trả lời chuẩn phỏng vấn:} Trình duyệt mặc định có phông chữ gốc là $16\text{px}$. Nếu đặt \codeinline{font-size: 62.5\%} trên thẻ \codeinline{<html>}, kích thước gốc sẽ trở thành $16\text{px} \times 62.5\% = 10\text{px}$. Việc này giúp lập trình viên cực kỳ dễ nhẩm tính chuyển đổi từ \codeinline{px} sang \codeinline{rem}: $1.4\text{rem} = 14\text{px}$, $1.6\text{rem} = 16\text{px}$, $2.4\text{rem} = 24\text{px}$, trong khi vẫn bảo toàn tính năng hỗ trợ truy cập Zoom phông chữ của người dùng.
\end{tipBox}

% ==============================================================================
% MỤC 4.1: HỆ THỐNG MÀU SẮC TRONG CSS (CSS COLORS MASTERCLASS TỪ A-Z)
% ==============================================================================
\subsection{Chuyên Đề Hệ Thống Màu Sắc Trong CSS (CSS Colors Masterclass Từ A đến Z)}
\label{sec:css_colors_masterclass}

Màu sắc là yếu tố sống còn quyết định tính thẩm mỹ, khả năng truyền tải thông điệp, trải nghiệm người dùng (\term{UX}) và khả năng truy cập (\term{Accessibility - a11y}) của mọi trang web. Trong CSS, lập trình viên có thể biểu diễn màu sắc qua nhiều hệ thống biểu diễn khác nhau từ đơn giản đến nâng cao.

\paragraph{1. Mục Đích Và Vai Trò Của Màu Sắc Trong CSS:}

\begin{prettyTableBox}[Bảng Phân Tích Vai Trò Của Màu Sắc Trên Giao Diện Web]
\rowcolors{2}{white}{primaryBlue!4}
\sffamily\small
\begin{tabularx}{\linewidth}{|>{\raggedright\arraybackslash \bfseries\color{primaryBlue}}p{3.8cm}|>{\raggedright\arraybackslash}X|>{\raggedright\arraybackslash}X|}
\hline
\rowcolor{primaryBlue} \thd{\sffamily Thành phần giao diện} & \thd{\sffamily Vai trò màu sắc đảm nhận} & \thd{\sffamily Mục tiêu trải nghiệm (UX)} \\
\hline
Văn bản (Text) & Định kiểu chữ chính, chữ phụ, tiêu đề. & Tăng độ đọc (\term{Readability}), phân cấp thông tin rõ ràng. \\
\hline
Nút bấm \& Liên kết & Nhấn mạnh các điểm tương tác chính (CTA). & Hướng người dùng thực hiện hành động chính. \\
\hline
Nền (Background) & Phân tách khối nội dung, tạo khoảng không. & Tạo bố cục cân đối, dễ chịu cho mắt. \\
\hline
Cảnh báo \& Thông báo & Thể hiện trạng thái hệ thống. & Báo lỗi (Đỏ), Cảnh báo (Vàng), Thành công (Xanh lá). \\
\hline
\end{tabularx}
\end{prettyTableBox}

\paragraph{2. Phân Loại \& Tổng Quan Các Kiểu Màu Trong CSS:}

\begin{prettyTableBox}[Bảng Tổng Quan Chi Tiết Các Kiểu Biểu Diễn Màu CSS]
\rowcolors{2}{white}{primaryBlue!4}
\sffamily\small
\begin{tabularx}{\linewidth}{|>{\raggedright\arraybackslash \ttfamily\color{primaryBlue}\bfseries}p{2.6cm}|>{\ttfamily\small}p{3.2cm}|>{\ttfamily\small}p{3.2cm}|>{\raggedright\arraybackslash}X|}
\hline
\rowcolor{primaryBlue} \thd{\sffamily Loại màu} & \thd{\sffamily Cú pháp} & \thd{\sffamily Ví dụ minh họa} & \thd{\sffamily Đặc tính \& Ghi chú} \\
\hline
Color Name & tên\_màu & color: red; & Có sẵn 147 tên màu chuẩn HTML/CSS. \\
\hline
Hexadecimal & \#RRGGBB & color: \#ff0000; & Mã 6 ký tự Hex, chính xác tuyệt đối. \\
\hline
Shorthand Hex & \#RGB & color: \#f00; & Rút gọn 3 ký tự khi từng cặp trùng nhau. \\
\hline
RGB & rgb(R, G, B) & rgb(255, 0, 0) & Dải giá trị $0 \rightarrow 255$ cho 3 kênh Red, Green, Blue. \\
\hline
RGBA & rgba(R, G, B, A) & rgba(255, 0, 0, 0.5) & Thêm kênh Alpha độ trong suốt ($0 \rightarrow 1$). \\
\hline
HSL & hsl(H, S\%, L\%) & hsl(0, 100\%, 50\%) & Hue (góc màu), Saturation (độ bão hòa), Lightness (độ sáng). \\
\hline
HSLA & hsla(H, S\%, L\%, A) & hsla(0, 100\%, 50\%, 0.5) & Thêm kênh Alpha độ trong suốt ($0 \rightarrow 1$). \\
\hline
Transparent & transparent & background: transparent; & Từ khóa đặc biệt tạo độ trong suốt tuyệt đối ($100\%$). \\
\hline
CurrentColor & currentColor & border-color: currentColor; & Từ khóa đặc biệt tự động kế thừa màu chữ (\codeinline{color}) hiện tại. \\
\hline
\end{tabularx}
\end{prettyTableBox}

\paragraph{3. Chi Tiết Từng Hệ Thống Biểu Diễn Màu Sắc:}

\paragraph{3.1. Color Name (Tên Màu Định Danh Trong CSS):}

\begin{conceptBox}[Khái Niệm Cơ Bản Về Tên Màu Định Danh (Color Name)]
\textbf{1. Color Name Là Gì?}:
\begin{itemize}[leftmargin=1.5em, itemsep=3pt, parsep=0pt]
    \item \textbf{Color Name}: Là tên màu sắc đã được định nghĩa sẵn trong tiêu chuẩn CSS. Người lập trình chỉ cần gọi trực tiếp tên tiếng Anh của màu sắc đó mà không cần sử dụng mã Hex hay hàm toán học nào (không cần dấu \codeinline{\#}, không cần đối số $RGB$).
    \item \textbf{Số lượng màu chuẩn}: CSS hỗ trợ tổng cộng \textbf{147 tên màu chuẩn} được quy định chính thức bởi W3C.
\end{itemize}

\textbf{2. Cú Pháp Sử Dụng Color Name trong CSS}:
\begin{codeWindow}[Cú Pháp Định Kiểu Bằng Tên Màu Định Danh]
\begin{lstlisting}[language=CSS3]
/* Cu phap tong quat */
element {
  color: <color-name>;
  background-color: <color-name>;
}

/* Vi du thuc te */
p {
  color: tomato;
  background-color: lightgray;
}
\end{lstlisting}
\end{codeWindow}
\end{conceptBox}

\textbf{3. Phân Nhóm Các Tên Màu Phổ Biến Trong CSS:}

\begin{prettyTableBox}[Bảng Phân Nhóm Các Tên Màu Từ Khóa Thường Dùng]
\rowcolors{2}{white}{primaryBlue!4}
\sffamily\small
\begin{tabularx}{\linewidth}{|>{\raggedright\arraybackslash \bfseries\color{primaryBlue}}p{3.2cm}|>{\ttfamily\color{accentOrange}}p{4.5cm}|>{\raggedright\arraybackslash}X|}
\hline
\rowcolor{primaryBlue} \thd{\sffamily Nhóm màu} & \thd{\sffamily Ví dụ tên màu định sẵn} & \thd{\sffamily Đặc điểm \& Cảm giác thị giác} \\
\hline
Màu Đỏ (Red) & red, crimson, tomato, firebrick & Ấm áp, nổi bật, dùng cho cảnh báo hoặc điểm nhấn. \\
\hline
Màu Xanh lá (Green) & green, lime, olive, seagreen & Tươi mát, tự nhiên, trạng thái thành công. \\
\hline
Màu Xanh dương (Blue) & blue, navy, skyblue, dodgerblue & Tin cậy, chuyên nghiệp, liên kết và điều hướng. \\
\hline
Màu Trung tính & black, white, gray, silver & Màu văn bản chính, màu nền và đường viền phụ. \\
\hline
Màu Vàng / Cam & yellow, gold, orange, khaki & Năng động, rực rỡ, điểm nhấn lưu ý nhẹ. \\
\hline
Màu Tím (Purple) & purple, violet, indigo, orchid & Sang trọng, sáng tạo, thương hiệu hiện đại. \\
\hline
\end{tabularx}
\end{prettyTableBox}

\textbf{4. Các Ví Dụ Mã Nguồn Thực Tế Với Color Name:}

\begin{codeWindow}[Ví Dụ 1: Định Kiểu Thẻ Bài Viết Tin Tức (Article Card) Bằng Color Name]
\begin{lstlisting}[language=HTML5]
<!-- Structure HTML -->
<div class="article-card">
  <h2 class="title">Tin Tuc Cong Nghe Moi Nhat</h2>
  <p class="desc">Kham pha cac tinh nang CSS3 hien dai giup toi uu giao dien Web.</p>
  <span class="tag">HOT</span>
</div>

<!-- Styling CSS -->
<style>
.article-card {
  background-color: whitesmoke;
  border: 2px solid navy;
  padding: 16px;
  border-radius: 8px;
}
.title {
  color: crimson;
  font-size: 18px;
}
.desc {
  color: darkslategrey;
}
.tag {
  background-color: darkred;
  color: white;
  padding: 4px 8px;
}
</style>
\end{lstlisting}
\end{codeWindow}

\begin{codeWindow}[Ví Dụ 2: Định Kiểu Huy Hiệu Trạng Thái Hệ Thống (Notification Badges)]
\begin{lstlisting}[language=CSS3]
/* Huy hieu thanh cong */
.badge-success {
  background-color: seagreen;
  color: white;
}

/* Huy hieu loi khan cap */
.badge-error {
  background-color: firebrick;
  color: white;
}

/* Huy hieu canh bao */
.badge-warning {
  background-color: darkorange;
  color: black;
}
\end{lstlisting}
\end{codeWindow}

\begin{browserWindow}[Kết Quả Hiển Thị Của Thẻ Bài Viết \& Huy Hiệu Trạng Thái Bằng Color Name]
\vspace{2pt}
\begin{tcolorbox}[
    enhanced,
    colback=gray!5,
    colframe=blue!80!black,
    arc=6pt, boxrule=1.2pt,
    top=10pt, bottom=10pt, left=12pt, right=12pt
]
  {\sffamily\bfseries\large\color{red!70!black} Tin Tức Công Nghệ Mới Nhất (\codeinline{color: crimson;})}\\[4pt]
  {\sffamily\small\color{gray!40!black} Khám phá các tính năng CSS3 hiện đại giúp tối ưu giao diện Web. (\codeinline{color: darkslategrey;})}\\[8pt]
  
  % Badges preview
  \begin{tcolorbox}[nobeforeafter, colback=green!60!black, colframe=green!60!black, arc=3pt, boxrule=0.4pt, top=3pt, bottom=3pt, left=6pt, right=6pt, fontupper=\sffamily\bfseries\scriptsize\color{white}]
    THÀNH CÔNG (seagreen)
  \end{tcolorbox}
  \quad
  \begin{tcolorbox}[nobeforeafter, colback=red!70!black, colframe=red!70!black, arc=3pt, boxrule=0.4pt, top=3pt, bottom=3pt, left=6pt, right=6pt, fontupper=\sffamily\bfseries\scriptsize\color{white}]
    LỖI HỆ THỐNG (firebrick)
  \end{tcolorbox}
  \quad
  \begin{tcolorbox}[nobeforeafter, colback=orange!90!black, colframe=orange!90!black, arc=3pt, boxrule=0.4pt, top=3pt, bottom=3pt, left=6pt, right=6pt, fontupper=\sffamily\bfseries\scriptsize\color{black}]
    CẢNH BÁO (darkorange)
  \end{tcolorbox}
\end{tcolorbox}
\end{browserWindow}

\textbf{5. Đánh Giá Ưu Điểm Và Hạn Chế Của Color Name:}

\begin{prettyTableBox}[Bảng So Sánh Ưu Điểm Và Hạn Chế Của Tên Màu Định Danh]
\rowcolors{2}{white}{primaryBlue!4}
\sffamily\small
\begin{tabularx}{\linewidth}{|>{\raggedright\arraybackslash}X|>{\raggedright\arraybackslash}X|}
\hline
\rowcolor{primaryBlue} \thd{\sffamily Ưu điểm (Đúng)} & \thd{\sffamily Hạn chế (Không)} \\
\hline
\textcolor{green!60!black}{\textbf{Cực kỳ dễ nhớ, dễ viết}}: Đọc tên tiếng Anh là viết được ngay. & \textcolor{red!80!black}{\textbf{Rất ít tùy biến}}: Bị giới hạn cố định trong 147 tên màu. \\
\hline
\textcolor{green!60!black}{\textbf{Không cần tính toán}}: Không phải ghi nhớ hay chuyển đổi mã Hex/RGB. & \textcolor{red!80!black}{\textbf{Không có kênh Alpha}}: Không hỗ trợ tạo độ mờ trong suốt. \\
\hline
\textcolor{green!60!black}{\textbf{Tương thích tuyệt đối}}: Hỗ trợ $100\%$ trên tất cả trình duyệt web. & \textcolor{red!80!black}{\textbf{Thiếu chính xác}}: Không đáp ứng được các thiết kế UI/UX phức tạp. \\
\hline
\end{tabularx}
\end{prettyTableBox}

\textbf{6. Danh Sách Một Số Tên Màu Đẹp \& Phổ Biến NÊN Biết:}

\begin{prettyTableBox}[Bảng Tra Cứu Danh Sách Tên Màu Đẹp \& Mã HEX Tương Đương]
\rowcolors{2}{white}{primaryBlue!4}
\sffamily\small
\begin{tabularx}{\linewidth}{|>{\raggedright\arraybackslash \ttfamily\color{primaryBlue}\bfseries}p{3cm}|>{\ttfamily\color{accentOrange}\bfseries}p{2.5cm}|>{\raggedright\arraybackslash}X|}
\hline
\rowcolor{primaryBlue} \thd{\sffamily Tên màu (Color Name)} & \thd{\sffamily Mã HEX tương đương} & \thd{\sffamily Mô tả sắc thái màu} \\
\hline
tomato & \#FF6347 & Màu cam đỏ tươi tắn, năng động. \\
\hline
dodgerblue & \#1E90FF & Màu xanh lam sáng rực rỡ, hiện đại. \\
\hline
gold & \#FFD700 & Màu vàng ánh kim sang trọng. \\
\hline
slategray & \#708090 & Màu xám xanh trung tính dịu mắt. \\
\hline
darkorchid & \#9932CC & Màu tím lan đậm sáng tạo. \\
\hline
lightcoral & \#F08080 & Màu hồng san hô nhạt nữ tính. \\
\hline
indigo & \#4B0082 & Màu tím chàm sâu thẳm. \\
\hline
\end{tabularx}
\end{prettyTableBox}

\textbf{7. Hướng Dẫn Ra Quyết Định: Lưu Ý Khi Dùng Color Name:}

\begin{tipBox}[Quy Tắc Sử Dụng Tên Màu Trong Dự Án Web]
\begin{itemize}[leftmargin=1.5em, itemsep=3pt, parsep=0pt]
    \item \textbf{Không phân biệt chữ hoa / chữ thường}: CSS không phân biệt hoa thường với tên màu. Các cách viết \codeinline{red}, \codeinline{Red}, \codeinline{RED} đều cho kết quả hoàn toàn giống nhau.
    \item \textbf{Nên dùng khi}: Dựng khung demo nhanh, tạo ví dụ giảng dạy đơn giản hoặc thử nghiệm ý tưởng giao diện sơ khởi.
    \item \textbf{Không nên dùng khi}: Thiết kế giao diện sản phẩm chuyên nghiệp (Figma/Adobe), cần độ chính xác cao theo bộ nhận diện thương hiệu, hoặc khi cần độ mờ trong suốt (Alpha) và dải chuyển màu (Gradient).
\end{itemize}
\end{tipBox}

\paragraph{3.2. Hệ Màu HEX Trong CSS (Hexadecimal Color Model):}

\begin{conceptBox}[Khái Niệm Cơ Bản Về Hệ Màu HEX (Hexadecimal)]
\textbf{1. HEX Là Gì?}:
\begin{itemize}[leftmargin=1.5em, itemsep=3pt, parsep=0pt]
    \item \textbf{HEX (viết tắt của Hexadecimal)}: Là hệ màu biểu diễn giá trị màu sắc bằng hệ số \term{Thập lục phân} (cơ số 16).
    \item \textbf{Cấu trúc tổng quát}: Một mã màu HEX chuẩn có dạng: \codeinline{\#RRGGBB}. Trong đó:
    \begin{itemize}[leftmargin=1.2em, itemsep=2pt]
        \item \textbf{RR}: Giá trị thuộc về kênh màu Đỏ (\term{Red}).
        \item \textbf{GG}: Giá trị thuộc về kênh màu Xanh lá (\term{Green}).
        \item \textbf{BB}: Giá trị thuộc về kênh màu Xanh dương (\term{Blue}).
    \end{itemize}
    \item \textbf{Bản chất kênh màu}: Mỗi cặp giá trị ($RR, GG, BB$) gồm 2 chữ số thập lục phân ($00 \rightarrow \text{FF}$), tương đương với giá trị từ $0$ đến $255$ trong hệ thập phân (Decimal).
\end{itemize}

\textbf{2. Cách Đọc Và Giải Mã Một Mã Màu HEX}:
Ví dụ phân tích mã màu \codeinline{\#FF5733}:
\begin{center}
\small\sffamily
\begin{tabular}{|c|c|c|L{5.5cm}|}
\hline
\rowcolor{primaryBlue} \thd{\sffamily Thành phần} & \thd{\sffamily Giá trị HEX} & \thd{\sffamily Giá trị thập phân} & \thd{\sffamily Giải thích mức độ phát sáng} \\
\hline
Red ($RR$) & \codeinline{FF} & $15 \times 16^1 + 15 \times 16^0 = \mathbf{255}$ & Cường độ phát sáng màu Đỏ cực đại ($100\%$). \\
\hline
Green ($GG$) & \codeinline{57} & $5 \times 16^1 + 7 \times 16^0 = \mathbf{87}$ & Cường độ màu Xanh lá trung bình ($34\%$). \\
\hline
Blue ($BB$) & \codeinline{33} & $3 \times 16^1 + 3 \times 16^0 = \mathbf{51}$ & Cường độ màu Xanh dương thấp ($20\%$). \\
\hline
\end{tabular}
\end{center}
$\rightarrow$ Kết quả phối hợp 3 ánh sáng tạo ra màu \textbf{Cam Đỏ (Orange-Red)} rực rỡ và ấm áp.
\end{conceptBox}

\textbf{3. Cú Pháp Khai Báo Chuẩn Trong CSS \& Ví Dụ Minh Họa:}
\begin{codeWindow}[Cú Pháp Định Dạng Mã Màu HEX Trong CSS]
\begin{lstlisting}[language=CSS3]
/* Cu phap chung */
color: #RRGGBB;

/* Cac mau co ban pho bien */
.black-text { color: #000000; } /* Den tuyen (R=0, G=0, B=0) */
.white-text { color: #ffffff; } /* Trang tinh (R=255, G=255, B=255) */
.red-text   { color: #ff0000; } /* Do thuan (R=255, G=0, B=0) */
.green-text { color: #00ff00; } /* Xanh la thuan (R=0, G=255, B=0) */
.blue-text  { color: #0000ff; } /* Xanh duong thuan (R=0, G=0, B=255) */
\end{lstlisting}
\end{codeWindow}

\begin{browserWindow}[Kết Quả Hiển Thị Trực Quan Của Các Mã Màu HEX Cơ Bản]
\vspace{2pt}
\begin{tcolorbox}[
    enhanced,
    colback=white,
    colframe=gray!30,
    arc=4pt, boxrule=0.6pt,
    top=8pt, bottom=8pt, left=10pt, right=10pt
]
  \begin{tcolorbox}[colback=black, colframe=black, arc=3pt, boxrule=0.4pt, top=4pt, bottom=4pt, left=8pt, right=8pt, fontupper=\sffamily\bfseries\color{white}]
    .black-text (\#000000): Chữ màu Đen tuyền tuyệt đối
  \end{tcolorbox}
  \vspace{3pt}
  \begin{tcolorbox}[colback=white, colframe=gray!40, arc=3pt, boxrule=0.4pt, top=4pt, bottom=4pt, left=8pt, right=8pt, fontupper=\sffamily\bfseries\color{black}]
    .white-text (\#ffffff): Chữ màu Trắng tinh nét trên nền xám nhẹ
  \end{tcolorbox}
  \vspace{3pt}
  \begin{tcolorbox}[colback=red!10, colframe=red!60, arc=3pt, boxrule=0.4pt, top=4pt, bottom=4pt, left=8pt, right=8pt, fontupper=\sffamily\bfseries\color{red!80!black}]
    .red-text (\#ff0000): Chữ màu Đỏ thuần tươi tắn
  \end{tcolorbox}
  \vspace{3pt}
  \begin{tcolorbox}[colback=green!10, colframe=green!60!black, arc=3pt, boxrule=0.4pt, top=4pt, bottom=4pt, left=8pt, right=8pt, fontupper=\sffamily\bfseries\color{green!50!black}]
    .green-text (\#00ff00): Chữ màu Xanh lá thuần chuẩn
  \end{tcolorbox}
  \vspace{3pt}
  \begin{tcolorbox}[colback=blue!10, colframe=blue!60, arc=3pt, boxrule=0.4pt, top=4pt, bottom=4pt, left=8pt, right=8pt, fontupper=\sffamily\bfseries\color{blue!80!black}]
    .blue-text (\#0000ff): Chữ màu Xanh dương thuần đậm
  \end{tcolorbox}
\end{tcolorbox}
\end{browserWindow}

\textbf{4. Cú Pháp Viết Gọn Mã HEX (Shorthand Hex Syntax):}

\begin{conceptBox}[Quy Trình Rút Gọn Mã Màu HEX]
\begin{itemize}[leftmargin=1.5em, itemsep=3pt, parsep=0pt]
    \item \textbf{Điều kiện rút gọn}: Nếu mỗi cặp ký tự đại diện cho từng kênh màu hoàn toàn giống nhau (\codeinline{\#AABBCC} trong đó $AA$, $BB$, $CC$), bạn có thể rút gọn mã Hex 6 ký tự thành mã Hex 3 ký tự \codeinline{\#ABC}.
    \item \textbf{Công thức quy đổi}: \codeinline{\#ABC} $\Leftrightarrow$ \codeinline{\#AABBCC}.
\end{itemize}
\end{conceptBox}

\begin{prettyTableBox}[Bảng So Sánh Các Mã HEX Đầy Đủ Và Dạng Viết Gọn]
\rowcolors{2}{white}{primaryBlue!4}
\sffamily\small
\begin{tabularx}{\linewidth}{|>{\raggedright\arraybackslash \ttfamily\color{primaryBlue}\bfseries}p{3cm}|>{\ttfamily\color{accentOrange}\bfseries}p{2.5cm}|>{\centering\arraybackslash}p{2.2cm}|>{\raggedright\arraybackslash}X|}
\hline
\rowcolor{primaryBlue} \thd{\sffamily Viết đầy đủ (6 ký tự)} & \thd{\sffamily Viết gọn (3 ký tự)} & \thd{\sffamily Hợp lệ?} & \thd{\sffamily Ghi chú chi tiết} \\
\hline
\#FFFFFF & \#FFF & \textcolor{green!60!black}{\textbf{Đúng}} & Cả 3 cặp FF, FF, FF đều trùng ký tự. \\
\hline
\#000000 & \#000 & \textcolor{green!60!black}{\textbf{Đúng}} & Cả 3 cặp 00, 00, 00 đều trùng ký tự. \\
\hline
\#FF0000 & \#F00 & \textcolor{green!60!black}{\textbf{Đúng}} & Tương ứng FF $\rightarrow$ F, 00 $\rightarrow$ 0, 00 $\rightarrow$ 0. \\
\hline
\#112233 & \#123 & \textcolor{green!60!black}{\textbf{Đúng}} & Tương ứng 11 $\rightarrow$ 1, 22 $\rightarrow$ 2, 33 $\rightarrow$ 3. \\
\hline
\#123456 & Không thể viết gọn & \textcolor{red!80!black}{\textbf{Không}} & Các cặp 12, 34, 56 không có ký tự trùng lặp trong cặp! \\
\hline
\end{tabularx}
\end{prettyTableBox}

\begin{warnBox}[Lưu Ý Quan Trọng Khi Viết Gọn Mã HEX]
Bạn chỉ được phép rút gọn mã Hex khi \textbf{từng cặp ký tự đơn lẻ} trùng nhau (ví dụ: \codeinline{FF}, \codeinline{00}, \codeinline{33}). Tuyệt đối không được rút gọn nếu hai ký tự trong cùng một cặp khác nhau (ví dụ: \codeinline{\#FF5733} không thể viết gọn thành \codeinline{\#F53}).
\end{warnBox}

\textbf{5. Danh Sách Các Mã Màu HEX Phổ Biến Trong Thiết Kế Web:}

\begin{prettyTableBox}[Bảng Tra Cứu Danh Sách Mã Màu HEX Phổ Biến]
\rowcolors{2}{white}{primaryBlue!4}
\sffamily\small
\begin{tabularx}{\linewidth}{|>{\raggedright\arraybackslash \bfseries\color{primaryBlue}}p{2.8cm}|>{\ttfamily\color{accentOrange}\bfseries}p{2.5cm}|>{\ttfamily}p{2.5cm}|>{\raggedright\arraybackslash}X|}
\hline
\rowcolor{primaryBlue} \thd{\sffamily Tên màu} & \thd{\sffamily Mã HEX} & \thd{\sffamily Mã RGB tương đương} & \thd{\sffamily Gợi ý ứng dụng thực tế} \\
\hline
Đen (Black) & \#000000 & rgb(0, 0, 0) & Màu tối hoàn toàn, màu chữ chính trên nền sáng. \\
\hline
Trắng (White) & \#FFFFFF & rgb(255, 255, 255) & Màu sáng hoàn toàn, màu nền trang chính. \\
\hline
Đỏ (Red) & \#FF0000 & rgb(255, 0, 0) & RGB Full Red -- Thể hiện trạng thái lỗi khẩn cấp. \\
\hline
Xanh lá (Green) & \#00FF00 & rgb(0, 255, 0) & RGB Full Green -- Trạng thái thành công, hoàn thành. \\
\hline
Xanh dương (Blue) & \#0000FF & rgb(0, 0, 255) & RGB Full Blue -- Liên kết (Link) và thông tin. \\
\hline
Cam (Orange) & \#FFA500 & rgb(255, 165, 0) & Nút cảnh báo nhẹ, nút hành động nhấn mạnh. \\
\hline
Vàng (Yellow) & \#FFFF00 & rgb(255, 255, 0) & Kết hợp Red + Green -- Điểm nhấn lưu ý. \\
\hline
Tím (Purple) & \#800080 & rgb(128, 0, 128) & Kết hợp Red + Blue -- Thương hiệu sáng tạo. \\
\hline
\end{tabularx}
\end{prettyTableBox}

\textbf{6. Thuật Toán Chuyển Đổi Giữa Mã HEX Và Hệ RGB:}

\begin{conceptBox}[Quy Trình Chuyển Đổi Mã HEX Sang RGB]
Chuyển đổi mã \codeinline{\#FF5733} sang cú pháp \codeinline{rgb(...)}:
\begin{enumerate}[leftmargin=1.5em, itemsep=2pt]
    \item Tách mã Hex thành 3 cặp byte: $\text{Red} = \text{FF}_{16}$, $\text{Green} = 57_{16}$, $\text{Blue} = 33_{16}$.
    \item Đổi từ hệ Hex (cơ số 16) sang hệ Thập phân (cơ số 10):
    \begin{itemize}[leftmargin=1.2em]
        \item Kênh Red ($\text{FF}_{16}$): $\rightarrow \mathbf{255}$.
        \item Kênh Green ($57_{16}$): $\rightarrow \mathbf{87}$.
        \item Kênh Blue ($33_{16}$): $\rightarrow \mathbf{51}$.
    \end{itemize}
    \item Kết quả cú pháp tương đương: $\rightarrow$ \codeinline{rgb(255, 87, 51)}.
\end{enumerate}
\textit{Mẹo nhỏ}: Bạn có thể tra cứu và chuyển đổi nhanh chóng tại trang web \codeinline{https://www.colorhexa.com}.
\end{conceptBox}

\textbf{7. So Sánh Chi Tiết Giữa Mã HEX, RGB Và RGBA:}

\begin{prettyTableBox}[Bảng So Sánh Toàn Diện Giữa Mã HEX, RGB Và RGBA]
\rowcolors{2}{white}{primaryBlue!4}
\sffamily\small
\begin{tabularx}{\linewidth}{|>{\raggedright\arraybackslash \bfseries\color{primaryBlue}}p{2.6cm}|>{\raggedright\arraybackslash}X|>{\raggedright\arraybackslash}X|>{\raggedright\arraybackslash}X|}
\hline
\rowcolor{primaryBlue} \thd{\sffamily Tiêu chí} & \thd{\sffamily Mã HEX (\#RRGGBB)} & \thd{\sffamily Hệ RGB (rgb(r,g,b))} & \thd{\sffamily Hệ RGBA (rgba(r,g,b,a))} \\
\hline
Độ dễ nhớ & Khó ghi nhớ bằng mắt thường. & Dễ hiểu trực quan theo giá trị số từ $0 \rightarrow 255$. & Dễ hiểu trực quan theo giá trị số từ $0 \rightarrow 255$. \\
\hline
Dạng viết gọn & \textcolor{green!60!black}{\textbf{Có hỗ trợ}} (mã 3 ký tự \codeinline{\#FFF}). & \textcolor{red!80!black}{\textbf{Không}} (phải viết đầy đủ 3 tham số). & \textcolor{red!80!black}{\textbf{Không}} (phải viết đầy đủ 4 tham số). \\
\hline
Kênh Alpha (Trong suốt) & \textcolor{red!80!black}{\textbf{Không}} (ở chuẩn truyền thống 6 ký tự). & \textcolor{red!80!black}{\textbf{Không}} (chỉ màu đặc $100\%$). & \textcolor{green!60!black}{\textbf{Có hỗ trợ}} (có kênh Alpha $0.0 \rightarrow 1.0$). \\
\hline
Thân thiện thiết kế & Rất phổ biến, xuất chuẩn từ Figma / Adobe. & Dùng nhiều khi tính toán màu động trong JavaScript. & Hoàn hảo cho hiệu ứng Overlay, bóng mờ. \\
\hline
\end{tabularx}
\end{prettyTableBox}

\textbf{8. Hướng Dẫn Ra Quyết Định: Khi Nào Nên Và Không Nên Dùng HEX?}

\begin{conceptBox}[Ma Trận Hướng Dẫn Trường Hợp Sử Dụng Mã HEX]
\begin{itemize}[leftmargin=1.5em, itemsep=3pt, parsep=0pt]
    \item \textbf{Nên dùng HEX khi (\textcolor{green!60!black}{\textbf{Khuyên dùng}})}:
    \begin{itemize}[leftmargin=1.2em, itemsep=2pt]
        \item Cần chọn màu tĩnh nhanh chóng cho các phần tử UI cố định.
        \item Làm việc trực tiếp với file thiết kế giao diện (Figma, Sketch, Adobe XD).
        \item Ghi mã màu ngắn gọn để tối ưu dung lượng file CSS nhỏ gọn.
    \end{itemize}
    \item \textbf{Không nên dùng HEX khi (\textcolor{red!80!black}{\textbf{Hạn chế}})}:
    \begin{itemize}[leftmargin=1.2em, itemsep=2pt]
        \item Cần tạo màu nền bán trong suốt (\term{Alpha opacity}) cho Modal/Popup.
        \item Cần thay đổi màu sắc động tăng giảm mượt mà bằng lập trình JavaScript.
        \item Cần hòa trộn dải màu (\term{Overlay Blend Mode}) dịu mắt với hình ảnh nền bên dưới.
    \end{itemize}
\end{itemize}
\end{conceptBox}

\textbf{9. Cơ Chế Hoạt Động Nội Bộ Của Máy Tính \& Trình Duyệt Khi Xử Lý HEX:}

\begin{conceptBox}[Quy Trình Xử Lý Nội Bội Của Trình Duyệt Khi Gặp Mã HEX]
\textbf{Các bước xử lý từ mã nguồn đến mắt người xem}:
\begin{enumerate}[leftmargin=1.5em, itemsep=3pt, parsep=0pt]
    \item \textbf{Người lập trình}: Khai báo mã màu trong CSS (ví dụ: \codeinline{color: \#FF5733;}).
    \item \textbf{Trình duyệt Web}: Giải mã chuỗi Hex 16 sang hệ RGB 10: $\text{FF} \rightarrow 255, 57 \rightarrow 87, 33 \rightarrow 51$.
    \item \textbf{Hệ thống phần cứng (Màn hình)}: Truyền tín hiệu điện điều khiển 3 điểm ảnh phụ (\term{Subpixels}) Red, Green, Blue phát sáng theo đúng tỷ lệ cường độ.
    \item \textbf{Mắt người xem}: Mắt tiếp nhận chùm ánh sáng cộng phát ra từ màn hình và cảm nhận được màu Cam Đỏ.
\end{enumerate}

\textbf{Mở rộng chuyên sâu -- Mã HEX 8 Ký Tự (CSS Color Module Level 4)}:
Trong tiêu chuẩn CSS hiện đại, bạn có thể bổ sung kênh Alpha trực tiếp vào mã Hex bằng cú pháp 8 ký tự \codeinline{\#RRGGBBAA} (ví dụ: \codeinline{\#FF573380} tương đương với độ mờ $50\%$).
\end{conceptBox}

\textbf{10. Ví Dụ Thực Tế Component UI Với Mã Màu HEX:}

\begin{codeWindow}[Ví Dụ 1: Định Kiểu Thẻ Sản Phẩm E-Commerce Với Mã HEX 6 Ký Tự Chuẩn Figma]
\begin{lstlisting}[language=HTML5]
<!-- Product Card Structure -->
<div class="product-card">
  <span class="sale-badge">-25%</span>
  <h3 class="product-name">Tai Nghe Bluetooth Wireless</h3>
  <p class="price">1.250.000 VND</p>
  <button class="btn-buy">Them Vao Gio Hang</button>
</div>

<!-- Styling with HEX 6-digit colors -->
<style>
.product-card {
  background-color: #f8fafc; /* Nen xam trang sang diu */
  border: 1px solid #e2e8f0;    /* Vien xam nhat */
  border-radius: 8px;
  padding: 16px;
}
.sale-badge {
  background-color: #ef4444; /* Do tuoi giam gia */
  color: #ffffff;            /* Trang tinh */
  padding: 2px 6px;
}
.product-name {
  color: #1e293b;            /* Xam den dam de doc */
}
.price {
  color: #059669;            /* Xanh la dam gia tien */
  font-weight: bold;
}
.btn-buy {
  background-color: #3b82f6; /* Xanh duong chinh (Primary) */
  color: #ffffff;
  border: none;
}
</style>
\end{lstlisting}
\end{codeWindow}

\begin{codeWindow}[Ví Dụ 2: Sử Dụng Mã HEX 8 Ký Tự Có Kênh Alpha (\#RRGGBBAA) Trong CSS4]
\begin{lstlisting}[language=CSS3]
/* Lop phu Modal nen den mo 50% opacity (#00000080) */
.modal-backdrop {
  background-color: #00000080; /* 80 trong Hex = 128 trong Decimal = 50% Alpha */
}

/* Tag thong bao xanh duong nhat mo 20% (#3b82f633) */
.info-tag {
  background-color: #3b82f633; /* 33 trong Hex = 51 trong Decimal = 20% Alpha */
  color: #1d4ed8;
  border: 1px solid #93c5fd;
}
\end{lstlisting}
\end{codeWindow}

\begin{browserWindow}[Kết Quả Hiển Thị Của Thẻ Sản Phẩm \& Mã HEX 8 Ký Tự Alpha]
\vspace{2pt}
\begin{tcolorbox}[
    enhanced,
    colback=white,
    colframe=gray!30,
    arc=6pt, boxrule=0.8pt,
    top=10pt, bottom=10pt, left=12pt, right=12pt
]
  % E-Commerce Product Card Preview
  \begin{tcolorbox}[
      enhanced,
      colback=gray!8,
      colframe=gray!30,
      arc=4pt, boxrule=0.6pt,
      top=8pt, bottom=8pt, left=10pt, right=10pt
  ]
    \begin{tcolorbox}[nobeforeafter, colback=red!80!black, colframe=red!80!black, arc=2pt, boxrule=0.3pt, top=2pt, bottom=2pt, left=5pt, right=5pt, fontupper=\sffamily\bfseries\scriptsize\color{white}]
      -25\% (\#ef4444)
    \end{tcolorbox}\\[4pt]
    {\sffamily\bfseries\large\color{gray!90!black} Tai Nghe Bluetooth Wireless (\#1e293b)}\\[2pt]
    {\sffamily\bfseries\color{green!60!black} Giá: 1.250.000 VNĐ (\#059669)}\\[6pt]
    \begin{tcolorbox}[nobeforeafter, colback=blue!70!black, colframe=blue!70!black, arc=3pt, boxrule=0.4pt, top=3pt, bottom=3pt, left=8pt, right=8pt, fontupper=\sffamily\bfseries\small\color{white}]
      Thêm Vào Giỏ Hàng (\#3b82f6)
    \end{tcolorbox}
  \end{tcolorbox}
  \vspace{4pt}
  
  % HEX 8-digit Alpha Preview
  \begin{tcolorbox}[colback=blue!12, colframe=blue!40, arc=3pt, boxrule=0.4pt, top=4pt, bottom=4pt, left=8pt, right=8pt, fontupper=\sffamily\small\color{blue!80!black}]
    \textbf{Tag Nổi Bật với Mã HEX 8 Ký Tự}: \codeinline{background-color: \#3b82f633;} (Nền Xanh dương mờ 20\%)
  \end{tcolorbox}
\end{tcolorbox}
\end{browserWindow}

\paragraph{3.3. Hệ Màu RGB Và RGBA (Red, Green, Blue, Alpha):}

\begin{conceptBox}[Khái Niệm Cơ Bản Về Mô Hình Màu RGB \& Kênh RGBA]
\textbf{1. Bản Chất Kỹ Thuật Của Hệ Màu RGB}:
\begin{itemize}[leftmargin=1.5em, itemsep=3pt, parsep=0pt]
    \item \textbf{RGB (Red - Green - Blue)}: Là mô hình biểu diễn màu dựa trên \term{Nguyên lý ánh sáng cộng} (\term{Additive Color Model}). Trong mô hình này, các màu sắc trên màn hình được tạo thành từ việc phát chiếu và phối hợp 3 nguồn chùm ánh sáng màu cơ bản: Đỏ (Red), Xanh lá (Green) và Xanh dương (Blue). Đây là tiêu chuẩn mô hình màu được sử dụng phổ biến nhất trên các thiết bị điện tử như màn hình máy tính, TV, điện thoại và giao diện Web.
    \item \textbf{Thang cường độ phát sáng}: Mỗi kênh màu $R, G, B$ nhận một giá trị số nguyên từ \codeinline{0} đến \codeinline{255} ($0 \rightarrow 255$), đại diện cho $256$ mức cường độ sáng phát ra từ điểm ảnh ($0$ là tắt hoàn toàn, $255$ là phát sáng cực đại). Khi phối hợp cả 3 kênh, hệ màu RGB tạo ra tổng cộng $256 \times 256 \times 256 = 16.777.216$ màu sắc khác nhau (đạt chuẩn màu \term{True Color 24-bit}).
\end{itemize}

\textbf{2. Sự Mở Rộng Kênh Alpha ($A$) Trong RGBA}:
\begin{itemize}[leftmargin=1.5em, itemsep=3pt, parsep=0pt]
    \item \textbf{RGBA (Red - Green - Blue - Alpha)}: Là biến thể mở rộng trực tiếp từ RGB, bổ sung thêm kênh thứ tư là \textbf{Alpha ($A$)} đại diện cho chỉ số \term{Độ trong suốt} (\term{Opacity / Transparency}).
    \item \textbf{Dải giá trị Alpha}: Kênh Alpha nhận giá trị số thực (\term{Float}) trong khoảng từ \codeinline{0.0} (hoàn toàn trong suốt $0\%$, đè lên nội dung khác nhưng nhìn xuyên thấu hoàn toàn) đến \codeinline{1.0} (đặc nguyên bản $100\%$, không trong suốt).
\end{itemize}
\end{conceptBox}

\textbf{2. Cú Pháp Định Dạng Chuẩn Trong CSS \& Bảng Tham Số Chi Tiết:}
\begin{codeWindow}[Cú Pháp Khai Báo RGB \& RGBA Trong CSS]
\begin{lstlisting}[language=CSS3]
/* Cu phap RGB truyen thong (3 tham so kenh mau) */
color: rgb(R, G, B);

/* Cu phap RGBA bo sung kenh Alpha (4 tham so) */
color: rgba(R, G, B, A);
\end{lstlisting}
\end{codeWindow}

\begin{prettyTableBox}[Bảng Giải Thích Chi Tiết Các Tham Số RGB / RGBA]
\rowcolors{2}{white}{primaryBlue!4}
\sffamily\small
\begin{tabularx}{\linewidth}{|>{\raggedright\arraybackslash \ttfamily\color{primaryBlue}\bfseries}c|>{\raggedright\arraybackslash}X|>{\ttfamily\color{accentOrange}}c|}
\hline
\rowcolor{primaryBlue} \thd{\sffamily Tham số} & \thd{\sffamily Ý nghĩa \& Vai trò chuyên môn} & \thd{\sffamily Dải giá trị chuẩn} \\
\hline
R & Cường độ phát sáng của kênh màu Đỏ (Red). & $0 \rightarrow 255$ \\
\hline
G & Cường độ phát sáng của kênh màu Xanh lá (Green). & $0 \rightarrow 255$ \\
\hline
B & Cường độ phát sáng của kênh màu Xanh dương (Blue). & $0 \rightarrow 255$ \\
\hline
A & Kênh độ trong suốt (Alpha Channel) kiểm soát mức độ nhìn xuyên qua. & $0.0 \rightarrow 1.0$ \\
\hline
\end{tabularx}
\end{prettyTableBox}

\textbf{3. Quy Tắc Sử Dụng Dấu Phẩy \& Sự Khác Biệt Giữa Cú Pháp Truyền Thống vs CSS Color Level 4:}

\begin{conceptBox}[Phân Tích Quy Tắc Sử Dụng Dấu Phẩy Trong Cú Pháp RGB / RGBA]
\textbf{a. Cú Pháp Truyền Thống (Phải Có Dấu Phẩy `,`)}:
\begin{itemize}[leftmargin=1.5em, itemsep=3pt, parsep=0pt]
    \item Khi sử dụng cú pháp hàm \codeinline{rgb(R, G, B)} hoặc \codeinline{rgba(R, G, B, A)} tiêu chuẩn truyền thống, \textbf{bắt buộc phải sử dụng dấu phẩy} giữa các tham số.
    \item \textcolor{red!80!black}{\textbf{Sai cú pháp (Không có dấu phẩy)}}:
    \codeinline{color: rgba(255 0 0 0.5);} $\rightarrow$ \textcolor{red!80!black}{\textbf{KHÔNG hợp lệ}} trong chuẩn CSS thông thường. Trình duyệt sẽ coi đây là lỗi cú pháp và không áp dụng thuộc tính màu này.
    \item \textcolor{green!60!black}{\textbf{Đúng cú pháp (Có dấu phẩy chuẩn)}}:
    \codeinline{color: rgba(255, 0, 0, 0.5);} $\rightarrow$ \textcolor{green!60!black}{\textbf{Đúng chuẩn}}, hiển thị màu đỏ trong suốt 50\%.
\end{itemize}

\textbf{b. Cú Pháp Hiện Đại CSS Color Module Level 4 (Dùng Khoảng Trắng \& Dấu Gạch Chéo `/`)}:
\begin{itemize}[leftmargin=1.5em, itemsep=3pt, parsep=0pt]
    \item Theo tiêu chuẩn mới (\term{CSS Color Module Level 4}), trình duyệt cho phép dùng khoảng trắng thay dấu phẩy và dùng dấu gạch chéo \codeinline{/} để ngăn kênh Alpha với cú pháp \codeinline{rgb()}:
    \item \textcolor{green!60!black}{\textbf{Cú pháp hợp lệ trong trình duyệt hiện đại}}:
    \codeinline{color: rgb(255 0 0 / 0.5);} $\rightarrow$ \textcolor{green!60!black}{\textbf{Hợp lệ chuẩn CSS4}}, tương đương với \codeinline{rgba(255, 0, 0, 0.5)}.
\end{itemize}
\end{conceptBox}

\begin{tipBox}[Ghi Nhớ Mẹo Sử Dụng Cú Pháp Khai Báo Màu]
\begin{itemize}[leftmargin=1.5em, itemsep=3pt, parsep=0pt]
    \item Với cú pháp \codeinline{rgb(...)} hoặc \codeinline{rgba(...)} kiểu truyền thống $\rightarrow$ **Dùng dấu phẩy `,`**.
    \item Với cú pháp kiểu mới CSS4 $\rightarrow$ **Dùng dấu cách và dấu gạch chéo `/`** để ngăn kênh độ trong suốt Alpha.
\end{itemize}
\end{tipBox}

\begin{prettyTableBox}[Bảng Tóm Tắt Tính Hợp Lệ Của Các Cách Viết Cú Pháp RGBA / RGB]
\rowcolors{2}{white}{primaryBlue!4}
\sffamily\small
\begin{tabularx}{\linewidth}{|>{\raggedright\arraybackslash \ttfamily\color{primaryBlue}\bfseries}p{4.0cm}|>{\centering\arraybackslash}p{2.2cm}|>{\raggedright\arraybackslash}X|}
\hline
\rowcolor{primaryBlue} \thd{\sffamily Cách viết cú pháp} & \thd{\sffamily Hợp lệ?} & \thd{\sffamily Ghi chú \& Độ tương thích trình duyệt} \\
\hline
rgba(255, 0, 0, 0.5) & \textcolor{green!60!black}{\textbf{Hợp lệ}} & Chuẩn truyền thống, hỗ trợ 100\% trên tất cả mọi trình duyệt. \\
\hline
rgb(255 0 0 / 0.5) & \textcolor{green!60!black}{\textbf{Hợp lệ}} & Chuẩn mới CSS4 -- Cần trình duyệt phiên bản hiện đại. \\
\hline
rgba(255 0 0 0.5) & \textcolor{red!80!black}{\textbf{Không hợp lệ}} & Sai cú pháp hoàn toàn -- Trình duyệt không chấp nhận. \\
\hline
\end{tabularx}
\end{prettyTableBox}

\textbf{4. Ví Dụ Cụ Thể \& Kết Quả Hiển Thị Trực Quan:}

\begin{codeWindow}[Khai Báo Định Dạng RGB Cơ Bản Và RGBA Trong Suốt]
\begin{lstlisting}[language=CSS3]
/* RGB co ban: do dac 100% va nen xanh la dac 100% */
.basic-text {
  color: rgb(255, 0, 0);          /* Mau do thuan */
  background-color: rgb(0, 255, 0); /* Mau nen xanh la */
  padding: 10px;
}

/* RGBA co do trong suot: xanh duong mo 50% */
.translucent-box {
  background-color: rgba(0, 0, 255, 0.5); /* Alpha = 0.5 */
  color: white;
  padding: 12px;
}
\end{lstlisting}
\end{codeWindow}

\begin{browserWindow}[Kết Quả Hiển Thị Trực Quan Của Mã CSS Trên]
\vspace{2pt}
% 1. Preview .basic-text
\begin{tcolorbox}[
    enhanced,
    colback=green!80!black,
    colframe=green!60!black,
    arc=4pt, boxrule=0.6pt,
    top=6pt, bottom=6pt, left=10pt, right=10pt,
    fontupper=\sffamily\bfseries\color{red}
]
.basic-text: Chữ màu Đỏ rgb(255, 0, 0) trên nền Xanh lá rgb(0, 255, 0)
\end{tcolorbox}

\vspace{4pt}

% 2. Preview .translucent-box trên lớp nền phía dưới
\begin{tcolorbox}[
    enhanced,
    colback=primaryBlue!5,
    colframe=primaryBlue!20,
    arc=4pt, boxrule=0.6pt,
    top=8pt, bottom=8pt, left=10pt, right=10pt
]
  {\sffamily\small\color{gray!70}\itshape Dòng văn bản nền phía dưới trang web... (Lớp giao diện nằm ở phía đằng sau)}\\[4pt]
  \begin{tcolorbox}[
      enhanced,
      colback=blue!45!white,
      colframe=blue!60!black,
      arc=4pt, boxrule=0.6pt,
      top=6pt, bottom=6pt, left=10pt, right=10pt,
      fontupper=\sffamily\bfseries\color{white}
  ]
  .translucent-box: Nền Xanh dương mờ 50\% rgba(0, 0, 255, 0.5)\\
  \normalfont\small(Hiệu ứng độ mờ bán trong suốt giúp nhìn xuyên qua nền)
  \end{tcolorbox}
\end{tcolorbox}
\end{browserWindow}

\textbf{5. Nguyên Lý Phối Màu Tuỳ Chỉnh Bằng RGB:}

\begin{prettyTableBox}[Bảng Nguyên Lý Phối Màu Tuỳ Chỉnh Phổ Biến Bằng RGB]
\rowcolors{2}{white}{primaryBlue!4}
\sffamily\small
\begin{tabularx}{\linewidth}{|>{\raggedright\arraybackslash \bfseries\color{primaryBlue}}p{2.5cm}|>{\ttfamily\color{accentOrange}}p{3.5cm}|>{\raggedright\arraybackslash}X|}
\hline
\rowcolor{primaryBlue} \thd{\sffamily Tên màu} & \thd{\sffamily Cú pháp RGB} & \thd{\sffamily Mô tả nguyên lý phối ánh sáng} \\
\hline
Đen (Black) & rgb(0, 0, 0) & Không có ánh sáng phát ra ở cả 3 kênh ($R=G=B=0$). \\
\hline
Trắng (White) & rgb(255, 255, 255) & Tất cả 3 kênh màu phát sáng ở mức tối đa ($R=G=B=255$). \\
\hline
Xám (Gray) & rgb(128, 128, 128) & Cả 3 kênh màu có giá trị bằng nhau ($R=G=B$) tạo ra màu xám trung tính. \\
\hline
Vàng (Yellow) & rgb(255, 255, 0) & Kết hợp ánh sáng Đỏ + Xanh lá cực đại ($R=255, G=255, B=0$). \\
\hline
Tím (Purple) & rgb(128, 0, 128) & Kết hợp ánh sáng Đỏ + Xanh dương ($R=128, G=0, B=128$). \\
\hline
Cam (Orange) & rgb(255, 165, 0) & Kết hợp Đỏ tối đa + một ít Xanh lá ($R=255, G=165, B=0$). \\
\hline
\end{tabularx}
\end{prettyTableBox}

\textbf{6. Đánh Giá Ưu Điểm Và Hạn Chế Của Hệ Màu RGB:}

\begin{conceptBox}[Phân Tích Ưu Điểm \& Hạn Chế Của Hệ Màu RGB / RGBA]
\begin{itemize}[leftmargin=1.5em, itemsep=3pt, parsep=0pt]
    \item \textbf{Ưu điểm (\textcolor{green!60!black}{\textbf{Nổi bật}})}:
    \begin{itemize}[leftmargin=1.2em, itemsep=2pt]
        \item Hỗ trợ độ trong suốt (\term{Alpha Channel}) linh hoạt với cú pháp \codeinline{rgba()}.
        \item Cực kỳ dễ dàng điều chỉnh màu sắc động bằng JavaScript (ví dụ: \codeinline{element.style.backgroundColor = `rgba(255, 100, 0, 0.8)`}).
        \item Hỗ trợ tạo dải màu chuyển tiếp (\codeinline{linear-gradient}) và hiệu ứng lớp phủ (\term{Overlay}) mượt mà.
    \end{itemize}
    \item \textbf{Hạn chế (\textcolor{red!80!black}{\textbf{Hạn chế}})}:
    \begin{itemize}[leftmargin=1.2em, itemsep=2pt]
        \item Trực quan không dễ đọc hoặc đoán màu bằng mắt thường so với mã Hex hoặc HSL.
        \item Cú pháp viết dài hơn so với mã Hex rút gọn hoặc tên từ khóa màu có sẵn.
        \item Khó ghi nhớ giá trị số chính xác nếu không sử dụng công cụ Color Picker chuyên dụng.
    \end{itemize}
\end{itemize}
\end{conceptBox}

\textbf{7. Hiệu Ứng Khi Thay Đổi Kênh Alpha ($A$) Trong RGBA:}

\begin{prettyTableBox}[Bảng Mức Độ Mờ Quan Sát Được Theo Kênh Alpha]
\rowcolors{2}{white}{primaryBlue!4}
\sffamily\small
\begin{tabularx}{\linewidth}{|>{\raggedright\arraybackslash \ttfamily\color{primaryBlue}\bfseries}c|>{\raggedright\arraybackslash}X|}
\hline
\rowcolor{primaryBlue} \thd{\sffamily Giá trị Alpha (A)} & \thd{\sffamily Hiệu ứng quan sát trực quan trên giao diện} \\
\hline
1.0 & Màu đậm đặc nguyên bản, rõ ràng 100\%, không xuyên thấu. \\
\hline
0.5 & Màu trong mờ 50\%, cho phép nhìn xuyên qua 50\% nội dung lớp nền bên dưới. \\
\hline
0.1 & Gần như trong suốt hoàn toàn, chỉ còn một vệt màu cực kỳ nhẹ. \\
\hline
0.0 & Hoàn toàn trong suốt 100\% (ẩn màu hoàn toàn trên giao diện). \\
\hline
\end{tabularx}
\end{prettyTableBox}

\begin{conceptBox}[Ứng Dụng Phổ Biến Của RGBA Trong Thiết Kế Giao Diện UI/UX]
\begin{itemize}[leftmargin=1.5em, itemsep=3pt, parsep=0pt]
    \item \textbf{Overlay (Lớp phủ bán trong suốt)}: Phủ một lớp màu đen/tối mờ che phủ ảnh hoặc video phía dưới khi mở cửa sổ Modal/Popup.
    \item \textbf{Hover nhẹ (Hiệu ứng di chuột)}: Tạo màu nền mờ mịn khi người dùng di chuột qua thẻ, nút bấm.
    \item \textbf{Box Shadow (Bóng đổ mịn)}: Khai báo màu bóng đổ có kênh Alpha giúp phần bóng mềm mại, dịu mắt và chuyên nghiệp hơn so với bóng đặc.
    \item \textbf{Gradient (Dải chuyển màu trong suốt)}: Kết hợp RGBA với dải màu chuyển tiếp để tạo hiệu ứng chuyển mờ dần vào ảnh nền.
\end{itemize}
\end{conceptBox}

\begin{codeWindow}[Ví Dụ Ứng Dụng RGBA Làm Lớp Phủ Overlay Modal]
\begin{lstlisting}[language=CSS3]
.overlay {
  background-color: rgba(0, 0, 0, 0.5); /* Nen den mo 50% che phut trang */
  position: absolute;
  top: 0;
  left: 0;
  width: 100%;
  height: 100%;
}
\end{lstlisting}
\end{codeWindow}

\begin{browserWindow}[Kết Quả Mô Phỏng Lớp Phủ Trong Suốt RGBA (Modal Overlay)]
\begin{tcolorbox}[
    enhanced,
    colback=black!65,
    colframe=black!80,
    arc=6pt, boxrule=0.8pt,
    top=14pt, bottom=14pt, left=14pt, right=14pt
]
  \centering
  \begin{tcolorbox}[
      enhanced,
      width=0.88\linewidth,
      colback=white,
      colframe=primaryBlue!40,
      arc=6pt, boxrule=1pt,
      top=10pt, bottom=10pt, left=12pt, right=12pt,
      drop shadow=black!40!white
  ]
    \centering
    {\sffamily\bfseries\large\color{primaryBlue} Cửa Sổ Modal / Popup Dialog}\\[6pt]
    {\sffamily\small\color{darkText} Lớp phủ xung quanh: \codeinline{background-color: rgba(0, 0, 0, 0.5);}}\\[4pt]
    {\sffamily\small\color{gray} Lớp màu tối bán trong suốt RGBA giúp làm nổi bật Modal trung tâm mà vẫn giữ được bối cảnh trang web đằng sau.}
  \end{tcolorbox}
\end{tcolorbox}
\end{browserWindow}

\textbf{8. Ví Dụ Thực Tế Component UI \& Kết Hợp Nâng Cao Của RGBA:}

\begin{codeWindow}[Ví Dụ Thực Tế Thẻ Thông Báo Error Box \& Gradient RGBA]
\begin{lstlisting}[language=HTML5]
<!-- Component HTML -->
<div class="note">Thong bao loi he thong</div>
<div class="gradient-banner">Banner chuyen mau RGBA</div>

<!-- CSS Dinh dang -->
<style>
.note {
  padding: 15px 20px;
  background-color: rgba(255, 0, 0, 0.1); /* Do nhat 10% lam nen */
  border: 1px solid rgb(255, 0, 0);       /* Vien do dac 100% */
  color: rgb(200, 0, 0);                   /* Mau chu do sam de doc */
  border-radius: 6px;
}

/* Ket hop RGBA voi dai mau chuyen tiep */
.gradient-banner {
  padding: 15px 20px;
  background: linear-gradient(to right, rgba(255, 0, 0, 0.5), rgba(0, 0, 255, 0.5));
  color: white;
  border-radius: 6px;
}
</style>
\end{lstlisting}
\end{codeWindow}

\begin{browserWindow}[Kết Quả Hiển Thị Thực Tế Thẻ Error Box \& Banner Gradient]
\vspace{2pt}
% 1. Error Note Box Component
\begin{tcolorbox}[
    enhanced,
    colback=red!8,
    colframe=red!80!black,
    arc=5pt, boxrule=0.8pt,
    top=8pt, bottom=8pt, left=12pt, right=12pt,
    fontupper=\sffamily\color{red!80!black}
]
  \textbf{.note}: Thông báo lỗi hệ thống\\
  {\small\codeinline{background-color: rgba(255, 0, 0, 0.1);} \quad|\quad \codeinline{border: 1px solid rgb(255, 0, 0);}}
\end{tcolorbox}

\vspace{4pt}

% 2. Linear Gradient Banner Component
\begin{tcolorbox}[
    enhanced,
    interior style={left color=red!60, right color=blue!60},
    colframe=primaryBlue!40,
    arc=5pt, boxrule=0.6pt,
    top=8pt, bottom=8pt, left=12pt, right=12pt,
    fontupper=\sffamily\bfseries\color{white}
]
  \centering
  .gradient-banner: Dải Chuyển Màu RGBA (Red 50\% $\rightarrow$ Blue 50\%)\\
  {\small\normalfont\color{white!95}\codeinline{background: linear-gradient(to right, rgba(255, 0, 0, 0.5), rgba(0, 0, 255, 0.5));}}
\end{tcolorbox}
\end{browserWindow}

\textbf{9. So Sánh Chi Tiết Giữa Cú Pháp `rgba()` Và Thuộc Tính `opacity`:}

\begin{prettyTableBox}[Bảng So Sánh Toàn Diện Giữa rgba() Và Thuộc Tính opacity]
\rowcolors{2}{white}{primaryBlue!4}
\sffamily\small
\begin{tabularx}{\linewidth}{|>{\raggedright\arraybackslash \bfseries\color{primaryBlue}}p{2.2cm}|>{\raggedright\arraybackslash}X|>{\raggedright\arraybackslash}X|}
\hline
\rowcolor{primaryBlue} \thd{\sffamily Tiêu chí} & \thd{\sffamily Hàm rgba()} & \thd{\sffamily Thuộc tính opacity} \\
\hline
Phạm vi áp dụng & Chỉ ảnh hưởng duy nhất màu được chỉ định (màu nền, màu chữ hoặc đường viền). & Làm mờ toàn bộ phần tử bao gồm cả tất cả phần tử con bên trong. \\
\hline
Tính linh hoạt & \textcolor{green!60!black}{\textbf{Linh hoạt cao}}: Có thể làm mờ nền mà chữ vẫn đặc $100\%$. & \textcolor{red!80!black}{\textbf{Không chọn lọc}}: Mờ cả văn bản, hình ảnh, icon bên trong. \\
\hline
Ứng dụng Overlay & \textcolor{green!60!black}{\textbf{Hoàn hảo cho Modal/Popup}}: Tạo nền tối mờ nhưng chữ hiển thị sắc nét. & \textcolor{red!80!black}{\textbf{Hạn chế}}: Khiến nội dung trên Modal cũng bị mờ mờ khó đọc. \\
\hline
\end{tabularx}
\end{prettyTableBox}

\textbf{10. Các Lưu Ý Quan Trọng Khi Sử Dụng RGB / RGBA:}

\begin{tipBox}[Lưu Ý Best Practices Khi Dùng RGBA Trong Thiết Kế Web]
\begin{itemize}[leftmargin=1.5em, itemsep=3pt, parsep=0pt]
    \item \textbf{Nên sử dụng \codeinline{rgba()}} khi chỉ muốn màu nền hoặc đường viền mờ nhẹ mà chữ và biểu tượng bên trong vẫn giữ độ rõ $100\%$.
    \item \textbf{Tránh dùng Alpha quá thấp}: Chỉ số Alpha quá nhỏ ($< 0.1$) khiến màu sắc quá nhạt, khó nhìn và dễ vi phạm tiêu chuẩn độ tương phản khả năng truy cập (WCAG Contrast Ratio).
    \item \textbf{Chú ý thứ tự đối số}: Cú pháp \codeinline{rgba(R, G, B, A)} luôn bắt buộc thứ tự 3 kênh màu $R, G, B$ trước, kênh $A$ đứng cuối cùng.
    \item \textbf{Điều chỉnh động bằng JavaScript}: Bạn dễ dàng dùng JS để thay đổi màu sắc bán trong suốt linh hoạt:
    \codeinline{element.style.backgroundColor = `rgba(255, 100, 0, 0.8)`;}
\end{itemize}
\end{tipBox}

\begin{warnBox}[Các Lỗi Thường Gặp Khi Làm Việc Với RGB / RGBA]
\begin{itemize}[leftmargin=1.5em, itemsep=3pt, parsep=0pt]
    \item \textbf{Nhầm lẫn giữa `rgba()` và `opacity`}: Dùng `opacity` cho container khiến chữ bên trong bị mờ nhạt khó đọc. Giải pháp: Thay thế `background-color: rgba(...)`.
    \item \textbf{Nhập sai dải giá trị}: Nhập giá trị kênh màu vượt quá 255 (ví dụ: `rgb(300, 0, 0)`) sẽ bị trình duyệt bỏ qua hoặc tự động ép về 255.
    \item \textbf{Cơ chế thừa kế màu trong DOM}: `rgba()` không bị thừa kế xuống con, nhưng nếu gán `opacity` lên cha thì con bắt buộc bị mờ theo mà không có cách nào bù đắp lại bằng CSS.
\end{itemize}
\end{warnBox}

\textbf{11. Tổng Kết Nhanh So Sánh RGB Và RGBA:}

\begin{prettyTableBox}[Bảng Tổng Kết Nhanh Phân Biệt RGB Và RGBA]
\rowcolors{2}{white}{primaryBlue!4}
\sffamily\small
\begin{tabularx}{\linewidth}{|>{\raggedright\arraybackslash \bfseries\color{primaryBlue}}p{2.2cm}|>{\ttfamily\color{accentOrange}}p{3.2cm}|>{\raggedright\arraybackslash}X|>{\raggedright\arraybackslash}X|}
\hline
\rowcolor{primaryBlue} \thd{\sffamily Thuộc tính} & \thd{\sffamily Cú pháp} & \thd{\sffamily Công dụng chính} & \thd{\sffamily Trường hợp sử dụng phổ biến} \\
\hline
RGB & rgb(R, G, B) & Màu sắc đặc 100\% cơ bản. & Màu chữ chính, đường viền đặc, màu nền cố định. \\
\hline
RGBA & rgba(R, G, B, A) & Màu sắc kết hợp độ mờ trong suốt. & Nền mờ, overlay popup, hiệu ứng hover, bóng đổ mịn. \\
\hline
\end{tabularx}
\end{prettyTableBox}

\textbf{12. Tra Cứu Danh Sách Các Mã Màu RGB Phổ Biến Trong Giao Diện Web:}

\paragraph{I. Nhóm Màu Cơ Bản:}
\begin{prettyTableBox}[Bảng Mã Màu RGB Phổ Biến - Nhóm Màu Cơ Bản]
\rowcolors{2}{white}{primaryBlue!4}
\sffamily\small
\begin{tabularx}{\linewidth}{|>{\raggedright\arraybackslash \bfseries\color{primaryBlue}}p{3cm}|>{\ttfamily\color{accentOrange}}p{3.2cm}|>{\ttfamily}p{2.5cm}|>{\centering\arraybackslash}X|}
\hline
\rowcolor{primaryBlue} \thd{\sffamily Tên màu} & \thd{\sffamily Cú pháp RGB} & \thd{\sffamily Mã Hex} & \thd{\sffamily Minh họa màu} \\
\hline
Red (Đỏ) & rgb(255, 0, 0) & \#FF0000 & \textcolor{red}{\rule{1.2ex}{1.2ex}} Đỏ thuần \\
\hline
Green (Xanh lá) & rgb(0, 128, 0) & \#008000 & \textcolor{green!50!black}{\rule{1.2ex}{1.2ex}} Xanh lá chuẩn \\
\hline
Blue (Xanh dương) & rgb(0, 0, 255) & \#0000FF & \textcolor{blue}{\rule{1.2ex}{1.2ex}} Xanh dương thuần \\
\hline
Yellow (Vàng) & rgb(255, 255, 0) & \#FFFF00 & \textcolor{yellow!80!black}{\rule{1.2ex}{1.2ex}} Vàng rực \\
\hline
Cyan / Aqua & rgb(0, 255, 255) & \#00FFFF & \textcolor{cyan}{\rule{1.2ex}{1.2ex}} Xanh lơ \\
\hline
Magenta / Fuchsia & rgb(255, 0, 255) & \#FF00FF & \textcolor{magenta}{\rule{1.2ex}{1.2ex}} Hồng tím \\
\hline
White (Trắng) & rgb(255, 255, 255) & \#FFFFFF & Trắng tinh \\
\hline
Black (Đen) & rgb(0, 0, 0) & \#000000 & Đen tuyền \\
\hline
\end{tabularx}
\end{prettyTableBox}

\paragraph{II. Màu Sắc Phổ Biến Dùng Trong Thiết Kế Giao Diện (UI):}
\begin{prettyTableBox}[Bảng Mã Màu RGB Phổ Biến - Nhóm Màu Giao Diện UI]
\rowcolors{2}{white}{primaryBlue!4}
\sffamily\small
\begin{tabularx}{\linewidth}{|>{\raggedright\arraybackslash \bfseries\color{primaryBlue}}p{2.8cm}|>{\ttfamily\color{accentOrange}}p{3.2cm}|>{\ttfamily}p{2.2cm}|>{\raggedright\arraybackslash}X|}
\hline
\rowcolor{primaryBlue} \thd{\sffamily Tên màu} & \thd{\sffamily Cú pháp RGB} & \thd{\sffamily Mã Hex} & \thd{\sffamily Gợi ý sử dụng thực tế} \\
\hline
Dark Gray & rgb(64, 64, 64) & \#404040 & Văn bản chính trên nền sáng hoặc văn bản phụ nền tối. \\
\hline
Light Gray & rgb(211, 211, 211) & \#D3D3D3 & Nền phụ, đường phân cách (divider), viền thẻ nhẹ. \\
\hline
Orange & rgb(255, 165, 0) & \#FFA500 & Nút cảnh báo, nút kêu gọi hành động (CTA). \\
\hline
Tomato & rgb(255, 99, 71) & \#FF6347 & Nút bấm nổi bật, trạng thái nhấn mạnh. \\
\hline
Lime & rgb(0, 255, 0) & \#00FF00 & Màu sáng tươi, điểm nhấn nổi bật trên giao diện tối. \\
\hline
Sky Blue & rgb(135, 206, 235) & \#87CEEB & Màu nền dịu nhẹ, banner thông tin mượt mà. \\
\hline
Violet & rgb(238, 130, 238) & \#EE82EE & Màu nhấn accent sáng tạo, thương hiệu hiện đại. \\
\hline
\end{tabularx}
\end{prettyTableBox}

\paragraph{III. Nhóm Màu Pastel (Mềm, Nhẹ, Dễ Nhìn):}
\begin{prettyTableBox}[Bảng Mã Màu RGB Phổ Biến - Nhóm Màu Pastel Dịu Mắt]
\rowcolors{2}{white}{primaryBlue!4}
\sffamily\small
\begin{tabularx}{\linewidth}{|>{\raggedright\arraybackslash \bfseries\color{primaryBlue}}p{2.8cm}|>{\ttfamily\color{accentOrange}}p{3.2cm}|>{\ttfamily}p{2.2cm}|>{\raggedright\arraybackslash}X|}
\hline
\rowcolor{primaryBlue} \thd{\sffamily Tên màu} & \thd{\sffamily Cú pháp RGB} & \thd{\sffamily Mã Hex} & \thd{\sffamily Mô tả \& Phong cách} \\
\hline
Pastel Red & rgb(255, 179, 186) & \#FFB3BA & Đỏ nhạt mềm mại, dùng cho thẻ cảnh báo nhẹ. \\
\hline
Pastel Green & rgb(186, 255, 201) & \#BAFFC9 & Xanh lá nhạt, màu thành công dịu mắt. \\
\hline
Pastel Blue & rgb(179, 217, 255) & \#B3D9FF & Xanh dương nhạt, dùng làm nền card hoặc tag. \\
\hline
Light Pink & rgb(255, 182, 193) & \#FFB6C1 & Hồng nhẹ dễ chịu cho giao diện thẩm mỹ. \\
\hline
Peach & rgb(255, 229, 180) & \#FFE5B4 & Màu đào nhạt, rất hợp làm màu nền section. \\
\hline
Mint & rgb(189, 252, 201) & \#BDFCC9 & Màu bạc hà hiện đại, giao diện tươi sáng. \\
\hline
\end{tabularx}
\end{prettyTableBox}

\paragraph{IV. Nhóm Màu Dành Cho Nền Tối (Dark Theme):}
\begin{prettyTableBox}[Bảng Mã Màu RGB Phổ Biến - Nhóm Màu Dark Theme]
\rowcolors{2}{white}{primaryBlue!4}
\sffamily\small
\begin{tabularx}{\linewidth}{|>{\raggedright\arraybackslash \bfseries\color{primaryBlue}}p{2.8cm}|>{\ttfamily\color{accentOrange}}p{3.2cm}|>{\ttfamily}p{2.2cm}|>{\raggedright\arraybackslash}X|}
\hline
\rowcolor{primaryBlue} \thd{\sffamily Tên màu} & \thd{\sffamily Cú pháp RGB} & \thd{\sffamily Mã Hex} & \thd{\sffamily Gợi ý sử dụng Dark Theme} \\
\hline
Dark Blue & rgb(18, 32, 47) & \#12202F & Nền chính cho giao diện Dark Mode sang trọng. \\
\hline
Gray-900 & rgb(33, 37, 41) & \#212529 & Nền card hoặc thanh điều hướng (Navbar) tối. \\
\hline
Slate & rgb(108, 117, 125) & \#6C757D & Văn bản phụ, icon xám trung tính trên nền tối. \\
\hline
Teal & rgb(0, 128, 128) & \#008080 & Màu nhấn Accent xanh ngọc hiện đại. \\
\hline
Charcoal & rgb(54, 69, 79) & \#36454F & Nền trung tính cao cấp, hạn chế mỏi mắt. \\
\hline
\end{tabularx}
\end{prettyTableBox}

\paragraph{V. Thuật Toán Tự Chuyển Đổi Mã Hex $\rightarrow$ RGB:}
\begin{conceptBox}[Quy Trình Chuyển Đổi Từ Mã Hex Sang Cú Pháp RGB]
Ví dụ: Chuyển đổi mã Hex \codeinline{\#FF5733} sang cú pháp \codeinline{rgb(R, G, B)}.

\textbf{Các bước thực hiện}:
\begin{enumerate}[leftmargin=1.5em, itemsep=2pt]
    \item Tách mã Hex 6 ký tự thành 3 cặp ký tự đại diện cho 3 kênh:
    $\text{Red} = \text{FF}_{16}$, $\text{Green} = 57_{16}$, $\text{Blue} = 33_{16}$.
    \item Chuyển đổi từng cặp số Thập lục phân (Hex - cơ số 16) sang Thập phân (Decimal - cơ số 10):
    \begin{itemize}[leftmargin=1.2em]
        \item Kênh Red ($\text{FF}_{16}$): $15 \times 16^1 + 15 \times 16^0 = 240 + 15 = \mathbf{255}$.
        \item Kênh Green ($57_{16}$): $5 \times 16^1 + 7 \times 16^0 = 80 + 7 = \mathbf{87}$.
        \item Kênh Blue ($33_{16}$): $3 \times 16^1 + 3 \times 16^0 = 48 + 3 = \mathbf{51}$.
    \end{itemize}
    \item Ghép lại vào cú pháp chuẩn CSS: $\rightarrow$ \codeinline{rgb(255, 87, 51)}.
\end{enumerate}
\end{conceptBox}

\begin{conceptBox}[Câu Hỏi Phỏng Vấn Thường Gặp Về Hệ Màu Trong CSS]
\begin{itemize}[leftmargin=1.5em, itemsep=4pt, parsep=0pt]
    \item \textbf{Câu hỏi 1}: Sự khác biệt cốt lõi giữa `background: rgba(0,0,0,0.5)` và `background: black; opacity: 0.5` là gì?
    \begin{itemize}[leftmargin=1.2em]
        \item \textit{Trả lời}: `rgba()` chỉ áp dụng độ mờ 50\% duy nhất lên màu nền, giữ cho chữ và icon bên trong rõ đặc 100\%. Ngược lại, `opacity: 0.5` làm mờ toàn bộ phần tử và tất cả phần tử con bên trong.
    \end{itemize}
    \item \textbf{Câu hỏi 2}: Khi nào nên dùng mã Hex và khi nào nên dùng `rgb()` / `rgba()`?
    \begin{itemize}[leftmargin=1.2em]
        \item \textit{Trả lời}: Mã Hex phù hợp nhất cho màu cố định 100\% vì ngắn gọn và chuẩn Figma. `rgba()` bắt buộc phải dùng khi cần màu nền mờ trong suốt, đổ bóng mịn hoặc thao tác đổi màu động bằng JavaScript.
    \end{itemize}
\end{itemize}
\end{conceptBox}

\begin{prettyTableBox}[Đọc Thêm: Mở Rộng Kiến Thức Chuyên Sâu Về RGB / RGBA]
\textbf{1. Cú Pháp Hiện Đại CSS Color Module Level 4 (Space-separated Syntax):}
Trong các tiêu chuẩn CSS hiện đại (CSS Color Module Level 4), W3C hỗ trợ cú pháp không sử dụng dấu phẩy và dùng dấu gạch chéo \codeinline{/} cho kênh Alpha:
\begin{codeWindow}[Cú Pháp CSS Color Level 4 Không Dấu Phẩy]
\begin{lstlisting}[language=CSS3]
/* Tuong duong rgba(255, 87, 51, 0.5) */
.modern-box {
  color: rgb(255 87 51 / 50%);
  background: rgb(0 0 0 / 0.7);
}
\end{lstlisting}
\end{codeWindow}

\textbf{2. CSS Color Module Level 5 \& Hàm `color-mix()`:}
CSS Color Level 5 giới thiệu hàm hòa trộn màu mạnh mẽ trực tiếp bằng CSS thuần:
\begin{codeWindow}[Ví Dụ Hòa Trộn Màu Với color-mix()]
\begin{lstlisting}[language=CSS3]
/* Hoa tron 80% mau do voi 20% mau trang */
.soft-red-mix {
  background-color: color-mix(in srgb, red 80%, white 20%);
}
\end{lstlisting}
\end{codeWindow}

\textbf{3. Nguyên Lý Xử Lý Alpha Blending Trên Phần Cứng Card Đồ Họa (GPU):}
Khi một phần tử có màu RGBA bán trong suốt đè lên phần tử khác, trình duyệt sẽ thực hiện thuật toán tính toán hòa trộn điểm ảnh (\term{Alpha Blending}):
\[
C_{\text{final}} = C_{\text{src}} \times A_{\text{src}} + C_{\text{dst}} \times (1 - A_{\text{src}})
\]
Quá trình này được tăng tốc phần cứng (\term{GPU Hardware Acceleration}) thông qua việc tạo ra một lớp tổng hợp (\term{Compositor Layer}) riêng biệt trên card đồ họa.

\textbf{4. Lưu Ý Khả Năng Truy Cập (Accessibility \& WCAG Contrast):}
Khi thiết kế chữ mờ hoặc background bán trong suốt bằng RGBA, cần đảm bảo tỷ lệ độ tương phản (\term{Contrast Ratio}) đạt tối thiểu $4.5:1$ đối với văn bản thường và $3.0:1$ đối với văn bản lớn theo tiêu chuẩn WCAG 2.1 AA.
\end{prettyTableBox}

\paragraph{3.4. Hệ Màu HSL Và HSLA (Hue, Saturation, Lightness, Alpha) Masterclass:}

\begin{conceptBox}[Khái Niệm Cơ Bản Về Mô Hình Màu HSL \& HSLA]
\textbf{1. Bản Chất Kỹ Thuật Của Hệ Màu HSL}:
\begin{itemize}[leftmargin=1.5em, itemsep=3pt, parsep=0pt]
    \item \textbf{HSL (Hue - Saturation - Lightness)}: Là hệ thống biểu diễn màu sắc dựa trên góc màu và cách mắt người cảm nhận trực quan trong thực tế (thay vì dựa trên cường độ phát sáng máy tính như RGB/HEX).
    \item \textbf{4 Kênh thông số của HSL / HSLA}:
    \begin{itemize}[leftmargin=1.2em, itemsep=2pt]
        \item \textbf{Hue (Góc màu - $H$)}: Góc xoay trên vòng tròn màu sắc từ $0^\circ$ đến $360^\circ$. Trong đó: $0^\circ$ (hoặc $360^\circ$) là Đỏ, $60^\circ$ là Vàng, $120^\circ$ là Xanh lá, $180^\circ$ là Xanh lơ, $240^\circ$ là Xanh dương, $300^\circ$ là Hồng tím.
        \item \textbf{Saturation (Độ bão hòa - $S\%$)}: Mức độ tươi rực rỡ của màu ($0\%$ là màu xám xịt trung tính, $100\%$ là sắc màu tươi nguyên bản).
        \item \textbf{Lightness (Độ sáng - $L\%$)}: Mức độ sáng/tối của màu ($0\%$ là màu đen tuyền, $50\%$ là tông màu chuẩn nguyên bản, $100\%$ là màu trắng tinh).
        \item \textbf{Alpha (Độ trong suốt - $A$)}: Chỉ số độ mờ nhìn xuyên thấu từ $0.0$ (trong suốt $0\%$) đến $1.0$ (đặc nguyên bản $100\%$).
    \end{itemize}
\end{itemize}
\end{conceptBox}

\textbf{2. Cú Pháp Định Dạng Chuẩn Trong CSS \& Bảng Giá Trị:}

\begin{codeWindow}[Cú Pháp Khai Báo HSL \& HSLA Trong CSS]
\begin{lstlisting}[language=CSS3]
/* Cu phap HSL truyen thong */
color: hsl(H, S%, L%);

/* Cu phap HSLA co kenh Alpha */
color: hsla(H, S%, L%, A);
\end{lstlisting}
\end{codeWindow}

\textbf{3. Các Ví Dụ Mã Nguồn Thực Tế Về Hệ Màu HSL \& HSLA:}

\begin{codeWindow}[Ví Dụ 1: Tạo Dải Màu Đồng Bộ (Color Palette) Chỉ Bằng Việc Đổi Chỉ Số Lightness L\%]
\begin{lstlisting}[language=CSS3]
/* Giu nguyen goc Hue=217 (Xanh duong) va Saturation=91% */
.bg-lightest  { background-color: hsl(217, 91%, 95%); color: #1e293b; } /* Nen xanh cuc nhat */
.btn-primary  { background-color: hsl(217, 91%, 60%); color: white; }   /* Mau nut nhan */
.btn-hover    { background-color: hsl(217, 91%, 50%); color: white; }   /* Mau nut hover dam hon */
.btn-active   { background-color: hsl(217, 91%, 40%); color: white; }   /* Mau nut active dam nhat */
\end{lstlisting}
\end{codeWindow}

\begin{codeWindow}[Ví Dụ 2: Định Kiểu Thẻ Alert Trong Suốt Nhẹ Bằng HSLA]
\begin{lstlisting}[language=CSS3]
/* Mau tim trong suot mo 15% bang HSLA */
.alert-purple {
  background-color: hsla(280, 80%, 60%, 0.15); /* Alpha = 0.15 */
  border: 1px solid hsl(280, 80%, 50%);
  color: hsl(280, 80%, 30%);
  padding: 12px 16px;
  border-radius: 6px;
}
\end{lstlisting}
\end{codeWindow}

\begin{browserWindow}[Kết Quả Hiển Thị Của Dải Màu HSL Palette \& Thẻ Alert HSLA]
\vspace{2pt}
\begin{tcolorbox}[
    enhanced,
    colback=white,
    colframe=gray!30,
    arc=6pt, boxrule=0.8pt,
    top=10pt, bottom=10pt, left=12pt, right=12pt
]
  % HSL Palette Preview
  {\sffamily\bfseries\color{primaryBlue} Dải Màu Đồng Bộ Tạo Từ HSL (Xoay quanh Góc Hue 217$^\circ$)}:\\[4pt]
  \begin{tcolorbox}[nobeforeafter, colback=blue!8, colframe=blue!20, arc=3pt, boxrule=0.4pt, top=4pt, bottom=4pt, left=6pt, right=6pt, fontupper=\sffamily\bfseries\scriptsize\color{black}]
    hsl(217, 91\%, 95\%) - Nền nhạt
  \end{tcolorbox}
  \quad
  \begin{tcolorbox}[nobeforeafter, colback=blue!65!white, colframe=blue!65!white, arc=3pt, boxrule=0.4pt, top=4pt, bottom=4pt, left=6pt, right=6pt, fontupper=\sffamily\bfseries\scriptsize\color{white}]
    hsl(217, 91\%, 60\%) - Nút chính
  \end{tcolorbox}
  \quad
  \begin{tcolorbox}[nobeforeafter, colback=blue!85!black, colframe=blue!85!black, arc=3pt, boxrule=0.4pt, top=4pt, bottom=4pt, left=6pt, right=6pt, fontupper=\sffamily\bfseries\scriptsize\color{white}]
    hsl(217, 91\%, 40\%) - Active
  \end{tcolorbox}
  \vspace{6pt}
  
  % HSLA Alert Card Preview
  \begin{tcolorbox}[
      enhanced,
      colback=violet!12,
      colframe=violet!60,
      arc=4pt, boxrule=0.8pt,
      top=6pt, bottom=6pt, left=10pt, right=10pt,
      fontupper=\sffamily\bfseries\color{violet!80!black}
  ]
    .alert-purple: Thẻ Alert Tím nhẹ bằng HSLA \codeinline{hsla(280, 80\%, 60\%, 0.15);}
  \end{tcolorbox}
\end{tcolorbox}
\end{browserWindow}

\textbf{4. Lợi Thế Vượt Trội Của HSL Trong Lập Trình Hệ Màu Động (Theme System):}

\begin{prettyTableBox}[Bảng Phân Tích Lợi Thế Vượt Trội Của Hệ Màu HSL]
\rowcolors{2}{white}{primaryBlue!4}
\sffamily\small
\begin{tabularx}{\linewidth}{|>{\raggedright\arraybackslash \bfseries\color{primaryBlue}}p{3.5cm}|>{\raggedright\arraybackslash}X|}
\hline
\rowcolor{primaryBlue} \thd{\sffamily Ưu điểm cốt lõi} & \thd{\sffamily Giải thích \& Ứng dụng thực tế} \\
\hline
Tạo Palette màu đồng bộ & Chỉ cần giữ nguyên góc Hue $H$, thay đổi chỉ số $L\%$ để tạo ra dải sắc thái từ nhạt đến đậm nhất quán. \\
\hline
Tối ưu hiệu ứng Hover / Focus & Dễ dàng tạo trạng thái di chuột nhẹ hơn hoặc đậm hơn mà không cần đoán mã Hex ngẫu nhiên. \\
\hline
Hợp tác tốt với Biến CSS & Khai báo góc màu trong biến: \codeinline{--brand-hue: 217; background: hsl(var(--brand-hue), 90\%, 50\%);}. \\
\hline
\end{tabularx}
\end{prettyTableBox}

\paragraph{3.5. Từ Khóa Đặc Biệt: transparent Và currentColor Masterclass:}

\begin{conceptBox}[Tổng Quan Về Hai Từ Khóa Đặc Biệt Trong CSS]
Trong hệ thống màu sắc của CSS, ngoài các mô hình mã màu số (HEX, RGB, HSL), W3C cung cấp hai từ khóa đặc biệt cực kỳ linh hoạt và có sức mạnh lớn trong thiết kế giao diện UI/UX hiện đại:
\begin{enumerate}[leftmargin=1.5em, itemsep=3pt, parsep=0pt]
    \item \textbf{\codeinline{transparent}}: Từ khóa đại diện cho màu sắc hoàn toàn trong suốt ($100\%$ transparency / $0\%$ opacity).
    \item \textbf{\codeinline{currentColor}}: Từ khóa động đại diện cho giá trị màu sắc của thuộc tính \codeinline{color} hiện tại.
\end{enumerate}
\end{conceptBox}

% ==========================================
% PHẦN I: TỪ KHÓA TRANSPARENT
% ==========================================
\subsubsection*{I. Chuyên Đề Từ Khóa Đặc Biệt \codeinline{transparent}:}

\begin{conceptBox}[Bản Chất Kỹ Thuật Của Từ Khóa transparent]
\textbf{1. transparent Là Gì?}:
\begin{itemize}[leftmargin=1.5em, itemsep=3pt, parsep=0pt]
    \item Từ khóa \codeinline{transparent} tạo ra một màu hoàn toàn nhìn xuyên thấu ($100\%$ trong suốt), cho phép các phần tử nền hoặc nội dung phía dưới hiển thị trọn vẹn.
    \item \textbf{Bản chất toán học trong hệ màu}: Từ khóa \codeinline{transparent} tương đương hoàn toàn với giá trị \codeinline{rgba(0, 0, 0, 0)} hoặc \codeinline{hsla(0, 0\%, 0\%, 0)}.
    \item Trong tiêu chuẩn CSS Color Module Level 3 trở đi, \codeinline{transparent} được coi là một tên màu chuẩn hợp lệ và có thể áp dụng cho bất kỳ thuộc tính màu sắc nào (\codeinline{color}, \codeinline{background-color}, \codeinline{border-color}, \codeinline{shadow}, \codeinline{stroke}).
\end{itemize}
\end{conceptBox}

\textbf{2. Các Ứng Dụng Kỹ Thuật Kinh Điển Của \codeinline{transparent} Trong Thiết Kế Web:}

\begin{prettyTableBox}[Bảng Tổng Hợp 4 Ứng Dụng Kỹ Thuật Thực Tế Của transparent]
\rowcolors{2}{white}{primaryBlue!4}
\sffamily\small
\begin{tabularx}{\linewidth}{|>{\raggedright\arraybackslash \bfseries\color{primaryBlue}}p{3.8cm}|>{\raggedright\arraybackslash}X|}
\hline
\rowcolor{primaryBlue} \thd{\sffamily Ứng dụng kỹ thuật} & \thd{\sffamily Mô tả \& Kỹ thuật triển khai trong CSS} \\
\hline
Vẽ hình khối tam giác / Tooltip & Kết hợp các đường viền \codeinline{border} với $3$ cạnh mang màu \codeinline{transparent} và $1$ cạnh có màu thực tế để tạo mũi tên Tooltip. \\
\hline
Nút bấm viền mỏng (Ghost Button) & Khai báo \codeinline{background-color: transparent;} để lộ nền trang web, chỉ giữ đường viền \codeinline{border} tinh tế. \\
\hline
Phủ dải mờ Gradient đè ảnh nền & Sử dụng \codeinline{linear-gradient(to top, \#1e293b, transparent)} giúp văn bản hiển thị rõ ràng trên nền hình ảnh complex. \\
\hline
Reset nền mặc định ô nhập liệu & Tháo bỏ màu nền xám mặc định của trình duyệt cho thẻ \codeinline{<input>}, \codeinline{<select>} hoặc đường ray thanh cuộn \codeinline{::-webkit-scrollbar-track}. \\
\hline
\end{tabularx}
\end{prettyTableBox}

\begin{codeWindow}[Ví Dụ 1: Định Kiểu Nút Bấm Ghost Button \& Mũi Tên Tooltip Bằng transparent]
\begin{lstlisting}[language=CSS3]
/* (*@ 1. Nút bấm Ghost Button sang trọng @*) */
.btn-ghost {
  background-color: transparent; /* (*@ Nền trong suốt nhìn xuyên qua @*) */
  border: 2px solid #3b82f6;
  color: #3b82f6;
  padding: 8px 16px;
  border-radius: 6px;
  transition: all 0.3s ease;
}

.btn-ghost:hover {
  background-color: #3b82f6;     /* Khi hover moi to mau dac */
  color: #ffffff;
}

/* 2. Mui ten tam giac Tooltip (Triangle Arrow) */
.tooltip-arrow {
  width: 0;
  height: 0;
  border-left: 8px solid transparent;  /* Vien trai trong suot */
  border-right: 8px solid transparent; /* Vien phai trong suot */
  border-bottom: 10px solid #1e293b;   /* Vien duoi tao thanh dinh tam giac */
}

/* 3. Gradient mo dan vao anh nen (Fading Gradient Overlay) */
.hero-card {
  background: linear-gradient(to top, rgba(15, 23, 42, 0.9), transparent);
}
\end{lstlisting}
\end{codeWindow}

\begin{browserWindow}[Kết Quả Hiển Thị Của Nút Bấm Ghost Button \& Mũi Tên Tooltip]
\vspace{2pt}
\begin{tcolorbox}[
    enhanced,
    colback=gray!5,
    colframe=gray!30,
    arc=6pt, boxrule=0.8pt,
    top=10pt, bottom=10pt, left=12pt, right=12pt
]
  % Ghost Button Preview
  {\sffamily\bfseries\color{primaryBlue} Nút Bấm Ghost Button (\codeinline{background: transparent;}):}\\[4pt]
  \begin{tcolorbox}[nobeforeafter, colback=gray!5, colframe=blue!70!black, arc=4pt, boxrule=1.2pt, top=4pt, bottom=4pt, left=10pt, right=10pt, fontupper=\sffamily\bfseries\small\color{blue!70!black}]
    Khám Phá Khóa Học (Ghost Button)
  \end{tcolorbox}
  \vspace{6pt}
  
  % Tooltip Triangle Preview
  {\sffamily\bfseries\color{primaryBlue} Mũi Tên Tooltip Tam Giác (Kết hợp \codeinline{transparent}):}\\[4pt]
  \begin{tcolorbox}[nobeforeafter, colback=gray!90!black, colframe=gray!90!black, arc=3pt, boxrule=0.4pt, top=4pt, bottom=4pt, left=10pt, right=10pt, fontupper=\sffamily\bfseries\scriptsize\color{white}]
    ▲ \quad Thẻ Tooltip Thông Tin (\codeinline{border-left/right: transparent;})
  \end{tcolorbox}
\end{tcolorbox}
\end{browserWindow}

\begin{prettyTableBox}[Bảng So Sánh Chi Tiết: transparent vs opacity: 0 vs visibility: hidden]
\rowcolors{2}{white}{primaryBlue!4}
\sffamily\small
\begin{tabularx}{\linewidth}{|>{\raggedright\arraybackslash \ttfamily\color{primaryBlue}\bfseries}p{3.2cm}|>{\raggedright\arraybackslash}p{2.8cm}|>{\raggedright\arraybackslash}p{3.0cm}|>{\raggedright\arraybackslash}X|}
\hline
\rowcolor{primaryBlue} \thd{\sffamily Tiêu chí} & \thd{\sffamily transparent} & \thd{\sffamily opacity: 0} & \thd{\sffamily visibility: hidden} \\
\hline
Bản chất tác động & Chỉ làm mờ màu sắc nền/viền được gán. & Làm mờ toàn bộ phần tử và tất cả con bên trong. & Ẩn hiển thị phần tử nhưng vẫn giữ nguyên khoảng trống DOM. \\
\hline
Sự tương tác chuột (Click/Hover) & \textcolor{green!60!black}{\textbf{Vẫn tương tác tốt}} (Pointer events hoạt động). & \textcolor{green!60!black}{\textbf{Vẫn nhận click}} (Rất dễ gây bẫy bấm nhầm). & \textcolor{red!80!black}{\textbf{Vô hiệu hóa click}} (Không nhận sự kiện chuột). \\
\hline
Khả năng kế thừa phần tử con & Phần tử con giữ màu sắc riêng biệt $100\%$. & Tất cả phần tử con bị mờ theo $100\%$. & Phần tử con có thể hiện lại bằng \codeinline{visibility: visible}. \\
\hline
\end{tabularx}
\end{prettyTableBox}

% ==========================================
% PHẦN II: TỪ KHÓA CURRENTCOLOR
% ==========================================
\subsubsection*{II. Chuyên Đề Từ Khóa Đặc Biệt \codeinline{currentColor}:}

\begin{conceptBox}[Bản Chất Kỹ Thuật Của Từ Khóa Đặc Biệt currentColor]
\textbf{1. currentColor Là Gì?}:
\begin{itemize}[leftmargin=1.5em, itemsep=3pt, parsep=0pt]
    \item \textbf{\codeinline{currentColor}}: Là một giá trị từ khóa động trong CSS, đại diện cho giá trị màu sắc hiện tại của thuộc tính \codeinline{color} đang áp dụng trực tiếp hoặc kế thừa lên phần tử đó.
    \item \textbf{Nguyên lý cốt lõi}: $\mathbf{\text{currentColor} \equiv \text{giá trị hiện tại của thuộc tính color}}$.
    \item \textbf{Công dụng chính}: Đồng bộ tự động màu sắc cho \codeinline{border-color}, \codeinline{background-color}, \codeinline{box-shadow}, \codeinline{outline-color}, và đồ họa SVG \codeinline{fill} / \codeinline{stroke}.
\end{itemize}
\end{conceptBox}

\begin{codeWindow}[Ví Dụ 2: Hiệu Ứng Hover Nút Bấm \& Đồng Bộ Icon SVG Với currentColor]
\begin{lstlisting}[language=CSS3]
/* 1. Nut bam dong bo mau vien khi Hover */
button.btn-sync {
  color: #2563eb;                /* Mau chu xanh duong */
  border: 2px solid currentColor; /* Vien tu dong xanh duong */
  background-color: transparent;
  padding: 8px 16px;
  border-radius: 6px;
  transition: color 0.3s ease;
}

button.btn-sync:hover {
  color: #d97706;                /* Chi can doi color, vien tu dong sang mau cam! */
}

/* 2. Dong bo mau khung nhap lieu va Icon SVG khi Focus */
.input-group {
  color: #64748b;                /* Mau icon va text phu xam */
  border: 2px solid currentColor; /* Vien xam theo color */
  padding: 8px 12px;
  border-radius: 6px;
}

.input-group:focus-within {
  color: #2563eb;                /* Ca vien va Icon SVG tu dong sang xanh duong! */
}
\end{lstlisting}
\end{codeWindow}

\begin{browserWindow}[Kết Quả Hiển Thị Đồng Bộ Trạng Thái Focus \& Hover Bằng currentColor]
\vspace{2pt}
\begin{tcolorbox}[
    enhanced,
    colback=white,
    colframe=gray!30,
    arc=6pt, boxrule=0.8pt,
    top=10pt, bottom=10pt, left=12pt, right=12pt
]
  % Input Focus Sync Preview
  {\sffamily\bfseries\color{primaryBlue} Trạng Thái Focus Input Đồng Bộ Icon SVG (\codeinline{currentColor}):}\\[4pt]
  \begin{tcolorbox}[
      enhanced,
      colback=blue!5,
      colframe=blue!70!black,
      arc=4pt, boxrule=1.2pt,
      top=6pt, bottom=6pt, left=10pt, right=10pt,
      fontupper=\sffamily\bfseries\color{blue!80!black}
  ]
    [Icon SVG] \quad Nhập email... \quad (\codeinline{border} và Icon tự động chuyển màu Xanh dương khi Focus)
  \end{tcolorbox}
\end{tcolorbox}
\end{browserWindow}

\begin{prettyTableBox}[Tóm Tắt Tổng Kết Nhanh So Sánh transparent Và currentColor]
\rowcolors{2}{white}{primaryBlue!4}
\sffamily\small
\begin{tabularx}{\linewidth}{|>{\raggedright\arraybackslash \ttfamily\color{primaryBlue}\bfseries}p{3.0cm}|>{\bfseries\color{accentOrange}}p{3.0cm}|>{\raggedright\arraybackslash}X|}
\hline
\rowcolor{primaryBlue} \thd{\sffamily Từ khóa} & \thd{\sffamily Bản chất biểu diễn} & \thd{\sffamily Trường hợp ứng dụng hàng đầu (Use Cases)} \\
\hline
transparent & Màu trong suốt 100\% (Alpha = 0). & Ghost Button, Mũi tên tam giác Tooltip, Gradient mờ dần vào ảnh nền, Reset nền mặc định. \\
\hline
currentColor & Giá trị màu động đại diện cho \codeinline{color}. & Đồng bộ màu viền/bóng đổ với chữ, icon SVG tự đổi màu theo theme, hiệu ứng Hover/Focus linh hoạt. \\
\hline
\end{tabularx}
\end{prettyTableBox}

\textbf{5. Các Lưu Ý Kỹ Thuật Quan Trọng Khi Sử Dụng currentColor:}

\begin{warnBox}[Các Cảnh Báo \& Bẫy Thường Gặp Với currentColor]
\begin{itemize}[leftmargin=1.5em, itemsep=3pt, parsep=0pt]
    \item \textbf{Phải có thuộc tính \codeinline{color} xác định}: Nếu phần tử không gán \codeinline{color} và không kế thừa từ cha, trình duyệt sẽ tự động lấy giá trị mặc định của trình duyệt (\codeinline{initial}).
    \item \textbf{Không dùng thay thế thuộc tính \codeinline{color}}: Tuyệt đối không khai báo \codeinline{color: currentColor;} trên cùng phần tử vì sẽ gây ra vòng lặp tham chiếu vô nghĩa.
    \item \textbf{Hạn chế trong dải dốc màu (\codeinline{linear-gradient()})}: Một số phiên bản trình duyệt cũ không hỗ trợ tốt việc dùng \codeinline{currentColor} làm điểm mốc màu (\term{Color Stop}) trong gradient.
\end{itemize}
\end{warnBox}

\textbf{6. Tóm Tắt Nhanh Bản Chất Của Từ Khóa currentColor:}

\begin{prettyTableBox}[Bảng Tổng Kết Nhanh Đặc Tính Của currentColor]
\rowcolors{2}{white}{primaryBlue!4}
\sffamily\small
\begin{tabularx}{\linewidth}{|>{\raggedright\arraybackslash \bfseries\color{primaryBlue}}p{3.0cm}|>{\raggedright\arraybackslash}X|}
\hline
\rowcolor{primaryBlue} \thd{\sffamily Tính năng} & \thd{\sffamily Giá trị \& Mô tả} \\
\hline
Loại giá trị & Giá trị từ khóa đặc biệt trong tiêu chuẩn CSS. \\
\hline
Đại diện cho & Giá trị màu chữ (\codeinline{color}) hiện tại của phần tử. \\
\hline
Ứng dụng phổ biến & \codeinline{border-color}, \codeinline{outline-color}, \codeinline{box-shadow}, SVG \codeinline{fill}/\codeinline{stroke}, màu hiệu ứng \codeinline{:hover}. \\
\hline
Tính kế thừa & Tự động kế thừa theo màu chữ của phần tử cha trong cây DOM. \\
\hline
Gợi ý sử dụng Best Practice & Tạo nút bấm (Button), biểu tượng Icon SVG, đường viền nhất quán màu sắc. Khi chuyển đổi theme chỉ cần đổi duy nhất 1 biến \codeinline{color}. \\
\hline
\end{tabularx}
\end{prettyTableBox}

\paragraph{3.7. Mẹo Chọn Màu Sắc Chuyên Nghiệp Cho Giao Diện Web (UI/UX Best Practices):}

\begin{prettyTableBox}[Bảng Gợi Ý Phối Màu Thực Tế Chuẩn Chuyên Nghiệp]
\rowcolors{2}{white}{primaryBlue!4}
\sffamily\small
\begin{tabularx}{\linewidth}{|>{\raggedright\arraybackslash \bfseries\color{primaryBlue}}p{4.2cm}|>{\ttfamily\color{accentOrange}\bfseries}p{4.2cm}|>{\raggedright\arraybackslash}X|}
\hline
\rowcolor{primaryBlue} \thd{\sffamily Mục đích sử dụng} & \thd{\sffamily Mã màu gợi ý} & \thd{\sffamily Nguyên tắc \& Lý do chọn} \\
\hline
Văn bản chính (Primary text) & \#212529, \#333333 & Dùng xám đậm thay vì đen tuyền \codeinline{\#000} để dịu mắt. \\
\hline
Văn bản phụ (Secondary text) & \#6c757d, \#888888 & Xám nhạt hơn để tạo tương phản phân cấp thông tin. \\
\hline
Nền chính (Main Background) & \#ffffff, \#f8f9fa & Trắng hoặc xám kem rất nhẹ tạo cảm giác sạch sẽ. \\
\hline
Nút gọi hành động (CTA Button) & \#007bff, \#ff6347 & Màu xanh/cam tươi nổi bật kích thích hành động nhấp. \\
\hline
Cảnh báo lỗi (Danger / Error) & \#dc3545, \#ff4d4f & Đỏ cảnh báo nổi bật đạt chuẩn tương phản W3C. \\
\hline
Giao diện tối (Dark Mode) & \#121212, \#1a1a1a & Nền tối kết hợp văn bản xám sáng (\codeinline{\#e0e0e0}). \\
\hline
\end{tabularx}
\end{prettyTableBox}

\paragraph{3.8. Ứng Dụng Màu Động Với CSS Variables \& JavaScript:}

Trong phát triển ứng dụng Web thực tế, việc quản lý hệ thống màu sắc bằng Biến CSS (\term{CSS Variables}) và điều khiển linh hoạt qua JavaScript giúp hỗ trợ tính năng chuyển đổi giao diện sáng/tối (\term{Dark Mode / Theme Switch}) cực kỳ dễ dàng.

\begin{codeWindow}[Khai Báo Hệ Thống Biến Màu Trong CSS]
\begin{lstlisting}[language=CSS3]
:root {
  --primary-color: #3498db;
  --text-color: #333333;
  --bg-color: #ffffff;
}

/* Ap dung bien mau cho giao dien */
body {
  background-color: var(--bg-color);
  color: var(--text-color);
}

.button-main {
  background-color: var(--primary-color);
  color: #ffffff;
}
\end{lstlisting}
\end{codeWindow}

\begin{codeWindow}[Cập Nhật Màu Động Bằng JavaScript]
\begin{lstlisting}[language=JavaScript]
// Thao tac thay doi mau nen bang JavaScript
const heroElement = document.querySelector('.hero-section');
heroElement.style.backgroundColor = 'rgba(255, 0, 0, 0.6)';

// Thay doi gia tri bien CSS truc tiep qua JavaScript
document.documentElement.style.setProperty('--primary-color', '#2ecc71');
\end{lstlisting}
\end{codeWindow}

\paragraph{3.9. Bảng So Sánh Tổng Quát Các Hệ Màu Trong CSS:}

\begin{prettyTableBox}[Bảng So Sánh Toàn Diện Các Hệ Biểu Diễn Màu]
\rowcolors{2}{white}{primaryBlue!4}
\sffamily\small
\begin{tabularx}{\linewidth}{|>{\raggedright\arraybackslash \bfseries\color{primaryBlue}\ttfamily}p{2.4cm}|c|c|c|Z|}
\hline
\rowcolor{primaryBlue} \thd{\sffamily Hệ màu} & \thd{\sffamily Độ phổ biến} & \thd{\sffamily Dễ viết tay?} & \thd{\sffamily Hỗ trợ Alpha?} & \thd{\sffamily Ứng dụng phù hợp nhất} \\
\hline
Color Name & Thấp & \textcolor{green!60!black}{\textbf{Rất dễ}} & \textcolor{red}{\textbf{Không}} & Học tập, ví dụ nhanh, chạy thử nghiệm code. \\
\hline
Hexadecimal & \textcolor{green!60!black}{\textbf{Cao nhất}} & \textcolor{orange!90!black}{\textbf{Khá dễ}} & \textcolor{red}{\textbf{Không}} & Dự án web thực tế, sao chép mã từ Figma/Photoshop. \\
\hline
RGB / RGBA & \textcolor{green!60!black}{\textbf{Rất cao}} & Dài hơn & \textcolor{green!60!black}{\textbf{Có (RGBA)}} & Hiệu ứng mờ overlay, đổ bóng, di chuột hover. \\
\hline
HSL / HSLA & Trung bình & Khó ban đầu & \textcolor{green!60!black}{\textbf{Có (HSLA)}} & Lập trình hệ màu động, theme sáng/tối, dải màu sắc. \\
\hline
CSS Variables & \textcolor{green!60!black}{\textbf{Rất cao}} & Chuẩn mực & \textcolor{green!60!black}{\textbf{Tùy thuộc biến}} & Xây dựng hệ thiết kế Design System chuyên nghiệp. \\
\hline
\end{tabularx}
\end{prettyTableBox}

\paragraph{3.10. Kết Luận – Khi Nào Nên Dùng Hệ Màu Nào?}

\begin{conceptBox}[Quy Tắc Vàng Lựa Chọn Hệ Màu Cho Dự Án]
\begin{enumerate}[leftmargin=1.5em, itemsep=4pt, parsep=0pt]
    \item \textbf{Dự án nhỏ, làm bài tập nhanh}: Sử dụng \codeinline{Color Name} hoặc \codeinline{Hex}.
    \item \textbf{Giao diện trang web thực tế / Sản phẩm nghiệp dư \& chuyên nghiệp}: Ưu tiên kết hợp \textbf{Mã Hex} với \textbf{CSS Variables} để quản lý tập trung.
    \item \textbf{Cần hiệu ứng trong suốt, lớp phủ overlay, bóng mờ}: Sử dụng \textbf{RGBA} hoặc \textbf{HSLA}.
    \item \textbf{Xây dựng dải màu chủ đạo, hệ màu động sáng/tối (\term{Dark Mode Switcher})}: Sử dụng \textbf{HSL kết hợp Biến CSS (\term{CSS Variables}) và JavaScript}.
\end{enumerate}
\end{conceptBox}

% ==========================================
% PHẦN 3.6: TÍNH KẾ THỪA CỦA THUỘC TÍNH COLOR
% ==========================================
\paragraph{3.6. Tính Kế Thừa Của Thuộc Tính Color (Color Property Inheritance Masterclass):}

\begin{conceptBox}[Khái Niệm Cơ Bản Về Tính Kế Thừa Của color]
\textbf{1. Tính Kế Thừa Mặc Định (Inherited Property)}:
\begin{itemize}[leftmargin=1.5em, itemsep=3pt, parsep=0pt]
    \item Thuộc tính \codeinline{color} trong CSS là một \term{thuộc tính kế thừa mặc định} (\term{Inherited Property}).
    \item \textbf{Nguyên lý hoạt động}: Nếu bạn không khai báo thuộc tính \codeinline{color} trực tiếp cho một phần tử con, phần tử đó sẽ tự động kế thừa giá trị \codeinline{color} từ phần tử cha chứa nó trong cây DOM.
\end{itemize}
\end{conceptBox}

\textbf{2. Phân Loại Thuộc Tính CSS Theo Tính Kế Thừa:}

\begin{prettyTableBox}[Bảng So Sánh Thuộc Tính CSS Có Kế Thừa Và Không Kế Thừa]
\rowcolors{2}{white}{primaryBlue!4}
\sffamily\small
\begin{tabularx}{\linewidth}{|>{\raggedright\arraybackslash \bfseries\color{primaryBlue}}p{3.5cm}|>{\ttfamily\color{accentOrange}}p{4.5cm}|>{\centering\arraybackslash}p{2.2cm}|>{\raggedright\arraybackslash}X|}
\hline
\rowcolor{primaryBlue} \thd{\sffamily Phân loại} & \thd{\sffamily Các thuộc tính tiêu biểu} & \thd{\sffamily Có kế thừa?} & \thd{\sffamily Ghi chú đặc tính} \\
\hline
Kế thừa (Inherited) & color, font-family, line-height, visibility, text-align & \textcolor{green!60!black}{\textbf{Có}} & Tự động lan truyền xuống các thẻ con trong cây DOM. \\
\hline
Không kế thừa (Non-inherited) & margin, padding, border, width, height, background-color & \textcolor{red!80!black}{\textbf{Không}} & Chỉ có tác dụng duy nhất trên phần tử được khai báo. \\
\hline
\end{tabularx}
\end{prettyTableBox}

\textbf{3. Các Ví Dụ Mã Nguồn Về Kế Thừa Màu Sắc:}

\begin{codeWindow}[Ví Dụ 1: Kế Thừa Màu Sắc Đơn Giản Từ Thẻ Cha]
\begin{lstlisting}[language=HTML5]
<!-- HTML Structure -->
<div class="parent">
  <p>(*@ Đoạn văn này tự động kế thừa màu đỏ từ thẻ div.parent. @*)</p>
</div>

<!-- CSS Styling -->
<style>
.parent {
  color: red; /* Thiet lap mau do cho the cha */
}
</style>
\end{lstlisting}
\end{codeWindow}

\begin{codeWindow}[Ví Dụ 2: Kế Thừa Màu Sắc Lồng Nhau Qua Nhiều Cấp Thẻ DOM]
\begin{lstlisting}[language=HTML5]
<!-- Structure DOM long nhau 3 cap -->
<div style="color: blue;">
  <section>
    <p>Color nay la blue do ke thua tu div cha ngoai cung quan section!</p>
  </section>
</div>
\end{lstlisting}
\end{codeWindow}

\begin{browserWindow}[Kết Quả Hiển Thị Kế Thừa Màu Sắc Đơn Giản \& Lồng Nhau]
\vspace{2pt}
\begin{tcolorbox}[
    enhanced,
    colback=white,
    colframe=gray!30,
    arc=6pt, boxrule=0.8pt,
    top=10pt, bottom=10pt, left=12pt, right=12pt
]
  % Preview Example 1
  {\sffamily\bfseries\color{red!80!black} Đoạn văn tự động có màu Đỏ do kế thừa từ \codeinline{.parent \{ color: red; \}}}\\[6pt]
  
  % Preview Example 2
  \begin{tcolorbox}[colback=blue!5, colframe=blue!40, arc=4pt, boxrule=0.6pt, top=6pt, bottom=6pt, left=8pt, right=8pt, fontupper=\sffamily\bfseries\color{blue!80!black}]
    Đoạn văn trong \codeinline{<section>} kế thừa màu Xanh dương (\codeinline{color: blue;}) từ thẻ \codeinline{<div>} ngoài cùng.
  \end{tcolorbox}
\end{tcolorbox}
\end{browserWindow}

\textbf{4. Ép Bắt Buộc Kế Thừa Hoặc Reset Màu Với \codeinline{inherit} Và \codeinline{initial}:}

\begin{conceptBox}[Ý Nghĩa Của Hai Giá Trị Từ Khóa inherit Và initial]
Trong các trường hợp đặc biệt (như thẻ liên kết \codeinline{<a>} hoặc nút bấm \codeinline{<button>} bị trình duyệt gán kiểu mặc định User Agent Style), màu chữ có thể không tự kế thừa. Bạn có thể điều khiển trực tiếp bằng hai từ khóa:
\begin{itemize}[leftmargin=1.5em, itemsep=3pt, parsep=0pt]
    \item \textbf{\codeinline{color: inherit;}}: Ép buộc phần tử phải nhận giá trị màu chữ từ phần tử cha (vượt qua mọi thiết lập mặc định của trình duyệt).
    \item \textbf{\codeinline{color: initial;}}: Reset giá trị màu chữ về lại giá trị mặc định ban đầu của CSS (thường là màu đen \codeinline{\#000000}).
\end{itemize}
\end{conceptBox}

\begin{codeWindow}[Ví Dụ 3: Ép Thẻ Liên Kết <a> Và Button Kế Thừa Màu Bằng color: inherit]
\begin{lstlisting}[language=CSS3]
/* Custom Reset cho the a va button de ke thua mau tu the cha */
.navbar {
  color: darkblue;
}

/* Thong thuong the <a> co mau xanh tim va gach chan mac dinh cua trinh duyet */
.navbar a {
  color: inherit; /* Ep the <a> ke thua mau darkblue tu .navbar! */
  text-decoration: none;
}
\end{lstlisting}
\end{codeWindow}

\textbf{5. Kết Hợp Tính Kế Thừa Với Từ Khóa \codeinline{currentColor}:}

\begin{codeWindow}[Ví Dụ 4: Phối Hợp Kế Thừa color Và currentColor Đồng Bộ Viền]
\begin{lstlisting}[language=HTML5]
<!-- HTML Structure -->
<div class="parent">
  <div class="child">
    <p>Doan van ben trong the child.</p>
  </div>
</div>

<!-- CSS Styling -->
<style>
.parent {
  color: darkblue; /* Mau chu darkblue lan truyen xuong child va p */
}

.child {
  border: 2px solid currentColor; /* Vien tu dong mang mau darkblue do ke thua tu parent! */
  padding: 12px;
}
</style>
\end{lstlisting}
\end{codeWindow}

\begin{browserWindow}[Kết Quả Hiển Thị Kết Hợp Tính Kế Thừa Và currentColor]
\vspace{2pt}
\begin{tcolorbox}[
    enhanced,
    colback=blue!5,
    colframe=blue!80!black,
    arc=6pt, boxrule=1.5pt,
    top=8pt, bottom=8pt, left=10pt, right=10pt,
    fontupper=\sffamily\bfseries\color{blue!80!black}
]
  .child: Đường viền tự động mang màu \codeinline{darkblue} nhờ sự phối hợp giữa tính kế thừa \codeinline{color} từ \codeinline{.parent} và từ khóa \codeinline{border: 2px solid currentColor;}
\end{tcolorbox}
\end{browserWindow}

\textbf{6. Các Lưu Ý Kỹ Thuật Quan Trọng Khi Làm Việc Với Tính Kế Thừa:}

\begin{warnBox}[Những Cảnh Báo Cần Nhớ Về Kế Thừa Màu Sắc]
\begin{itemize}[leftmargin=1.5em, itemsep=3pt, parsep=0pt]
    \item \textbf{Cây DOM quyết định cấp kế thừa}: Phần tử con luôn nhận màu từ phần tử cha trực tiếp gần nhất có khai báo thuộc tính \codeinline{color}.
    \item \textbf{Bị ghi đè ngay khi con có \codeinline{color} riêng}: Nếu bạn khai báo \codeinline{color} trực tiếp trên phần tử con, thuộc tính kế thừa từ cha sẽ lập tức bị vô hiệu hóa.
    \item \textbf{Ngoại lệ User Agent Style của trình duyệt}: Các thẻ HTML như \codeinline{<a>}, \codeinline{<button>}, \codeinline{<input>}, \codeinline{<select>} có quy tắc màu mặc định riêng của trình duyệt. Bạn bắt buộc phải khai báo \codeinline{color: inherit;} nếu muốn chúng kế thừa màu từ thẻ cha.
\end{itemize}
\end{warnBox}

\textbf{7. Bảng Tổng Kết Đặc Tính Kế Thừa Của Thuộc Tính Color:}

\begin{prettyTableBox}[Bảng Tổng Kết Các Giá Trị Kế Thừa Của color]
\rowcolors{2}{white}{primaryBlue!4}
\sffamily\small
\begin{tabularx}{\linewidth}{|>{\raggedright\arraybackslash \ttfamily\color{primaryBlue}\bfseries}p{3.2cm}|>{\centering\arraybackslash}p{2.5cm}|>{\raggedright\arraybackslash}X|}
\hline
\rowcolor{primaryBlue} \thd{\sffamily Cú pháp giá trị} & \thd{\sffamily Có kế thừa?} & \thd{\sffamily Ý nghĩa \& Mục đích sử dụng} \\
\hline
color: <value> & \textcolor{green!60!black}{\textbf{Có (Mặc định)}} & Mặc định tự động lan truyền màu xuống tất cả thẻ con. \\
\hline
color: inherit & \textcolor{green!60!black}{\textbf{Có (Ép buộc)}} & Ép phần tử nhận màu từ cha (khắc phục CSS Reset / User Agent Style). \\
\hline
color: initial & \textcolor{red}{\textbf{Không}} & Reset màu chữ về mặc định của trình duyệt (thường là đen \codeinline{\#000000}). \\
\hline
\end{tabularx}
\end{prettyTableBox}

% ==========================================
% PHẦN 3.7: TÍNH KẾ THỪA CỦA THUỘC TÍNH COLOR
% ==========================================
% ==============================================================================
% CHUYÊN ĐỀ: CÁC THUỘC TÍNH FONT & TYPOGRAPHY CHUYÊN SÂU TRONG CSS
% ==============================================================================
% ==============================================================================
% CHUYÊN ĐỀ: CÁC THUỘC TÍNH FONT & TYPOGRAPHY CHUYÊN SÂU TRONG CSS
% ==============================================================================
\section{Chuyên Đề Các Thuộc Tính Font \& Typography Chuyên Sâu Trong CSS (CSS Font Masterclass)}
\label{sec:css_font_properties_masterclass}

CSS cung cấp một bộ các thuộc tính để kiểm soát mọi khía cạnh của văn bản, từ loại font đến kích thước, kiểu dáng, độ đậm (trọng lượng) và chiều cao dòng. Việc tùy chỉnh font chữ là một phần quan trọng trong việc tạo ra giao diện người dùng hấp dẫn, hiện đại và dễ đọc. Bằng cách hiểu rõ các thuộc tính và kỹ thuật nhúng font, bạn có thể kiểm soát hoàn toàn typography trên trang web của mình.

\subsection{Chuyên Đề Phông Chữ \& Typography Trong CSS (Generic Font Families Masterclass)}
\label{sec:css_typography_foundations}

\begin{conceptBox}[Khái Niệm Cơ Bản Về Generic Font Families]
\textbf{1. Generic Font Families Là Gì?}:
\begin{itemize}[leftmargin=1.5em, itemsep=3pt, parsep=0pt]
    \item \textbf{\term{Generic Font Families} (Nhóm phông chữ tổng quát)}: Là các nhóm phông chữ mặc định chuẩn hóa bởi W3C mà tất cả trình duyệt Web trên mọi hệ điều hành (Windows, macOS, Linux, Android, iOS) luôn luôn hiểu và hỗ trợ sẵn.
    \item \textbf{Vai trò cốt lõi trong cơ chế Fallback Font}: Khi người lập trình khai báo một font chữ cụ thể (ví dụ: \codeinline{"Times New Roman"}), nếu thiết bị người dùng chưa cài đặt font đó, trình duyệt sẽ tự động dự phòng (\term{Fallback}) sang nhóm font tổng quát tương ứng được chỉ định ở cuối danh sách (ví dụ: \codeinline{serif}).
\end{itemize}
\end{conceptBox}

\textbf{2. Sự Khác Biệt Trực Quan Giữa Phông Chữ Có Chân (Serif) Và Không Chân (Sans-Serif):}

\begin{browserWindow}[Minh Họa Trực Quan Sự Khác Biệt Giữa Serif Và Sans-serif]
\vspace{2pt}
\begin{tcolorbox}[
    enhanced,
    colback=white,
    colframe=gray!30,
    arc=6pt, boxrule=0.8pt,
    top=12pt, bottom=12pt, left=14pt, right=14pt
]
  \begin{center}
    {\sffamily\bfseries\Large\color{primaryBlue} Phân Biệt Nét Chân Chữ (Serifs) Trực Quan}\\[10pt]
    
    % Top Part: 3 F's comparison
    \begin{tabular}{c @{\hspace{2.2cm}} c @{\hspace{2.2cm}} c}
      {\sffamily\fontsize{42}{50}\selectfont\bfseries F} &
      {\rmfamily\fontsize{42}{50}\selectfont\bfseries F} &
      {\rmfamily\fontsize{42}{50}\selectfont\bfseries\color{red!80!black} F} \\[4pt]
      {\sffamily\bfseries Sans-serif} &
      {\rmfamily\bfseries Serif} &
      {\rmfamily\bfseries Serif}\\[-2pt]
      & & {\sffamily\small\bfseries\color{red!80!black}(red serifs)}
    \end{tabular}
    
    \vspace{12pt}
    
    % Bottom Part: Outer Rectangular Box matching the image layout
    \begin{tcolorbox}[
      colback=white,
      colframe=black!80,
      arc=0pt, boxrule=0.8pt,
      top=10pt, bottom=10pt, left=15pt, right=15pt
    ]
      \begin{tabularx}{\linewidth}{@{}l >{\raggedleft\arraybackslash}X@{}}
        {\sffamily\Huge\bfseries AaBbCc} & {\sffamily\large\color{black!80} Sans-serif font} \\[14pt]
        {\rmfamily\Huge\bfseries AaBbCc} & {\rmfamily\large\color{black!80} Serif font} \\[14pt]
        {\rmfamily\Huge\bfseries\color{red!80!black} AaBbCc} & {\rmfamily\large\color{red!80!black} Serif font (red serifs)}
      \end{tabularx}
    \end{tcolorbox}
  \end{center}
\end{tcolorbox}
\end{browserWindow}

\begin{prettyTableBox}[Bảng So Sánh Nhanh Giữa Serif Và Sans-serif]
\rowcolors{2}{white}{primaryBlue!4}
\sffamily\small
\begin{tabularx}{\linewidth}{|>{\raggedright\arraybackslash \bfseries\color{primaryBlue}}p{3.2cm}|>{\raggedright\arraybackslash}X|>{\raggedright\arraybackslash}X|}
\hline
\rowcolor{primaryBlue} \thd{\sffamily Tiêu chí so sánh} & \thd{\sffamily Phông chữ Serif (Có chân)} & \thd{\sffamily Phông chữ Sans-serif (Không chân)} \\
\hline
Nét chữ & Có các nét gạch nhỏ ở đầu và chân ký tự. & Không có chân chữ, nét chữ đơn giản gọn phẳng. \\
\hline
Phong cách thiết kế & Cổ điển, sang trọng, nghiêm túc, trang trọng. & Hiện đại, sạch sẽ, gọn gàng, thân thiện. \\
\hline
Độ dễ đọc trên giấy in & Dễ đọc hơn nhờ các chân dẫn hướng mắt khi đọc bài dài. & Đọc ổn nhưng có thể thiếu điểm nhấn phân dòng. \\
\hline
Độ dễ đọc trên màn hình & Có thể mỏi mắt/khó đọc ở độ phân giải màn hình thấp. & Rất dễ đọc, đặc biệt trên màn hình di động nhỏ. \\
\hline
Ứng dụng phổ biến & Sách in, báo chí, luận văn, tài liệu chính thức. & Trang web, ứng dụng di động, giao diện UI/UX. \\
\hline
\end{tabularx}
\end{prettyTableBox}

\textbf{3. Bảng Phân Tích Chi Tiết 5 Nhóm Generic Font Families Chuẩn W3C:}

\begin{prettyTableBox}[Bảng Phân Tích Chi Tiết 5 Nhóm Phông Chữ Tổng Quát (Generic Font Families)]
\rowcolors{2}{white}{primaryBlue!4}
\sffamily\small
\begin{tabularx}{\linewidth}{|>{\raggedright\arraybackslash \ttfamily\color{primaryBlue}\bfseries}p{2.2cm}|>{\raggedright\arraybackslash}p{3.2cm}|>{\raggedright\arraybackslash}p{3.0cm}|>{\ttfamily\color{accentOrange}}p{3.0cm}|>{\raggedright\arraybackslash}X|}
\hline
\rowcolor{primaryBlue} \thd{\sffamily Tên nhóm} & \thd{\sffamily Đặc điểm kỹ thuật} & \thd{\sffamily Cảm nhận thị giác} & \thd{\sffamily Ví dụ font cụ thể} & \thd{\sffamily Ứng dụng phù hợp nhất} \\
\hline
serif & Nét gạch nhỏ (serif) ở chân/đầu ký tự. & Truyền thống, nghiêm túc, chuyên nghiệp. & Times New Roman, Georgia, Garamond & Văn bản dài, sách báo, luận văn, tiêu đề trang trọng. \\
\hline
sans-serif & Không chân chữ, nét phẳng đều gọn gàng. & Hiện đại, sạch sẽ, dễ nhìn kỹ thuật số. & Arial, Helvetica, Roboto, Open Sans & Giao diện người dùng Web/App, email, blog cá nhân. \\
\hline
monospace & Chiều rộng các ký tự bằng nhau tuyệt đối. & Rõ ràng, đồng đều, dễ soi từng ký tự. & Courier New, Consolas, Lucida Console & Code editor, terminal, log hệ thống, bảng dữ liệu. \\
\hline
cursive & Nét nghiêng mô phỏng chữ viết tay. & Cá tính, nghệ thuật, mềm mại, trang nhã. & Brush Script, Comic Sans, Pacifico & Thiệp mời, chữ ký, tiêu đề trang trí nhẹ, blog cá nhân. \\
\hline
fantasy & Kiểu chữ cách điệu trang trí độc đáo. & Độc lạ, vui nhộn, gây chú ý mạnh mẽ. & Impact, Papyrus, Jokerman, Lobster & Poster quảng cáo, game, tiêu đề thu hút sự chú ý. \\
\hline
\end{tabularx}
\end{prettyTableBox}

\textbf{4. Các Phông Chữ Hay Dùng Nhất Trong Thiết Kế Web Theo Từng Nhóm:}

\begin{prettyTableBox}[Danh Sách Các Phông Chữ Tiêu Chuẩn Phổ Biến Trong Thiết Kế Web]
\sffamily\small
\begin{tabularx}{\linewidth}{|>{\raggedright\arraybackslash \bfseries\color{primaryBlue}}p{3.2cm}|>{\ttfamily\color{accentOrange}\bfseries}p{3.8cm}|>{\raggedright\arraybackslash}X|}
\hline
\rowcolor{primaryBlue} \thd{\sffamily Nhóm họ phông} & \thd{\sffamily Tên font cụ thể} & \thd{\sffamily Đặc điểm \& Trường hợp ứng dụng thực tế} \\
\hline
\rowcolor{white} \textbf{Serif (Có chân)} & Times New Roman & Cổ điển, trang trọng (dùng nhiều trong báo chí, tài liệu in). \\
\cline{2-3}
\rowcolor{primaryBlue!4} & Georgia & Thanh lịch, thiết kế tối ưu dễ đọc trên màn hình máy tính hơn Times. \\
\cline{2-3}
\rowcolor{white} & Merriweather & Thân thiện, nét chữ rõ ràng lý tưởng cho blog đọc nội dung dài. \\
\hline
\rowcolor{primaryBlue!4} \textbf{Sans-serif (Không chân)} & Arial & Đơn giản, phổ biến nhất thế giới trên mọi hệ điều hành. \\
\cline{2-3}
\rowcolor{white} & Roboto & Font mặc định của hệ điều hành Android, hiện đại, sạch sẽ. \\
\cline{2-3}
\rowcolor{primaryBlue!4} & Open Sans & Dễ đọc tuyệt vời, phù hợp cho hệ thống Web/App quy mô lớn. \\
\cline{2-3}
\rowcolor{white} & Lato & Mềm mại, cân bằng, hợp cho cả tiêu đề lẫn văn bản chính. \\
\cline{2-3}
\rowcolor{primaryBlue!4} & Montserrat & Đậm chất thiết kế hiện đại, tạo cảm giác chuyên nghiệp sáng tạo. \\
\hline
\rowcolor{white} \textbf{Monospace (Rộng đều)} & Courier New & Phong cách máy đánh chữ cổ điển truyền thống. \\
\cline{2-3}
\rowcolor{primaryBlue!4} & Consolas & Font chuẩn mực khi viết mã nguồn code, gọn gàng dễ đọc. \\
\hline
\rowcolor{white} \textbf{Cursive (Viết tay)} & Pacifico & Nét chữ viết tay đáng yêu, hợp cho tiêu đề nghệ thuật hoặc Logo. \\
\cline{2-3}
\rowcolor{primaryBlue!4} & Dancing Script & Nét chữ uyển chuyển, trang nhã cho thiệp mừng / banner. \\
\hline
\rowcolor{white} \textbf{Fantasy (Cách điệu)} & Impact & Nét chữ cực đậm, thu hút sự chú ý ngay lập tức cho tiêu đề. \\
\cline{2-3}
\rowcolor{primaryBlue!4} & Playfair Display & Sang trọng, phù hợp cho giao diện thời trang cao cấp. \\
\cline{2-3}
\rowcolor{white} & Lobster & Độc đáo, nghệ thuật, tạo điểm nhấn bắt mắt sơ khởi. \\
\hline
\end{tabularx}
\end{prettyTableBox}

\textbf{5. Cú Pháp Sử Dụng Trong CSS \& Cơ Chế Fallback Font List:}

\begin{codeWindow}[Cú Pháp Khai Báo Danh Sách Phông Chữ Dự Phòng (Fallback Chain)]
\begin{lstlisting}[language=CSS3]
/* 1. Dinh kieu van ban chinh voi chuoi Fallback Sans-serif */
body {
  font-family: "Open Sans", Arial, Helvetica, sans-serif;
}

/* 2. Dinh kieu ma nguoi va khoi Code voi Monospace */
code, pre {
  font-family: Consolas, "Courier New", monospace;
}

/* 3. Dinh kieu tieu de trang trong voi Serif */
h1.academic-title {
  font-family: "Playfair Display", Georgia, "Times New Roman", serif;
}
\end{lstlisting}
\end{codeWindow}

\begin{browserWindow}[Kết Quả Hiển Thị Của Các Kiểu Phông Chữ Khác Nhau Trên Giao Diện]
\vspace{2pt}
\begin{tcolorbox}[
    enhanced,
    colback=white,
    colframe=gray!30,
    arc=6pt, boxrule=0.8pt,
    top=10pt, bottom=10pt, left=12pt, right=12pt
]
  % Sans-serif preview
  {\sffamily\bfseries\color{primaryBlue} 1. Sans-serif (Roboto / Arial)}:\\[2pt]
  {\sffamily Đây là đoạn văn bản không chân hiện đại, gọn gàng và rất dễ đọc trên màn hình kỹ thuật số.}\\[8pt]
  
  % Serif preview
  {\sffamily\bfseries\color{primaryBlue} 2. Serif (Georgia / Times New Roman)}:\\[2pt]
  {\rmfamily Đây là đoạn văn bản có chân trang trọng, mang phong cách cổ điển báo chí và tài liệu chính thức.}\\[8pt]
  
  % Monospace preview
  {\sffamily\bfseries\color{primaryBlue} 3. Monospace (Consolas / Courier New)}:\\[2pt]
  {\ttfamily const greeting = "Hello World!"; // Ký tự rộng bằng nhau tuyệt đối}\\[8pt]
  
  % Cursive / Decorative text description
  {\sffamily\bfseries\color{primaryBlue} 4. Cursive \& Fantasy (Chữ nghệ thuật \& Tiêu đề cách điệu)}:\\[2pt]
  {\sffamily\small\color{purple!80!black}\textit{Chữ viết tay nghệ thuật (Cursive) và chữ cách điệu đậm nét (Fantasy) dùng tạo điểm nhấn Tooltip / Banner.}}
\end{tcolorbox}
\end{browserWindow}

\textbf{6. Tóm Tắt Nhanh Cheatsheet Lựa Chọn Phông Chữ:}

\begin{prettyTableBox}[Bảng Tóm Tắt Cheatsheet Lựa Chọn Loại Font Cho Dự Án]
\rowcolors{2}{white}{primaryBlue!4}
\sffamily\small
\begin{tabularx}{\linewidth}{|>{\raggedright\arraybackslash \ttfamily\color{primaryBlue}\bfseries}p{2.5cm}|>{\bfseries\color{accentOrange}}p{3.8cm}|>{\raggedright\arraybackslash}X|}
\hline
\rowcolor{primaryBlue} \thd{\sffamily Loại font} & \thd{\sffamily Cảm giác chính thị giác} & \thd{\sffamily Trường hợp ứng dụng hàng đầu} \\
\hline
serif & Trang trọng, cổ điển, trang nhã. & Ấn phẩm in ấn, sách báo, luận văn học thuật. \\
\hline
sans-serif & Gọn gàng, hiện đại, thân thiện. & Trang web, giao diện UI/UX, ứng dụng Mobile. \\
\hline
monospace & Kỹ thuật, rõ ràng, chính xác. & Trình soạn thảo Code, terminal, bảng số liệu. \\
\hline
cursive & Nghệ thuật, mềm mại, tinh tế. & Thiệp mừng, chữ ký, banner trang trí nhẹ. \\
\hline
fantasy & Kỳ lạ, độc đáo, thu hút mạnh. & Tiêu đề quảng cáo đặc biệt, Poster, Game UI. \\
\hline
\end{tabularx}
\end{prettyTableBox}

\subsection{Vấn Đề 1: Phông Chữ Mặc Định (System Fonts \& Web-Safe Fonts)}
\label{sec:system_font_stack_van_de_1}

\begin{conceptBox}[VẤN ĐỀ 1: PHÔNG CHỮ MẶC ĐỊNH (SYSTEM FONTS \& WEB-SAFE FONTS)]
\textbf{Khái niệm cốt lõi}: Đối với các phông chữ mặc định (system fonts hoặc web-safe fonts), bạn \textbf{không cần nhúng font} hay dùng \codeinline{@font-face} hoặc Google Fonts -- chỉ cần gọi đúng tên trong \codeinline{font-family} là được. Trình duyệt sẽ tự tìm font có sẵn trên hệ điều hành của người dùng.
\end{conceptBox}

\paragraph{7. Chuyên Đề System Font Stack Trong CSS (System-UI Masterclass):}
\label{sec:system_font_stack}

System Font (phông chữ hệ thống) là phông chữ mặc định được cài đặt sẵn trên hệ điều hành mà người dùng đang sử dụng (như Windows, macOS, iOS, Android, Linux). Khi bạn khai báo system font trong CSS, trình duyệt sẽ hiển thị văn bản bằng đúng phông chữ hệ thống đó, giúp giao diện trang web đồng bộ tuyệt đối với các thành phần của hệ điều hành (menu, tiêu đề cửa sổ, nút bấm hệ thống...).

\begin{conceptBox}[Khái Niệm Cốt Lõi Về System Font]
\textbf{1. System Font Là Gì?}:
\begin{itemize}[leftmargin=1.5em, itemsep=3pt, parsep=0pt]
    \item \textbf{Khái niệm}: System Font là bộ phông chữ giao diện người dùng mặc định (\term{Native UI Font}) của hệ điều hành.
    \item \textbf{Cơ chế hiển thị}: Trình duyệt truy xuất trực tiếp font chữ của hệ điều hành mà không cần tải bất kỳ tệp font nào từ máy chủ web hoặc dịch vụ bên ngoài (như Google Fonts).
\end{itemize}
\textbf{2. Tài Liệu \& Nguồn Tham Khảo Chuẩn Kỹ Thuật}:
\begin{itemize}[leftmargin=1.5em, itemsep=3pt, parsep=0pt]
    \item Tham khảo cấu trúc chuỗi System Font Stack tại: \href{https://css-tricks.com/snippets/css/system-font-stack/}{\codeinline{CSS-Tricks System Font Stack}}.
    \item Kiểm tra mức độ phổ biến của các phông chữ tại: \href{https://www.cssfontstack.com/}{\codeinline{CSS Font Stack}}.
\end{itemize}
\end{conceptBox}

\paragraph{Ưu Điểm Và Lý Do Nên Sử Dụng System Font:}

\begin{prettyTableBox}[Bảng Phân Tích Chi Tiết Các Ưu Điểm Của System Font]
\rowcolors{2}{white}{primaryBlue!4}
\sffamily\small
\begin{tabularx}{\linewidth}{|>{\raggedright\arraybackslash \bfseries\color{primaryBlue}}p{3.2cm}|>{\bfseries\color{accentOrange}}p{4.2cm}|>{\raggedright\arraybackslash}X|}
\hline
\rowcolor{primaryBlue} \thd{\sffamily Ưu điểm cốt lõi} & \thd{\sffamily Lý do kỹ thuật} & \thd{\sffamily Lợi ích trải nghiệm người dùng (UX)} \\
\hline
Tốc độ cực nhanh (Fast) & Không cần tải font từ web (Google Fonts / Custom Web Fonts). & Trang web hiển thị tức thì, loại bỏ hoàn toàn hiện tượng nhấp nháy chữ (\term{FOUT / FOIT}), tối ưu chỉ số Core Web Vitals. \\
\hline
Tương thích cao (High Compatibility) & Sử dụng phông chữ hệ thống được tối ưu hóa cho màn hình thiết bị. & Hiển thị chuẩn xác, thân thuộc và thân thiện với người dùng trên mọi hệ điều hành. \\
\hline
Thống nhất giao diện (UI Consistency) & Phông chữ giống hệt với giao diện hệ điều hành (OS/UI). & Giao diện trang web hòa nhập tự nhiên như một ứng dụng Native App của thiết bị. \\
\hline
Responsive tốt (High Adaptability) & Tự động tối ưu trên nhiều thiết bị và độ phân giải khác nhau. & Hiển thị sắc nét trên cả màn hình Retina / High-DPI và màn hình tiêu chuẩn. \\
\hline
\end{tabularx}
\end{prettyTableBox}

\paragraph{Cách Sử Dụng system-ui Trong CSS (Cú Pháp Đơn Giản vs Fallback Đầy Đủ):}

CSS Fonts Module Level 4 giới thiệu giá trị từ khóa \codeinline{system-ui}, cho phép gọi trực tiếp phông chữ mặc định của hệ điều hành hiện tại chỉ với một dòng mã.

\begin{codeWindow}[1. Cú Pháp Đơn Giản Nhất Với system-ui (CSS Fonts Level 4)]
\begin{lstlisting}[language=CSS3]
/* Tu dong goi font mac dinh cua he dieu hanh hien tai */
body {
  font-family: system-ui;
}
\end{lstlisting}
\end{codeWindow}

\begin{prettyTableBox}[Danh Sách Phông Chữ Mặc Định (UI Font) Theo Từng Hệ Điều Hành]
\rowcolors{2}{white}{primaryBlue!4}
\sffamily\small
\begin{tabularx}{\linewidth}{|>{\raggedright\arraybackslash \bfseries\color{primaryBlue}}p{4.2cm}|>{\ttfamily\color{accentOrange}\bfseries}p{4.2cm}|>{\raggedright\arraybackslash}X|}
\hline
\rowcolor{primaryBlue} \thd{\sffamily Hệ điều hành (OS)} & \thd{\sffamily Font mặc định (UI Font)} & \thd{\sffamily Ghi chú đặc tính thị giác} \\
\hline
macOS & San Francisco & Phông chữ tối giản, sắc nét trên màn hình Retina của Apple. \\
\hline
iOS & San Francisco & Font chuẩn mực trên iPhone / iPad. \\
\hline
Windows 10+ & Segoe UI & Phông chữ giao diện chuẩn mực của Microsoft Windows. \\
\hline
Android & Roboto & Phông chữ mặc định hệ sinh thái Google / Android. \\
\hline
Ubuntu / Linux & Ubuntu & Phông chữ tự do chuẩn mực của hệ điều hành Ubuntu. \\
\hline
Older Windows (7 / XP / Vista) & Tahoma, Arial & Các phông chữ dự phòng cổ điển trên hệ điều hành Windows cũ. \\
\hline
\end{tabularx}
\end{prettyTableBox}

\begin{codeWindow}[2. Chuỗi Khai Báo Dự Phòng Đầy Đủ (Full System Font Stack Fallback)]
\begin{lstlisting}[language=CSS3]
/* Chuoi Fallback System Font Stack giup tuong thich tuyet doi moi trinh duyet & OS */
body {
  font-family:
    -apple-system,       /* macOS & iOS (Safari / Apple WebKit) */
    BlinkMacSystemFont,  /* Chrome & Opera tren macOS (Blink Engine) */
    "Segoe UI",          /* Windows 10 / 11 */
    Roboto,              /* Android */
    Ubuntu,              /* Ubuntu Linux */
    "Helvetica Neue",    /* macOS / iOS phien ban cu */
    sans-serif;          /* Generic Fallback cuoi cung */
}
\end{lstlisting}
\end{codeWindow}

\begin{tipBox}[Cơ Chế Hoạt Động Của Chuỗi Fallback Font]
Trình duyệt sẽ duyệt danh sách từ trên xuống dưới và chọn phông chữ đầu tiên có sẵn trên hệ điều hành của người dùng. Nếu không tìm thấy font nào trong danh sách, trình duyệt sẽ tự động dự phòng xuống nhóm font tổng quát \codeinline{sans-serif}.
\end{tipBox}

\paragraph{Ví Dụ Thực Tế \& Kết Quả Hiển Thị System Font:}

\begin{codeWindow}[Ví Dụ Mã Nguồn HTML \& CSS Sử Dụng System Font]
\begin{lstlisting}[language=HTML5]
<!DOCTYPE html>
<html lang="vi">
<head>
  <meta charset="UTF-8">
  <title>Vi Du System Font</title>
  <style>
    body {
      font-family: system-ui, -apple-system, BlinkMacSystemFont, "Segoe UI", Roboto, sans-serif;
      font-size: 16px;
      line-height: 1.6;
      color: #212529;
      padding: 20px;
    }
  </style>
</head>
<body>
  <h1>(*@ Tiêu đề theo system font @*)</h1>
  <p>(*@ Đây là đoạn văn bản sử dụng font mặc định của hệ điều hành bạn đang dùng. @*)</p>
</body>
</html>
\end{lstlisting}
\end{codeWindow}

\begin{browserWindow}[Kết Quả Hiển Thị Thực Tế Của System Font Trên Trình Duyệt]
\sffamily\small
\begin{tcolorbox}[arc=4pt, colback=white, colframe=primaryBlue!30, boxrule=0.8pt, top=12pt, bottom=12pt, left=14pt, right=14pt]
  {\fontsize{18}{22}\selectfont\bfseries\color{primaryBlue} Tiêu đề theo system font} \par\vspace{6pt}
  {\fontsize{12}{16}\selectfont Đây là đoạn văn bản sử dụng font mặc định của hệ điều hành bạn đang dùng. Giao diện tự động hòa nhập với phông chữ hệ thống (San Francisco trên macOS/iOS, Segoe UI trên Windows, Roboto trên Android).}
\end{tcolorbox}
\end{browserWindow}

\paragraph{Lưu Ý Quan Trọng Khi Sử Dụng System Font:}

\begin{prettyTableBox}[Bảng Phân Tích Lưu Ý \& Giải Thích Kỹ Thuật Khi Dùng System Font]
\rowcolors{2}{white}{primaryBlue!4}
\sffamily\small
\begin{tabularx}{\linewidth}{|>{\raggedright\arraybackslash \bfseries\color{primaryBlue}}p{4.2cm}|>{\raggedright\arraybackslash}X|}
\hline
\rowcolor{primaryBlue} \thd{\sffamily Lưu ý kỹ thuật} & \thd{\sffamily Giải thích chi tiết \& Giải pháp} \\
\hline
Không đồng nhất giữa các hệ điều hành & Mỗi hệ điều hành (macOS, Windows, Android) sử dụng một phông chữ hệ thống riêng, khiến giao diện văn bản có sự chênh lệch nhẹ về nét chữ giữa các thiết bị khác nhau. \\
\hline
Không thể chọn chính xác từng glyph / font cụ thể & Phụ thuộc hoàn toàn vào hệ điều hành của người dùng, nhà phát triển không thể can thiệp tùy chỉnh chi tiết từng nét glyph cố định. \\
\hline
Không hỗ trợ tốt các hệ thống cũ & Từ khóa \codeinline{system-ui} đòi hỏi trình duyệt hiện đại (Chrome 56+, Firefox 43+, Safari 11+). Bắt buộc phải kết hợp chuỗi fallback đầy đủ (\codeinline{-apple-system, "Segoe UI", Roboto...}) để đảm bảo hiển thị ổn định trên thiết bị cũ. \\
\hline
\end{tabularx}
\end{prettyTableBox}

\paragraph{Tổng Kết Chuyên Đề System Font trong CSS:}

\begin{prettyTableBox}[Bảng Tổng Kết Kiến Thức Cốt Lõi Về System Font]
\rowcolors{2}{white}{primaryBlue!4}
\sffamily\small
\begin{tabularx}{\linewidth}{|>{\raggedright\arraybackslash \ttfamily\color{primaryBlue}\bfseries}p{3.5cm}|>{\raggedright\arraybackslash}X|}
\hline
\rowcolor{primaryBlue} \thd{\sffamily Thuộc tính / Khái niệm} & \thd{\sffamily Ý nghĩa \& Giá trị sử dụng} \\
\hline
system-ui & Gọi trực tiếp phông chữ giao diện mặc định của hệ điều hành hiện tại. \\
\hline
font-family fallback & Chuỗi phông chữ dự phòng đảm bảo hiển thị ổn định khi trình duyệt cũ không hỗ trợ \codeinline{system-ui}. \\
\hline
Ưu điểm vượt trội & Tốc độ tải cực nhanh (không tốn dung lượng mạng), giao diện thân thuộc với OS, hiệu suất cao. \\
\hline
Khuyết điểm hạn chế & Không kiểm soát chi tiết 100\% về kiểu dáng chữ cố định trên mọi thiết bị người dùng. \\
\hline
\end{tabularx}
\end{prettyTableBox}


% ==============================================================================


\subsection{Chi Tiết 7 Thuộc Tính Font Cốt Lõi Trong CSS}

\subsubsection{1. Thuộc Tính font-family \& Quy Tắc Khai Báo Chuẩn (Vấn Đề 2: Quy Tắc Viết font-family Trong CSS)}

\begin{conceptBox}[VẤN ĐỀ 2: QUY TẮC VIẾT FONT-FAMILY TRONG CSS]
\textbf{Khái niệm cốt lõi}: Thuộc tính \codeinline{font-family} nhận một danh sách các phông chữ ưu tiên từ trái qua phải, phân tách bằng dấu phẩy và bắt buộc kết thúc bằng nhóm phông chữ tổng quát (\term{generic font family}). Tên font có khoảng trắng phải được bọc trong dấu ngoặc kép.
\end{conceptBox}

\label{sec:font_family}

Thuộc tính \codeinline{font-family} trong CSS là một trong những thuộc tính quan trọng nhất để kiểm soát giao diện văn bản trên trang web của bạn. Nó cho phép bạn chỉ định một hoặc nhiều họ phông chữ (\term{font family}) mà trình duyệt nên sử dụng để hiển thị văn bản. Việc hiểu rõ cách hoạt động của nó giúp bạn đảm bảo trang web của mình trông nhất quán và chuyên nghiệp trên nhiều thiết bị và trình duyệt khác nhau.


\begin{tipBox}[Ẩn Dụ Thực Tế Dễ Nhớ: Chuỗi Fallback Giống Như "Gọi Món Tại Quán Cà Phê"]
Để dễ hình dung cách trình duyệt xử lý chuỗi \codeinline{font-family: "CustomFont", Arial, sans-serif;}, hãy tưởng tượng bạn đi gọi nước uống:
\begin{enumerate}[leftmargin=1.5em, itemsep=3pt, parsep=0pt]
    \item \textbf{Món ưu tiên 1 ("CustomFont")}: Bạn gọi món sinh tố đặc biệt của quán. Nếu quán có sẵn tệp font này (được nhúng hoặc có trong máy), trình duyệt sẽ phục vụ ngay.
    \item \textbf{Món dự phòng 2 ("Arial")}: Nếu quán hết món đặc biệt, bạn bảo nhân viên đổi sang ly nước cam \codeinline{Arial} phổ biến mà quán nào cũng sẵn có.
    \item \textbf{Món mặc định nhóm ("sans-serif")}: Nếu quán cũng không có sẵn \codeinline{Arial}, bạn bảo: "Lấy cho tôi bất kỳ loại nước không chân (\codeinline{sans-serif}) nào quán đang có!". Trình duyệt sẽ dùng font mặc định của hệ điều hành.
\end{enumerate}
\end{tipBox}

\begin{conceptBox}[Sơ Đồ Luồng Quyết Định (Decision Flowchart) Của Trình Duyệt Khi Render Font]
\sffamily\small
\vspace{2pt}

% Start Header
\begin{center}
\tcbox[on line, arc=4pt, colback=primaryBlue!10, colframe=primaryBlue!50, boxrule=0.8pt, top=4pt, bottom=4pt, left=12pt, right=12pt]{%
  \bfseries\color{primaryBlue} [ Bắt đầu: Trình duyệt bắt đầu vẽ ký tự văn bản trên màn hình ]%
}
\end{center}
\vspace{-2pt}

% Step 1 Card
\begin{tcolorbox}[arc=5pt, colback=white, colframe=primaryBlue!30, boxrule=0.8pt, top=6pt, bottom=6pt, left=10pt, right=10pt]
  \textbf{\color{primaryBlue}Bước 1}: Kiểm tra phông chữ đầu tiên trong danh sách (Ví dụ: \codeinline{"Roboto"})\\[4pt]
  \begin{tabular}{@{}l l@{}}
    \tcbox[on line, arc=3pt, colback=green!10, colframe=green!60!black, boxsep=1pt, left=4pt, right=4pt, top=1pt, bottom=1pt]{\bfseries\color{green!60!black}\scriptsize ĐÚNG (Có sẵn)} & $\rightarrow$ Dùng font \codeinline{"Roboto"} để hiển thị văn bản ngay lập tức. \\ [3pt]
    \tcbox[on line, arc=3pt, colback=red!10, colframe=red!70!black, boxsep=1pt, left=4pt, right=4pt, top=1pt, bottom=1pt]{\bfseries\color{red!70!black}\scriptsize SAI (Không có)} & $\rightarrow$ Bỏ qua, chuyển sang kiểm tra phông chữ thứ 2.
  \end{tabular}
\end{tcolorbox}

\begin{center}
\textcolor{primaryBlue!70}{\textbf{$\downarrow$}}
\end{center}
\vspace{-4pt}

% Step 2 Card
\begin{tcolorbox}[arc=5pt, colback=white, colframe=primaryBlue!30, boxrule=0.8pt, top=6pt, bottom=6pt, left=10pt, right=10pt]
  \textbf{\color{primaryBlue}Bước 2}: Kiểm tra phông chữ thứ hai trong danh sách (Ví dụ: \codeinline{Arial})\\[4pt]
  \begin{tabular}{@{}l l@{}}
    \tcbox[on line, arc=3pt, colback=green!10, colframe=green!60!black, boxsep=1pt, left=4pt, right=4pt, top=1pt, bottom=1pt]{\bfseries\color{green!60!black}\scriptsize ĐÚNG (Có sẵn)} & $\rightarrow$ Dùng font \codeinline{Arial} để hiển thị văn bản. \\ [3pt]
    \tcbox[on line, arc=3pt, colback=red!10, colframe=red!70!black, boxsep=1pt, left=4pt, right=4pt, top=1pt, bottom=1pt]{\bfseries\color{red!70!black}\scriptsize SAI (Không có)} & $\rightarrow$ Bỏ qua, chuyển sang bước lấy nhóm phông tổng quát.
  \end{tabular}
\end{tcolorbox}

\begin{center}
\textcolor{primaryBlue!70}{\textbf{$\downarrow$}}
\end{center}
\vspace{-4pt}

% Step 3 Card
\begin{tcolorbox}[arc=5pt, colback=white, colframe=primaryBlue!30, boxrule=0.8pt, top=6pt, bottom=6pt, left=10pt, right=10pt]
  \textbf{\color{primaryBlue}Bước 3 (Dự phòng cuối)}: Lấy phông chữ mặc định thuộc nhóm tổng quát (Ví dụ: \codeinline{sans-serif}) từ Hệ Điều Hành.\\[4pt]
  \tcbox[on line, arc=3pt, colback=primaryBlue!10, colframe=primaryBlue!50, boxsep=1pt, left=4pt, right=4pt, top=1pt, bottom=1pt]{\bfseries\color{primaryBlue}\scriptsize Kết Quả Tổng Kết} $\rightarrow$ \textbf{Văn bản luôn luôn hiển thị ổn định, không bao giờ bị vỡ hoặc mất chữ!}
\end{tcolorbox}

\end{conceptBox}
\subsection*{I. Cú Pháp và Thứ Tự Ưu Tiên Của Chuỗi Fallback}

Thuộc tính \codeinline{font-family} nhận một danh sách các tên font, được phân tách bằng dấu phẩy. Thứ tự này rất quan trọng vì nó xác định thứ tự ưu tiên mà trình duyệt sẽ sử dụng để tìm và áp dụng font chữ:

\begin{codeWindow}[Cú Pháp Định Dạng Họ Phông Chữ Chuẩn]
\begin{lstlisting}[language=CSS3]
p {
  font-family: "Ten Font 1", "Ten Font 2", Ten_Font_Chung;
}
\end{lstlisting}
\end{codeWindow}

\begin{conceptBox}[Cơ Chế Xử Lý Thứ Tự Ưu Tiên Của Trình Duyệt]
Trình duyệt sẽ duyệt qua danh sách từ trái sang phải. Nó sẽ cố gắng sử dụng:
\begin{enumerate}[leftmargin=1.5em, itemsep=3pt, parsep=0pt]
    \item \textbf{"Tên Font 1"}: Nếu font này có sẵn trên hệ thống của người dùng hoặc được nhúng vào trang web.
    \item \textbf{"Tên Font 2"}: Nếu "Tên Font 1" không có, trình duyệt sẽ chuyển sang tìm font này.
    \item \textbf{Tên\_Font\_Chung}: Nếu không có font cụ thể nào trong danh sách được tìm thấy, trình duyệt sẽ sử dụng một font chung từ hệ thống của người dùng thuộc loại đó. Đây là một font dự phòng (\term{fallback font}) quan trọng.
\end{enumerate}
\end{conceptBox}

\paragraph{Ví Dụ Cụ Thể Phân Tích Chuỗi Dự Phòng:}
\begin{codeWindow}[Khai Báo Chuỗi Fallback Cho Body]
\begin{lstlisting}[language=CSS3]
body {
  font-family: "Helvetica Neue", Arial, sans-serif;
}
\end{lstlisting}
\end{codeWindow}

\begin{browserWindow}[Kết Quả Phân Tích Hiển Thị Trực Quan]
\begin{tcolorbox}[arc=4pt, colback=white, colframe=primaryBlue!30, boxrule=0.8pt, top=10pt, bottom=10pt, left=12pt, right=12pt]
  {\sffamily\bfseries\color{primaryBlue} Phân tích luồng ưu tiên trình duyệt:}\par\vspace{4pt}
  \begin{itemize}[leftmargin=1.2em, itemsep=3pt, parsep=0pt]
    \item Trình duyệt sẽ ưu tiên sử dụng \codeinline{Helvetica Neue}.
    \item Nếu \codeinline{Helvetica Neue} không có sẵn, nó sẽ sử dụng \codeinline{Arial}.
    \item Cuối cùng, nếu cả hai đều không có, nó sẽ chọn một font \codeinline{sans-serif} bất kỳ có sẵn trên hệ thống.
  \end{itemize}
\end{tcolorbox}
\end{browserWindow}

\subsection*{II. Các Loại Tên Font Trong CSS}

Có hai loại tên font bạn có thể sử dụng trong thuộc tính \codeinline{font-family}:

\subsubsection*{1. Tên Font Cụ Thể (Specific Font Names)}
Đây là tên chính xác của một font chữ cụ thể. Ví dụ như \codeinline{"Roboto"}, \codeinline{"Times New Roman"}, \codeinline{"Arial"}, \codeinline{"Open Sans"}.

\begin{conceptBox}[Quy Tắc Đặt Dấu Ngoặc Kép \& Đánh Giá]
\begin{itemize}[leftmargin=1.5em, itemsep=3pt, parsep=0pt]
    \item \textbf{Quy tắc ngoặc kép}:
    \begin{itemize}[leftmargin=1.2em, itemsep=2pt]
        \item Nếu tên font có chứa khoảng trắng, bạn \textbf{bắt buộc phải đặt nó trong dấu ngoặc kép} (đơn hoặc kép). Ví dụ: \codeinline{"Times New Roman"}, \codeinline{'Courier New'}.
        \item Nếu tên font không có khoảng trắng (ví dụ: \codeinline{Arial}, \codeinline{Verdana}, \codeinline{Roboto}), việc có hoặc không có dấu ngoặc kép đều được chấp nhận, nhưng khuyến nghị nên dùng dấu ngoặc kép để thống nhất.
    \end{itemize}
    \item \textbf{Ưu điểm}: Cho phép bạn sử dụng chính xác font chữ mà bạn mong muốn, mang lại sự độc đáo cho thiết kế.
    \item \textbf{Nhược điểm}: Font đó cần phải có sẵn trên hệ thống của người dùng hoặc phải được nhúng vào trang web (qua Google Fonts, \codeinline{@font-face}, v.v.). Nếu không, trang web sẽ hiển thị bằng font dự phòng.
\end{itemize}
\end{conceptBox}

\subsubsection*{2. Tên Font Chung (Generic Font Families)}
Đây là các từ khóa đại diện cho một nhóm các font có đặc điểm hình dạng tương tự nhau. Chúng đóng vai trò là font dự phòng cuối cùng, đảm bảo rằng văn bản luôn hiển thị bằng một loại font dễ đọc, ngay cả khi các font cụ thể không tải được.

\begin{warnBox}[Quy Tắc Vàng Với Generic Font Family]
\textbf{Luôn luôn đặt một tên font chung ở cuối danh sách \codeinline{font-family} của bạn!}
\end{warnBox}

\paragraph{5 Nhóm Generic Font Family Phổ Biến:}
\begin{prettyTableBox}[Bảng Tra Cứu 5 Nhóm Generic Font Family]
\sffamily\small
\begin{tabularx}{\linewidth}{|>{\raggedright\arraybackslash\ttfamily\color{primaryBlue}\bfseries}p{2.6cm}|>{\raggedright\arraybackslash}p{4.5cm}|>{\raggedright\arraybackslash}X|}
\hline
\rowcolor{primaryBlue} \thd{\sffamily Loại font} & \thd{\sffamily Đặc điểm hình thái} & \thd{\sffamily Ví dụ font cụ thể tiêu biểu} \\
\hline
serif & Các font có "chân" hoặc nét trang trí ở cuối nét chữ. Giúp mắt dễ dàng theo dõi dòng chữ dài. & Times New Roman, Georgia, Palatino. \\
\hline
sans-serif & Các font không có "chân". Trông hiện đại, sạch sẽ và dễ đọc trên màn hình kỹ thuật số. & Arial, Helvetica, Verdana, Roboto, Open Sans. \\
\hline
monospace & Tất cả các ký tự có cùng chiều rộng cố định. Thường dùng cho mã lập trình hoặc căn cột. & Courier New, Consolas, Monaco. \\
\hline
cursive & Font trông giống chữ viết tay, có nét nối hoặc trang trí uốn lượn. & Comic Sans MS, Pacifico, Brush Script MT. \\
\hline
fantasy & Font mang tính trang trí cao, kiểu dáng độc đáo, không hợp cho văn bản dài. & Impact, Copperplate, Papyrus. \\
\hline
\end{tabularx}
\end{prettyTableBox}

\subsection*{III. Font An Toàn Trên Web (Web-Safe Fonts)}

Đây là các font được cài đặt sẵn trên hầu hết các hệ điều hành và thiết bị. Sử dụng chúng đảm bảo rằng trang web của bạn sẽ hiển thị nhất quán mà không cần phải nhúng thêm font từ bên ngoài.

\begin{prettyTableBox}[Bảng Tổng Hợp Các Nhóm Web-Safe Fonts Phổ Biến]
\sffamily\small
\begin{tabularx}{\linewidth}{|>{\raggedright\arraybackslash\bfseries\color{primaryBlue}}p{3.2cm}|>{\raggedright\arraybackslash\ttfamily\color{accentOrange}}X|}
\hline
\rowcolor{primaryBlue} \thd{\sffamily Phân loại} & \thd{\sffamily Danh sách chuỗi font-family web-safe chuẩn} \\
\hline
Sans-serif (Không chân) & Arial, Helvetica, sans-serif \\
 & Verdana, Geneva, sans-serif \\
 & Tahoma, Geneva, sans-serif \\
 & Trebuchet MS, Helvetica, sans-serif \\
\hline
Serif (Có chân) & Times New Roman, Times, serif \\
 & Georgia, serif \\
 & Palatino Linotype, Book Antiqua, Palatino, serif \\
\hline
Monospace (Đơn cách) & Courier New, Courier, monospace \\
 & Lucida Console, Monaco, monospace \\
\hline
\end{tabularx}
\end{prettyTableBox}

\subsection*{IV. Kỹ Thuật Nhúng Font Ngoại (Web Fonts)}

Để sử dụng các font chữ không phải là web-safe (như các font từ Google Fonts hoặc font tùy chỉnh của bạn), bạn cần nhúng chúng vào trang web. Điều này đảm bảo rằng font sẽ được tải xuống cho người dùng nếu họ không có font đó trên hệ thống của mình.

\paragraph{1. Nhúng Qua Google Fonts (Phổ Biến Nhất):}
Dễ dàng nhất. Bạn thêm một thẻ \codeinline{<link>} vào \codeinline{<head>} của HTML hoặc dùng \codeinline{@import} vào CSS.

\begin{codeWindow}[Ví Dụ Nhúng Google Fonts Và Áp Dụng]
\begin{lstlisting}[language=CSS3]
/* Sau khi nhung Roboto tu Google Fonts qua <link> hoac @import */
body {
  font-family: 'Roboto', sans-serif;
}
\end{lstlisting}
\end{codeWindow}

\paragraph{2. Nhúng Qua Quy Tắc @font-face (Tự Host Font Trên Server):}
Dùng để tự host font trên máy chủ của bạn. Bạn định nghĩa font và đường dẫn tới các tệp font (\codeinline{.woff2}, \codeinline{.woff}, \codeinline{.ttf}, v.v.).

\begin{codeWindow}[Cú Pháp Khai Báo Quy Tắc @font-face]
\begin{lstlisting}[language=CSS3]
@font-face {
  font-family: 'MyCustomFont';
  src: url('fonts/MyCustomFont.woff2') format('woff2');
  font-weight: normal;
  font-style: normal;
}

h1 {
  font-family: 'MyCustomFont', serif;
}
\end{lstlisting}
\end{codeWindow}

\begin{browserWindow}[Kết Quả: Tự Host Font Ngoại Với @font-face]
\begin{tcolorbox}[arc=4pt, colback=white, colframe=primaryBlue!30, boxrule=0.8pt, top=10pt, bottom=10pt, left=12pt, right=12pt]
  {\sffamily\bfseries\color{primaryBlue} Kết quả hiển thị phông chữ tự host trên máy chủ:}\par\vspace{4pt}
  {\sffamily Trình duyệt chủ động tải file \codeinline{MyCustomFont.woff2} từ máy chủ web và hiển thị nét chữ đúng theo ý đồ thiết kế UI.}
\end{tcolorbox}
\end{browserWindow}

\subsection*{V. Các Lưu Ý Kỹ Thuật Quan Trọng}

\begin{prettyTableBox}[Bảng Lưu Ý Quan Trọng Khi Khai Báo font-family]
\sffamily\small
\begin{tabularx}{\linewidth}{|>{\raggedright\arraybackslash\bfseries\color{primaryBlue}}p{4.2cm}|>{\raggedright\arraybackslash}X|}
\hline
\rowcolor{primaryBlue} \thd{\sffamily Lưu ý} & \thd{\sffamily Mô tả chi tiết \& Hướng xử lý} \\
\hline
Không bắt buộc dấu ngoặc cho tên font không khoảng trắng & Ví dụ: \codeinline{Arial}, không nhất thiết phải viết \codeinline{"Arial"}. Tuy nhiên khuyến nghị nên dùng dấu ngoặc kép đồng bộ. \\
\hline
Nên luôn luôn có fallback font & Tránh văn bản bị lỗi hoặc vỡ giao diện nếu font chính không tồn tại trên máy người dùng. \\
\hline
font-family không tự tải font & Khai báo \codeinline{font-family} không tự động tải font về; bạn bắt buộc phải import qua Google Fonts hoặc \codeinline{@font-face}. \\
\hline
Trình duyệt render font khác nhau trên các HĐH & Dù cùng một tên font, cách hiển thị trên Windows và macOS có thể hơi khác nhau do engine font rendering. \\
\hline
Không dùng quá nhiều font trên 1 trang & Làm tăng thời gian tải trang, gây nặng mạng và giảm tính thống nhất của thiết kế UI. \\
\hline
\end{tabularx}
\end{prettyTableBox}

\subsection*{VI. Bảng Tổng Kết Cú Pháp Cấu Trúc Thuộc Tính font-family}

\begin{prettyTableBox}[Bảng Tổng Kết Các Thành Phần Cấu Thành font-family]
\sffamily\small
\begin{tabularx}{\linewidth}{|>{\raggedright\arraybackslash\ttfamily\color{primaryBlue}\bfseries}p{3.8cm}|>{\raggedright\arraybackslash}X|}
\hline
\rowcolor{primaryBlue} \thd{\sffamily Thành phần} & \thd{\sffamily Ý nghĩa \& Vai trò} \\
\hline
"Font Name" & Font chính bạn muốn dùng (nên để trong dấu ngoặc nếu có khoảng trắng). \\
\hline
Fallback fonts & Font dự phòng thứ cấp, nên chọn những font phổ biến có sẵn trên hệ thống. \\
\hline
generic-family & Nhóm font tổng quát cuối cùng: \codeinline{serif}, \codeinline{sans-serif}, \codeinline{monospace}, v.v. \\
\hline
Cách chọn nhiều font & Khai báo viết cách nhau bằng dấu phẩy \codeinline{,} ghi theo thứ tự ưu tiên từ trái qua phải. \\
\hline
Yêu cầu dấu ngoặc kép & Bắt buộc áp dụng nếu tên font có chứa khoảng trắng (ví dụ: \codeinline{"Times New Roman"}). \\
\hline
\end{tabularx}
\end{prettyTableBox}

\subsection*{VII. Gợi Ý Các Phông Chữ Đẹp \& Phổ Biến Cho Web}

\begin{prettyTableBox}[Bảng Gợi Ý Font Chữ Theo Mục Đích Thiết Kế Web]
\sffamily\small
\begin{tabularx}{\linewidth}{|>{\raggedright\arraybackslash\bfseries\color{primaryBlue}}p{4.0cm}|>{\raggedright\arraybackslash\ttfamily\color{accentOrange}}X|}
\hline
\rowcolor{primaryBlue} \thd{\sffamily Mục đích thiết kế} & \thd{\sffamily Font chữ gợi ý hàng đầu} \\
\hline
Website hiện đại (Modern UI) & Roboto, Open Sans, Inter, Lato \\
\hline
Website cổ điển (Editorial/Classic) & Georgia, Times New Roman \\
\hline
Blog cá nhân / Đọc báo dài & Merriweather, Playfair Display \\
\hline
Lập trình / Công nghệ (Tech/Code) & Fira Code, Source Code Pro, Menlo \\
\hline
Trang trí nhẹ / Tiêu đề nghệ thuật & Pacifico, Lobster, Dancing Script \\
\hline
\end{tabularx}
\end{prettyTableBox}

\subsection*{VIII. Lời Khuyên \& Best Practices Khi Sử Dụng font-family}

\begin{tipBox}[Thực Hành Tốt Nhất (Best Practices) Cho font-family]
\begin{enumerate}[leftmargin=1.5em, itemsep=3pt, parsep=0pt]
    \item \textbf{Luôn có font dự phòng}: Đảm bảo danh sách \codeinline{font-family} của bạn kết thúc bằng một tên font chung (\codeinline{serif}, \codeinline{sans-serif}, \codeinline{monospace}, v.v.) để trang web luôn hiển thị văn bản, ngay cả khi các font ưu tiên không tải được.
    \item \textbf{Cân nhắc hiệu suất}: Nhúng quá nhiều font hoặc các font có kích thước lớn có thể làm chậm thời gian tải trang. Hãy chọn font cẩn thận và chỉ nhúng những style (đậm, nghiêng) mà bạn thực sự cần.
    \item \textbf{Tính nhất quán}: Sử dụng một số lượng font hạn chế (thường là 1--3 font) trên toàn bộ trang web để tạo sự nhất quán và chuyên nghiệp trong thiết kế.
    \item \textbf{Kiểm tra trên nhiều trình duyệt và thiết bị}: Đảm bảo font hiển thị đúng như mong đợi trên các môi trường khác nhau (Windows, macOS, iOS, Android).
    \item \textbf{Ưu tiên khả năng đọc}: Chọn font dễ đọc, đặc biệt cho các khối văn bản dài.
\end{enumerate}
\end{tipBox}

Việc làm chủ thuộc tính \codeinline{font-family} không chỉ là về việc chọn font đẹp mà còn là về việc đảm bảo trải nghiệm người dùng tối ưu và hiệu suất trang web tốt.

\subsubsection{2. Thuộc Tính font-size (Kích Thước Phông Chữ MASTERCLASS)}
\label{sec:font_size}

Thuộc tính \codeinline{font-size} trong CSS được sử dụng để thiết lập kích thước của phông chữ (văn bản) trong một phần tử HTML. Đây là một trong những thuộc tính cơ bản và quan trọng nhất trong việc kiểm soát typography trên trang web, ảnh hưởng trực tiếp đến khả năng đọc và thẩm mỹ tổng thể.

\begin{conceptBox}[Ý Nghĩa \& Mục Đích Của Thuộc Tính font-size]
\begin{itemize}[leftmargin=1.5em, itemsep=4pt, parsep=0pt]
    \item \textbf{Khả năng đọc (Readability):} Kích thước phông chữ phù hợp là yếu tố then chốt giúp người dùng đọc hiểu nội dung dễ dàng. Phông chữ quá nhỏ có thể gây khó khăn khi đọc, đặc biệt trên các thiết bị di động, trong khi phông chữ quá lớn có thể chiếm quá nhiều không gian và làm giảm mật độ thông tin.
    \item \textbf{Phân cấp trực quan (Visual Hierarchy):} Bằng cách sử dụng các kích thước phông chữ khác nhau cho tiêu đề, phụ đề, nội dung chính và chú thích, bạn có thể tạo ra một hệ thống phân cấp rõ ràng, giúp người dùng nhanh chóng nhận biết các phần quan trọng của trang.
    \item \textbf{Thẩm mỹ và Thương hiệu:} Kích thước phông chữ góp phần vào phong cách tổng thể của trang web, phản ánh tính cách và thương hiệu.
    \item \textbf{Chức năng cốt lõi:} Xác định chính xác kích thước hiển thị của chữ trên màn hình.
\end{itemize}
\end{conceptBox}

\paragraph*{I. Cú Pháp Khai Báo Standard}
\begin{codeWindow}[Cú Pháp Định Dạng font-size]
\begin{lstlisting}[language=CSS3]
selector {
  font-size: value;
}

/* Vi du khai bao co ban */
p {
  font-size: 16px;
}
\end{lstlisting}
\end{codeWindow}

\paragraph*{II. Các Nhóm Giá Trị Của font-size}
Thuộc tính \codeinline{font-size} có thể nhận nhiều loại giá trị khác nhau, được chia thành 4 nhóm chính với đặc tính và trường hợp sử dụng riêng biệt:

\begin{conceptBox}[1. Nhóm Từ Khóa Tuyệt Đối \& Tương Đối (Absolute \& Relative Keywords)]
Các từ khóa tham chiếu đến kích thước mặc định của trình duyệt. Không khuyến khích dùng rộng rãi cho văn bản chính thiết kế web hiện đại do thiếu tính linh hoạt:
\begin{itemize}[leftmargin=1.5em, itemsep=4pt, parsep=0pt]
    \item \textbf{Từ khóa tuyệt đối (Absolute Keywords)}: 
    \par\vspace{4pt}
    \badge{xx-small} \badge{x-small} \badge{small} \badge{medium (mặc định ~16px)} \badge{large} \badge{x-large} \badge{xx-large} \badge{xxx-large}
    \par\vspace{4pt}
    \item \textbf{Từ khóa tương đối (Relative Keywords)}: 
    \par\vspace{4pt}
    \badge{smaller} (Nhỏ hơn 1 cấp so với cha) \qquad \badge{larger} (Lớn hơn 1 cấp so với cha)
\end{itemize}
\end{conceptBox}

\begin{codeWindow}[Ví Dụ Sử Dụng Giá Trị Từ Khóa Của font-size]
\begin{lstlisting}[language=CSS3]
h1 { font-size: xx-large; } /* Cuc to */
p  { font-size: medium;   } /* Mac dinh ~16px */
span { font-size: smaller; } /* Nho hon phan tu cha */
\end{lstlisting}
\end{codeWindow}

\begin{conceptBox}[2. Nhóm Đơn Vị Độ Dài Tuyệt Đối (Absolute Length Units -- Cố Định)]
Xác định kích thước phông chữ cố định, không co giãn theo cài đặt người dùng hay kích thước màn hình:
\begin{itemize}[leftmargin=1.5em, itemsep=4pt, parsep=0pt]
    \item \textbf{\codeinline{px} (Pixels)}: Đơn vị phổ biến nhất ($16\text{px}$ là mặc định trình duyệt).
    \begin{tcolorbox}[arc=3pt, colback=white, colframe=primaryBlue!20, boxrule=0.5pt, top=4pt, bottom=4pt, left=8pt, right=8pt]
      \textcolor{green!60!black}{\textbf{Ưu điểm}}: Dễ kiểm soát chính xác, ổn định. \\
      \textcolor{red!80!black}{\textbf{Nhược điểm}}: Không co giãn theo cài đặt zoom chữ trình duyệt (ảnh hưởng accessibility). \\
      \textbf{Ví dụ}: \codeinline{font-size: 16px;}
    \end{tcolorbox}
    \item \textbf{\codeinline{pt} (Points)}: Chủ yếu dùng cho in ấn ($1\text{pt} = 1/72\text{ inch}$).
    \item \textbf{\codeinline{cm}, \codeinline{mm}, \codeinline{in}}: Đơn vị vật lý, không phù hợp cho màn hình hiển thị.
\end{itemize}
\end{conceptBox}

\begin{tipBox}[3. Nhóm Đơn Vị Độ Dài Tương Đối (Relative Length Units -- Khuyên Dùng Hàng Đầu)]
Các đơn vị được khuyến nghị hàng đầu trong thiết kế Web Responsive hiện đại nhờ tính linh hoạt cao:

\begin{tcolorbox}[arc=4pt, colback=white, colframe=tipBorder!30, boxrule=0.6pt, top=6pt, bottom=6pt, left=10pt, right=10pt]
  {\sffamily\bfseries\color{tipBorder} a. Đơn vị \codeinline{rem} (Root em -- Lựa chọn Top 1 cho Web Hiện Đại):}\par\vspace{3pt}
  Tương đối với \codeinline{font-size} của phần tử gốc (\codeinline{<html>}). $1\text{rem} = \text{font-size gốc}$ (thường là $16\text{px}$).
  \begin{itemize}[leftmargin=1.2em, itemsep=2pt, parsep=0pt]
      \item \textcolor{green!60!black}{\textbf{Ưu điểm}}: Đồng bộ toàn bộ website, rất dễ dự đoán và tôn trọng cài đặt người dùng.
      \item \textbf{Công thức}: $\text{Kích thước px} = \text{giá trị rem} \times \text{font-size của <html>}$.
  \end{itemize}
\end{tcolorbox}

\par\vspace{4pt}

\begin{tcolorbox}[arc=4pt, colback=white, colframe=secondaryTeal!30, boxrule=0.6pt, top=6pt, bottom=6pt, left=10pt, right=10pt]
  {\sffamily\bfseries\color{secondaryTeal} b. Đơn vị \codeinline{em} (Element em -- Tương đối theo phần tử cha):}\par\vspace{3pt}
  Tương đối với \codeinline{font-size} của phần tử cha trực tiếp. Nếu cha có \codeinline{font-size: 16px}, thì \codeinline{1.5em} $= 16 \times 1.5 = 24\text{px}$.
  \begin{itemize}[leftmargin=1.2em, itemsep=2pt, parsep=0pt]
      \item \textcolor{green!60!black}{\textbf{Ưu điểm}}: Thích hợp định dạng icon, padding nút bấm linh hoạt theo cỡ chữ của nút.
      \item \textcolor{red!80!black}{\textbf{Nhược điểm}}: Dễ bị cộng dồn chồng kích thước khi lồng nhiều cấp DOM.
  \end{itemize}
\end{tcolorbox}

\par\vspace{4pt}

\begin{tcolorbox}[arc=4pt, colback=white, colframe=accentOrange!30, boxrule=0.6pt, top=6pt, bottom=6pt, left=10pt, right=10pt]
  {\sffamily\bfseries\color{accentOrange} c. Đơn vị Viewport (\codeinline{vw}, \codeinline{vh}) \& Phần Trăm (\codeinline{\%}):}\par\vspace{3pt}
  \begin{itemize}[leftmargin=1.2em, itemsep=2pt, parsep=0pt]
      \item \textbf{\codeinline{vw} / \codeinline{vh}}: $1\text{vw} = 1\%$ chiều rộng Viewport; $1\text{vh} = 1\%$ chiều cao Viewport. Phông chữ co giãn tự động theo màn hình.
      \item \textbf{\codeinline{\%}}: Tỷ lệ phần trăm so với font-size của phần tử cha (ví dụ: $120\%$ của $20\text{px} = 24\text{px}$).
  \end{itemize}
\end{tcolorbox}
\end{tipBox}

\paragraph*{III. Bảng Tra Cứu Chi Tiết Tất Cả Các Đơn Vị Của font-size}
\begin{prettyTableBox}[Bảng Tra Cứu Toàn Diện Các Đơn Vị \& Giá Trị Của font-size]
\sffamily\small
\begin{tabularx}{\linewidth}{|>{\ttfamily\color{accentOrange}\bfseries}p{2.2cm}|>{\raggedright\arraybackslash\bfseries\color{primaryBlue}}p{2.8cm}|>{\raggedright\arraybackslash}X|>{\ttfamily\color{darkText}}p{3.2cm}|}
\hline
\rowcolor{primaryBlue} \thd{\sffamily Đơn vị} & \thd{\sffamily Ý nghĩa / Loại} & \thd{\sffamily Đặc điểm \& Cơ chế tính toán} & \thd{\sffamily Ví dụ cú pháp} \\
\hline
px & Tuyệt đối (Pixel) & Cố định, chính xác, không thay đổi theo phần tử cha hay zoom trình duyệt. & font-size: 16px; \\
\hline
pt, cm, mm, in & Tuyệt đối (Vật lý) & 1pt = 1/72 inch, thích hợp in ấn, không dùng cho web responsive. & font-size: 12pt; \\
\hline
em & Tương đối (Parent) & Tỷ lệ với font-size của phần tử cha trực tiếp (1.5em = 150\% cha). & font-size: 1.5em; \\
\hline
rem & Tương đối (Root) & Tỷ lệ với font-size của thẻ html gốc (1rem = html font-size = 16px). & font-size: 1.2rem; \\
\hline
\% & Tương đối (Percentage) & Tỷ lệ phần trăm so với font-size của cha (120\% của 16px là 19.2px). & font-size: 120\%; \\
\hline
vw / vh & Responsive (Viewport) & Tỷ lệ 1\% theo chiều rộng (vw) hoặc chiều cao (vh) màn hình trình duyệt. & font-size: 5vw; \\
\hline
keywords & Từ khóa chuẩn & xx-small đến xxx-large, smaller, larger (medium = mặc định trình duyệt ~16px). & font-size: medium; \\
\hline
clamp() & Hàm Responsive & Kết hợp min, val, max làm chữ co giãn linh hoạt giới hạn an toàn. & font-size: clamp(1rem, 2vw, 3rem); \\
\hline
\end{tabularx}
\end{prettyTableBox}

\paragraph*{IV. Minh Họa Mã Nguồn \& So Sánh Thực Tế Các Đơn Vị}

\begin{codeWindow}[Mã Nguồn HTML \& CSS Minh Họa font-size Với Các Đơn Vị]
\begin{lstlisting}[language=HTML5]
<!DOCTYPE html>
<html lang="vi">
<head>
<style>
  html {
    font-size: 16px; /* (*@ Kích thước gốc (Root font-size) cho rem @*) */
  }

  body {
    font-family: Arial, sans-serif;
    font-size: 1rem; /* (*@ Kế thừa 16px từ html @*) */
    line-height: 1.6;
  }

  /* (*@ Tiêu đề sử dụng rem (Khuyên dùng) @*) */
  h1 {
    font-size: 2.5rem; /* 2.5 * 16px = 40px */
    font-weight: bold;
  }

  h2 {
    font-size: 2rem; /* 2 * 16px = 32px */
  }

  /* (*@ Ví dụ sử dụng em (Tăng kế thừa theo cha) @*) */
  .box {
    font-size: 18px; /* (*@ font-size riêng cho box @*) */
  }
  .box p {
    font-size: 0.9em; /* (*@ 0.9 * 18px = 16.2px @*) */
  }

  /* (*@ Ví dụ sử dụng px (Cố định) @*) */
  .fixed-text {
    font-size: 14px;
  }

  /* (*@ Ví dụ sử dụng viewport units (Responsive) @*) */
  .responsive-title {
    font-size: 4vw; /* (*@ Co giãn theo width màn hình @*) */
  }

  /* (*@ Ví dụ sử dụng clamp() hiện đại @*) */
  .fluid-title {
    font-size: clamp(1.2rem, 3vw, 2.5rem);
  }
</style>
</head>
<body>
  <h1>(*@ Tiêu Đề Lớn (2.5rem = 40px) @*)</h1>
  <p>(*@ Đoạn văn bình thường (1rem = 16px). @*)</p>
  
  <div class="box">
    <p>(*@ Đoạn văn trong .box (0.9em của 18px = 16.2px). @*)</p>
  </div>
  
  <p class="fixed-text">(*@ Chữ cố định px (14px). @*)</p>
  <h2 class="responsive-title">(*@ Tiêu Đề Responsive (4vw) @*)</h2>
  <h2 class="fluid-title">(*@ Tiêu Đề Fluid clamp(1.2rem, 3vw, 2.5rem) @*)</h2>
</body>
</html>
\end{lstlisting}
\end{codeWindow}

\begin{browserWindow}[Kết Quả Hiển Thị Mô Phỏng Các Loại Đơn Vị font-size]
\begin{tcolorbox}[arc=4pt, colback=white, colframe=primaryBlue!30, boxrule=0.8pt, top=10pt, bottom=10pt, left=12pt, right=12pt]
  {\sffamily\bfseries\Huge\color{primaryBlue} 1. Tiêu Đề Lớn (\codeinline{2.5rem} / $40\text{px}$)}\par\vspace{6pt}
  {\sffamily\large\color{secondaryTeal}\bfseries 2. Tiêu Đề Phụ (\codeinline{2rem} / $32\text{px}$)}\par\vspace{6pt}
  {\sffamily\small\color{darkText} 3. Đoạn văn tiêu chuẩn (\codeinline{1rem} / $16\text{px}$): Hiển thị rõ ràng, dễ đọc trên mọi màn hình.}\par\vspace{6pt}
  {\sffamily\fontsize{9}{11}\selectfont\color{gray!80} 4. Đoạn văn nhỏ (\codeinline{0.875rem} / $14\text{px}$): Thích hợp làm chú thích hoặc thông tin phụ.}\par\vspace{6pt}
  \begin{tcolorbox}[arc=3pt, colback=primaryBlue!5, colframe=primaryBlue!30, boxrule=0.6pt, top=4pt, bottom=4pt, left=8pt, right=8pt]
    {\sffamily\bfseries\color{accentOrange} 5. Fluid Typography (\codeinline{clamp(1.2rem, 3vw, 2.5rem)}):} \\
    {\sffamily\small Chữ tự động co giãn từ $1.2\text{rem}$ ($19.2\text{px}$) trên Mobile đến $2.5\text{rem}$ ($40\text{px}$) trên Desktop!}
  \end{tcolorbox}
\end{tcolorbox}
\end{browserWindow}

\paragraph*{V. Mẹo \& Hướng Dẫn Thực Hành Tốt Nhất (Best Practices)}
\begin{tipBox}[Lời Khuyên Khi Thiết Kế Hệ Thống Typography]
\begin{itemize}[leftmargin=1.5em, itemsep=4pt, parsep=0pt]
    \item \textbf{Sử dụng rem làm đơn vị chính cho font-size}: Nó cung cấp sự linh hoạt và khả năng đáp ứng tốt nhất. Bằng cách điều chỉnh \codeinline{font-size} trên thẻ \codeinline{<html>}, bạn có thể dễ dàng thay đổi tỷ lệ của toàn bộ typography trên trang web.
    \item \textbf{Sử dụng em cho các thành phần con linh hoạt}: Đôi khi, \codeinline{em} vẫn hữu ích cho các thành phần phụ thuộc vào kích thước phông chữ của phần tử cha trực tiếp (ví dụ: kích thước biểu tượng icon, padding/margin bên trong một nút bấm dựa trên kích thước phông chữ của nút đó).
    \item \textbf{Tránh px cho văn bản chính}: Mặc dù dễ hiểu, nhưng \codeinline{px} kém linh hoạt và có thể gây ra vấn đề về khả năng tiếp cận nếu người dùng phóng to/thu nhỏ văn bản trong trình duyệt.
    \item \textbf{Sử dụng clamp() cho Responsive Typography}: Kết hợp đơn vị viewport (\codeinline{vw}) với \codeinline{rem} trong hàm \codeinline{clamp(MIN, VAL, MAX)} giúp phông chữ co giãn mượt mà mà không bao giờ bị quá nhỏ hoặc quá to.
    \item \textbf{Kiểm tra trên nhiều thiết bị}: Luôn kiểm tra giao diện trên các kích thước màn hình và thiết bị khác nhau để đảm bảo phông chữ hiển thị dễ đọc và đẹp mắt.
\end{itemize}
\end{tipBox}

\begin{conceptBox}[Quy Tắc Đặt Kích Thước Tiêu Chuẩn Cho Website]
Để giữ kích thước chữ cân đối và dễ đọc, các chuyên gia khuyến nghị hệ thống kích thước chuẩn:
\begin{itemize}[leftmargin=1.5em, itemsep=3pt, parsep=0pt]
    \item \textbf{Root (\codeinline{<html>})}: \codeinline{16px} ($100\%$)
    \item \textbf{Văn bản thân bài (Body Text)}: \codeinline{1rem} ($16\text{px}$)
    \item \textbf{Tiêu đề lớn (Heading H1)}: \codeinline{2.5rem} -- \codeinline{3rem} ($40\text{px}$ -- $48\text{px}$)
    \item \textbf{Tiêu đề phụ (Heading H2)}: \codeinline{2rem} ($32\text{px}$)
    \item \textbf{Tiêu đề H3}: \codeinline{1.5rem} ($24\text{px}$)
    \item \textbf{Văn bản chú thích nhỏ (Captions / Small Text)}: \codeinline{0.875rem} ($14\text{px}$)
\end{itemize}
\end{conceptBox}

\begin{warnBox}[Cảnh Báo Về Khả Năng Tiếp Cận (Accessibility \& WCAG 2.1)]
\begin{itemize}[leftmargin=1.5em, itemsep=3pt, parsep=0pt]
    \item \textbf{Chuẩn WCAG 1.4.4 (Resize Text)}: Đảm bảo giao diện cho phép phóng to chữ tới $200\%$ mà không bị tràn khung hoặc mất nội dung. Sử dụng đơn vị \codeinline{rem} giúp trình duyệt tự động mở rộng theo thiết lập của thị lực người dùng.
    \item \textbf{Bẫy Tự Động Zoom Trên iOS Safari}: Nếu phần tử nhập liệu \codeinline{<input>}, \codeinline{<select>}, \codeinline{<textarea>} có \codeinline{font-size < 16px}, iOS Safari sẽ tự động phóng to toàn màn hình khi người dùng chạm vào ô nhập. Luôn duy trì tối thiểu $16\text{px}$ cho ô nhập liệu trên di động.
    \item \textbf{Hạn chế Trick font-size: 62.5\%}: Kỹ thuật gán \codeinline{html \{ font-size: 62.5\%; \}} (để $1\text{rem} = 10\text{px}$) có thể làm sai lệch cỡ chữ mặc định của các widget bên thứ ba. Hãy cân nhắc dùng biến CSS (\codeinline{--step-0: 1rem}) để quản lý hệ thống typography hiện đại.
\end{itemize}
\end{warnBox}

\begin{conceptBox}[Đọc Thêm: Thuộc Tính Bổ Trợ font-size-adjust]
Khi phông chữ chính chưa tải xong hoặc không khả dụng, trình duyệt sử dụng phông chữ dự phòng (\term{fallback font}). Tuy nhiên, nếu hai phông chữ có tỷ lệ chiều cao chữ thường (\term{x-height}) khác nhau, trang web sẽ bị vỡ khung hình hoặc nhảy chữ (\term{Layout Shift}).

Thuộc tính \codeinline{font-size-adjust} giúp duy trì tỷ lệ \term{x-height} nhất quán giữa phông chính và phông dự phòng:
\begin{codeWindow}[Cú Pháp Định Dạng font-size-adjust]
\begin{lstlisting}[language=CSS3]
p {
  font-family: "Futura", Arial, sans-serif;
  font-size: 16px;
  font-size-adjust: 0.5; /* (*@ Giữ nguyên tỉ lệ x-height chuẩn @*) */
}
\end{lstlisting}
\end{codeWindow}
\end{conceptBox}

\subsubsection{3. Thuộc Tính font-weight (Trọng Lượng / Độ Đậm Phông Chữ MASTERCLASS)}
\label{sec:font_weight}

\begin{conceptBox}[Khái Niệm Cốt Lõi Về Thuộc Tính font-weight]
Thuộc tính \codeinline{font-weight} trong CSS dùng để xác định \textbf{độ đậm (trọng lượng) của chữ}. Nó giúp bạn linh hoạt biến đổi văn bản từ mỏng nhẹ, bình thường, đậm đến siêu đậm tùy thuộc vào mục đích thiết kế giao diện và phân cấp thị giác (\term{Visual Hierarchy}).

\paragraph{Cú Pháp Khai Báo Standard:}
\begin{codeWindow}[Cú Pháp Định Dạng font-weight]
\begin{lstlisting}[language=CSS3]
selector {
  font-weight: value;
}
\end{lstlisting}
\end{codeWindow}
\end{conceptBox}

\subsection*{I. Các Nhóm Giá Trị Của font-weight}

Trong CSS, \codeinline{font-weight} nhận hai nhóm giá trị chính: \textbf{Giá trị từ khóa (Keyword Values)} và \textbf{Giá trị số (Numeric Values)}.

\paragraph{1. Giá Trị Từ Khóa (Keyword Values):}

\begin{prettyTableBox}[Bảng Tra Cứu Các Giá Trị Từ Khóa Của font-weight]
\sffamily\small
\begin{tabularx}{\linewidth}{|>{\raggedright\arraybackslash\ttfamily\color{primaryBlue}\bfseries}p{2.6cm}|>{\raggedright\arraybackslash}X|>{\raggedright\arraybackslash\bfseries\color{accentOrange}}p{3.8cm}|}
\hline
\rowcolor{primaryBlue} \thd{\sffamily Giá trị từ khóa} & \thd{\sffamily Ý nghĩa \& Cảm nhận thị giác} & \thd{\sffamily Giá trị số tương đương} \\
\hline
normal & Độ đậm mặc định của font chữ. & 400 \\
\hline
bold & Chữ in đậm tiêu chuẩn. & 700 \\
\hline
bolder & Làm chữ \textbf{đậm hơn} so với phần tử cha trực tiếp. & Tùy thuộc vào phần tử cha \\
\hline
lighter & Làm chữ \textbf{nhẹ (mỏng) hơn} so với phần tử cha trực tiếp. & Tùy thuộc vào phần tử cha \\
\hline
\end{tabularx}
\end{prettyTableBox}

\paragraph{2. Giá Trị Số (Numeric Values Từ 100 Đến 900):}

CSS quy định thang độ đậm theo 9 mức số nguyên (bội số của 100 từ $100 \rightarrow 900$):

\begin{prettyTableBox}[Bảng Thang Trọng Lượng Số (100 - 900) \& Ứng Dụng Thực Tế]
\sffamily\small
\begin{tabularx}{\linewidth}{|>{\raggedright\arraybackslash\ttfamily\color{primaryBlue}\bfseries}p{1.8cm}|>{\raggedright\arraybackslash\bfseries}p{3.2cm}|>{\raggedright\arraybackslash}X|>{\raggedright\arraybackslash\color{secondaryTeal}}p{4.2cm}|}
\hline
\rowcolor{primaryBlue} \thd{\sffamily Mức số} & \thd{\sffamily Tên gọi chuẩn (Name)} & \thd{\sffamily Đánh giá mức độ đậm} & \thd{\sffamily Thường dùng cho thành phần} \\
\hline
100 & Thin & Mỏng nhất (cực mỏng) & Giao diện nhẹ nhàng, thiết kế sang trọng \\
\hline
200 & Extra Light / Ultra Light & Rất mỏng & Thẻ tiêu đề phụ trang nhã \\
\hline
300 & Light & Mỏng & Phụ đề, đoạn văn mô tả nhỏ \\
\hline
400 & Normal / Regular & Bình thường (mặc định) & Văn bản nội dung chính (Body text) \\
\hline
500 & Medium & Trung bình (hơi đậm) & Danh mục thanh điều hướng, tiêu đề nhỏ \\
\hline
600 & Semi Bold / Demi Bold & Đậm vừa & Nhấn mạnh văn bản quan trọng nhẹ \\
\hline
700 & Bold & Đậm & Tiêu đề chính (\codeinline{<h1>} - \codeinline{<h3>}), in đậm nút bấm \\
\hline
800 & Extra Bold & Rất đậm & Tiêu đề nổi bật, bài viết đặc biệt \\
\hline
900 & Black / Heavy & Siêu đậm (rất dày) & Logo, Banner, tiêu điểm thị giác mạnh \\
\hline
\end{tabularx}
\end{prettyTableBox}

\begin{warnBox}[Lưu Ý Quan Trọng Về Sự Hỗ Trợ Phông Chữ]
\textbf{Không phải phông chữ nào cũng hỗ trợ đầy đủ 9 mức trọng lượng!} Nhiều phông chữ hệ thống mặc định (như Arial, Times New Roman) chỉ có 2 mức: \codeinline{normal} (400) và \codeinline{bold} (700). Để sử dụng mượt mà các độ đậm từ 100 đến 900, bạn bắt buộc phải chọn các font đa trọng lượng (\term{Multi-weight fonts}) như: \codeinline{Roboto}, \codeinline{Poppins}, \codeinline{Inter}, \codeinline{Open Sans}, \codeinline{Montserrat}, \codeinline{Lato}.
\end{warnBox}

\subsection*{II. Ví Dụ Minh Họa Tất Cả Các Giá Trị font-weight}

\begin{codeWindow}[File style.css -- Định Kiểu Thang Độ Đậm Tất Cả Các Lớp]
\begin{lstlisting}[language=CSS3]
p.normal     { font-weight: normal; }   /* 400 */
p.bold       { font-weight: bold; }     /* 700 */
p.bolder     { font-weight: bolder; }   /* Dam hon cha */
p.lighter    { font-weight: lighter; }  /* Nhe hon cha */

p.thin       { font-weight: 100; }     /* Thin */
p.extra-light{ font-weight: 200; }     /* Extra Light */
p.light      { font-weight: 300; }     /* Light */
p.medium     { font-weight: 500; }     /* Medium */
p.semi-bold  { font-weight: 600; }     /* Semi Bold */
p.extra-bold { font-weight: 800; }     /* Extra Bold */
p.black      { font-weight: 900; }     /* Black */
\end{lstlisting}
\end{codeWindow}

\begin{codeWindow}[File index.html -- Cấu Trúc Khung Thẻ Văn Bản]
\begin{lstlisting}[language=HTML5]
<p class="normal">Chuu binh thuong (normal - 400).</p>
<p class="bold">Chuu dam (bold - 700).</p>
<p class="bolder">Chuu dam hon cha (bolder).</p>
<p class="lighter">Chuu nhat hon cha (lighter).</p>
<p class="thin">Chuu rat mong (100).</p>
<p class="extra-light">Chuu cuc nhe (200).</p>
<p class="light">Chuu mong (300).</p>
<p class="medium">Chuu trung binh (500).</p>
<p class="semi-bold">Chuu dam vua (600).</p>
<p class="extra-bold">Chuu rat dam (800).</p>
<p class="black">Chuu sieu dam (900).</p>
\end{lstlisting}
\end{codeWindow}

\begin{browserWindow}[Kết Quả Hiển Thị: Minh Họa Trực Quan Phân Cấp Sắc Độ Đậm Nhạt (font-weight)]
\begin{tcolorbox}[arc=6pt, colback=white, colframe=blue!30, boxrule=1pt, top=12pt, bottom=12pt, left=16pt, right=16pt]
  {\sffamily\bfseries\color{primaryBlue}\large Kết quả mô phỏng độ đậm nhạt từng dòng văn bản:}\par\vspace{10pt}
  \begin{tabular}{@{}l l@{}}
    {\sffamily\fontsize{13}{17}\selectfont\color{blue!80!black} Chữ bình thường (normal).}  & {\sffamily\small\color{gray!70} \codeinline{font-weight: normal;}} \\ [5pt]
    {\sffamily\bfseries\fontsize{13}{17}\selectfont\color{blue!90!black} Chữ đậm (bold).}    & {\sffamily\small\color{gray!70} \codeinline{font-weight: bold;}} \\ [5pt]
    {\sffamily\bfseries\fontsize{13}{17}\selectfont\color{blue!95!black} Chữ đậm hơn cha (bolder).}  & {\sffamily\small\color{gray!70} \codeinline{font-weight: bolder;}} \\ [5pt]
    {\sffamily\fontsize{13}{17}\selectfont\color{blue!30!white} Chữ nhạt hơn cha (lighter).} & {\sffamily\small\color{gray!70} \codeinline{font-weight: lighter;}} \\ [5pt]
    {\sffamily\fontsize{13}{17}\selectfont\color{blue!35!white} Chữ rất mỏng (100).}     & {\sffamily\small\color{gray!70} \codeinline{font-weight: 100;}} \\ [5pt]
    {\sffamily\fontsize{13}{17}\selectfont\color{blue!50!white} Chữ cực nhẹ (200).}     & {\sffamily\small\color{gray!70} \codeinline{font-weight: 200;}} \\ [5pt]
    {\sffamily\fontsize{13}{17}\selectfont\color{blue!70!black} Chữ mỏng (300).}     & {\sffamily\small\color{gray!70} \codeinline{font-weight: 300;}} \\ [5pt]
    {\sffamily\bfseries\fontsize{13}{17}\selectfont\color{blue!85!black} Chữ trung bình (500).}     & {\sffamily\small\color{gray!70} \codeinline{font-weight: 500;}} \\ [5pt]
    {\sffamily\bfseries\fontsize{13}{17}\selectfont\color{blue!95!black} Chữ đậm vừa (600).}     & {\sffamily\small\color{gray!70} \codeinline{font-weight: 600;}} \\ [5pt]
    {\sffamily\bfseries\fontsize{13}{17}\selectfont\color{blue!100!black} Chữ rất đậm (800).}     & {\sffamily\small\color{gray!70} \codeinline{font-weight: 800;}} \\ [5pt]
    {\sffamily\bfseries\fontsize{13}{17}\selectfont\color{black} Chữ siêu đậm (900).}     & {\sffamily\small\color{gray!70} \codeinline{font-weight: 900;}} \\
  \end{tabular}\par\vspace{10pt}
  {\sffamily\small\color{gray!80} \textbf{Ghi chú}: Khi áp dụng font chữ đa biến thể như \codeinline{Roboto}, \codeinline{Inter}, hoặc \codeinline{Poppins}, các mức độ đậm nhạt từ siêu mỏng (100) đến siêu đậm (900) sẽ hiển thị có sự phân cấp nét chữ cực kỳ rõ rệt.}
\end{tcolorbox}
\end{browserWindow}

\subsection*{III. Kỹ Thuật Kết Hợp font-weight Với Font Ngoại Đa Trọng Lượng}

Nếu bạn muốn sử dụng linh hoạt các độ đậm khác nhau, hãy nhúng đầy đủ biến thể \codeinline{wght} từ Google Fonts.

\begin{codeWindow}[Mã HTML Nhúng Google Fonts Đa Trọng Lượng Đầy Đủ]
\begin{lstlisting}[language=HTML5]
<!-- Nhung Google Fonts Roboto voi day du cac muc wght tu 100 den 900 -->
<link rel="stylesheet" href="https://fonts.googleapis.com/css2?family=Roboto:wght@100;200;300;400;500;600;700;800;900&display=swap">
\end{lstlisting}
\end{codeWindow}

\begin{codeWindow}[File style.css -- Áp Dụng Trọng Lượng Cho Các Thẻ Văn Bản]
\begin{lstlisting}[language=CSS3]
body {
  font-family: 'Roboto', sans-serif;
}

h1 {
  font-weight: 900; /* Sieu dam cho Tieu de lon */
}

h2 {
  font-weight: 700; /* Dam cho Tieu de phu */
}

p {
  font-weight: 400; /* Binh thuong cho Noi dung van ban */
}

span.light {
  font-weight: 300; /* Nhe cho Ghi chu nho */
}
\end{lstlisting}
\end{codeWindow}

\subsection*{IV. Sự Khác Biệt Giữa Giá Trị bolder Và bold trong CSS}

\begin{conceptBox}[So Sánh Cơ Chế Hoạt Động Của bold và bolder]
\begin{itemize}[leftmargin=1.5em, itemsep=3pt, parsep=0pt]
    \item \textbf{\codeinline{bold}}: Chỉ định độ đậm cố định tại mức \textbf{700} (không phụ thuộc vào phần tử cha).
    \item \textbf{\codeinline{bolder}}: Tính toán độ đậm tương đối dựa trên phần tử cha. Ví dụ:
    \begin{itemize}[leftmargin=1.2em, itemsep=2pt]
        \item Nếu thẻ cha có độ đậm \codeinline{400} (Normal) $\rightarrow$ Thẻ con mang \codeinline{bolder} sẽ tăng lên thành \codeinline{700} (Bold).
        \item Nếu thẻ cha có độ đậm \codeinline{700} (Bold) $\rightarrow$ Thẻ con mang \codeinline{bolder} sẽ tăng lên thành \codeinline{900} (Black).
    \end{itemize}
\end{itemize}
\end{conceptBox}

\begin{codeWindow}[Ví Dụ Minh Họa Cơ Chế Kế Thừa Tính Tương Đối Của bolder]
\begin{lstlisting}[language=CSS3]
div {
  font-weight: 400; /* Normal */
}

span {
  font-weight: bolder; /* Tu dong tang len 700 (dam hon div) */
}
\end{lstlisting}
\end{codeWindow}

\subsection*{V. Hướng Dẫn Ra Quyết Định: Khi Nào Nên Dùng Giá Trị font-weight Nào?}

\begin{prettyTableBox}[Bảng Quy Chuẩn Lựa Chọn Trọng Lượng Phông Chữ Trong Thiết Kế UI/UX]
\sffamily\small
\begin{tabularx}{\linewidth}{|>{\raggedright\arraybackslash\ttfamily\color{primaryBlue}\bfseries}p{2.8cm}|>{\raggedright\arraybackslash\bfseries}p{3.2cm}|>{\raggedright\arraybackslash}X|}
\hline
\rowcolor{primaryBlue} \thd{\sffamily Giá trị font-weight} & \thd{\sffamily Tên gọi chuẩn} & \thd{\sffamily Tình huống \& Trường hợp sử dụng tối ưu} \\
\hline
400 (normal) & Normal & Văn bản nội dung chính (Body text), đọc thông tin dài. \\
\hline
700 (bold) & Bold & Tiêu đề bài viết (\codeinline{<h1>} - \codeinline{<h3>}), nút kêu gọi hành động (CTA). \\
\hline
300 (light) & Light & Phụ đề trang trí, dòng ghi chú tác giả, khoảng thời gian đăng bài. \\
\hline
900 (black) & Black / Heavy & Logo thương hiệu, Banner quảng cáo chính, tiêu điểm gây chú ý mạnh. \\
\hline
\end{tabularx}
\end{prettyTableBox}

\begin{tipBox}[Mẹo Thiết Kế Giao Diện Chuyên Nghiệp]
Sử dụng độ đậm hợp lý giúp tạo ra sự \term{phân cấp thị giác rõ ràng (Visual Hierarchy)}, giúp mắt người đọc dễ dàng quét qua các phần quan trọng mà không bị mỏi mắt.
\end{tipBox}

\subsection*{VI. Bảng Tổng Kết Các Lưu Ý Quan Trọng Khi Dùng font-weight}

\begin{prettyTableBox}[Bảng Tổng Kết Các Lưu Ý Kỹ Thuật Cần Nhớ]
\sffamily\small
\begin{tabularx}{\linewidth}{|>{\raggedright\arraybackslash\bfseries\color{primaryBlue}}p{4.2cm}|>{\raggedright\arraybackslash}X|}
\hline
\rowcolor{primaryBlue} \thd{\sffamily Nguyên tắc / Lưu ý} & \thd{\sffamily Giải thích kỹ thuật chi tiết} \\
\hline
Không phải font nào cũng hỗ trợ nhiều trọng số & Cần kiểm tra kỹ danh sách biến thể font trên Google Fonts trước khi nhúng. \\
\hline
Kiểm tra trên nhiều hệ điều hành & Font có thể hiển thị nét chữ dày hơn nhẹ trên macOS so với Windows. \\
\hline
Chọn font có phân cấp rõ giữa 400, 500, 700 & Giúp người dùng phân biệt ngay tiêu đề và đoạn văn bản. \\
\hline
Kết hợp với font-size, color, letter-spacing & Độ đậm chữ cần đi kèm với kích thước và độ tương phản màu sắc chuẩn. \\
\hline
\end{tabularx}
\end{prettyTableBox}

\begin{conceptBox}[Tổng Kết 4 Điểm Cốt Lõi Về Thuộc Tính font-weight]
\begin{enumerate}[leftmargin=1.5em, itemsep=3pt, parsep=0pt]
    \item \codeinline{font-weight} giúp điều chỉnh độ đậm của chữ từ mỏng ($100$) đến siêu đậm ($900$).
    \item Bạn có thể sử dụng từ khóa (\codeinline{bold}, \codeinline{bolder}) hoặc giá trị số để kiểm soát chính xác.
    \item Nhúng font đa trọng lượng giúp tận dụng tối đa sức mạnh của thuộc tính này.
    \item Kết hợp thông minh độ đậm và kiểu chữ giúp nội dung nổi bật, trực quan và dễ đọc.
\end{enumerate}
\end{conceptBox}

\subsection*{VII. Chuyên Đề Xử Lý Sự Cố \& Tình Huống Thực Tế Khi Dùng font-weight}

Trong thực tế phát triển Web, bạn sẽ gặp một số tình huống sự cố làm cho thuộc tính \codeinline{font-weight} hiển thị không như mong muốn. Dưới đây là phân tích chi tiết 6 tình huống thực tế thường gặp và hướng dẫn khắc phục từng bước:

% ------------------------------------------------------------------------------
\paragraph{Tình Huống 1: Phông Chữ Không Hỗ Trợ Nhiều Độ Đậm (Limited Weight Support):}

\begin{warnBox}[Nguyên Nhân Kỹ Thuật Tình Huống 1]
Không phải font nào cũng hỗ trợ đầy đủ tất cả các mức độ đậm từ $100 \rightarrow 900$. Ví dụ, nhiều font hệ thống mặc định (như \codeinline{Arial}, \codeinline{Times New Roman}) chỉ có 2 biến thể:
\begin{itemize}[leftmargin=1.5em, itemsep=2pt, parsep=0pt]
    \item \codeinline{Normal (400)}
    \item \codeinline{Bold (700)}
\end{itemize}
Do đó, khi bạn thiết lập \codeinline{font-weight: 200;} hoặc \codeinline{font-weight: 900;} cho font Arial, trình duyệt chỉ có thể hiển thị nó giống như \codeinline{normal} hoặc \codeinline{bold}.
\end{warnBox}

\begin{tipBox}[$\rightarrow$ Giải Pháp Khắc Phục Tình Huống 1]
Chuyển sang sử dụng các bộ phông chữ hiện đại hỗ trợ đa trọng lượng (\term{Multi-weight Fonts}) từ Google Fonts như: \codeinline{Roboto}, \codeinline{Poppins}, \codeinline{Inter}, \codeinline{Open Sans}.
\end{tipBox}

\begin{codeWindow}[Mã HTML Nhúng Google Fonts Đa Trọng Lượng Đầy Đủ]
\begin{lstlisting}[language=HTML5]
<!-- Nhung Google Fonts Roboto ho tro day du cac muc wght tu 100 den 900 -->
<link rel="stylesheet" href="https://fonts.googleapis.com/css2?family=Roboto:wght@100;200;300;400;500;600;700;800;900&display=swap">
\end{lstlisting}
\end{codeWindow}

\begin{codeWindow}[File style.css -- Áp Dụng Font Đa Trọng Lượng]
\begin{lstlisting}[language=CSS3]
body {
  font-family: 'Roboto', sans-serif;
}
\end{lstlisting}
\end{codeWindow}

% ------------------------------------------------------------------------------
\paragraph{Tình Huống 2: Trình Duyệt Tự Thay Thế Phông Chữ Không Hỗ Trợ (Weight Fallback):}

\begin{conceptBox}[Cơ Chế Weight Fallback Của Trình Duyệt]
Nếu phông chữ bạn chọn không có biến thể độ đậm cụ thể mà bạn khai báo, trình duyệt sẽ tự động chọn độ đậm gần nhất có sẵn trong tệp font để hiển thị thay thế:
\begin{center}
\codeinline{font-weight: 600;} $\rightarrow$ Font không có biến thể 600? $\rightarrow$ Trình duyệt tự chọn \codeinline{700} (Bold) thay thế.
\end{center}
\end{conceptBox}

\begin{tipBox}[$\rightarrow$ Giải Pháp Khắc Phục Tình Huống 2]
Đảm bảo kiểm tra kỹ URL nhúng font và tải bổ sung đúng biến thể độ đậm cần dùng (như biến thể \codeinline{600}) hoặc chọn phông chữ khác có hỗ trợ.
\end{tipBox}

% ------------------------------------------------------------------------------
\paragraph{Tình Huống 3: Font Quá Mỏng Hoặc Quá Dày Trên Màn Hình Nhỏ (Low Resolution Displays):}

\begin{warnBox}[Ảnh Hưởng Của Màn Hình Đến Nét Chữ]
Trên một số màn hình có độ phân giải thấp hoặc máy di động kích thước nhỏ, các độ đậm cực nhẹ ($100 \rightarrow 300$) dễ bị đứt nét, mờ nhạt, trong khi các độ đậm siêu dày ($800 \rightarrow 900$) dễ bị dính nét, bết chữ.
\end{warnBox}

\begin{tipBox}[$\rightarrow$ 3 Giải Pháp Khắc Phục Tình Huống 3]
\begin{enumerate}[leftmargin=1.5em, itemsep=3pt, parsep=0pt]
    \item \textbf{Tăng kích thước chữ}: Khai báo cỡ chữ lớn hơn để nét mỏng hiển thị rõ hơn (ví dụ: \codeinline{font-size: 20px;}).
    \item \textbf{Sử dụng font tối ưu cho màn hình}: Chọn các font thiết kế chuyên biệt cho màn hình kỹ thuật số như \codeinline{Inter} hoặc \codeinline{Verdana}.
    \item \textbf{Tăng độ tương phản màu sắc}: Sử dụng màu chữ đậm tương phản mạnh với màu nền để các nét mỏng không bị chìm.
\end{enumerate}
\end{tipBox}

% ------------------------------------------------------------------------------
\paragraph{Tình Huống 4: Cú Pháp Không Đúng Hoặc Thuộc Tính CSS Bị Ghi Đè (Specificity Conflict):}

\begin{conceptBox}[Hiện Tượng Specificity Trong CSS Cascade]
Đôi khi bạn đặt thuộc tính \codeinline{font-weight} nhưng giao diện không hề thay đổi do quy tắc CSS bị ghi đè bởi một Selector có độ ưu tiên cao hơn hoặc nằm ở vị trí phía sau trong file CSS.
\end{conceptBox}

\begin{tipBox}[$\rightarrow$ Giải Pháp Khắc Phục Tình Huống 4]
Sử dụng phím \codeinline{F12} (DevTools) trên trình duyệt để kiểm tra xem dòng CSS \codeinline{font-weight} có bị gạch ngang hay không. Nếu bị ghi đè, hãy tăng độ ưu tiên của bộ chọn hoặc tạm thời dùng từ khóa \codeinline{!important}:
\end{tipBox}

\begin{codeWindow}[Sử Dụng Từ Khóa !important Để Ép Ưu Tiên]
\begin{lstlisting}[language=CSS3]
p {
  font-weight: 700 !important; /* Ep buoc uu tien do dam 700 */
}
\end{lstlisting}
\end{codeWindow}

% ------------------------------------------------------------------------------
\paragraph{Tình Huống 5: Khoảng Cách Chữ Hoặc Chiều Cao Dòng Làm Mờ Sự Khác Biệt:}

\begin{conceptBox}[Tác Động Của letter-spacing Và line-height]
Nếu khoảng cách giữa các chữ cái (\codeinline{letter-spacing}) hoặc chiều cao dòng (\codeinline{line-height}) quá lớn hoặc quá nhỏ, thị giác người xem sẽ khó phân biệt được các mức độ đậm nhạt khác nhau.
\end{conceptBox}

\begin{tipBox}[$\rightarrow$ Giải Pháp Khắc Phục Tình Huống 5]
Điều chỉnh khoảng cách ký tự và chiều cao dòng cân đối hợp lý:
\end{tipBox}

\begin{codeWindow}[Thiết Lập Tối Ưu Khoảng Cách Cho Phông Chữ Đậm]
\begin{lstlisting}[language=CSS3]
.custom-heading {
  font-weight: 700;
  letter-spacing: 0.5px; /* Khoang cach ky tu Vua phai */
  line-height: 1.6;      /* Chieu cao dong Chuan 1.6 */
}
\end{lstlisting}
\end{codeWindow}

% ------------------------------------------------------------------------------
\paragraph{Tình Huống 6: Một Số Phông Chữ Được Tối Ưu Hóa Để Ít Khác Biệt Nét Chữ:}

\begin{warnBox}[Đặc Điểm Thiết Kế Của Font Cổ Điển]
Có những font chữ được thiết kế với độ dày ít thay đổi giữa các mức trọng lượng, ngay cả khi bạn thay đổi thuộc tính \codeinline{font-weight}. Ví dụ: Font \codeinline{Times New Roman} không có sự chênh lệch nét chữ quá rõ rệt giữa hai mức \codeinline{500} và \codeinline{600}.
\end{warnBox}

\begin{tipBox}[$\rightarrow$ Giải Pháp Khắc Phục Tình Huống 6]
Chuyển sang sử dụng các phông chữ thiết kế hiện đại có sự phân cấp độ đậm cực kỳ rõ ràng như: \codeinline{Poppins}, \codeinline{Montserrat}, hoặc \codeinline{Lato}.
\end{tipBox}

% ------------------------------------------------------------------------------
\begin{tipBox}[Mẹo Nhanh Để Thấy Sự Khác Biệt Độ Đậm Rõ Nhất]
\begin{enumerate}[leftmargin=1.5em, itemsep=3pt, parsep=0pt]
    \item \textbf{Chọn font đa biến thể}: Luôn chọn phông chữ hỗ trợ nhiều biến thể trọng lượng (ví dụ: \codeinline{Roboto}, \codeinline{Inter}).
    \item \textbf{Tăng kích thước chữ khi thử nghiệm}: Dùng kích thước chữ lớn hơn (\codeinline{font-size: 24px;}) để dễ quan sát chi tiết nét chữ trên màn hình.
    \item \textbf{Kiểm thử đa trình duyệt}: Test trên các trình duyệt khác nhau (Chrome, Firefox, Safari) hoặc trên màn hình chất lượng cao.
    \item \textbf{Đảm bảo không bị ghi đè}: Mở DevTools (\codeinline{F12}) để xác nhận quy tắc \codeinline{font-weight} đang thực sự được áp dụng.
\end{enumerate}
\end{tipBox}

\subsubsection{4. Thuộc Tính font-style (Kiểu Phông Chữ Nghiêng MASTERCLASS)}
\label{sec:font_style}

Thuộc tính \codeinline{font-style} trong CSS được sử dụng để xác định kiểu dáng của chữ (văn bản), chủ yếu là chuyển đổi giữa chữ thường (\codeinline{normal}), chữ nghiêng nghệ thuật (\codeinline{italic}) hoặc chữ nghiêng mô phỏng (\codeinline{oblique}). Nó giúp bạn thêm sắc thái và sự nhấn mạnh vào nội dung mà không cần phải thay đổi họ phông chữ (\codeinline{font-family}).

\begin{conceptBox}[Ý Nghĩa \& Mục Đích Của Thuộc Tính font-style]
\begin{itemize}[leftmargin=1.5em, itemsep=4pt, parsep=0pt]
    \item \textbf{Nhấn mạnh \& Tạo điểm nhấn}: Khai báo \codeinline{font-style: italic;} là cách phổ biến nhất để nhấn mạnh một từ, cụm từ, hoặc đoạn trích ngắn, thường dùng cho tên sách, tên tác phẩm, phim ảnh, thuật ngữ kỹ thuật, hoặc lời trích dẫn.
    \item \textbf{Phân biệt loại nội dung}: Giúp phân biệt các loại nội dung đặc thù trên trang web như chú thích hình ảnh, lời thoại nhân vật, hoặc đoạn ghi chú nằm ngoài văn bản chính.
    \item \textbf{Thẩm mỹ Typography}: Góp phần tạo sự đa dạng, hấp dẫn thị giác và tăng nét trang trọng cho thiết kế tổng thể.
\end{itemize}
\end{conceptBox}

\paragraph*{I. Cú Pháp Khai Báo Standard}
\begin{codeWindow}[Cú Pháp Định Dạng font-style]
\begin{lstlisting}[language=CSS3]
selector {
  font-style: value;
}

/* Vi du khai bao co ban */
p {
  font-style: italic;
}
\end{lstlisting}
\end{codeWindow}

\paragraph*{II. Các Nhóm Giá Trị Của font-style}

\begin{prettyTableBox}[Bảng Tra Cứu Toàn Diện Các Giá Trị Của font-style \& Độ Hỗ Trợ Trình Duyệt]
\sffamily\small
\begin{tabularx}{\linewidth}{|>{\ttfamily\color{accentOrange}\bfseries}p{2.2cm}|>{\raggedright\arraybackslash}X|>{\raggedright\arraybackslash\bfseries\color{green!60!black}}p{3.8cm}|}
\hline
\rowcolor{primaryBlue} \thd{\sffamily Giá trị} & \thd{\sffamily Ý nghĩa \& Cơ chế hiển thị kỹ thuật} & \thd{\sffamily Hỗ trợ trình duyệt} \\
\hline
normal & Chữ thẳng đứng bình thường, không có hiệu ứng nghiêng (giá trị mặc định cho hầu hết các font). & \textcolor{green!60!black}{\textbf{Đầy đủ}} (Chrome, Firefox, Safari, Edge) \\
\hline
italic & Chữ nghiêng mượt mà (True Italic), trình duyệt sử dụng phiên bản nghiêng do nhà thiết kế font vẽ riêng. & \textcolor{green!60!black}{\textbf{Đầy đủ}} (Chrome, Firefox, Safari, Edge) \\
\hline
oblique & Chữ nghiêng cưỡng ép (Synthetic Skewing), trình duyệt tự kéo nghiêng ký tự thẳng bằng thuật toán. Hỗ trợ kèm góc chỉ định (ví dụ: \codeinline{oblique 14deg}). & \textcolor{green!60!black}{\textbf{Đầy đủ}} (Chrome, Firefox, Safari, Edge) \\
\hline
inherit & Kế thừa giá trị kiểu chữ nghiêng từ phần tử cha trực tiếp. & \textcolor{green!60!black}{\textbf{Đầy đủ}} (Chrome, Firefox, Safari, Edge) \\
\hline
initial & Đặt lại kiểu chữ về giá trị mặc định ban đầu của CSS (\codeinline{normal}). & \textcolor{green!60!black}{\textbf{Đầy đủ}} (Chrome, Firefox, Safari, Edge) \\
\hline
unset & Gỡ bỏ kiểu chữ đã đặt, trở về giá trị mặc định hoặc kế thừa tùy theo ngữ cảnh DOM. & \textcolor{green!60!black}{\textbf{Đầy đủ}} (Chrome, Firefox, Safari, Edge) \\
\hline
\end{tabularx}
\end{prettyTableBox}

\paragraph*{III. Phân Biệt Chi Tiết Giữa italic Và oblique}

Sự khác biệt quan trọng nhất giữa \codeinline{italic} và \codeinline{oblique} nằm ở \textbf{cách thức ký tự được tạo ra}:

\begin{prettyTableBox}[Bảng So Sánh Chi Tiết Giữa True Italic Và Synthetic Oblique]
\sffamily\small
\begin{tabularx}{\linewidth}{|>{\raggedright\arraybackslash\bfseries\color{primaryBlue}}p{2.6cm}|>{\raggedright\arraybackslash}X|>{\raggedright\arraybackslash}X|}
\hline
\rowcolor{primaryBlue} \thd{\sffamily Tiêu chí} & \thd{\sffamily Giá trị italic (True Italic)} & \thd{\sffamily Giá trị oblique (Synthetic Oblique)} \\
\hline
Nguồn gốc ký tự & Là một bản thiết kế phông chữ riêng biệt, được nhà tạo phông vẽ tay các đường uốn lượn tinh tế. & Trình duyệt tự nghiêng phông chữ thẳng đứng bằng cách lệch góc ký tự cơ học. \\
\hline
Chất lượng thị giác & Chữ rất mượt mà, nghệ thuật, đẹp và cân đối tuyệt đối. & Chữ có thể hơi thô, méo hoặc mất cân bằng nét do kéo nghiêng máy móc. \\
\hline
Yêu cầu file font & Phụ thuộc vào việc file phông chữ có chứa biến thể Italic hay không. & Không phụ thuộc, trình duyệt luôn tự tạo được chữ nghiêng cho mọi phông chữ. \\
\hline
Tùy chỉnh góc & Không tùy chỉnh được góc nghiêng (theo thiết kế font). & Cho phép tùy chỉnh góc nghiêng cụ thể (ví dụ: \codeinline{oblique 10deg}, \codeinline{oblique 14deg}). \\
\hline
\end{tabularx}
\end{prettyTableBox}

\begin{conceptBox}[Cơ Chế Fallback Tự Động Của Trình Duyệt Khi Khai Báo italic]
Khi bạn khai báo \codeinline{font-style: italic;}:
\begin{itemize}[leftmargin=1.5em, itemsep=3pt, parsep=0pt]
    \item Nếu phông chữ có phiên bản \codeinline{italic} thực sự, trình duyệt sẽ ưu tiên sử dụng bản \codeinline{italic} mượt mà đó.
    \item Nếu phông chữ không có phiên bản \codeinline{italic} thực sự, trình duyệt sẽ tự động cố gắng tạo hiệu ứng nghiêng bằng cách "làm xiên" phiên bản \codeinline{normal} (tức là tự chuyển sang cơ chế tương tự \codeinline{oblique}).
\end{itemize}
Do đó, các nhà thiết kế Web luôn được khuyến nghị khai báo \codeinline{font-style: italic;} và để trình duyệt tự động đưa ra quyết định tốt nhất!
\end{conceptBox}

\paragraph*{IV. Minh Họa Mã Nguồn HTML/CSS \& Kết Quả Hiển Thị}

\begin{codeWindow}[File style.css -- Áp Dụng Tất Cả Giá Trị Của font-style]
\begin{lstlisting}[language=CSS3]
/* (*@ Văn bản normal @*) */
p.normal  { font-style: normal; }

/* (*@ Văn bản nghiêng Italic thực sự @*) */
p.italic  { font-style: italic; }

/* (*@ Văn bản nghiêng Oblique cưỡng ép @*) */
p.oblique { font-style: oblique; }

/* (*@ Văn bản nghiêng Oblique với góc custom 14 độ @*) */
p.skewed  { font-style: oblique 14deg; }

/* (*@ Các từ khóa điều khiển DOM @*) */
p.inherit { font-style: inherit; }
p.initial { font-style: initial; }
p.unset   { font-style: unset; }
\end{lstlisting}
\end{codeWindow}

\begin{codeWindow}[File index.html -- Khung Thẻ HTML Minh Họa]
\begin{lstlisting}[language=HTML5]
<p class="normal">(*@ Chữ bình thường (normal). @*)</p>
<p class="italic">(*@ Chữ nghiêng nhẹ mượt mà (italic). @*)</p>
<p class="oblique">(*@ Chữ nghiêng cưỡng ép trình duyệt (oblique). @*)</p>
<p class="skewed">(*@ Chữ nghiêng tùy chỉnh góc (oblique 14deg). @*)</p>
<p class="inherit">(*@ Chữ kế thừa từ phần tử cha (inherit). @*)</p>
<p class="initial">(*@ Chữ trở về mặc định (initial). @*)</p>
<p class="unset">(*@ Chữ về mặc định hoặc kế thừa (unset). @*)</p>
\end{lstlisting}
\end{codeWindow}

\begin{browserWindow}[Kết Quả Hiển Thị Trực Quan Các Giá Trị font-style]
\begin{tcolorbox}[arc=4pt, colback=white, colframe=primaryBlue!30, boxrule=0.8pt, top=10pt, bottom=10pt, left=12pt, right=12pt]
  {\sffamily\bfseries\color{primaryBlue} Kết quả hiển thị mô phỏng thực tế trên trình duyệt:}\par\vspace{6pt}
  \begin{tabular}{@{}l l@{}}
    {\sffamily\fontsize{12}{15}\selectfont Chữ bình thường (normal).} & {\sffamily\small\color{gray!70} \codeinline{font-style: normal;}} \\ [5pt]
    {\sffamily\itshape\fontsize{12}{15}\selectfont Chữ nghiêng nhẹ mượt mà (italic).} & {\sffamily\small\color{gray!70} \codeinline{font-style: italic;}} \\ [5pt]
    {\sffamily\slshape\fontsize{12}{15}\selectfont Chữ nghiêng cưỡng ép (oblique).} & {\sffamily\small\color{gray!70} \codeinline{font-style: oblique;}} \\ [5pt]
    {\sffamily\slshape\fontsize{12}{15}\selectfont Chữ nghiêng góc $14^\circ$ (oblique 14deg).} & {\sffamily\small\color{gray!70} \codeinline{font-style: oblique 14deg;}} \\ [5pt]
    {\sffamily\fontsize{12}{15}\selectfont Chữ kế thừa từ cha (inherit).} & {\sffamily\small\color{gray!70} \codeinline{font-style: inherit;}} \\ [5pt]
    {\sffamily\fontsize{12}{15}\selectfont Chữ trở về mặc định (initial).} & {\sffamily\small\color{gray!70} \codeinline{font-style: initial;}} \\
  \end{tabular}
\end{tcolorbox}
\end{browserWindow}

\paragraph*{V. Hướng Dẫn Nhúng Font Có Hỗ Trợ Italic Chuẩn \& Khắc Phục Lỗi}

Để đảm bảo chữ nghiêng hiển thị sắc nét và mượt mà nhất, bạn cần nhúng đầy đủ biến thể \codeinline{italic} từ Google Fonts hoặc file ngoài.

\begin{codeWindow}[Mã HTML Nhúng Google Fonts Chuẩn Biến Thể Italic]
\begin{lstlisting}[language=HTML5]
<!-- Nhung font Roboto voi ca bien the normal (400) va italic (ital,wght@1,400) -->
<link rel="stylesheet" href="https://fonts.googleapis.com/css2?family=Roboto:ital,wght@0,400;1,400&display=swap">
\end{lstlisting}
\end{codeWindow}

\begin{warnBox}[Ví Dụ Lỗi Thường Gặp \& Cách Khắc Phục Chuẩn]
\textbf{\color{red!80!black}Ví dụ lỗi thường gặp}: Bạn dùng phông chữ tùy chỉnh nhưng quên không tải file biến thể nghiêng:
\begin{codeWindow}[Mã CSS Lỗi -- Không Có File Font Italic]
\begin{lstlisting}[language=CSS3]
.custom-text {
  font-family: 'MyCustomFont';
  font-style: italic; /* (*@ Nhưng 'MyCustomFont' không có bản Italic! @*) */
}
/* (*@ KẾT QUẢ: Trình duyệt tự nghiêng (oblique) làm chữ bị méo và xấu. @*) */
\end{lstlisting}
\end{codeWindow}

\textbf{\color{green!60!black}Cách khắc phục chuẩn với @font-face}:
\begin{codeWindow}[Mã CSS Khắc Phục Chuẩn Qua @font-face]
\begin{lstlisting}[language=CSS3]
/* Khai bao file font nghieng rieng */
@font-face {
  font-family: 'MyCustomFont';
  src: url('fonts/MyCustomFont-Italic.woff2') format('woff2');
  font-style: italic; /* Gan kieu nghieng cho file nay */
}

.custom-text {
  font-family: 'MyCustomFont';
  font-style: italic; /* hien thi dep mat tu file WOFF2 nghieng */
}
\end{lstlisting}
\end{codeWindow}
\end{warnBox}

\paragraph*{VI. Các Thẻ HTML Mặc Định Mang Kiểu font-style: italic}

Trong HTML, có một số thẻ mang ngữ nghĩa chuyên biệt được trình duyệt gán mặc định kiểu chữ nghiêng:

\begin{prettyTableBox}[Danh Sách Thẻ HTML Mặc Định Mang Style Italic]
\sffamily\small
\begin{tabularx}{\linewidth}{|>{\ttfamily\color{primaryBlue}\bfseries}p{2.0cm}|>{\raggedright\arraybackslash}X|>{\centering\arraybackslash\ttfamily\color{accentOrange}}p{3.2cm}|}
\hline
\rowcolor{primaryBlue} \thd{\sffamily Thẻ HTML} & \thd{\sffamily Ý nghĩa ngữ nghĩa chuẩn (Semantic Meaning)} & \thd{\sffamily Style mặc định} \\
\hline
\codeinline{<em>} & Nhấn mạnh văn bản mang tính ngữ nghĩa (\term{Emphasis}). & font-style: italic; \\
\hline
\codeinline{<i>} & Thể hiện thuật ngữ, ý nghĩ nhân vật hoặc tên gọi mang tính hình thức. & font-style: italic; \\
\hline
\codeinline{<cite>} & Trích dẫn tên tác phẩm nghệ thuật, tên sách, tên bài báo hoặc phim. & font-style: italic; \\
\hline
\codeinline{<dfn>} & Định nghĩa thuật ngữ mới xuất hiện lần đầu trong văn bản. & font-style: italic; \\
\hline
\end{tabularx}
\end{prettyTableBox}

\paragraph*{VII. Ứng Dụng Thực Tế \& Lời Khuyên (Best Practices)}

\begin{tipBox}[Ứng Dụng Thực Tế Của font-style]
\begin{itemize}[leftmargin=1.5em, itemsep=4pt, parsep=0pt]
    \item \textbf{Nhấn mạnh từ hoặc câu trong đoạn văn}:
    \par\vspace{2pt}
    \codeinline{em \{ font-style: italic; \}} $\rightarrow$ Khi dùng thẻ \codeinline{<em>}, văn bản sẽ nghiêng nhẹ để thu hút sự chú ý.
    \item \textbf{Định dạng khối trích dẫn (Blockquote)}:
    \par\vspace{2pt}
    \codeinline{blockquote \{ font-style: italic; \}} $\rightarrow$ Giúp phần trích dẫn nổi bật và phân biệt với nội dung chính.
    \item \textbf{Giao diện nút bấm sáng tạo}:
    \par\vspace{2pt}
    \codeinline{button \{ font-style: oblique 10deg; \}} $\rightarrow$ Nút bấm nghiêng nhẹ tạo nét cá tính độc đáo cho UI.
\end{itemize}
\end{tipBox}

\begin{conceptBox}[Tổng Kết Các Điểm Cốt Lõi Về font-style]
\begin{itemize}[leftmargin=1.5em, itemsep=3pt, parsep=0pt]
    \item \codeinline{font-style} giúp định dạng kiểu chữ dễ dàng (\codeinline{normal}, \codeinline{italic}, \codeinline{oblique}).
    \item \codeinline{italic} đẹp và nghệ thuật hơn nhưng cần file phông chữ hỗ trợ; \codeinline{oblique} linh hoạt hơn nhưng có thể kém tinh tế.
    \item \codeinline{inherit}, \codeinline{initial}, \codeinline{unset} giúp kiểm soát kiểu chữ linh hoạt theo từng ngữ cảnh DOM.
    \item \textbf{Lưu ý đọc}: Tránh sử dụng \codeinline{italic} cho các đoạn văn bản quá dài vì chữ nghiêng liên tục sẽ làm giảm khả năng đọc và gây mỏi mắt người dùng. Chỉ nên dùng để nhấn mạnh các cụm từ ngắn gọn.
\end{itemize}
\end{conceptBox}

\subsubsection{5. Thuộc Tính font-variant (Biến Thể Phông Chữ Small-Caps)}
\label{sec:font_variant}

Thuộc tính \codeinline{font-variant} dùng để đặt các biến thể chữ hoa nhỏ (\term{small-caps}).

\begin{conceptBox}[Giá Trị Của font-variant]
\begin{itemize}[leftmargin=1.5em, itemsep=3pt, parsep=0pt]
    \item \codeinline{normal}: Văn bản bình thường.
    \item \codeinline{small-caps}: Tất cả các chữ cái viết thường sẽ được hiển thị dưới dạng chữ hoa nhưng với kích thước nhỏ hơn so với các chữ hoa ban đầu.
\end{itemize}
\end{conceptBox}

\begin{codeWindow}[Ví Dụ Sử Dụng Thuộc Tính font-variant]
\begin{lstlisting}[language=CSS3]
.small-caps-text {
  font-variant: small-caps; /* Hien thi chu hoa nho (Small Caps) */
}
\end{lstlisting}
\end{codeWindow}

\begin{browserWindow}[Kết Quả: Hiển Thị font-variant Small-Caps]
\begin{tcolorbox}[arc=4pt, colback=white, colframe=primaryBlue!30, boxrule=0.8pt, top=10pt, bottom=10pt, left=12pt, right=12pt]
  {\sffamily\bfseries\color{primaryBlue} So sánh chữ thường và Small-Caps:}\par\vspace{6pt}
  \begin{tabular}{@{}l l@{}}
    {\sffamily\fontsize{14}{18}\selectfont Đây là chữ thường bình thường} & {\sffamily\small\color{gray!70} font-variant: normal} \\ [8pt]
    {\sffamily\fontsize{14}{18}\selectfont\textsc{Đây là chữ Small-Caps}} & {\sffamily\small\color{gray!70} font-variant: small-caps} \\
  \end{tabular}\par\vspace{6pt}
  {\sffamily\small\color{gray!70} → Chữ viết thường được đổi thành dạng CHỮ HOA nhỏ hơn, tạo hiệu ứng trang trọng.}
\end{tcolorbox}
\end{browserWindow}

\subsubsection{6. Thuộc Tính line-height (Chiều Cao Dòng \& Khoảng Cách Dòng MASTERCLASS)}
\label{sec:line_height}

\paragraph*{I. Khái Niệm Cốt Lõi Về line-height}

Thuộc tính \codeinline{line-height} trong CSS xác định \textbf{chiều cao của một dòng văn bản}. Nó quyết định trực tiếp:
\begin{itemize}[leftmargin=1.5em, itemsep=3pt, parsep=0pt]
    \item Khoảng cách từ dòng trên đến dòng dưới trong cùng một khối văn bản.
    \item Độ cao mà mỗi dòng chữ chiếm giữ trên giao diện trang web.
    \item \textbf{Không làm thay đổi kích thước chữ} (\codeinline{font-size}), mà tạo không gian đứng (\term{vertical space}) cho mỗi dòng chữ.
\end{itemize}

\begin{conceptBox}[Hình Dung Đơn Giản Về Hộp Dòng (Line Box)]
Hãy tưởng tượng mỗi dòng chữ nằm gọn trong một chiếc hộp vô hình -- \codeinline{line-height} chính là \textbf{chiều cao của chiếc hộp đó}.
\par\vspace{4pt}
Ví dụ với khai báo:
\begin{codeWindow}[Ví Dụ Khai Báo line-height Cơ Bản]
\begin{lstlisting}[language=CSS3]
p {
  font-size: 16px;
  line-height: 24px;
}
\end{lstlisting}
\end{codeWindow}
$\rightarrow$ Dòng chữ có kích thước $16\text{px}$, nhưng mỗi dòng sẽ chiếm tổng cộng $24\text{px}$ chiều cao trên màn hình (có $8\text{px}$ khoảng trống được chia đều: $4\text{px}$ phía trên và $4\text{px}$ phía dưới chữ).
\end{conceptBox}

\paragraph*{II. Cơ Chế Hoạt Động Chi Tiết Của line-height (Line Box \& Leading)}

Mỗi dòng văn bản được chứa trong một hộp dòng (\term{line box}). Khoảng không gian thừa giữa \codeinline{line-height} và \codeinline{font-size} được gọi là \term{leading}. Trình duyệt sẽ tự động chia đôi lượng leading này để thêm vào phía trên và phía dưới văn bản.

\begin{browserWindow}[Sơ Đồ Minh Họa Trực Quan Cấu Trúc Hộp Dòng (Line Box \& Leading)]
\begin{tcolorbox}[arc=6pt, colback=white, colframe=primaryBlue!30, boxrule=1pt, top=12pt, bottom=12pt, left=8pt, right=8pt]
  \centering
  \begin{tikzpicture}[x=1cm, y=1cm, >=stealth]
    % Outer Line Box boundary (Width 11.8cm, Height 4.4cm)
    \draw[line width=1.5pt, primaryBlue, fill=accentOrange!8, rounded corners=8pt] (0,0) rectangle (11.8, 4.4);
    
    % Half-leading top region
    \fill[accentOrange!25, rounded corners=5pt] (0.2, 3.0) rectangle (11.6, 4.1);
    \node[anchor=west, font=\sffamily\small\bfseries\color{accentOrange!90!black}] at (0.4, 3.55) {Half-leading (Khoảng trống phía trên)};
    \node[anchor=east, font=\sffamily\small\bfseries\color{gray!80!black}] at (11.4, 3.55) {+4px (Leading / 2)};

    % Content area (font-size)
    \fill[primaryBlue!15, rounded corners=5pt] (0.2, 1.4) rectangle (11.6, 2.8);
    \node[anchor=west, font=\sffamily\fontsize{11}{13}\selectfont\bfseries\color{primaryBlue}] at (0.4, 2.1) {Văn bản hiển thị (Content Area)};
    \node[anchor=east, font=\sffamily\fontsize{11}{13}\selectfont\bfseries\color{primaryBlue}] at (11.4, 2.1) {font-size = 16px};

    % Half-leading bottom region
    \fill[accentOrange!25, rounded corners=5pt] (0.2, 0.3) rectangle (11.6, 1.2);
    \node[anchor=west, font=\sffamily\small\bfseries\color{accentOrange!90!black}] at (0.4, 0.75) {Half-leading (Khoảng trống phía dưới)};
    \node[anchor=east, font=\sffamily\small\bfseries\color{gray!80!black}] at (11.4, 0.75) {+4px (Leading / 2)};

    % Height Indicator Arrow on the left
    \draw[<->, line width=1.5pt, primaryBlue] (-0.4, 0) -- (-0.4, 4.4);
    \node[rotate=90, anchor=center, font=\sffamily\small\bfseries\color{primaryBlue}] at (-0.8, 2.2) {line-height = 24px};
  \end{tikzpicture}
  \par\vspace{8pt}
  {\sffamily\small\color{gray!70} \textbf{Giải thích}: Khi \codeinline{line-height: 24px} và \codeinline{font-size: 16px}, lượng dư $8\text{px}$ được trình duyệt tự động chia đôi ($4\text{px}$ phía trên và $4\text{px}$ phía dưới) giúp chữ luôn nằm chính giữa hộp dòng.}
\end{tcolorbox}
\end{browserWindow}

\begin{conceptBox}[Công Thức Tính Leading Chi Tiết]
\begin{center}
\codeinline{Leading = line-height - font-size}\\[4pt]
\codeinline{Half-leading (trên) = Leading / 2}\\[4pt]
\codeinline{Half-leading (dưới) = Leading / 2}
\end{center}
\end{conceptBox}

\paragraph*{III. Các Cách Khai Báo \& Bảng So Sánh Các Kiểu Giá Trị}

\begin{prettyTableBox}[Bảng Tổng Quan Các Cách Khai Báo line-height]
\sffamily\small
\begin{tabularx}{\linewidth}{|>{\raggedright\arraybackslash\bfseries\color{primaryBlue}}p{3.0cm}|>{\ttfamily\color{accentOrange}\bfseries}p{2.6cm}|>{\raggedright\arraybackslash}X|}
\hline
\rowcolor{primaryBlue} \thd{\sffamily Cách viết} & \thd{\sffamily Ví dụ} & \thd{\sffamily Ý nghĩa \& Đặc tính kỹ thuật} \\
\hline
Số tuyệt đối (px) & line-height: 24px; & Dòng cao cố định đúng 24px (không tự thay đổi khi đổi font-size). \\
\hline
Số tương đối (em) & line-height: 1.5em; & Chiều cao dòng bằng 1.5 lần font-size của chính phần tử đó. \\
\hline
Số không đơn vị (tỉ lệ) & line-height: 1.5; & Bằng 1.5 lần font-size. \textbf{Đây là cách được khuyên dùng nhất} vì kế thừa linh hoạt cho các con. \\
\hline
Phần trăm (\%) & line-height: 150\%; & Chiều cao dòng bằng 150\% font-size. Giá trị đã tính sẽ được kế thừa bởi phần tử con. \\
\hline
\end{tabularx}
\end{prettyTableBox}

\begin{tipBox}[Lời Khuyên Hàng Đầu]
Khai báo \codeinline{line-height: 1.5;} (số không đơn vị) là cách thông dụng, an toàn và chuyên nghiệp nhất vì nó luôn giữ tỉ lệ chuẩn theo \codeinline{font-size} của mọi phần tử con.
\end{tipBox}

\paragraph*{IV. Kế Thừa Khi Phần Tử KHÔNG Khai Báo font-size}

\begin{conceptBox}[Phần Tử Không Khai Báo font-size Thì Lấy Giá Trị Ở Đâu?]
Trong trường hợp phần tử có khai báo \codeinline{line-height: <number>} nhưng \textbf{không có \codeinline{font-size} riêng}, nó sẽ:
\begin{enumerate}[leftmargin=1.5em, itemsep=3pt, parsep=0pt]
    \item Kế thừa \codeinline{font-size} từ phần tử cha gần nhất (hoặc lấy mặc định $16\text{px}$ từ trình duyệt).
    \item Sau đó, giá trị \codeinline{line-height: <number>} được tính toán dựa trên chính \codeinline{font-size} kế thừa đó.
\end{enumerate}
\end{conceptBox}

\begin{codeWindow}[Ví Dụ Thực Tế Kế Thừa line-height Với Số Không Đơn Vị]
\begin{lstlisting}[language=HTML5]
<style>
  body {
    font-size: 20px; /* Font-size goc tai body */
  }
  .box {
    line-height: 1.5; /* (*@ Khong khai bao font-size o day @*) */
  }
</style>
<div class="box">(*@ Van ban thu nghiem trong .box @*)</div>
\end{lstlisting}
\end{codeWindow}

\begin{conceptBox}[Giải Thích Chi Tiết Ví Dụ Kế Thừa]
\begin{itemize}[leftmargin=1.5em, itemsep=3pt, parsep=0pt]
    \item Phần tử \codeinline{.box} không khai báo \codeinline{font-size}, nên nó kế thừa \codeinline{font-size: 20px} từ \codeinline{body}.
    \item Khai báo \codeinline{line-height: 1.5} được tính toán là: $1.5 \times 20\text{px} = 30\text{px}$.
\end{itemize}
\end{conceptBox}

\begin{prettyTableBox}[Bảng Tóm Tắt Quy Tắc Tính line-height Theo font-size]
\sffamily\small
\begin{tabularx}{\linewidth}{|>{\raggedright\arraybackslash\bfseries\color{primaryBlue}}p{5.0cm}|>{\raggedright\arraybackslash}X|}
\hline
\rowcolor{primaryBlue} \thd{\sffamily Ngữ cảnh phần tử} & \thd{\sffamily Căn cứ tính toán của line-height: <number>} \\
\hline
Phần tử có \codeinline{font-size} riêng & Dựa vào \codeinline{font-size} của chính phần tử đó. \\
\hline
Phần tử không có \codeinline{font-size} & Dựa vào \codeinline{font-size} kế thừa từ phần tử cha gần nhất. \\
\hline
\end{tabularx}
\end{prettyTableBox}

\paragraph*{V. Phân Tích Chuyên Sâu Các Ví Dụ Thực Tế Của line-height}

\begin{codeWindow}[Ví Dụ Thực Tế 1: Phân Tích Khai Báo line-height: 100px]
\begin{lstlisting}[language=HTML5]
<h1>
  (*@ Write, edit and run HTML, CSS and JavaScript code online. @*)
</h1>
<style>
  h1 {
    line-height: 100px;
  }
</style>
\end{lstlisting}
\end{codeWindow}

\begin{conceptBox}[Tác Dụng Thật Sự Của Khai Báo line-height: 100px]
Khai báo này tăng chiều cao dòng chữ lên $100\text{px}$ (khoảng cách từ mép trên đến mép dưới của dòng là $100\text{px}$, chứ không làm chữ to lên $100\text{px}$):
\begin{enumerate}[leftmargin=1.5em, itemsep=3pt, parsep=0pt]
    \item \textbf{Nếu \codeinline{h1} chỉ có 1 dòng chữ}: Chữ sẽ trôi vào giữa khoảng $100\text{px}$ đó, tạo hiệu ứng giống như được căn giữa theo chiều dọc.
    \item \textbf{Nếu \codeinline{h1} có nhiều dòng chữ} (do màn hình hẹp): Các dòng chữ sẽ cách xa nhau đúng $100\text{px}$.
    \item \textbf{\codeinline{line-height} không trực tiếp thay đổi \codeinline{height} phần tử} nếu có khai báo chiều cao cố định.
\end{enumerate}
\end{conceptBox}

\begin{browserWindow}[Mô Phỏng Trực Quan Tác Động Của line-height: 100px]
\begin{tcolorbox}[arc=6pt, colback=white, colframe=primaryBlue!30, boxrule=1pt, top=14pt, bottom=14pt, left=10pt, right=10pt]
  \centering
  % CASE 1: Single Line Text
  {\sffamily\bfseries\color{primaryBlue} Trường Hợp 1: Văn bản 1 dòng (Chữ tự động trôi vào giữa khoảng 100px chiều cao dòng)}\par\vspace{6pt}
  \begin{tikzpicture}[x=1cm, y=1cm, >=stealth]
    % Line box (height = 100px => 2.8cm)
    \draw[line width=1.5pt, primaryBlue, fill=primaryBlue!5, rounded corners=6pt] (0,0) rectangle (11.2, 2.8);
    \node[anchor=center, font=\sffamily\fontsize{11}{11}\selectfont\bfseries\color{darkText}] at (5.6, 1.4) {Write, edit and run HTML, CSS and JavaScript code online.};
    
    % Top/Bottom padding markers
    \node[anchor=center, font=\sffamily\tiny\color{gray!70}] at (5.6, 2.4) {Khoảng trống phía trên (+40px)};
    \node[anchor=center, font=\sffamily\tiny\color{gray!70}] at (5.6, 0.4) {Khoảng trống phía dưới (+40px)};

    % Height arrow on left
    \draw[<->, line width=1.3pt, primaryBlue] (-0.3, 0) -- (-0.3, 2.8);
    \node[rotate=90, anchor=center, font=\sffamily\small\bfseries\color{primaryBlue}] at (-0.7, 1.4) {line-height = 100px};
  \end{tikzpicture}
  
  \par\vspace{14pt}
  % CASE 2: Multi-line Text
  {\sffamily\bfseries\color{accentOrange!90!black} Trường Hợp 2: Văn bản nhiều dòng (Các dòng cách xa nhau đúng 100px chiều cao)}\par\vspace{6pt}
  \begin{tikzpicture}[x=1cm, y=1cm, >=stealth]
    % Line 1 box
    \draw[line width=1.2pt, accentOrange, fill=accentOrange!8, rounded corners=5pt] (0, 2.2) rectangle (11.2, 4.2);
    \node[anchor=west, font=\sffamily\small\bfseries\color{darkText}] at (0.4, 3.2) {Dòng 1: Write, edit and run HTML, CSS...};

    % Line 2 box
    \draw[line width=1.2pt, accentOrange, fill=accentOrange!8, rounded corners=5pt] (0, 0) rectangle (11.2, 2.0);
    \node[anchor=west, font=\sffamily\small\bfseries\color{darkText}] at (0.4, 1.0) {Dòng 2: ...and JavaScript code online.};

    % Baseline gap indicator on right
    \draw[<->, line width=1.5pt, accentOrange!90!black] (11.5, 1.0) -- (11.5, 3.2);
    \node[rotate=-90, anchor=center, font=\sffamily\small\bfseries\color{accentOrange!90!black}] at (11.9, 2.1) {Khoảng cách 2 dòng = 100px};
  \end{tikzpicture}
  \par\vspace{10pt}
  {\sffamily\small\color{gray!70} $\rightarrow$ \textbf{Kết luận}: \codeinline{line-height: 100px} tạo khoảng đứng $100\text{px}$ cho mỗi dòng. Nếu 1 dòng chữ thì trôi vào giữa, nếu nhiều dòng chữ thì các dòng đẩy xa nhau đúng $100\text{px}$.}
\end{tcolorbox}
\end{browserWindow}

\begin{codeWindow}[Kỹ Thuật Căn Giữa Chữ 1 Dòng Với Container Chuẩn]
\begin{lstlisting}[language=HTML5]
<div class="box">
  <h1>(*@ Write, edit and run HTML, CSS and JavaScript code online. @*)</h1>
</div>
<style>
  .box {
    height: 100px;
    background: #eee;
  }
  h1 {
    margin: 0;
    font-size: 20px;
    line-height: 100px; /* (*@ Bang voi chieu cao height cua .box @*) */
    text-align: center;
  }
</style>
\end{lstlisting}
\end{codeWindow}

\begin{browserWindow}[Mô Phỏng Trực Quan Kỹ Thuật Căn Giữa Chữ 1 Dòng Theo Chiều Dọc]
\begin{tcolorbox}[arc=6pt, colback=white, colframe=primaryBlue!30, boxrule=1pt, top=14pt, bottom=14pt, left=10pt, right=10pt]
  \centering
  \begin{tikzpicture}[x=1cm, y=1cm, >=stealth]
    % Outer .box Container (height = 100px => 3.2cm, width = 11.2cm)
    \draw[line width=1.5pt, gray!60, fill=gray!15, rounded corners=6pt] (0,0) rectangle (11.2, 3.2);
    
    % Top half-leading space (40px => 1.2cm)
    \node[anchor=center, font=\sffamily\small\color{gray!70}] at (5.6, 2.6) {Half-leading (Khoảng trống phía trên): 40px};
    
    % Centered text line area (20px font => 0.8cm)
    \fill[primaryBlue!15, rounded corners=4pt] (0.4, 1.2) rectangle (10.8, 2.0);
    \node[anchor=center, font=\sffamily\fontsize{11}{11}\selectfont\bfseries\color{primaryBlue}] at (5.6, 1.6) {Write, edit and run HTML, CSS and JavaScript code online.};

    % Bottom half-leading space (40px => 1.2cm)
    \node[anchor=center, font=\sffamily\small\color{gray!70}] at (5.6, 0.6) {Half-leading (Khoảng trống phía dưới): 40px};

    % Left Arrow: Container Height (.box height = 100px)
    \draw[<->, line width=1.3pt, darkText] (-0.3, 0) -- (-0.3, 3.2);
    \node[rotate=90, anchor=center, font=\sffamily\small\bfseries\color{darkText}] at (-0.7, 1.6) {.box height = 100px};

    % Right Arrow: Line-height (h1 line-height = 100px)
    \draw[<->, line width=1.3pt, primaryBlue] (11.5, 0) -- (11.5, 3.2);
    \node[rotate=-90, anchor=center, font=\sffamily\small\bfseries\color{primaryBlue}] at (11.9, 1.6) {line-height = 100px};
  \end{tikzpicture}
  \par\vspace{10pt}
  {\sffamily\small\color{gray!70} \textbf{Giải thích kỹ thuật}: Khi \codeinline{line-height} của chữ ($100\text{px}$) bằng đúng \codeinline{height} của khung chứa ($100\text{px}$), $40\text{px}$ nửa khoảng trống trên và $40\text{px}$ nửa khoảng trống dưới sẽ đẩy chữ trôi vào chính giữa khung.}
\end{tcolorbox}
\end{browserWindow}

\begin{codeWindow}[Ví Dụ Thực Tế 2: Phân Tích Lỗi Xung Đột height: 50px VÀ line-height: 100px]
\begin{lstlisting}[language=HTML5]
<h1>
  (*@ Write, edit and run HTML, CSS and JavaScript code online. @*)
</h1>
<style>
  h1 {
    height: 50px;       /* (*@ Gioi han chieu cao khung chua chi 50px @*) */
    line-height: 100px; /* (*@ Yeu cau dong chu cao toi 100px @*) */
  }
</style>
\end{lstlisting}
\end{codeWindow}

\begin{browserWindow}[Mô Phỏng Trực Quan Lỗi Tràn/Cắt Chữ Khi height: 50px VÀ line-height: 100px]
\begin{tcolorbox}[arc=6pt, colback=white, colframe=red!30, boxrule=1pt, top=14pt, bottom=16pt, left=12pt, right=12pt]
  \centering
  \begin{tikzpicture}[x=1cm, y=1cm, >=stealth]
    % Outer Container Box (height = 50px => 1.6cm)
    \draw[line width=1.5pt, red!80!black, fill=red!5, rounded corners=6pt] (0,0) rectangle (11.5, 1.6);
    \node[anchor=north west, font=\sffamily\small\bfseries\color{red!80!black}] at (0.3, 1.5) {Khung phần tử h1 (giới hạn height: 50px)};

    % Line Box phantom outline (height = 100px => 3.2cm)
    \draw[dashed, line width=1.2pt, primaryBlue!70] (0.4, -0.9) rectangle (11.1, 2.3);
    \node[anchor=east, font=\sffamily\footnotesize\bfseries\color{primaryBlue}] at (10.9, 2.05) {Line box (line-height: 100px)};

    % Text overflowing out of the bottom border!
    \node[anchor=center, font=\sffamily\fontsize{11}{11}\selectfont\bfseries\color{darkText}] at (5.75, 0.0) {Write, edit and run HTML, CSS and JavaScript code online.};

    % Red Warning Arrow pointing to overflow area
    \draw[<-, line width=1.5pt, red!80!black] (8.2, -0.1) -- (9.2, -0.6);
    \node[anchor=west, font=\sffamily\small\bfseries\color{red!80!black}] at (9.3, -0.6) {Tràn khỏi đáy khung!};
  \end{tikzpicture}
  \par\vspace{10pt}
  {\sffamily\small\color{red!80!black} \textbf{Giải thích trực quan}: Khung \codeinline{h1} chỉ cao $50\text{px}$ (khung đỏ), nhưng dòng chữ có \codeinline{line-height: 100px} (khung nét đứt xanh) nên tâm dòng bị kéo xuống thấp, khiến toàn bộ văn bản bị trôi rớt ra khỏi mép đáy!}
\end{tcolorbox}
\end{browserWindow}

\begin{warnBox}[Phân Tích Kết Quả Lỗi Không Như Mong Muốn]
\begin{itemize}[leftmargin=1.5em, itemsep=3pt, parsep=0pt]
    \item Trình duyệt cố gắng hiển thị dòng chữ cao $100\text{px}$.
    \item Nhưng chiều cao phần tử \codeinline{h1} lại bị giới hạn cứng chỉ $50\text{px}$.
    \item $\rightarrow$ \textbf{Kết quả}: Văn bản bị tràn ra ngoài (\term{overflow}) hoặc bị cắt mất không hiển thị đủ!
\end{itemize}
\end{warnBox}

\begin{prettyTableBox}[Bảng Giải Thích Chi Tiết Xung Đột Giữa height Và line-height]
\sffamily\small
\begin{tabularx}{\linewidth}{|>{\ttfamily\color{primaryBlue}\bfseries}p{3.5cm}|>{\raggedright\arraybackslash}X|}
\hline
\rowcolor{primaryBlue} \thd{\sffamily Thuộc tính} & \thd{\sffamily Tác động thực tế \& Kết quả} \\
\hline
height: 50px & Gán cố định chiều cao khung chứa phần tử là $50\text{px}$. \\
\hline
line-height: 100px & Dòng chữ chiếm $100\text{px}$ chiều cao $\rightarrow$ Bị tràn/cắt (\term{overflow}). \\
\hline
\end{tabularx}
\end{prettyTableBox}

\begin{tipBox}[Giải Pháp Sử Dụng Chuẩn Cho Từng Nhu Cầu]
\begin{itemize}[leftmargin=1.5em, itemsep=3pt, parsep=0pt]
    \item \textbf{Trường hợp 1: Muốn chữ nằm giữa chiều cao 100px (văn bản 1 dòng)}:
    \par\vspace{2pt}
    Đặt \codeinline{height: 100px;} kết hợp với \codeinline{line-height: 100px;} (bằng chiều cao).
    \item \textbf{Trường hợp 2: Tránh dùng height nhỏ hơn line-height}:
    \par\vspace{2pt}
    Không bao giờ đặt \codeinline{height < line-height} trừ khi bạn cố ý cắt xén văn bản.
\end{itemize}
\end{tipBox}

\begin{browserWindow}[So Sánh Mô Phỏng Khoảng Cách Dòng line-height]
\begin{tcolorbox}[arc=6pt, colback=white, colframe=primaryBlue!30, boxrule=1pt, top=12pt, bottom=12pt, left=6pt, right=6pt]
  \centering
  \begin{tikzpicture}[x=1cm, y=1cm, >=stealth]
    % --- CASE 1: line-height 1.0 ---
    \draw[line width=1.2pt, red!70!black, fill=red!5, rounded corners=5pt] (0, 8.3) rectangle (11.8, 10.7);
    \node[anchor=north west, font=\sffamily\small\bfseries\color{red!80!black}] at (0.3, 10.55) {line-height: 1.0 (Chật hẹp -- các dòng sát rạt nhau):};
    \node[anchor=west, font=\sffamily\small\color{darkText}] at (0.4, 9.85) {Dòng 1: Thuộc tính line-height quyết định chiều cao mỗi dòng chữ.};
    \node[anchor=west, font=\sffamily\small\color{darkText}] at (0.4, 9.38) {Dòng 2: Khi đặt 1.0, khoảng cách giữa các dòng bị ép sát nhau.};
    \node[anchor=west, font=\sffamily\small\color{darkText}] at (0.4, 8.90) {Dòng 3: Gây mỏi mắt và khó theo dõi khi đọc văn bản dài.};

    % --- CASE 2: line-height 1.4 ---
    \draw[line width=1.2pt, accentOrange, fill=accentOrange!5, rounded corners=5pt] (0, 5.5) rectangle (11.8, 8.0);
    \node[anchor=north west, font=\sffamily\small\bfseries\color{accentOrange!90!black}] at (0.3, 7.85) {line-height: 1.4 (Vừa đủ -- phù hợp bài viết ngắn):};
    \node[anchor=west, font=\sffamily\small\color{darkText}] at (0.4, 7.15) {Dòng 1: Thuộc tính line-height quyết định chiều cao mỗi dòng chữ.};
    \node[anchor=west, font=\sffamily\small\color{darkText}] at (0.4, 6.60) {Dòng 2: Khi đặt 1.4, khoảng cách dòng có độ giãn vừa phải.};
    \node[anchor=west, font=\sffamily\small\color{darkText}] at (0.4, 6.05) {Dòng 3: Đọc tốt trên giao diện di động hoặc danh sách ngắn.};

    % --- CASE 3: line-height 1.6 ---
    \draw[line width=1.2pt, green!60!black, fill=green!5, rounded corners=5pt] (0, 2.4) rectangle (11.8, 5.2);
    \node[anchor=north west, font=\sffamily\small\bfseries\color{green!70!black}] at (0.3, 5.05) {line-height: 1.6 (Thoải mái, chuẩn UI/UX -- Khuyên dùng):};
    \node[anchor=west, font=\sffamily\small\color{darkText}] at (0.4, 4.35) {Dòng 1: Thuộc tính line-height quyết định chiều cao mỗi dòng chữ.};
    \node[anchor=west, font=\sffamily\small\color{darkText}] at (0.4, 3.65) {Dòng 2: Đặt 1.6 giúp mắt lướt qua các dòng cực kỳ tự nhiên.};
    \node[anchor=west, font=\sffamily\small\color{darkText}] at (0.4, 2.95) {Dòng 3: Đây là tiêu chuẩn vàng được áp dụng phổ biến nhất.};

    % --- CASE 4: line-height 2.0 ---
    \draw[line width=1.2pt, primaryBlue, fill=primaryBlue!5, rounded corners=5pt] (0, -0.9) rectangle (11.8, 2.1);
    \node[anchor=north west, font=\sffamily\small\bfseries\color{primaryBlue}] at (0.3, 1.95) {line-height: 2.0 (Rộng rãi -- Tạo khoảng thở thiết kế):};
    \node[anchor=west, font=\sffamily\small\color{darkText}] at (0.4, 1.20) {Dòng 1: Thuộc tính line-height quyết định chiều cao mỗi dòng chữ.};
    \node[anchor=west, font=\sffamily\small\color{darkText}] at (0.4, 0.35) {Dòng 2: Đặt 2.0 tạo cảm giác thoáng đãng, tự do tối đa.};
    \node[anchor=west, font=\sffamily\small\color{darkText}] at (0.4, -0.50) {Dòng 3: Phù hợp cho thiết kế hiện đại chú trọng khoảng trắng.};
  \end{tikzpicture}
  \par\vspace{8pt}
  {\sffamily\small\color{gray!70} $\rightarrow$ \textbf{Khuyên dùng}: Đặt \codeinline{line-height: 1.5} đến \codeinline{1.8} cho văn bản bài viết dài (\codeinline{body text}).}
\end{tcolorbox}
\end{browserWindow}

\paragraph*{VI. Ứng Dụng Thực Tế \& Những Điều Cần Tránh (Best Practices)}

\begin{prettyTableBox}[Bảng Ứng Dụng Chi Tiết Của Thuộc Tính line-height]
\sffamily\small
\begin{tabularx}{\linewidth}{|>{\raggedright\arraybackslash\bfseries\color{primaryBlue}}p{4.2cm}|>{\raggedright\arraybackslash}p{3.2cm}|>{\raggedright\arraybackslash}X|}
\hline
\rowcolor{primaryBlue} \thd{\sffamily Mục đích thiết kế} & \thd{\sffamily Cách dùng line-height} & \thd{\sffamily Hiệu quả mang lại} \\
\hline
Căn giữa chữ theo chiều dọc & \codeinline{line-height = height} & Áp dụng cho văn bản 1 dòng (button, badge, menu item). \\
\hline
Tăng khoảng cách các dòng & \codeinline{line-height: 1.5 -- 1.8} & Tăng sự thông thoáng cho văn bản đọc dài (body text). \\
\hline
Tối ưu tiêu đề lớn (Heading) & \codeinline{line-height: 1.1 -- 1.3} & Tránh tiêu đề bị xa nhau quá mức do font to. \\
\hline
Tạo cảm giác hiện đại (Spacious) & \codeinline{line-height: 1.8 -- 2.0} & Tạo ``khoảng thở'' tinh tế cho giao diện UI/UX. \\
\hline
\end{tabularx}
\end{prettyTableBox}

\begin{warnBox}[Những Điều Cần Tránh Tuyệt Đối Với line-height]
\begin{enumerate}[leftmargin=1.5em, itemsep=3pt, parsep=0pt]
    \item \textbf{Không đặt \codeinline{line-height} quá nhỏ} (ví dụ $< 1.0$ hoặc $0$): Chữ sẽ dính chặt vào nhau hoặc chồng đè lên nhau gây mất nét.
    \item \textbf{Hạn chế dùng đơn vị \codeinline{px} cố định}: Khi thay đổi \codeinline{font-size}, \codeinline{line-height} không tự điều chỉnh theo, có thể làm vỡ layout.
    \item \textbf{Không dùng \codeinline{line-height} để căn giữa đoạn văn nhiều dòng}: Sẽ khiến các dòng bị đẩy xa nhau cực kỳ bất hợp lý. Căn giữa nhiều dòng nên dùng Flexbox (\codeinline{align-items: center}) hoặc CSS Grid.
\end{enumerate}
\end{warnBox}

\paragraph*{VII. Kết Luận Tổng Kết Ngắn Gọn}

\begin{prettyTableBox}[Bảng Tổng Kết Ngắn Gọn Về Thuộc Tính line-height]
\sffamily\small
\begin{tabularx}{\linewidth}{|>{\raggedright\arraybackslash\bfseries\color{primaryBlue}}p{4.0cm}|>{\raggedright\arraybackslash}X|}
\hline
\rowcolor{primaryBlue} \thd{\sffamily Tiêu chí} & \thd{\sffamily Tóm tắt bản chất} \\
\hline
Là gì? & Chiều cao của mỗi dòng văn bản (line box). \\
\hline
Không phải là gì? & Không phải là kích thước phông chữ (\codeinline{font-size}). \\
\hline
Dùng để làm gì? & Căn giữa chữ (1 dòng), tạo khoảng cách dòng thoáng mát, tăng khả năng đọc. \\
\hline
Cách dùng khuyên dùng? & \codeinline{line-height: 1.5} (số không đơn vị). \\
\hline
\end{tabularx}
\end{prettyTableBox}

\paragraph*{VIII. Tình Huống Thực Tế 1: Xử Lý Căn Giữa Nút Bấm Khi Height Xung Đột line-height}

Hãy cùng phân tích chi tiết từng bước vì sao nút bấm (\codeinline{.btn}) của bạn không căn giữa theo chiều dọc và cách sửa triệt để mà không cần tới Flexbox.

\begin{conceptBox}[Mục Tiêu Thiết Kế Nút Bấm]
Bạn muốn tạo một nút bấm (\codeinline{See Courses}) nằm chính giữa khung chứa cả theo chiều dọc lẫn chiều ngang.
\end{conceptBox}

\begin{codeWindow}[Mã Nguồn Ban Đầu Gặp Lỗi Căn Giữa]
\begin{lstlisting}[language=HTML5]
<div class="hero-cta">
  <div class="btn">(*@ See Courses @*)</div>
</div>

<style>
  .hero-cta {
    height: 64px; /* Khung cha cao 64px */
  }
  .btn {
    line-height: 50px; /* (*@ LỖI: Dòng chữ chỉ cao 50px @*) */
  }
</style>
\end{lstlisting}
\end{codeWindow}

\begin{warnBox}[Vấn Đề Gì Đang Xảy Ra?]
\begin{itemize}[leftmargin=1.5em, itemsep=3pt, parsep=0pt]
    \item \codeinline{.btn} có \codeinline{line-height: 50px} làm cho dòng chữ ``See Courses'' chỉ chiếm $50\text{px}$ chiều cao.
    \item Nhưng khung chứa \codeinline{.hero-cta} lại cao tới $64\text{px}$.
    \item $\rightarrow$ \textbf{Kết quả}: Dòng chữ bị lệch lên phía trên vì không chiếm hết chiều cao của khung \codeinline{.hero-cta}.
\end{itemize}
\end{warnBox}

\begin{codeWindow}[Cách Sửa Chuẩn Kỹ Thuật (Không Dùng Flexbox)]
\begin{lstlisting}[language=CSS3]
.btn {
  height: 64px;          /* Dat chieu cao bang .hero-cta */
  line-height: 64px;     /* Dat line-height bang height */
  text-align: center;    /* Can giua theo chieu ngang */
  padding: 0 16px;
  color: #ffffff;
  font-size: 1rem;
  font-weight: 600;
  border-radius: 999px;
  background: #171100;
  display: inline-block;
}
.hero-cta {
  width: 180px;
}
\end{lstlisting}
\end{codeWindow}

\begin{browserWindow}[Mô Phỏng Trực Quan Kết Quả Nút Bấm Căn Giữa Hoàn Hảo]
\begin{tcolorbox}[arc=6pt, colback=white, colframe=primaryBlue!30, boxrule=1pt, top=14pt, bottom=14pt, left=10pt, right=10pt]
  \centering
  \begin{tikzpicture}[x=1cm, y=1cm, >=stealth]
    % Container Label at top
    \node[anchor=south west, font=\sffamily\small\bfseries\color{gray!80!black}] at (0, 2.9) {.hero-cta (Khung chứa: height = 64px)};

    % Outer .hero-cta container outline (width = 11.4cm, height = 64px => 2.6cm)
    \draw[dashed, line width=1.5pt, gray!50, fill=gray!6, rounded corners=8pt] (0, 0.1) rectangle (11.4, 2.7);

    % Inner .btn Pill (height = 64px => 2.6cm, centered horizontally between 3.2cm and 8.2cm)
    \fill[black!90, rounded corners=14pt] (3.2, 0.1) rectangle (8.2, 2.7);
    \node[anchor=center, font=\sffamily\fontsize{11}{11}\selectfont\bfseries\color{white}] at (5.7, 1.4) {See Courses};

    % Left Arrow: Line-height indicator
    \draw[<->, line width=1.3pt, primaryBlue] (1.4, 0.1) -- (1.4, 2.7);
    \node[rotate=90, anchor=center, font=\sffamily\small\bfseries\color{primaryBlue}] at (0.9, 1.4) {line-height = 64px};

    % Right Arrow: Height indicator
    \draw[<->, line width=1.3pt, darkText] (10.0, 0.1) -- (10.0, 2.7);
    \node[rotate=-90, anchor=center, font=\sffamily\small\bfseries\color{darkText}] at (10.5, 1.4) {height = 64px};
  \end{tikzpicture}
  \par\vspace{10pt}
  {\sffamily\small\color{gray!70} $\rightarrow$ \textbf{Giải thích}: Khi đặt \codeinline{line-height = height = 64px}, chữ trôi vào chính giữa nút bấm mà không cần Flexbox.}
\end{tcolorbox}
\end{browserWindow}

\begin{conceptBox}[Vì Sao Giải Pháp Này Hoạt Động?]
Khi bạn đặt \codeinline{line-height = height}, văn bản giống như nằm trong một ô vuông và tự động trôi vào chính giữa ô đó theo chiều dọc. Kết hợp \codeinline{text-align: center} giúp chữ căn giữa theo chiều ngang.
\end{conceptBox}

\begin{warnBox}[Hạn Chế Của Kỹ Thuật Này]
\begin{itemize}[leftmargin=1.5em, itemsep=3pt, parsep=0pt]
    \item Chỉ sử dụng được khi văn bản chỉ có \textbf{1 dòng}.
    \item Nếu chữ quá dài xuống hàng hoặc có thêm Icon, nút bấm sẽ bị tràn vỡ giao diện. Khi đó bắt buộc dùng Flexbox (\codeinline{display: flex; align-items: center;}).
\end{itemize}
\end{warnBox}

\begin{prettyTableBox}[Tóm Tắt Phương Pháp Căn Giữa Dòng Văn Bản]
\sffamily\small
\begin{tabularx}{\linewidth}{|>{\raggedright\arraybackslash\bfseries\color{primaryBlue}}p{3.5cm}|>{\raggedright\arraybackslash}X|}
\hline
\rowcolor{primaryBlue} \thd{\sffamily Tình huống} & \thd{\sffamily Cách giải quyết tối ưu} \\
\hline
Chữ 1 dòng (Button, Badge) & Đặt \codeinline{line-height = height} kết hợp \codeinline{text-align: center}. \\
\hline
Chữ nhiều dòng / Có Icon & Dùng Flexbox (\codeinline{display: flex; align-items: center; gap: 8px;}) hoặc \codeinline{display: table-cell}. \\
\hline
\end{tabularx}
\end{prettyTableBox}

\paragraph*{IX. Chuyên Đề Thuộc Tính Khoảng Cách gap (Flexbox, Grid \& Multi-column)}

Trong CSS hiện đại, \codeinline{gap} (trước đây là \codeinline{grid-gap}) là thuộc tính cực kỳ tiện lợi dùng để tạo khoảng cách giữa các phần tử con bên trong một khối container sử dụng \textbf{Flexbox}, \textbf{Grid}, hoặc \textbf{Multi-column}.

\begin{conceptBox}[Tổng Quan Về Thuộc Tính gap]
\begin{codeWindow}[Khai Báo gap Cơ Bản]
\begin{lstlisting}[language=CSS3]
.container {
  display: grid; /* Hoac flex */
  gap: 20px;     /* Khoang cach 20px giua cac item con */
}
\end{lstlisting}
\end{codeWindow}
$\rightarrow$ \codeinline{gap: 20px;} tạo khoảng cách $20\text{px}$ giữa các phần tử con với nhau, \textbf{không áp dụng ra mép ngoài} của container (giúp loại bỏ hoàn toàn bẫy \codeinline{margin-right} bị thừa ở ô cuối cùng).
\end{conceptBox}

\begin{prettyTableBox}[Cú Pháp Đầy Đủ Của Thuộc Tính gap]
\sffamily\small
\begin{tabularx}{\linewidth}{|>{\ttfamily\color{accentOrange}\bfseries}p{4.2cm}|>{\raggedright\arraybackslash}X|}
\hline
\rowcolor{primaryBlue} \thd{\sffamily Cách viết cú pháp} & \thd{\sffamily Ý nghĩa kỹ thuật} \\
\hline
gap: 10px; & Hàng và cột đều có khoảng cách đúng $10\text{px}$. \\
\hline
gap: 10px 20px; & Cách hàng (\term{row}) $10\text{px}$, cách cột (\term{column}) $20\text{px}$. \\
\hline
row-gap: 10px; & Chỉ tạo khoảng cách giữa các hàng. \\
\hline
column-gap: 20px; & Chỉ tạo khoảng cách giữa các cột. \\
\hline
\end{tabularx}
\end{prettyTableBox}

\begin{codeWindow}[1. Sử Dụng gap Với CSS Grid Layout]
\begin{lstlisting}[language=CSS3]
.container {
  display: grid;
  grid-template-columns: repeat(3, 1fr);
  gap: 20px; /* Tao luoi 3 cot, moi o cach nhau 20px ca hang va cot */
}
\end{lstlisting}
\end{codeWindow}

\begin{codeWindow}[2. Sử Dụng gap Với Flexbox Layout (Modern CSS)]
\begin{lstlisting}[language=CSS3]
.container {
  display: flex;
  flex-wrap: wrap; /* (*@ Giúp phần tử tự xuống dòng khi thiếu chỗ @*) */
  gap: 16px;      /* (*@ Tạo khoảng cách 16px giữa các ô flex item @*) */
}
\end{lstlisting}
\end{codeWindow}

\begin{tipBox}[Sự Khác Biệt Khi Dùng gap Trong Flexbox So Với Cách Cũ]
Trước đây khi chưa có \codeinline{gap} cho Flexbox, lập trình viên phải xử lý thủ công bằng cách tạo margin và xóa margin ở phần tử cuối cùng. Giờ đây với \codeinline{gap: 16px;}, bạn không cần viết thêm các bộ chọn phức tạp nữa!
\end{tipBox}

\begin{codeWindow}[Kỹ Thuật Xử Lý Thủ Công Cũ Khi Chưa Có gap Trong Flexbox]
\begin{lstlisting}[language=CSS3]
/* (*@ Cách cũ vất vả: Phải dùng :not(:last-child) để bỏ margin ô cuối @*) */
.item:not(:last-child) {
  margin-right: 16px;
}
\end{lstlisting}
\end{codeWindow}

\begin{codeWindow}[3. Sử Dụng gap Với Multi-column Layout]
\begin{lstlisting}[language=CSS3]
.columns {
  column-count: 3;
  column-gap: 40px; /* (*@ Khoảng cách 40px giữa các cột văn bản @*) */
}
\end{lstlisting}
\end{codeWindow}

\begin{prettyTableBox}[Bảng Tổng Quan Khả Năng Hỗ Trợ Layout Của gap]
\sffamily\small
\begin{tabularx}{\linewidth}{|>{\ttfamily\color{primaryBlue}\bfseries}p{3.5cm}|>{\center\arraybackslash}p{2.5cm}|>{\center\arraybackslash}p{2.5cm}|>{\center\arraybackslash}X|}
\hline
\rowcolor{primaryBlue} \thd{\sffamily Thuộc tính} & \thd{\sffamily CSS Grid} & \thd{\sffamily Flexbox} & \thd{\sffamily Multi-column} \\
\hline
gap & \textcolor{green!60!black}{\textbf{Có}} & \textcolor{green!60!black}{\textbf{Có}} & \textcolor{green!60!black}{\textbf{Có}} \\
\hline
row-gap & \textcolor{green!60!black}{\textbf{Có}} & \textcolor{green!60!black}{\textbf{Có}} & \color{gray}{\textbf{Không}} \\
\hline
column-gap & \textcolor{green!60!black}{\textbf{Có}} & \textcolor{green!60!black}{\textbf{Có}} & \textcolor{green!60!black}{\textbf{Có}} \\
\hline
\end{tabularx}
\end{prettyTableBox}

\paragraph*{X. Phân Tích Chi Tiết Tình Huống Áp Dụng .line-clamp \& .break-all}

Trong thiết kế giao diện Web hiện đại, xử lý văn bản dài bị tràn khung là một yêu cầu thường xuyên. Ba class dưới đây giải quyết triệt để các bài toán rút gọn văn bản và chống vỡ layout.

\begin{codeWindow}[Khai Báo Class .line-clamp (Rút Gọn Văn Bản 1 Dòng Mặc Định)]
\begin{lstlisting}[language=CSS3]
.line-clamp {
  display: -webkit-box;
  -webkit-line-clamp: var(--line-clamp, 1);
  -webkit-box-orient: vertical;
  overflow: hidden;
}
\end{lstlisting}
\end{codeWindow}

\begin{conceptBox}[Giải Thích Chi Tiết Các Thuộc Tính Trong .line-clamp]
\begin{itemize}[leftmargin=1.5em, itemsep=3pt, parsep=0pt]
    \item \codeinline{display: -webkit-box;}: Sử dụng mô hình hộp Flexbox cũ của Webkit để kết hợp tính năng cắt dòng.
    \item \codeinline{-webkit-line-clamp: var(--line-clamp, 1);}: Giới hạn số dòng hiển thị. Sử dụng biến CSS \codeinline{--line-clamp}, mặc định là 1 dòng.
    \item \codeinline{-webkit-box-orient: vertical;}: Định hướng hộp theo chiều dọc (bắt buộc phải có để \codeinline{line-clamp} hoạt động).
    \item \codeinline{overflow: hidden;}: Ẩn phần văn bản vượt quá số dòng quy định và tự động thêm dấu ba chấm (\codeinline{...}).
\end{itemize}
\end{conceptBox}

\begin{codeWindow}[Khai Báo Class .line-clamp.line-2 (Rút Gọn Văn Bản 2 Dòng)]
\begin{lstlisting}[language=CSS3]
.line-clamp.line-2 {
  --line-clamp: 2; /* (*@ Gán lại biến --line-clamp thành 2 dòng @*) */
}
\end{lstlisting}
\end{codeWindow}

\begin{codeWindow}[Khai Báo Class .break-all (Chống Vỡ Layout Khi Chuỗi Không Có Khoảng Trắng)]
\begin{lstlisting}[language=CSS3]
.break-all {
  word-break: break-all; /* (*@ Ép ngắt từ ở bất kỳ ký tự nào khi chạm mép khung @*) */
}
\end{lstlisting}
\end{codeWindow}

\begin{codeWindow}[Ví Dụ Thực Tế Minh Họa Đầy Đủ 3 Tình Huống]
\begin{lstlisting}[language=HTML5]
<!DOCTYPE html>
<html lang="vi">
<head>
  <meta charset="UTF-8" />
  <style>
    body { font-family: sans-serif; line-height: 1.5; padding: 20px; }
    .box { width: 300px; border: 1px solid #ccc; padding: 10px; margin-bottom: 15px; }
    
    .line-clamp {
      display: -webkit-box;
      -webkit-line-clamp: var(--line-clamp, 1);
      -webkit-box-orient: vertical;
      overflow: hidden;
    }
    .line-clamp.line-2 { --line-clamp: 2; }
    .break-all { word-break: break-all; }
  </style>
</head>
<body>
  <h2>(*@ 1. Giới hạn 1 dòng (.line-clamp) @*)</h2>
  <div class="box line-clamp">
    (*@ Đây là nội dung mô tả sản phẩm rất dài, nó sẽ tự động bị cắt và thêm dấu ba chấm sau 1 dòng. @*)
  </div>

  <h2>(*@ 2. Giới hạn 2 dòng (.line-clamp.line-2) @*)</h2>
  <div class="box line-clamp line-2">
    (*@ Đây là nội dung mô tả sản phẩm rất dài. Nó sẽ bị cắt sau 2 dòng hiển thị. Nhấn Xem thêm để đọc tiếp. @*)
  </div>

  <h2>(*@ 3. Bẻ từ mọi vị trí (.break-all) @*)</h2>
  <div class="box break-all">
    https://example.com/chuoi-url-sieu-sieu-sieu-dai-khong-co-khoang-trang-123456789
  </div>
</body>
</html>
\end{lstlisting}
\end{codeWindow}

\begin{browserWindow}[Mô Phỏng Trực Quan 3 Class Rút Gọn \& Bẻ Từ]
\begin{tcolorbox}[arc=6pt, colback=white, colframe=primaryBlue!30, boxrule=1pt, top=16pt, bottom=16pt, left=10pt, right=10pt]
  \centering
  \begin{tikzpicture}[x=1cm, y=1cm, >=stealth]

    % === BOX 1: .line-clamp (1 dong) ===
    % Outer frame
    \draw[line width=1.4pt, primaryBlue, fill=primaryBlue!4, rounded corners=6pt] (0, 7.0) rectangle (11.4, 9.6);
    % Title badge
    \fill[primaryBlue, rounded corners=4pt] (0.3, 8.9) rectangle (7.8, 9.4);
    \node[anchor=west, font=\sffamily\small\bfseries\color{white}] at (0.5, 9.15) {1. Class .line-clamp (Giới hạn 1 dòng)};
    % Content area with gray bg
    \fill[gray!8, rounded corners=4pt] (0.4, 7.2) rectangle (11.0, 8.6);
    \draw[gray!30, rounded corners=4pt] (0.4, 7.2) rectangle (11.0, 8.6);
    \node[anchor=west, font=\sffamily\small\color{darkText}] at (0.7, 7.9)
      {Đây là nội dung mô tả sản phẩm dài, rất dài, có thể dài đến\textcolor{gray!50}{\textbf{...}}};

    % === BOX 2: .line-clamp.line-2 (2 dong) ===
    % Outer frame
    \draw[line width=1.4pt, accentOrange, fill=accentOrange!4, rounded corners=6pt] (0, 3.2) rectangle (11.4, 6.6);
    % Title badge
    \fill[accentOrange!90!black, rounded corners=4pt] (0.3, 5.9) rectangle (8.6, 6.4);
    \node[anchor=west, font=\sffamily\small\bfseries\color{white}] at (0.5, 6.15) {2. Class .line-clamp.line-2 (Giới hạn 2 dòng)};
    % Content area with gray bg
    \fill[gray!8, rounded corners=4pt] (0.4, 3.4) rectangle (11.0, 5.6);
    \draw[gray!30, rounded corners=4pt] (0.4, 3.4) rectangle (11.0, 5.6);
    \node[anchor=west, font=\sffamily\small\color{darkText}] at (0.7, 5.0)
      {Dòng 1: Đây là nội dung mô tả sản phẩm dài, rất dài.};
    \node[anchor=west, font=\sffamily\small\color{darkText}] at (0.7, 4.1)
      {Dòng 2: Nó sẽ bị cắt sau 2 dòng hiển thị, phần sau\textcolor{gray!50}{\textbf{...}}};

    % === BOX 3: .break-all ===
    % Outer frame
    \draw[line width=1.4pt, green!60!black, fill=green!4, rounded corners=6pt] (0, -1.0) rectangle (11.4, 2.8);
    % Title badge
    \fill[green!50!black, rounded corners=4pt] (0.3, 2.1) rectangle (9.6, 2.6);
    \node[anchor=west, font=\sffamily\small\bfseries\color{white}] at (0.5, 2.35) {3. Class .break-all (Ngắt từ tự động chống vỡ khung)};
    % Content area with gray bg
    \fill[gray!8, rounded corners=4pt] (0.4, -0.8) rectangle (11.0, 1.8);
    \draw[gray!30, rounded corners=4pt] (0.4, -0.8) rectangle (11.0, 1.8);
    \node[anchor=west, font=\ttfamily\small\color{primaryBlue!80!black}] at (0.7, 1.2)
      {https://example.com/chuoi-url-sieu-sieu-};
    \node[anchor=west, font=\ttfamily\small\color{primaryBlue!80!black}] at (0.7, 0.4)
      {sieu-dai-khong-co-khoang-trang-1234567};
    \node[anchor=west, font=\ttfamily\small\color{primaryBlue!80!black}] at (0.7, -0.4)
      {89-mama-mamamamaaaaaaaaaa};

  \end{tikzpicture}
\end{tcolorbox}
\end{browserWindow}

\begin{prettyTableBox}[Bảng Tóm Tắt Tình Huống Sử Dụng 3 Class Rút Gọn \& Bẻ Từ]
\sffamily\small
\begin{tabularx}{\linewidth}{|>{\ttfamily\color{primaryBlue}\bfseries}p{3.2cm}|>{\raggedright\arraybackslash}p{4.2cm}|>{\raggedright\arraybackslash}X|}
\hline
\rowcolor{primaryBlue} \thd{\sffamily Class CSS} & \thd{\sffamily Hiệu ứng hiển thị} & \thd{\sffamily Tình huống sử dụng thực tế} \\
\hline
.line-clamp & Chỉ hiển thị 1 dòng, phần sau ẩn và có dấu \codeinline{...}. & Rút gọn tựa đề bài viết, tên sản phẩm ngắn gọn. \\
\hline
.line-clamp.line-2 & Hiển thị đúng 2 dòng, phần sau ẩn và có dấu \codeinline{...}. & Rút gọn mô tả sản phẩm, trích dẫn bình luận. \\
\hline
.break-all & Bẻ từ tại đúng mép khung mà không chờ khoảng trắng. & Chống vỡ khung cho URL dài, email, mã token, tên file. \\
\hline
\end{tabularx}
\end{prettyTableBox}

\paragraph*{XI. Đọc Thêm: Mở Rộng Kiến Thức Chuyên Sâu \& Phỏng Vấn (Deep-Dive \& Interview Tips)}

\begin{conceptBox}[1. Cơ Chế Inline Formatting Context (IFC) \& Strut Trong Trình Duyệt]
Khi trình duyệt xuất một đoạn văn bản, nó khởi tạo một ngữ cảnh định dạng dòng (\term{Inline Formatting Context - IFC}). 
\begin{itemize}[leftmargin=1.5em, itemsep=3pt, parsep=0pt]
    \item Mỗi dòng được bắt đầu bằng một phần tử ảo không có chiều rộng gọi là \term{strut}.
    \item \term{Strut} mang thuộc tính \codeinline{font-size} và \codeinline{line-height} của phần tử cha, đặt ra chiều cao tối thiểu cho toàn bộ hộp dòng (\term{line box}).
    \item Do đó, dù phần tử con có \codeinline{font-size: 0;}, hộp dòng vẫn giữ chiều cao tối thiểu từ \term{strut} của cha trừ khi cha cũng có \codeinline{line-height: 0;}.
\end{itemize}
\end{conceptBox}

\begin{tipBox}[2. Tiêu Chuẩn Khả Năng Truy Cập (Accessibility - WCAG 2.1 SC 1.4.12)]
Theo tiêu chuẩn quốc tế về khả năng truy cập Web (\term{WCAG 2.1 Success Criterion 1.4.12 - Text Spacing}):
\begin{itemize}[leftmargin=1.5em, itemsep=3pt, parsep=0pt]
    \item Khoảng cách dòng (\codeinline{line-height}) cho văn bản bài viết đọc dài (\codeinline{body text}) phải đạt ít nhất \textbf{1.5 lần} kích thước chữ (\codeinline{font-size}).
    \item Điều này đảm bảo người dùng bị thị lực kém, rối loạn đọc (\term{dyslexia}) hoặc đọc trên thiết bị di động nhỏ có thể theo dõi văn bản mượt mà không bị lộn dòng.
\end{itemize}
\end{tipBox}

\begin{warnBox}[3. Mẹo Xử Lý Lỗi Thẻ Ảnh <img> Bị Hở Khoảng Trắng Đáy]
Mặc định thẻ \codeinline{<img>} là phần tử nội dòng (\term{inline element}), nó nằm trên đường cơ sở (\term{baseline}) của chữ. Khoảng trống phía dưới đường cơ sở (\term{descender space}) do \codeinline{line-height} tạo ra sẽ khiến đáy ảnh bị hở một khe trắng khoảng $3\text{px} - 5\text{px}$.
\end{warnBox}

\begin{codeWindow}[Giải Pháp Khắc Phục Triệt Để Lỗi Hở Đáy Thẻ Ảnh]
\begin{lstlisting}[language=CSS3]
img {
  display: block;          /* (*@ Chuyển thành khối block @*) */
  /* (*@ Hoặc dùng: @*) */
  vertical-align: middle;  /* (*@ Căn giữa theo chiều dọc @*) */
}
\end{lstlisting}
\end{codeWindow}

\begin{prettyTableBox}[4. Bộ Câu Hỏi Phỏng Vấn Thường Gặp Về line-height]
\sffamily\small
\begin{tabularx}{\linewidth}{|>{\raggedright\arraybackslash\bfseries\color{primaryBlue}}p{4.0cm}|>{\raggedright\arraybackslash}X|}
\hline
\rowcolor{primaryBlue} \thd{\sffamily Câu hỏi phỏng vấn} & \thd{\sffamily Câu trả lời chuẩn Senior/Expert} \\
\hline
Phân biệt \codeinline{line-height: 1.5} và \codeinline{line-height: 150\%} khi kế thừa? & \codeinline{150\%} sẽ tính toán giá trị \codeinline{px} cố định ở phần tử cha rồi ép con kế thừa con số đó (gây dính chữ nếu con có font to). Trong khi \codeinline{1.5} ép phần tử con tự nhân lại với chính \codeinline{font-size} của nó (luôn an toàn). \\
\hline
Tại sao \codeinline{line-height = height} thất bại với chữ 2 dòng? & Khi chữ xuống 2 dòng, mỗi dòng sẽ chiếm chiều cao bằng đúng \codeinline{height} của khung, làm cho 2 dòng bị đẩy xa nhau cực kỳ bất hợp lý và tràn vỡ khung. \\
\hline
Bản chất của \codeinline{-webkit-line-clamp} là gì? & Là thuộc tính CSS Webkit độc quyền kết hợp với \codeinline{display: -webkit-box} để tính toán điểm cắt văn bản dựa trên số dòng và tự động chèn ký tự \codeinline{...} ở cuối. \\
\hline
\end{tabularx}
\end{prettyTableBox}

\subsubsection{7. Thuộc Tính Viết Tắt font (Shorthand Property)}
\label{sec:font_shorthand}

Thuộc tính viết tắt \codeinline{font} cho phép bạn đặt nhiều thuộc tính font trong một khai báo duy nhất. Thứ tự các giá trị là rất quan trọng và một số giá trị là bắt buộc.

\begin{warnBox}[Quy Tắc Cú Pháp Bắt Buộc Khi Dùng Shorthand font]
\textbf{Cú pháp chuẩn}:
\begin{center}
\codeinline{font: [font-style] [font-variant] [font-weight] font-size[/line-height] font-family;}
\end{center}
\textbf{Các lưu ý kỹ thuật quan trọng}:
\begin{enumerate}[leftmargin=1.5em, itemsep=3pt, parsep=0pt]
    \item \codeinline{font-size} và \codeinline{font-family} là \textbf{bắt buộc phải có}. Nếu thiếu một trong hai, toàn bộ khai báo \codeinline{font} sẽ bị trình duyệt coi là vô hiệu (invalid).
    \item Nếu \codeinline{line-height} được chỉ định, nó bắt buộc phải nằm ngay sau \codeinline{font-size} và được phân tách bằng dấu gạch chéo (\codeinline{/}) (ví dụ: \codeinline{16px/1.5}).
    \item Các giá trị tùy chọn (\codeinline{font-style}, \codeinline{font-variant}, \codeinline{font-weight}) bắt buộc phải đứng trước \codeinline{font-size}. Nếu bỏ qua, chúng sẽ tự động đặt về giá trị mặc định (\codeinline{normal}).
\end{enumerate}
\end{warnBox}

\begin{codeWindow}[Ví Dụ Cú Pháp Shorthand font]
\begin{lstlisting}[language=CSS3]
.shorthand-text {
  font: italic small-caps bold 16px/1.5 "Times New Roman", serif;
  /* Tuong duong vo'i khai bao dai:
     font-style: italic;
     font-variant: small-caps;
     font-weight: bold;
     font-size: 16px;
     line-height: 1.5;
     font-family: "Times New Roman", serif;
  */
}

.another-shorthand {
  font: 24px Arial, sans-serif; /* Chi co font-size va font-family bat buoc */
}
\end{lstlisting}
\end{codeWindow}

\begin{browserWindow}[Kết Quả: Áp Dụng font Shorthand Property]
\begin{tcolorbox}[arc=4pt, colback=white, colframe=primaryBlue!30, boxrule=0.8pt, top=10pt, bottom=10pt, left=12pt, right=12pt]
  {\sffamily\bfseries\color{primaryBlue} Kết quả \codeinline{.shorthand-text}:}\par\vspace{4pt}
  {\rmfamily\itshape\bfseries\fontsize{13}{20}\selectfont\textsc{Chữ Này Là Small-Caps Đậm Nghiêng, Kích Thước 16px, Dòng 1.5.}}\par\vspace{8pt}
  {\sffamily\bfseries\color{primaryBlue} Kết quả \codeinline{.another-shorthand}:}\par\vspace{4pt}
  {\sffamily\fontsize{16}{20}\selectfont Chữ bình thường 24px dùng Arial, đơn giản và hiện đại.}\par\vspace{4pt}
  {\sffamily\small\color{gray!70} → Một dòng shorthand thay thế 6 dòng khai báo riêng lẻ.}
\end{tcolorbox}
\end{browserWindow}


\subsection*{IX. Chuyên Đề Bắt Lỗi Thực Chiến (CSS Font Debugging Masterclass)}

Trong thực tế lập trình Web, lập trình viên thường gặp một số lỗi phổ biến liên quan đến phông chữ. Dưới đây là bảng phân tích nguyên nhân và cách khắc phục chi tiết:

\begin{prettyTableBox}[Bảng Phân Tích \& Khắc Phục Các Lỗi Phông Chữ Thường Gặp]
\sffamily\small
\begin{tabularx}{\linewidth}{|>{\raggedright\arraybackslash\bfseries\color{red!80!black}}p{3.5cm}|>{\raggedright\arraybackslash}X|>{\raggedright\arraybackslash\bfseries\color{green!60!black}}p{3.8cm}|}
\hline
\rowcolor{primaryBlue} \thd{\sffamily Biểu hiện lỗi} & \thd{\sffamily Nguyên nhân kỹ thuật cốt lõi} & \thd{\sffamily Giải pháp khắc phục chuẩn} \\
\hline
Chữ bị vỡ hoặc hiển thị font mặc định xấu & Quên bọc dấu ngoặc kép cho tên font có khoảng trắng (ví dụ viết \codeinline{Times New Roman} thay vì \codeinline{"Times New Roman"}). & Bắt buộc dùng dấu ngoặc kép: \codeinline{"Times New Roman"}. \\
\hline
Chữ Bold bị nhòe/mờ (\term{Faux Bold}) & Khai báo \codeinline{font-weight: 700} nhưng tệp nhúng chỉ tải biến thể Regular (400). Trình duyệt tự bôi đậm giả lập. & Nhúng thêm tệp font biến thể Bold 700 từ Google Fonts/Server. \\
\hline
Quy tắc \codeinline{@import} bị trình duyệt bỏ qua & Đặt quy tắc \codeinline{@import} phía dưới các khai báo CSS khác thay vì ở dòng đầu tiên của file CSS. & Di chuyển \codeinline{@import} lên ngay dòng 1 ở đầu tệp CSS. \\
\hline
Hiện tượng chữ bị ẩn 3s (\term{FOIT}) & Thiếu thuộc tính \codeinline{font-display: swap;} khiến trình duyệt ẩn chữ trong khi chờ tải tệp font qua mạng. & Thêm \codeinline{font-display: swap;} vào khối \codeinline{@font-face} hoặc URL nhúng. \\
\hline
\end{tabularx}
\end{prettyTableBox}

\subsection{Kỹ Thuật Nhúng Font Ngoại Vào CSS (External Web Fonts)}
\label{sec:web_fonts_embedding}

Trong thiết kế web hiện đại, việc sử dụng các phông chữ độc đáo (không nằm trong danh sách \term{Web-Safe Fonts} mặc định) là yếu tố quan trọng giúp định hình nhận diện thương hiệu và tạo nên phong cách thị giác ấn tượng. Để hiển thị phông chữ tùy chỉnh đồng bộ trên mọi thiết bị và hệ điều hành, lập trình viên cần thực hiện \term{nhúng font} (\term{Font Embedding}).

\begin{conceptBox}[Tổng Quan 3 Phương Pháp Nhúng Font Ngoại]
Có 3 phương pháp nhúng phông chữ chính trong thiết kế web:
\begin{enumerate}[leftmargin=1.5em, itemsep=3pt, parsep=0pt]
    \item \textbf{Sử dụng Google Fonts}: Phương pháp miễn phí, phổ biến và dễ triển khai nhất thông qua hệ thống CDN toàn cầu.
    \item \textbf{Sử dụng Quy tắc \codeinline{@font-face}}: Tự host tệp font trực tiếp trên máy chủ riêng, kiểm soát $100\%$ dữ liệu và bản quyền.
    \item \textbf{Sử dụng Adobe Fonts (Typekit)}: Dịch vụ font chữ cao cấp dành cho nhà thiết kế chuyên nghiệp trong hệ sinh thái Adobe.
\end{enumerate}
\end{conceptBox}

% ------------------------------------------------------------------------------
\subsubsection{1. Sử Dụng Google Fonts (Phương Pháp Phổ Biến \& Dễ Dùng Nhất)}

Google Fonts là thư viện phông chữ miễn phí khổng lồ do Google tối ưu hóa cho môi trường web.

\paragraph{Quy Trình 4 Bước Nhúng Google Fonts:}
\begin{itemize}[leftmargin=1.5em, itemsep=4pt, parsep=0pt]
    \item \textbf{Bước 1 (Truy cập)}: Mở trình duyệt và truy cập vào \codeinline{fonts.google.com}.
    \item \textbf{Bước 2 (Chọn Font \& Style)}: Lựa chọn phông chữ và các biến thể cần dùng (ví dụ: Regular 400, Bold 700, Italic). \textbf{Lưu ý}: Chỉ chọn những biến thể thực sự sử dụng để tối ưu dung lượng tải trang.
    \item \textbf{Bước 3 (Lấy mã nhúng)}: Lựa chọn 1 trong 2 hình thức nhúng do Google cung cấp:
    \begin{itemize}[leftmargin=1.2em, itemsep=2pt]
        \item \textbf{Nhúng qua thẻ \codeinline{<link>} trong HTML} (Khuyên dùng cho hiệu suất tối ưu).
        \item \textbf{Nhúng qua quy tắc \codeinline{@import} trong CSS} (Tiện lợi khi viết thuần tệp CSS).
    \end{itemize}
    \item \textbf{Bước 4 (Áp dụng CSS)}: Khai báo tên font trong thuộc tính \codeinline{font-family} của bộ chọn CSS tương ứng.
\end{itemize}

\paragraph{Hình Thức A: Nhúng Qua Thẻ <link> Trong HTML (Khuyến Nghị Hàng Đầu):}

\begin{codeWindow}[File index.html -- Khai Báo Thẻ Link Trong Head]
\begin{lstlisting}[language=HTML5]
<!DOCTYPE html>
<html lang="vi">
<head>
    <meta charset="UTF-8">
    <meta name="viewport" content="width=device-width, initial-scale=1.0">
    <title>Vi du Google Fonts qua Link</title>
    <!-- Thiet lap ket noi som voi may chu Google Fonts -->
    <link rel="preconnect" href="https://fonts.googleapis.com">
    <link rel="preconnect" href="https://fonts.gstatic.com" crossorigin>
    <!-- Nhung font Roboto voi bien the Regular 400 va Bold 700 -->
    <link href="https://fonts.googleapis.com/css2?family=Roboto:wght@400;700&display=swap" rel="stylesheet">
    <link rel="stylesheet" href="style.css">
</head>
<body>
    <h1>Tieu de su dung Roboto</h1>
    <p>Doan van ban nay cung su dung Roboto.</p>
</body>
</html>
\end{lstlisting}
\end{codeWindow}

\begin{codeWindow}[File style.css -- Áp Dụng Font Cho Các Thẻ HTML]
\begin{lstlisting}[language=CSS3]
/* Trong file style.css */
body {
  font-family: 'Roboto', sans-serif; /* 'Roboto' la ten font dinh nghia tu Google Fonts */
}

h1 {
  font-weight: 700; /* Ap dung style Bold 700 */
}
\end{lstlisting}
\end{codeWindow}

\begin{noteBox}[Ý Nghĩa Của Thẻ Preconnect]
Các thẻ \codeinline{<link rel="preconnect">} thông báo cho trình duyệt chủ động thiết lập kết nối sớm (DNS lookup, TCP handshake, TLS negotiation) với máy chủ Google Fonts (\codeinline{fonts.googleapis.com} và \codeinline{fonts.gstatic.com}), giúp quá trình tải font diễn ra nhanh hơn đáng kể.
\end{noteBox}

\paragraph{Hình Thức B: Nhúng Qua Quy Tắc @import Trong CSS:}

\begin{codeWindow}[File style.css -- Đặt Quy Tắc @import Ở Đầu File CSS]
\begin{lstlisting}[language=CSS3]
/* Dat quy tac @import o dau file style.css (truoc moi quy tac CSS khac) */
@import url('https://fonts.googleapis.com/css2?family=Roboto:wght@400;700&display=swap');

body {
  font-family: 'Roboto', sans-serif;
}

h1 {
  font-weight: 700; /* Ap dung style Bold 700 */
}
\end{lstlisting}
\end{codeWindow}

\begin{warnBox}[Lưu Ý Về Hạn Chế Hiệu Suất Của @import]
Phương pháp \codeinline{@import} có thể làm chậm quá trình tải font một chút so với \codeinline{<link>} vì trình duyệt phải tải và đọc xong file CSS chính mới phát hiện được yêu cầu tải font.
\end{warnBox}

\begin{browserWindow}[Kết Quả Hiển Thị Trực Quan: Phông Chữ Google Fonts Roboto]
\begin{tcolorbox}[arc=4pt, colback=white, colframe=primaryBlue!30, boxrule=0.8pt, top=10pt, bottom=10pt, left=12pt, right=12pt]
  {\sffamily\bfseries\color{primaryBlue}\Large Tiêu đề sử dụng Roboto (font-weight: 700)}\par\vspace{6pt}
  {\sffamily\small Đoạn văn bản này cũng sử dụng Roboto (\codeinline{font-family: 'Roboto', sans-serif;}). Văn bản hiển thị sắc nét, hiện đại trên trình duyệt web.}
\end{tcolorbox}
\end{browserWindow}

\paragraph{So Sánh Chi Tiết Giữa 2 Phương Pháp Nhúng Google Fonts:}

\begin{prettyTableBox}[Bảng So Sánh Chi Tiết Hai Phương Pháp Nhúng Font: <link> vs @import]
\sffamily\small
\begin{tabularx}{\linewidth}{|>{\raggedright\arraybackslash\bfseries\color{primaryBlue}}p{3.2cm}|>{\raggedright\arraybackslash}X|>{\raggedright\arraybackslash}X|}
\hline
\rowcolor{primaryBlue} \thd{\sffamily Tiêu chí so sánh} & \thd{\sffamily Nhúng qua <link> trong HTML} & \thd{\sffamily Nhúng qua @import trong CSS} \\
\hline
\textbf{Vị trí đặt mã} & Đặt thẻ \codeinline{<link>} vào bên trong phần \codeinline{<head>} của file HTML. & Đặt quy tắc \codeinline{@import} ở đầu file CSS (trước mọi quy tắc CSS khác trừ \codeinline{@charset}). \\
\hline
\textbf{Cơ chế tác động} & Trình duyệt đọc HTML gặp thẻ \codeinline{<link>} lập tức gửi yêu cầu tải font song song. & Trình duyệt tải và đọc xong file CSS chính mới phát hiện và gửi yêu cầu tải font. \\
\hline
\textbf{Độc lập tệp mã} & Cần file CSS riêng để viết các quy tắc \codeinline{font-family} áp dụng. & Quy tắc \codeinline{@import} nằm trực tiếp trong file CSS hiện có, áp dụng ngay tại cùng file. \\
\hline
\textbf{Đánh giá hiệu suất} & \textcolor{green!60!black}{\textbf{Ưu điểm: Khuyến nghị hàng đầu}} vì hiệu suất tốt hơn, tải font sớm hơn. & \textcolor{red!80!black}{\textbf{Nhược điểm: Chậm hơn}}, dễ gây ra hiện tượng giật font (\term{FOUT/FOIT}). \\
\hline
\end{tabularx}
\end{prettyTableBox}

\paragraph{Ưu Điểm Nổi Bật Của Thư Viện Google Fonts:}
\begin{itemize}[leftmargin=1.5em, itemsep=3pt, parsep=0pt]
    \item \textbf{Miễn phí \& Dễ sử dụng}: Không cần tải tệp font hay cấu hình máy chủ phức tạp.
    \item \textbf{Hiệu suất tối ưu}: Tệp font được nén tối đa và phân phối qua hệ thống CDN toàn cầu của Google.
    \item \textbf{Đa dạng lựa chọn}: Thư viện khổng lồ với hàng ngàn họ phông chữ độc đáo.
    \item \textbf{Tương thích rộng rãi}: Hỗ trợ mượt mà trên tất cả các trình duyệt và thiết bị di động.
\end{itemize}

% ------------------------------------------------------------------------------
\subsubsection{2. Sử Dụng Quy Tắc @font-face (Tự Host Font Trên Server)}

Phương pháp \codeinline{@font-face} cho phép bạn tự host các tệp font trực tiếp trên máy chủ của dự án. Đây là giải pháp bắt buộc khi sử dụng phông chữ độc quyền, phông chữ bản quyền trả phí, hoặc khi dự án yêu cầu bảo mật và độc lập dữ liệu tuyệt đối.

\paragraph{Quy Trình 3 Bước Tự Host Font Với @font-face:}

\begin{itemize}[leftmargin=1.5em, itemsep=4pt, parsep=0pt]
    \item \textbf{Bước 1 (Chuẩn bị tệp font)}: Chuẩn bị tệp font ở các định dạng tương thích web.
    \begin{itemize}[leftmargin=1.2em, itemsep=3pt, parsep=0pt]
        \item \codeinline{.woff2} (\term{Web Open Font Format 2.0}): \textbf{Khuyên dùng hàng đầu} cho web hiện đại với thuật toán nén Brotli tốt nhất.
        \item \codeinline{.woff} (\term{Web Open Font Format}): Chuẩn font web truyền thống được hỗ trợ rộng rãi.
        \item \codeinline{.ttf} (\term{TrueType Font}): Hỗ trợ tốt trên các thiết bị và trình duyệt cũ.
        \item \codeinline{.otf} (\term{OpenType Font}): Tương tự TTF với các tính năng typography nâng cao.
        \item \codeinline{.eot} (\term{Embedded OpenType}): Định dạng riêng dành cho Internet Explorer (IE6--IE11).
        \item \codeinline{.svg} (\term{SVG Font}): Định dạng vector cũ cho iOS di động đời đầu.
    \end{itemize}
    \textit{Mẹo}: Bạn có thể dùng công cụ \term{Font Squirrel's Webfont Generator} để tự động chuyển đổi và tối ưu hóa tệp font sang các định dạng web.

    \item \textbf{Bước 2 (Lưu trữ)}: Lưu các tệp font vào thư mục nội bộ dự án (ví dụ: \codeinline{fonts/}).
    \item \textbf{Bước 3 (Khai báo @font-face)}: Khai báo quy tắc \codeinline{@font-face} trong file CSS cho từng biến thể style.
\end{itemize}

\begin{codeWindow}[File style.css -- Khai Báo Các Khối @font-face Cho Regular Và Bold]
\begin{lstlisting}[language=CSS3]
/* Khai bao cho style Regular (Chua chu thuong) */
@font-face {
  font-family: 'MyCustomFont'; /* Ten custom font ban tu dat */
  src: url('../fonts/MyCustomFont-Regular.woff2') format('woff2'),
       url('../fonts/MyCustomFont-Regular.woff') format('woff'),
       url('../fonts/MyCustomFont-Regular.ttf') format('truetype');
  font-weight: normal; /* Trong luong thuong */
  font-style: normal; /* Kieu thuong */
  font-display: swap; /* Toi uu hien thi khi cho tai */
}

/* Khai bao cho style Bold (Chua chu in dam) */
@font-face {
  font-family: 'MyCustomFont'; /* Cung ten family, nhung dinh nghia style Bold */
  src: url('../fonts/MyCustomFont-Bold.woff2') format('woff2'),
       url('../fonts/MyCustomFont-Bold.woff') format('woff'),
       url('../fonts/MyCustomFont-Bold.ttf') format('truetype');
  font-weight: bold; /* Trong luong dam */
  font-style: normal;
  font-display: swap;
}

/* Ap dung font trong CSS */
body {
  font-family: 'MyCustomFont', sans-serif;
}

h1 {
  font-weight: bold; /* Trinh duyet tu dong chon MyCustomFont-Bold */
}
\end{lstlisting}
\end{codeWindow}

\begin{browserWindow}[Kết Quả Hiển Thị Custom Font Tự Host Với @font-face]
\begin{tcolorbox}[arc=4pt, colback=white, colframe=primaryBlue!30, boxrule=0.8pt, top=10pt, bottom=10pt, left=12pt, right=12pt]
  {\sffamily\bfseries\color{primaryBlue}\Large Tiêu đề đậm dùng MyCustomFont-Bold}\par\vspace{6pt}
  {\sffamily\small Văn bản thân dùng MyCustomFont-Regular. Tệp \codeinline{.woff2} tải trực tiếp từ server nội bộ, hiển thị mượt mà không bị trễ trang.}
\end{tcolorbox}
\end{browserWindow}

\paragraph{Giải Thích Chi Tiết Các Thuộc Tính Trong Khối @font-face:}
\begin{itemize}[leftmargin=1.5em, itemsep=3pt, parsep=0pt]
    \item \codeinline{font-family}: Tên định danh phông chữ do bạn tự đặt để tái sử dụng trong các quy tắc CSS khác.
    \item \codeinline{src}: Đường dẫn tới các tệp font. Liệt kê nhiều định dạng kèm hàm \codeinline{format()} để trình duyệt tự ưu tiên chọn định dạng tối ưu nhất (như \codeinline{.woff2}).
    \item \codeinline{font-weight} \& \codeinline{font-style}: Đạt vai trò phân loại các biến thể. Mỗi biến thể (Regular, Bold, Italic) cần được khai báo trong một khối \codeinline{@font-face} riêng biệt. Khi bạn viết \codeinline{font-weight: bold;}, trình duyệt sẽ khớp đúng tệp font được gán \codeinline{font-weight: bold;}.
    \item \codeinline{font-display}: Điều khiển cách trình duyệt hiển thị chữ trong quá trình chờ tải tệp font qua mạng.
\end{itemize}

\begin{conceptBox}[Phân Tích Chi Tiết Các Giá Trị Của Thuộc Tính font-display]
\begin{itemize}[leftmargin=1.5em, itemsep=3pt, parsep=0pt]
    \item \codeinline{auto}: Mặc định của trình duyệt (thường ẩn chữ tạm thời).
    \item \codeinline{block}: Ẩn văn bản tối đa 3 giây chờ font tải xong (hiện tượng \term{FOIT - Flash of Invisible Text}).
    \item \codeinline{swap} (\textbf{Khuyên dùng}): Hiển thị ngay văn bản bằng phông chữ dự phòng (\term{fallback font}), sau đó hoán đổi (\term{swap}) sang phông tùy chỉnh khi đã tải xong. Giúp nâng cao trải nghiệm người dùng (\term{UX}).
    \item \codeinline{fallback}: Ẩn văn bản cực ngắn (100ms), nếu font chưa tải xong sẽ dùng font dự phòng mà không hoán đổi về sau.
    \item \codeinline{optional}: Tương tự \codeinline{fallback}, nhưng trình duyệt có thể bỏ qua không tải font nếu kết nối mạng quá chậm.
\end{itemize}
\end{conceptBox}

\paragraph{Đánh Giá Ưu Điểm Và Nhược Điểm Của Phương Pháp @font-face:}

\begin{prettyTableBox}[Bảng So Sánh Ưu Điểm Và Nhược Điểm Của Tự Host Font (@font-face)]
\sffamily\small
\begin{tabularx}{\linewidth}{|>{\raggedright\arraybackslash}X|>{\raggedright\arraybackslash}X|}
\hline
\rowcolor{primaryBlue} \thd{\sffamily Ưu điểm} & \thd{\sffamily Nhược điểm} \\
\hline
\textcolor{green!60!black}{\textbf{Kiểm soát hoàn toàn}}: Sử dụng được mọi font tùy chỉnh hay bản quyền độc quyền. & \textcolor{red!80!black}{\textbf{Cần tự quản lý tệp font}}: Phải tự lưu trữ, tổ chức thư mục và quản lý tài nguyên font trên server. \\
\hline
\textcolor{green!60!black}{\textbf{Độc lập dữ liệu}}: Không bị ảnh hưởng nếu các dịch vụ bên ngoài (Google Fonts) bị gián đoạn. & \textcolor{red!80!black}{\textbf{Tương thích trình duyệt}}: Cần khai báo nhiều định dạng tệp font (\codeinline{.woff2}, \codeinline{.woff}, \codeinline{.ttf}). \\
\hline
\textcolor{green!60!black}{\textbf{Bảo mật cao}}: Thích hợp cho các dự án doanh nghiệp yêu cầu quy định khắt khe về quyền riêng tư. & \textcolor{red!80!black}{\textbf{Yêu cầu tối ưu hiệu suất}}: Phải cấu hình \codeinline{font-display: swap} để tránh làm chậm hiển thị trang. \\
\hline
\end{tabularx}
\end{prettyTableBox}

% ------------------------------------------------------------------------------
\subsubsection{3. Sử Dụng Adobe Fonts (Trước Đây Là Typekit)}

Adobe Fonts là dịch vụ phông chữ cao cấp dành cho người dùng Adobe Creative Cloud, cung cấp thư viện font chuyên nghiệp độc quyền cho các dự án thiết kế thương mại.

\paragraph{Quy Trình 4 Bước Sử Dụng Adobe Fonts:}
\begin{itemize}[leftmargin=1.5em, itemsep=3pt, parsep=0pt]
    \item \textbf{Bước 1 (Kích hoạt font)}: Đăng nhập vào Adobe Creative Cloud, truy cập Adobe Fonts và kích hoạt các phông chữ mong muốn.
    \item \textbf{Bước 2 (Tạo Web Project)}: Tạo một Web Project trên Adobe Fonts và thêm các phông chữ đã chọn vào dự án.
    \item \textbf{Bước 3 (Lấy mã nhúng)}: Chèn đoạn mã \codeinline{<link>} do Adobe cấp vào thẻ \codeinline{<head>} của HTML.
\end{itemize}

\begin{codeWindow}[Mã Nhúng HTML Từ Adobe Fonts]
\begin{lstlisting}[language=HTML5]
<!-- Nhung tep CSS kit tu Adobe Fonts vao head -->
<link rel="stylesheet" href="https://use.typekit.net/your-kit-id.css">
\end{lstlisting}
\end{codeWindow}

\begin{itemize}[leftmargin=1.5em, itemsep=3pt, parsep=0pt]
    \item \textbf{Bước 4 (Áp dụng CSS)}: Sử dụng tên font được cấp trong thuộc tính \codeinline{font-family} của tệp CSS.
\end{itemize}

\paragraph{Ưu Điểm Và Nhược Điểm Của Adobe Fonts:}
\begin{itemize}[leftmargin=1.5em, itemsep=3pt, parsep=0pt]
    \item \textbf{Ưu điểm}: Bộ phông chữ thương mại cao cấp độc quyền; nhúng nhanh qua CDN; tương thích tốt.
    \item \textbf{Nhược điểm}: Yêu cầu tài khoản trả phí Adobe Creative Cloud.
\end{itemize}

\begin{browserWindow}[Kết Quả Hiển Thị: Nhúng Font Cao Cấp Adobe Fonts]
\begin{tcolorbox}[arc=4pt, colback=white, colframe=primaryBlue!30, boxrule=0.8pt, top=10pt, bottom=10pt, left=12pt, right=12pt]
  {\sffamily\bfseries\color{primaryBlue} Kết quả hiển thị phông chữ Adobe Fonts (Typekit):}\par\vspace{6pt}
  {\sffamily\fontsize{13}{17}\selectfont\color{darkgray} Font chữ cao cấp tải nhanh qua mạng phân phối CDN Adobe Typekit, đảm bảo nét chữ chuẩn mực cho các trang web thương mại chuyên nghiệp.}
\end{tcolorbox}
\end{browserWindow}

% ------------------------------------------------------------------------------
\subsubsection{4. Kỹ Thuật Nhúng Nhiều Phông Chữ Trong CSS (Vấn Đề 3: Nhúng Nhiều Font trong CSS)}

\begin{conceptBox}[VẤN ĐỀ 3: NHÚNG NHIỀU FONT TRONG CSS]
\textbf{Khái niệm cốt lõi}: Khi nhúng nhiều phông chữ tùy chỉnh, sử dụng tham số \codeinline{\&family=} để gộp URL trong Google Fonts hoặc khai báo riêng từng khối \codeinline{@font-face} cho mỗi biến thể style khi tự host. Bắt buộc cấu hình \codeinline{font-display: swap;} và giới hạn $1 \rightarrow 3$ họ font để đảm bảo hiệu suất.
\end{conceptBox}

\label{sec:multi_font_embedding}

Khi xây dựng giao diện web phức tạp, nhà thiết kế thường kết hợp nhiều hơn một họ phông chữ tùy chỉnh (ví dụ: dùng 1 font hiện đại cho Tiêu đề và 1 font dễ đọc cho Văn bản thân). Việc nhúng nhiều phông chữ đòi hỏi chiến lược quản lý mã nguồn hợp lý để không làm giảm hiệu suất tải trang.

\paragraph{4.1. Nhúng Nhiều Phông Chữ Từ Google Fonts (Gộp URL Tối Ưu):}

Google Fonts cho phép gộp tất cả các phông chữ và biến thể style đã chọn vào một đường dẫn liên kết \codeinline{<link>} hoặc \codeinline{@import} duy nhất thông qua tham số \codeinline{\&family=}.

\begin{conceptBox}[3 Bước Chọn \& Lấy Mã Nhúng Gộp Nhiều Font Trên Giao Diện Google Fonts]
\begin{enumerate}[leftmargin=1.5em, itemsep=3pt, parsep=0pt]
    \item \textbf{Bước 1 (Chọn font thứ 1)}: Truy cập \codeinline{fonts.google.com}, tìm font \textbf{Roboto} và chọn các kiểu dáng cần dùng (ví dụ: \codeinline{Regular 400} và \codeinline{Bold 700}).
    \item \textbf{Bước 2 (Chọn font thứ 2)}: Không cần tắt bảng điều khiển, tiếp tục tìm font \textbf{Open Sans} và chọn thêm các kiểu dáng (\codeinline{Regular 400} và \codeinline{Bold 700}).
    \item \textbf{Bước 3 (Lấy 1 đoạn mã duy nhất)}: Nhìn sang bảng \term{Selected families} bên phải, Google Fonts tự động tạo cho bạn \textbf{duy nhất 1 thẻ \codeinline{<link>}} gộp cả 2 font.
\end{enumerate}
\end{conceptBox}

\begin{noteBox}[Giải Thích Chi Tiết Cấu Trúc Đường Dẫn URL Gộp Nhiều Font]
Hãy quan sát đoạn đường dẫn \codeinline{href} do Google Fonts tự động tạo ra:
\begin{codeWindow}[Đường Dẫn URL Gộp Nhúng Nhiều Font Từ Google Fonts]
\begin{lstlisting}[language=HTML5]
https://fonts.googleapis.com/css2?family=Open+Sans:wght@400;700&family=Roboto:wght@400;700&display=swap
\end{lstlisting}
\end{codeWindow}
\begin{itemize}[leftmargin=1.5em, itemsep=2pt, parsep=0pt]
    \item \codeinline{family=Open+Sans:wght@400;700} $\rightarrow$ Yêu cầu tải font \textbf{Open Sans} với 2 độ đậm $400$ (Regular) và $700$ (Bold).
    \item Dấu \codeinline{\&} $\rightarrow$ Ký tự kết nối nhiều tham số trong đường dẫn URL.
    \item \codeinline{family=Roboto:wght@400;700} $\rightarrow$ Yêu cầu tải thêm font thứ 2 là \textbf{Roboto} với 2 độ đậm $400$ và $700$.
    \item \codeinline{display=swap} $\rightarrow$ Áp dụng cơ chế hoán đổi font dự phòng ngay lập tức để tránh mờ/ẩn chữ.
\end{itemize}
\end{noteBox}

\begin{codeWindow}[File index.html -- Mã Nhúng Gộp Nhiều Font Trong Head HTML]
\begin{lstlisting}[language=HTML5]
<!DOCTYPE html>
<html lang="vi">
<head>
    <meta charset="UTF-8">
    <meta name="viewport" content="width=device-width, initial-scale=1.0">
    <title>Vi Du Nhung Nhieu Font Google Fonts</title>
    
    <!-- 1. Thiet lap ket noi som (Preconnect) -->
    <link rel="preconnect" href="https://fonts.googleapis.com">
    <link rel="preconnect" href="https://fonts.gstatic.com" crossorigin>
    
    <!-- 2. Nhung DUY NHAT 1 link hop nhat ca Roboto va Open Sans -->
    <link href="https://fonts.googleapis.com/css2?family=Open+Sans:wght@400;700&family=Roboto:wght@400;700&display=swap" rel="stylesheet">
    
    <link rel="stylesheet" href="style.css">
</head>
<body>
    <h1 class="header-text">Tieu De Su Dung Phong Chu Roboto (Bold 700)</h1>
    <p class="body-text">Doan van ban nay su dung phong chu Open Sans (Regular 400) giup de doc tren man hinh.</p>
</body>
</html>
\end{lstlisting}
\end{codeWindow}

\begin{codeWindow}[File style.css -- Áp Dụng Nhiều Font Cho Các Thẻ Giao Diện]
\begin{lstlisting}[language=CSS3]
/* Nhung qua @import o dau file style.css (Neu dung trong CSS) */
@import url('https://fonts.googleapis.com/css2?family=Open+Sans:wght@400;700&family=Roboto:wght@400;700&display=swap');

/* Ap dung phong chu Roboto 700 cho tieu de */
.header-text {
  font-family: 'Roboto', sans-serif;
  font-weight: 700;
}

/* Ap dung phong chu Open Sans 400 va font du phong Arial cho van ban */
.body-text {
  font-family: 'Open Sans', Arial, sans-serif;
  font-weight: 400;
}
\end{lstlisting}
\end{codeWindow}

\begin{browserWindow}[Kết Quả Hiển Thị: Kết Hợp Nhiều Phông Chữ Google Fonts]
\begin{tcolorbox}[arc=4pt, colback=white, colframe=primaryBlue!30, boxrule=0.8pt, top=10pt, bottom=10pt, left=12pt, right=12pt]
  {\sffamily\bfseries\color{primaryBlue}\Large Đây là tiêu đề Roboto Bold (\codeinline{font-family: 'Roboto', sans-serif;})}\par\vspace{6pt}
  {\sffamily\small\color{darkText} Đây là nội dung văn bản sử dụng Open Sans Regular (\codeinline{font-family: 'Open Sans', Arial, sans-serif;}). Việc kết hợp hai phông chữ mang lại độ tương phản thị giác rõ ràng và hiện đại.}
\end{tcolorbox}
\end{browserWindow}

\begin{prettyTableBox}[Lợi Ích Của Việc Gộp Nhiều Font Trong Google Fonts]
\sffamily\small
\begin{tabularx}{\linewidth}{|>{\raggedright\arraybackslash\bfseries\color{primaryBlue}}p{3.8cm}|>{\raggedright\arraybackslash}X|}
\hline
\rowcolor{primaryBlue} \thd{\sffamily Ưu điểm} & \thd{\sffamily Ý nghĩa \& Giá trị hiệu suất} \\
\hline
Giảm số lượng HTTP Request & Gộp nhiều phông chữ vào 1 yêu cầu duy nhất giúp trình duyệt không phải mở nhiều kết nối mạng song song. \\
\hline
Dễ dàng quản lý mã nguồn & Chỉ cần một đoạn mã nhúng duy nhất ở phần \codeinline{<head>} cho toàn bộ hệ thống phông chữ của website. \\
\hline
Tự động nén tối ưu CDN & Google tự động phát hiện trình duyệt và gửi định dạng font nén tốt nhất (như \codeinline{.woff2}). \\
\hline
\end{tabularx}
\end{prettyTableBox}

\paragraph{4.2. Nhúng Nhiều Phông Chữ Với Quy Tắc @font-face (Tự Host Multi-Font):}

Khi tự host nhiều phông chữ trên máy chủ riêng, bạn phải tạo từng khối \codeinline{@font-face} độc lập tương ứng với từng họ font và từng biến thể trọng lượng (\codeinline{font-weight}) / kiểu dáng (\codeinline{font-style}).

\begin{codeWindow}[File style.css -- Khai Báo Nhiều Khối @font-face Độc Lập]
\begin{lstlisting}[language=CSS3]
/* 1. Khai bao cho font MyBrandFont - Bien the Regular (400) */
@font-face {
  font-family: 'MyBrandFont'; /* Ten family chung */
  src: url('../fonts/MyBrandFont-Regular.woff2') format('woff2'),
       url('../fonts/MyBrandFont-Regular.woff') format('woff');
  font-weight: 400;
  font-style: normal;
  font-display: swap;
}

/* 2. Khai bao cho font MyBrandFont - Bien the Bold (700) */
@font-face {
  font-family: 'MyBrandFont'; /* Cung ten family, nhung dinh nghia weight 700 */
  src: url('../fonts/MyBrandFont-Bold.woff2') format('woff2'),
       url('../fonts/MyBrandFont-Bold.woff') format('woff');
  font-weight: 700;
  font-style: normal;
  font-display: swap;
}

/* 3. Khai bao cho font MyHeadlineFont - Bien the Regular */
@font-face {
  font-family: 'MyHeadlineFont'; /* Ten family khac cho tieu de */
  src: url('../fonts/MyHeadlineFont.woff2') format('woff2'),
       url('../fonts/MyHeadlineFont.woff') format('woff');
  font-weight: normal;
  font-style: normal;
  font-display: swap;
}

/* 4. Ap dung trong CSS */
h1 {
  font-family: 'MyHeadlineFont', serif;
}

p {
  font-family: 'MyBrandFont', sans-serif;
  font-weight: 400; /* Trinh duyet tu dong chon tesp MyBrandFont-Regular */
}

strong {
  font-family: 'MyBrandFont', sans-serif;
  font-weight: 700; /* Trinh duyet tu dong chon tesp MyBrandFont-Bold */
}
\end{lstlisting}
\end{codeWindow}

\begin{browserWindow}[Kết Quả Hiển Thị: Áp Dụng Nhiều Custom Font Tự Host]
\begin{tcolorbox}[arc=4pt, colback=white, colframe=primaryBlue!30, boxrule=0.8pt, top=10pt, bottom=10pt, left=12pt, right=12pt]
  {\sffamily\bfseries\color{primaryBlue}\Large Tiêu đề h1 dùng MyHeadlineFont (serif)}\par\vspace{6pt}
  {\sffamily\small Đoạn văn p dùng \codeinline{MyBrandFont} weight 400. Nhấn mạnh \textbf{văn bản bold dùng MyBrandFont weight 700} hoạt động chuẩn xác nhờ phân loại \codeinline{font-weight} trong từng khối \codeinline{@font-face}.}
\end{tcolorbox}
\end{browserWindow}

\begin{warnBox}[Lưu Ý Quan Trọng Khi Nhúng Nhiều Font Qua @font-face]
\begin{itemize}[leftmargin=1.5em, itemsep=3pt, parsep=0pt]
    \item \textbf{Đồng bộ tên \codeinline{font-family}}: Đặt tên \codeinline{font-family} giống nhau cho các biến thể cùng một họ phông chữ (ví dụ: \codeinline{'MyBrandFont'}), nhưng phân biệt bằng \codeinline{font-weight} và \codeinline{font-style}.
    \item \textbf{Khai báo riêng từng khối}: Mỗi sự kết hợp giữa trọng lượng (Thin, Regular, Bold) và kiểu dáng (Normal, Italic) phải có một khối \codeinline{@font-face} riêng biệt.
    \item \textbf{Cân nhắc tải mạng}: Mỗi tệp font tự host là một yêu cầu HTTP riêng. Nhúng quá 3 phông chữ có thể làm giảm đáng kể tốc độ tải trang.
\end{itemize}
\end{warnBox}

\paragraph{4.3. Mẹo \& Thực Hành Tốt Nhất Khi Nhúng Nhiều Phông Chữ (Best Practices):}

\begin{tipBox}[6 Quy Tắc Vàng Khi Nhúng Nhiều Phông Chữ Trên Web]
\begin{enumerate}[leftmargin=1.5em, itemsep=3pt, parsep=0pt]
    \item \textbf{Chỉ nhúng những gì thực sự cần}: Giới hạn số lượng font family từ $1 \rightarrow 3$ họ phông chữ trên một website. Chỉ chọn các biến thể weight thực sự sử dụng (như 400 và 700).
    \item \textbf{Ưu tiên Google Fonts hoặc CDN}: Sử dụng hạ tầng CDN toàn cầu giúp tăng tốc độ tải tệp font qua bộ nhớ đệm trình duyệt (\term{browser cache}).
    \item \textbf{Bắt buộc dùng \codeinline{font-display: swap;}}: Giúp hiển thị văn bản ngay bằng font dự phòng trong khi chờ tải font chính, loại bỏ hoàn toàn hiện tượng chữ bị ẩn (\term{FOIT}).
    \item \textbf{Sắp xếp thứ tự thẻ \codeinline{<link>} chuẩn}: Đặt thẻ \codeinline{<link>} nhúng font lên đầu phần \codeinline{<head>}, trước tệp CSS chính để trình duyệt gửi yêu cầu tải sớm nhất.
    \item \textbf{Luôn cung cấp font dự phòng (Fallback Fonts)}: Kết thúc chuỗi \codeinline{font-family} bằng các từ khóa chung (\codeinline{sans-serif}, \codeinline{serif}, \codeinline{monospace}).
    \item \textbf{Kiểm tra hiệu suất với Google PageSpeed Insights \& Lighthouse}: Đánh giá chỉ số \term{CLS (Cumulative Layout Shift)} và \term{FCP} sau khi nhúng để đảm bảo không bị giật trang.
\end{enumerate}
\end{tipBox}



% ------------------------------------------------------------------------------
\paragraph{Lời Khuyên Chung Về Hiệu Suất Khi Nhúng Font (Font Performance Best Practices):}

\begin{tipBox}[Các Quy Tắc Vàng Tối Ưu Hiệu Suất Tải Font Trên Web]
\begin{enumerate}[leftmargin=1.5em, itemsep=3pt, parsep=0pt]
    \item \textbf{Giảm số lượng font và style}: Mỗi biến thể phông chữ đều là một tệp riêng cần tải. Chỉ nhúng những style thực sự sử dụng trên giao diện.
    \item \textbf{Sử dụng \codeinline{font-display: swap;}}: Đây là giải pháp quan trọng nhất giúp loại bỏ hiện tượng chữ bị ẩn (\term{FOIT}) và tăng chỉ số hiển thị nội dung đầu tiên (\term{FCP}).
    \item \textbf{Tối ưu hóa tệp font}:
    \begin{itemize}[leftmargin=1.2em, itemsep=2pt]
        \item Ưu tiên sử dụng định dạng \codeinline{.woff2} để đạt hiệu quả nén cao nhất.
        \item Áp dụng kỹ thuật \term{"subsetting"} (loại bỏ các bộ ký tự không sử dụng đến) để giảm $50\% \rightarrow 80\%$ dung lượng tệp font.
    \end{itemize}
    \item \textbf{Sử dụng kỹ thuật Preconnect và Preload}:
    \begin{itemize}[leftmargin=1.2em, itemsep=2pt]
        \item \codeinline{rel="preconnect"}: Chủ động thiết lập kết nối mạng sớm tới máy chủ font.
        \item \codeinline{rel="preload"}: Yêu cầu trình duyệt ưu tiên tải tệp font stylesheet càng sớm càng tốt.
    \end{itemize}
    \begin{codeWindow}[Cú Pháp Preconnect \& Preload Chuẩn]
\begin{lstlisting}[language=HTML5]
<link rel="preconnect" href="https://fonts.gstatic.com" crossorigin>
<link rel="preload" href="https://fonts.googleapis.com/css2?family=Roboto:wght@400;700&display=swap" as="style">
\end{lstlisting}
    \end{codeWindow}
    \textit{Lưu ý}: Thuộc tính \codeinline{as="style"} rất quan trọng để trình duyệt nhận biết tài nguyên tải trước là một stylesheet.
\end{enumerate}
\end{tipBox}

Bằng cách lựa chọn phương pháp nhúng font phù hợp và áp dụng các mẹo tối ưu hiệu suất trên, bạn có thể đảm bảo rằng trang web của mình không chỉ đẹp mắt với các phông chữ tùy chỉnh độc đáo mà còn có tốc độ tải nhanh chóng, mang lại trải nghiệm người dùng tuyệt vời!

\subsection{Mẹo \& Best Practices Tối Ưu Typography Trong CSS}

\begin{prettyTableBox}[Tổng Hợp Thực Hành Tốt Nhất (Best Practices) Cho CSS Typography]
\sffamily\small
\begin{tabularx}{\linewidth}{|>{\raggedright\arraybackslash \bfseries\color{primaryBlue}}p{3.8cm}|>{\raggedright\arraybackslash}X|}
\hline
\rowcolor{primaryBlue} \thd{\sffamily Tiêu chí tối ưu} & \thd{\sffamily Giải pháp \& Kỹ thuật khuyến nghị} \\
\hline
\rowcolor{primaryBlue!10} \multicolumn{2}{|l|}{\sffamily\bfseries\color{primaryBlue} 1. Tối Ưu Hiệu Suất (Performance Optimization)} \\
\hline
Giảm số lượng font & Chỉ nhúng những font và style (Regular 400, Bold 700, Italic) thực sự cần thiết trong thiết kế. \\
\hline
Dùng font-display: swap & Tránh hiện tượng nhấp nháy chữ trắng (\term{FOIT}) và cải thiện điểm chỉ số Core Web Vitals (LCP, CLS). \\
\hline
Nén tệp font (Subsetting) & Sử dụng công cụ như Font Squirrel Webfont Generator để nén tệp font WOFF2 và loại bỏ các ký tự không dùng đến. \\
\hline
Kỹ thuật Preconnect / Preload & Thiết lập kết nối sớm bằng \codeinline{<link rel="preconnect" href="https://fonts.gstatic.com" crossorigin>} và \codeinline{<link rel="preload" href="..." as="style">}. \\
\hline
\rowcolor{primaryBlue!10} \multicolumn{2}{|l|}{\sffamily\bfseries\color{primaryBlue} 2. Tối Ưu Khả Năng Đọc (Readability \& Ergonomics)} \\
\hline
Kích thước font phù hợp & Đảm bảo kích thước chữ đủ lớn trên mọi thiết bị. Kích thước cơ bản tiêu chuẩn trên Desktop thường là \codeinline{16px}. \\
\hline
Chiều cao dòng (line-height) & Đặt \codeinline{line-height} từ \codeinline{1.5} đến \codeinline{1.8} để tạo khoảng cách đọc thoải mái, tránh mỏi mắt. \\
\hline
Độ tương phản màu sắc & Đảm bảo màu chữ có đủ độ tương phản (\term{Contrast Ratio}) với màu nền theo tiêu chuẩn WCAG. \\
\hline
Text rendering engine & Cải thiện hiển thị nét chữ bằng cách áp dụng \codeinline{text-rendering: optimizeLegibility;}. \\
\hline
\rowcolor{primaryBlue!10} \multicolumn{2}{|l|}{\sffamily\bfseries\color{primaryBlue} 3. Thiết Kế Đáp Ứng (Responsive Typography)} \\
\hline
Đơn vị tương đối linh hoạt & Sử dụng \codeinline{rem}, \codeinline{em}, \codeinline{vw} cho \codeinline{font-size} để font tự động co giãn linh hoạt theo màn hình. \\
\hline
Media Queries co giãn & Điều chỉnh \codeinline{font-size} cho từng mốc màn hình cụ thể bằng \codeinline{@media (max-width: 768px) \{ ... \}}. \\
\hline
\end{tabularx}
\end{prettyTableBox}

\subsection{Bảng Cheatsheet Tóm Tắt Nhanh Các Thuộc Tính Font Trong CSS}

\begin{prettyTableBox}[Bảng Tra Cứu Cheatsheet Nhanh Các Thuộc Tính Font Trong CSS]
\sffamily\small
\begin{tabularx}{\linewidth}{|>{\raggedright\arraybackslash \ttfamily\color{primaryBlue}\bfseries}p{2.6cm}|>{\raggedright\arraybackslash}p{3.2cm}|>{\raggedright\arraybackslash}p{4.0cm}|>{\ttfamily\color{accentOrange}}X|}
\hline
\rowcolor{primaryBlue} \thd{\sffamily Thuộc tính} & \thd{\sffamily Chức năng chính} & \thd{\sffamily Các giá trị chấp nhận} & \thd{\sffamily Ví dụ minh họa CSS} \\
\hline
font-family & Chọn họ phông chữ cho văn bản. & Danh sách font chữ (ưu tiên từ trái qua phải, kết thúc bằng generic family). & font-family: Arial, sans-serif; \\
\hline
font-size & Đặt kích thước chữ. & px, em, rem, \%, vw, vh, larger, smaller, xx-small...xx-large. & font-size: 16px; \\
\hline
font-weight & Đặt độ đậm của chữ. & normal, bold, lighter, bolder, hoặc số từ 100 đến 900. & font-weight: bold; \\
\hline
font-style & Đặt kiểu chữ nghiêng. & normal, italic, oblique. & font-style: italic; \\
\hline
font-variant & Đặt biến thể chữ hoa nhỏ. & normal, small-caps. & font-variant: small-caps; \\
\hline
line-height & Đặt chiều cao dòng văn bản. & normal, số không đơn vị (1.5), px, em, rem, \%. & line-height: 1.6; \\
\hline
font & Thuộc tính viết tắt shorthand. & [style] [variant] [weight] size[/line-height] family. & font: 16px/1.5 Arial, sans-serif; \\
\hline
@font-face & Quy tắc nhúng font ngoài tự host. & Tệp font với src, url(), format(), font-display. & @font-face \{ font-family: 'Custom'; src: url('font.woff2') format('woff2'); \} \\
\hline
font-kerning & Điều chỉnh khoảng cách giữa các cặp ký tự. & auto, normal, none. & font-kerning: normal; \\
\hline
\end{tabularx}
\end{prettyTableBox}

\subsection{Đọc Thêm: Variable Fonts, CSS Font Loading API \& Text Rendering Engine}
\label{sec:doc_them_css_fonts}

\begin{conceptBox}[Mở Rộng Kiến Thức Nâng Cao Về Typography Trong CSS]
\begin{enumerate}[leftmargin=1.5em, itemsep=4pt, parsep=0pt]
    \item \textbf{Variable Fonts (Font Biến Thể CSS Fonts Level 4)}: Cho phép đóng gói toàn bộ các trọng lượng (\codeinline{wght}), độ rộng (\codeinline{wdth}), độ nghiêng (\codeinline{slnt}) vào trong \textbf{chỉ 1 tệp font duy nhất} (.woff2), giúp tiết kiệm tới 80\% dung lượng tải mạng!
    \item \textbf{CSS Font Loading API Trong JavaScript}: Giúp lập trình viên chủ động kiểm soát trạng thái tải font bằng JavaScript thông qua đối tượng \codeinline{document.fonts}.
    \item \textbf{Ngăn Trình Duyệt Tự Sinh Font Giả (\codeinline{font-synthesis})}: Khi font không có bản Bold hoặc Italic thực sự, trình duyệt có thể tự bóp méo ký tự để làm nghiêng/đậm giả. Dùng \codeinline{font-synthesis: none;} để ngăn chặn.
\end{enumerate}
\end{conceptBox}

\begin{codeWindow}[1. Cú Pháp Khai Báo Variable Fonts Trong CSS]
\begin{lstlisting}[language=CSS3]
h1 {
  font-family: 'Roboto Flex', sans-serif;
  font-variation-settings: 'wght' 650, 'wdth' 115, 'slnt' -5;
}
\end{lstlisting}
\end{codeWindow}

\begin{codeWindow}[2. Kiểm Soát Trạng Thái Tải Font Bằng JavaScript Font Loading API]
\begin{lstlisting}[language=JavaScript]
document.fonts.load('1em Roboto').then(() => {
  console.log('Font Roboto da san sang hien thi!');
});
\end{lstlisting}
\end{codeWindow}

\begin{codeWindow}[3. Cấu Hình font-synthesis Tránh Sinh Chữ Giả]
\begin{lstlisting}[language=CSS3]
body {
  font-synthesis: none; /* Khong cho trinh duyet tu lam gia chu dam / nghieng */
}
\end{lstlisting}
\end{codeWindow}

\begin{browserWindow}[Kết Quả: Tối Ưu Font Cao Cấp Với Variable Fonts \& font-synthesis]
\begin{tcolorbox}[arc=4pt, colback=white, colframe=primaryBlue!30, boxrule=0.8pt, top=10pt, bottom=10pt, left=12pt, right=12pt]
  {\sffamily\bfseries\color{primaryBlue} Minh họa kết quả Variable Font \& font-synthesis:}\par\vspace{6pt}
  {\sffamily\fontsize{14}{18}\selectfont\fontseries{b}\selectfont Tiêu đề co giãn độ đậm 650 mượt mà từ 1 tệp .woff2 duy nhất.}\par\vspace{4pt}
  {\sffamily\small\color{gray!70} → \codeinline{font-synthesis: none;} bảo vệ nét thiết kế chuẩn của phông chữ gốc, loại bỏ hoàn toàn hiện tượng nghiêng/đậm giả lập.}
\end{tcolorbox}
\end{browserWindow}
% MỤC 5: PHÂN LOẠI BỘ CHỌN CSS (CSS SELECTORS MASTERCLASS)
% ==============================================================================
\section{Các Loại Selectors Trong CSS (CSS Selectors Masterclass)}


Bộ chọn (\term{Selector}) là cách để xác định các phần tử HTML mà bạn muốn áp dụng kiểu dáng (\term{CSS styles}). CSS cung cấp nhiều loại selectors để bạn có thể chọn phần tử theo tên thẻ, ID, class, trạng thái, hoặc cấu trúc của chúng.

\subsection{1. Bảng Tổng Quan Chi Tiết Các Bộ Chọn CSS Quan Trọng}

\begin{prettyTableBox}[Bảng Tra Cứu Danh Sách Đầy Đủ Các Bộ Chọn CSS Quan Trọng]
\rowcolors{2}{white}{primaryBlue!4}
\sffamily\small
\begin{tabularx}{\linewidth}{|>{\raggedright\arraybackslash \bfseries\color{primaryBlue}}K{3.8cm}|>{\ttfamily\color{darkText}\bfseries}K{3.5cm}|Z|}
\hline
\rowcolor{primaryBlue} \thd{Loại Selector} & \thd{Cú pháp} & \thd{Công dụng} \\
\hline
Bộ chọn theo thẻ (Type Selector) & tagname \{\} & Chọn tất cả phần tử có tên thẻ cụ thể. \\
\hline
Bộ chọn theo class (Class Selector) & .classname \{\} & Chọn tất cả phần tử có class chỉ định. \\
\hline
Bộ chọn theo ID (ID Selector) & \#idname \{\} & Chọn phần tử có ID cụ thể (duy nhất). \\
\hline
Bộ chọn nhóm (Group Selector) & tag1, tag2 \{\} & Áp dụng chung CSS cho nhiều phần tử. \\
\hline
Bộ chọn chung (Universal Selector) & * \{\} & Chọn tất cả các thành phần HTML trên trang. \\
\hline
Bộ chọn con trực tiếp (Child Selector) & parent > child \{\} & Chọn phần tử con trực tiếp của phần tử cha. \\
\hline
Bộ chọn con (mọi cấp) (Descendant Selector) & parent child \{\} & Chọn tất cả phần tử con bên trong phần tử cha. \\
\hline
Bộ chọn anh em liền kề (Adjacent Sibling) & element1 + element2 \{\} & Chọn phần tử ngay sau một phần tử khác cùng cấp. \\
\hline
Bộ chọn anh em chung (General Sibling) & element1 \textasciitilde{} element2 \{\} & Chọn tất cả phần tử cùng cấp phía sau một phần tử. \\
\hline
Bộ chọn thuộc tính (Attribute Selector) & [attribute] \{\} & Chọn phần tử có một thuộc tính cụ thể. \\
\hline
Bộ chọn giá trị thuộc tính cụ thể & [attribute="value"] \{\} & Chọn phần tử có thuộc tính bằng giá trị cụ thể. \\
\hline
Bộ chọn giá trị thuộc tính chứa đoạn & [attribute*="value"] \{\} & Chọn phần tử có giá trị thuộc tính chứa đoạn text nhất định. \\
\hline
Pseudo-class (Lớp giả) & selector:pseudo-class \{\} & Chọn phần tử dựa vào trạng thái (hover, focus, checked, first-child, last-child, nth-child,...). \\
\hline
Pseudo-element (Phần tử giả) & selector::pseudo-element \{\} & Chọn một phần của nội dung (first-letter, first-line, before, after). \\
\hline
\end{tabularx}
\end{prettyTableBox}

\subsection{2. Chi Tiết Các Bộ Chọn Cơ Bản (Basic Selectors)}

\begin{prettyTableBox}[Bảng Tổng Hợp Chi Tiết Các Bộ Chọn Cơ Bản (Basic Selectors)]
\rowcolors{2}{white}{primaryBlue!4}
\sffamily\small
\begin{tabularx}{\linewidth}{|>{\raggedright\arraybackslash \bfseries\color{primaryBlue}}K{3.8cm}|>{\ttfamily\color{darkText}\bfseries}K{3.2cm}|Z|}
\hline
\rowcolor{primaryBlue} \thd{Loại Selector} & \thd{Cú pháp} & \thd{Công dụng} \\
\hline
Bộ chọn theo thẻ (Type Selector) & tagname \{\} & Chọn tất cả phần tử có tên thẻ cụ thể. \\
\hline
Bộ chọn theo class (Class Selector) & .classname \{\} & Chọn tất cả phần tử có class chỉ định. \\
\hline
Bộ chọn theo ID (ID Selector) & \#idname \{\} & Chọn phần tử có ID cụ thể (duy nhất). \\
\hline
Bộ chọn nhóm (Group Selector) & tag1, tag2 \{\} & Áp dụng chung CSS cho nhiều phần tử. \\
\hline
Bộ chọn chung (Universal Selector) & * \{\} & Chọn tất cả các thành phần HTML trên trang. \\
\hline
\end{tabularx}
\end{prettyTableBox}

\subsubsection{2.1. Bộ Chọn Theo Thẻ (Type Selector / Tag Selector)}

\paragraph{a. Khái niệm:}
Bộ chọn theo thẻ (Type Selector) là bộ chọn cơ bản trong CSS, dùng để chọn tất cả các phần tử có cùng tên thẻ HTML và áp dụng chung một quy tắc CSS cho chúng.

\paragraph{b. Cú pháp:}
\begin{codeWindow}[Cú Pháp Type Selector]
\begin{lstlisting}[language=CSS3]
tagname {
    thuoc_tinh: gia_tri;
}
\end{lstlisting}
\end{codeWindow}
\begin{itemize}[leftmargin=1.5em, itemsep=2pt, parsep=0pt]
    \item \codeinline{tagname}: Tên thẻ HTML cần chọn.
    \item \codeinline{thuoc\_tinh}: Thuộc tính CSS sẽ được áp dụng.
    \item \codeinline{gia\_tri}: Giá trị của thuộc tính CSS.
\end{itemize}

\paragraph{c. Các ví dụ minh họa:}
\begin{codeWindow}[Ví Dụ 1: Định Dạng Tất Cả Đoạn Văn <p>]
\begin{lstlisting}[language=CSS3]
p {
    color: blue;
    font-size: 18px;
}
\end{lstlisting}
\end{codeWindow}
\textcolor{tipBorder}{\textbf{[OK] Kết quả}}: Tất cả các thẻ \codeinline{<p>} trong trang web sẽ có màu chữ xanh và kích thước chữ 18px.

\begin{codeWindow}[Ví Dụ 2: Định Dạng Các Tiêu Đề h1, h2, h3]
\begin{lstlisting}[language=CSS3]
h1 {
    color: red;
    text-align: center;
}
h2 {
    color: green;
}
h3 {
    color: orange;
}
\end{lstlisting}
\end{codeWindow}
\textcolor{tipBorder}{\textbf{[OK] Kết quả}}: \codeinline{<h1>} màu đỏ căn giữa, \codeinline{<h2>} màu xanh, \codeinline{<h3>} màu cam.

\begin{codeWindow}[Ví Dụ 3: Áp Dụng Font Chữ Cho Nhiều Thẻ Cùng Lúc]
\begin{lstlisting}[language=CSS3]
h1, h2, h3 {
    font-family: Arial, sans-serif;
}
\end{lstlisting}
\end{codeWindow}
\textcolor{tipBorder}{\textbf{[OK] Kết quả}}: Các tiêu đề \codeinline{<h1>}, \codeinline{<h2>}, \codeinline{<h3>} đều sử dụng font chữ Arial.

\paragraph{d. Cách hoạt động:}
Trình duyệt quét HTML để tìm tất cả các phần tử có tên thẻ phù hợp $\rightarrow$ Áp dụng quy tắc CSS $\rightarrow$ Hiển thị kết quả đồng bộ.

\begin{codeWindow}[Mã HTML Và CSS Kiểm Thử Cách Hoạt Động Thẻ]
\begin{lstlisting}[language=HTML5]
<h1>Header Title</h1>
<h2>Header Title</h2>
<p>Paragraph text</p>
<p>Paragraph text</p>
\end{lstlisting}

\begin{lstlisting}[language=CSS3]
p {
    color: blue;
    font-size: 16px;
}
h1 {
    text-align: center;
    color: red;
}
\end{lstlisting}
\end{codeWindow}

\begin{browserWindow}[Kết Quả Hiển Thị Cách Hoạt Động Type Selector]
\sffamily\small
\centering{\color{red}\bfseries\Large Tiêu đề chính} \par\vspace{2pt}
\flushleft
\textcolor{darkText}{\Large\bfseries Tiêu đề phụ} \par\vspace{2pt}
\textcolor{primaryBlue}{Đây là một đoạn văn bản.} \par\vspace{2pt}
\textcolor{primaryBlue}{Đây là một đoạn khác.}
\end{browserWindow}

\paragraph{e. Ưu điểm \& Hạn chế của Type Selector:}

\begin{tipBox}[Ưu Điểm \& Khi Nào Nên Dùng Type Selector]
\begin{itemize}[leftmargin=1.5em, itemsep=3pt, parsep=0pt]
    \item \textbf{Khi nào nên dùng}: Áp dụng quy tắc CSS chung cho tất cả phần tử cùng loại mà không cần gán class hay ID (ví dụ: tất cả \codeinline{<p>}, \codeinline{<h1>}, \codeinline{<table>}, \codeinline{<button>}).
    \item \textbf{Ưu điểm}: Dễ sử dụng, tiết kiệm thời gian, tăng tính nhất quán giao diện và tạo bố cục nhanh chóng.
\end{itemize}
\end{tipBox}

\begin{codeWindow}[Ví Dụ Thiết Lập Kiểu Dáng Đồng Bộ Cho Nút Bấm <button>]
\begin{lstlisting}[language=CSS3]
button {
    background-color: #ff6600;
    color: white;
    padding: 10px 20px;
    border: none;
    cursor: pointer;
    border-radius: 4px;
}
\end{lstlisting}
\end{codeWindow}

\begin{browserWindow}[Kết Quả Hiển Thị: Nút Bấm <button> Đồng Bộ Giao Diện]
\sffamily\small\centering
\tcbox[on line, arc=4pt, colback=orange!90!black, colframe=orange!90!black, boxsep=4pt, left=12pt, right=12pt]{\color{white}\bfseries Đăng Nhập} \quad
\tcbox[on line, arc=4pt, colback=orange!90!black, colframe=orange!90!black, boxsep=4pt, left=12pt, right=12pt]{\color{white}\bfseries Đăng Ký} \par\vspace{4pt}
\textcolor{gray!70}{\small\textit{($\rightarrow$ Kết quả: Tất cả các nút bấm <button> trên trang sẽ có cùng kiểu dáng da cam đồng bộ)}}
\end{browserWindow}

\begin{warnBox}[Hạn Chế Của Type Selector]
\textbf{Hạn chế}: Nếu chỉ muốn định dạng một số phần tử cụ thể thay vì tất cả, nên dùng \textbf{Class Selector} (\codeinline{.classname}) hoặc \textbf{ID Selector} (\codeinline{\#idname}) để tránh ảnh hưởng ngoài ý muốn đến các phần tử cùng loại khác trên trang.
\end{warnBox}

\subsubsection{2.2. Bộ Chọn Theo Lớp (Class Selector)}

\paragraph{a. Khái niệm \& Cú pháp:}
Bộ chọn class chọn các phần tử HTML có thuộc tính \codeinline{class} nhất định.
\begin{codeWindow}[Cú Pháp Class Selector]
\begin{lstlisting}[language=CSS3]
.classname {
    thuoc_tinh: gia_tri;
}
\end{lstlisting}
\end{codeWindow}
\begin{noteBox}[Lưu Ý Đặt Tên Class]
Tên class không được bắt đầu bằng số và không chứa ký tự đặc biệt (trừ dấu gạch ngang \codeinline{-} và gạch dưới \codeinline{\_}). Một phần tử có thể có nhiều class, ngăn cách bằng dấu cách.
\end{noteBox}

\paragraph{b. Ví dụ minh họa HTML + CSS:}
\begin{codeWindow}[Mã HTML Sử Dụng Nhiều Class]
\begin{lstlisting}[language=HTML5]
<h1 class="title">Welcome</h1>
<p class="content">Paragraph text</p>
<p class="content highlight">Paragraph text</p>
<button class="btn">Sample Text</button>
\end{lstlisting}

\begin{lstlisting}[language=CSS3]
/* Select element */
.title {
    color: #ff6600;
    text-align: center;
    font-size: 28px;
}
/* Select element */
.content {
    color: #333;
    font-size: 16px;
    line-height: 1.5;
}
/* Select element */
.highlight {
    background-color: yellow;
    font-weight: bold;
}
/* Select element */
.btn {
    background-color: #007bff;
    color: white;
    padding: 10px 20px;
    border: none;
    cursor: pointer;
    border-radius: 5px;
}
\end{lstlisting}
\end{codeWindow}

\begin{browserWindow}[Kết Quả Hiển Thị Class Selector]
\sffamily\small
\centering{\color{orange!90!black}\bfseries\Huge Chào mừng bạn đến với trang web!} \par\vspace{4pt}
\flushleft
\textcolor{darkText}{Đây là đoạn văn bản nội dung.} \par\vspace{2pt}
\colorbox{yellow!80}{\textcolor{darkText}{\textbf{Đây là đoạn văn bản nổi bật.}}} \par\vspace{4pt}
\tcbox[on line, arc=5pt, colback=blue!80!white, colframe=blue!80!white, boxsep=4pt, left=10pt, right=10pt]{\color{white}\bfseries Nhấn vào đây}
\end{browserWindow}

\paragraph{c. Ưu/Nhược điểm \& Tóm tắt nhanh:}

\begin{tipBox}[Ưu Điểm \& Khi Nào Nên Dùng Class Selector]
\begin{itemize}[leftmargin=1.5em, itemsep=2pt, parsep=0pt]
    \item \textbf{Tái sử dụng cao}: Một class có thể áp dụng cho nhiều phần tử khác nhau trên toàn bộ trang web.
    \item \textbf{Dễ bảo trì}: Thay đổi kiểu dáng ở một nơi duy nhất trong file CSS, tất cả thẻ HTML có class đó sẽ tự động cập nhật.
    \item \textbf{Tránh ảnh hưởng toàn cục}: Không như bộ chọn theo thẻ, class chỉ tác động đến các phần tử cụ thể được khai báo.
\end{itemize}
\end{tipBox}

\begin{warnBox}[Hạn Chế Của Class Selector]
\begin{itemize}[leftmargin=1.5em, itemsep=2pt, parsep=0pt]
    \item \textbf{Độ ưu tiên thấp hơn ID}: Nếu một phần tử có cả class và ID, kiểu dáng trong ID sẽ ghi đè class.
    \item \textbf{Có thể làm CSS phức tạp}: Nếu lạm dụng quá nhiều class chồng chéo, mã CSS sẽ trở nên khó đọc và khó bảo trì.
\end{itemize}
\end{warnBox}

\begin{prettyTableBox}[Tóm Tắt Nhanh Về Class Selector]
\rowcolors{2}{white}{primaryBlue!4}
\sffamily\small
\begin{tabularx}{\linewidth}{|>{\raggedright\arraybackslash \bfseries\color{primaryBlue}}K{3.5cm}|Z|}
\hline
\rowcolor{primaryBlue} \thd{Đặc điểm} & \thd{Chi tiết} \\
\hline
Cú pháp & \codeinline{.classname \{ thuộc\_tính: giá\_trị; \}} \\
\hline
Chọn phần tử & Chọn tất cả các phần tử có class tương ứng \\
\hline
Ký hiệu & Dấu chấm \codeinline{.} đứng trước tên class \\
\hline
Độ ưu tiên & Thấp hơn ID, nhưng cao hơn bộ chọn theo thẻ \\
\hline
Tính tái sử dụng & Cao -- có thể áp dụng cùng class cho nhiều thẻ \\
\hline
Ứng dụng thực tế & Thiết kế layout linh hoạt, tái sử dụng kiểu dáng, giảm lặp mã \\
\hline
\end{tabularx}
\end{prettyTableBox}

\subsubsection{2.3. Bộ Chọn Theo ID (ID Selector)}

\paragraph{a. Khái niệm \& Cú pháp:}
ID là một định danh duy nhất cho mỗi phần tử HTML. Bộ chọn ID dùng dấu thăng \codeinline{\#}.
\begin{codeWindow}[Cú Pháp ID Selector]
\begin{lstlisting}[language=CSS3]
#idname {
    thuoc_tinh: gia_tri;
}
\end{lstlisting}
\end{codeWindow}

\paragraph{b. Ví dụ minh họa HTML + CSS:}
\begin{codeWindow}[Mã HTML Sử Dụng ID Duy Nhất]
\begin{lstlisting}[language=HTML5]
<h1 id="main-title">Welcome</h1>
<p id="intro">Paragraph text</p>
<button id="submit-btn">Submit</button>
\end{lstlisting}

\begin{lstlisting}[language=CSS3]
/* Select element */
#main-title {
    color: #ff6600;
    font-size: 32px;
    text-align: center;
}
/* Select element */
#intro {
    color: #444;
    font-style: italic;
    font-size: 18px;
}
/* Select element */
#submit-btn {
    background-color: #007bff;
    color: white;
    padding: 10px 20px;
    border: none;
    border-radius: 5px;
    cursor: pointer;
}
\end{lstlisting}
\end{codeWindow}

\begin{browserWindow}[Kết Quả Hiển Thị ID Selector]
\sffamily\small
\centering{\color{orange!90!black}\bfseries\Huge Chào mừng bạn đến với trang web!} \par\vspace{4pt}
\flushleft
\textcolor{gray!80}{\large\textit{Đây là đoạn giới thiệu về trang web.}} \par\vspace{4pt}
\tcbox[on line, arc=5pt, colback=blue!80!white, colframe=blue!80!white, boxsep=4pt, left=10pt, right=10pt]{\color{white}\bfseries Gửi}
\end{browserWindow}

\paragraph{c. Ưu/Nhược điểm \& Khi nào nên dùng ID Selector:}

\begin{tipBox}[Ưu Điểm \& Khi Nào Nên Dùng ID Selector]
\begin{itemize}[leftmargin=1.5em, itemsep=2pt, parsep=0pt]
    \item \textbf{Chính xác tuyệt đối}: Định dạng chính xác một phần tử duy nhất (như tiêu đề chính, header, footer).
    \item \textbf{Độ ưu tiên cao}: Dễ dàng ghi đè các kiểu dáng class và tag khi cần thiết.
    \item \textbf{Tối ưu với JavaScript}: Rất tiện lợi khi kết hợp với \codeinline{document.getElementById()}.
\end{itemize}
\end{tipBox}

\begin{warnBox}[Hạn Chế Của ID Selector]
\begin{itemize}[leftmargin=1.5em, itemsep=2pt, parsep=0pt]
    \item \textbf{Không tái sử dụng}: Mỗi ID là duy nhất trên trang, dùng nhiều lần là vi phạm chuẩn W3C HTML.
    \item \textbf{Khó bảo trì}: Lạm dụng ID sẽ làm tăng độ phức tạp khi muốn mở rộng giao diện.
\end{itemize}
\end{warnBox}

\begin{prettyTableBox}[Tóm Tắt Nhanh Về ID Selector]
\rowcolors{2}{white}{primaryBlue!4}
\sffamily\small
\begin{tabularx}{\linewidth}{|>{\raggedright\arraybackslash \bfseries\color{primaryBlue}}K{3.5cm}|Z|}
\hline
\rowcolor{primaryBlue} \thd{Đặc điểm} & \thd{Chi tiết} \\
\hline
Cú pháp & \codeinline{\#idname \{ thuộc\_tính: giá\_trị; \}} \\
\hline
Ký hiệu & Dấu thăng \codeinline{\#} đứng trước tên ID \\
\hline
Độ ưu tiên & Rất cao -- cao hơn class và type selector \\
\hline
Tính tái sử dụng & Không -- mỗi ID chỉ dùng cho 1 phần tử duy nhất \\
\hline
Ứng dụng thực tế & Định dạng phần tử duy nhất, làm việc với JavaScript \\
\hline
\end{tabularx}
\end{prettyTableBox}

\subsubsection{2.4. Bộ Chọn Chung (Universal Selector - *)}

Bộ chọn chung dùng dấu sao \codeinline{*} giúp bạn chọn tất cả các phần tử HTML trên trang web.

\begin{codeWindow}[Cú Pháp Universal Selector]
\begin{lstlisting}[language=CSS3]
* {
    thuoc_tinh: gia_tri;
}
\end{lstlisting}
\end{codeWindow}

\paragraph{Các ví dụ ứng dụng phổ biến:}
\begin{codeWindow}[Ví Dụ 1: Reset Margin, Padding Và Thiết Lập Box-Sizing]
\begin{lstlisting}[language=CSS3]
* {
    margin: 0;
    padding: 0;
    box-sizing: border-box;
}
\end{lstlisting}
\end{codeWindow}

\begin{codeWindow}[Ví Dụ 2: Đặt Font Chữ Và Màu Nền Mặc Định Toàn Trang]
\begin{lstlisting}[language=CSS3]
* {
    font-family: Arial, sans-serif;
    background-color: #f0f0f0;
    color: #333;
}
\end{lstlisting}
\end{codeWindow}

\begin{codeWindow}[Ví Dụ 3: Thêm Viền Đỏ Để Kiểm Tra Layout (Debug)]
\begin{lstlisting}[language=CSS3]
* {
    border: 1px solid red;
}
\end{lstlisting}
\end{codeWindow}

\begin{browserWindow}[Kết Quả: Hiệu Ứng Bộ Chọn Chung Universal Selector (*)]
\begin{tcolorbox}[arc=4pt, colback=white, colframe=primaryBlue!30, boxrule=0.8pt, top=10pt, bottom=10pt, left=12pt, right=12pt]
  {\sffamily\bfseries\color{primaryBlue} Tác động toàn diện của \codeinline{*}:}\par\vspace{6pt}
  \begin{tcolorbox}[colback=red!5, colframe=red!50, arc=2pt, boxrule=0.8pt, top=4pt, bottom=4pt, left=6pt, right=6pt]
    {\sffamily\small\color{red!80!black} Thẻ \codeinline{<div>} có viền đỏ}
    \begin{tcolorbox}[colback=red!10, colframe=red!60, arc=2pt, boxrule=0.8pt, top=4pt, bottom=4pt, left=6pt, right=6pt]
      {\sffamily\small\color{red!80!black} Thẻ \codeinline{<p>} bên trong cũng có viền đỏ}
    \end{tcolorbox}
  \end{tcolorbox}\par\vspace{4pt}
  {\sffamily\small\color{gray!70} → Bộ chọn \codeinline{*} nhắm trúng mọi thẻ HTML không ngoại lệ.}
\end{tcolorbox}
\end{browserWindow}

\begin{noteBox}[Mẹo Giới Hạn Phạm Vi Của Universal Selector]
Để tránh ảnh hưởng toàn trang, bạn có thể giới hạn phạm vi bằng cách kết hợp với bộ chọn cha:
\codeinline{div * \{ margin: 0; padding: 0; \}} $\rightarrow$ Chỉ áp dụng reset cho các phần tử nằm bên trong thẻ \codeinline{<div>}.
\end{noteBox}

\begin{prettyTableBox}[Tóm Tắt Nhanh Về Universal Selector]
\rowcolors{2}{white}{primaryBlue!4}
\sffamily\small
\begin{tabularx}{\linewidth}{|>{\raggedright\arraybackslash \bfseries\color{primaryBlue}}K{3.5cm}|Z|}
\hline
\rowcolor{primaryBlue} \thd{Đặc điểm} & \thd{Chi tiết} \\
\hline
Cú pháp & \codeinline{* \{ thuộc\_tính: giá\_trị; \}} \\
\hline
Chức năng & Chọn và áp dụng quy tắc CSS cho tất cả phần tử \\
\hline
Ưu điểm & Nhanh, đơn giản, tiện lợi khi cần áp dụng toàn trang \\
\hline
Nhược điểm & Ảnh hưởng quá rộng, giảm hiệu suất, dễ gây xung đột \\
\hline
Trường hợp dùng & Reset CSS, cấu hình toàn trang, kiểm tra layout \\
\hline
\end{tabularx}
\end{prettyTableBox}

\subsection{3. Bộ Chọn Kết Hợp (Combinator Selectors) Chi Tiết \& Ví Dụ Thực Tế}

\begin{conceptBox}[Khái Niệm Về Bộ Chọn Kết Hợp (Combinator Selectors)]
Trong CSS, \textbf{combinator selectors (bộ chọn kết hợp)} là các bộ chọn dùng để kết nối hoặc liên kết các phần tử dựa trên mối quan hệ của chúng trong cây DOM. Chúng giúp bạn chọn phần tử con, phần tử liền kề, hoặc phần tử nằm cùng cấp cực kỳ linh hoạt!
\end{conceptBox}

\begin{prettyTableBox}[Bảng Tổng Quan Các Loại Bộ Chọn Kết Hợp (Combinator Selectors)]
\rowcolors{2}{white}{primaryBlue!4}
\sffamily\small
\begin{tabularx}{\linewidth}{|>{\raggedright\arraybackslash \bfseries\color{primaryBlue}}K{4.2cm}|>{\ttfamily\color{darkText}\bfseries}K{3.0cm}|Z|}
\hline
\rowcolor{primaryBlue} \thd{Loại Selector} & \thd{Cú pháp} & \thd{Công dụng} \\
\hline
Bộ chọn con trực tiếp (Child Selector) & parent > child \{\} & Chọn phần tử con trực tiếp của phần tử cha. \\
\hline
Bộ chọn con (mọi cấp) (Descendant Selector) & parent child \{\} & Chọn tất cả phần tử con bên trong phần tử cha (mọi cấp). \\
\hline
Bộ chọn anh em liền kề (Adjacent Sibling) & element1 + element2 \{\} & Chọn phần tử ngay sau một phần tử khác cùng cấp. \\
\hline
Bộ chọn anh em chung (General Sibling) & element1 \textasciitilde{} element2 \{\} & Chọn tất cả phần tử cùng cấp phía sau một phần tử. \\
\hline
\end{tabularx}
\end{prettyTableBox}

\subsubsection{3.1. Descendant Selector (Bộ Chọn Hậu Duệ) -- Chi Tiết \& Ví Dụ Thực Tế}

\begin{conceptBox}[Khái Niệm Descendant Selector (Bộ Chọn Hậu Duệ)]
\textbf{Descendant Selector (Bộ chọn hậu duệ)} trong CSS dùng để chọn \textbf{tất cả các phần tử con, cháu, chắt\ldots} nằm bên trong một phần tử cha, bất kể cấp độ lồng nhau sâu bao nhiêu trong cây DOM.
\end{conceptBox}

\paragraph{a. Cú pháp và nguyên tắc hoạt động:}
\begin{codeWindow}[Cú Pháp Descendant Selector (Khoảng Trắng)]
\begin{lstlisting}[language=CSS3]
cha con {
    thuoc_tinh: gia_tri;
}
\end{lstlisting}
\end{codeWindow}

\begin{noteBox}[Khoảng Trắng -- Ký Hiệu Kỳ Diệu Của Descendant Selector]
\textbf{Khoảng trắng} giữa hai bộ chọn chính là ký hiệu của Descendant Selector.
\begin{enumerate}[leftmargin=1.5em, itemsep=2pt, parsep=0pt]
    \item \textbf{Bước 1}: Tìm phần tử cha (\codeinline{cha}).
    \item \textbf{Bước 2}: Tìm tất cả phần tử hậu duệ bên trong nó (\codeinline{con}), bất kể cấp độ lồng nhau.
    \item \textbf{Bước 3}: Áp dụng quy tắc CSS cho tất cả các phần tử hậu duệ thỏa mãn.
\end{enumerate}
\end{noteBox}

\paragraph{b. Ví dụ 1: Cơ bản về mức độ lồng nhau (Con, Cháu, Chắt):}
\begin{codeWindow}[Ví Dụ 1: Định Dạng Tất Cả Thẻ p Nằm Trong .container]
\begin{lstlisting}[language=HTML5]
<div class="container">
    <p>Sample Text</p>
    <div class="box">
        <p>Sample Text</p>
        <span>
            <p>Sample Text</p>
        </span>
    </div>
</div>
\end{lstlisting}

\begin{lstlisting}[language=CSS3]
.container p {
    color: blue;
    font-weight: bold;
}
\end{lstlisting}
\end{codeWindow}

\begin{browserWindow}[Kết Quả Hiển Thị Ví Dụ 1 (.container p)]
\sffamily\small
\begin{tcolorbox}[arc=2pt, colback=primaryBlue!5, colframe=primaryBlue!40, boxrule=1pt, title={\sffamily\small\color{white}\textbf{.container}}, colbacktitle=primaryBlue!70]
\textcolor{primaryBlue}{\textbf{Noi dung 1}} {\small\textcolor{gray!60}{$\leftarrow$ Con trực tiếp (Được chọn)}} \par\vspace{3pt}
\begin{tcolorbox}[arc=2pt, colback=white, colframe=primaryBlue!30, boxrule=1pt, title={\sffamily\scriptsize\color{darkText}.box}]
\textcolor{primaryBlue}{\textbf{Noi dung 2}} {\small\textcolor{gray!60}{$\leftarrow$ Cháu (Được chọn)}} \par\vspace{2pt}
\tcbox[on line, arc=2pt, colback=gray!10, colframe=gray!40, boxrule=1pt, boxsep=2pt]{%
\textcolor{primaryBlue}{\textbf{<span>}} \textcolor{primaryBlue}{\textbf{Noi dung 3}} {\small\textcolor{gray!60}{$\leftarrow$ Chắt (Được chọn)}}%
}
\end{tcolorbox}
\end{tcolorbox}
\end{browserWindow}

\begin{conceptBox}[Sơ Đồ Phân Cấp Cây DOM Trực Quan Của Ví Dụ 1]
\begin{center}
\begin{tikzpicture}[
  font=\sffamily\small,
  node distance=0.45cm and 0.6cm,
  root_box/.style={draw=primaryBlue, fill=primaryBlue!10, rounded corners=3pt, line width=1pt, inner sep=4pt, font=\sffamily\bfseries\color{primaryBlue}},
  selected_box/.style={draw=primaryBlue, fill=primaryBlue!15, rounded corners=3pt, line width=1pt, inner sep=4pt},
  normal_box/.style={draw=gray!40, fill=gray!5, rounded corners=3pt, line width=0.8pt, inner sep=4pt},
  badge_sel/.style={fill=primaryBlue, text=white, font=\sffamily\tiny\bfseries, rounded corners=2pt, inner sep=2pt},
  badge_mid/.style={fill=gray!50, text=white, font=\sffamily\tiny\bfseries, rounded corners=2pt, inner sep=2pt}
]
  % Root
  \node[root_box] (r) {.container \quad \scriptsize\textcolor{gray!70}{(Gốc DOM)}};

  % Level 1
  \node[selected_box, below right=0.5cm and 0.6cm of r.south west, anchor=north west] (n1) {\codeinline{<p>} Noi dung 1 \codeinline{</p>}};
  \node[badge_sel, right=4pt of n1] {Con trực tiếp};

  \node[normal_box, below=0.35cm of n1.south west, anchor=north west] (n2) {\codeinline{<div class="box">}};
  \node[badge_mid, right=4pt of n2] {Trung gian (Cấp 1)};

  % Level 2
  \node[selected_box, below right=0.5cm and 0.6cm of n2.south west, anchor=north west] (n3) {\codeinline{<p>} Noi dung 2 \codeinline{</p>}};
  \node[badge_sel, right=4pt of n3] {Cháu (Cấp 2)};

  \node[normal_box, below=0.35cm of n3.south west, anchor=north west] (n4) {\codeinline{<span>}};
  \node[badge_mid, right=4pt of n4] {Trung gian (Cấp 2)};

  % Level 3
  \node[selected_box, below right=0.5cm and 0.6cm of n4.south west, anchor=north west] (n5) {\codeinline{<p>} Noi dung 3 \codeinline{</p>}};
  \node[badge_sel, right=4pt of n5] {Chắt (Cấp 3)};

  % Connectors
  \draw[primaryBlue!70, line width=1pt] (r.south west) ++(0.3,0) |- (n1.west);
  \draw[primaryBlue!70, line width=1pt] (r.south west) ++(0.3,0) |- (n2.west);
  \draw[accentOrange!80, line width=1pt] (n2.south west) ++(0.3,0) |- (n3.west);
  \draw[accentOrange!80, line width=1pt] (n2.south west) ++(0.3,0) |- (n4.west);
  \draw[red!70, line width=1pt] (n4.south west) ++(0.3,0) |- (n5.west);

\end{tikzpicture}
\end{center}
\vspace{2pt}
\textbf{Phân tích}: \codeinline{.container p} quét từ gốc \codeinline{.container} và chọn tất cả thẻ \codeinline{<p>} bên trong. Hậu duệ bao gồm con trực tiếp (\codeinline{<p>Noi dung 1</p>}), cháu (\codeinline{<p>Noi dung 2</p>}) và chắt (\codeinline{<p>Noi dung 3</p>}).
\end{conceptBox}

\paragraph{c. Ví dụ 2: Phân tích chi tiết các cấp P1, P2, P3 trong cây DOM:}
\begin{codeWindow}[Ví Dụ 2: Thử Nghiệm P1, P2, P3 Với .parent p]
\begin{lstlisting}[language=HTML5]
<div class="parent">
    <p>Sample Text</p>
    <div class="child">
        <p>Sample Text</p>
        <span>
            <p>Sample Text</p>
        </span>
    </div>
</div>
\end{lstlisting}

\begin{lstlisting}[language=CSS3]
.parent p {
    color: blue;
    font-weight: bold;
}
\end{lstlisting}
\end{codeWindow}

\begin{browserWindow}[Kết Quả Hiển Thị Của Ví Dụ 2]
\sffamily\small
\textcolor{primaryBlue}{\textbf{P1 (trong .parent)}} $\rightarrow$ Màu xanh, in đậm \par\vspace{2pt}
\textcolor{primaryBlue}{\textbf{P2 (trong .child, nằm trong .parent)}} $\rightarrow$ Màu xanh, in đậm \par\vspace{2pt}
\textcolor{primaryBlue}{\textbf{P3 (trong <span>, nằm trong .child, nằm trong .parent)}} $\rightarrow$ Màu xanh, in đậm
\end{browserWindow}

\begin{conceptBox}[Sơ Đồ Phân Cấp Cây DOM Trực Quan Của Ví Dụ 2]
\begin{center}
\begin{tikzpicture}[
  font=\sffamily\small,
  node distance=0.45cm and 0.6cm,
  root_box/.style={draw=primaryBlue, fill=primaryBlue!10, rounded corners=3pt, line width=1pt, inner sep=4pt, font=\sffamily\bfseries\color{primaryBlue}},
  selected_box/.style={draw=primaryBlue, fill=primaryBlue!15, rounded corners=3pt, line width=1pt, inner sep=4pt},
  normal_box/.style={draw=gray!40, fill=gray!5, rounded corners=3pt, line width=0.8pt, inner sep=4pt},
  badge_sel/.style={fill=primaryBlue, text=white, font=\sffamily\tiny\bfseries, rounded corners=2pt, inner sep=2pt},
  badge_mid/.style={fill=gray!50, text=white, font=\sffamily\tiny\bfseries, rounded corners=2pt, inner sep=2pt}
]
  % Root
  \node[root_box] (r2) {.parent \quad \scriptsize\textcolor{gray!70}{(Gốc DOM)}};

  % Level 1
  \node[selected_box, below right=0.5cm and 0.6cm of r2.south west, anchor=north west] (p1) {\codeinline{<p>} P1 \codeinline{</p>}};
  \node[badge_sel, right=4pt of p1] {Con trực tiếp};

  \node[normal_box, below=0.35cm of p1.south west, anchor=north west] (c1) {\codeinline{<div class="child">}};
  \node[badge_mid, right=4pt of c1] {Trung gian (Cấp 1)};

  % Level 2
  \node[selected_box, below right=0.5cm and 0.6cm of c1.south west, anchor=north west] (p2) {\codeinline{<p>} P2 \codeinline{</p>}};
  \node[badge_sel, right=4pt of p2] {Cháu (Cấp 2)};

  \node[normal_box, below=0.35cm of p2.south west, anchor=north west] (sp) {\codeinline{<span>}};
  \node[badge_mid, right=4pt of sp] {Trung gian (Cấp 2)};

  % Level 3
  \node[selected_box, below right=0.5cm and 0.6cm of sp.south west, anchor=north west] (p3) {\codeinline{<p>} P3 \codeinline{</p>}};
  \node[badge_sel, right=4pt of p3] {Chắt (Cấp 3)};

  % Connectors
  \draw[primaryBlue!70, line width=1pt] (r2.south west) ++(0.3,0) |- (p1.west);
  \draw[primaryBlue!70, line width=1pt] (r2.south west) ++(0.3,0) |- (c1.west);
  \draw[accentOrange!80, line width=1pt] (c1.south west) ++(0.3,0) |- (p2.west);
  \draw[accentOrange!80, line width=1pt] (c1.south west) ++(0.3,0) |- (sp.west);
  \draw[red!70, line width=1pt] (sp.south west) ++(0.3,0) |- (p3.west);

\end{tikzpicture}
\end{center}
\vspace{2pt}
\textbf{Giải thích}: Không quan trọng \codeinline{<p>} nằm ở cấp độ lồng sâu bao nhiêu, chỉ cần nó là hậu duệ của \codeinline{.parent}, nó sẽ bị ảnh hưởng bởi CSS!
\end{conceptBox}

\paragraph{d. Ví dụ 3: So sánh chuyên sâu giữa Selector cụ thể (.parent .child) vs Selector toàn cục (p):}

\begin{codeWindow}[Mã HTML Thử Nghiệm So Sánh Mức Độ Tác Động]
\begin{lstlisting}[language=HTML5]
<p>Sample Text</p>
<div class="parent">
    <p class="child">Sample Text</p>
    <p class="child">Sample Text</p>
    <div class="child">
        <p>Sample Text</p>
        <p>Sample Text</p>
    </div>
    <span class="child">
        <p>Sample Text</p>
    </span>
</div>
\end{lstlisting}
\end{codeWindow}

\begin{codeWindow}[Trường hợp 1: Sử Dụng Selector Hậu Duệ Cụ Thể .parent .child]
\begin{lstlisting}[language=CSS3]
/* Select element */
.parent .child {
    color: red;
    font-weight: bold;
}
\end{lstlisting}
\end{codeWindow}

\begin{browserWindow}[Kết Quả Trường Hợp 1 (.parent .child)]
\sffamily\small
\textcolor{darkText}{Noi dung 0} {\small\textcolor{gray!60}{$\leftarrow$ Không ảnh hưởng (ngoài .parent)}} \par\vspace{4pt}
\begin{tcolorbox}[arc=2pt, colback=primaryBlue!5, colframe=primaryBlue!40, boxrule=1pt, title={\sffamily\small\color{white}\textbf{.parent}}, colbacktitle=primaryBlue!70]
\textcolor{red}{\textbf{Noi dung 1}} {\small\textcolor{gray!60}{$\leftarrow$ Đỏ, đậm (thẻ p có class .child)}} \par\vspace{2pt}
\textcolor{red}{\textbf{Noi dung 2}} {\small\textcolor{gray!60}{$\leftarrow$ Đỏ, đậm (thẻ p có class .child)}} \par\vspace{3pt}
\begin{tcolorbox}[arc=2pt, colback=red!5, colframe=red!40, boxrule=1pt, title={\sffamily\scriptsize\color{white}\textbf{div.child (Đổi màu toàn bộ nội dung trong div)}}, colbacktitle=red!70]
\textcolor{red}{\textbf{Noi dung 3}} {\small\textcolor{gray!60}{$\leftarrow$ Đỏ, đậm (kế thừa từ div.child)}} \par\vspace{2pt}
\textcolor{red}{\textbf{Noi dung 4}} {\small\textcolor{gray!60}{$\leftarrow$ Đỏ, đậm (kế thừa từ div.child)}}
\end{tcolorbox}\vspace{3pt}
\tcbox[on line, arc=2pt, colback=red!10, colframe=red!50, boxrule=1pt, boxsep=3pt]{%
\textcolor{red}{\textbf{span.child:}} \textcolor{red}{\textbf{Noi dung 5}} {\small\textcolor{gray!60}{$\leftarrow$ Đỏ, đậm (nằm trong span.child)}}%
}
\end{tcolorbox}
\end{browserWindow}

\begin{noteBox}[Mở Rộng Thẻ Con Cụ Thể Và Selector Toàn Bộ *]
\begin{itemize}[leftmargin=1.5em, itemsep=2pt, parsep=0pt]
    \item \codeinline{.parent .child h3 \{ color: red; \}} $\rightarrow$ Chỉ chọn thẻ \codeinline{<h3>} nằm trong phần tử \codeinline{.child} của \codeinline{.parent}.
    \item \codeinline{.parent .child * \{ color: red; \}} $\rightarrow$ Chọn \textbf{tất cả} các thẻ con mọi cấp bên trong \codeinline{.child}.
\end{itemize}
\end{noteBox}

\begin{codeWindow}[Trường hợp 2: Sử Dụng Selector Thẻ Toàn Cục p]
\begin{lstlisting}[language=CSS3]
/* Select element */
p {
    color: red;
    font-weight: bold;
}
\end{lstlisting}
\end{codeWindow}

\begin{prettyTableBox}[So Sánh Tiêu Chí Giữa Selector Hậu Duệ Cụ Thể vs Selector Toàn Cục]
\rowcolors{2}{white}{primaryBlue!4}
\sffamily\small
\begin{tabularx}{\linewidth}{|>{\raggedright\arraybackslash \bfseries\color{primaryBlue}}K{3.0cm}|Z|Z|}
\hline
\rowcolor{primaryBlue} \thd{Tiêu chí} & \thd{Trường hợp 1 (.parent .child)} & \thd{Trường hợp 2 (p)} \\
\hline
Phạm vi tác động & Chỉ các phần tử có class \codeinline{.child} nằm trong \codeinline{.parent}. & Tất cả thẻ \codeinline{<p>} trên toàn bộ trang web. \\
\hline
Tính linh hoạt & Cụ thể hơn, dễ quản lý theo cấu trúc HTML từng mô-đun. & Bao trùm toàn bộ, dễ gây tác động ngoài ý muốn. \\
\hline
Tính bảo trì & High (Cao), dễ dàng chỉnh sửa theo từng thành phần nhỏ. & Low (Thấp), vì mọi thay đổi đều ảnh hưởng tới toàn bộ thẻ \codeinline{<p>}. \\
\hline
Độ ưu tiên (Specificity) & Cao hơn (\textbf{20 điểm}: 2 class). & Thấp hơn (\textbf{1 điểm}: 1 element). \\
\hline
\end{tabularx}
\end{prettyTableBox}

\begin{tipBox}[Khi Nào Nên Dùng Selector Nào?]
\begin{itemize}[leftmargin=1.5em, itemsep=3pt, parsep=0pt]
    \item \textbf{Nên dùng \codeinline{.parent .child}}: Khi muốn quản lý phần tử theo cấu trúc cụ thể, định dạng chỉ nội dung trong 1 khu vực nhất định, tăng độ chính xác và tránh vô tình làm hỏng các phần tử ngoài ý muốn.
    \item \textbf{Nên dùng \codeinline{p}}: Khi muốn định dạng nhanh toàn bộ văn bản chung cho toàn trang hoặc khi viết CSS reset / style tổng quát ban đầu.
\end{itemize}
\end{tipBox}

\paragraph{e. So sánh Descendant Selector với các loại Combinators khác:}

\begin{prettyTableBox}[Bảng So Sánh Chi Tiết Các Loại Bộ Chọn Kết Hợp Combinators]
\rowcolors{2}{white}{primaryBlue!4}
\sffamily\small
\begin{tabularx}{\linewidth}{|>{\raggedright\arraybackslash \ttfamily\bfseries\color{accentOrange}}K{3.2cm}|Z|K{4.2cm}|}
\hline
\rowcolor{primaryBlue} \thd{Cú pháp Selector} & \thd{Cách hoạt động} & \thd{Ví dụ minh họa} \\
\hline
.cha con & Chọn \textbf{tất cả hậu duệ} của \codeinline{.cha} (mọi cấp). & \codeinline{.parent p} (Chọn P1, P2, P3) \\
\hline
.cha > con & Chỉ chọn \textbf{con trực tiếp} của \codeinline{.cha}. & \codeinline{.parent > p} (Chỉ chọn P1) \\
\hline
element1 + element2 & Chọn phần tử \codeinline{element2} \textbf{ngay liền kề} sau \codeinline{element1}. & \codeinline{div + p} (Chọn p đứng ngay sau div) \\
\hline
element1 \textasciitilde{} element2 & Chọn \textbf{tất cả anh chị em} \codeinline{element2} nằm phía sau \codeinline{element1}. & \codeinline{h2 \textasciitilde{} p} (Chọn tất cả p sau h2) \\
\hline
\end{tabularx}
\end{prettyTableBox}

\begin{codeWindow}[So Sánh Trực Quan: Descendant (.parent p) vs Child (.parent > p)]
\begin{lstlisting}[language=CSS3]
/* Select element */
.parent p {
    color: blue;
}

/* Sample Text */
.parent > p {
    color: red;
}
\end{lstlisting}
\end{codeWindow}

\begin{browserWindow}[Kết Quả So Sánh Descendant vs Child]
\sffamily\small
P1 (trong .parent) $\rightarrow$ \textcolor{red}{\textbf{Đỏ}} (với \codeinline{.parent > p} vì là con trực tiếp) \par\vspace{2pt}
P2, P3 $\rightarrow$ \textcolor{primaryBlue}{\textbf{Xanh}} (với \codeinline{.parent p} vì là hậu duệ)
\end{browserWindow}

\begin{codeWindow}[So Sánh Trực Quan: Descendant (.parent p) vs Adjacent Sibling (div + p)]
\begin{lstlisting}[language=CSS3]
/* Select element */
.parent p {
    color: blue;
}

/* Sample Text */
div + p {
    color: red;
}
\end{lstlisting}
\end{codeWindow}

\paragraph{f. 3 Ứng dụng thực tế phổ biến:}

\begin{tipBox}[3 Ứng Dụng Thực Tế Hàng Đầu Của Descendant Selector]
\begin{itemize}[leftmargin=1.5em, itemsep=3pt, parsep=0pt]
    \item \textbf{Ứng dụng 1}: Định dạng tất cả tiêu đề \codeinline{<h2>} nằm bên trong bài viết \codeinline{<article>}, bất kể cấp độ lồng nhau.
    \item \textbf{Ứng dụng 2}: Thiết kế liên kết menu điều hướng nhiều cấp \codeinline{nav ul li a}.
    \item \textbf{Ứng dụng 3}: Định dạng đồng bộ tất cả các ô nhập liệu \codeinline{input} nằm trong thẻ \codeinline{<label>} của \codeinline{<form>}.
\end{itemize}
\end{tipBox}

\begin{codeWindow}[Ứng Dụng 1: Định Dạng Tất Cả Tiêu Đề h2 Trong Bài Viết article]
\begin{lstlisting}[language=CSS3]
article h2 {
    color: #ff6347; /* Sample Text */
    font-size: 24px;
}
\end{lstlisting}
\end{codeWindow}

\begin{browserWindow}[Kết Quả: Tiêu Đề h2 Trong Article Đổi Màu Đỏ Cà Chua]
\sffamily\small
\begin{tcolorbox}[arc=3pt, colback=white, colframe=gray!30, boxrule=1pt, title={\sffamily\scriptsize\color{darkText}<article> (Khối bài viết)}]
{\color{red!80!black}\bfseries\Large Tiêu Đề Bài Viết h2} {\small\textcolor{gray!60}{$\leftarrow$ được chọn (màu \#ff6347, 24px)}} \par\vspace{4pt}
\textcolor{darkText}{Nội dung chi tiết của bài viết được trình bày tại đây...}
\end{tcolorbox}
\end{browserWindow}

\begin{codeWindow}[Ứng Dụng 2: Thiết Kế Menu Điều Hướng Nhiều Cấp (nav ul li a)]
\begin{lstlisting}[language=CSS3]
nav ul li a {
    text-decoration: none;
    color: #333;
    font-weight: bold;
}
\end{lstlisting}
\end{codeWindow}

\begin{browserWindow}[Kết Quả: Các Thẻ a Trong Menu Điều Hướng Phân Cấp]
\sffamily\small
\tcbox[on line, arc=3pt, colback=gray!15, colframe=gray!40, boxrule=1pt, boxsep=4pt, left=8pt, right=8pt]{\textcolor{darkText}{\bfseries Trang chủ}} \quad
\tcbox[on line, arc=3pt, colback=gray!15, colframe=gray!40, boxrule=1pt, boxsep=4pt, left=8pt, right=8pt]{\textcolor{darkText}{\bfseries Khóa học $\downarrow$}} \par\vspace{4pt}
\textcolor{gray!70}{\small\textit{($\rightarrow$ Chọn tất cả thẻ <a> trong menu, bất kể lồng sâu bao nhiêu cấp trong <ul> và <li>)}}
\end{browserWindow}

\begin{codeWindow}[Ứng Dụng 3: Định Dạng Đồng Bộ Tất Cả Ô Input Trong Form]
\begin{lstlisting}[language=CSS3]
form label input {
    border: 1px solid #ccc;
    padding: 8px;
    border-radius: 4px;
}
\end{lstlisting}
\end{codeWindow}

\begin{browserWindow}[Kết Quả: Ô Input Nằm Trong Thẻ Label Của Form]
\sffamily\small
\textbf{Họ và tên:} \tcbox[on line, arc=3pt, colback=white, colframe=gray!40, boxrule=1pt, boxsep=3pt, left=8pt, right=20pt]{\textcolor{gray!60}{Nhập họ tên...}} \par\vspace{4pt}
\textcolor{gray!70}{\small\textit{($\rightarrow$ Chọn tất cả thẻ <input> nằm bên trong <label> của <form>)}}
\end{browserWindow}

\paragraph{g. Lưu ý quan trọng \& Tóm tắt:}

\begin{warnBox}[2 Lưu Ý Cực Kỳ Quan Trọng Khi Sử Dụng Descendant Selector]
\begin{itemize}[leftmargin=1.5em, itemsep=3pt, parsep=0pt]
    \item \textbf{Hiệu suất (Performance)}: Nếu cây DOM quá lớn và lồng nhau quá sâu (ví dụ: \codeinline{body div main section article div p}), trình duyệt phải quét nhiều cấp để so khớp, có thể làm chậm tốc độ render trang web.
    \item \textbf{Độ ưu tiên (Specificity)}: Nếu có nhiều bộ chọn cùng tác động lên một phần tử, trình duyệt sẽ tính điểm Specificity để quyết định kiểu dáng nào thắng.
\end{itemize}
\end{warnBox}

\begin{tipBox}[Tóm Tắt Về Descendant Selector (Khoảng Trắng)]
\begin{itemize}[leftmargin=1.5em, itemsep=3pt, parsep=0pt]
    \item \textbf{Descendant Selector} rất mạnh mẽ và tiện lợi khi cần áp dụng style cho nhiều phần tử con và hậu duệ.
    \item Chỉ cần một \textbf{khoảng trắng}, bạn đã có thể chọn mọi phần tử nằm bên trong phần tử cha, bất kể cấp độ sâu.
    \item \textbf{Nhớ rằng}: Hãy kiểm soát độ sâu của bộ chọn để tránh tác động ngoài mong muốn đến các thành phần UI khác!
\end{itemize}
\end{tipBox}

\subsubsection{3.2. Child Selector (Bộ Chọn Con Trực Tiếp) -- Chi Tiết \& Cơ Chế Xử Lý}

\begin{conceptBox}[Khái Niệm Child Selector (bộ chọn con trực tiếp)]
\textbf{Child Selector (Bộ chọn con trực tiếp)} trong CSS là một bộ chọn dùng để chọn các \textbf{phần tử con trực tiếp} của phần tử cha, tức là các phần tử nằm ngay bên trong phần tử cha mà không đi qua trung gian. Nếu phần tử con nằm sâu hơn một cấp (là cháu, chắt\ldots), thì \textbf{không} bị ảnh hưởng bởi bộ chọn này.
\end{conceptBox}

\paragraph{a. Cú pháp và cách hoạt động:}
\begin{codeWindow}[Cú Pháp Child Selector (>)]
\begin{lstlisting}[language=CSS3]
A > B {
    /* CSS Commentp dCommentng cho B */
}
\end{lstlisting}
\end{codeWindow}

\begin{itemize}[leftmargin=1.5em, itemsep=2pt, parsep=0pt]
    \item \codeinline{A}: Phần tử cha.
    \item \codeinline{B}: Phần tử con trực tiếp của A (ngay bên trong A, không lồng thêm cấp nào).
\end{itemize}

\begin{noteBox}[Cơ Chế Hoạt Động Của Trình Duyệt Khi Xử Lý A > B]
Khi bạn viết \codeinline{A > B} trong CSS, trình duyệt sẽ thực hiện theo 2 bước:
\begin{enumerate}[leftmargin=1.5em, itemsep=3pt, parsep=0pt]
    \item \textbf{Bước 1}: Tìm phần tử \codeinline{A} (phần tử cha).
    \item \textbf{Bước 2}: Kiểm tra các phần tử bên trong \codeinline{A}:
    \begin{itemize}[leftmargin=1.2em, itemsep=2pt, parsep=0pt]
        \item Nếu phần tử \codeinline{B} là \textbf{con trực tiếp} của \codeinline{A} $\rightarrow$ \textcolor{tipBorder}{\textbf{Áp dụng CSS}}.
        \item Nếu phần tử \codeinline{B} \textbf{không} phải con trực tiếp (là cháu, chắt) $\rightarrow$ \textcolor{warnBorder}{\textbf{Bỏ qua, không áp dụng CSS}}.
    \end{itemize}
\end{enumerate}
\textit{Ghi chú: Con trực tiếp là phần tử nằm ngay bên trong phần tử cha, không qua trung gian.}
\end{noteBox}

\paragraph{b. Trình tự xử lý của trình duyệt (Browser Pipeline):}

\begin{conceptBox}[3 Giai Đoạn Trình Duyệt Xử Lý CSS]
Trình duyệt xử lý CSS theo thứ tự duyệt cây DOM:
\begin{enumerate}[leftmargin=1.5em, itemsep=3pt, parsep=0pt]
    \item \textbf{1. Phân tích cấu trúc HTML} $\rightarrow$ Xây dựng \term{DOM Tree} (Cây DOM).
    \item \textbf{2. Áp dụng quy tắc CSS}:
    \begin{itemize}[leftmargin=1.2em, itemsep=2pt, parsep=0pt]
        \item Kiểm tra bộ chọn (\codeinline{A > B}).
        \item So khớp phần tử: Nếu \codeinline{B} là con trực tiếp của \codeinline{A}, áp dụng CSS.
        \item Bỏ qua phần tử không khớp quan hệ trực tiếp.
    \end{itemize}
    \item \textbf{3. Kết xuất giao diện (Render)} $\rightarrow$ Hiển thị kết quả ra màn hình.
\end{enumerate}
\end{conceptBox}

\paragraph{c. Ví dụ minh họa cơ bản (Con vs Cháu):}
\begin{codeWindow}[Ví Dụ 1: Phân Biệt Con Trực Tiếp Và Cháu]
\begin{lstlisting}[language=HTML5]
<div class="parent">
    <p>Sample Text</p>
    <div>
        <p>Sample Text</p>
    </div>
</div>
\end{lstlisting}

\begin{lstlisting}[language=CSS3]
.parent > p {
    color: red;
    font-weight: bold;
}
\end{lstlisting}
\end{codeWindow}

\begin{browserWindow}[Kết Quả Hiển Thị Ví Dụ 1 (.parent > p)]
\sffamily\small
\begin{tcolorbox}[arc=2pt, colback=primaryBlue!5, colframe=primaryBlue!40, boxrule=1pt, title={\sffamily\small\color{white}\textbf{.parent}}, colbacktitle=primaryBlue!70]
\textcolor{red}{\textbf{Con trực tiếp}} {\small\textcolor{gray!60}{$\leftarrow$ được chọn (màu đỏ, in đậm)}} \par\vspace{3pt}
\begin{tcolorbox}[arc=2pt, colback=white, colframe=gray!30, boxrule=1pt, title={\sffamily\scriptsize\color{darkText}div lồng bên trong}]
\textcolor{darkText}{Cháu (không phải con trực tiếp)} {\small\textcolor{gray!60}{$\leftarrow$ không đổi màu}}
\end{tcolorbox}
\end{tcolorbox}
\end{browserWindow}

\begin{noteBox}[Phân Tích Chi Tiết Quá Trình So Khớp Ví Dụ 1]
Trình duyệt đọc CSS và tìm phần tử \codeinline{.parent}, sau đó kiểm tra các thẻ \codeinline{<p>}:
\begin{itemize}[leftmargin=1.5em, itemsep=2pt, parsep=0pt]
    \item \textcolor{tipBorder}{\textbf{Thẻ <p> đầu tiên}}: Thỏa mãn (là con trực tiếp) $\rightarrow$ Áp dụng CSS (màu đỏ, in đậm).
    \item \textcolor{warnBorder}{\textbf{Thẻ <p> thứ hai}}: Không thỏa mãn (là cháu, nằm trong thẻ \codeinline{<div>}) $\rightarrow$ Bỏ qua, không áp dụng CSS.
\end{itemize}
\end{noteBox}

\paragraph{d. Ví dụ nâng cao (Phân tích lồng ghép nhiều cấp):}

\begin{codeWindow}[Ví Dụ 2: Phân Tích Cấu Trúc Lồng Nhiều Cấp VỚI div > p]
\begin{lstlisting}[language=HTML5]
<p>Sample Text</p>
<div>
    <p>Sample Text</p>
    <p>Sample Text</p>
    <div>
        <p>Sample Text</p>
        <p>Sample Text</p>
    </div>
    <span>
        <p>Sample Text</p>
    </span>
</div>
\end{lstlisting}

\begin{lstlisting}[language=CSS3]
div > p {
    color: red;
    font-weight: bold;
}
\end{lstlisting}
\end{codeWindow}

\begin{noteBox}[Phân Tích Từng Phần Tử Trong Ví Dụ 2]
\begin{enumerate}[leftmargin=1.5em, itemsep=3pt, parsep=0pt]
    \item \codeinline{<p>Noi dung 0</p>}: \textbf{Không được chọn} vì đứng độc lập, không nằm trong \codeinline{<div>}.
    \item \codeinline{<p>Noi dung 1</p>} \& \codeinline{<p>Noi dung 2</p>}: \textcolor{red}{\textbf{Được chọn}} vì là con trực tiếp của thẻ \codeinline{<div>} ngoài.
    \item \codeinline{<p>Noi dung 3</p>} \& \codeinline{<p>Noi dung 4</p>}: \textcolor{red}{\textbf{Được chọn}} vì là con trực tiếp của thẻ \codeinline{<div>} trong.
    \item \codeinline{<p>Noi dung 5</p>}: \textbf{Không được chọn} vì nó là con trực tiếp của \codeinline{<span>}, không phải con trực tiếp của \codeinline{<div>}.
\end{enumerate}
\end{noteBox}

\begin{browserWindow}[Kết Quả Hiển Thị Chi Tiết Ví Dụ 2 (div > p)]
\sffamily\small
\textcolor{darkText}{Noi dung 0} {\small\textcolor{gray!60}{$\leftarrow$ Bình thường (ngoài div)}} \par\vspace{4pt}
\begin{tcolorbox}[arc=2pt, colback=primaryBlue!5, colframe=primaryBlue!40, boxrule=1pt, title={\sffamily\small\color{white}\textbf{div ngoài}}, colbacktitle=primaryBlue!70]
\textcolor{red}{\textbf{Noi dung 1}} {\small\textcolor{gray!60}{$\leftarrow$ màu đỏ, in đậm}} \par\vspace{2pt}
\textcolor{red}{\textbf{Noi dung 2}} {\small\textcolor{gray!60}{$\leftarrow$ màu đỏ, in đậm}} \par\vspace{3pt}
\begin{tcolorbox}[arc=2pt, colback=white, colframe=accentOrange!50, boxrule=1pt, title={\sffamily\scriptsize\color{white}\textbf{div trong}}, colbacktitle=accentOrange!80]
\textcolor{red}{\textbf{Noi dung 3}} {\small\textcolor{gray!60}{$\leftarrow$ màu đỏ, in đậm (con trực tiếp của div trong)}} \par\vspace{2pt}
\textcolor{red}{\textbf{Noi dung 4}} {\small\textcolor{gray!60}{$\leftarrow$ màu đỏ, in đậm (con trực tiếp của div trong)}}
\end{tcolorbox}\vspace{3pt}
\tcbox[on line, arc=2pt, colback=gray!10, colframe=gray!40, boxrule=1pt, boxsep=3pt]{%
\textcolor{darkText}{\textbf{<span>}} \textcolor{darkText}{Noi dung 5} {\small\textcolor{gray!60}{$\leftarrow$ Bình thường (con của span)}}%
}
\end{tcolorbox}
\end{browserWindow}

\paragraph{e. Ví dụ chuỗi Selector 3 cấp (.container > .box > p):}

\begin{codeWindow}[Ví Dụ 3: Chuỗi Selector 3 Cấp Độ Lồng Nhau]
\begin{lstlisting}[language=HTML5]
<div class="container">
    <div class="box">
        <p>Sample Text</p>
        <div class="inner-box">
            <p>Sample Text</p>
        </div>
        <p>Sample Text</p>
    </div>
</div>
\end{lstlisting}

\begin{lstlisting}[language=CSS3]
.container > .box > p {
    color: blue;
    font-weight: bold;
}
\end{lstlisting}
\end{codeWindow}

\begin{noteBox}[Phân Tích Chuỗi Selector .container > .box > p]
\begin{enumerate}[leftmargin=1.5em, itemsep=2pt, parsep=0pt]
    \item Tìm \codeinline{.container}.
    \item Tìm \codeinline{.box} là con trực tiếp của \codeinline{.container}.
    \item Tìm các thẻ \codeinline{<p>} là con trực tiếp của \codeinline{.box}.
    \item Áp dụng màu xanh và in đậm cho \codeinline{<p>} con trực tiếp của \codeinline{.box} (bỏ qua \codeinline{<p>} nằm trong \codeinline{.inner-box}).
\end{enumerate}
\end{noteBox}

\begin{browserWindow}[Kết Quả Hiển Thị Chuỗi Selector 3 Cấp]
\sffamily\small
\begin{tcolorbox}[arc=2pt, colback=primaryBlue!5, colframe=primaryBlue!40, boxrule=1pt, title={\sffamily\small\color{white}\textbf{.container}}, colbacktitle=primaryBlue!70]
\begin{tcolorbox}[arc=2pt, colback=white, colframe=primaryBlue!70, boxrule=1pt, title={\sffamily\scriptsize\color{white}\textbf{.box}}, colbacktitle=primaryBlue!90]
\textcolor{primaryBlue}{\textbf{Con trực tiếp 1}} {\small\textcolor{gray!60}{$\leftarrow$ màu xanh, in đậm}} \par\vspace{2pt}
\begin{tcolorbox}[arc=2pt, colback=gray!10, colframe=gray!40, boxrule=1pt, title={\sffamily\tiny\color{darkText}.inner-box}]
\textcolor{darkText}{Cháu 1} {\small\textcolor{gray!60}{$\leftarrow$ không bị ảnh hưởng}}
\end{tcolorbox}\vspace{2pt}
\textcolor{primaryBlue}{\textbf{Con trực tiếp 2}} {\small\textcolor{gray!60}{$\leftarrow$ màu xanh, in đậm}}
\end{tcolorbox}
\end{tcolorbox}
\end{browserWindow}

\paragraph{f. So sánh với Descendant Selector (Bộ chọn hậu duệ):}

\begin{prettyTableBox}[So Sánh Child Selector (>) vs Descendant Selector (Khoảng Trắng)]
\rowcolors{2}{white}{primaryBlue!4}
\sffamily\small
\begin{tabularx}{\linewidth}{|>{\raggedright\arraybackslash \ttfamily\bfseries\color{accentOrange}}K{2.0cm}|K{4.0cm}|Z|}
\hline
\rowcolor{primaryBlue} \thd{Bộ chọn} & \thd{Ý nghĩa} & \thd{Phạm vi chọn} \\
\hline
A > B & Chọn con trực tiếp của A & Chỉ phần tử con ngay bên trong A \\
\hline
A B & Chọn tất cả hậu duệ của A & Chọn tất cả phần tử bên trong A (bao gồm con, cháu, chắt\ldots) \\
\hline
\end{tabularx}
\end{prettyTableBox}

\begin{codeWindow}[So Sánh Trực Quan: .parent p vs .parent > p]
\begin{lstlisting}[language=CSS3]
/* Select element */
.parent p { 
    color: blue; 
}

/* Sample Text */
.parent > p { 
    color: red; 
}
\end{lstlisting}
\end{codeWindow}

\begin{browserWindow}[Kết Quả So Sánh: Con Trực Tiếp vs Cháu]
\sffamily\small
\textcolor{red}{\textbf{Con trực tiếp}} {\small\textcolor{gray!60}{$\leftarrow$ Màu đỏ (với .parent > p)}} \par\vspace{2pt}
\textcolor{primaryBlue}{\textbf{Cháu}} {\small\textcolor{gray!60}{$\leftarrow$ Màu xanh (với .parent p)}}
\end{browserWindow}

\paragraph{g. 3 Quy tắc quan trọng cốt lõi:}

\begin{noteBox}[3 Quy Tắc Quan Trọng Khi Sử Dụng Child Selector]
\begin{enumerate}[leftmargin=1.5em, itemsep=3pt, parsep=0pt]
    \item \textbf{Chọn đúng con trực tiếp}: Chỉ áp dụng CSS cho phần tử nằm ngay bên trong phần tử cha.
    \item \textbf{Không ảnh hưởng đến hậu duệ sâu hơn}: Nếu phần tử là cháu, chắt (không trực tiếp), trình duyệt tự động bỏ qua.
    \item \textbf{Tính thứ tự trong DOM}: Chỉ xét theo cấu trúc HTML, không quan tâm vị trí hiển thị trực quan trên giao diện màn hình.
\end{enumerate}
\end{noteBox}

\paragraph{h. Trường hợp không hoạt động:}

\begin{warnBox}[Trường Hợp Child Selector KHÔNG Hoạt Động]
\textbf{Không có quan hệ cha-con trực tiếp}: Nếu phần tử không nằm ngay bên trong phần tử cha (bị bọc bởi một thẻ khác), bộ chọn \codeinline{>} sẽ \textbf{không} áp dụng.
\end{warnBox}

\begin{codeWindow}[Ví Dụ Sai: Thẻ p Nằm Trong article Nên Không Phải Con Trực Tiếp Của section]
\begin{lstlisting}[language=HTML5]
<section>
    <article>
        <p>Sample Text</p>
    </article>
</section>
\end{lstlisting}

\begin{lstlisting}[language=CSS3]
section > p {
    color: green;
}
\end{lstlisting}
\end{codeWindow}

\begin{browserWindow}[Kết Quả: Không Có Thay Đổi]
\sffamily\small
\textcolor{darkText}{Không phải con trực tiếp} {\small\textcolor{gray!60}{$\leftarrow$ Không có gì thay đổi vì thẻ \codeinline{<p>} nằm trong \codeinline{<article>}, không phải con trực tiếp của \codeinline{<section>}.}}
\end{browserWindow}

\paragraph{i. Các ứng dụng thực tế phong phú:}

\begin{tipBox}[5 Ứng Dụng Thực Tế Mạnh Mẽ Của Child Selector]
\begin{itemize}[leftmargin=1.5em, itemsep=3pt, parsep=0pt]
    \item \textbf{1. Cấu trúc menu điều hướng (\codeinline{.nav > li > a})}: Chỉ chọn thẻ \codeinline{<a>} là con trực tiếp của \codeinline{<li>}, tránh áp dụng CSS lên các thẻ con lồng sâu hơn trong menu dropdown.
    \item \textbf{2. Layout card sản phẩm (\codeinline{.card > .title})}: Chỉ chọn tiêu đề là con trực tiếp của thẻ card, giữ cấu trúc gọn gàng.
    \item \textbf{3. Bố cục lưới Grid / Flexbox (\codeinline{.grid > .item})}: Chọn phần tử item trực tiếp trong grid, giúp dễ dàng căn chỉnh và quản lý khoảng cách.
    \item \textbf{4. Ẩn/hiện phần tử footer con trực tiếp (\codeinline{.card > .footer})}: Ẩn phần footer nếu nó là con trực tiếp của card.
    \item \textbf{5. Định dạng cấu trúc bảng (\codeinline{table > tbody > tr})}: Thêm đường viền dưới mỗi hàng của bảng, chỉ khi là trực tiếp trong \codeinline{<tbody>}.
\end{itemize}
\end{tipBox}

\begin{codeWindow}[Ứng Dụng 1: Cấu Trúc Menu Điều Hướng (.nav > li > a)]
\begin{lstlisting}[language=CSS3]
.nav > li > a {
    text-decoration: none;
    padding: 10px;
    color: #007bff;
}
\end{lstlisting}
\end{codeWindow}

\begin{codeWindow}[Ứng Dụng 2: Card Sản Phẩm (.card > .title)]
\begin{lstlisting}[language=CSS3]
.card > .title {
    font-size: 20px;
    font-weight: bold;
    color: #333;
}
\end{lstlisting}
\end{codeWindow}

\begin{codeWindow}[Ứng Dụng 3: Bố Cục Lưới Flexbox/Grid (.grid > .item)]
\begin{lstlisting}[language=CSS3]
.grid > .item {
    flex: 1;
    margin: 10px;
}
\end{lstlisting}
\end{codeWindow}

\begin{tipBox}[Tóm Tắt Về Child Selector (>)]
\begin{itemize}[leftmargin=1.5em, itemsep=3pt, parsep=0pt]
    \item \textbf{Child Selector (\codeinline{>})} giúp chọn chính xác phần tử con trực tiếp của một phần tử cha mà không ảnh hưởng đến phần tử lồng sâu hơn.
    \item \textbf{Hoạt động theo cấu trúc DOM}: Phân tích, so khớp, và áp dụng CSS theo thứ tự cây DOM.
    \item \textbf{Ứng dụng rộng rãi} trong thiết kế layout, menu điều hướng, card sản phẩm, và cấu trúc trang web phức tạp.
\end{itemize}
\end{tipBox}

\subsubsection{3.3. Adjacent Sibling Selector (Bộ Chọn Anh Chị Liền Kề) -- Chi Tiết}

\begin{conceptBox}[Khái Niệm Adjacent Sibling Selector]
\textbf{Adjacent Sibling Selector (Bộ chọn anh chị liền kề)} là một loại bộ chọn trong CSS dùng để chọn phần tử ngay liền kề sau một phần tử khác trong cây DOM, cùng cấp cha. Bộ chọn này rất hữu ích khi bạn muốn tác động chỉ đến phần tử đứng ngay sau một phần tử cụ thể, thay vì chọn tất cả các phần tử sau đó.
\end{conceptBox}

\paragraph{a. Cú pháp:}
\begin{codeWindow}[Cú Pháp Adjacent Sibling Selector (+)]
\begin{lstlisting}[language=CSS3]
element1 + element2 {
    /* CSS styles */
}
\end{lstlisting}
\end{codeWindow}

\begin{itemize}[leftmargin=1.5em, itemsep=2pt, parsep=0pt]
    \item \codeinline{element1}: Phần tử đứng trước.
    \item \codeinline{element2}: Phần tử liền kề ngay sau phần tử 1.
\end{itemize}

\begin{noteBox}[Lưu Ý Quan Trọng]
Hai phần tử phải \textbf{cùng cấp} (sibling elements) và \textbf{liền kề nhau} trong cây DOM. Nếu có bất kỳ phần tử nào chen giữa, bộ chọn sẽ không hoạt động.
\end{noteBox}

\paragraph{b. Ví dụ cơ bản:}
\begin{codeWindow}[Ví Dụ Cơ Bản Adjacent Sibling Selector]
\begin{lstlisting}[language=HTML5]
<div class="box">Box</div>
<div class="text">Sample Text</div>
<div class="text">Sample Text</div>
\end{lstlisting}

\begin{lstlisting}[language=CSS3]
.box + .text {
    color: red;
}
\end{lstlisting}
\end{codeWindow}

\begin{browserWindow}[Kết Quả Hiển Thị Ví Dụ Cơ Bản (.box + .text)]
\sffamily\small
\tcbox[on line, arc=3pt, colback=codeBg, colframe=darkText, boxrule=1pt, boxsep=2pt]{\textbf{Box}} \par\vspace{4pt}
\textcolor{red}{\textbf{Văn bản 1}} {\small\textcolor{gray!60}{$\leftarrow$ được chọn (liền kề ngay sau .box)}} \par\vspace{2pt}
\textcolor{darkText}{Văn bản 2} {\small\textcolor{gray!60}{$\leftarrow$ không bị ảnh hưởng (không nằm ngay sát .box)}} \par\vspace{4pt}
\end{browserWindow}

\paragraph{c. Nguyên tắc hoạt động:}

\begin{noteBox}[3 Nguyên Tắc Hoạt Động Của Adjacent Sibling Selector]
\begin{enumerate}[leftmargin=1.5em, itemsep=3pt, parsep=0pt]
    \item \textbf{Cùng cấp (siblings)}: Hai phần tử phải cùng nằm trong một phần tử cha.
    \item \textbf{Liền kề nhau (adjacent)}: Phần tử thứ hai phải ngay sau phần tử thứ nhất (không có phần tử nào chen giữa).
    \item \textbf{Một phần tử duy nhất}: CSS sẽ chỉ chọn phần tử liền kề đầu tiên, không áp dụng cho các phần tử sau đó.
\end{enumerate}
\end{noteBox}

\paragraph{d. Phân tích Element 1 và Element 2:}
Với Adjacent Sibling Selector (\codeinline{+}), chỉ phần tử thứ hai (\codeinline{element2}) --- nếu nằm ngay sau \codeinline{element1} --- mới bị ảnh hưởng bởi CSS. \textbf{Element 1 hoàn toàn không bị tác động.}

\begin{codeWindow}[Phân Tích Chi Tiết: Element 1 vs Element 2]
\begin{lstlisting}[language=HTML5]
<div class="box">Box</div>
<div class="text">Sample Text</div>
<div class="text">Sample Text</div>
\end{lstlisting}

\begin{lstlisting}[language=CSS3]
.box + .text {
    color: red;
}
\end{lstlisting}
\end{codeWindow}

\begin{browserWindow}[Phân Tích Kết Quả: Phần Tử Nào Bị Ảnh Hưởng?]
\sffamily\small
\tcbox[on line, arc=3pt, colback=codeBg, colframe=darkText, boxrule=1pt, boxsep=2pt]{\textbf{Box}} {\small\textcolor{gray!60}{$\leftarrow$ .box (element1) $\rightarrow$ \textbf{không bị ảnh hưởng}}} \par\vspace{4pt}
\textcolor{red}{\textbf{Văn bản 1}} {\small\textcolor{gray!60}{$\leftarrow$ .text ngay sau .box $\rightarrow$ \textbf{đỏ}}} \par\vspace{2pt}
\textcolor{darkText}{Văn bản 2} {\small\textcolor{gray!60}{$\leftarrow$ không liền kề trực tiếp .box $\rightarrow$ \textbf{không thay đổi}}} \par\vspace{4pt}
\end{browserWindow}

\paragraph{e. Khi có phần tử chen giữa:}
\begin{codeWindow}[Trường Hợp Phần Tử Chen Giữa -- Adjacent Sibling Thất Bại]
\begin{lstlisting}[language=HTML5]
<div class="box">Box</div>
<span>Sample Text</span>
<div class="text">Sample Text</div>
\end{lstlisting}

\begin{lstlisting}[language=CSS3]
.box + .text {
    color: red;
}
\end{lstlisting}
\end{codeWindow}

\begin{browserWindow}[Kết Quả Khi Có Phần Tử Chen Giữa]
\sffamily\small
\tcbox[on line, arc=3pt, colback=codeBg, colframe=darkText, boxrule=1pt, boxsep=2pt]{\textbf{Box}} \par\vspace{4pt}
\textcolor{darkText}{\textit{Chen giữa}} \par\vspace{2pt}
\textcolor{darkText}{Văn bản 1} {\small\textcolor{gray!60}{$\leftarrow$ không đổi màu (bị \codeinline{<span>} chen giữa)}} \par\vspace{4pt}
\end{browserWindow}

\begin{warnBox}[Khi Nào Adjacent Sibling KHÔNG Hoạt Động?]
Nếu có \textbf{bất kỳ phần tử nào chen giữa} \codeinline{element1} và \codeinline{element2}, bộ chọn sẽ không áp dụng. Nếu bạn muốn chọn tất cả các phần tử phía sau mà không cần liền kề, hãy dùng \textbf{General Sibling Selector (\textasciitilde{})}!
\end{warnBox}

\paragraph{f. Khái niệm ``Cùng cấp'' (Sibling Elements):}

\begin{conceptBox}[Hiểu Rõ ``Cùng Cấp'' Trong Cây DOM]
Trong HTML, khi nói \textbf{cùng cấp}, ta đang nói về các \textbf{sibling elements} (phần tử anh chị em). Chúng là những phần tử nằm cùng một cấp bậc trong cây DOM và có cùng phần tử cha. Hãy hình dung cấu trúc như một cây gia đình:
\begin{itemize}[leftmargin=1.5em, itemsep=2pt, parsep=0pt]
    \item \textbf{Phần tử cha (parent element)}: Giống như bố mẹ.
    \item \textbf{Phần tử con (child element)}: Là con cái.
    \item \textbf{Phần tử cùng cấp (sibling element)}: Là anh chị em ruột --- các phần tử con của cùng một phần tử cha.
\end{itemize}
\textit{Cùng nằm trên 1 đường thẳng sẽ là cùng cấp với nhau!}
\end{conceptBox}

\begin{codeWindow}[Ví Dụ Minh Họa Khái Niệm Sibling Elements]
\begin{lstlisting}[language=HTML5]
<div class="parent">
    <div class="child-1">Child 1</div>
    <div class="child-2">Child 2</div>
    <div class="child-3">Child 3</div>
</div>
\end{lstlisting}

\begin{lstlisting}[language=CSS3]
.child-1 + .child-2 {
    color: red;
}
\end{lstlisting}
\end{codeWindow}

\begin{browserWindow}[Kết Quả Minh Họa Sibling Elements]
\sffamily\small
\begin{tcolorbox}[arc=2pt, colback=primaryBlue!5, colframe=primaryBlue!40, boxrule=1pt, title={\sffamily\small\color{white}\textbf{.parent (Phần tử cha)}}, colbacktitle=primaryBlue!70]
\textcolor{darkText}{Child 1} {\small\textcolor{gray!60}{$\leftarrow$ .child-1 (element1, không bị ảnh hưởng)}} \par\vspace{3pt}
\textcolor{red}{\textbf{Child 2}} {\small\textcolor{gray!60}{$\leftarrow$ .child-2 (liền kề, đổi đỏ)}} \par\vspace{3pt}
\textcolor{darkText}{Child 3} {\small\textcolor{gray!60}{$\leftarrow$ .child-3 (không liền kề .child-1, không đổi)}}
\end{tcolorbox}
\textcolor{gray!70}{\small\textit{Child 1, Child 2, Child 3 là các phần tử cùng cấp (siblings) vì chúng đều là con trực tiếp của .parent.}}
\end{browserWindow}

\paragraph{g. Ví dụ chi tiết hơn:}
\begin{codeWindow}[Ví Dụ Chi Tiết: h1 + p]
\begin{lstlisting}[language=HTML5]
<h1>Header Title</h1>
<p>Sample Text</p>
<p>Sample Text</p>
\end{lstlisting}

\begin{lstlisting}[language=CSS3]
h1 + p {
    font-weight: bold;
    color: blue;
}
\end{lstlisting}
\end{codeWindow}

\begin{browserWindow}[Kết Quả Hiển Thị h1 + p]
\sffamily\small
{\bfseries\Large\textcolor{darkText}{Tiêu đề}} \par\vspace{4pt}
\textcolor{primaryBlue}{\textbf{Đoạn văn đầu tiên}} {\small\textcolor{gray!60}{$\leftarrow$ in đậm + xanh dương}} \par\vspace{2pt}
\textcolor{darkText}{Đoạn văn thứ hai} {\small\textcolor{gray!60}{$\leftarrow$ không bị ảnh hưởng}}
\end{browserWindow}

\paragraph{h. Trường hợp không hoạt động:}
\begin{codeWindow}[Trường Hợp Không Hoạt Động: Phần Tử Chen Giữa h1 và p]
\begin{lstlisting}[language=HTML5]
<h1>Header Title</h1>
<span>Sample Text</span>
<p>Sample Text</p>
\end{lstlisting}

\begin{lstlisting}[language=CSS3]
h1 + p {
    color: green;
}
\end{lstlisting}
\end{codeWindow}

\begin{browserWindow}[Kết Quả: Không Có Phần Tử Nào Đổi Màu]
\sffamily\small
{\bfseries\Large\textcolor{darkText}{Tiêu đề}} \par\vspace{4pt}
\textcolor{darkText}{\textit{Không tính là liền kề}} \par\vspace{2pt}
\textcolor{darkText}{Đoạn văn} {\small\textcolor{gray!60}{$\leftarrow$ không đổi màu vì \codeinline{<span>} chen giữa \codeinline{<h1>} và \codeinline{<p>}}}
\end{browserWindow}

\begin{tipBox}[Ứng Dụng Thực Tế Của Adjacent Sibling Selector]
\begin{itemize}[leftmargin=1.5em, itemsep=3pt, parsep=0pt]
    \item \textbf{Ứng dụng 1}: Tạo khoảng cách hoặc định dạng đặc biệt cho phần tử ngay sau tiêu đề.
    \item \textbf{Ứng dụng 2}: Định dạng menu hoặc nút bấm liền kề nhau.
\end{itemize}
\end{tipBox}

\begin{codeWindow}[Ứng Dụng 1: Định Dạng Đoạn Văn Ngay Sau Tiêu Đề]
\begin{lstlisting}[language=CSS3]
h2 + p {
    margin-top: 20px;
    font-style: italic;
}
\end{lstlisting}
\end{codeWindow}

\begin{browserWindow}[Kết Quả: Đoạn Văn Đầu Tiên Sau h2 Được In Nghiêng]
\sffamily\small
{\bfseries\large\textcolor{darkText}{Tiêu đề h2}} \par\vspace{4pt}
\textcolor{darkText}{\textit{Đoạn văn đầu tiên --- in nghiêng, cách trên 20px}} {\small\textcolor{gray!60}{$\leftarrow$ được chọn}} \par\vspace{2pt}
\textcolor{darkText}{Đoạn văn thứ hai --- bình thường} {\small\textcolor{gray!60}{$\leftarrow$ không bị ảnh hưởng}}
\end{browserWindow}

\begin{codeWindow}[Ứng Dụng 2: Định Dạng Nút Bấm Liền Kề]
\begin{lstlisting}[language=CSS3]
button.primary + button.secondary {
    margin-left: 10px;
    background-color: gray;
}
\end{lstlisting}
\end{codeWindow}

\begin{browserWindow}[Kết Quả: Nút Primary và Secondary Liền Kề]
\sffamily\small
\tcbox[on line, arc=5pt, colback=primaryBlue, colframe=primaryBlue, boxsep=4pt, left=10pt, right=10pt]{\color{white}\bfseries Primary} \hspace{10pt}
\tcbox[on line, arc=5pt, colback=gray!60, colframe=gray!60, boxsep=4pt, left=10pt, right=10pt]{\color{white}\bfseries Secondary} \par\vspace{4pt}
\textcolor{gray!70}{\small\textit{($\rightarrow$ Nút Secondary (liền kề sau Primary) có nền xám và cách trái 10px.)}}
\end{browserWindow}

\paragraph{j. Tóm tắt nhanh \& Kết luận:}

\begin{prettyTableBox}[Tóm Tắt Nhanh Về Adjacent Sibling Selector]
\rowcolors{2}{white}{primaryBlue!4}
\sffamily\small
\begin{tabularx}{\linewidth}{|>{\raggedright\arraybackslash \bfseries\color{primaryBlue}}K{3.5cm}|Z|}
\hline
\rowcolor{primaryBlue} \thd{Đặc điểm} & \thd{Chi tiết} \\
\hline
Cú pháp & \codeinline{element1 + element2 \{ \}} \\
\hline
Ký hiệu & Dấu cộng \codeinline{+} giữa 2 phần tử \\
\hline
Phạm vi chọn & Chỉ chọn phần tử liền kề ngay sau (1 phần tử duy nhất) \\
\hline
Yêu cầu & Hai phần tử phải cùng cấp và liền kề, không có phần tử chen giữa \\
\hline
Element 1 & Không bị ảnh hưởng bởi CSS \\
\hline
Element 2 & Bị ảnh hưởng bởi CSS (nếu liền kề) \\
\hline
So sánh với \textasciitilde{} & \codeinline{+} chọn 1 phần tử liền kề, \codeinline{\textasciitilde{}} chọn tất cả phía sau \\
\hline
\end{tabularx}
\end{prettyTableBox}

\begin{tipBox}[Kết Luận Về Adjacent Sibling Selector]
\begin{itemize}[leftmargin=1.5em, itemsep=3pt, parsep=0pt]
    \item \textbf{Adjacent sibling selector (+)} cực kỳ mạnh khi bạn muốn kiểm soát phần tử ngay sau một phần tử khác.
    \item Tuy nhiên, nếu bạn muốn chọn \textbf{tất cả phần tử phía sau}, hãy dùng \textbf{General Sibling Selector (\textasciitilde{})}!
\end{itemize}
\end{tipBox}


\subsubsection{3.4. General Sibling Selector (Bộ Chọn Anh Chị Tổng Quát) -- Chi Tiết}

\begin{conceptBox}[Khái Niệm General Sibling Selector]
\textbf{General Sibling Selector (bộ chọn anh chị tổng quát)} là một công cụ mạnh trong CSS giúp bạn chọn \textbf{tất cả} các phần tử cùng cấp nằm sau phần tử được chọn, không cần phải liền kề như với Adjacent Sibling Selector (\codeinline{+}). Nó hoạt động bằng cách quét từ trên xuống dưới trong DOM và chọn tất cả các phần tử phù hợp nằm sau phần tử gốc.
\end{conceptBox}

\paragraph{a. Cú pháp:}
\begin{codeWindow}[Cú Pháp General Sibling Selector (\textasciitilde{})]
\begin{lstlisting}[language=CSS3]
A ~ B {
    /* CSS styles */
}
\end{lstlisting}
\end{codeWindow}

\begin{itemize}[leftmargin=1.5em, itemsep=2pt, parsep=0pt]
    \item \codeinline{A}: Phần tử gốc (anchor element).
    \item \codeinline{B}: Phần tử anh chị em tổng quát của A (cùng cấp và nằm sau A).
\end{itemize}

\begin{noteBox}[Điều Kiện Hoạt Động]
\begin{itemize}[leftmargin=1.5em, itemsep=2pt, parsep=0pt]
    \item A và B phải \textbf{cùng cấp} (cùng cha) trong cây DOM.
    \item B phải \textbf{nằm sau} A, nhưng \textbf{không cần liền kề} ngay lập tức.
\end{itemize}
\end{noteBox}

\paragraph{b. Ví dụ cơ bản:}
\begin{codeWindow}[Ví Dụ Cơ Bản: Chọn Tất Cả Thẻ p Sau h2]
\begin{lstlisting}[language=HTML5]
<h2>Header Title</h2>
<p>Sample Text</p>
<p>Sample Text</p>
<p>Sample Text</p>
\end{lstlisting}

\begin{lstlisting}[language=CSS3]
h2 ~ p {
    color: blue;
    font-weight: bold;
}
\end{lstlisting}
\end{codeWindow}

\begin{browserWindow}[Kết Quả Hiển Thị h2 \textasciitilde{} p]
\sffamily\small
{\bfseries\Large\textcolor{darkText}{Tiêu đề}} \par\vspace{4pt}
\textcolor{primaryBlue}{\textbf{Đoạn 1}} {\small\textcolor{gray!60}{$\leftarrow$ được chọn}} \par\vspace{2pt}
\textcolor{primaryBlue}{\textbf{Đoạn 2}} {\small\textcolor{gray!60}{$\leftarrow$ được chọn}} \par\vspace{2pt}
\textcolor{primaryBlue}{\textbf{Đoạn 3}} {\small\textcolor{gray!60}{$\leftarrow$ được chọn}} \par\vspace{4pt}
\textcolor{gray!70}{\small\textit{($\rightarrow$ Tất cả các thẻ \codeinline{<p>} sau \codeinline{<h2>} đều có màu xanh và chữ đậm.)}}
\end{browserWindow}

\paragraph{c. Ứng dụng thực tế:}

\begin{tipBox}[4 Ứng Dụng Thực Tế Mạnh Mẽ Của General Sibling Selector]
\begin{itemize}[leftmargin=1.5em, itemsep=3pt, parsep=0pt]
    \item \textbf{Ứng dụng 1}: Ẩn/hiện nội dung bằng checkbox (không cần JavaScript!).
    \item \textbf{Ứng dụng 2}: Hover làm sáng các mục menu phía sau.
    \item \textbf{Ứng dụng 3}: Đánh dấu viền đỏ cho ghi chú sau tiêu đề.
    \item \textbf{Ứng dụng 4}: Ẩn tất cả thông báo khi nhấn nút.
\end{itemize}
\end{tipBox}

\begin{codeWindow}[Ứng Dụng 1: Toggle Ẩn/Hiện Bằng Checkbox]
\begin{lstlisting}[language=HTML5]
<input type="checkbox" id="toggle">
<p>Content text</p>
<p>Content text</p>
\end{lstlisting}

\begin{lstlisting}[language=CSS3]
/* Paragraph text */
#toggle:checked ~ p {
    display: none;
}
\end{lstlisting}
\end{codeWindow}

\begin{browserWindow}[Kết Quả: Toggle Checkbox Ẩn/Hiện Nội Dung]
\sffamily\small
\tcbox[on line, arc=2pt, colback=white, colframe=darkText, boxrule=1pt, boxsep=1pt]{\small\textsf{x}} {\small\textbf{Ẩn nội dung}} \par\vspace{4pt}
\textcolor{gray!50}{\sout{Nội dung 1}} {\small\textcolor{gray!60}{$\leftarrow$ ẩn đi}} \par\vspace{2pt}
\textcolor{gray!50}{\sout{Nội dung 2}} {\small\textcolor{gray!60}{$\leftarrow$ ẩn đi}} \par\vspace{4pt}
\textcolor{gray!70}{\small\textit{($\rightarrow$ Khi check vào ô, tất cả các đoạn văn sau checkbox sẽ biến mất!)}}
\end{browserWindow}

\begin{codeWindow}[Ứng Dụng 2: Hover Làm Sáng Các Mục Phía Sau]
\begin{lstlisting}[language=HTML5]
<div class="menu">Sample Text</div>
<div class="menu">Sample Text</div>
<div class="menu">Sample Text</div>
\end{lstlisting}

\begin{lstlisting}[language=CSS3]
.menu:hover ~ .menu {
    background-color: yellow;
}
\end{lstlisting}
\end{codeWindow}

\begin{browserWindow}[Kết Quả: Hover Vào Mục 1 Làm Sáng Mục 2 và 3]
\sffamily\small
\tcbox[on line, arc=3pt, colback=primaryBlue!10, colframe=primaryBlue!40, boxrule=1pt, boxsep=3pt, left=8pt, right=8pt]{\textbf{Mục 1}} {\small\textcolor{gray!60}{$\leftarrow$ đang hover}} \par\vspace{3pt}
\tcbox[on line, arc=3pt, colback=yellow!60, colframe=yellow!80!black, boxrule=1pt, boxsep=3pt, left=8pt, right=8pt]{\textbf{Mục 2}} {\small\textcolor{gray!60}{$\leftarrow$ đổi nền vàng}} \par\vspace{3pt}
\tcbox[on line, arc=3pt, colback=yellow!60, colframe=yellow!80!black, boxrule=1pt, boxsep=3pt, left=8pt, right=8pt]{\textbf{Mục 3}} {\small\textcolor{gray!60}{$\leftarrow$ đổi nền vàng}}
\end{browserWindow}

\begin{codeWindow}[Ứng Dụng 3: Đánh Dấu Viền Đỏ Cho Ghi Chú Sau Tiêu Đề]
\begin{lstlisting}[language=HTML5]
<h3>Header Title</h3>
<div class="note">Sample Text</div>
<div class="note">Sample Text</div>
\end{lstlisting}

\begin{lstlisting}[language=CSS3]
h3 ~ .note {
    border-left: 3px solid red;
    padding-left: 10px;
}
\end{lstlisting}
\end{codeWindow}

\begin{browserWindow}[Kết Quả: Viền Đỏ Bên Trái Cho Ghi Chú]
\sffamily\small
{\bfseries\large\textcolor{darkText}{Tiêu đề chính}} \par\vspace{4pt}
\begin{tcolorbox}[arc=0pt, colback=white, colframe=white, boxrule=0pt, boxsep=2pt, left=10pt, right=4pt, borderline west={3pt}{0pt}{red}]
\textcolor{darkText}{\textbf{Lưu ý 1}}
\end{tcolorbox}\vspace{2pt}
\begin{tcolorbox}[arc=0pt, colback=white, colframe=white, boxrule=0pt, boxsep=2pt, left=10pt, right=4pt, borderline west={3pt}{0pt}{red}]
\textcolor{darkText}{\textbf{Lưu ý 2}}
\end{tcolorbox}\vspace{4pt}
\textcolor{gray!70}{\small\textit{($\rightarrow$ Tất cả các phần tử .note sau thẻ \codeinline{<h3>} sẽ có viền đỏ bên trái.)}}
\end{browserWindow}

\begin{codeWindow}[Ứng Dụng 4: Ẩn Tất Cả Thông Báo Khi Nhấn Nút]
\begin{lstlisting}[language=HTML5]
<button id="hideBtn">Sample Text</button>
<div class="alert">Sample Text</div>
<div class="alert">Sample Text</div>
\end{lstlisting}

\begin{lstlisting}[language=CSS3]
#hideBtn:active ~ .alert {
    display: none;
}
\end{lstlisting}
\end{codeWindow}

\begin{browserWindow}[Kết Quả: Nhấn Giữ Nút Ẩn Thông Báo]
\sffamily\small
\tcbox[on line, arc=5pt, colback=red!70, colframe=red!70, boxsep=4pt, left=10pt, right=10pt]{\color{white}\bfseries Ẩn thông báo} \par\vspace{4pt}
\textcolor{gray!50}{\sout{Thông báo 1}} {\small\textcolor{gray!60}{$\leftarrow$ biến mất}} \par\vspace{2pt}
\textcolor{gray!50}{\sout{Thông báo 2}} {\small\textcolor{gray!60}{$\leftarrow$ biến mất}} \par\vspace{4pt}
\textcolor{gray!70}{\small\textit{($\rightarrow$ Khi nhấn giữ nút, tất cả thông báo phía sau sẽ biến mất!)}}
\end{browserWindow}

\paragraph{d. Phân biệt với Adjacent Sibling Selector (+):}

\begin{prettyTableBox}[So Sánh Adjacent Sibling (+) vs General Sibling (\textasciitilde{})]
\rowcolors{2}{white}{primaryBlue!4}
\sffamily\small
\begin{tabularx}{\linewidth}{|>{\raggedright\arraybackslash \ttfamily\bfseries\color{accentOrange}}K{2.0cm}|K{4.0cm}|Z|}
\hline
\rowcolor{primaryBlue} \thd{Bộ chọn} & \thd{Ý nghĩa} & \thd{Hoạt động} \\
\hline
A + B & Chọn phần tử B liền kề ngay sau A & Chỉ chọn phần tử đầu tiên nằm ngay sau A \\
\hline
A \textasciitilde{} B & Chọn tất cả phần tử B nằm sau A & Chọn tất cả các phần tử sau A, không cần liền kề \\
\hline
\end{tabularx}
\end{prettyTableBox}

\begin{codeWindow}[So Sánh Trực Quan: + vs \textasciitilde{}]
\begin{lstlisting}[language=CSS3]
/* Paragraph text */
h2 + p { color: red; }

/* Paragraph text */
h2 ~ p { color: blue; }
\end{lstlisting}
\end{codeWindow}

\paragraph{e. Nguyên tắc hoạt động:}

\begin{noteBox}[3 Nguyên Tắc Hoạt Động Của General Sibling Selector]
\begin{enumerate}[leftmargin=1.5em, itemsep=3pt, parsep=0pt]
    \item \textbf{Cùng cấp trong cây DOM}: Phần tử được chọn phải cùng cấp (sibling) với phần tử gốc.
    \item \textbf{Xét theo thứ tự xuất hiện}: CSS sẽ quét từ trên xuống dưới trong DOM và chọn tất cả các phần tử phù hợp.
    \item \textbf{Không chọn phần tử phía trước}: Chỉ các phần tử nằm \textbf{sau} phần tử gốc mới bị ảnh hưởng.
\end{enumerate}
\end{noteBox}

\paragraph{f. Ví dụ nâng cao (phần tử xen kẽ):}
\begin{codeWindow}[Ví Dụ Nâng Cao: Phần Tử Xen Kẽ Không Ảnh Hưởng]
\begin{lstlisting}[language=HTML5]
<div class="container">
    <h2>Header Title</h2>
    <p>Sample Text</p>
    <span>Sample Text</span>
    <p>Sample Text</p>
    <p>Sample Text</p>
</div>
\end{lstlisting}

\begin{lstlisting}[language=CSS3]
h2 ~ p {
    background-color: yellow;
}
\end{lstlisting}
\end{codeWindow}

\begin{browserWindow}[Kết Quả: Chỉ Thẻ p Được Chọn, span Không Bị Ảnh Hưởng]
\sffamily\small
\begin{tcolorbox}[arc=2pt, colback=primaryBlue!5, colframe=primaryBlue!40, boxrule=1pt, title={\sffamily\small\color{white}\textbf{.container}}, colbacktitle=primaryBlue!70]
{\bfseries\large\textcolor{darkText}{Tiêu đề}} \par\vspace{4pt}
\colorbox{yellow!60}{\textbf{Đoạn văn 1}} {\small\textcolor{gray!60}{$\leftarrow$ nền vàng}} \par\vspace{3pt}
\textcolor{darkText}{\textit{Span không bị ảnh hưởng}} {\small\textcolor{gray!60}{$\leftarrow$ không phải \codeinline{<p>}}} \par\vspace{3pt}
\colorbox{yellow!60}{\textbf{Đoạn văn 2}} {\small\textcolor{gray!60}{$\leftarrow$ nền vàng}} \par\vspace{3pt}
\colorbox{yellow!60}{\textbf{Đoạn văn 3}} {\small\textcolor{gray!60}{$\leftarrow$ nền vàng}}
\end{tcolorbox}
\end{browserWindow}

\paragraph{g. Trường hợp không hoạt động:}

\begin{warnBox}[2 Trường Hợp General Sibling KHÔNG Hoạt Động]
\begin{enumerate}[leftmargin=1.5em, itemsep=3pt, parsep=0pt]
    \item \textbf{Không cùng cấp}: Nếu A và B không cùng cấp (khác phần tử cha), thì \codeinline{\textasciitilde{}} sẽ không hoạt động.
    \item \textbf{B nằm trước A}: Nếu phần tử B xuất hiện \textbf{trước} A trong DOM, bộ chọn cũng không áp dụng CSS cho B.
\end{enumerate}
\end{warnBox}

\begin{codeWindow}[Ví Dụ Sai: B Nằm Trước A]
\begin{lstlisting}[language=HTML5]
<p>Sample Text</p>
<h2>Header Title</h2>
<p>Sample Text</p>
\end{lstlisting}

\begin{lstlisting}[language=CSS3]
/* Header Title */
p ~ h2 {
    color: red;
}
\end{lstlisting}
\end{codeWindow}

\begin{browserWindow}[Kết Quả: Thứ Tự Trong DOM Rất Quan Trọng]
\sffamily\small
\textcolor{darkText}{Đoạn văn 1} {\small\textcolor{gray!60}{$\leftarrow$ nằm trước, không bị ảnh hưởng}} \par\vspace{2pt}
\textcolor{red}{\bfseries\large Tiêu đề} {\small\textcolor{gray!60}{$\leftarrow$ nằm sau \codeinline{<p>}, bị đổi đỏ}} \par\vspace{2pt}
\textcolor{darkText}{Đoạn văn 2} {\small\textcolor{gray!60}{$\leftarrow$ không phải \codeinline{<h2>}, không bị chọn}}
\end{browserWindow}

\paragraph{h. Mẹo sử dụng hiệu quả:}

\begin{tipBox}[Mẹo Sử Dụng General Sibling Selector Hiệu Quả]
\begin{itemize}[leftmargin=1.5em, itemsep=3pt, parsep=0pt]
    \item \textbf{Kết hợp với pseudo-classes} như \codeinline{:hover}, \codeinline{:checked}, hoặc \codeinline{:nth-child} để tạo hiệu ứng phức tạp mà không cần JavaScript.
    \item \textbf{Kết hợp với position và z-index} để xây dựng giao diện động, ví dụ như tooltip, popup menu, hoặc progress bar.
    \item \textbf{Khi nào nên dùng?} Khi bạn muốn chọn \textbf{nhiều phần tử cùng cấp} nằm sau một phần tử cụ thể. Rất hữu ích trong tạo menu động, ẩn/hiện nội dung, progress bar, và form validation.
\end{itemize}
\end{tipBox}

\paragraph{i. Tóm tắt nhanh \& Kết luận:}

\begin{prettyTableBox}[Tóm Tắt Nhanh Về General Sibling Selector]
\rowcolors{2}{white}{primaryBlue!4}
\sffamily\small
\begin{tabularx}{\linewidth}{|>{\raggedright\arraybackslash \bfseries\color{primaryBlue}}K{3.5cm}|Z|}
\hline
\rowcolor{primaryBlue} \thd{Đặc điểm} & \thd{Chi tiết} \\
\hline
Cú pháp & \codeinline{A \textasciitilde{} B \{ \}} \\
\hline
Ký hiệu & Dấu ngã \codeinline{\textasciitilde{}} giữa 2 phần tử \\
\hline
Phạm vi chọn & Chọn tất cả phần tử B cùng cấp nằm sau A \\
\hline
Yêu cầu & A và B phải cùng cha, B phải nằm sau A \\
\hline
Liền kề? & Không cần liền kề, có phần tử xen kẽ vẫn hoạt động \\
\hline
So sánh với + & \codeinline{+} chọn 1 phần tử liền kề, \codeinline{\textasciitilde{}} chọn tất cả phía sau \\
\hline
Ứng dụng & Toggle checkbox, hover nhóm, viền ghi chú, ẩn thông báo \\
\hline
\end{tabularx}
\end{prettyTableBox}

\begin{tipBox}[Kết Luận Về General Sibling Selector]
\begin{itemize}[leftmargin=1.5em, itemsep=3pt, parsep=0pt]
    \item \textbf{General Sibling Selector (\textasciitilde{})} giúp chọn tất cả phần tử cùng cấp nằm sau phần tử gốc, không cần liền kề.
    \item \textbf{Thứ tự trong DOM rất quan trọng}: Chỉ chọn các phần tử xuất hiện sau phần tử gốc.
    \item \textbf{Ứng dụng cực kỳ linh hoạt}: Tạo hiệu ứng hover nhóm, ẩn/hiện nội dung bằng checkbox, làm nổi bật các phần tử liên quan\ldots{} mà không cần JavaScript!
\end{itemize}
\end{tipBox}

\subsubsection{3.5. Bảng Tóm Tắt Dễ Nhớ Về Combinators}

\begin{prettyTableBox}[Bảng Tóm Tắt Dễ Nhớ Các Bộ Chọn Kết Hợp Combinator]
\rowcolors{2}{white}{primaryBlue!4}
\sffamily\small
\begin{tabularx}{\linewidth}{|>{\raggedright\arraybackslash \bfseries\color{primaryBlue}}K{2.8cm}|>{\centering\arraybackslash\ttfamily\bfseries\color{accentOrange}}C{2.0cm}|Z|>{\ttfamily\small}K{3.5cm}|}
\hline
\rowcolor{primaryBlue} \thd{Combinator} & \thd{Ký hiệu} & \thd{Ý nghĩa} & \thd{Ví dụ} \\
\hline
Descendant & (Khoảng trắng) & Chọn tất cả hậu duệ (con/cháu mọi cấp) & .parent .child \\
\hline
Child & > & Chọn con trực tiếp & .parent > .child \\
\hline
Adjacent & + & Chọn phần tử liền kề ngay sau & .box + .text \\
\hline
General & \textasciitilde{} & Chọn tất cả anh chị phía sau & .box \textasciitilde{} .text \\
\hline
\end{tabularx}
\end{prettyTableBox}

\subsubsection{3.6. Ví Dụ Thực Tế: Tạo Menu Điều Hướng Với Trạng Thái Hover}

\begin{codeWindow}[Tạo Menu Dropdown Điều Hướng Bằng Child Selector \& Hover]
\begin{lstlisting}[language=CSS3]
/* Sample Text */
nav ul > li:hover > ul {
    display: block;
}
\end{lstlisting}

\begin{lstlisting}[language=HTML5]
<nav>
    <ul>
        <li>Menu 1
            <ul>
                <li>Sample Text</li>
                <li>Sample Text</li>
            </ul>
        </li>
    </ul>
</nav>
\end{lstlisting}
\end{codeWindow}

\begin{browserWindow}[Kết Quả Hiển Thị: Menu Dropdown Xuất Hiện Khi Hover]
\sffamily\small
\tcbox[on line, arc=3pt, colback=primaryBlue, colframe=primaryBlue, boxsep=4pt, left=10pt, right=10pt]{\color{white}\bfseries Menu 1 $\downarrow$} \par\vspace{2pt}
\begin{tcolorbox}[arc=2pt, colback=white, colframe=primaryBlue!40, boxrule=1pt, width=4.0cm, boxsep=4pt, left=8pt, right=8pt]
\textcolor{darkText}{\textbf{Mục 1}} \par\vspace{2pt}
\textcolor{darkText}{\textbf{Mục 2}}
\end{tcolorbox} \par\vspace{4pt}
\textcolor{gray!70}{\small\textit{($\rightarrow$ Kết quả: Khi di chuột (hover) vào Menu 1, danh sách menu con dropdown tự động xuất hiện!)}}
\end{browserWindow}

\begin{tipBox}[Kết Luận Về Bộ Chọn Kết Hợp Combinator Selectors]
\begin{itemize}[leftmargin=1.5em, itemsep=3pt, parsep=0pt]
    \item \textbf{Linh hoạt mã nguồn}: Combinator selectors giúp bạn viết CSS linh hoạt hơn, chọn đúng phần tử mà không cần thêm class hoặc ID thừa vào HTML.
    \item \textbf{Tối ưu hiệu suất}: Hiểu sâu về các selector này sẽ giúp bạn xây dựng giao diện gọn gàng, dễ bảo trì và tối ưu hiệu suất tải trang!
\end{itemize}
\end{tipBox}


%% ===================================================================
%% MỤC MỚI: CHUYÊN ĐỀ PHẦN TỬ CÙNG CẤP \& BỘ CHỌN COMBINATOR STRUCTURAL
%% ===================================================================

\subsection{4. Chuyên Đề Phân Biệt Chi Tiết: "Phần Tử Cùng Cấp" \& "Phần Tử Con Trực Tiếp" (Masterclass Structural Combinators)}

\begin{conceptBox}[Khái Niệm Phần Tử Cùng Cấp (Sibling Element) Là Gì?]
\textbf{Phần tử cùng cấp (\term{Sibling Element})} là những phần tử nằm trong \textbf{cùng một phần tử cha (parent)}.
\par\vspace{3pt}
\textbf{Nói cách khác}: Chúng có cùng "cha trực tiếp" trong cây DOM.
\par\vspace{3pt}
\begin{tipBox}[Mẹo Dễ Nhớ Về Phần Tử Cùng Cấp]
Hai phần tử HTML là cùng cấp (\term{sibling}) nếu chúng \textbf{đứng ngang hàng nhau} trong cùng một khối (\codeinline{<div>}, \codeinline{<section>}, \codeinline{<body>}\ldots).
\end{tipBox}
\end{conceptBox}

\paragraph{a. Ví dụ minh họa phần tử cùng cấp vs KHÔNG cùng cấp:}

\begin{codeWindow}[Ví Dụ 1: Các Phần Tử Cùng Cấp (Sibling Elements)]
\begin{lstlisting}[language=HTML5]
<div class="parent">
    <h3>Header Title</h3>        <!-- A -->
    <p>Sample Text</p>        <!-- B -->
    <p>Sample Text</p>        <!-- C -->
</div>
\end{lstlisting}
\end{codeWindow}

\begin{noteBox}[Giải Thích Ví Dụ 1 (Cùng Cấp)]
Ở đây: \codeinline{<h3>}, \codeinline{<p>} (Đoạn 1), và \codeinline{<p>} (Đoạn 2) đều nằm trong cùng 1 thẻ \codeinline{<div class="parent">}.
\par
$\rightarrow$ \textbf{Vì vậy, chúng là phần tử cùng cấp (sibling elements).}
\end{noteBox}

\begin{codeWindow}[Ví Dụ 2: Các Phần Tử KHÔNG Cùng Cấp]
\begin{lstlisting}[language=HTML5]
<div class="parent">
    <h3>Header Title</h3>
    <div>
        <p>Sample Text</p>  <!-- Sample Text -->
    </div>
</div>
\end{lstlisting}
\end{codeWindow}

\begin{warnBox}[Giải Thích Ví Dụ 2 (KHÔNG Cùng Cấp)]
Ở đây:
\begin{itemize}[leftmargin=1.5em, itemsep=2pt, parsep=0pt]
    \item \codeinline{<h3>} là con của \codeinline{div.parent}.
    \item \codeinline{<p>} là con của \codeinline{<div>} khác nằm bên trong \codeinline{div.parent}.
\end{itemize}
$\rightarrow$ \textbf{\codeinline{<h3>} và \codeinline{<p>} KHÔNG cùng cấp.}
\end{warnBox}

\paragraph{b. Phân biệt rõ ràng: "Phần tử con trực tiếp" và "Phần tử cùng cấp":}

\begin{prettyTableBox}[Bảng Phân Biệt Khái Niệm: Phần Tử Con Trực Tiếp vs Phần Tử Cùng Cấp]
\rowcolors{2}{white}{primaryBlue!4}
\sffamily\small
\begin{tabularx}{\linewidth}{|>{\raggedright\arraybackslash \bfseries\color{primaryBlue}}K{3.2cm}|Z|Z|}
\hline
\rowcolor{primaryBlue} \thd{Khái niệm} & \thd{Phần tử con trực tiếp (Direct Child)} & \thd{Phần tử cùng cấp (Sibling Element)} \\
\hline
\textbf{Là gì?} & Là phần tử nằm bên trong phần tử cha, ngay trực tiếp -- không qua trung gian. & Là hai phần tử cùng nằm trong một phần tử cha -- tức là anh em với nhau. \\
\hline
\end{tabularx}
\end{prettyTableBox}

\begin{codeWindow}[Ví Dụ Minh Họa Phân Biệt Cấu Trúc Cha - Con - Cùng Cấp]
\begin{lstlisting}[language=HTML5]
<div class="cha">
    <h3>Header Title</h3>         <!-- Sample Text -->
    <p>Sample Text</p>            <!-- Sample Text -->
    <div>
        <p>Sample Text</p>          <!-- Sample Text -->
    </div>
</div>
\end{lstlisting}
\end{codeWindow}

\begin{noteBox}[Giải Thích Chi Tiết Cấu Trúc Trong HTML Trên]
\begin{itemize}[leftmargin=1.5em, itemsep=3pt, parsep=0pt]
    \item \codeinline{<h3>} và \codeinline{<p>Đoạn 1} $\rightarrow$ \textcolor{tipBorder}{\textbf{Cùng cấp}}, vì cùng nằm trực tiếp trong \codeinline{<div class="cha">}.
    \item \codeinline{<p>Đoạn 1} là \textcolor{tipBorder}{\textbf{con trực tiếp}} của \codeinline{div.cha}.
    \item \codeinline{<p>Đoạn 2} \textcolor{warnBorder}{\textbf{KHÔNG phải con trực tiếp}} của \codeinline{div.cha} mà là con của một thẻ \codeinline{<div>} khác bên trong.
    \item \codeinline{<p>Đoạn 2} \textcolor{warnBorder}{\textbf{KHÔNG cùng cấp với <h3>}}, vì nằm sâu hơn một cấp.
\end{itemize}
\end{noteBox}

\begin{prettyTableBox}[Tóm Lại So Sánh Đặc Điểm]
\rowcolors{2}{white}{primaryBlue!4}
\sffamily\small
\begin{tabularx}{\linewidth}{|>{\raggedright\arraybackslash \bfseries\color{primaryBlue}}K{4.5cm}|>{\centering\arraybackslash}C{3.5cm}|>{\centering\arraybackslash}C{3.5cm}|}
\hline
\rowcolor{primaryBlue} \thd{Tiêu chí so sánh} & \thd{Phần tử con trực tiếp} & \thd{Phần tử cùng cấp} \\
\hline
Cùng nằm trong phần tử cha & \textbf{Có} & \textbf{Có} \\
\hline
Phải là con của cha & \textbf{Có} & \textbf{Có} \\
\hline
Phải là ngay bên trong & \textbf{Có} & \textbf{Không bắt buộc} (chỉ cần chung cha) \\
\hline
\end{tabularx}
\end{prettyTableBox}

\begin{tipBox}[Mẹo Ghi Nhớ Cốt Lõi]
\begin{itemize}[leftmargin=1.5em, itemsep=3pt, parsep=0pt]
    \item \textbf{Cùng cấp} = "Ngang hàng" (Anh em ruột cùng cha).
    \item \textbf{Con trực tiếp} = "Con ruột" của phần tử cha.
\end{itemize}
\end{tipBox}

\paragraph{c. Chuyên đề thực nghiệm 1: Bộ chọn anh chị em tổng quát (\codeinline{h3 \textasciitilde{} p}):}

\begin{codeWindow}[Đoạn Mã HTML Và CSS Thực Nghiệm Bộ Chọn h3 \textasciitilde{} p]
\begin{lstlisting}[language=HTML5]
<div class='Parent'>
    <h3>lorem different</h3>
    <p>lorem is apple </p>       <!-- 1 -->
    <p>lorem is apple </p>       <!-- 2 -->
    <div class='child_1'>
        <p>lorem is apple </p>     <!-- 3 -->
        <p>lorem is apple </p>     <!-- 4 -->
    </div>
    <p>lorem is apple </p>       <!-- 5 -->
    <div class='child_2'>
        <p>lorem is apple </p>     <!-- 6 -->
        <p>lorem is apple </p>     <!-- 7 -->
    </div>
</div>
\end{lstlisting}

\begin{lstlisting}[language=CSS3]
h3 ~ p {
    color: red;
}
\end{lstlisting}
\end{codeWindow}

\begin{conceptBox}[Giải Thích Cơ Chế Hoạt Động Của Selector h3 \textasciitilde{} p]
Selector \codeinline{h3 \textasciitilde{} p} là \term{General Sibling Combinator}, nghĩa là tất cả các phần tử \codeinline{<p>} \textbf{cùng cấp (siblings)} với \codeinline{<h3>} và \textbf{đứng sau} \codeinline{<h3>} trong cùng một phần tử cha sẽ được áp dụng style.
\end{conceptBox}

\begin{browserWindow}[Kết Quả Tô Màu Đỏ Cho Selector h3 \textasciitilde{} p]
\sffamily\small
\begin{tcolorbox}[arc=2pt, colback=primaryBlue!5, colframe=primaryBlue!40, boxrule=1pt, title={\sffamily\small\color{white}\textbf{.Parent}}, colbacktitle=primaryBlue!70]
{\bfseries\large\textcolor{darkText}{<h3> lorem different </h3>}} \par\vspace{4pt}
\textcolor{red}{\textbf{<p> lorem is apple </p> (Số 1)}} {\small\textcolor{gray!60}{$\leftarrow$ Đỏ (Sibling trực tiếp sau h3)}} \par\vspace{2pt}
\textcolor{red}{\textbf{<p> lorem is apple </p> (Số 2)}} {\small\textcolor{gray!60}{$\leftarrow$ Đỏ (Sibling trực tiếp sau h3)}} \par\vspace{3pt}
\begin{tcolorbox}[arc=2pt, colback=white, colframe=gray!30, boxrule=1pt, title={\sffamily\scriptsize\color{darkText}.child\_1}]
\textcolor{darkText}{<p> lorem is apple </p> (Số 3)} {\small\textcolor{gray!60}{$\leftarrow$ Không đổi màu (nằm trong .child\_1)}} \par\vspace{2pt}
\textcolor{darkText}{<p> lorem is apple </p> (Số 4)} {\small\textcolor{gray!60}{$\leftarrow$ Không đổi màu (nằm trong .child\_1)}}
\end{tcolorbox}\vspace{3pt}
\textcolor{red}{\textbf{<p> lorem is apple </p> (Số 5)}} {\small\textcolor{gray!60}{$\leftarrow$ Đỏ (Sibling trực tiếp sau h3)}} \par\vspace{3pt}
\begin{tcolorbox}[arc=2pt, colback=white, colframe=gray!30, boxrule=1pt, title={\sffamily\scriptsize\color{darkText}.child\_2}]
\textcolor{darkText}{<p> lorem is apple </p> (Số 6)} {\small\textcolor{gray!60}{$\leftarrow$ Không đổi màu (nằm trong .child\_2)}} \par\vspace{2pt}
\textcolor{darkText}{<p> lorem is apple </p> (Số 7)} {\small\textcolor{gray!60}{$\leftarrow$ Không đổi màu (nằm trong .child\_2)}}
\end{tcolorbox}
\end{tcolorbox}
\end{browserWindow}

\begin{noteBox}[Chi Tiết Phân Tích Kết Quả Áp Dụng]
\begin{itemize}[leftmargin=1.5em, itemsep=3pt, parsep=0pt]
    \item \textcolor{tipBorder}{\textbf{Áp dụng CSS color: red cho}}: \codeinline{<p>} số 1, \codeinline{<p>} số 2, và \codeinline{<p>} số 5. Vì cả 3 thẻ này là sibling trực tiếp của \codeinline{<h3>}, nằm trong cùng thẻ cha \codeinline{.Parent} và đứng sau \codeinline{<h3>}.
    \item \textcolor{warnBorder}{\textbf{KHÔNG áp dụng cho}}:
    \begin{itemize}[leftmargin=1.2em, itemsep=2pt, parsep=0pt]
        \item \codeinline{<p>} số 3 và số 4: nằm bên trong \codeinline{.child\_1}.
        \item \codeinline{<p>} số 6 và số 7: nằm bên trong \codeinline{.child\_2}.
    \end{itemize}
    Vì chúng không phải sibling trực tiếp của \codeinline{<h3>}, mà là con của các thẻ \codeinline{<div>} khác.
\end{itemize}
\end{noteBox}

\begin{conceptBox}[Sơ Đồ Phân Cấp Cây DOM Tổng Kết Cho h3 \textasciitilde{} p]
\begin{center}
\begin{tikzpicture}[
  font=\sffamily\small,
  node distance=0.4cm and 0.6cm,
  root_box/.style={draw=primaryBlue, fill=primaryBlue!10, rounded corners=3pt, line width=1pt, inner sep=4pt, font=\sffamily\bfseries\color{primaryBlue}},
  h3_box/.style={draw=gray!50, fill=gray!8, rounded corners=3pt, line width=0.8pt, inner sep=4pt},
  p_red/.style={draw=red!70, fill=red!10, rounded corners=3pt, line width=1pt, inner sep=4pt, font=\sffamily\bfseries\color{red!80!black}},
  p_norm/.style={draw=gray!40, fill=gray!5, rounded corners=3pt, line width=0.8pt, inner sep=4pt, font=\sffamily\color{darkText}},
  div_box/.style={draw=gray!40, fill=gray!5, rounded corners=3pt, line width=0.8pt, inner sep=4pt},
  badge_red/.style={fill=red!70, text=white, font=\sffamily\tiny\bfseries, rounded corners=2pt, inner sep=2pt},
  badge_gray/.style={fill=gray!50, text=white, font=\sffamily\tiny\bfseries, rounded corners=2pt, inner sep=2pt}
]
  % Root h3 anchor
  \node[h3_box] (h3) {\codeinline{<h3>} Tiêu đề \codeinline{</h3>}};

  % Siblings under parent
  \node[p_red, below=0.35cm of h3.south west, anchor=north west] (p1) {\codeinline{<p>} (1) lorem is apple \codeinline{</p>}};
  \node[badge_red, right=4pt of p1] {Đỏ (Sibling)};

  \node[p_red, below=0.35cm of p1.south west, anchor=north west] (p2) {\codeinline{<p>} (2) lorem is apple \codeinline{</p>}};
  \node[badge_red, right=4pt of p2] {Đỏ (Sibling)};

  \node[div_box, below=0.35cm of p2.south west, anchor=north west] (c1) {\codeinline{<div class="child\_1">}};

  \node[p_norm, below right=0.35cm and 0.5cm of c1.south west, anchor=north west] (p3) {\codeinline{<p>} (3) và (4) lorem is apple \codeinline{</p>}};
  \node[badge_gray, right=4pt of p3] {Bỏ qua (Con của child\_1)};

  \node[p_red, below=1.25cm of c1.south west, anchor=north west] (p5) {\codeinline{<p>} (5) lorem is apple \codeinline{</p>}};
  \node[badge_red, right=4pt of p5] {Đỏ (Sibling)};

  \node[div_box, below=0.35cm of p5.south west, anchor=north west] (c2) {\codeinline{<div class="child\_2">}};

  \node[p_norm, below right=0.35cm and 0.5cm of c2.south west, anchor=north west] (p6) {\codeinline{<p>} (6) và (7) lorem is apple \codeinline{</p>}};
  \node[badge_gray, right=4pt of p6] {Bỏ qua (Con của child\_2)};

  % Connectors
  \draw[primaryBlue!70, line width=1pt] (h3.south west) ++(-0.2,0) |- (p1.west);
  \draw[primaryBlue!70, line width=1pt] (h3.south west) ++(-0.2,0) |- (p2.west);
  \draw[primaryBlue!70, line width=1pt] (h3.south west) ++(-0.2,0) |- (c1.west);
  \draw[accentOrange!80, line width=1pt] (c1.south west) ++(0.3,0) |- (p3.west);
  \draw[primaryBlue!70, line width=1pt] (h3.south west) ++(-0.2,0) |- (p5.west);
  \draw[primaryBlue!70, line width=1pt] (h3.south west) ++(-0.2,0) |- (c2.west);
  \draw[accentOrange!80, line width=1pt] (c2.south west) ++(0.3,0) |- (p6.west);

\end{tikzpicture}
\end{center}
\end{conceptBox}

\begin{warnBox}[Lưu Ý Nâng Cao \& Hướng Xử Lý Trong Thực Tế]
Trong CSS, với bộ chọn \codeinline{h3 \textasciitilde{} p}, chỉ các phần tử \codeinline{<p>} nào là \textbf{"cùng cấp"} (tức là sibling trực tiếp của \codeinline{<h3>}) và nằm sau \codeinline{<h3>} mới bị ảnh hưởng. Còn những phần tử \codeinline{<p>} không cùng cấp (ví dụ như nằm trong \codeinline{div.child\_1}, \codeinline{div.child\_2}) thì \textbf{KHÔNG} bị ảnh hưởng bởi selector này.
\par\vspace{3pt}
\begin{codeWindow}[Minh Họa Nhanh Khác Biệt Cấp Độ]
\begin{lstlisting}[language=HTML5]
<div class="Parent">
    <h3>Header Title</h3>          <!-- Sample Text -->
    <p>Sample Text</p>        <!-- Sample Text -->
    <div>
        <p>Sample Text</p>  <!-- Sample Text -->
    </div>
</div>
\end{lstlisting}
\end{codeWindow}
\textbf{Giải thích}: Cùng cấp nghĩa là 2 phần tử có chung cha trực tiếp. Nếu phần tử đó nằm lồng bên trong một thẻ khác $\rightarrow$ không còn cùng cấp.
\par\vspace{3pt}
\textbf{Giải pháp}: Nếu bạn muốn tất cả các \codeinline{<p>} nằm sau \codeinline{<h3>} đều bị ảnh hưởng, bất kể cấp lồng nào, bạn không thể làm điều đó chỉ bằng CSS hiện tại (vì CSS không có bộ chọn "mọi phần tử nằm sau ở mọi cấp"). Bạn cần dùng JavaScript để quét DOM hoặc thay đổi cấu trúc HTML / thêm class phù hợp.
\end{warnBox}

\paragraph{d. Chuyên đề thực nghiệm 2: Bộ chọn con trực tiếp (\codeinline{.Parent > p} và \codeinline{div > p}):}

\begin{codeWindow}[Ví Dụ 1: Sử Dụng Selector .Parent > p]
\begin{lstlisting}[language=CSS3]
.Parent > p {
    color: red;
}
\end{lstlisting}
\end{codeWindow}

\begin{conceptBox}[Ý Nghĩa Selector .Parent > p]
Selector \codeinline{.Parent > p} chọn tất cả các phần tử \codeinline{<p>} là \textbf{con trực tiếp (direct children)} của phần tử có class \codeinline{.Parent}.
\end{conceptBox}

\begin{prettyTableBox}[Phân Tích Từng Thẻ Trong HTML Khi Áp Dụng .Parent > p]
\rowcolors{2}{white}{primaryBlue!4}
\sffamily\small
\begin{tabularx}{\linewidth}{|>{\raggedright\arraybackslash \bfseries}K{3.5cm}|>{\centering\arraybackslash}C{2.5cm}|Z|}
\hline
\rowcolor{primaryBlue} \thd{Thẻ HTML} & \thd{Kết quả} & \thd{Giải thích} \\
\hline
\codeinline{<h3>lorem different</h3>} & Không tác động & Không phải thẻ \codeinline{<p>}, bỏ qua. \\
\hline
\codeinline{<p>lorem is apple </p>} (Số 1) & \textcolor{red}{\textbf{Tô đỏ}} & Con trực tiếp của \codeinline{.Parent}. \\
\hline
\codeinline{<p>lorem is apple </p>} (Số 2) & \textcolor{red}{\textbf{Tô đỏ}} & Con trực tiếp của \codeinline{.Parent}. \\
\hline
Thẻ \codeinline{<p>} số 3 \& 4 trong \codeinline{.child\_1} & Bỏ qua & Con của \codeinline{div.child\_1}, KHÔNG phải con trực tiếp của \codeinline{.Parent}. \\
\hline
\codeinline{<p>lorem is apple </p>} (Số 5) & \textcolor{red}{\textbf{Tô đỏ}} & Con trực tiếp của \codeinline{.Parent}. \\
\hline
Thẻ \codeinline{<p>} số 6 \& 7 trong \codeinline{.child\_2} & Bỏ qua & Con của \codeinline{div.child\_2}, KHÔNG phải con trực tiếp của \codeinline{.Parent}. \\
\hline
\end{tabularx}
\end{prettyTableBox}

\begin{browserWindow}[Kết Luận Cho Ví Dụ 1 (.Parent > p): 3 Thẻ p Bị Tô Đỏ]
\sffamily\small
\begin{tcolorbox}[arc=2pt, colback=primaryBlue!5, colframe=primaryBlue!40, boxrule=1pt, title={\sffamily\small\color{white}\textbf{.Parent}}, colbacktitle=primaryBlue!70]
\textcolor{red}{\textbf{<p> lorem is apple </p> (Số 1 sau <h3>)}} $\rightarrow$ Tô đỏ \par\vspace{2pt}
\textcolor{red}{\textbf{<p> lorem is apple </p> (Số 2 trước div.child\_1)}} $\rightarrow$ Tô đỏ \par\vspace{3pt}
\begin{tcolorbox}[arc=2pt, colback=white, colframe=gray!30, boxrule=1pt, title={\sffamily\scriptsize\color{darkText}.child\_1}]
<p> lorem is apple </p> (Số 3 \& 4) $\rightarrow$ Không bị ảnh hưởng
\end{tcolorbox}\vspace{3pt}
\textcolor{red}{\textbf{<p> lorem is apple </p> (Số 5 sau div.child\_1)}} $\rightarrow$ Tô đỏ \par\vspace{3pt}
\begin{tcolorbox}[arc=2pt, colback=white, colframe=gray!30, boxrule=1pt, title={\sffamily\scriptsize\color{darkText}.child\_2}]
<p> lorem is apple </p> (Số 6 \& 7) $\rightarrow$ Không bị ảnh hưởng
\end{tcolorbox}
\end{tcolorbox}
\end{browserWindow}

\begin{codeWindow}[Ví Dụ 2: Sử Dụng Selector div > p]
\begin{lstlisting}[language=CSS3]
div > p {
    color: red;
}
\end{lstlisting}
\end{codeWindow}

\begin{conceptBox}[Ý Nghĩa Selector div > p]
Selector \codeinline{div > p} sẽ chọn tất cả các thẻ \codeinline{<p>} là \textbf{con trực tiếp} của \textbf{bất kỳ thẻ \codeinline{<div>} nào}.
\end{conceptBox}

\begin{noteBox}[Phân Tích Tác Động Của div > p Lên Mã HTML]
\begin{itemize}[leftmargin=1.5em, itemsep=3pt, parsep=0pt]
    \item Thẻ \codeinline{<p>} 1, 2, 5 $\rightarrow$ là con trực tiếp của \codeinline{div.Parent} $\rightarrow$ \textcolor{red}{\textbf{Tô đỏ}}.
    \item Thẻ \codeinline{<p>} 3, 4 $\rightarrow$ là con trực tiếp của \codeinline{div.child\_1} $\rightarrow$ \textcolor{red}{\textbf{Tô đỏ}}.
    \item Thẻ \codeinline{<p>} 6, 7 $\rightarrow$ là con trực tiếp của \codeinline{div.child\_2} $\rightarrow$ \textcolor{red}{\textbf{Tô đỏ}}.
    \item Thẻ \codeinline{<h3>} $\rightarrow$ Không phải thẻ \codeinline{<p>}, bỏ qua.
\end{itemize}
\textbf{Kết luận}: \textcolor{red}{\textbf{TẤT CẢ 7 THẺ <p> ĐỀU BỊ TÔ ĐỎ!}}
\end{noteBox}

\begin{warnBox}[Cảnh Báo Về Bộ Chọn Quá Tổng Quát]
Selector càng tổng quát (\codeinline{div > p}) thì càng dễ "vô tình" áp dụng lên nhiều phần tử hơn mong muốn. Nếu bạn chỉ muốn định dạng các thẻ \codeinline{<p>} trực thuộc \codeinline{.Parent} mà không làm ảnh hưởng tới \codeinline{.child\_1} hay \codeinline{.child\_2}, bắt buộc phải dùng bộ chọn cụ thể hơn: \codeinline{.Parent > p}.
\end{warnBox}

\paragraph{e. Chuyên đề thực nghiệm 3: Bộ chọn hậu duệ (\codeinline{.Parent p}):}

\begin{codeWindow}[Đoạn Mã Selector Hậu Duệ .Parent p]
\begin{lstlisting}[language=CSS3]
.Parent p {
    color: red;
}
\end{lstlisting}
\end{codeWindow}

\begin{conceptBox}[Ý Nghĩa Selector .Parent p]
Selector \codeinline{.Parent p} chọn \textbf{tất cả các thẻ \codeinline{<p>} nằm bên trong (bất kỳ cấp lồng nào)} của phần tử có class \codeinline{.Parent}.
\par\vspace{2pt}
\textbf{Lưu ý}: Không cần phải là con trực tiếp -- chỉ cần nằm lồng trong \codeinline{.Parent} là được!
\end{conceptBox}

\begin{browserWindow}[Kết Quả Cho .Parent p: TẤT CẢ 7 Thẻ p Đều Tô Đỏ]
\sffamily\small
\textcolor{red}{\textbf{TẤT CẢ 7 THẺ <p> ĐỀU BỊ TÔ ĐỎ}} vì chúng đều nằm bên trong phần tử \codeinline{.Parent}, dù là con trực tiếp (P1, P2, P5) hay con gián tiếp/cháu (P3, P4, P6, P7).
\end{browserWindow}

\paragraph{f. Bảng tổng hợp so sánh sự khác biệt giữa các Selector:}

\begin{prettyTableBox}[Bảng So Sánh Tổng Hợp Sự Khác Biệt Giữa Các Structural Selectors]
\rowcolors{2}{white}{primaryBlue!4}
\sffamily\small
\begin{tabularx}{\linewidth}{|>{\raggedright\arraybackslash \ttfamily\bfseries\color{primaryBlue}}K{2.8cm}|Z|>{\centering\arraybackslash\bfseries}C{3.2cm}|}
\hline
\rowcolor{primaryBlue} \thd{Selector} & \thd{Mô tả phạm vi chọn} & \thd{Số thẻ <p> bị ảnh hưởng} \\
\hline
.Parent p & Tất cả \codeinline{<p>} nằm bên trong \codeinline{.Parent} (mọi cấp lồng) & \textcolor{red}{\textbf{7 thẻ}} (Tất cả) \\
\hline
.Parent > p & Chỉ các \codeinline{<p>} là \textbf{con trực tiếp} của \codeinline{.Parent} & \textcolor{tipBorder}{\textbf{3 thẻ}} (P1, P2, P5) \\
\hline
div > p & Bất kỳ \codeinline{<p>} nào là \textbf{con trực tiếp} của bất kỳ \codeinline{<div>} nào & \textcolor{red}{\textbf{7 thẻ}} (Tất cả) \\
\hline
h3 \textasciitilde{} p & Các \codeinline{<p>} \textbf{cùng cấp (siblings)}, nằm phía sau \codeinline{<h3>} & \textcolor{tipBorder}{\textbf{3 thẻ}} (P1, P2, P5) \\
\hline
\end{tabularx}
\end{prettyTableBox}

\begin{tipBox}[Kết Luận Lựa Chọn Selector Phù Hợp]
Việc lựa chọn bộ chọn chặt chẽ (\codeinline{.Parent > p}, \codeinline{h3 \textasciitilde{} p}) hay tổng quát (\codeinline{.Parent p}, \codeinline{div > p}) tùy thuộc hoàn toàn vào mục tiêu thiết kế giao diện của bạn. Hãy ưu tiên các bộ chọn cụ thể và có phạm vi kiểm soát tốt để tránh lỗi giao diện phát sinh ngoài ý muốn!
\end{tipBox}

%% ===================================================================
%% MỤC MỚI: SO SÁNH TỔNG HỢP TẤT CẢ CÁC BỘ CHỌN KẾT HỢP TRONG CSS
%% ===================================================================

\subsection{5. So Sánh Tổng Hợp Chi Tiết Và Dễ Hiểu Nhất Về Tất Cả Các Bộ Chọn Kết Hợp Trong CSS}


\begin{conceptBox}[Tổng Quan: 4 Loại Bộ Chọn Kết Hợp (Combinator Selectors)]
CSS cung cấp 4 loại bộ chọn kết hợp (Combinator Selectors) giúp bạn chọn phần tử dựa trên \textbf{mối quan hệ phân cấp} và \textbf{vị trí tương đối} trong cây DOM. Hiểu rõ từng loại, cách hoạt động và ứng dụng thực tế sẽ giúp bạn viết CSS tinh gọn, dễ bảo trì và hiệu suất cao.
\end{conceptBox}

\subsubsection{5.1. Bảng So Sánh Chi Tiết 4 Bộ Chọn Kết Hợp}

\begin{prettyTableBox}[Bảng So Sánh Tổng Hợp 4 Bộ Chọn Kết Hợp CSS Combinator]
\rowcolors{2}{white}{primaryBlue!4}
\sffamily\small
\begin{tabularx}{\linewidth}{|>{\raggedright\arraybackslash \bfseries\color{primaryBlue}}K{2.4cm}|>{\centering\arraybackslash\ttfamily\bfseries\color{accentOrange}}C{1.8cm}|Z|>{\ttfamily\small}K{3.0cm}|}
\hline
\rowcolor{primaryBlue} \thd{Bộ chọn} & \thd{Cú pháp} & \thd{Ý nghĩa} & \thd{Ví dụ} \\
\hline
Descendant Selector & cha con & Chọn \textbf{tất cả phần tử con cháu} bên trong phần tử cha, không phân biệt cấp độ lồng nhau. & .container p \{ color: blue; \} \\
\hline
Child Selector & cha > con & Chọn \textbf{chỉ con trực tiếp} của phần tử cha (cháu, chắt sẽ bị bỏ qua). & .container > p \{ color: red; \} \\
\hline
Adjacent Sibling & e1 + e2 & Chọn phần tử \textbf{anh chị em liền kề ngay sau} phần tử đầu tiên, cùng cấp trong DOM. & h1 + p \{ color: green; \} \\
\hline
General Sibling & e1 \textasciitilde{} e2 & Chọn \textbf{tất cả phần tử anh chị em cùng cấp} nằm sau phần tử đầu tiên. & h1 \textasciitilde{} p \{ color: purple; \} \\
\hline
\end{tabularx}
\end{prettyTableBox}

\subsubsection{5.2. Ví Dụ Minh Họa Trực Quan -- Áp Dụng Đồng Thời Cả 4 Bộ Chọn}

\begin{codeWindow}[HTML Mẫu: Cấu Trúc DOM Dùng Chung Cho Cả 4 Bộ Chọn Kết Hợp]
\begin{lstlisting}[language=HTML5]
<div class="container">
    <h1>Header Title</h1>
    <p>Sample Text</p>
    <p>Sample Text</p>
    <div>
        <p>Sample Text</p>
    </div>
    <p>Sample Text</p>
</div>
\end{lstlisting}
\end{codeWindow}

\begin{codeWindow}[CSS: Áp Dụng Đồng Thời Tất Cả 4 Bộ Chọn Kết Hợp]
\begin{lstlisting}[language=CSS3]
/* Select element */
.container p {
    color: blue;
}

/* Sample Text */
.container > p {
    color: red;
}

/* Select element */
h1 + p {
    color: green;
}

/* Select element */
h1 ~ p {
    color: purple;
}
\end{lstlisting}
\end{codeWindow}

\begin{browserWindow}[Kết Quả Hiển Thị: Tổng Hợp 4 Bộ Chọn Kết Hợp Trong Cây DOM]
\begin{tcolorbox}[arc=4pt, colback=white, colframe=primaryBlue!30, boxrule=0.8pt, top=10pt, bottom=10pt, left=12pt, right=12pt]
  {\sffamily\bfseries\color{primaryBlue} Kết quả tương tác màu sắc của 4 bộ chọn:}\par\vspace{6pt}
  \begin{itemize}[leftmargin=1.2em, itemsep=3pt, parsep=0pt]
    \item \codeinline{h1 + p}: Nhắm vào thẻ \codeinline{<p>} ngay sau \codeinline{<h1>} $\rightarrow$ \textbf{\color{green!70!black}Màu Xanh Lá (Green)}.
    \item \codeinline{.container > p}: Nhắm vào thẻ \codeinline{<p>} con trực tiếp $\rightarrow$ \textbf{\color{red!80!black}Màu Đỏ (Red)} đè lên màu tím.
    \item \codeinline{.container p}: Nhắm vào thẻ \codeinline{<p>} nằm sâu bên trong \codeinline{<div>} con $\rightarrow$ \textbf{\color{blue}Màu Xanh Dương (Blue)}.
  \end{itemize}
\end{tcolorbox}
\end{browserWindow}

\begin{conceptBox}[Sơ Đồ Phân Cấp Cây DOM Cho Ví Dụ Tổng Hợp]
\begin{center}
\begin{tikzpicture}[
  font=\sffamily\small,
  node distance=0.4cm and 0.6cm,
  root_box/.style={draw=primaryBlue, fill=primaryBlue!10, rounded corners=3pt, line width=1pt, inner sep=4pt, font=\sffamily\bfseries\color{primaryBlue}},
  h1_box/.style={draw=gray!50, fill=gray!8, rounded corners=3pt, line width=0.8pt, inner sep=4pt},
  p_box/.style={draw=primaryBlue, fill=primaryBlue!12, rounded corners=3pt, line width=1pt, inner sep=4pt},
  div_box/.style={draw=gray!40, fill=gray!5, rounded corners=3pt, line width=0.8pt, inner sep=4pt},
  badge_green/.style={fill=tipBorder, text=white, font=\sffamily\tiny\bfseries, rounded corners=2pt, inner sep=2pt},
  badge_purple/.style={fill=purple!70, text=white, font=\sffamily\tiny\bfseries, rounded corners=2pt, inner sep=2pt},
  badge_blue/.style={fill=primaryBlue, text=white, font=\sffamily\tiny\bfseries, rounded corners=2pt, inner sep=2pt},
  badge_gray/.style={fill=gray!50, text=white, font=\sffamily\tiny\bfseries, rounded corners=2pt, inner sep=2pt}
]
  % Root
  \node[root_box] (root) {.container \quad \scriptsize\textcolor{gray!70}{(Gốc DOM)}};

  % Level 1: h1
  \node[h1_box, below right=0.45cm and 0.6cm of root.south west, anchor=north west] (h1) {\codeinline{<h1>} Tiêu đề chính \codeinline{</h1>}};
  \node[badge_gray, right=4pt of h1] {Không bị chọn};

  % Level 1: p1
  \node[p_box, below=0.3cm of h1.south west, anchor=north west] (p1) {\codeinline{<p>} Đoạn văn 1 \codeinline{</p>}};
  \node[badge_green, right=4pt of p1] {h1 + p (liền kề)};

  % Level 1: p2
  \node[p_box, below=0.3cm of p1.south west, anchor=north west] (p2) {\codeinline{<p>} Đoạn văn 2 \codeinline{</p>}};
  \node[badge_purple, right=4pt of p2] {h1 \textasciitilde{} p (tổng quát)};

  % Level 1: div
  \node[div_box, below=0.3cm of p2.south west, anchor=north west] (div1) {\codeinline{<div>}};
  \node[badge_gray, right=4pt of div1] {Trung gian};

  % Level 2: p3 (con của div1)
  \node[p_box, below right=0.4cm and 0.6cm of div1.south west, anchor=north west] (p3) {\codeinline{<p>} Đoạn văn 3 \codeinline{</p>}};
  \node[badge_blue, right=4pt of p3] {Hậu duệ (Cấp 2)};

  % Level 1: p4 (cùng cấp với div1, h1)
  \node[p_box, below=1.25cm of div1.south west, anchor=north west] (p4) {\codeinline{<p>} Đoạn văn 4 \codeinline{</p>}};
  \node[badge_purple, right=4pt of p4] {h1 \textasciitilde{} p (tổng quát)};

  % Connectors
  \draw[primaryBlue!70, line width=1pt] (root.south west) ++(0.3,0) |- (h1.west);
  \draw[primaryBlue!70, line width=1pt] (root.south west) ++(0.3,0) |- (p1.west);
  \draw[primaryBlue!70, line width=1pt] (root.south west) ++(0.3,0) |- (p2.west);
  \draw[primaryBlue!70, line width=1pt] (root.south west) ++(0.3,0) |- (div1.west);
  \draw[primaryBlue!70, line width=1pt] (root.south west) ++(0.3,0) |- (p4.west);
  \draw[accentOrange!80, line width=1pt] (div1.south west) ++(0.3,0) |- (p3.west);

\end{tikzpicture}
\end{center}
\end{conceptBox}

\subsubsection{5.3. Bảng Kết Quả Trực Quan -- Phân Tích Từng Phần Tử}

\begin{prettyTableBox}[Kết Quả Trực Quan: Màu Sắc Của Từng Phần Tử Khi Áp Dụng Cả 4 Selector]
\rowcolors{2}{white}{primaryBlue!4}
\sffamily\small
\begin{tabularx}{\linewidth}{|>{\raggedright\arraybackslash \bfseries}K{3.5cm}|>{\centering\arraybackslash}C{2.8cm}|Z|}
\hline
\rowcolor{primaryBlue} \thd{Nội dung} & \thd{Màu sắc} & \thd{Giải thích} \\
\hline
Tiêu đề chính & Mặc định (đen) & Không bị ảnh hưởng vì không có bộ chọn nào tác động đến \codeinline{<h1>}. \\
\hline
Đoạn văn 1 & \textcolor{tipBorder}{\textbf{Xanh lá (green)}} & Bộ chọn \codeinline{h1 + p} (Adjacent Sibling) tác động, vì nó là thẻ \codeinline{<p>} nằm \textbf{ngay liền sau} \codeinline{<h1>}. \\
\hline
Đoạn văn 2 & \textcolor{purple!70}{\textbf{Tím (purple)}} & Bộ chọn \codeinline{h1 \textasciitilde{} p} (General Sibling) tác động, vì nó nằm \textbf{sau} \codeinline{<h1>} và \textbf{cùng cấp}. \\
\hline
Đoạn văn 3 (trong \codeinline{<div>}) & \textcolor{primaryBlue}{\textbf{Xanh dương (blue)}} & Bộ chọn \codeinline{.container p} (Descendant) tác động, vì nó là \textbf{hậu duệ} của \codeinline{.container} (dù lồng cấp 2 trong \codeinline{<div>}). \\
\hline
Đoạn văn 4 & \textcolor{purple!70}{\textbf{Tím (purple)}} & Bộ chọn \codeinline{h1 \textasciitilde{} p} (General Sibling) tác động, vì nó nằm \textbf{sau} \codeinline{<h1>} và \textbf{cùng cấp}. \\
\hline
\end{tabularx}
\end{prettyTableBox}

\begin{browserWindow}[Kết Quả Hiển Thị Trực Quan Trên Trình Duyệt]
\sffamily\small
{\bfseries\Large Tiêu đề chính} \par\vspace{4pt}
\textcolor{tipBorder}{\textbf{Đoạn văn 1}} \textcolor{gray!60}{\scriptsize $\leftarrow$ h1 + p (green)} \par\vspace{2pt}
\textcolor{purple!70}{\textbf{Đoạn văn 2}} \textcolor{gray!60}{\scriptsize $\leftarrow$ h1 \textasciitilde{} p (purple)} \par\vspace{2pt}
\begin{tcolorbox}[arc=2pt, colback=gray!3, colframe=gray!25, boxrule=0.5pt, boxsep=2pt, left=10pt]
\textcolor{primaryBlue}{\textbf{Đoạn văn 3}} \textcolor{gray!60}{\scriptsize $\leftarrow$ .container p (blue) -- lồng trong <div>}
\end{tcolorbox}\vspace{2pt}
\textcolor{purple!70}{\textbf{Đoạn văn 4}} \textcolor{gray!60}{\scriptsize $\leftarrow$ h1 \textasciitilde{} p (purple)}
\end{browserWindow}

\begin{noteBox}[Tại Sao Đoạn Văn 1 Là Xanh Lá Chứ Không Phải Tím Hoặc Đỏ?]
Đoạn văn 1 thỏa mãn \textbf{cả 4 bộ chọn} cùng lúc:
\begin{enumerate}[leftmargin=1.5em, itemsep=2pt, parsep=0pt]
    \item \codeinline{.container p} (Descendant) -- vì P1 là hậu duệ của \codeinline{.container}.
    \item \codeinline{.container > p} (Child) -- vì P1 là con trực tiếp.
    \item \codeinline{h1 + p} (Adjacent Sibling) -- vì P1 liền kề ngay sau \codeinline{<h1>}.
    \item \codeinline{h1 \textasciitilde{} p} (General Sibling) -- vì P1 nằm sau \codeinline{<h1>} cùng cấp.
\end{enumerate}
Khi nhiều bộ chọn cùng tác động lên 1 phần tử, CSS sẽ áp dụng theo quy tắc \textbf{Cascade (thứ tự xuất hiện)}: bộ chọn được viết \textbf{sau cùng} trong file CSS sẽ thắng. Vì \codeinline{h1 + p} nằm sau \codeinline{.container > p}, nên màu xanh lá (green) được hiển thị. Nếu \codeinline{h1 \textasciitilde{} p} (purple) được viết cuối cùng, nó sẽ đè lên. Nhưng vì \codeinline{h1 + p} cụ thể hơn \codeinline{h1 \textasciitilde{} p} trong trường hợp liền kề, P1 vẫn nhận màu green!
\end{noteBox}

\subsubsection{5.4. Ứng Dụng Thực Tế Của Từng Bộ Chọn}

\begin{prettyTableBox}[Bảng Ứng Dụng Thực Tế Của 4 Bộ Chọn Kết Hợp]
\rowcolors{2}{white}{primaryBlue!4}
\sffamily\small
\begin{tabularx}{\linewidth}{|>{\raggedright\arraybackslash \bfseries\color{primaryBlue}}K{3.5cm}|Z|}
\hline
\rowcolor{primaryBlue} \thd{Bộ chọn} & \thd{Ứng dụng thực tế} \\
\hline
Descendant Selector \newline \codeinline{cha con} & Styling toàn bộ nội dung trong section: Định dạng tất cả tiêu đề, đoạn văn, nút\ldots{} bên trong một khu vực cụ thể. \\
\hline
Child Selector \newline \codeinline{cha > con} & Menu đa cấp: Chọn chỉ các mục cấp 1 của menu để định dạng đặc biệt, bỏ qua các mục lồng cấp 2, 3\ldots{} \\
\hline
Adjacent Sibling \newline \codeinline{e1 + e2} & Tạo khoảng cách sau tiêu đề: Định dạng đoạn văn ngay sau tiêu đề để cách dòng, tạo khoảng thở cho bố cục. \\
\hline
General Sibling \newline \codeinline{e1 \textasciitilde{} e2} & Styling danh sách thông báo: Tô màu hoặc highlight tất cả các thông báo sau tiêu đề ``Thông báo mới''. \\
\hline
\end{tabularx}
\end{prettyTableBox}

\subsubsection{5.5. Khi Nào Nên Dùng Bộ Chọn Nào?}

\begin{tipBox}[Hướng Dẫn Chọn Đúng Bộ Chọn Kết Hợp Cho Từng Tình Huống]
\begin{itemize}[leftmargin=1.5em, itemsep=4pt, parsep=0pt]
    \item Cần định dạng \textbf{toàn bộ phần tử con} (bất kể cấp độ) $\rightarrow$ \textbf{Descendant Selector} (\codeinline{cha con}).
    \item Chỉ muốn tác động đến \textbf{con trực tiếp}, tránh ảnh hưởng sâu $\rightarrow$ \textbf{Child Selector} (\codeinline{cha > con}).
    \item Muốn định dạng phần tử \textbf{liền kề} để tạo bố cục đẹp $\rightarrow$ \textbf{Adjacent Sibling Selector} (\codeinline{e1 + e2}).
    \item Muốn tác động đến \textbf{nhiều phần tử cùng cấp} mà không cần liền kề $\rightarrow$ \textbf{General Sibling Selector} (\codeinline{e1 \textasciitilde{} e2}).
\end{itemize}
\end{tipBox}

\subsubsection{5.6. Mẹo Tối Ưu CSS Với Bộ Chọn Kết Hợp}

\begin{warnBox}[3 Mẹo Tối Ưu Hiệu Suất CSS Khi Dùng Bộ Chọn Kết Hợp]
\begin{enumerate}[leftmargin=1.5em, itemsep=4pt, parsep=0pt]
    \item \textbf{Tránh lạm dụng bộ chọn quá sâu}: \codeinline{div ul li a span \{\}} có thể làm trình duyệt chậm vì phải tìm kiếm qua nhiều cấp. Hãy tối giản bộ chọn xuống tối đa 3--4 cấp.
    \item \textbf{Ưu tiên bộ chọn cụ thể hơn}: Child Selector (\codeinline{>}) nhanh hơn Descendant Selector vì chỉ tìm con trực tiếp, không phải quét toàn bộ cây DOM con.
    \item \textbf{Kết hợp thông minh}: Dùng kết hợp Sibling Selector với Pseudo-class (ví dụ: \codeinline{:hover}, \codeinline{:nth-child}) để tạo hiệu ứng đẹp mắt mà không cần JavaScript!
\end{enumerate}
\end{warnBox}

\subsubsection{5.7. Kết Luận Tổng Quan}

\begin{tipBox}[Kết Luận Tổng Hợp Về Bộ Chọn Kết Hợp (Combinator Selectors)]
Bộ chọn kết hợp là công cụ mạnh mẽ giúp CSS linh hoạt và chính xác hơn. Hiểu rõ từng loại, cách hoạt động và ứng dụng thực tế sẽ giúp bạn:
\begin{itemize}[leftmargin=1.5em, itemsep=3pt, parsep=0pt]
    \item \textbf{Viết CSS tinh gọn}: Chọn đúng phần tử mà không cần thêm class hoặc ID thừa vào HTML.
    \item \textbf{Tối ưu hiệu suất}: Hiểu sâu về các selector này giúp xây dựng giao diện gọn gàng, dễ bảo trì và tối ưu tốc độ render trang!
    \item \textbf{Tạo hiệu ứng không cần JavaScript}: Kết hợp Sibling Selectors với Pseudo-class để tạo menu dropdown, toggle panel, highlight nhóm phần tử\ldots{} hoàn toàn bằng CSS thuần!
\end{itemize}
\end{tipBox}


%% ===================================================================
%% MỤC MỚI: CHUYÊN ĐỀ BỘ CHỌN THUỘC TÍNH (ATTRIBUTE SELECTORS MASTERCLASS)
%% ===================================================================

\subsection{6. Chuyên Đề Phân Tích Chi Tiết \& Thực Nghiệm Về Bộ Chọn Thuộc Tính (Attribute Selectors Masterclass)}

\begin{conceptBox}[Khái Niệm Về Bộ Chọn Thuộc Tính (Attribute Selectors)]
\textbf{Bộ chọn thuộc tính (\term{Attribute Selector})} trong CSS cho phép bạn chọn và áp dụng định kiểu cho các phần tử HTML dựa trên sự \textbf{tồn tại của một thuộc tính} hoặc \textbf{giá trị cụ thể của thuộc tính đó} (ví dụ: \codeinline{href}, \codeinline{type}, \codeinline{src}, \codeinline{title}, \codeinline{target}, \codeinline{disabled}, \codeinline{required}, các thuộc tính tùy chỉnh \codeinline{data-*}, hay thuộc tính trợ năng \codeinline{aria-*}).
\par\vspace{3pt}
\textbf{Lợi ích cốt lõi}: Giúp bạn định dạng giao diện linh hoạt trực tiếp từ ngữ nghĩa HTML mà không cần phải gán thêm thẻ \codeinline{class} hay \codeinline{id} dư thừa vào mã nguồn!
\end{conceptBox}

\paragraph{a. Bảng tổng quan 7 loại bộ chọn thuộc tính \& Cờ phân biệt hoa/thường:}

\begin{prettyTableBox}[Bảng Tổng Quan Chi Tiết 7 Loại Bộ Chọn Thuộc Tính Trong CSS]
\rowcolors{2}{white}{primaryBlue!4}
\sffamily\small
\begin{tabularx}{\linewidth}{|>{\raggedright\arraybackslash \ttfamily\bfseries\color{primaryBlue}}K{2.8cm}|>{\ttfamily\color{accentOrange}\bfseries}K{2.8cm}|Z|}
\hline
\rowcolor{primaryBlue} \thd{Tên Bộ Chọn} & \thd{Cú Pháp} & \thd{Mô Tả Quy Tắc So Khớp} \\
\hline
Tồn tại thuộc tính & [attr] & Chọn phần tử \textbf{có chứa thuộc tính} \codeinline{attr}, bất kể giá trị là gì. \\
\hline
Giá trị chính xác & [attr="val"] & Chọn phần tử có giá trị thuộc tính \textbf{bằng đúng tuyệt đối} \codeinline{val}. \\
\hline
Chứa từ độc lập & [attr\textasciitilde{}="val"] & Giá trị thuộc tính chứa \codeinline{val} như \textbf{một từ độc lập} (phân cách bởi khoảng trắng). \\
\hline
Bắt đầu bằng tiền tố & [attr|="val"] & Giá trị bằng đúng \codeinline{val} hoặc \textbf{bắt đầu bằng \codeinline{val-}} (thường dùng cho mã ngôn ngữ). \\
\hline
Bắt đầu bằng chuỗi & [attr\textasciicircum{}="val"] & Giá trị thuộc tính \textbf{bắt đầu bằng} chuỗi \codeinline{val}. \\
\hline
Kết thúc bằng chuỗi & [attr\$="val"] & Giá trị thuộc tính \textbf{kết thúc bằng} chuỗi \codeinline{val}. \\
\hline
Chứa chuỗi con & [attr*="val"] & Giá trị thuộc tính \textbf{chứa chuỗi con \codeinline{val}} ở bất kỳ vị trí nào. \\
\hline
Không phân biệt hoa/thường & [attr="val" i] & Cờ \codeinline{i} giúp quy tắc so khớp \textbf{không phân biệt chữ hoa / chữ thường}. \\
\hline
\end{tabularx}
\end{prettyTableBox}

\paragraph{b. Phân tích chi tiết 7 loại bộ chọn thuộc tính kèm ví dụ thực tế:}

\subsubsection{6.1. Bộ Chọn Sự Tồn Tại Của Thuộc Tính ([attr])}

\begin{conceptBox}[{Ý Nghĩa Selector [attr]}]
Bộ chọn \codeinline{[attr]} sẽ khớp với tất cả các phần tử có chứa thuộc tính được chỉ định, không quan tâm giá trị của thuộc tính đó là gì hay có rỗng hay không.
\end{conceptBox}

\begin{codeWindow}[Ví Dụ 1: Định Dạng Thẻ Input Có Thuộc Tính required]
\begin{lstlisting}[language=CSS3]
/* Sample Text */
input[required] {
    border: 2px solid red;
    background-color: #fff0f0;
}
\end{lstlisting}

\begin{lstlisting}[language=HTML5]
<form>
    <input type="text" placeholder="Sample Text">
    <input type="email" required placeholder="Sample Text"> <!-- Sample Text -->
</form>
\end{lstlisting}
\end{codeWindow}

\begin{browserWindow}[Kết Quả Hiển Thị: Ô Email Bắt Buộc Được Viền Đỏ]
\sffamily\small
\begin{tcolorbox}[arc=2pt, colback=white, colframe=gray!30, boxrule=1pt, width=7.5cm]
\textcolor{gray!70}{Họ và tên (Tùy chọn)} \par\vspace{4pt}
\begin{tcolorbox}[arc=2pt, colback=red!5, colframe=red!70, boxrule=1.5pt, boxsep=2pt]
\textcolor{red!80!black}{\textbf{Email (Bắt buộc)}} \hfill \textcolor{red!70}{\scriptsize $\leftarrow$ input[required] (Viền đỏ)}
\end{tcolorbox}
\end{tcolorbox}
\end{browserWindow}

\subsubsection{6.2. Bộ Chọn Giá Trị Chính Xác ([attr="val"])}

\begin{conceptBox}[{Ý Nghĩa Selector [attr="val"]}]
So khớp chính xác từng ký tự của giá trị thuộc tính. Chỉ khi giá trị khớp 100\% thì CSS mới được áp dụng.
\end{conceptBox}

\begin{codeWindow}[Ví Dụ 2: Định Kiểu Nút Bấm Submit Trong Form]
\begin{lstlisting}[language=CSS3]
/* Submit */
input[type="submit"] {
    background-color: #28a745;
    color: white;
    font-weight: bold;
    border-radius: 4pt;
    padding: 6px 12px;
}
\end{lstlisting}

\begin{lstlisting}[language=HTML5]
<input type="text" placeholder="Sample Text">
<input type="submit" value="Submit"> <!-- Sample Text -->
\end{lstlisting}
\end{codeWindow}

\begin{browserWindow}[Kết Quả Hiển Thị: Nút Submit Được Tô Màu Xanh Lá]
\sffamily\small
\tcbox[on line, arc=3pt, colback=white, colframe=gray!40, boxsep=3pt]{\color{gray!70}Nhập tên...} \quad
\tcbox[on line, arc=3pt, colback=tipBorder, colframe=tipBorder, boxsep=4pt, left=8pt, right=8pt]{\color{white}\bfseries Gửi Dữ Liệu} \textcolor{gray!60}{\scriptsize $\leftarrow$ input[type="submit"]}
\end{browserWindow}

\subsubsection{6.3. Bộ Chọn Chứa Từ Độc Lập ([attr\textasciitilde{}="val"])}

\begin{conceptBox}[{Ý Nghĩa Selector [attr\textasciitilde{}="val"]}]
Bộ chọn này tìm giá trị thuộc tính chứa từ \codeinline{val} như \textbf{một từ riêng biệt hoàn chỉnh} (được phân cách bởi khoảng trắng xung quanh). Thường dùng khi một thuộc tính chứa danh sách nhiều từ (như \codeinline{class="btn primary active"}).
\end{conceptBox}

\begin{codeWindow}[Ví Dụ 3: Tìm Thẻ Có Class Chứa Từ "active"]
\begin{lstlisting}[language=CSS3]
/* Select element */
[class~="active"] {
    color: #007bff;
    font-weight: bold;
}
\end{lstlisting}

\begin{lstlisting}[language=HTML5]
<div class="card active main">Sample Text</div>
<div class="card interactive">Sample Text</div>
\end{lstlisting}
\end{codeWindow}

\begin{browserWindow}[Kết Quả Hiển Thị: Thẻ Chứa Từ Độc Lập "active" Được Tô Màu]
\sffamily\small
\tcbox[on line, arc=3pt, colback=primaryBlue!10, colframe=primaryBlue, boxsep=3pt]{\textcolor{primaryBlue}{\textbf{Khớp - Được chọn (Class chứa từ "active")}}} \quad
\tcbox[on line, arc=3pt, colback=gray!5, colframe=gray!30, boxsep=3pt]{\textcolor{darkText}{Không khớp (Class interactive)}}
\end{browserWindow}

\subsubsection{6.4. Bộ Chọn Bắt Đầu Tiền Tố / Mã Ngôn Ngữ ([attr|="val"])}

\begin{conceptBox}[{Ý Nghĩa Selector [attr|="val"]}]
Khớp với các phần tử có giá trị thuộc tính \textbf{chính xác là \codeinline{val}} hoặc \textbf{bắt đầu bằng \codeinline{val-}} (theo sau là dấu gạch nối). Bộ chọn này sinh ra đặc biệt để chọn mã ngôn ngữ chuẩn ISO (như \codeinline{lang="en"}, \codeinline{lang="en-US"}, \codeinline{lang="en-GB"}).
\end{conceptBox}

\begin{codeWindow}[Ví Dụ 4: Định Dạng Thẻ Có Thuộc Tính Ngôn Ngữ Tiếng Anh]
\begin{lstlisting}[language=CSS3]
/* Sample Text */
p[lang|="en"] {
    font-style: italic;
    color: #4a5568;
}
\end{lstlisting}

\begin{lstlisting}[language=HTML5]
<p lang="en">Hello World</p>       <!-- Sample Text -->
<p lang="en-US">Hello America</p> <!-- Sample Text -->
<p lang="fr">Bonjour</p>          <!-- Sample Text -->
\end{lstlisting}
\end{codeWindow}

\begin{browserWindow}[Kết Quả Hiển Thị: Định Dạng Nghiêng Cho Thẻ Thuộc Nhóm Tiếng Anh]
\sffamily\small
\textit{\textcolor{gray!80}{\textbf{Hello World}}} \textcolor{gray!60}{\scriptsize $\leftarrow$ lang="en" (In nghiêng)} \par\vspace{2pt}
\textit{\textcolor{gray!80}{\textbf{Hello America}}} \textcolor{gray!60}{\scriptsize $\leftarrow$ lang="en-US" (In nghiêng)} \par\vspace{2pt}
Bonjour \textcolor{gray!60}{\scriptsize $\leftarrow$ lang="fr" (Mặc định)}
\end{browserWindow}

\subsubsection{6.5. Bộ Chọn Bắt Đầu Chuỗi ([attr\textasciicircum{}="val"])}

\begin{conceptBox}[{Ý Nghĩa Selector [attr\textasciicircum{}="val"]}]
Khớp với tất cả các phần tử có giá trị thuộc tính \textbf{bắt đầu bằng chuỗi \codeinline{val}} (Prefix Match).
\end{conceptBox}

\begin{codeWindow}[Ví Dụ 5: Tự Động Thêm Khung Cho Đường Dẫn Bảo Mật HTTPS]
\begin{lstlisting}[language=CSS3]
/* Sample Text */
a[href^="https://"] {
    color: #2b6cb0;
    font-weight: bold;
}
\end{lstlisting}

\begin{lstlisting}[language=HTML5]
<a href="https://google.com">Sample Text</a> <!-- Sample Text -->
<a href="http://unsafe.com">Sample Text</a> <!-- Sample Text -->
\end{lstlisting}
\end{codeWindow}

\begin{browserWindow}[Kết Quả Hiển Thị: Link HTTPS Tự Động Được Nổi Bật]
\sffamily\small
\textcolor{primaryBlue}{\textbf{\underline{Trang an toàn HTTPS}}} \textcolor{gray!60}{\scriptsize $\leftarrow$ a[href\textasciicircum{}="https://"]} \quad | \quad
\textcolor{darkText}{\underline{Trang không an toàn HTTP}}
\end{browserWindow}

\subsubsection{6.6. Bộ Chọn Kết Thúc Chuỗi ([attr\$="val"])}

\begin{conceptBox}[{Ý Nghĩa Selector [attr\$="val"]}]
Khớp với tất cả các phần tử có giá trị thuộc tính \textbf{kết thúc bằng chuỗi \codeinline{val}} (Suffix Match). Ứng dụng phổ biến nhất là nhận diện đuôi mở rộng của tệp tin tải về (\codeinline{.pdf}, \codeinline{.docx}, \codeinline{.zip}).
\end{conceptBox}

\begin{codeWindow}[Ví Dụ 6: Tự Động Nhận Diện Tệp Tải Về PDF \& Thêm Icon]
\begin{lstlisting}[language=CSS3]
/* Sample Text */
a[href$=".pdf"]::after {
    content: " [PDF]";
    color: red;
    font-weight: bold;
    font-size: 0.8em;
}
\end{lstlisting}

\begin{lstlisting}[language=HTML5]
<a href="tailieu.pdf">Sample Text</a> <!-- Sample Text -->
<a href="trangchu.html">Sample Text</a>          <!-- Sample Text -->
\end{lstlisting}
\end{codeWindow}

\begin{browserWindow}[Kết Quả Hiển Thị: Tự Động Chèn Nhãn [PDF] Đỏ Mắt Nhìn]
\sffamily\small
\textcolor{primaryBlue}{\underline{Tải Báo Cáo Tài Chính}} \textcolor{red}{\textbf{[PDF]}} \textcolor{gray!60}{\scriptsize $\leftarrow$ a[href\$=".pdf"]::after}
\end{browserWindow}

\subsubsection{6.7. Bộ Chọn Chứa Chuỗi Con ([attr*="val"])}

\begin{conceptBox}[{Ý Nghĩa Selector [attr*="val"]}]
Khớp với các phần tử có giá trị thuộc tính \textbf{chứa chuỗi \codeinline{val} ở bất kỳ vị trí nào} (Substring Match).
\end{conceptBox}

\begin{codeWindow}[Ví Dụ 7: Định Kiểu Cho Tất Cả Đường Dẫn Đến Mạng Xã Hội Facebook]
\begin{lstlisting}[language=CSS3]
/* Select element */
a[href*="facebook"] {
    color: #1877f2;
    font-weight: bold;
}
\end{lstlisting}

\begin{lstlisting}[language=HTML5]
<a href="https://www.facebook.com/profile">Sample Text</a> <!-- Sample Text -->
\end{lstlisting}
\end{codeWindow}

\begin{browserWindow}[Kết Quả Hiển Thị: Link Facebook Tự Động Được Tô Màu Xanh Nhận Diện]
\sffamily\small
\textcolor{primaryBlue}{\textbf{\underline{Trang Fanpage}}} \textcolor{gray!60}{\scriptsize $\leftarrow$ a[href*="facebook"] (Tô màu xanh Facebook)}
\end{browserWindow}

\subsubsection{6.8. Cờ Phân Biệt Chữ Hoa / Chữ Thường ([attr="val" i])}

\begin{conceptBox}[{Ý Nghĩa Cờ Case-Insensitive Flag "i"}]
Mặc định trong HTML5, một số thuộc tính phân biệt hoa thường. Thêm cờ \codeinline{i} trước dấu đóng ngoặc \codeinline{]} giúp bộ chọn \textbf{bỏ qua phân biệt hoa/thường}, giúp khớp cả \codeinline{.pdf}, \codeinline{.PDF}, \codeinline{.Pdf}.
\end{conceptBox}

\begin{codeWindow}[Ví Dụ 8: Bắt File PDF Không Phân Biệt Chữ Hoa/Thường]
\begin{lstlisting}[language=CSS3]
/* Sample Text */
a[href$=".pdf" i] {
    color: red;
}
\end{lstlisting}
\end{codeWindow}

\begin{browserWindow}[Kết Quả Hiển Thị: Bắt Cả .PDF và .Pdf Nhờ Cờ "i"]
\sffamily\small
\textcolor{red}{\textbf{\underline{tailieu\_baocao.PDF}}} \textcolor{red}{\scriptsize $\leftarrow$ Khớp nhờ cờ "i"} \par\vspace{2pt}
\textcolor{red}{\textbf{\underline{huongdan\_sudung.Pdf}}} \textcolor{red}{\scriptsize $\leftarrow$ Khớp nhờ cờ "i"}
\end{browserWindow}

\paragraph{c. Bảng tra cứu tóm tắt các ký hiệu so khớp Attribute Selectors:}

\begin{prettyTableBox}[Bảng Tra Cứu So Sánh Các Ký Hiệu Attribute Selectors]
\rowcolors{2}{white}{primaryBlue!4}
\sffamily\small
\begin{tabularx}{\linewidth}{|>{\raggedright\arraybackslash \ttfamily\bfseries\color{accentOrange}}C{1.8cm}|>{\bfseries\color{primaryBlue}}K{3.0cm}|Z|Z|}
\hline
\rowcolor{primaryBlue} \thd{Ký Hiệu} & \thd{Loại So Khớp} & \thd{Trường Hợp KHỚP (Match)} & \thd{Trường Hợp BỎ QUA (No match)} \\
\hline
[attr] & Tồn tại & \codeinline{<div title="hello">} & \codeinline{<div>} (không có title) \\
\hline
= & Bằng đúng tuyệt đối & \codeinline{<a target="\_blank">} & \codeinline{<a target="\_self">} \\
\hline
\textasciitilde{}= & Từ độc lập & \codeinline{<div class="btn active">} & \codeinline{<div class="btn-active">} \\
\hline
|= & Tiền tố ngôn ngữ & \codeinline{<p lang="en-US">} & \codeinline{<p lang="us-en">} \\
\hline
\textasciicircum{}= & Bắt đầu chuỗi & \codeinline{<a href="https://a.com">} & \codeinline{<a href="http://a.com">} \\
\hline
\$= & Kết thúc chuỗi & \codeinline{<a href="file.pdf">} & \codeinline{<a href="file.pdf.png">} \\
\hline
*= & Chứa chuỗi con & \codeinline{<img src="user\_avatar.jpg">} & \codeinline{<img src="user\_photo.jpg">} \\
\hline
\end{tabularx}
\end{prettyTableBox}

\paragraph{d. 4 Ứng dụng thực tế đỉnh cao của Attribute Selectors trong Frontend Development:}

\begin{prettyTableBox}[4 Pattern Thực Tế Sử Dụng Attribute Selectors Vô Cùng Hiệu Quả]
\rowcolors{2}{white}{primaryBlue!4}
\sffamily\small
\begin{tabularx}{\linewidth}{|>{\raggedright\arraybackslash \bfseries\color{primaryBlue}}K{3.8cm}|Z|}
\hline
\rowcolor{primaryBlue} \thd{Mẫu Thiết Kế (UI Pattern)} & \thd{Mô Tả \& Cách Triển Khai Thực Tế} \\
\hline
1. Form Validation \& UI States & Định kiểu trạng thái nhập liệu trực tiếp từ HTML5 attributes: \newline
\codeinline{input[disabled]} $\rightarrow$ Nền xám nhạt, cấm trỏ chuột (\codeinline{cursor: not-allowed}). \newline
\codeinline{input[readonly]} $\rightarrow$ Không cho chỉnh sửa nhưng vẫn chọn được text. \newline
\codeinline{input:invalid} $\rightarrow$ Hiển thị khung cảnh báo viền đỏ khi nhập sai format. \\
\hline
2. Tự Động Nhận Diện Link Ngoại Mạng \& File Tải Về & Tăng trải nghiệm người dùng (UX) tự động: \newline
\codeinline{a[target="\_blank"]::after} $\rightarrow$ Tự chèn icon mũi tên chỉ ra ngoài. \newline
\codeinline{a[href\$=".zip"]} $\rightarrow$ Tự thêm icon tập tin nén ZIP. \\
\hline
3. Custom Data Attributes (\codeinline{data-*}) & Quản lý State giao diện trong React / Vue / Vanilla JS: \newline
\codeinline{div[data-status="success"]} $\rightarrow$ Tô màu xanh thông báo thành công. \newline
\codeinline{body[data-theme="dark"]} $\rightarrow$ Tự động chuyển đổi toàn bộ màu giao diện Dark Mode! \\
\hline
4. Web Accessibility ARIA Attributes (\codeinline{aria-*}) & Đảm bảo trợ năng chuẩn WCAG cho người khuyết tật: \newline
\codeinline{button[aria-disabled="true"]} $\rightarrow$ Giảm độ đục opacity, vô hiệu hóa nút bấm chuẩn a11y. \newline
\codeinline{div[aria-expanded="true"]} $\rightarrow$ Hiển thị menu accordion đang mở. \\
\hline
\end{tabularx}
\end{prettyTableBox}

\begin{tipBox}[Best Practices Checklist Cho Attribute Selectors]
\begin{itemize}[leftmargin=1.5em, itemsep=3pt, parsep=0pt]
    \item \textbf{Khi nào nên dùng Attribute Selector?} Dùng khi thuộc tính đó mang ý nghĩa HTML gốc (như \codeinline{type}, \codeinline{href}, \codeinline{disabled}, \codeinline{data-*}, \codeinline{aria-*}) để tránh phải viết thêm class thừa.
    \item \textbf{Khi nào nên dùng Class Selector?} Dùng Class cho các thành phần giao diện tái sử dụng linh hoạt nhiều nơi (như \codeinline{.card}, \codeinline{.btn-primary}).
    \item \textbf{Lưu ý hiệu năng}: Attribute selector rất nhanh trên các trình duyệt hiện đại, nhưng hãy hạn chế dùng selector chứa chuỗi quá phức tạp \codeinline{[attr*="..."]} trên cây DOM quá lớn hàng nghìn node.
\end{itemize}
\end{tipBox}

\subsection{7. Chuyên Đề Phân Tích Chi Tiết \& Thực Nghiệm Về Bộ Chọn Lớp Giả (Pseudo-class Selectors Masterclass)}

Pseudo-class selectors (Bộ chọn lớp giả) là một phần cực kỳ mạnh mẽ trong CSS3 giúp bạn chọn và định kiểu các phần tử HTML dựa trên \textbf{trạng thái tương tác đặc biệt} hoặc \textbf{vị trí vị thế cấu trúc} của chúng trong cây DOM -- mà \textbf{không cần phải thêm bất kỳ class hay ID thủ công nào vào file HTML}!

\begin{conceptBox}[{Bản Chất \& Cú Pháp Tổng Quát Về Pseudo-Class}]
\begin{itemize}[leftmargin=1.5em, itemsep=4pt, parsep=0pt]
    \item \textbf{Bản chất}: Chọn và định kiểu phần tử dựa trên \textbf{trạng thái tương tác đặc biệt} (như di chuột, click bấm) hoặc \textbf{vị trí vị thế cấu trúc} trong cây DOM mà \textbf{không cần thêm class/ID vào HTML}.
    \item \textbf{Cú pháp chuẩn}: Bắt đầu bằng \textbf{một dấu hai chấm} (\codeinline{:}).
    \item \textbf{Công thức tổng quát}:
    \par\vspace{3pt}
    \tcbox[on line, arc=3pt, colback=primaryBlue!6, colframe=primaryBlue!40, boxrule=0.8pt, boxsep=3pt, left=8pt, right=8pt]{%
        \ttfamily\color{primaryBlue}\bfseries selector:pseudo-class \{ property: value; \}%
    }
    \par\vspace{3pt}
    \item \textbf{Ví dụ minh họa}: \codeinline{button:hover \{ color: red; \}} (Đổi chữ sang màu đỏ khi rê chuột vào nút bấm).
    \item \textbf{Ứng dụng thực tế}: Phản hồi hành vi rê chuột, con trỏ focus nhập liệu, trạng thái checkbox/radio được chọn, hoặc tự động định kiểu dòng chẵn/lẻ cho bảng dữ liệu.
\end{itemize}
\end{conceptBox}

\subsubsection{7.1. Bảng Tra Cứu Tổng Hợp Toàn Bộ Pseudo-Classes}

Dưới đây là phân tích chi tiết toàn bộ các bộ chọn lớp giả (Pseudo-classes) được nhóm theo 4 phân loại chức năng thực tế trong thiết kế giao diện Web:

\subsubsection{7.1. Nhóm Trạng Thái Tương Tác \& Điều Hướng (User Action \& Links)}

Nhóm này kiểm soát các hành vi động của người dùng khi di chuột, click bấm, focus ô nhập liệu hoặc truy cập đường dẫn:

\begin{prettyTableBox}[Bảng Tra Cứu Pseudo-Classes Trạng Thái Tương Tác \& Liên Kết]
\rowcolors{2}{white}{primaryBlue!4}
\sffamily\small
\begin{tabularx}{\linewidth}{|>{\raggedright\arraybackslash \ttfamily\color{primaryBlue}\bfseries}K{2.2cm}|>{\sffamily\small\bfseries}K{2.8cm}|>{\ttfamily\small}K{3.6cm}|Z|}
\hline
\rowcolor{primaryBlue} \thd{\sffamily Selector} & \thd{\sffamily Phạm vi áp dụng} & \thd{\sffamily Ví dụ cú pháp} & \thd{\sffamily Chức năng \& Ý nghĩa} \\
\hline
:hover & Tất cả phần tử & button:hover \{ color: red; \} & Kích hoạt đổi màu đỏ khi trỏ chuột vào phần tử. \\
\hline
:active & Tất cả phần tử & a:active \{ color: green; \} & Kích hoạt chuyển xanh khi đang giữ chuột bấm xuống. \\
\hline
:focus & Input, Button, Link & input:focus \{ border: 2px solid blue; \} & Kích hoạt viền xanh khi nhận tiêu điểm nhập liệu. \\
\hline
:link & Thẻ \codeinline{<a href="...">} & a:link \{ color: blue; \} & Áp dụng cho liên kết chưa từng được truy cập. \\
\hline
:visited & Thẻ \codeinline{<a href="...">} & a:visited \{ color: purple; \} & Áp dụng cho liên kết đã từng lưu trong lịch sử web. \\
\hline
:target & Phần tử có \codeinline{id} & \#section:target \{ background: yellow; \} & Khớp với id trên URL anchor (\codeinline{\#section}) đổi nền vàng. \\
\hline
\end{tabularx}
\end{prettyTableBox}

\paragraph{a. Bộ Chọn :link (Liên Kết Chưa Truy Cập)}

Bộ chọn \codeinline{:link} trong CSS được sử dụng để chọn các liên kết (\codeinline{<a>}) chưa được truy cập. Điều này có nghĩa là liên kết đó chưa được người dùng click vào hoặc trình duyệt chưa ghi nhận nó là đã truy cập.

\begin{conceptBox}[Đặc Điểm \& Cú Pháp Của :link]
\begin{itemize}[leftmargin=1.5em, itemsep=3pt, parsep=0pt]
    \item \textbf{Cú pháp}:
    \begin{codeWindow}[Mã CSS Định Kiểu Liên Kết Chưa Truy Cập]
\begin{lstlisting}[language=CSS3]
a:link {
    color: blue;
    text-decoration: none;
}
\end{lstlisting}
    \end{codeWindow}
    \item \textbf{Ý nghĩa}:
    \begin{itemize}[leftmargin=1.2em, itemsep=2pt]
        \item \codeinline{a:link} chỉ áp dụng cho các thẻ \codeinline{<a>} có thuộc tính \codeinline{href}.
        \item Liên kết chưa truy cập sẽ được áp dụng style này.
    \end{itemize}
\end{itemize}
\end{conceptBox}

\begin{codeWindow}[Ví Dụ Minh Họa Đầy Đủ HTML \& CSS Cho :link Và :visited]
\begin{lstlisting}[language=HTML5]
<!DOCTYPE html>
<html lang="vi">
<head>
    <meta charset="UTF-8">
    <meta name="viewport" content="width=device-width, initial-scale=1.0">
    <title>Sample Text</title>
    <style>
        a:link {
            color: blue;
            font-weight: bold;
            text-decoration: none;
        }
        a:visited {
            color: purple;
        }
    </style>
</head>
<body>
    <h2>Sample Text</h2>
    <a href="https://www.example.com">Sample Text</a>  
    <br>
    <a href="https://www.google.com">Sample Text</a>
</body>
</html>
\end{lstlisting}
\end{codeWindow}

\begin{browserWindow}[Kết Quả Hiển Thị Trực Quan Trạng Thái :link Và :visited]
\sffamily\small
\begin{itemize}[leftmargin=1.5em, itemsep=3pt, parsep=0pt]
    \item \textbf{Liên kết chưa truy cập (\codeinline{a:link})}:
    \par\vspace{2pt}
    \tcbox[on line, arc=3pt, colback=blue!6, colframe=blue!40, boxrule=0.8pt, boxsep=3pt]{%
        \small\sffamily\bfseries\color{blue} Link chưa truy cập%
    }
    \quad (Hiển thị chữ in đậm, màu xanh dương rực rỡ).
    \par\vspace{4pt}
    \item \textbf{Liên kết đã từng nhấp truy cập (\codeinline{a:visited})}:
    \par\vspace{2pt}
    \tcbox[on line, arc=3pt, colback=purple!6, colframe=purple!40, boxrule=0.8pt, boxsep=3pt]{%
        \small\sffamily\bfseries\color{purple} Link đã truy cập%
    }
    \quad (Trình duyệt tự động đổi sang màu tím dịu mắt dựa trên lịch sử duyệt web).
\end{itemize}
\end{browserWindow}

\begin{tipBox}[Ứng Dụng Thực Tế Của :link]
\begin{itemize}[leftmargin=1.5em, itemsep=2pt, parsep=0pt]
    \item Tạo giao diện rõ ràng cho người dùng biết link nào đã hoặc chưa click.
    \item Giúp người dùng dễ dàng điều hướng và tránh nhầm lẫn khi duyệt trang web.
\end{itemize}
\end{tipBox}

\paragraph{b. Bộ Chọn :visited (Liên Kết Đã Truy Cập)}

Trong CSS, \codeinline{a:visited} là một pseudo-class được sử dụng để chọn các liên kết (thẻ \codeinline{<a>}) đã được người dùng truy cập. Điều này giúp bạn tùy chỉnh giao diện của liên kết sau khi người dùng đã nhấp vào và truy cập trang đích.

\begin{conceptBox}[Cú Pháp \& Giải Thích :visited]
\begin{codeWindow}[Cú Pháp Định Kiểu :visited]
\begin{lstlisting}[language=CSS3]
a:visited {
    color: purple; /* Sample Text */
    text-decoration: underline; /* Sample Text */
}
\end{lstlisting}
\end{codeWindow}
\textbf{Giải thích}:
\begin{itemize}[leftmargin=1.5em, itemsep=2pt, parsep=0pt]
    \item \codeinline{a}: Chọn thẻ liên kết (anchor).
    \item \codeinline{:visited}: Chọn các liên kết đã được nhấp vào và truy cập trước đó (dựa trên lịch sử trình duyệt).
\end{itemize}
\end{conceptBox}

\begin{warnBox}[{Lưu Ý Bảo Mật Cực Kỳ Quan Trọng Cho a:visited}]
Vì lý do bảo mật, trình duyệt chỉ cho phép bạn thay đổi các thuộc tính không ảnh hưởng đến bố cục khi sử dụng \codeinline{:visited} để tránh cuộc tấn công lộ lịch sử duyệt web (\term{History Sniffing Attack}):
\begin{itemize}[leftmargin=1.5em, itemsep=2pt, parsep=0pt]
    \item \textbf{Các thuộc tính ĐƯỢC CHO PHÉP}: \codeinline{color}, \codeinline{background-color}, \codeinline{border-color}, \codeinline{outline-color}, \codeinline{text-decoration}, \codeinline{cursor}, \codeinline{fill}, \codeinline{stroke}.
    \item \textbf{Các thuộc tính BỊ CẤM HOÀN TOÀN}: Các thuộc tính làm thay đổi kích thước hoặc bố cục như \codeinline{display}, \codeinline{visibility}, \codeinline{content}, \codeinline{font-size}, \codeinline{padding}, \codeinline{margin} sẽ không có tác dụng với \codeinline{:visited}.
\end{itemize}
\end{warnBox}

\begin{codeWindow}[Ví Dụ Hoàn Chỉnh Kết Hợp a, a:hover Và a:visited]
\begin{lstlisting}[language=CSS3]
a {
    color: blue; /* Sample Text */
    text-decoration: none;
}
a:hover {
    color: red; /* Sample Text */
    text-decoration: underline;
}
a:visited {
    color: purple; /* Sample Text */
}
\end{lstlisting}
\end{codeWindow}

\begin{browserWindow}[Kết Quả Hiển Thị Trực Quan Chu Trình Chuyển Đổi Trạng Thái Link]
\sffamily\small
\begin{itemize}[leftmargin=1.5em, itemsep=3pt, parsep=0pt]
    \item \textbf{Trạng thái 1 -- Ban đầu (Chưa click)}:
    \par\vspace{2pt}
    \tcbox[on line, arc=3pt, colback=blue!6, colframe=blue!40, boxrule=0.8pt, boxsep=3pt]{%
        \small\sffamily\color{blue} Trang chủ Website%
    }
    \quad (Chữ màu xanh dương mặc định).
    \par\vspace{4pt}
    \item \textbf{Trạng thái 2 -- Di chuột vào (\codeinline{a:hover})}:
    \par\vspace{2pt}
    \tcbox[on line, arc=3pt, colback=red!6, colframe=red!40, boxrule=0.8pt, boxsep=3pt]{%
        \small\sffamily\bfseries\color{red}\underline{Trang chủ Website}%
    }
    \quad (Lập tức đổi sang màu đỏ rực kèm đường gạch chân).
    \par\vspace{4pt}
    \item \textbf{Trạng thái 3 -- Sau khi đã nhấp vào xem (\codeinline{a:visited})}:
    \par\vspace{2pt}
    \tcbox[on line, arc=3pt, colback=purple!6, colframe=purple!40, boxrule=0.8pt, boxsep=3pt]{%
        \small\sffamily\color{purple} Trang chủ Website%
    }
    \quad (Đổi sang màu tím đại diện cho liên kết trong quá khứ).
\end{itemize}
\end{browserWindow}

\paragraph{c. Bộ Chọn :hover (Trạng Thái Di Chuột Vào Phần Tử)}

Pseudo-class \codeinline{:hover} là một công cụ cực kỳ phổ biến và hữu ích trong CSS! Nó giúp bạn thay đổi giao diện của phần tử khi người dùng di chuột vào, tạo ra trải nghiệm tương tác trực quan và sinh động hơn.

\begin{conceptBox}[1. :hover Là Gì \& Cú Pháp Tổng Quát]
\textbf{:hover là gì?} \codeinline{:hover} là một pseudo-class trong CSS, áp dụng cho một phần tử khi người dùng di chuột qua nó (không cần nhấp chuột, chỉ cần đưa chuột vào).
\par\vspace{4pt}
\textbf{Cú pháp}:
\begin{codeWindow}[Cú Pháp Khai Báo :hover]
\begin{lstlisting}[language=CSS3]
/* selector:hover { Sample Text: Sample Text; } */
selector:hover {
    property: value;
}
\end{lstlisting}
\end{codeWindow}
\textbf{Ý nghĩa}: Khi con trỏ chuột nằm trên phần tử được chọn, các thuộc tính CSS bên trong \codeinline{:hover} sẽ được kích hoạt.
\end{conceptBox}

\begin{codeWindow}[2. Ví Dụ Cơ Bản HTML \& CSS Với :hover]
\begin{lstlisting}[language=HTML5]
<!DOCTYPE html>
<html lang="vi">
<head>
    <meta charset="UTF-8">
    <meta name="viewport" content="width=device-width, initial-scale=1.0">
    <title>Sample Text</title>
    <style>
        a:hover {
            color: red;
            text-decoration: underline;
        }
        button:hover {
            background-color: #007bff;
            color: #fff;
            cursor: pointer;
        }
    </style>
</head>
<body>
    <h2>Sample Text</h2>
    <a href="#">Sample Text</a>
    <br><br>
    <button>Sample Text</button>
</body>
</html>
\end{lstlisting}
\end{codeWindow}

\begin{browserWindow}[Kết Quả Hiển Thị Trực Quan Trạng Thái Di Chuột :hover]
\sffamily\small
\begin{itemize}[leftmargin=1.5em, itemsep=3pt, parsep=0pt]
    \item \textbf{Khi rê con trỏ chuột vào liên kết <a>}:
    \par\vspace{2pt}
    \tcbox[on line, arc=3pt, colback=red!6, colframe=red!40, boxrule=0.8pt, boxsep=3pt]{%
        \small\sffamily\bfseries\color{red}\underline{Di chuột vào đây để thấy hiệu ứng!}%
    }
    \quad (Chữ lập tức chuyển màu đỏ rực rỡ kèm đường gạch chân).
    \par\vspace{4pt}
    \item \textbf{Khi rê con trỏ chuột vào nút bấm <button>}:
    \par\vspace{2pt}
    \tcbox[on line, arc=3pt, colback=primaryBlue, colframe=primaryBlue!80, boxrule=0.8pt, boxsep=4pt]{%
        \small\sffamily\bfseries\color{white} Di chuột vào nút này%
    }
    \quad (Nền chuyển sang xanh dương đậm \codeinline{\#007bff}, chữ trắng nổi bật, con trỏ đổi thành hình bàn tay \codeinline{pointer}).
\end{itemize}
\end{browserWindow}

\begin{conceptBox}[3. Các Phần Tử Có Thể Dùng Với :hover]
Hầu như mọi phần tử HTML đều có thể áp dụng \codeinline{:hover}, không chỉ các thẻ tương tác: \codeinline{<a>}, \codeinline{<button>}, \codeinline{<div>}, \codeinline{<p>}, \codeinline{<img>}, \codeinline{<span>}, v.v.

\textbf{Ví dụ với hình ảnh}:
\begin{codeWindow}[Hiệu Ứng Phóng To Ảnh Khi Hover]
\begin{lstlisting}[language=CSS3]
img:hover {
    transform: scale(1.1);
    transition: 0.3s ease;
}
\end{lstlisting}
\end{codeWindow}

\begin{browserWindow}[Kết Quả Hiển Thị Hiệu Ứng Phóng To Ảnh]
\sffamily\small
\begin{itemize}[leftmargin=1.5em, itemsep=3pt, parsep=0pt]
    \item \textbf{Ban đầu (Normal State)}:
    \par\vspace{2pt}
    \tcbox[on line, arc=3pt, colback=gray!10, colframe=gray!40, boxrule=0.8pt, boxsep=3pt]{%
        \small\sffamily  Hình ảnh ($100\%$)%
    }
    \quad (Ảnh hiển thị với tỷ lệ mặc định $100\%$).
    \par\vspace{4pt}
    \item \textbf{Khi di chuột vào (Hover State)}:
    \par\vspace{2pt}
    \tcbox[on line, arc=3pt, colback=primaryBlue!10, colframe=primaryBlue!50, boxrule=0.8pt, boxsep=5pt]{%
        \small\sffamily\bfseries\color{primaryBlue}  Hình ảnh ($110\%$)%
    }
    \quad (Nhờ thuộc tính \codeinline{transform: scale(1.1)}, ảnh tự động phóng to mượt mà $0.3$s).
\end{itemize}
\end{browserWindow}
\end{conceptBox}

\begin{tipBox}[4. Ứng Dụng Thực Tế Của :hover]
\begin{itemize}[leftmargin=1.5em, itemsep=2pt, parsep=0pt]
    \item \textbf{Thiết kế menu điều hướng}: Khi hover vào mục menu sẽ đổi màu hoặc hiển thị menu thả xuống (dropdown).
    \item \textbf{Nút bấm nổi bật}: Nút thay đổi màu sắc hoặc hiệu ứng khi người dùng di chuột qua.
    \item \textbf{Thẻ sản phẩm (hover card)}: Khi hover vào thẻ sản phẩm sẽ hiển thị thêm thông tin chi tiết.
    \item \textbf{Hiệu ứng ảnh}: Phóng to ảnh, làm mờ hoặc thêm bộ lọc khi hover.
    \item \textbf{Tooltip}: Hiển thị chú thích hoặc thông tin bổ sung khi di chuột vào phần tử.
\end{itemize}
\end{tipBox}

\begin{codeWindow}[Ví Dụ Menu Với Dropdown Khi Hover]
\begin{lstlisting}[language=HTML5]
<style>
    .menu {
        position: relative;
        display: inline-block;
    }
    .dropdown {
        display: none;
        position: absolute;
        background-color: #f9f9f9;
        padding: 10px;
        box-shadow: 0 2px 10px rgba(0, 0, 0, 0.1);
    }
    .menu:hover .dropdown {
        display: block;
    }
</style>
<div class="menu">
    <button>Sample Text</button>
    <div class="dropdown">
        <p>Sample Text</p>
        <p>Sample Text</p>
        <p>Sample Text</p>
    </div>
</div>
\end{lstlisting}
\end{codeWindow}

\begin{browserWindow}[Minh Họa Trực Quan Trạng Thái Giao Diện Dropdown Menu]
\begin{itemize}[leftmargin=1.5em, itemsep=3pt, parsep=0pt]
    \item \textbf{Trạng thái 1 -- Ban đầu (Chưa di chuột vào nút Menu)}:
    \par\vspace{2pt}
    \tcbox[on line, arc=3pt, colback=gray!10, colframe=gray!40, boxrule=0.8pt, boxsep=3pt]{%
        \small\sffamily\bfseries [ Menu ]%
    }
    \quad (Thẻ \codeinline{.dropdown} ở trạng thái \codeinline{display: none} nên hoàn toàn bị ẩn).
    \par\vspace{5pt}
    \item \textbf{Trạng thái 2 -- Khi di chuột vào nút Menu (:hover)}:
    \par\vspace{2pt}
    \tcbox[on line, arc=3pt, colback=primaryBlue!10, colframe=primaryBlue!50, boxrule=0.8pt, boxsep=4pt]{%
        \small\sffamily\bfseries\color{primaryBlue} [ Menu ] $\rightarrow$ Menu thả xuống hiện ra ngay phía dưới:%
    }
    \par\vspace{4pt}
    \begin{tcolorbox}[arc=4pt, colback=primaryBlue!5, colframe=primaryBlue!30, boxrule=0.8pt, top=6pt, bottom=6pt, left=10pt, right=10pt]
      \sffamily\small\bfseries\color{primaryBlue} Nội Dung Dropdown Xuất Hiện:
      \begin{itemize}[leftmargin=1.2em, itemsep=1pt]
          \item Mục 1
          \item Mục 2
          \item Mục 3
      \end{itemize}
    \end{tcolorbox}
\end{itemize}
\end{browserWindow}

\begin{warnBox}[5. Lưu Ý Quan Trọng Khi Sử Dụng :hover]
\begin{itemize}[leftmargin=1.5em, itemsep=3pt, parsep=0pt]
    \item \textbf{Hover không hoạt động trên thiết bị cảm ứng}: Trên điện thoại hoặc tablet không có con trỏ chuột, nên bạn có thể kết hợp với \codeinline{:focus} hoặc JavaScript để hỗ trợ trải nghiệm người dùng.
    \item \textbf{Thứ tự viết CSS (Quy tắc LVHA)}: Khi kết hợp với các pseudo-class khác (như \codeinline{:link}, \codeinline{:visited}, \codeinline{:active}), hãy viết đúng thứ tự:
    \par\vspace{3pt}
    \tcbox[on line, arc=3pt, colback=primaryBlue!6, colframe=primaryBlue!40, boxrule=0.8pt, boxsep=3pt, left=6pt, right=6pt]{%
        \ttfamily\color{primaryBlue}\bfseries :link $\rightarrow$ :visited $\rightarrow$ :hover $\rightarrow$ :active%
    }
\end{itemize}
\end{warnBox}

\begin{conceptBox}[6. Tổng Kết Về :hover]
\codeinline{:hover} là một trong những pseudo-class mạnh mẽ nhất trong CSS, giúp bạn:
\begin{itemize}[leftmargin=1.5em, itemsep=2pt, parsep=0pt]
    \item \textbf{Tăng tính tương tác}: Làm nổi bật các phần tử khi người dùng di chuột qua.
    \item \textbf{Cải thiện UX/UI}: Tạo hiệu ứng hình ảnh, nút bấm, và dropdown mượt mà.
    \item \textbf{Dễ dàng triển khai}: Chỉ cần vài dòng CSS là đã có thể làm giao diện sinh động hơn!
\end{itemize}
\end{conceptBox}

\paragraph{d. Bộ Chọn :active (Trạng Thái Khi Nhấp Bấm Giữ Chuột)}

Trong CSS, \codeinline{:active} là một pseudo-class được sử dụng để chọn một phần tử khi nó đang được nhấp vào (tức là khi bạn giữ chuột trái trên phần tử đó nhưng chưa thả ra). Nó thường được dùng để tạo hiệu ứng nhấn, giúp giao diện sinh động hơn!

\begin{conceptBox}[1. Cú Pháp Tổng Quát]
\begin{codeWindow}[Cú Pháp Định Kiểu :active]
\begin{lstlisting}[language=CSS3]
selector:active {
    /* Sample Text */
}
\end{lstlisting}
\end{codeWindow}
\end{conceptBox}

\begin{conceptBox}[2. Ví Dụ Cơ Bản Với Nút Bấm]
\begin{codeWindow}[Mã Nguồn HTML \& CSS Tạo Hiệu Ứng Nút Bấm Với :active]
\begin{lstlisting}[language=HTML5]
<button class="btn">Sample Text</button>
<style>
    .btn {
        padding: 10px 20px;
        background: blue;
        color: white;
        border: none;
        border-radius: 5px;
    }
    .btn:active {
        background: red;
        transform: scale(0.95);
    }
</style>
\end{lstlisting}
\end{codeWindow}
\textbf{Giải thích}:
\begin{itemize}[leftmargin=1.5em, itemsep=2pt, parsep=0pt]
    \item Khi nhấn giữ nút, nền sẽ chuyển sang màu đỏ và nút sẽ hơi thu nhỏ lại (\codeinline{transform: scale(0.95)} -- hiệu ứng nhấn).
    \item Khi thả chuột, nút quay về trạng thái bình thường.
\end{itemize}
\end{conceptBox}

\begin{browserWindow}[Kết Quả Hiển Thị Demo Nút Bấm Cơ Bản Với :active]
\begin{itemize}[leftmargin=1.5em, itemsep=3pt, parsep=0pt]
    \item \textbf{Trạng thái ban đầu}: Nút màu xanh dương (\codeinline{blue}), chữ trắng, kích thước $100\%$.
    \item \textbf{Khi nhấp giữ chuột trái (:active)}: Nút lập tức đổi nền sang màu đỏ (\codeinline{red}) và thu nhỏ về $95\%$ kích thước ban đầu.
    \item \textbf{Khi thả chuột ra}: Nút quay về nền màu xanh và kích thước bình thường.
\end{itemize}
\end{browserWindow}

\begin{tipBox}[3. Ứng Dụng Thực Tế Của :active]
\begin{itemize}[leftmargin=1.5em, itemsep=3pt, parsep=0pt]
    \item \textbf{Tạo hiệu ứng bấm nút}: Đổi màu hoặc kích thước khi nhấn nút.
    \item \textbf{Hiệu ứng nhấn menu hoặc link}: Làm nổi bật liên kết khi nhấp vào.
    \item \textbf{Phản hồi hành động người dùng}: Giúp người dùng thấy ngay lập tức họ đã nhấn vào đâu.
\end{itemize}
\end{tipBox}

\paragraph{4. Ví Dụ Định Kiểu Tương Tác Với Thẻ Liên Kết (Link):}

\begin{codeWindow}[Mã Nguồn CSS Kết Hợp :hover Và :active Cho Thẻ <a>]
\begin{lstlisting}[language=HTML5]
<a href="#" class="link">Sample Text</a>
<style>
    .link {
        color: blue;
        text-decoration: none;
        padding: 5px;
    }
    .link:hover {
        color: green;
    }
    .link:active {
        color: red;
    }
</style>
\end{lstlisting}
\end{codeWindow}

\begin{conceptBox}[Giải Thích Chi Tiết Các Trạng Thái Tương Tác]
\begin{itemize}[leftmargin=1.5em, itemsep=2pt, parsep=0pt]
    \item Khi di chuột vào link $\rightarrow$ màu xanh lá (\codeinline{green}).
    \item Khi nhấn giữ link $\rightarrow$ màu đỏ (\codeinline{red} -- trạng thái active).
    \item Khi thả chuột $\rightarrow$ màu trở về bình thường (\codeinline{blue}).
\end{itemize}
\end{conceptBox}

\begin{browserWindow}[Kết Quả Hiển Thị Minh Họa Tương Tác Link Với :active]
\begin{itemize}[leftmargin=1.5em, itemsep=3pt, parsep=0pt]
    \item \textbf{Di chuột vào link}: Chữ đổi màu xanh lá (\codeinline{green}).
    \item \textbf{Nhấn giữ chuột trái}: Chữ chuyển ngay sang màu đỏ (\codeinline{red}).
    \item \textbf{Thả chuột ra}: Chữ trở về màu xanh dương (\codeinline{blue}) mặc định.
\end{itemize}
\end{browserWindow}

\begin{warnBox}[5. Lưu Ý Quan Trọng Khi Sử Dụng :active]
\begin{enumerate}[leftmargin=1.5em, itemsep=3pt, parsep=0pt]
    \item \textbf{Thời gian hiệu lực}: Chỉ hoạt động trong lúc nhấn giữ chuột. Khi thả ra, trạng thái \codeinline{:active} sẽ mất ngay lập tức.
    \item \textbf{Thứ tự ưu tiên khi viết CSS (Quy tắc LVHA)}:
    \par\vspace{2pt}
    Nếu kết hợp nhiều pseudo-class cho liên kết, thứ tự đúng bắt buộc phải là:
    \par\vspace{3pt}
    \tcbox[on line, arc=3pt, colback=primaryBlue!6, colframe=primaryBlue!40, boxrule=0.8pt, boxsep=3pt, left=6pt, right=6pt]{%
        \ttfamily\color{primaryBlue}\bfseries a:link $\rightarrow$ a:visited $\rightarrow$ a:hover $\rightarrow$ a:active%
    }
    \par\vspace{3pt}
    \textit{Lý do}: Vì trình duyệt áp dụng theo thứ tự này, nếu bạn viết \codeinline{:active} trước \codeinline{:hover}, hiệu ứng \codeinline{:active} có thể không hiển thị đúng do bị \codeinline{:hover} ghi đè.
    \item \textbf{Khả năng kích hoạt của các phần tử}:
    \begin{itemize}[leftmargin=1.2em, itemsep=2pt]
        \item Phần tử có thể nhấp mặc định: \codeinline{button}, \codeinline{a}, \codeinline{input}, \codeinline{label}, v.v.
        \item Với thẻ không tương tác (ví dụ \codeinline{div}, \codeinline{span}), bạn cần thêm thuộc tính \codeinline{tabindex="0"} hoặc dùng JavaScript để bắt sự kiện.
    \end{itemize}
\end{enumerate}
\end{warnBox}

\begin{conceptBox}[6. Ví Dụ Nâng Cao: Tạo Hiệu Ứng Nhấn Cho Khối Div Bằng Tabindex]
\begin{codeWindow}[Mã Nguồn HTML \& CSS Sử Dụng Tabindex Cho Card Div]
\begin{lstlisting}[language=HTML5]
<div class="card" tabindex="0">
    Sample Text
</div>
<style>
    .card {
        width: 200px;
        height: 100px;
        background: lightblue;
        display: flex;
        align-items: center;
        justify-content: center;
        border: 2px solid blue;
    }
    .card:active {
        background: orange;
        border: 2px solid red;
    }
</style>
\end{lstlisting}
\end{codeWindow}
\textbf{Giải thích}:
\begin{itemize}[leftmargin=1.5em, itemsep=2pt, parsep=0pt]
    \item \codeinline{tabindex="0"} giúp phần tử \codeinline{<div>} có thể nhận focus, từ đó sử dụng được \codeinline{:active}.
    \item Khi nhấn vào, nền đổi thành màu cam (\codeinline{orange}), viền đổi thành màu đỏ (\codeinline{red}).
\end{itemize}
\end{conceptBox}

\begin{browserWindow}[Kết Quả Hiển Thị Khối Div Nâng Cao Với Tabindex]
\sffamily\small
\begin{itemize}[leftmargin=1.5em, itemsep=3pt, parsep=0pt]
    \item \textbf{Trạng thái ban đầu}:
    \par\vspace{2pt}
    \tcbox[on line, arc=3pt, colback=blue!10, colframe=blue, boxrule=1.5pt, boxsep=4pt, left=10pt, right=10pt]{%
        \small\sffamily\bfseries\color{blue} Nhấn vào tôi!%
    }
    \quad (Khối \codeinline{card} có nền xanh nhạt, viền xanh dương).
    \par\vspace{4pt}
    \item \textbf{Khi nhấn giữ chuột vào khối card (:active)}:
    \par\vspace{2pt}
    \tcbox[on line, arc=3pt, colback=orange!20, colframe=red, boxrule=1.5pt, boxsep=4pt, left=10pt, right=10pt]{%
        \small\sffamily\bfseries\color{red} Nhấn vào tôi!%
    }
    \quad (Nền đổi ngay sang màu cam \codeinline{orange}, viền chuyển thành màu đỏ \codeinline{red}).
\end{itemize}
\end{browserWindow}

\begin{tipBox}[7. Khi Nào Nên Dùng :active]
\begin{itemize}[leftmargin=1.5em, itemsep=3pt, parsep=0pt]
    \item \textbf{Nút bấm, liên kết}: Tạo phản hồi nhanh khi người dùng nhấn vào.
    \item \textbf{Trò chơi hoặc ứng dụng web}: Phát hiện hành động nhấn giữ để thực hiện thao tác (ví dụ như giữ nút bắn trong game).
    \item \textbf{Giao diện kéo thả}: Làm nổi bật phần tử đang được nhấn giữ.
\end{itemize}
\end{tipBox}

\begin{conceptBox}[8. Tóm Lại Về :active]
\begin{itemize}[leftmargin=1.5em, itemsep=2pt, parsep=0pt]
    \item \codeinline{:active} giúp chọn phần tử trong lúc nhấn giữ chuột.
    \item \textbf{Ứng dụng}: Tạo hiệu ứng bấm nút, nhấn link, hoặc làm nổi bật phần tử đang được chọn.
    \item Hoạt động tốt với các phần tử tương tác và có thể mở rộng với thuộc tính \codeinline{tabindex}.
\end{itemize}
\end{conceptBox}

\paragraph{e. Chuyên Đề Phân Tích Sâu: Bộ Chọn :focus, :focus-visible \& :focus-within (Trạng Thái Tiêu Điểm Nhập Liệu \& Bàn Phím)}

Trạng thái tiêu điểm (\term{Focus State}) là một trong những khía cạnh quan trọng nhất trong thiết kế giao diện Web hiện đại và Trợ năng (\term{Web Accessibility - WCAG}). Nó giúp người dùng nhận biết chính xác phần tử nào đang nhận tương tác từ bàn phím hoặc thiết bị nhập liệu.

\begin{conceptBox}[1. Bản Chất \& Cú Pháp Của Pseudo-Class :focus]
\begin{itemize}[leftmargin=1.5em, itemsep=3pt, parsep=0pt]
    \item \textbf{Bản chất}: Pseudo-class \codeinline{:focus} kích hoạt khi một phần tử nhận tiêu điểm điều hướng -- bằng cách nhấp chuột vào, chạm ngón tay trên thiết bị cảm ứng, dùng phím \codeinline{Tab} di chuyển tới, hoặc gọi hàm \codeinline{element.focus()} trong JavaScript.
    \item \textbf{Các phần tử có khả năng nhận focus mặc định}: Thẻ ô nhập liệu (\codeinline{<input>}, \codeinline{<textarea>}, \codeinline{<select>}), nút bấm (\codeinline{<button>}), liên kết (\codeinline{<a>}), thẻ tóm tắt (\codeinline{<summary>}) và bất kỳ phần tử nào có thuộc tính \codeinline{tabindex}.
    \item \textbf{Cú pháp tổng quát}:
    \begin{codeWindow}[Cú Pháp Định Kiểu :focus]
\begin{lstlisting}[language=CSS3]
selector:focus {
    outline: none; /* Sample Text */
    border-color: #007bff; /* Sample Text */
    box-shadow: 0 0 8px rgba(0, 123, 255, 0.4); /* Sample Text */
}
\end{lstlisting}
    \end{codeWindow}
\end{itemize}
\end{conceptBox}

\begin{conceptBox}[2. Ví Dụ Tùy Biến Giao Diện Form Chuẩn UI/UX]
\begin{codeWindow}[Mã Nguồn HTML \& CSS Định Kiểu Ô Nhập Liệu Khi Focus]
\begin{lstlisting}[language=HTML5]
<style>
    .form-input {
        width: 100%;
        padding: 12px 16px;
        font-size: 14px;
        border: 2px solid #ccc;
        border-radius: 8px;
        outline: none;
        transition: all 0.3s ease;
    }
    .form-input:focus {
        border-color: #007bff;
        box-shadow: 0 0 10px rgba(0, 123, 255, 0.35);
        background-color: #f8fbff;
    }
</style>

<label for="username">Sample Text</label>
<input type="text" id="username" class="form-input" placeholder="Sample Text">
\end{lstlisting}
\end{codeWindow}
\end{conceptBox}

\begin{browserWindow}[Kết Quả Hiển Thị Trực Quan Khi Focus Ô Nhập Liệu]
\begin{itemize}[leftmargin=1.5em, itemsep=3pt, parsep=0pt]
    \item \textbf{Trạng thái ban đầu (Unfocused)}: Ô nhập liệu hiển thị viền xám nhạt (\codeinline{\#ccc}), nền màu trắng phẳng.
    \item \textbf{Khi nhấp chuột hoặc bấm phím Tab vào ô (Focused)}: Viền ô lập tức sáng lên màu xanh dương đậm (\codeinline{\#007bff}), nền chuyển sang màu nhạt (\codeinline{\#f8fbff}), kèm theo quầng sáng tỏa mượt (\codeinline{box-shadow glow}) xung quanh khung nhập.
\end{itemize}
\end{browserWindow}

\begin{conceptBox}[3. Chuyên Đề Phân Biệt :focus Và :focus-visible (Modern Accessibility)]
\begin{itemize}[leftmargin=1.5em, itemsep=3pt, parsep=0pt]
    \item \textbf{Vấn đề của :focus truyền thống}: Khi người dùng nhấp chuột vào nút bấm (\codeinline{<button>}), trình duyệt sẽ vẽ một đường viền đen (\term{Focus Ring / Outline}) gây mất thẩm mỹ giao diện. Nếu lập trình viên xóa viền này bằng \codeinline{outline: none}, người dùng phím \codeinline{Tab} (người khuyết tật dùng thiết bị trợ năng) sẽ hoàn toàn không thấy tiêu điểm ở đâu!
    \item \textbf{Giải pháp Modern CSS :focus-visible}: Chỉ hiển thị đường viền tiêu điểm khi phần tử được điều hướng bằng \textbf{bàn phím (Phím Tab)}, còn khi nhấp bằng chuột thì trình duyệt sẽ \textbf{tự động ẩn đường viền}!
\end{itemize}

\begin{codeWindow}[Mã Nguồn CSS Tối Ưu Tiêu Điểm Chuẩn WCAG Với :focus-visible]
\begin{lstlisting}[language=CSS3]
/* Sample Text */
button:focus {
    outline: none;
}

/* Sample Text */
button:focus-visible {
    outline: 3px solid #ff9800;
    outline-offset: 3px;
}
\end{lstlisting}
\end{codeWindow}
\end{conceptBox}

\begin{browserWindow}[So Sánh Trải Nghiệm Chuột Vs Bàn Phím Với :focus-visible]
\sffamily\small
\begin{itemize}[leftmargin=1.5em, itemsep=3pt, parsep=0pt]
    \item \textbf{Khi nhấp chuột vào nút bấm (Mouse User)}:
    \par\vspace{2pt}
    \tcbox[on line, arc=3pt, colback=primaryBlue, colframe=primaryBlue, boxrule=0.8pt, boxsep=4pt]{%
        \small\sffamily\bfseries\color{white} Nút Bấm%
    }
    \quad (Nút kích hoạt bình thường mà \textbf{KHÔNG xuất hiện viền đen} — \codeinline{outline} bị ẩn).
    \par\vspace{4pt}
    \item \textbf{Khi dùng phím Tab di chuyển tới nút (Keyboard User)}:
    \par\vspace{2pt}
    \tcbox[on line, arc=3pt, colback=primaryBlue, colframe=orange!90, boxrule=2pt, boxsep=4pt]{%
        \small\sffamily\bfseries\color{white} Nút Bấm%
    }
    \quad (Xuất hiện viền cam rực rỡ \codeinline{\#ff9800} cách mép nút $3\text{px}$ giúp định vị tiêu điểm).
\end{itemize}
\end{browserWindow}

\begin{conceptBox}[4. Chuyên Đề Phân Tích :focus-within (Tiêu Điểm Thẻ Cha Khi Con Focus)]
\begin{itemize}[leftmargin=1.5em, itemsep=3pt, parsep=0pt]
    \item \textbf{Bản chất}: Pseudo-class \codeinline{:focus-within} kích hoạt trên phần tử cha nếu \textbf{chính nó} HOẶC \textbf{bất kỳ phần tử con/cháu nào nằm bên trong nó} nhận tiêu điểm focus!
    \item \textbf{Ứng dụng thực tế}: Định kiểu lại khung chứa Search Bar, Card Form, hoặc nhóm nút bấm khi người dùng tương tác với bất kỳ ô nhập liệu nào bên trong.
\end{itemize}

\begin{codeWindow}[Mã Nguồn CSS Tạo Thanh Tìm Kiếm Đổi Màu Khung Với :focus-within]
\begin{lstlisting}[language=HTML5]
<style>
    .search-box {
        display: flex;
        align-items: center;
        padding: 8px 16px;
        border: 2px solid #ddd;
        border-radius: 25px;
        background: #f9f9f9;
        transition: all 0.3s ease;
    }
    /* Sample Text */
    .search-box:focus-within {
        border-color: #28a745;
        background: #fff;
        box-shadow: 0 4px 12px rgba(40, 167, 69, 0.25);
    }
    .search-input {
        border: none;
        outline: none;
        background: transparent;
        width: 100%;
    }
</style>

<div class="search-box">
    <span class="icon">\textbf{[KiTextm Tra]}</span>
    <input type="text" class="search-input" placeholder="Sample Text">
</div>
\end{lstlisting}
\end{codeWindow}
\end{conceptBox}

\begin{browserWindow}[Minh Họa Trực Quan Thanh Tìm Kiếm Với :focus-within]
\begin{itemize}[leftmargin=1.5em, itemsep=3pt, parsep=0pt]
    \item \textbf{Khi chưa nhấp vào ô nhập (Normal State)}: Khung chứa \codeinline{.search-box} có viền màu xám (\codeinline{\#ddd}) và nền xám nhạt (\codeinline{\#f9f9f9}).
    \item \textbf{Khi nhấp vào ô input tìm kiếm (:focus-within)}: Nhờ bộ chọn \codeinline{:focus-within}, toàn bộ khung bao ngoài \codeinline{.search-box} (chứa cả icon kính lúp và input) đồng loạt chuyển nền sang trắng rực rỡ, viền đổi thành màu xanh lá (\codeinline{\#28a745}) kèm hiệu ứng bóng đổ nổi khối nổi bật.
\end{itemize}
\end{browserWindow}

\begin{prettyTableBox}[Bảng So Sánh Bộ Chọn :focus vs :focus-visible vs :focus-within]
\rowcolors{2}{white}{primaryBlue!4}
\sffamily\small
\begin{tabularx}{\linewidth}{|>{\raggedright\arraybackslash \ttfamily\color{primaryBlue}\bfseries\raggedright\arraybackslash}m{2.5cm}|>{\raggedright\arraybackslash}m{3.2cm}|>{\ttfamily\small\raggedright\arraybackslash}m{4.0cm}|Z|}
\hline
\rowcolor{primaryBlue} \thd{\sffamily Selector} & \thd{\sffamily Điều kiện kích hoạt} & \thd{\sffamily Ví dụ ứng dụng} & \thd{\sffamily Trải nghiệm UX / a11y} \\
\hline
:focus & \sffamily\bfseries Bất kỳ khi nào phần tử nhận tiêu điểm (Chuột, Tab, JS). & input:focus \{ border-color: blue; \} & Tuyệt vời cho Form input nhưng có thể hiện viền đen xấu trên Nút bấm khi nhấp chuột. \\
\hline
:focus-visible & \sffamily\bfseries CHỈ khi phần tử nhận tiêu điểm từ bàn phím (Phím Tab). & button:focus-visible \{ outline: 2px solid orange; \} & Chuẩn WCAG hiện đại: Ẩn viền khi dùng chuột, hiện viền rõ nét khi dùng bàn phím. \\
\hline
:focus-within & \sffamily\bfseries Khi CHÍNH NÓ hoặc BẤT KỲ THẺ CON NÀO nhận focus. & .form-group:focus-within \{ background: lightyellow; \} & Tô sáng toàn bộ khối khung chứa (Form Group, Search Box) khi con bên trong tương tác. \\
\hline
\end{tabularx}
\end{prettyTableBox}

\begin{warnBox}[Lưu Ý Quan Trọng Cho Focus Selectors (Accessibility Checklist)]
\begin{itemize}[leftmargin=1.5em, itemsep=3pt, parsep=0pt]
    \item \textbf{TUYỆT ĐỐI KHÔNG DÙNG \codeinline{outline: none} ĐƠN ĐỘC}: Nếu bạn xóa \codeinline{outline}, bắt buộc phải cung cấp style thay thế (\codeinline{border-color}, \codeinline{box-shadow}, hoặc \codeinline{:focus-visible}) để người dùng bàn phím không bị mất tiêu điểm.
    \item \textbf{Thứ tự ưu tiên khi viết CSS (Quy tắc LVHFA)}: Khi viết cùng lúc cho thẻ liên kết \codeinline{<a>}, thứ tự chuẩn là:
    \par\vspace{3pt}
    \tcbox[on line, arc=3pt, colback=primaryBlue!6, colframe=primaryBlue!40, boxrule=0.8pt, boxsep=3pt, left=6pt, right=6pt]{%
        \ttfamily\color{primaryBlue}\bfseries :link $\rightarrow$ :visited $\rightarrow$ :hover $\rightarrow$ :focus $\rightarrow$ :active%
    }
\end{itemize}
\end{warnBox}

\begin{conceptBox}[Tóm Tắt Cốt Lõi Về Focus Selectors]
\begin{itemize}[leftmargin=1.5em, itemsep=2pt, parsep=0pt]
    \item \codeinline{:focus}: Kiểm soát tiêu điểm cá nhân cho ô nhập liệu.
    \item \codeinline{:focus-visible}: Tối ưu trải nghiệm bàn phím chuẩn Accessibility.
    \item \codeinline{:focus-within}: Quản lý tiêu điểm cho khối khung chứa cha từ phần tử con.
\end{itemize}
\end{conceptBox}

\paragraph{f. Bộ Chọn Lớp Giả :target (Trạng Thái Định Vị Thẻ Theo URL Anchor ID)}

Bộ chọn lớp giả \codeinline{:target} trong CSS được dùng để chọn phần tử HTML được nhắm đến (target) bởi liên kết có chứa \codeinline{\#id} trong URL, tức là khi phần tử đó được "kích hoạt" qua một đường dẫn có gắn \codeinline{\#id}.

\begin{codeWindow}[Cú Pháp Cơ Bản Của Bộ Chọn :target]
\begin{lstlisting}[language=CSS3]
#section1:target {
    background: yellow;
}
\end{lstlisting}
\end{codeWindow}

\begin{conceptBox}[Cách Hoạt Động Cơ Bản Của :target]
\begin{enumerate}[leftmargin=1.5em, itemsep=3pt, parsep=0pt]
    \item \textbf{Bước 1}: Khi bạn click vào một liên kết như: \codeinline{<a href="\#section1">Xem mục 1</a>}.
    \item \textbf{Bước 2}: Trình duyệt sẽ tự động cuộn đến phần tử có thuộc tính \codeinline{id="section1"} trong trang web.
    \item \textbf{Bước 3}: CSS \codeinline{\#section1:target} sẽ lập tức áp dụng lên phần tử này ngay tại thời điểm nó được "target".
\end{enumerate}
\end{conceptBox}

\paragraph{Ví Dụ Đầy Đủ \& Phân Tích So Sánh Trước / Sau Khi Dùng :target:}

\begin{codeWindow}[Trường Hợp 1 -- Khi KHÔNG Dùng :target (Chưa Nhấp Link Đã Có Style)]
\begin{lstlisting}[language=HTML5]
<a href="#about">Sample Text</a>
<div id="about">Sample Text</div>

<style>
  #about {
    background-color: yellow;
    padding: 10px;
    border: 2px solid red;
  }
</style>
\end{lstlisting}
\end{codeWindow}

\begin{warnBox}[Phân Tích Lỗi Khi KHÔNG Dùng :target]
CSS ở ví dụ trên đang áp dụng mặc định cho phần tử có \codeinline{id="about"}. Điều này \textbf{không liên quan gì đến :target cả} -- tức là khối \codeinline{\#about} luôn luôn có màu nền vàng và viền đỏ ngay từ đầu, ngay cả khi người dùng chưa click vào link!
\end{warnBox}

\begin{codeWindow}[Trường Hợp 2 -- Khi DÙNG :target (Chỉ Kích Hoạt Style Khi Click Link)]
\begin{lstlisting}[language=HTML5]
<a href="#about">Sample Text</a>
<div id="about">Sample Text</div>

<style>
  #about:target {
    background-color: yellow;
    padding: 10px;
    border: 2px solid red;
  }
</style>
\end{lstlisting}
\end{codeWindow}

\begin{browserWindow}[Kết Quả So Sánh Trước Và Sau Khi Nhấp Link]
\sffamily\small
\begin{itemize}[leftmargin=1.5em, itemsep=4pt, parsep=0pt]
    \item \textbf{Trước khi nhấp vào link}: Phần tử \codeinline{\#about} hiển thị hoàn toàn bình thường (không có nền vàng, không có viền đỏ).
    \item \textbf{Khi bạn click vào liên kết "Xem giới thiệu"}: Phần tử \codeinline{\#about} sẽ:
    \begin{itemize}[leftmargin=1.2em, itemsep=2pt]
        \item Được trình duyệt tự động cuộn tới.
        \item Lập tức chuyển sang màu nền vàng kẽm viền đỏ nổi bật nhờ kích hoạt bộ chọn \codeinline{\#about:target}.
    \end{itemize}
\end{itemize}
\end{browserWindow}

\begin{tipBox}[Lưu Ý Cốt Lõi Về :target]
\begin{itemize}[leftmargin=1.5em, itemsep=2pt, parsep=0pt]
    \item \codeinline{:target} \textbf{chỉ áp dụng khi URL có chứa \#id}.
    \item Khi bạn truy cập lại URL gốc (không chứa \codeinline{\#}) hoặc click sang phần tử khác, thì CSS của \codeinline{:target} trên phần tử cũ sẽ tự động mất tác dụng.
\end{itemize}
\end{tipBox}

\begin{prettyTableBox}[Các Ứng Dụng Thực Tế Của Bộ Chọn :target]
\rowcolors{2}{white}{primaryBlue!4}
\sffamily\small
\begin{tabularx}{\linewidth}{|>{\raggedright\arraybackslash \bfseries\color{primaryBlue}}K{4.2cm}|Z|}
\hline
\rowcolor{primaryBlue} \thd{\sffamily Ứng dụng thực tế} & \thd{\sffamily Mô tả chi tiết} \\
\hline
Điều hướng trong 1 trang & Highlight mục đang xem trong trang (kiểu như Table of Contents / Mục lục bài viết). \\
\hline
Accordion / Toggle đơn giản & Mở / đóng khối nội dung thả xuống mà không cần sử dụng JavaScript. \\
\hline
Hiệu ứng nhấn nút cuộn xuống & Nhấn nút rồi scroll + đổi màu mục tiêu được cuộn đến. \\
\hline
\end{tabularx}
\end{prettyTableBox}

\paragraph{So Sánh :target Với Các Selector Trạng Thái Khác:}

\begin{prettyTableBox}[Bảng So Sánh :hover vs :focus vs :target]
\rowcolors{2}{white}{primaryBlue!4}
\sffamily\small
\begin{tabularx}{\linewidth}{|>{\raggedright\arraybackslash \ttfamily\color{primaryBlue}\bfseries}K{2.5cm}|Z|}
\hline
\rowcolor{primaryBlue} \thd{\sffamily Selector} & \thd{\sffamily Ý nghĩa \& Thời điểm kích hoạt} \\
\hline
:hover & Kích hoạt khi con trỏ chuột rê vào phần tử. \\
\hline
:focus & Kích hoạt khi phần tử nhận tiêu điểm nhập liệu (thường dùng cho input, button). \\
\hline
:target & Kích hoạt khi URL có chứa mã băm \codeinline{\#id} trùng khớp hoàn toàn với \codeinline{id} của phần tử đó. \\
\hline
\end{tabularx}
\end{prettyTableBox}

\paragraph{Chi Tiết Về Cách Hoạt Động Của :target:}

\begin{conceptBox}[1. Cơ Chế Xử Lý Của Trình Duyệt Khi Duyệt URL]
\begin{itemize}[leftmargin=1.5em, itemsep=3pt, parsep=0pt]
    \item \textbf{Khi URL có chứa \#id} (ví dụ: \codeinline{example.com/page.html\#about}):
    \begin{itemize}[leftmargin=1.2em, itemsep=2pt]
        \item Trình duyệt tự động cuộn màn hình đến phần tử có \codeinline{id="about"}.
        \item Nếu CSS có khai báo \codeinline{\#about:target}, quy tắc CSS này sẽ được áp dụng.
    \end{itemize}
    \item \textbf{Khi KHÔNG có \#id trong URL}:
    \begin{itemize}[leftmargin=1.2em, itemsep=2pt]
        \item Không có phần tử nào bị xem là \codeinline{:target}.
        \item CSS \codeinline{:target} sẽ không áp dụng cho bất kỳ phần tử nào trên trang.
    \end{itemize}
\end{itemize}
\end{conceptBox}

\begin{codeWindow}[Ví Dụ Minh Họa Đơn Giản URL Target]
\begin{lstlisting}[language=HTML5]
<!-- Sample Text -->
<a href="#section1">Sample Text</a>
<div id="section1">Sample Text</div>

<!-- Sample Text -->
<style>
  #section1:target {
    background-color: yellow;
    border: 2px solid red;
  }
</style>
\end{lstlisting}
\end{codeWindow}

\begin{browserWindow}[Kết Quả Khi Click Và Khi Không Có Target]
\sffamily\small
\begin{itemize}[leftmargin=1.5em, itemsep=3pt, parsep=0pt]
    \item \textbf{Khi click vào liên kết}:
    \begin{enumerate}[leftmargin=1.2em, itemsep=1pt]
        \item URL đổi thành: \codeinline{.../page.html\#section1}.
        \item Thẻ \codeinline{<div id="section1">} trở thành target hiện tại.
        \item CSS \codeinline{:target} được áp dụng $\rightarrow$ Hiển thị nền vàng, viền đỏ.
    \end{enumerate}
    \item \textbf{Khi bạn Refresh lại trang (mất \#), Click sang mục khác, hoặc Chạy trang lần đầu}:
    \begin{itemize}[leftmargin=1.2em, itemsep=1pt]
        \item Không có phần tử nào là \codeinline{:target}, nên không có CSS \codeinline{:target} nào được áp dụng.
    \end{itemize}
\end{itemize}
\end{browserWindow}

\begin{prettyTableBox}[Bảng So Sánh Tình Huống: Không Dùng :target vs Dùng :target]
\rowcolors{2}{white}{primaryBlue!4}
\sffamily\small
\begin{tabularx}{\linewidth}{|>{\raggedright\arraybackslash \sffamily\bfseries}K{3.8cm}|>{\sffamily}K{4.2cm}|Z|}
\hline
\rowcolor{primaryBlue} \thd{\sffamily Tình huống tương tác} & \thd{\sffamily Không dùng :target} & \thd{\sffamily Dùng :target (Có CSS :target)} \\
\hline
Click link có \codeinline{href="\#abc"} & Trình duyệt chỉ cuộn đến phần tử \codeinline{id="abc"}. & Trình duyệt cuộn + áp dụng CSS \codeinline{:target}. \\
\hline
Click vào phần tử không có id & Không cuộn, không có gì xảy ra. & Không cuộn, không có gì xảy ra. \\
\hline
Không click gì & Không có thay đổi giao diện. & Không có thay đổi giao diện. \\
\hline
Muốn thay đổi giao diện khi cuộn & Không làm được bằng CSS thuần. & Có thể highlight tự động phần tử được nhắm. \\
\hline
\end{tabularx}
\end{prettyTableBox}

\begin{prettyTableBox}[Bảng Phân Tích Các Lưu Ý Khi Dùng :target]
\rowcolors{2}{white}{primaryBlue!4}
\sffamily\small
\begin{tabularx}{\linewidth}{|>{\raggedright\arraybackslash \ttfamily\color{primaryBlue}\bfseries}K{4.0cm}|Z|}
\hline
\rowcolor{primaryBlue} \thd{\sffamily Lưu ý kỹ thuật} & \thd{\sffamily Giải thích nguyên lý} \\
\hline
:target chỉ áp dụng khi URL chứa \#id trùng với id phần tử & CSS không tự động chạy nếu không có sự kiện dẫn tới mã băm \codeinline{\#id}. \\
\hline
Không dùng được để kiểm tra trạng thái thường xuyên & Nếu bạn muốn thay đổi giao diện dựa trên trạng thái khác (như toggle phức tạp), dùng JS sẽ linh hoạt hơn. \\
\hline
Có thể lợi dụng :target để tạo hiệu ứng toggle, accordion, modal... & Thực hiện hoàn toàn bằng CSS thuần mà không cần JavaScript nếu HTML được cấu trúc hợp lý. \\
\hline
Không nên lạm dụng :target cho UI phức tạp & Vì không kiểm soát được nhiều trạng thái như hover, active hoặc nhiều phần tử cùng lúc. \\
\hline
\end{tabularx}
\end{prettyTableBox}

\begin{prettyTableBox}[Tóm Tắt Nhanh Về Đặc Điểm Bộ Chọn :target]
\rowcolors{2}{white}{primaryBlue!4}
\sffamily\small
\begin{tabularx}{\linewidth}{|>{\raggedright\arraybackslash \sffamily\bfseries\color{primaryBlue}}K{3.5cm}|Z|}
\hline
\rowcolor{primaryBlue} \thd{\sffamily Đặc điểm} & \thd{\sffamily Chi tiết bộ chọn :target} \\
\hline
Điều kiện áp dụng & URL chứa \codeinline{\#id} trùng khớp với \codeinline{id} của phần tử. \\
\hline
Phạm vi ảnh hưởng & CHỈ 1 phần tử duy nhất tại 1 thời điểm. \\
\hline
Dùng để làm gì? & Highlight bài viết, Accordion toggle, Scroll-to-section. \\
\hline
Ưu điểm chính & Không cần JavaScript vẫn làm được hiệu ứng chuyển màu / highlight. \\
\hline
Nhược điểm chính & Không có trạng thái "tắt", khó linh hoạt khi cần chuyển trạng thái phức tạp. \\
\hline
\end{tabularx}
\end{prettyTableBox}

\paragraph{Mã Nguồn Hoàn Chỉnh Thử Nghiệm "Chuyện Tình Bông Hoa Màu Tím":}

\begin{codeWindow}[Mã Nguồn HTML5 \& CSS3 Đầy Đủ Thử Nghiệm :target]
\begin{lstlisting}[language=HTML5]
<!DOCTYPE html>
<html lang="vi">
<head>
  <meta charset="UTF-8" />
  <title>Sample Text</title>
  <style>
    h2:target {
      color: blue;
      background-color: #eef;
      padding: 10px;
      border-left: 5px solid blue;
      transition: all 0.3s ease;
    }
  </style>
</head>
<body>
  <h1>Sample Text</h1>
  <ul>
    <li><a href="#mobai">Sample Text</a></li>
    <li><a href="#thanbai">Sample Text</a></li>
    <li><a href="#ketbai">Sample Text</a></li>
  </ul>
  <h2 id="mobai">Sample Text</h2>
  <p>Paragraph text</p>
  <h2 id="thanbai">Sample Text</h2>
  <p>Paragraph text</p>
  <h2 id="ketbai">Sample Text</h2>
  <p>Paragraph text</p>
</body>
</html>
\end{lstlisting}
\end{codeWindow}

\begin{browserWindow}[Kết Quả Hiển Thị: Trạng Thái Target Khi Nhấp Liên Kết Anchor]
\sffamily\small
\begin{itemize}[leftmargin=1.5em, itemsep=3pt, parsep=0pt]
    \item \textbf{Thanh điều hướng liên kết}:
    \begin{itemize}[leftmargin=1.2em, itemsep=2pt]
        \item \codeinline{Chuyển đến Mở bài}
        \item \textcolor{blue}{\textbf{Chuyển đến Thân bài}} $\leftarrow$ \textit{(Đã nhấp liên kết)}
        \item \codeinline{Chuyển đến Kết bài}
    \end{itemize}
    \item \textbf{Giao diện trực quan bài viết}:
    \begin{itemize}[leftmargin=1.2em, itemsep=2pt]
        \item \textbf{Mở bài}: Tiêu đề \codeinline{<h2>} hiển thị màu đen mặc định.
        \item \textbf{Thân bài}: Kích hoạt \codeinline{h2:target} $\rightarrow$ Tiêu đề chuyển sang \textbf{chữ màu xanh dương}, viền trái màu xanh đậm \codeinline{5px solid blue}, nền xanh nhạt nổi bật thu hút ánh nhìn!
        \item \textbf{Kết bài}: Tiêu đề \codeinline{<h2>} hiển thị màu đen mặc định.
    \end{itemize}
\end{itemize}
\end{browserWindow}


\begin{conceptBox}[Khẳng Định Đúng Cốt Lõi Về :target]
\textbf{"Khi người dùng click vào link nào thì CSS với lớp giả :target được chỉ định sẽ hoạt động tại đó."}

\par\vspace{4pt}
\textbf{Giải thích chi tiết}:
\begin{enumerate}[leftmargin=1.5em, itemsep=2pt, parsep=0pt]
    \item Khi người dùng click vào thẻ \codeinline{<a>} có \codeinline{href="\#id"} (Ví dụ \codeinline{<a href="\#thanbai">}), trình duyệt cuộn xuống phần tử có \codeinline{id="thanbai"} và phần tử đó trở thành \codeinline{:target}.
    \item CSS \codeinline{h2:target \{ color: blue; \}} chỉ áp dụng duy nhất cho phần tử được nhắm đến. Các phần tử \codeinline{<h2>} khác không bị ảnh hưởng.
    \item Khi nhấp sang link khác, \codeinline{:target} chuyển sang phần tử mới. Phần tử cũ lập tức mất hiệu ứng. Chỉ duy nhất một phần tử là \codeinline{:target} tại một thời điểm.
\end{enumerate}
\end{conceptBox}

\paragraph{Ví Dụ Thực Tế Nhiều Box \& Bảng Phân Tích Ảnh Hưởng:}

\begin{codeWindow}[Mã Nguồn Thử Nghiệm 3 Khối Box Với :target]
\begin{lstlisting}[language=HTML5]
<!-- Sample Text -->
<a href="#box1">Sample Text</a>
<a href="#box2">Sample Text</a>

<!-- Sample Text -->
<div id="box1" class="box">Box 1</div>
<div id="box2" class="box">Box 2</div>
<div id="box3" class="box">Box 3</div>

<!-- Sample Text -->
<style>
  .box:target {
    background-color: yellow;
    border: 2px solid red;
  }
</style>
\end{lstlisting}
\end{codeWindow}

\begin{prettyTableBox}[Bảng Phân Tích Tác Động Tương Tác Khi Click Link]
\rowcolors{2}{white}{primaryBlue!4}
\sffamily\small
\begin{tabularx}{\linewidth}{|>{\raggedright\arraybackslash \sffamily\bfseries}K{4.2cm}|>{\centering\arraybackslash\sffamily\bfseries}K{3.0cm}|Z|}
\hline
\rowcolor{primaryBlue} \thd{\sffamily Thao tác Click link} & \thd{\sffamily Phần tử bị ảnh hưởng} & \thd{\sffamily Phân tích tác động} \\
\hline
Click link \codeinline{href="\#box1"} & Chỉ \codeinline{div\#box1} & \codeinline{div\#box1} có nền vàng, viền đỏ. Các box khác bình thường. \\
\hline
Click link \codeinline{href="\#box2"} & Chỉ \codeinline{div\#box2} & \codeinline{div\#box2} có nền vàng, viền đỏ. \codeinline{div\#box1} mất hiệu ứng. \\
\hline
Không click hoặc click ngoài & Không box nào & Không có phần tử nào là \codeinline{:target}. \\
\hline
Box 3 (\codeinline{id="box3"}) & Không bao giờ & Vì không có liên kết nào trỏ tới \codeinline{\#box3}. \\
\hline
\end{tabularx}
\end{prettyTableBox}

\begin{browserWindow}[Kết Quả Hiển Thị Trực Quan Giao Diện Thử Nghiệm 3 Khối Box]
\sffamily\small
\begin{itemize}[leftmargin=1.5em, itemsep=3pt, parsep=0pt]
    \item \textbf{Tình huống A (Khi chưa nhấp liên kết nào -- URL gốc)}:
    \begin{itemize}[leftmargin=1.2em, itemsep=1pt]
        \item Khối \codeinline{Box 1}, \codeinline{Box 2}, \codeinline{Box 3} hiển thị chữ màu xám đen bình thường, không có nền vàng hay viền đỏ.
    \end{itemize}
    \item \textbf{Tình huống B (Khi người dùng click vào link "Xem Box 1" $\rightarrow$ URL: \codeinline{.../index.html\#box1})}:
    \begin{itemize}[leftmargin=1.2em, itemsep=1pt]
        \item \textbf{Box 1} (\codeinline{div\#box1}): Lập tức chuyển sang \textbf{màu nền vàng}, viền phát sáng \textbf{màu đỏ} (\codeinline{2px solid red}).
        \item \textbf{Box 2} \& \textbf{Box 3}: Giữ nguyên trạng thái bình thường.
    \end{itemize}
    \item \textbf{Tình huống C (Khi nhấp tiếp vào link "Xem Box 2" $\rightarrow$ URL: \codeinline{.../index.html\#box2})}:
    \begin{itemize}[leftmargin=1.2em, itemsep=1pt]
        \item \textbf{Box 1}: Tự động ngắt hiệu ứng \codeinline{:target}, trở về trạng thái bình thường.
        \item \textbf{Box 2} (\codeinline{div\#box2}): Trở thành target mới $\rightarrow$ Chuyển sang nền vàng, viền đỏ!
        \item \textbf{Box 3} (\codeinline{div\#box3}): Không bao giờ bị tác động vì không có thẻ \codeinline{<a>} nào chứa \codeinline{href="\#box3"}.
    \end{itemize}
\end{itemize}
\end{browserWindow}

\begin{tipBox}[Kết Luận Quan Trọng Về :target]
Chỉ phần tử nào có \codeinline{id} trùng với \codeinline{\#} trong URL mới được CSS \codeinline{:target} áp dụng. Các phần tử khác dù giống nhau về kiểu hay class nhưng không được target thì không bị ảnh hưởng!
\end{tipBox}

\paragraph{g. Chuyên Đề Thực Nghiệm Sâu: Ẩn/Hiện Phần Tử Con Với :hover \& Bộ Chọn Hậu Duệ}

Dòng mã này tạo ra một hộp màu vàng có nội dung và một biểu tượng ẩn đi. Khi bạn di chuột vào hộp, nền sẽ chuyển sang màu cam và biểu tượng xuất hiện.

\begin{codeWindow}[Mã Nguồn Cấu Trúc HTML]
\begin{lstlisting}[language=HTML5]
<div class="box">Content text
    <i>Icon</i>
</div>
\end{lstlisting}
\end{codeWindow}
\begin{itemize}[leftmargin=1.5em, itemsep=2pt, parsep=0pt]
    \item \codeinline{<div class="box">}: Tạo một khối chứa nội dung và icon.
    \item \codeinline{<i>Icon</i>}: Thẻ \codeinline{<i>} đại diện cho icon hoặc nội dung bổ sung. Ban đầu bị ẩn nhờ CSS.
\end{itemize}

\begin{codeWindow}[Mã Nguồn CSS Định Kiểu Ẩn/Hiện Khi Hover]
\begin{lstlisting}[language=CSS3]
.box {
    background: yellow;
    width: 300px;
    height: 200px;
    margin-top: 20px;
}
.box:hover {
    background: orange;
}
.box i {
    display: none;
}
.box:hover i {
    display: inline;
}
\end{lstlisting}
\end{codeWindow}

\begin{conceptBox}[Phân Tích Chi Tiết Các Thẻ CSS]
\begin{itemize}[leftmargin=1.5em, itemsep=3pt, parsep=0pt]
    \item \codeinline{background: yellow}: Đặt nền màu vàng cho \codeinline{.box}.
    \item \codeinline{width: 300px} và \codeinline{height: 200px}: Kích thước của hộp là 300x200 pixel.
    \item \codeinline{margin-top: 20px}: Tạo khoảng cách 20px phía trên hộp.
    \item \codeinline{.box:hover}: Khi di chuột vào \codeinline{.box}, nền đổi sang màu cam.
    \item \codeinline{.box i}: Ban đầu, thẻ \codeinline{<i>} không hiển thị nhờ \codeinline{display: none}.
    \item \codeinline{.box:hover i}: Khi di chuột vào \codeinline{.box}, thẻ \codeinline{<i>} hiện ra dưới dạng \codeinline{inline} (nằm trên cùng dòng với nội dung khác).
    \item Khi di chuột vào phần tử có class \codeinline{.box}, tất cả thẻ \codeinline{<i>} bên trong sẽ được hiển thị theo kiểu inline.
\end{itemize}
\end{conceptBox}

\begin{conceptBox}[Giải Thích Sâu Kỹ Thuật: Bộ Chọn Hậu Duệ (Descendant Selector)]
Câu lệnh này viết theo bộ chọn hậu duệ (Descendant Selector):
\begin{codeWindow}[Cú Pháp Bộ Chọn Hậu Duệ]
\begin{lstlisting}[language=CSS3]
.box:hover i {
    display: inline;
}
\end{lstlisting}
\end{codeWindow}
\begin{itemize}[leftmargin=1.5em, itemsep=2pt, parsep=0pt]
    \item \codeinline{.box}: Phần tử cha.
    \item \codeinline{i}: Phần tử con nằm bên trong \codeinline{.box}.
    \item \codeinline{:hover}: Kích hoạt khi di chuột vào phần tử cha \codeinline{.box}.
\end{itemize}
$\rightarrow$ \textbf{Bộ chọn hậu duệ chọn tất cả thẻ \codeinline{<i>} nằm bên trong \codeinline{.box}, bất kể mức lồng nhau bao nhiêu.}

Ví dụ, cả hai thẻ \codeinline{<i>} dưới đây đều bị ảnh hưởng khi hover vào \codeinline{.box}:
\begin{codeWindow}[Ví Dụ Mức Lồng Nhau Của Bộ Chọn Hậu Duệ]
\begin{lstlisting}[language=HTML5]
<div class="box">
    <i>Icon 1</i> <!-- Sample Text -->
    <div>
        <i>Icon 2</i> <!-- Sample Text -->
    </div>
</div>
\end{lstlisting}
\end{codeWindow}

\textbf{Nếu bạn muốn chỉ ảnh hưởng trực tiếp lên thẻ \codeinline{<i>} con trực tiếp thì phải dùng Child Selector}:
\begin{codeWindow}[Sử Dụng Child Selector (>)]
\begin{lstlisting}[language=CSS3]
.box:hover > i {
    display: inline;
}
\end{lstlisting}
\end{codeWindow}
\end{conceptBox}

\begin{browserWindow}[So Sánh Kết Quả Minh Họa Thực Nghiệm Cho 2 Bộ Chọn]
\begin{itemize}[leftmargin=1.5em, itemsep=4pt, parsep=0pt]
    \item \textbf{Trường hợp 1 -- Sử dụng Descendant Selector (\codeinline{.box:hover i})}:
    \begin{itemize}[leftmargin=1.2em, itemsep=2pt]
        \item \textbf{Ban đầu}: Hộp màu vàng (\codeinline{background: yellow}), chỉ hiển thị chữ "Nội dung Box". Cả \codeinline{Icon 1} và \codeinline{Icon 2} đều ẩn.
        \item \textbf{Khi di chuột vào .box}: Nền chuyển sang màu cam (\codeinline{background: orange}). Cả hai thẻ \codeinline{<i>} (\textbf{Icon 1} và \textbf{Icon 2}) nằm bên trong khối \codeinline{.box} (dù ở bất kỳ cấp lồng nhau nào) đều chuyển thành \codeinline{display: inline} và \textbf{xuất hiện cùng lúc}.
    \end{itemize}
    \par\vspace{4pt}
    \item \textbf{Trường hợp 2 -- Sử dụng Child Selector (\codeinline{.box:hover > i})}:
    \begin{itemize}[leftmargin=1.2em, itemsep=2pt]
        \item \textbf{Ban đầu}: Hộp màu vàng, tất cả Icon đều ẩn.
        \item \textbf{Khi di chuột vào .box}: Nền chuyển sang màu cam. \textbf{CHỈ CÓ} thẻ \codeinline{<i>Icon 1</i>} (con trực tiếp của \codeinline{.box}) mới xuất hiện; còn \codeinline{<i>Icon 2</i>} (nằm trong thẻ \codeinline{<div>} cháu) \textbf{vẫn tiếp tục bị ẩn}.
    \end{itemize}
\end{itemize}
\end{browserWindow}

\begin{tipBox}[Ứng Dụng Thực Tế Nâng Cao]
\begin{itemize}[leftmargin=1.5em, itemsep=2pt, parsep=0pt]
    \item \textbf{Hiệu ứng tooltip}: Khi hover vào thẻ, hiển thị mô tả hoặc ghi chú.
    \item \textbf{Thông báo nhỏ}: Ẩn/hiện thông báo hoặc biểu tượng trạng thái.
    \item \textbf{Menu thả xuống}: Hiện menu con khi hover vào nút hoặc thanh điều hướng.
\end{itemize}
\end{tipBox}

\paragraph{e. Đọc Thêm: Chuyên Đề Nâng Cao Về Pseudo-Classes Trạng Thái Tương Tác}

Để giúp bạn làm chủ hoàn toàn các trạng thái tương tác trong môi trường sản xuất thực tế (Production), dưới đây là các phân tích chuyên sâu, kinh nghiệm thực chiến và câu hỏi phỏng vấn tuyển dụng thường gặp:

\begin{conceptBox}[1. Xử Lý Lỗi "Sticky Hover" Trên Thiết Bị Cảm Ứng (Touch Devices)]
\begin{itemize}[leftmargin=1.5em, itemsep=3pt, parsep=0pt]
    \item \textbf{Vấn đề thực tế}: Trên smartphone hoặc tablet (màn hình cảm ứng), khi người dùng chạm ngón tay vào nút bấm, trình duyệt sẽ \textbf{dính chặt trạng thái :hover} ngay cả sau khi người dùng đã nhấc tay ra, gây ra lỗi giao diện rất khó chịu.
    \item \textbf{Giải pháp Modern CSS}: Sử dụng Media Query kiểm tra khả năng rơ chuột của phần cứng (\codeinline{@media (hover: hover)}):
\end{itemize}
\begin{codeWindow}[Tối Ưu :hover Chỉ Cho Thiết Bị Dùng Chuột]
\begin{lstlisting}[language=CSS3]
/* Sample Text */
@media (hover: hover) and (pointer: fine) {
    button:hover {
        background-color: #0056b3;
        transform: translateY(-2px);
    }
}
\end{lstlisting}
\end{codeWindow}

\begin{browserWindow}[Minh Họa Trực Quan Trải Nghiệm Trên Các Màn Hình]
\begin{itemize}[leftmargin=1.5em, itemsep=3pt, parsep=0pt]
    \item \textbf{Trên máy tính Desktop / Laptop (Chuột cứng)}: Điều kiện \codeinline{(hover: hover)} thỏa mãn. Khi di chuột qua nút $\rightarrow$ Nút nhích lên $2\text{px}$ và đổi màu xanh thẫm. Nhấc con trỏ ra $\rightarrow$ Nút trở lại vị trí ban đầu.
    \item \textbf{Trên Điện thoại / Máy tính bảng (Touchscreen)}: Điều kiện \codeinline{(hover: hover)} không thỏa mãn $\rightarrow$ Trình duyệt bỏ qua khối lệnh hover này. Người dùng chạm ngón tay vào nút $\rightarrow$ Thực hiện thao tác Click mà \textbf{KHÔNG bị dính màu hover (Sticky Hover Bug)}.
\end{itemize}
\end{browserWindow}
\end{conceptBox}

\begin{tipBox}[2. Bí Quyết Ghi Nhớ Quy Tắc LVHA Đạt Chuẩn Phỏng Vấn]
\begin{itemize}[leftmargin=1.5em, itemsep=3pt, parsep=0pt]
    \item \textbf{Mẹo ghi nhớ từ khóa}: Hãy nhớ cụm từ \textbf{"LoVe HAte"} (\textbf{L}ink $\rightarrow$ \textbf{V}isited $\rightarrow$ \textbf{H}over $\rightarrow$ \textbf{A}ctive) hoặc \textbf{"Lord Vader Has Arrived"}.
    \item \textbf{Giải thích kỹ thuật sâu}: Do cả 4 pseudo-classes này có \textbf{Độ ưu tiên Specificity bằng nhau} ($0,1,0$). Nếu bạn viết \codeinline{:link} sau \codeinline{:hover}, khi di chuột vào liên kết chưa truy cập, thuộc tính của \codeinline{:link} xuất hiện sau sẽ ghi đè và làm mất hiệu ứng của \codeinline{:hover}!
    \item \textbf{Mở rộng thứ tự khi có :focus (Quy tắc LVHFA)}:
    \par\vspace{3pt}
    \tcbox[on line, arc=3pt, colback=primaryBlue!6, colframe=primaryBlue!40, boxrule=0.8pt, boxsep=3pt, left=6pt, right=6pt]{%
        \ttfamily\color{primaryBlue}\bfseries :link $\rightarrow$ :visited $\rightarrow$ :hover $\rightarrow$ :focus $\rightarrow$ :active%
    }
\end{itemize}
\end{tipBox}

\begin{warnBox}[3. Tối Ưu Hiệu Năng Render (60 FPS Performance Notes)]
Khi định kiểu hiệu ứng \codeinline{:hover}, hãy tuân thủ nguyên tắc tối ưu GPU:
\begin{itemize}[leftmargin=1.5em, itemsep=2pt, parsep=0pt]
    \item \textbf{KHÔNG NÊN}: Thay đổi các thuộc tính kích thước như \codeinline{width}, \codeinline{height}, \codeinline{margin}, \codeinline{padding}, \codeinline{top}, \codeinline{left} khi hover vì sẽ ép trình duyệt tính toán lại bố cục toàn trang (\term{Reflow / Layout Shift}), gây giật lag.
    \item \textbf{NÊN DÙNG}: Thay đổi \codeinline{transform} (\codeinline{scale}, \codeinline{translate}) và \codeinline{opacity}. Các thuộc tính này được đẩy thẳng lên card màn hình xử lý (\term{GPU Hardware Acceleration}), đảm bảo mượt mà 60 khung hình/giây (60fps).
\end{itemize}
\end{warnBox}

\begin{tipBox}[4. Tiêu Chuẩn Trợ Năng Bàn Phím (Web Accessibility - a11y)]
Để trang web thân thiện với người khuyết tật di chuyển bằng bàn phím (phím \codeinline{Tab}):
\begin{itemize}[leftmargin=1.5em, itemsep=2pt, parsep=0pt]
    \item Luôn đi kèm \codeinline{:hover} với \codeinline{:focus} hoặc \codeinline{:focus-visible}.
    \item Ví dụ: \codeinline{button:hover, button:focus-visible \{ outline: 2px solid blue; \}} giúp người dùng duyệt bằng phím Tab nhìn thấy rõ tiêu điểm đang nằm ở đâu.
\end{itemize}
\end{tipBox}

\subsubsection{7.2. Nhóm Vị Trí Vị Thế Cấu Trúc Cây DOM (Structural Selectors)}

Cho phép chọn phần tử dựa trên thứ tự xuất hiện hoặc vị thế anh chị em trong cây HTML mà không cần gán class:

\begin{prettyTableBox}[Bảng Tra Cứu Structural Pseudo-Classes Trong Cây DOM]
\rowcolors{2}{white}{primaryBlue!4}
\sffamily\small
\begin{tabularx}{\linewidth}{|>{\raggedright\arraybackslash \ttfamily\color{primaryBlue}\bfseries}K{3.6cm}|>{\sffamily\small\bfseries}K{2.2cm}|>{\ttfamily\small}K{4.5cm}|Z|}
\hline
\rowcolor{primaryBlue} \thd{\sffamily Selector} & \thd{\sffamily Phạm vi} & \thd{\sffamily Ví dụ cú pháp} & \thd{\sffamily Chức năng \& Ý nghĩa} \\
\hline
:first-child & Tất cả phần tử & li:first-child \{ color: red; \} & Chọn phần tử là con đầu tiên của phần tử cha. \\
\hline
:last-child & Tất cả phần tử & li:last-child \{ color: blue; \} & Chọn phần tử là con cuối cùng của phần tử cha. \\
\hline
:nth-child(n) & Tất cả phần tử & li:nth-child(2) \{ color: purple; \} & Chọn phần tử con thứ n (truyền số, odd, even, 2n+1). \\
\hline
:nth-last-child(n) & Tất cả phần tử & li:nth-last-child(1) \{ color: teal; \} & Chọn phần tử con thứ n đếm ngược từ dưới lên. \\
\hline
:first-of-type & Tất cả phần tử & p:first-of-type \{ color: green; \} & Chọn phần tử đầu tiên trong số các phần tử cùng loại thẻ. \\
\hline
:last-of-type & Tất cả phần tử & p:last-of-type \{ color: orange; \} & Chọn phần tử cuối cùng trong số các phần tử cùng loại thẻ. \\
\hline
:nth-of-type(n) & Tất cả phần tử & p:nth-of-type(1) \{ font-weight: bold; \} & Chọn phần tử thứ n trong số các phần tử cùng loại thẻ. \\
\hline
:nth-last-of-type(n) & Tất cả phần tử & p:nth-last-of-type(1) \{ color: brown; \} & Chọn phần tử thứ n đếm ngược từ dưới lên trong số các phần tử cùng loại thẻ. \\
\hline
:only-child & Tất cả phần tử & div:only-child \{ border: 1px solid black; \} & Chọn phần tử nếu nó là con độc nhất của thẻ cha. \\
\hline
:only-of-type & Tất cả phần tử & p:only-of-type \{ font-style: italic; \} & Chọn phần tử nếu nó là con duy nhất thuộc loại thẻ đó trong thẻ cha. \\
\hline
:empty & Tất cả phần tử & div:empty \{ display: none; \} & Chọn phần tử hoàn toàn rỗng (không con, không chữ). \\
\hline
:root & Thẻ gốc Document & :root \{ --main-color: red; \} & Chọn phần tử gốc của cây DOM (thường là thẻ html). \\
\hline
\end{tabularx}
\end{prettyTableBox}

\paragraph{a. Mã nguồn ví dụ cấu trúc DOM đính kèm kết quả visual:}

\begin{codeWindow}[Mã CSS Định Kiểu Danh Sách Theo Vị Trí Vị Thế DOM]
\begin{lstlisting}[language=CSS3]
/* Sample Text */
li:first-child {
    color: green;
    font-weight: bold;
}

/* Sample Text */
li:nth-child(2) {
    color: red;
}
\end{lstlisting}
\end{codeWindow}

\begin{browserWindow}[Kết Quả Hiển Thị: Danh Sách Được Định Kiểu Theo Vị Trí]
\sffamily\small
\begin{itemize}[leftmargin=1.5em, itemsep=2pt, parsep=0pt]
    \item \textcolor{green!70!black}{\textbf{Mục 1 (li:first-child -- Màu xanh lá)}}
    \item \textcolor{red}{Mục 2 (li:nth-child(2) -- Màu đỏ)}
    \item Mục 3 (Mặc định)
\end{itemize}
\end{browserWindow}

\paragraph{b. Tổng Quan \& Bản Chất Cốt Lõi Của Pseudo-Classes Cấu Trúc}

\begin{conceptBox}[Định Nghĩa \& Bản Chất Của Pseudo-Classes Cấu Trúc]
\term{Structural Pseudo-Classes} (Các lớp giả cấu trúc) là nhóm bộ chọn đặc biệt trong CSS3 cho phép bạn chọn và trang trí các phần tử HTML dựa hoàn toàn vào \textbf{vị trí tương đối của chúng trong cây DOM} (Document Object Model) mà không cần phải chèn thêm bất kỳ class hay id thủ công nào vào mã HTML.

\par\vspace{4pt}
\textbf{ Ẩn dụ tư duy dễ nhớ:} Hãy hình dung cây DOM HTML như một danh sách học sinh đang xếp hàng. Thay vì phải dán nhãn tên từng người (gán class), bạn có thể ra lệnh cho CSS: \textit{``Tô màu đỏ cho học sinh đứng đầu hàng''}, \textit{``Tô màu xanh cho học sinh đứng cuối hàng''}, hoặc \textit{``Tô màu vàng cho các học sinh ở vị trí số lẻ (1, 3, 5...)''}. Đó chính là sức mạnh của Pseudo-Classes Cấu trúc!

\par\vspace{4pt}
\textbf{ 3 Lợi ích cốt lõi trong phát triển Web hiện đại:}
\begin{itemize}[leftmargin=1.5em, itemsep=3pt, parsep=0pt]
    \item \textbf{Giữ mã HTML sạch tuyệt đối}: Tránh việc phình to mã nguồn HTML do gán hàng loạt class lặp lại (ví dụ dán \codeinline{class="first-item"}, \codeinline{class="even-row"} vào từng thẻ).
    \item \textbf{Tự động hóa giao diện động}: Khi dữ liệu gửi từ Server/API thay đổi số lượng dòng (thêm/bớt sản phẩm), giao diện vẫn tự động áp dụng đúng quy luật mà không cần can thiệp CSS.
    \item \textbf{Tối ưu bảo trì mã nguồn}: Thay đổi quy luật thiết kế giao diện tại một nơi duy nhất trong file CSS.
\end{itemize}
\end{conceptBox}

\begin{prettyTableBox}[Bảng Phân Loại 10 Structural Pseudo-Classes Theo Cơ Chế Hoạt Động]
\rowcolors{2}{white}{primaryBlue!4}
\sffamily\small
\begin{tabularx}{\linewidth}{|>{\raggedright\arraybackslash \ttfamily\color{accentOrange}\bfseries\raggedright\arraybackslash}m{4.2cm}|Z|}
\hline
\rowcolor{primaryBlue} \thd{\sffamily Bộ chọn CSS} & \thd{\sffamily Nguyên lý chọn phần tử trong DOM} \\
\hline
\rowcolor{primaryBlue!15} \multicolumn{2}{|l|}{\bfseries\sffamily\color{primaryBlue} 1. Nhóm Vị Trí Tuyệt Đối (Đếm tất cả các phần tử con bất kể loại thẻ)} \\
\hline
:first-child & Chọn phần tử nếu nó là con đầu tiên tuyệt đối của cha. \\
\hline
:last-child & Chọn phần tử nếu nó là con cuối cùng tuyệt đối của cha. \\
\hline
:nth-child(n) & Chọn phần tử đứng ở vị trí thứ $n$ từ trên xuống. \\
\hline
:nth-last-child(n) & Chọn phần tử đứng ở vị trí thứ $n$ đếm từ dưới lên. \\
\hline
:only-child & Chọn phần tử nếu nó là con duy nhất tuyệt đối trong cha. \\
\hline
\rowcolor{primaryBlue!15} \multicolumn{2}{|l|}{\bfseries\sffamily\color{primaryBlue} 2. Nhóm Lọc Theo Kiểu Thẻ (Chỉ lọc và đếm các phần tử cùng kiểu thẻ)} \\
\hline
:first-of-type & Chọn phần tử cùng kiểu thẻ đầu tiên trong cha. \\
\hline
:last-of-type & Chọn phần tử cùng kiểu thẻ cuối cùng trong cha. \\
\hline
:nth-of-type(n) & Chọn phần tử cùng kiểu thẻ thứ $n$ từ trên xuống. \\
\hline
:nth-last-of-type(n) & Chọn phần tử cùng kiểu thẻ thứ $n$ đếm từ dưới lên. \\
\hline
:only-of-type & Chọn phần tử nếu nó là thẻ duy nhất thuộc loại đó trong cha. \\
\hline
\end{tabularx}
\end{prettyTableBox}

\paragraph{c. Chuyên Đề 1: Bộ Chọn Phần Tử Con Đầu Tiên (:first-child vs :first-of-type)}

Trong phát triển giao diện Web, việc chọn phần tử con xuất hiện đầu tiên trong một khối container (ví dụ: tạo kiểu làm nổi bật tab "Trang chủ" ở menu, tạo chữ hoa lớn cho đoạn văn đầu bài viết \term{Drop Cap}, hoặc xóa margin-top thừa của phần tử đỉnh) là một thao tác căn bản. CSS cung cấp hai pseudo-class selector cho mục đích này: \codeinline{:first-child} và \codeinline{:first-of-type}. Phần này sẽ phân tích chi tiết và tách biệt hai bộ chọn.

\begin{conceptBox}[Tổng Quan Nhanh: 2 Cơ Chế Chọn Phần Tử Đỉnh]
\begin{itemize}[leftmargin=1.5em, itemsep=3pt, parsep=0pt]
    \item \codeinline{:first-child}: Chọn phần tử nếu nó là \textbf{con đầu tiên tuyệt đối} (ở vị trí số 1 trong danh sách tất cả các con của cha).
    \item \codeinline{:first-of-type}: Chọn phần tử nếu nó là \textbf{con đầu tiên thuộc loại thẻ đó} trong phần tử cha (bỏ qua các thẻ khác loại đứng phía trước).
\end{itemize}
\end{conceptBox}

%% ===================================================================
%% PHẦN C1: :first-child -- PHÂN TÍCH CHUYÊN SÂU
%% ===================================================================

\subparagraph{c1. Phân Tích Chuyên Sâu :first-child (Con Đầu Tiên Tuyệt Đối)}

\codeinline{:first-child} là một pseudo-class selector trong CSS dùng để chọn phần tử con đứng ở vị trí đầu tiên tuyệt đối (vị trí $index = 1$) trong danh sách tất cả phần tử con của cha.

\begin{conceptBox}[Định Nghĩa Chi Tiết \& Cơ Chế Hoạt Động Của :first-child]
\codeinline{:first-child} chọn phần tử nếu nó nằm ở vị trí con đầu tiên tuyệt đối trong phần tử cha:
\begin{itemize}[leftmargin=1.5em, itemsep=2pt, parsep=0pt]
    \item Phần tử con đứng đầu tiên trong cây DOM có vị trí $index = 1$.
    \item Trình duyệt đếm tất cả các loại thẻ con (element nodes).
    \item Nếu phần tử ở vị trí 1 không khớp với thẻ chỉ định trong CSS $\rightarrow$ \textcolor{red}{\textbf{Bỏ qua, không tô màu!}}
\end{itemize}
\par\vspace{4pt}
\textbf{Tương đương logic:} \codeinline{selector:first-child} $\equiv$ \codeinline{selector:nth-child(1)}.
\end{conceptBox}

\begin{codeWindow}[Cú Pháp Cơ Bản Của :first-child]
\begin{lstlisting}[language=CSS3]
selector:first-child {
    /* CSS properties apply here */
}

/* Equivalent to: */
selector:nth-child(1) {
    /* Same result */
}
\end{lstlisting}
\end{codeWindow}

\begin{codeWindow}[Ví Dụ 1: Chọn Thẻ <li> Đầu Tiên Trong Danh Sách]
\begin{lstlisting}[language=HTML5]
<ul>
    <li>Item 1</li> <!-- 1st child of <ul> -->
    <li>Item 2</li>
    <li>Item 3</li>
</ul>

<style>
    /* Select the 1st <li> element inside list */
    li:first-child {
        color: red;
        font-weight: bold;
    }
</style>
\end{lstlisting}
\end{codeWindow}

\begin{browserWindow}[Kết Quả Hiển Thị Trực Quan Ví Dụ 1]
\sffamily\small
\begin{itemize}[leftmargin=1.5em, itemsep=2pt, parsep=0pt]
    \item \textcolor{red}{\textbf{Item 1}} {\scriptsize\textcolor{gray!60}{($\leftarrow$ li:first-child khớp -- Chữ màu đỏ, in đậm vì là phần tử con đầu tiên)}}
    \item Item 2 {\scriptsize\textcolor{gray!60}{(Mặc định)}}
    \item Item 3 {\scriptsize\textcolor{gray!60}{(Mặc định)}}
\end{itemize}
\par\vspace{3pt}
{\scriptsize\textcolor{gray!60}{(\textit{Giải thích: li:first-child chọn phần tử <li> đầu tiên trong danh sách. Kết quả "Item 1" sẽ có màu đỏ và chữ in đậm. Nếu không áp dụng :first-child thì tất cả các thẻ <li> nằm trong <ul> đều bị ảnh hưởng giống nhau.})}}
\end{browserWindow}

\begin{codeWindow}[Ví Dụ 2: Chọn Phần Tử Con Đầu Tiên Trong Container Cụ Thể]
\begin{lstlisting}[language=HTML5]
<div class="container">
    <p>Paragraph 1</p> <!-- 1st child of .container -->
    <p>Paragraph 2</p>
</div>

<style>
    /* Select 1st <p> child inside .container */
    .container p:first-child {
        color: blue;
        font-weight: bold;
    }
</style>
\end{lstlisting}
\end{codeWindow}

\begin{browserWindow}[Kết Quả Hiển Thị Ví Dụ 2: Container Cụ Thể]
\sffamily\small
\textcolor{primaryBlue}{\textbf{Paragraph 1}} {\scriptsize\textcolor{gray!60}{($\leftarrow$ p:first-child khớp -- Chữ xanh dương, in đậm)}} \par\vspace{2pt}
Paragraph 2 {\scriptsize\textcolor{gray!60}{(Mặc định)}} \par\vspace{3pt}
{\scriptsize\textcolor{gray!60}{(\textit{Giải thích: Chọn <p> đầu tiên trực tiếp trong .container và đổi màu chữ thành xanh dương. Chỉ áp dụng cho phần tử con đầu tiên, không ảnh hưởng các phần tử khác.})}}
\end{browserWindow}

\begin{codeWindow}[Ví Dụ 3: Phối Hợp :first-child Với Pseudo-Class Hover]
\begin{lstlisting}[language=CSS3]
/* Highlight 1st item on mouse hover */
li:first-child:hover {
    color: purple;
    text-decoration: underline;
    cursor: pointer;
}
\end{lstlisting}
\end{codeWindow}

\begin{tipBox}[Phối Hợp :first-child Với Hover]
\textbf{[Hiệu Ứng]:} Phần tử đầu tiên trong danh sách khi người dùng rê chuột (hover) vào sẽ tự động chuyển sang màu tím, xuất hiện đường gạch chân và đổi con trỏ chuột dạng pointer.
\end{tipBox}

\begin{codeWindow}[Ví Dụ 4: Tình Huống Thất Bại Phổ Biến Của :first-child]
\begin{lstlisting}[language=HTML5]
<!-- Mixed container: <h1> is at position 1 -->
<div class="card">
    <h1>Header Title</h1> <!-- 1st child of .card -->
    <p>First paragraph text</p>
</div>

<style>
    /* FAILS: 1st child is <h1>, NOT <p> */
    .card p:first-child { color: green; }
</style>
\end{lstlisting}
\end{codeWindow}

\begin{warnBox}[Phân Tích Chi Tiết: Tại Sao :first-child Thất Bại \& Thuật Toán 2 Bước Của Trình Duyệt]
\textbf{Cơ chế duyệt 2 bước của Trình duyệt khi xử lý \codeinline{p:first-child}:}
\begin{enumerate}[leftmargin=1.5em, itemsep=2pt, parsep=0pt]
    \item \textbf{Bước 1 (Đếm vị trí):} Trình duyệt tìm phần tử con nằm ở vị trí số 1 tuyệt đối bên trong phần tử cha.
    \item \textbf{Bước 2 (Kiểm tra thẻ):} Trình duyệt kiểm tra xem phần tử ở vị trí số 1 đó có đúng là thẻ \codeinline{<p>} hay không.
\end{enumerate}

\par\vspace{4pt}
\textbf{Tại sao \codeinline{p:first-child} thất bại trong DOM hỗn hợp?}
\begin{itemize}[leftmargin=1.5em, itemsep=3pt, parsep=0pt]
    \item Nếu vị trí số 1 là thẻ \codeinline{<h1>} (Tiêu đề), còn thẻ \codeinline{<p>} nằm ở vị trí số 2 $\rightarrow$ Trình duyệt thấy vị trí số 1 là \codeinline{<h1>} chứ không phải \codeinline{<p>} $\rightarrow$ \textcolor{red}{\textbf{Điều kiện SAI! Bỏ qua, không tô màu đoạn văn nào!}}
\end{itemize}

\par\vspace{4pt}
\textbf{Giải pháp hoàn hảo -- Sử dụng \codeinline{p:first-of-type}:}
\begin{itemize}[leftmargin=1.5em, itemsep=3pt, parsep=0pt]
    \item \codeinline{p:first-of-type} thay đổi quy trình: Trình duyệt sẽ \textbf{lọc tất cả các thẻ \codeinline{<p>} thành một nhóm riêng trước}, sau đó chọn thẻ \codeinline{<p>} đầu tiên trong nhóm đó (bất chấp phía trước có bao nhiêu thẻ \codeinline{<h1>} hay \codeinline{<div>}) $\rightarrow$ \textcolor{green!70!black}{\textbf{Tô màu xanh thành công 100\%!}}
\end{itemize}
\end{warnBox}

\begin{conceptBox}[4 Ứng Dụng Thực Tế Của :first-child]
\begin{itemize}[leftmargin=1.5em, itemsep=3pt, parsep=0pt]
    \item \textbf{Tạo kiểu đặc biệt cho phần tử đầu tiên}: Làm nổi bật tiêu đề đầu tiên hoặc mục đầu tiên trong danh sách.
    \item \textbf{Thiết kế Menu điều hướng}: Nhấn mạnh tab đầu tiên (ví dụ nút ``Trang chủ'') hoặc thêm khoảng cách/màu sắc đặc biệt.
    \item \textbf{Căn chỉnh Bố cục trang}: Áp dụng padding/margin riêng cho phần tử đầu tiên để căn chỉnh layout cân đối, không bị dính lề trên.
    \item \textbf{Danh sách sản phẩm \& Bài viết}: Đánh dấu sản phẩm đầu tiên hoặc bài viết mới nhất trong danh sách.
\end{itemize}
\end{conceptBox}

%% ===================================================================
%% PHẦN C2: :first-of-type -- PHÂN TÍCH CHUYÊN SÂU
%% ===================================================================

\subparagraph{c2. Phân Tích Chuyên Sâu :first-of-type (Con Đầu Tiên Theo Loại Thẻ)}

\codeinline{:first-of-type} là pseudo-class selector chọn phần tử con là \textbf{phần tử đầu tiên thuộc loại thẻ đó} (cùng tag name) trong phần tử cha, bất kể phía trước nó có bao nhiêu thẻ thuộc loại khác.

\begin{conceptBox}[Định Nghĩa Chi Tiết \& Cú Pháp :first-of-type]
\codeinline{:first-of-type} chọn phần tử đầu tiên thuộc loại (thẻ) đó trong các phần tử con của một phần tử cha.
\begin{itemize}[leftmargin=1.5em, itemsep=2pt, parsep=0pt]
    \item Trình duyệt gom tất cả các thẻ cùng tag name (ví dụ tất cả các thẻ \codeinline{<p>}).
    \item Chọn phần tử xuất hiện \textbf{đầu tiên} trong nhóm phần tử cùng loại đó.
    \item Các thẻ thuộc loại khác đứng phía trước hoàn toàn \textbf{bị bỏ qua} và không ảnh hưởng đến kết quả.
\end{itemize}
\par\vspace{4pt}
\textbf{Khác với \codeinline{:first-child}:}
\begin{itemize}[leftmargin=1.5em, itemsep=2pt, parsep=0pt]
    \item \codeinline{:first-child}: Chọn phần tử đầu tiên bất kể loại nào.
    \item \codeinline{:first-of-type}: Chọn phần tử đầu tiên thuộc loại đó (thẻ giống nhau).
\end{itemize}
\par\vspace{4pt}
\textbf{Tương đương logic:} \codeinline{selector:first-of-type} $\equiv$ \codeinline{selector:nth-of-type(1)}.
\end{conceptBox}

\begin{codeWindow}[Cú Pháp Cơ Bản Của :first-of-type]
\begin{lstlisting}[language=CSS3]
selector:first-of-type {
    /* CSS properties apply here */
}

/* Equivalent to: */
selector:nth-of-type(1) {
    /* Same result */
}
\end{lstlisting}
\end{codeWindow}

\begin{codeWindow}[Ví Dụ 1: Mã Nguồn Chuẩn Chọn Thẻ <p> Đầu Tiên Trong .box]
\begin{lstlisting}[language=HTML5]
<div class="box">
  <h2>Tieu de</h2>
  <p>Doan 1</p> <!-- First <p> in .box -->
  <p>Doan 2</p>
</div>

<style>
  p:first-of-type {
    color: red;
  }
</style>
\end{lstlisting}
\end{codeWindow}

\begin{browserWindow}[Kết Quả Hiển Thị Ví Dụ 1: .box p:first-of-type]
\sffamily\small
Tieu de {\scriptsize\textcolor{gray!60}{(Thẻ <h2> đứng trước -- Bị bỏ qua)}} \par\vspace{2pt}
\textcolor{red}{\textbf{Doan 1}} {\scriptsize\textcolor{gray!60}{($\leftarrow$ p:first-of-type KHỚP -- Tô màu đỏ vì là <p> đầu tiên trong các thẻ con của .box)}} \par\vspace{2pt}
Doan 2 {\scriptsize\textcolor{gray!60}{(Thẻ <p> thứ 2 -- Mặc định)}}
\end{browserWindow}

\begin{codeWindow}[Ví Dụ 2: Trường Hợp Có Thẻ <span> Chen Giữa Phía Trước]
\begin{lstlisting}[language=HTML5]
<div class="box">
  <h2>Tieu de</h2>
  <span>span</span> <!-- Sibling <span> before <p> -->
  <p>Doan 1</p> <!-- Still first-of-type among <p> tags -->
  <p>Doan 2</p>
</div>

<style>
  p:first-of-type {
    color: red;
  }
</style>
\end{lstlisting}
\end{codeWindow}

\begin{browserWindow}[Kết Quả Hiển Thị Ví Dụ 2: Thẻ <span> Chen Giữa]
\sffamily\small
Tieu de {\scriptsize\textcolor{gray!60}{(Thẻ <h2>)}} \par\vspace{2pt}
span {\scriptsize\textcolor{gray!60}{(Thẻ <span>)}} \par\vspace{2pt}
\textcolor{red}{\textbf{Doan 1}} {\scriptsize\textcolor{gray!60}{($\leftarrow$ p:first-of-type VẪN KHỚP -- Chọn được <p>Doan 1</p> vì là thẻ <p> đầu tiên, dù không phải con đầu tiên của .box)}} \par\vspace{2pt}
Doan 2 {\scriptsize\textcolor{gray!60}{(Thẻ <p> thứ 2)}}
\end{browserWindow}

\begin{codeWindow}[Ví Dụ 3: Mỗi Loại Thẻ Đều Có :first-of-type Riêng (Chọn Thẻ <span> Đầu Tiên)]
\begin{lstlisting}[language=HTML5]
<div>
  <h2>Tieu de</h2>
  <p>Doan van</p>
  <span>Span 1</span> <!-- 1st <span> in <div> -->
  <p>Doan van khac</p>
</div>

<style>
  span:first-of-type {
    color: blue;
    font-weight: bold;
  }
</style>
\end{lstlisting}
\end{codeWindow}

\begin{browserWindow}[Kết Quả Hiển Thị Ví Dụ 3: span:first-of-type]
\sffamily\small
Tieu de {\scriptsize\textcolor{gray!60}{(Thẻ <h2>)}} \par\vspace{2pt}
Doan van {\scriptsize\textcolor{gray!60}{(Thẻ <p> thứ 1)}} \par\vspace{2pt}
\textcolor{primaryBlue}{\textbf{Span 1}} {\scriptsize\textcolor{gray!60}{($\leftarrow$ span:first-of-type KHỚP -- Chỉ tác động đến span đầu tiên (trong số các span) trong div)}} \par\vspace{2pt}
Doan van khac {\scriptsize\textcolor{gray!60}{(Thẻ <p> thứ 2)}}
\end{browserWindow}

\begin{warnBox}[5 Lưu Ý Cốt Lõi Khi Sử Dụng :first-of-type]
\begin{enumerate}[leftmargin=1.5em, itemsep=3pt, parsep=0pt]
    \item \textbf{1. Là phần tử đầu tiên cùng loại trong cha:} Chỉ áp dụng nếu phần tử đó là loại đầu tiên (cùng tên thẻ: \codeinline{p}, \codeinline{li}, \codeinline{div}...) trong danh sách phần tử con.
    \item \textbf{2. Không cần phải là phần tử con đầu tiên tuyệt đối:} Dù có thẻ \codeinline{<span>} hay \codeinline{<h2>} đứng trước, \codeinline{p:first-of-type} vẫn hoạt động thành công.
    \item \textbf{3. Không hoạt động cho phần tử cùng loại thứ 2 trở đi:} Khi có \codeinline{<p>Doan 1</p>} và \codeinline{<p>Doan 2</p>}, chỉ \codeinline{<p>Doan 1</p>} được tô đỏ, \codeinline{<p>Doan 2</p>} bị bỏ qua.
    \item \textbf{4. Khác biệt cốt lõi với :first-child:} \codeinline{:first-child} yêu cầu là phần tử đầu tiên tuyệt đối của cha (bất kể loại gì); \codeinline{:first-of-type} yêu cầu là phần tử đầu tiên theo loại (thẻ) của cha.
    \item \textbf{5. Thường dùng trong danh sách nhiều loại phần tử con:} Giúp định kiểu độc lập từng loại thẻ trong giao diện hỗn hợp.
\end{enumerate}
\end{warnBox}

\begin{prettyTableBox}[Bảng Kiểm Tra Tính Đúng / Sai Của Các Lưu Ý Về :first-of-type]
\rowcolors{2}{white}{primaryBlue!4}
\sffamily\small
\begin{tabularx}{\linewidth}{|Z|>{\centering\arraybackslash\bfseries}m{2.2cm}|>{\raggedright\arraybackslash}m{4.2cm}|}
\hline
\rowcolor{primaryBlue} \thd{\sffamily Phát biểu lưu ý} & \thd{\sffamily Đúng / Sai} & \thd{\sffamily Ghi chú ngắn} \\
\hline
Chọn phần tử đầu tiên theo loại & \textcolor{green!70!black}{Đúng} & Đúng loại mới được áp dụng \\
\hline
Bị chặn nếu có anh em cùng loại đứng trước & \textcolor{green!70!black}{Đúng} & Chỉ 1 phần tử đầu tiên được áp dụng \\
\hline
Không quan tâm đến vị trí tổng thể trong DOM & \textcolor{green!70!black}{Đúng} & Miễn là loại đầu tiên trong nhóm \\
\hline
Là phần tử đầu tiên bất kể loại thẻ? & \textcolor{red}{\textbf{Sai}} & Đó là \codeinline{:first-child}, không phải \codeinline{:first-of-type} \\
\hline
\end{tabularx}
\end{prettyTableBox}

\begin{conceptBox}[4 Ứng Dụng Thực Tế Của :first-of-type]
\begin{itemize}[leftmargin=1.5em, itemsep=3pt, parsep=0pt]
    \item \textbf{Đoạn văn mở đầu bài viết (Lead Paragraph / Drop Cap)}: Định dạng cỡ chữ lớn hoặc chữ hoa hoa mỹ cho đoạn văn đầu tiên trong bài blog.
    \item \textbf{Form ô nhập liệu}: Căn chỉnh lề trên (\codeinline{margin-top: 0}) cho ô \codeinline{<input>} đầu tiên trong Form.
    \item \textbf{Ảnh đại diện bài viết}: Tự động định dạng kích thước hoặc border cho hình ảnh \codeinline{<img>} đầu tiên.
    \item \textbf{Hero section}: Làm nổi bật thẻ tiêu đề \codeinline{<h1>} xuất hiện đầu tiên trong trang.
\end{itemize}
\end{conceptBox}

%% ===================================================================
%% PHẦN C3: BẢNG SO SÁNH TỔNG HỢP :first-child vs :first-of-type
%% ===================================================================

\subparagraph{c3. Bảng So Sánh Chi Tiết :first-child vs :first-of-type}
\mbox{}\par\vspace{2pt}

\begin{codeWindow}[Ví Dụ Trực Quan So Sánh :first-child vs :first-of-type]
\begin{lstlisting}[language=HTML5]
<div class="demo">
    <h1>Title</h1>
    <p>Paragraph 1</p>
    <p>Paragraph 2</p>
</div>

<style>
    /* FAILS: 1st child of .demo is <h1>, not <p> */
    .demo p:first-child { color: red; }

    /* WORKS: Selects 1st <p> (Paragraph 1), ignoring <h1> */
    .demo p:first-of-type { color: green; font-weight: bold; }
</style>
\end{lstlisting}
\end{codeWindow}

\begin{browserWindow}[Kết Quả So Sánh Trong Trình Duyệt]
\sffamily\small
Title {\scriptsize\textcolor{gray!60}{(Thẻ <h1> -- Con đầu tiên thực tế)}} \par\vspace{2pt}
\textcolor{green!70!black}{\textbf{Paragraph 1}} {\scriptsize\textcolor{gray!60}{($\leftarrow$ p:first-of-type KHỚP -- Tô xanh lá vì là <p> đầu tiên)}} \par\vspace{2pt}
Paragraph 2 {\scriptsize\textcolor{gray!60}{(Thẻ <p> thứ 2)}} \par\vspace{4pt}
{\scriptsize\textcolor{gray!60}{\textit{Lưu ý: p:first-child KHÔNG khớp vì con đầu tiên tuyệt đối là <h1>, không phải <p>.}}}
\end{browserWindow}

\begin{prettyTableBox}[Bảng So Sánh Chi Tiết Cụ Thể: :first-child vs :first-of-type]
\rowcolors{2}{white}{primaryBlue!4}
\sffamily\small
\begin{tabularx}{\linewidth}{|>{\raggedright\arraybackslash\bfseries\color{primaryBlue}}m{3.2cm}|Z|Z|}
\hline
\rowcolor{primaryBlue} \thd{\sffamily Pseudo-class} & \thd{\sffamily Quan tâm đến vị trí?} & \thd{\sffamily Quan tâm đến loại thẻ?} \\
\hline
:first-child & \textbf{Có} (Phải là con đầu tiên tuyệt đối của cha) & \textbf{Có} (Phải trùng khớp với thẻ chỉ định) \\
\hline
:first-of-type & \textbf{Không} (Không cần đứng đầu danh sách con) & \textbf{Có} (Là phần tử cùng loại đầu tiên trong cha) \\
\hline
\end{tabularx}
\end{prettyTableBox}

\begin{prettyTableBox}[Bảng So Sánh Toàn Diện :first-child vs :first-of-type]
\rowcolors{2}{white}{primaryBlue!4}
\sffamily\small
\begin{tabularx}{\linewidth}{|>{\raggedright\arraybackslash\bfseries\color{primaryBlue}}m{3.2cm}|Z|Z|}
\hline
\rowcolor{primaryBlue} \thd{\sffamily Tiêu chí so sánh} & \thd{\sffamily :first-child} & \thd{\sffamily :first-of-type} \\
\hline
Cơ chế đếm & Kiểm tra phần tử đứng ở \textbf{vị trí số 1 tuyệt đối} & Lọc riêng nhóm phần tử \textbf{cùng loại thẻ} và chọn phần tử đầu \\
\hline
Ảnh hưởng bởi thẻ khác loại? & \textcolor{red}{\textbf{Có}} (Bất kỳ thẻ khác loại nào đứng ở vị trí 1 đều làm thất bại) & \textcolor{green!70!black}{\textbf{Không}} (Thẻ khác loại đứng phía trước bị bỏ qua) \\
\hline
Tương đương logic & \codeinline{:nth-child(1)} & \codeinline{:nth-of-type(1)} \\
\hline
Đối ứng vị trí đáy & \codeinline{:last-child} & \codeinline{:last-of-type} \\
\hline
Mức độ nghiêm ngặt & Nghiêm ngặt (Cha phải chỉ có thẻ này ở vị trí đỉnh) & Linh hoạt (Phù hợp với container chứa nhiều thẻ hỗn hợp) \\
\hline
Khi nào nên dùng? & Khi danh sách thuần nhất (ví dụ toàn \codeinline{<li>}) & Khi container có cấu trúc hỗn hợp nhiều loại thẻ \\
\hline
\end{tabularx}
\end{prettyTableBox}

\begin{tipBox}[Mẹo Ghi Nhớ Nhanh Phân Biệt :first-child vs :first-of-type]
\rowcolors{2}{white}{primaryBlue!4}
\sffamily\small
\begin{tabularx}{\linewidth}{|>{\raggedright\arraybackslash \bfseries\color{primaryBlue}}m{3.0cm}|Z|Z|Z|}
\hline
\rowcolor{primaryBlue} \thd{\sffamily Bộ chọn} & \thd{\sffamily Vị trí đếm} & \thd{\sffamily Quan tâm loại thẻ?} & \thd{\sffamily Dùng khi nào?} \\
\hline
:first-child & Đầu tiên trong mọi con & \textcolor{red}{\textbf{Không}} & Style phần tử đầu tiên bất kỳ \\
\hline
:first-of-type & Đầu tiên trong cùng loại & \textcolor{green!70!black}{\textbf{Có}} & Style phần tử đầu tiên của loại thẻ cụ thể \\
\hline
\end{tabularx}
\end{tipBox}

\paragraph{d. Chuyên Đề 2: Bộ Chọn Phần Tử Con Cuối Cùng (:last-child vs :last-of-type)}

Trong thiết kế giao diện Web, việc chọn phần tử con cuối cùng trong một khối container (ví dụ: xóa đường gạch kẻ dưới của mục menu cuối cùng, căn chỉnh khoảng cách cho nút bấm cuối form, hoặc định dạng đoạn văn cuối bài) là nhu cầu diễn ra vô cùng thường xuyên. CSS cung cấp hai pseudo-class selector cho mục đích này: \codeinline{:last-child} và \codeinline{:last-of-type}. Phần này sẽ phân tích chi tiết và tách biệt hai bộ chọn.

\begin{conceptBox}[Tổng Quan Nhanh: 2 Cơ Chế Chọn Phần Tử Đáy]
\begin{itemize}[leftmargin=1.5em, itemsep=3pt, parsep=0pt]
    \item \codeinline{:last-child}: Chọn phần tử nếu nó là \textbf{con cuối cùng tuyệt đối} (ở vị trí cuối cùng trong danh sách tất cả các con của cha).
    \item \codeinline{:last-of-type}: Chọn phần tử nếu nó là \textbf{con cuối cùng thuộc loại thẻ đó} trong phần tử cha (bỏ qua các thẻ khác loại đứng phía sau).
\end{itemize}
\end{conceptBox}

%% ===================================================================
%% PHẦN D1: :last-child -- PHÂN TÍCH CHUYÊN SÂU
%% ===================================================================

\subparagraph{d1. Phân Tích Chuyên Sâu :last-child (Con Cuối Cùng Tuyệt Đối)}

\codeinline{:last-child} là một pseudo-class selector trong CSS dùng để chọn phần tử con đứng ở vị trí cuối cùng tuyệt đối trong danh sách tất cả phần tử con của cha.

\begin{conceptBox}[Định Nghĩa Chi Tiết \& Cơ Chế Hoạt Động Của :last-child]
\codeinline{:last-child} khớp với một phần tử khi và chỉ khi hai điều kiện sau được thỏa mãn đồng thời:
\begin{itemize}[leftmargin=1.5em, itemsep=2pt, parsep=0pt]
    \item Phần tử nằm ở \textbf{vị trí cuối cùng tuyệt đối} trong danh sách các phần tử con của cha.
    \item Phần tử đó khớp với bộ chọn \codeinline{selector} phía trước.
    \item \textbf{Lưu ý quan trọng:} Trình duyệt kiểm tra vị trí cuối cùng trong DOM trước. Nếu phần tử đứng cuối cùng là một loại thẻ khác với \codeinline{selector} $\rightarrow$ \textcolor{red}{\textbf{Bỏ qua, không áp dụng CSS!}}
\end{itemize}
\par\vspace{4pt}
\textbf{Tương đương logic:} \codeinline{selector:last-child} $\equiv$ \codeinline{selector:nth-last-child(1)}.
\end{conceptBox}

\begin{codeWindow}[Cú Pháp Cơ Bản Của :last-child]
\begin{lstlisting}[language=CSS3]
selector:last-child {
    /* CSS properties apply here */
}

/* Equivalent to: */
selector:nth-last-child(1) {
    /* Same result */
}
\end{lstlisting}
\end{codeWindow}

\begin{codeWindow}[Ví Dụ 1: Chọn Thẻ <li> Cuối Cùng Trong Danh Sách]
\begin{lstlisting}[language=HTML5]
<ul>
    <li>Item 1</li>
    <li>Item 2</li>
    <li>Item 3</li> <!-- Last child of <ul> -->
</ul>

<style>
    /* Select the last <li> element inside list */
    li:last-child {
        color: blue;
        font-weight: bold;
    }
</style>
\end{lstlisting}
\end{codeWindow}

\begin{browserWindow}[Kết Quả Hiển Thị Trực Quan Ví Dụ 1]
\sffamily\small
\begin{itemize}[leftmargin=1.5em, itemsep=2pt, parsep=0pt]
    \item Item 1 {\scriptsize\textcolor{gray!60}{(Mặc định)}}
    \item Item 2 {\scriptsize\textcolor{gray!60}{(Mặc định)}}
    \item \textcolor{primaryBlue}{\textbf{Item 3}} {\scriptsize\textcolor{gray!60}{($\leftarrow$ li:last-child khớp -- Chữ xanh, in đậm vì là phần tử cuối cùng trong <ul>)}}
\end{itemize}
\end{browserWindow}

\begin{codeWindow}[Ví Dụ 2: Chọn Đoạn Văn Cuối Cùng Trong Container Cụ Thể]
\begin{lstlisting}[language=HTML5]
<div class="container">
    <p>Paragraph 1</p>
    <p>Paragraph 2</p> <!-- Last child of .container -->
</div>

<style>
    /* Select last <p> child inside .container */
    .container p:last-child {
        color: red;
        font-weight: bold;
    }
</style>
\end{lstlisting}
\end{codeWindow}

\begin{browserWindow}[Kết Quả Hiển Thị Ví Dụ 2]
\sffamily\small
Paragraph 1 {\scriptsize\textcolor{gray!60}{(Mặc định)}} \par\vspace{2pt}
\textcolor{red}{\textbf{Paragraph 2}} {\scriptsize\textcolor{gray!60}{($\leftarrow$ p:last-child khớp -- Chữ đỏ vì là <p> đứng cuối cùng trong .container)}}
\end{browserWindow}

\begin{codeWindow}[Ví Dụ 3: Phối Hợp :last-child Với Pseudo-Class Hover]
\begin{lstlisting}[language=CSS3]
/* Highlight last item on mouse hover */
li:last-child:hover {
    color: purple;
    text-decoration: underline;
    cursor: pointer;
}
\end{lstlisting}
\end{codeWindow}

\begin{tipBox}[Phối Hợp :last-child Với Hover]
\textbf{[Hiệu Ứng]:} Phần tử cuối cùng trong danh sách khi di chuột vào (hover) sẽ tự động chuyển sang màu tím, xuất hiện đường gạch chân và đổi con trỏ chuột dạng pointer.
\end{tipBox}

\begin{codeWindow}[Ví Dụ 4: Tình Huống Thất Bại Phổ Biến Của :last-child]
\begin{lstlisting}[language=HTML5]
<!-- Mixed container: <span> is the actual last child -->
<div class="card">
    <p>First paragraph text</p>
    <span>Span element text</span> <!-- Actual last child -->
</div>

<style>
    /* FAILS: Last child of .card is <span>, NOT <p> */
    .card p:last-child {
        color: red;
    }
</style>
\end{lstlisting}
\end{codeWindow}

\begin{browserWindow}[Kết Quả Ví Dụ 4: :last-child Thất Bại Trong DOM Hỗn Hợp]
\sffamily\small
First paragraph text {\scriptsize\textcolor{gray!60}{(KHÔNG đỏ -- Thẻ <p> không phải là con cuối cùng trong .card)}} \par\vspace{2pt}
Span element text {\scriptsize\textcolor{gray!60}{(Thẻ <span> -- Con cuối cùng thực tế trong .card)}} \par\vspace{4pt}
{\scriptsize\textcolor{gray!60}{\textit{Giải thích: .card p:last-child kiểm tra con cuối cùng trong .card (là <span>). Vì <span> $\neq$ <p> nên không có thẻ <p> nào được áp dụng CSS.}}}
\end{browserWindow}

\begin{warnBox}[Các Lưu Ý Quan Trọng Khi Sử Dụng :last-child]
\textbf{1. Phải là con cuối cùng tuyệt đối:} Nếu đứng sau thẻ bạn chọn còn bất kỳ thẻ anh em nào khác (dù khác loại), \codeinline{:last-child} sẽ thất bại.
\par\vspace{4pt}
\textbf{2. Text node không bị đếm:} Khoảng trắng, xuống dòng trong HTML không bị tính là phần tử con đứng cuối.
\par\vspace{4pt}
\textbf{3. Giải pháp khi có thẻ khác loại đứng sau:} Chuyển sang sử dụng \codeinline{:last-of-type}.
\end{warnBox}

\begin{tipBox}[Mẹo Ghi Nhớ :last-child]
\textbf{Hình dung:} \codeinline{:last-child} giống như "người xếp hàng ở vị trí cuối cùng tuyệt đối". Nếu người đứng cuối cùng không mặc áo đỏ (không phải loại thẻ cần tìm), bộ chọn \codeinline{p:last-child} sẽ không chọn được ai.
\end{tipBox}

\begin{conceptBox}[5 Ứng Dụng Thực Tế Của :last-child]
\begin{enumerate}[leftmargin=1.5em, itemsep=3pt, parsep=0pt]
    \item \textbf{Tạo viền / Cách điệu phần tử cuối}: Bỏ đường viền dưới (\codeinline{border-bottom: none}) hoặc đường gạch kẻ ngang ở mục cuối cùng trong danh sách.
    \item \textbf{Điều hướng Menu}: Nhấn mạnh tab cuối cùng hoặc thêm biểu tượng đặc biệt (nút CTA) cho menu navigation.
    \item \textbf{Bố cục giao diện}: Điều chỉnh margin/padding dưới (\codeinline{margin-bottom: 0}) của phần tử cuối để tránh làm phình lề dưới của container.
    \item \textbf{Danh sách sản phẩm \& Footer}: Làm nổi bật phần tử sản phẩm cuối hoặc tạo kiểu đặc biệt cho liên kết cuối trong Footer.
    \item \textbf{Form nhập liệu}: Điều chỉnh khoảng cách hoặc căn chỉnh nút Submit ở cuối biểu mẫu.
\end{enumerate}
\end{conceptBox}

%% ===================================================================
%% PHẦN D2: :last-of-type -- PHÂN TÍCH CHUYÊN SÂU
%% ===================================================================

\subparagraph{d2. Phân Tích Chuyên Sâu :last-of-type (Con Cuối Cùng Theo Loại Thẻ)}

\codeinline{:last-of-type} là pseudo-class selector dùng để chọn phần tử con cuối cùng của một loại (tag name) trong một phần tử cha.

\begin{conceptBox}[Định Nghĩa Chi Tiết \& Cơ Chế Hoạt Động Của :last-of-type]
\codeinline{:last-of-type} lọc và hoạt động theo các bước:
\begin{itemize}[leftmargin=1.5em, itemsep=2pt, parsep=0pt]
    \item Trình duyệt lọc riêng danh sách tất cả phần tử con cùng loại thẻ (tag name) trong cha.
    \item Chọn phần tử xuất hiện \textbf{cuối cùng} trong nhóm phần tử cùng loại đó.
    \item Các phần tử khác loại thẻ đứng phía sau hoàn toàn \textbf{bị bỏ qua} và không ảnh hưởng đến kết quả.
\end{itemize}
\par\vspace{4pt}
\textbf{So sánh với \codeinline{:last-child}:}
\begin{itemize}[leftmargin=1.5em, itemsep=2pt, parsep=0pt]
    \item \codeinline{:last-child}: Phần tử con cuối cùng trong tất cả các con (Không quan tâm loại thẻ).
    \item \codeinline{:last-of-type}: Phần tử con cuối cùng thuộc một loại (thẻ) (Có quan tâm loại thẻ).
\end{itemize}
\par\vspace{4pt}
\textbf{Tương đương logic:} \codeinline{selector:last-of-type} $\equiv$ \codeinline{selector:nth-last-of-type(1)}.
\end{conceptBox}

\begin{codeWindow}[Cú Pháp Cơ Bản Của :last-of-type]
\begin{lstlisting}[language=CSS3]
selector:last-of-type {
    /* style apply */
}

/* Equivalent to: */
selector:nth-last-of-type(1) {
    /* Same result */
}
\end{lstlisting}
\end{codeWindow}

\begin{codeWindow}[Ví Dụ 1: Mã Nguồn Cơ Bản Với Thẻ <span> Đứng Sau]
\begin{lstlisting}[language=HTML5]
<div>
  <h2>Tieu de</h2>
  <p>Doan 1</p>
  <p>Doan 2</p> <!-- Applied: Last <p> tag -->
  <span>Chu thich</span>
</div>

<style>
  p:last-of-type {
    color: red;
  }
</style>
\end{lstlisting}
\end{codeWindow}

\begin{browserWindow}[Kết Quả Hiển Thị Ví Dụ 1: p:last-of-type]
\sffamily\small
Tieu de {\scriptsize\textcolor{gray!60}{(Thẻ <h2>)}} \par\vspace{2pt}
Doan 1 {\scriptsize\textcolor{gray!60}{(Thẻ <p> thứ 1)}} \par\vspace{2pt}
\textcolor{red}{\textbf{Doan 2}} {\scriptsize\textcolor{gray!60}{($\leftarrow$ p:last-of-type KHỚP -- Doan 2 bị tô đỏ vì nó là phần tử <p> cuối cùng trong cha div)}} \par\vspace{2pt}
Chu thich {\scriptsize\textcolor{gray!60}{(Thẻ <span> đứng sau -- Bị bỏ qua)}}
\end{browserWindow}

\begin{codeWindow}[Ví Dụ 2: Ví Dụ Phân Biệt Chi Tiết Với :last-child]
\begin{lstlisting}[language=HTML5]
<div>
  <p>Doan 1</p>
  <p>Doan 2</p>
  <span>Ghi chu</span>
</div>

<style>
  /* FAILS: p is not the absolute last child */
  p:last-child { color: red; }

  /* SUCCEEDS: Last <p> in <p> group */
  p:last-of-type { color: green; font-weight: bold; }
</style>
\end{lstlisting}
\end{codeWindow}

\begin{browserWindow}[Kết Quả Hiển Thị Ví Dụ 2: Phân Biệt Với :last-child]
\sffamily\small
Doan 1 {\scriptsize\textcolor{gray!60}{(Thẻ <p> thứ 1)}} \par\vspace{2pt}
\textcolor{green!70!black}{\textbf{Doan 2}} {\scriptsize\textcolor{gray!60}{($\leftarrow$ p:last-of-type KHỚP -- p:last-child THẤT BẠI vì p không phải con cuối cùng, còn p:last-of-type THÀNH CÔNG vì p là cuối cùng trong nhóm p)}} \par\vspace{2pt}
Ghi chu {\scriptsize\textcolor{gray!60}{(Thẻ <span> -- Con cuối cùng thực tế)}}
\end{browserWindow}

\begin{warnBox}[Các Lưu Ý Quan Trọng Khi Sử Dụng :last-of-type]
\textbf{1. Áp dụng cho thẻ cùng loại:} Không cần là phần tử con cuối cùng tuyệt đối trong container.
\par\vspace{4pt}
\textbf{2. Trường hợp phần tử đơn lẻ:} Nếu container chỉ có đúng 1 phần tử thuộc loại thẻ đó, \codeinline{:last-of-type} sẽ chọn đúng phần tử đó (trùng kết quả với \codeinline{:first-of-type} và \codeinline{:only-of-type}).
\end{warnBox}

\begin{tipBox}[Mẹo Ghi Nhớ :last-of-type]
\textbf{Hình dung:} \codeinline{:last-of-type} giống như "người cuối cùng của từng nhóm". Nếu có nhóm nam và nhóm nữ, \codeinline{:last-of-type} chọn người đứng cuối của nhóm nam và người đứng cuối của nhóm nữ.
\end{tipBox}

\begin{conceptBox}[4 Ứng Dụng Thực Tế Của :last-of-type]
\begin{itemize}[leftmargin=1.5em, itemsep=3pt, parsep=0pt]
    \item \textbf{Đoạn văn bài viết}: Xóa \codeinline{margin-bottom} cho đoạn văn \codeinline{<p>} cuối bài viết blog để lề bottom ôm sát wrapper.
    \item \textbf{Chữ ký / Nguồn tác giả}: Định dạng kiểu chữ nghiêng hoặc chèn đường kẻ phía trên thẻ \codeinline{<p>} cuối cùng.
    \item \textbf{Form ô nhập liệu}: Reset khoảng cách lề cho ô \codeinline{<input>} cuối cùng trước nút bấm.
    \item \textbf{Gallery ảnh bài viết}: Bo góc tròn hoặc thêm border đặc biệt cho hình ảnh \codeinline{<img>} cuối cùng.
\end{itemize}
\end{conceptBox}

%% ===================================================================
%% PHẦN D3: BẢNG SO SÁNH TỔNG HỢP :last-child vs :last-of-type
%% ===================================================================

\subparagraph{d3. Bảng So Sánh Chi Tiết :last-child vs :last-of-type}
\mbox{}\par\vspace{2pt}

\begin{prettyTableBox}[Bảng So Sánh Cốt Lõi: :last-child vs :last-of-type]
\rowcolors{2}{white}{primaryBlue!4}
\sffamily\small
\begin{tabularx}{\linewidth}{|>{\raggedright\arraybackslash\bfseries\color{primaryBlue}}m{3.2cm}|Z|Z|}
\hline
\rowcolor{primaryBlue} \thd{\sffamily Selector} & \thd{\sffamily Ý nghĩa} & \thd{\sffamily Quan tâm đến loại thẻ?} \\
\hline
:last-child & Phần tử con cuối cùng trong tất cả các con & \textcolor{red}{\textbf{Không}} \\
\hline
:last-of-type & Phần tử con cuối cùng thuộc một loại (thẻ) & \textcolor{green!70!black}{\textbf{Có}} \\
\hline
\end{tabularx}
\end{prettyTableBox}

\begin{prettyTableBox}[Bảng Tổng Kết Quy Tắc Tra Cứu :last-child \& :last-of-type]
\rowcolors{2}{white}{primaryBlue!4}
\sffamily\small
\begin{tabularx}{\linewidth}{|>{\raggedright\arraybackslash\bfseries\color{primaryBlue}}m{3.0cm}|Z|Z|>{\raggedright\arraybackslash}m{4.0cm}|}
\hline
\rowcolor{primaryBlue} \thd{\sffamily Thuộc tính} & \thd{\sffamily Vị trí} & \thd{\sffamily Quan tâm loại?} & \thd{\sffamily Ví dụ dùng} \\
\hline
:last-child & Con cuối cùng bất kỳ & \textcolor{red}{\textbf{Không}} & Kiểm tra phần tử kết thúc tuyệt đối \\
\hline
:last-of-type & Con cuối cùng trong nhóm cùng loại & \textcolor{green!70!black}{\textbf{Có}} & Kiểm tra phần tử cuối cùng của một loại thẻ cụ thể \\
\hline
\end{tabularx}
\end{prettyTableBox}

\begin{prettyTableBox}[Bảng So Sánh Toàn Diện :last-child vs :last-of-type]
\rowcolors{2}{white}{primaryBlue!4}
\sffamily\small
\begin{tabularx}{\linewidth}{|>{\raggedright\arraybackslash\bfseries\color{primaryBlue}}m{3.2cm}|Z|Z|}
\hline
\rowcolor{primaryBlue} \thd{\sffamily Tiêu chí so sánh} & \thd{\sffamily :last-child} & \thd{\sffamily :last-of-type} \\
\hline
Cơ chế đếm & Kiểm tra phần tử đứng ở \textbf{vị trí cuối cùng tuyệt đối} & Lọc riêng nhóm phần tử \textbf{cùng loại thẻ} và chọn phần tử cuối \\
\hline
Ảnh hưởng bởi thẻ khác loại? & \textcolor{red}{\textbf{Có}} (Bất kỳ thẻ khác loại nào đứng cuối đều làm thất bại) & \textcolor{green!70!black}{\textbf{Không}} (Thẻ khác loại đứng phía sau bị bỏ qua) \\
\hline
Tương đương logic & \codeinline{:nth-last-child(1)} & \codeinline{:nth-last-of-type(1)} \\
\hline
Đối ứng vị trí đầu & \codeinline{:first-child} & \codeinline{:first-of-type} \\
\hline
Mức độ nghiêm ngặt & Nghiêm ngặt (Cha phải chỉ có thẻ này ở vị trí đáy) & Linh hoạt (Phù hợp với container chứa nhiều thẻ hỗn hợp) \\
\hline
Khi nào nên dùng? & Khi danh sách thuần nhất (ví dụ toàn \codeinline{<li>}) & Khi container có cấu trúc hỗn hợp nhiều loại thẻ \\
\hline
\end{tabularx}
\end{prettyTableBox}

\begin{tipBox}[Tóm Kết 4 Bộ Chọn Đầu \& Cuối Cơ Bản]
\begin{itemize}[leftmargin=1.5em, itemsep=3pt, parsep=0pt]
    \item \codeinline{:first-child}: Con đầu tiên tuyệt đối.
    \item \codeinline{:first-of-type}: Con đầu tiên theo loại thẻ.
    \item \codeinline{:last-child}: Con cuối cùng tuyệt đối.
    \item \codeinline{:last-of-type}: Con cuối cùng theo loại thẻ.
\end{itemize}
\end{tipBox}

\paragraph{e. Chuyên Đề 3: Bộ Chọn Vị Trí Vạn Năng (:nth-child(n) vs :nth-of-type(n))}

\codeinline{:nth-child(n)} và \codeinline{:nth-of-type(n)} là hai pseudo-class selector mạnh mẽ nhất trong CSS3 để chọn các phần tử dựa trên vị trí thứ tự của chúng trong danh sách con của cha. Chúng cho phép lập trình viên định kiểu theo chỉ số cố định, theo chu kỳ lặp số học, hoặc theo vị trí chẵn/lẻ. Phần này sẽ phân tích tách biệt và chuyên sâu từng bộ chọn.

\begin{conceptBox}[Tổng Quan Nhanh: 2 Cơ Chế Đếm Vị Trí Từ Trên Xuống]
\begin{itemize}[leftmargin=1.5em, itemsep=3pt, parsep=0pt]
    \item \codeinline{:nth-child(an + b)}: Đếm vị trí \textbf{từ trên xuống dưới} (từ con đầu tiên = 1), đếm \textbf{tất cả phần tử con} không phân biệt loại thẻ.
    \item \codeinline{:nth-of-type(an + b)}: Đếm vị trí \textbf{từ trên xuống dưới}, nhưng \textbf{chỉ đếm các phần tử cùng loại thẻ} (cùng tag name).
\end{itemize}
\end{conceptBox}

%% ===================================================================
%% PHẦN E1: :nth-child(an + b) -- PHÂN TÍCH CHUYÊN SÂU
%% ===================================================================

\subparagraph{e1. Phân Tích Chuyên Sâu :nth-child(an + b) (Đếm Vị Trí Tuyệt Đối Trong DOM)}

\codeinline{:nth-child(n)} là một pseudo-class selector dùng để chọn phần tử con dựa trên vị trí thứ tự tuyệt đối của nó trong danh sách tất cả các phần tử con của cha (đếm từ trên xuống).

\begin{conceptBox}[Định Nghĩa Chi Tiết \& Cơ Chế Đếm Của :nth-child(n)]
\codeinline{:nth-child(an + b)} chọn phần tử con theo thứ tự chỉ định $n$ hoặc công thức đại số $an + b$:
\begin{itemize}[leftmargin=1.5em, itemsep=2pt, parsep=0pt]
    \item Phần tử con đầu tiên trong phần tử cha có chỉ số $index = 1$.
    \item Trình duyệt đếm \textbf{tất cả phần tử con} (element nodes) mà không phân biệt loại thẻ (\codeinline{<div>}, \codeinline{<p>}, \codeinline{<span>}, \codeinline{<h2>}...).
    \item \textbf{Lưu ý cốt lõi:} \codeinline{:nth-child(an + b)} \textbf{không chọn dựa trên loại thẻ}, mà chọn theo vị trí thứ tự tuyệt đối của phần tử con trong cây DOM.
    \item Nếu một phần tử ở đúng vị trí tính ra từ công thức $an + b$ và khớp với thẻ được chỉ định, nó sẽ bị áp dụng CSS. Ngược lại, dù cùng loại thẻ nhưng sai vị trí DOM $\rightarrow$ không bị ảnh hưởng.
\end{itemize}
\end{conceptBox}

\begin{codeWindow}[Cú Pháp Cơ Bản Của :nth-child(n)]
\begin{lstlisting}[language=CSS3]
/* Standard nth-child syntax */
selector:nth-child(n) {
    /* CSS properties apply here */
}

/* Formula syntax */
selector:nth-child(an + b) {
    /* CSS properties apply here */
}
\end{lstlisting}
\end{codeWindow}

\begin{conceptBox}[Giải Thích Thành Phần Cú Pháp :nth-child(an + b)]
\begin{itemize}[leftmargin=1.5em, itemsep=2pt, parsep=0pt]
    \item \codeinline{selector}: Là phần tử bạn muốn chọn (ví dụ \codeinline{li}, \codeinline{p}, \codeinline{tr}).
    \item \codeinline{a}: \textbf{Bước nhảy} (khoảng cách giữa các phần tử sau mỗi lần lặp).
    \item \codeinline{n}: \textbf{Biến lặp index} (Số nguyên không âm $n = 0, 1, 2, 3, 4, \dots$).
    \item \codeinline{b}: \textbf{Vị trí bắt đầu} (Offset ban đầu).
\end{itemize}
\end{conceptBox}

\begin{codeWindow}[Ví Dụ 1: Minh Họa Vị Trí Đếm DOM Của :nth-child(3)]
\begin{lstlisting}[language=HTML5]
<div class="container">
  <h2>Header Title</h2>   <!-- Child position 1 -->
  <p>Paragraph 1</p>      <!-- Child position 2 -->
  <p>Paragraph 2</p>      <!-- Child position 3 -->
  <span>Span Note</span>   <!-- Child position 4 -->
  <p>Paragraph 3</p>      <!-- Child position 5 -->
</div>

<style>
  /* Select 3rd child IF it is a <p> tag */
  p:nth-child(3) {
    color: red;
    font-weight: bold;
  }
</style>
\end{lstlisting}
\end{codeWindow}

\begin{browserWindow}[Kết Quả Hiển Thị Chi Tiết Ví Dụ 1: p:nth-child(3)]
\sffamily\small
Header Title {\scriptsize\textcolor{gray!60}{(Child vị trí 1 -- Thẻ <h2>)}} \par\vspace{2pt}
Paragraph 1 {\scriptsize\textcolor{gray!60}{(Child vị trí 2 -- Thẻ <p>, không bị tác động)}} \par\vspace{2pt}
\textcolor{red}{\textbf{Paragraph 2}} {\scriptsize\textcolor{gray!60}{($\leftarrow$ p:nth-child(3) KHỚP -- Vừa là <p> vừa là child vị trí thứ 3 trong DOM)}} \par\vspace{2pt}
Span Note {\scriptsize\textcolor{gray!60}{(Child vị trí 4 -- Thẻ <span>)}} \par\vspace{2pt}
Paragraph 3 {\scriptsize\textcolor{gray!60}{(Child vị trí 5 -- Thẻ <p>, không bị tác động)}} \par\vspace{4pt}
{\scriptsize\textcolor{gray!60}{(\textit{Giải thích: <p>Paragraph 2</p> bị tác động vì nó vừa là thẻ <p> vừa là phần tử con thứ 3 đúng yêu cầu. <p>Paragraph 1</p> và <p>Paragraph 3</p> không bị tác động vì vị trí đếm trong DOM không khớp 3.})}}
\end{browserWindow}

\begin{conceptBox}[Công Thức Tổng Quát: nth-child(an + b)]
Công thức tổng quát đại số có dạng:
\[ \text{\texttt{nth-child(}} a n + b \text{\texttt{)}} \]
\par\vspace{4pt}
\textbf{Phân tích chi tiết công thức $3n + 1$:}
\begin{itemize}[leftmargin=1.5em, itemsep=1pt, parsep=0pt]
    \item $a = 3$: Nhảy cách 3 phần tử mỗi lần.
    \item $b = 1$: Bắt đầu xuất phát từ phần tử thứ 1.
\end{itemize}
\end{conceptBox}

\begin{prettyTableBox}[Bảng Thử Thay Số n Vào Công Thức 3n + 1 Qua Các Bước Lặp]
\rowcolors{2}{white}{primaryBlue!4}
\sffamily\small
\begin{tabularx}{\linewidth}{|>{\centering\arraybackslash\bfseries}C{1.8cm}|>{\centering\arraybackslash\ttfamily\color{accentOrange}\bfseries}C{3.8cm}|Z|}
\hline
\rowcolor{primaryBlue} \thd{\sffamily Giá trị $n$} & \thd{\sffamily Tính toán biểu thức $3n + 1$} & \thd{\sffamily Vị trí phần tử trong DOM} \\
\hline
$n = 0$ & 3(0) + 1 = 1 & \textbf{\color{primaryBlue}Vị trí 1} \quad {\scriptsize\textcolor{gray!60}{(Phần tử con đầu tiên)}} \\
\hline
$n = 1$ & 3(1) + 1 = 4 & \textbf{\color{primaryBlue}Vị trí 4} \quad {\scriptsize\textcolor{gray!60}{(Nhảy tiếp 3 bước)}} \\
\hline
$n = 2$ & 3(2) + 1 = 7 & \textbf{\color{primaryBlue}Vị trí 7} \quad {\scriptsize\textcolor{gray!60}{(Nhảy tiếp 3 bước)}} \\
\hline
$n = 3$ & 3(3) + 1 = 10 & \textbf{\color{primaryBlue}Vị trí 10} \quad {\scriptsize\textcolor{gray!60}{(Nhảy tiếp 3 bước)}} \\
\hline
$n = 4$ & 3(4) + 1 = 13 & \textbf{\color{primaryBlue}Vị trí 13} \quad {\scriptsize\textcolor{gray!60}{(Nhảy tiếp 3 bước)}} \\
\hline
\end{tabularx}
\end{prettyTableBox}

\begin{tipBox}[Mô Hình Tư Duy Trực Quan Về Công Thức 3n + 1]
\begin{itemize}[leftmargin=1.5em, itemsep=3pt, parsep=0pt]
    \item \codeinline{3n} tạo ra bước nhảy 3 phần tử một lần.
    \item \codeinline{+1} đẩy điểm bắt đầu xuất phát về phần tử số 1.
    \item \textbf{Luồng nhảy trực quan:} Điểm xuất phát từ phần tử 1 $\rightarrow$ nhảy 3 bước đến 4 $\rightarrow$ nhảy 3 bước đến 7 $\rightarrow$ 10 $\rightarrow$ 13 $\dots$
\end{itemize}
\end{tipBox}

\begin{codeWindow}[Mã Nguồn Minh Họa Trực Quan Công Thức 3n + 1 Trên Danh Sách 10 Phần Tử]
\begin{lstlisting}[language=HTML5]
<ul>
    <li>Item 1</li>
    <li>Item 2</li>
    <li>Item 3</li>
    <li>Item 4</li>
    <li>Item 5</li>
    <li>Item 6</li>
    <li>Item 7</li>
    <li>Item 8</li>
    <li>Item 9</li>
    <li>Item 10</li>
</ul>

<style>
    /* Select items at positions 1, 4, 7, 10 */
    li:nth-child(3n + 1) {
        color: red;
        font-weight: bold;
    }
</style>
\end{lstlisting}
\end{codeWindow}

\begin{browserWindow}[Kết Quả Hiển Thị 10 Items Với 3n + 1]
\sffamily\small
\begin{itemize}[leftmargin=1.5em, itemsep=2pt, parsep=0pt]
    \item \textcolor{red}{\textbf{Item 1}} {\scriptsize\textcolor{gray!60}{($\rightarrow$ Đỏ, in đậm (Vị trí 1))}}
    \item Item 2 {\scriptsize\textcolor{gray!60}{(Mặc định)}}
    \item Item 3 {\scriptsize\textcolor{gray!60}{(Mặc định)}}
    \item \textcolor{red}{\textbf{Item 4}} {\scriptsize\textcolor{gray!60}{($\rightarrow$ Đỏ, in đậm (Vị trí 4))}}
    \item Item 5 {\scriptsize\textcolor{gray!60}{(Mặc định)}}
    \item Item 6 {\scriptsize\textcolor{gray!60}{(Mặc định)}}
    \item \textcolor{red}{\textbf{Item 7}} {\scriptsize\textcolor{gray!60}{($\rightarrow$ Đỏ, in đậm (Vị trí 7))}}
    \item Item 8 {\scriptsize\textcolor{gray!60}{(Mặc định)}}
    \item Item 9 {\scriptsize\textcolor{gray!60}{(Mặc định)}}
    \item \textcolor{red}{\textbf{Item 10}} {\scriptsize\textcolor{gray!60}{($\rightarrow$ Đỏ, in đậm (Vị trí 10))}}
\end{itemize}
\end{browserWindow}

\begin{codeWindow}[Ví Dụ Chi Tiết Chọn Phần Tử Lẻ (odd) Và Chẵn (even)]
\begin{lstlisting}[language=HTML5]
<style>
    /* Odd items get light gray background */
    li:nth-child(odd) {
        background-color: #f0f0f0;
    }
    /* Even items get darker gray background */
    li:nth-child(even) {
        background-color: #d3d3d3;
    }
</style>
\end{lstlisting}
\end{codeWindow}

\begin{browserWindow}[Kết Quả Hiển Thị Xen Kẽ Lẻ / Chẵn]
\sffamily\small
\begin{itemize}[leftmargin=1.5em, itemsep=3pt, parsep=0pt]
    \item \colorbox{gray!20}{\textbf{Item 1, 3, 5...}} {\scriptsize\textcolor{gray!60}{($\leftarrow$ Có màu nền xám nhạt \#f0f0f0)}}
    \item Item 2, 4, 6... {\scriptsize\textcolor{gray!60}{(Màu nền trắng mặc định)}}
\end{itemize}
\end{browserWindow}

\begin{codeWindow}[Ví Dụ Chọn Với Công Thức Bội Số (3n) \& Công Thức Offset (n + 4)]
\begin{lstlisting}[language=CSS3]
/* Example 1: Select multiples of 3 (Item 3, 6, 9...) */
li:nth-child(3n) {
    color: blue;
    font-weight: bold;
}

/* Example 2: Select all items after item 3 (Item 4, 5, 6...) */
li:nth-child(n + 4) {
    text-decoration: underline;
}
\end{lstlisting}
\end{codeWindow}

\begin{browserWindow}[Kết Quả Hiển Thị Công Thức 3n Và n + 4]
\sffamily\small
\begin{itemize}[leftmargin=1.5em, itemsep=2pt, parsep=0pt]
    \item \textcolor{primaryBlue}{\textbf{Item 3}} {\scriptsize\textcolor{gray!60}{($\rightarrow$ li:nth-child(3n) chọn màu xanh, in đậm)}}
    \item \underline{Item 4, 5, 6...} {\scriptsize\textcolor{gray!60}{($\rightarrow$ li:nth-child(n + 4) gạch chân từ phần tử 4 trở đi)}}
\end{itemize}
\end{browserWindow}

\begin{warnBox}[Phân Tích Chi Tiết: Các Trường Hợp :nth-child(an + b) KHÔNG Hoạt Động Như Kỳ Vọng]
Mặc dù \codeinline{:nth-child(an + b)} (bao gồm \codeinline{odd}, \codeinline{even}, \codeinline{n}, \codeinline{2n+1}, \codeinline{-n+3}, v.v.) rất mạnh trong việc chọn phần tử theo vị trí, nhưng có những trường hợp cụ thể nó \textbf{không thỏa mãn hoặc không hoạt động như bạn kỳ vọng}:

\par\vspace{5pt}
\textbf{1. Khi phần tử không nằm ở đúng vị trí tính từ đầu:}
\begin{itemize}[leftmargin=1.5em, itemsep=2pt, parsep=0pt]
    \item Bạn nghĩ \codeinline{p:nth-child(1)} sẽ tô đỏ đoạn \codeinline{<p>Paragraph 1</p>} vì nó là đoạn văn đầu tiên.
    \item \textcolor{red}{\textbf{SAI!}} Vì thẻ \codeinline{<p>} đứng sau thẻ \codeinline{<h2>} nên nó nằm ở vị trí con thứ 2 trong DOM. Do đó \codeinline{nth-child(1)} không áp dụng cho \codeinline{p}.
    \item \textbf{Cách sửa:} Dùng \codeinline{p:first-of-type \{ color: red; \}} hoặc \codeinline{p:nth-of-type(1) \{ color: red; \}}.
\end{itemize}

\par\vspace{5pt}
\textbf{2. Số lượng phần tử không đủ để khớp với công thức:}
\begin{itemize}[leftmargin=1.5em, itemsep=2pt, parsep=0pt]
    \item Bạn viết \codeinline{li:nth-child(3n)} với mong muốn chọn các phần tử ở vị trí chia hết cho 3 trong danh sách \codeinline{<ul>}.
    \item \textcolor{red}{\textbf{SAI!}} Nếu danh sách chỉ có 2 thẻ \codeinline{<li>}, \textbf{không có gì xảy ra} vì công thức $3n$ chỉ khớp tại các vị trí 3, 6, 9...
    \item \textbf{Cách hiểu đúng:} Công thức $3n$ chỉ tác động khi tồn tại phần tử ở vị trí index 3, 6, 9...
\end{itemize}

\par\vspace{5pt}
\textbf{3. Phối hợp odd, even sai đối tượng trong DOM hỗn hợp:}
\begin{itemize}[leftmargin=1.5em, itemsep=2pt, parsep=0pt]
    \item Bạn viết \codeinline{p:nth-child(odd)} với mong muốn tô đỏ thẻ \codeinline{<p>} ở vị trí lẻ.
    \item \textcolor{red}{\textbf{SAI!}} Nếu cấu trúc DOM xen kẽ \codeinline{<span>1</span>}, \codeinline{<p>2</p>}, \codeinline{<span>3</span>}, \codeinline{<p>4</p>}:
    \begin{itemize}[leftmargin=1em]
        \item \codeinline{<span>} là con vị trí 1 $\rightarrow$ odd (lẻ)
        \item \codeinline{<p>} là con vị trí 2 $\rightarrow$ even (chẵn)
        \item \codeinline{<span>} là con vị trí 3 $\rightarrow$ odd (lẻ)
        \item \codeinline{<p>} là con vị trí 4 $\rightarrow$ even (chẵn)
    \end{itemize}
    Do thẻ \codeinline{<p>} nằm ở vị trí 2 và 4 (chẵn), bộ chọn \codeinline{p:nth-child(odd)} tìm thẻ \codeinline{<p>} ở vị trí lẻ nhưng không có thẻ \codeinline{<p>} nào ở vị trí lẻ $\rightarrow$ \textbf{Không có thẻ p nào được tô màu!}
    \item \textbf{Cách đúng:} Dùng \codeinline{p:nth-of-type(odd) \{ color: red; \}} để chỉ đếm các thẻ \codeinline{<p>} ở vị trí lẻ trong tập hợp các thẻ \codeinline{<p>}.
\end{itemize}
\end{warnBox}

\begin{conceptBox}[Ứng Dụng Thực Tế Của :nth-child(n)]
\begin{enumerate}[leftmargin=1.5em, itemsep=3pt, parsep=0pt]
    \item \textbf{Danh sách sản phẩm}: Tô màu xen kẽ cho từng dòng sản phẩm, làm nổi bật sản phẩm theo nhóm.
    \item \textbf{Bảng dữ liệu (Tables)}: Làm nổi bật từng dòng hoặc cột theo thứ tự lẻ/chẵn (\term{Zebra Striping}).
    \item \textbf{Menu hoặc thanh điều hướng}: Thêm hiệu ứng đặc biệt cho nút đầu tiên, nút cuối hoặc các nút theo quy luật.
    \item \textbf{Thẻ bài viết hoặc ảnh (Grid Gallery)}: Chọn bài viết/ảnh theo quy luật để tạo bố cục lưới đẹp mắt, phân trang linh hoạt.
\end{enumerate}
\end{conceptBox}

%% ===================================================================
%% PHẦN E2: :nth-of-type(an + b) -- PHÂN TÍCH CHUYÊN SÂU
%% ===================================================================

\subparagraph{e2. Phân Tích Chuyên Sâu :nth-of-type(an + b) (Đếm Vị Trí Theo Loại Thẻ)}

\codeinline{:nth-of-type(an + b)} là pseudo-class selector chọn phần tử dựa trên vị trí thứ tự của nó khi \textbf{chỉ lọc đếm riêng tập hợp các phần tử cùng loại thẻ} (cùng tag name) trong phần tử cha (đếm từ trên xuống).

\begin{conceptBox}[Định Nghĩa Chi Tiết \& Cơ Chế Đếm Của :nth-of-type(n)]
Khác với \codeinline{:nth-child(n)}, \codeinline{:nth-of-type(n)} \textbf{bỏ qua các phần tử khác loại thẻ} và tạo một danh sách đếm độc lập:
\begin{itemize}[leftmargin=1.5em, itemsep=2pt, parsep=0pt]
    \item Trình duyệt gom tất cả các thẻ cùng tag name (ví dụ tất cả các thẻ \codeinline{<p>}).
    \item Thẻ \codeinline{<p>} đầu tiên xuất hiện trong container có chỉ số $index = 1$ (trong nhóm thẻ \codeinline{<p>}).
    \item Thẻ \codeinline{<p>} thứ hai xuất hiện trong container có chỉ số $index = 2$ (trong nhóm thẻ \codeinline{<p>}).
    \item Các thẻ khác loại (\codeinline{<h2>}, \codeinline{<span>}, \codeinline{<div>}...) nằm chen giữa hoàn toàn \textbf{không làm gián đoạn hay tăng chỉ số đếm} của nhóm thẻ \codeinline{<p>}.
\end{itemize}
\end{conceptBox}

\begin{codeWindow}[Cú Pháp Cơ Bản Của :nth-of-type(n)]
\begin{lstlisting}[language=CSS3]
/* Select specific tag by index in its tag group */
selector:nth-of-type(n) {
    /* CSS properties apply here */
}

/* Formula syntax for specific tag group */
selector:nth-of-type(an + b) {
    /* CSS properties apply here */
}
\end{lstlisting}
\end{codeWindow}

\begin{codeWindow}[Mã Nguồn Minh Họa Phân Biệt Chi Tiết: .container :nth-child(3) vs .container p:nth-of-type(3)]
\begin{lstlisting}[language=HTML5]
<div class="container">
  <p>Doan 1</p>        <!-- p thu 1 (STT DOM: 1) -->
  <span>Ghi chu 1</span> <!-- span thu 1 (STT DOM: 2) -->
  <p>Doan 2</p>        <!-- p thu 2 (STT DOM: 3) -->
  <p>Doan 3</p>        <!-- p thu 3 (STT DOM: 4) -->
  <span>Ghi chu 2</span> <!-- span thu 2 (STT DOM: 5) -->
</div>

<style>
  /* Select 3rd child overall in .container regardless of tag */
  .container :nth-child(3) {
    color: red;
  }

  /* Select 3rd <p> tag specifically among all <p> siblings */
  .container p:nth-of-type(3) {
    background-color: yellow;
  }
</style>
\end{lstlisting}
\end{codeWindow}

\begin{prettyTableBox}[Bảng Thống Kê Thứ Tự DOM Chi Tiết Cho .container :nth-child(3)]
\rowcolors{2}{white}{primaryBlue!4}
\sffamily\small
\begin{tabularx}{\linewidth}{|>{\centering\arraybackslash\bfseries}m{2.0cm}|Z|>{\raggedright\arraybackslash\bfseries\color{primaryBlue}}m{4.5cm}|}
\hline
\rowcolor{primaryBlue} \thd{\sffamily STT DOM} & \thd{\sffamily Phần tử HTML con} & \thd{\sffamily Trạng thái chọn với :nth-child(3)} \\
\hline
1 & \codeinline{<p>Doan 1</p>} & Bỏ qua (Vị trí 1) \\
\hline
2 & \codeinline{<span>Ghi chu 1</span>} & Bỏ qua (Vị trí 2) \\
\hline
3 & \codeinline{<p>Doan 2</p>} & \textcolor{red}{\textbf{KHỚP :nth-child(3)}} ($\leftarrow$ Tô màu chữ đỏ) \\
\hline
4 & \codeinline{<p>Doan 3</p>} & Bỏ qua (Vị trí 4) \\
\hline
5 & \codeinline{<span>Ghi chu 2</span>} & Bỏ qua (Vị trí 5) \\
\hline
\end{tabularx}
\end{prettyTableBox}

\begin{prettyTableBox}[Bảng Thống Kê Thứ Tự Nhóm Thẻ <p> Cho .container p:nth-of-type(3)]
\rowcolors{2}{white}{primaryBlue!4}
\sffamily\small
\begin{tabularx}{\linewidth}{|Z|>{\centering\arraybackslash\bfseries\color{accentOrange}}m{3.5cm}|>{\raggedright\arraybackslash\bfseries\color{green!70!black}}m{4.5cm}|}
\hline
\rowcolor{primaryBlue} \thd{\sffamily Thẻ <p> trong danh sách} & \thd{\sffamily Chỉ số :nth-of-type} & \thd{\sffamily Trạng thái chọn với p:nth-of-type(3)} \\
\hline
\codeinline{<p>Doan 1</p>} & \codeinline{p:nth-of-type(1)} & Bỏ qua (Chỉ số 1 trong nhóm \codeinline{<p>}) \\
\hline
\codeinline{<p>Doan 2</p>} & \codeinline{p:nth-of-type(2)} & Bỏ qua (Chỉ số 2 trong nhóm \codeinline{<p>}) \\
\hline
\codeinline{<p>Doan 3</p>} & \codeinline{p:nth-of-type(3)} & \textcolor{green!70!black}{\textbf{KHỚP p:nth-of-type(3)}} ($\leftarrow$ Tô màu nền vàng) \\
\hline
\end{tabularx}
\end{prettyTableBox}

\begin{browserWindow}[Kết Quả Hiển Thị Phân Biệt Chi Tiết Trong Trình Duyệt]
\sffamily\small
Doan 1 {\scriptsize\textcolor{gray!60}{(Thẻ <p> thứ 1 -- STT DOM 1)}} \par\vspace{2pt}
Ghi chu 1 {\scriptsize\textcolor{gray!60}{(Thẻ <span> thứ 1 -- STT DOM 2)}} \par\vspace{2pt}
\textcolor{red}{\textbf{Doan 2}} {\scriptsize\textcolor{gray!60}{($\leftarrow$ .container :nth-child(3) KHỚP -- Con thứ 3 tổng thể của .container)}} \par\vspace{2pt}
\colorbox{yellow!30}{\textbf{Doan 3}} {\scriptsize\textcolor{gray!60}{($\leftarrow$ .container p:nth-of-type(3) KHỚP -- Thẻ <p> thứ 3 trong các <p> con của .container)}} \par\vspace{2pt}
Ghi chu 2 {\scriptsize\textcolor{gray!60}{(Thẻ <span> thứ 2 -- STT DOM 5)}}
\end{browserWindow}

\begin{prettyTableBox}[Bảng Đánh Giá Kết Luận Chuẩn Xác Chọn Phần Tử]
\rowcolors{2}{white}{primaryBlue!4}
\sffamily\small
\begin{tabularx}{\linewidth}{|>{\raggedright\arraybackslash \ttfamily\color{accentOrange}\bfseries}m{4.0cm}|Z|>{\bfseries\color{green!70!black}}m{3.5cm}|}
\hline
\rowcolor{primaryBlue} \thd{\sffamily Selector} & \thd{\sffamily Tiêu chí chọn phần tử} & \thd{\sffamily Kết quả đúng chọn được} \\
\hline
:nth-child(3) & Con thứ 3 tổng thể trong container & \codeinline{<p>Doan 2</p>} \\
\hline
p:nth-of-type(3) & Thẻ \codeinline{<p>} thứ 3 trong nhóm \codeinline{<p>} con & \codeinline{<p>Doan 3</p>} \\
\hline
\end{tabularx}
\end{prettyTableBox}

\begin{warnBox}[Các Lưu Ý Quan Trọng Khi Sử Dụng :nth-of-type(n)]
\textbf{1. Đếm độc lập theo tag name:} Mỗi loại thẻ trong container có một bộ đếm chỉ số độc lập.
\par\vspace{4pt}
\textbf{2. An toàn hơn trong giao diện phức tạp:} Khi làm việc với HTML động có thẻ tiêu đề, quảng cáo hoặc icon chen giữa các bài viết, \codeinline{:nth-of-type(n)} là lựa chọn an toàn tuyệt đối.
\end{warnBox}

\begin{tipBox}[Mẹo Ghi Nhớ :nth-of-type(n)]
\textbf{Công thức tương đương:} \codeinline{p:nth-of-type(1)} $\equiv$ \codeinline{p:first-of-type} (đoạn văn đầu tiên).
\end{tipBox}

\begin{conceptBox}[Ứng Dụng Thực Tế Của :nth-of-type(n)]
\begin{itemize}[leftmargin=1.5em, itemsep=3pt, parsep=0pt]
    \item \textbf{Trang tin tức / Blog}: Định kiểu cho đoạn văn đầu tiên (\term{Drop cap}) hoặc các đoạn văn lẻ/chẵn khi có ảnh chen giữa.
    \item \textbf{Form phức tạp}: Chọn các ô \codeinline{<input type="text">} chẵn/lẻ để tô màu nền khi có thẻ \codeinline{<label>} chen giữa.
    \item \textbf{Gallery ảnh}: Style cho ảnh \codeinline{<img>} thứ N trong bài viết.
\end{itemize}
\end{conceptBox}

%% ===================================================================
%% PHẦN E3: BẢNG SO SÁNH TỔNG HỢP :nth-child(n) vs :nth-of-type(n)
%% ===================================================================

\subparagraph{e3. Bảng So Sánh Chi Tiết :nth-child(n) vs :nth-of-type(n)}
\mbox{}\par\vspace{2pt}

\begin{codeWindow}[Mã Nguồn Minh Họa 3 Trường Hợp Thất Bại \& Cách Khắc Phục]
\begin{lstlisting}[language=HTML5]
<!-- Case 1: Header at position 1 -->
<div class="case1">
  <h2>Header Title</h2>
  <p>Paragraph 1</p> <!-- Child 2 -->
</div>
<style>
  /* FAILS: p is child 2, not child 1 */
  .case1 p:nth-child(1) { color: red; }
  /* FIX: Select 1st <p> specifically */
  .case1 p:nth-of-type(1) { color: green; }
</style>

<!-- Case 2: Insufficient items -->
<ul>
  <li>Item 1</li>
  <li>Item 2</li>
</ul>
<style>
  /* FAILS: List only has 2 items, 3n matches position 3, 6, 9 */
  li:nth-child(3n) { color: blue; }
</style>

<!-- Case 3: Mixed odd/even -->
<div class="case3">
  <span>Span 1</span>  <!-- Child 1 (odd) -->
  <p>Paragraph 2</p>  <!-- Child 2 (even) -->
  <span>Span 3</span>  <!-- Child 3 (odd) -->
  <p>Paragraph 4</p>  <!-- Child 4 (even) -->
</div>
<style>
  /* FAILS: <p> tags are at even positions 2 \& 4, no <p> at odd position */
  .case3 p:nth-child(odd) { color: red; }
  /* FIX: Select odd <p> within <p> elements only */
  .case3 p:nth-of-type(odd) { color: green; }
</style>
\end{lstlisting}
\end{codeWindow}

\begin{browserWindow}[Kết Quả Trực Quan 3 Trường Hợp Thất Bại \& Cách Khắc Phục]
\sffamily\small
\textbf{Trường hợp 1 (Vị trí 1 không phải <p>):} \par\vspace{1pt}
\hspace{0.5em}p:nth-child(1) $\rightarrow$ \textcolor{red}{Thất bại (Không đỏ)}. p:nth-of-type(1) $\rightarrow$ \textcolor{green!70!black}{\textbf{Thành công (Tô xanh)}} \par\vspace{5pt}

\textbf{Trường hợp 2 (Không đủ 3 phần tử):} \par\vspace{1pt}
\hspace{0.5em}li:nth-child(3n) $\rightarrow$ \textcolor{gray!60}{Không có tác dụng vì danh sách chỉ có 2 items} \par\vspace{5pt}

\textbf{Trường hợp 3 (Odd/Even trong DOM hỗn hợp):} \par\vspace{1pt}
\hspace{0.5em}p:nth-child(odd) $\rightarrow$ \textcolor{red}{Thất bại (Vì <p> nằm ở vị trí chẵn 2 \& 4)}. p:nth-of-type(odd) $\rightarrow$ \textcolor{green!70!black}{\textbf{Thành công (Tô xanh Paragraph 2)}}
\end{browserWindow}

\begin{prettyTableBox}[Bảng So Sánh Toàn Diện Cơ Chế Đếm: :nth-child(n) vs :nth-of-type(n)]
\rowcolors{2}{white}{primaryBlue!4}
\sffamily\small
\begin{tabularx}{\linewidth}{|>{\raggedright\arraybackslash\bfseries\color{primaryBlue}}m{3.2cm}|Z|Z|}
\hline
\rowcolor{primaryBlue} \thd{\sffamily Tiêu chí so sánh} & \thd{\sffamily :nth-child(an + b)} & \thd{\sffamily :nth-of-type(an + b)} \\
\hline
Cơ chế đếm & Đếm \textbf{TẤT CẢ} phần tử con (từ trên xuống) & Chỉ đếm phần tử con \textbf{CÙNG loại thẻ} (từ trên xuống) \\
\hline
Ảnh hưởng bởi thẻ khác loại? & \textcolor{red}{\textbf{Có}} (Thẻ khác loại chen vào sẽ làm tăng vị trí index) & \textcolor{green!70!black}{\textbf{Không}} (Thẻ khác loại chen vào bị bỏ qua hoàn toàn) \\
\hline
Tương đương logic với $n=1$ & \codeinline{:first-child} & \codeinline{:first-of-type} \\
\hline
Mức độ nghiêm ngặt & Nghiêm ngặt (Yêu cầu đúng vị trí tuyệt đối trong DOM) & Linh hoạt (Phù hợp với HTML lộn xộn nhiều loại thẻ) \\
\hline
Khi nào nên dùng? & Khi danh sách đồng nhất (ví dụ toàn \codeinline{<li>} hoặc \codeinline{<tr>}) & Khi container hỗn hợp nhiều loại thẻ (\codeinline{<h2>}, \codeinline{<p>}, \codeinline{<span>}) \\
\hline
\end{tabularx}
\end{prettyTableBox}

\begin{prettyTableBox}[Bảng Tóm Tắt Các Trường Hợp Sai Phổ Biến \& Cách Khắc Phục]
\rowcolors{2}{white}{primaryBlue!4}
\sffamily\small
\begin{tabularx}{\linewidth}{|>{\raggedright\arraybackslash\bfseries\color{primaryBlue}}m{3.8cm}|Z|Z|}
\hline
\rowcolor{primaryBlue} \thd{\sffamily Trường hợp sai} & \thd{\sffamily Nguyên nhân cốt lõi} & \thd{\sffamily Cách khắc phục chuẩn} \\
\hline
nth-child(n) không áp dụng đúng & Không tính đúng thứ tự toàn bộ phần tử con trong DOM & Dùng :nth-of-type(n) \\
\hline
3n, 4n+1... không chọn gì & Không có phần tử ở vị trí khớp với công thức số học & Kiểm tra lại tổng số phần tử con \\
\hline
odd, even không chọn như ý & Đếm tính theo tất cả phần tử con (không phân biệt loại thẻ) & Dùng :nth-of-type(odd/even) \\
\hline
Kết hợp :nth-child() với selector thẻ & Chọn sai vì thẻ cụ thể không đứng ở đúng vị trí index trong DOM & Xem lại thứ tự DOM hoặc dùng :nth-of-type() \\
\hline
\end{tabularx}
\end{prettyTableBox}

\begin{prettyTableBox}[Bảng Tổng Hợp Các Biểu Thức n Và Cách Hoạt Động Chi Tiết]
\rowcolors{2}{white}{primaryBlue!4}
\sffamily\small
\begin{tabularx}{\linewidth}{|>{\raggedright\arraybackslash \ttfamily\color{accentOrange}\bfseries}m{3.8cm}|Z|}
\hline
\rowcolor{primaryBlue} \thd{\sffamily Biểu thức n} & \thd{\sffamily Ý nghĩa \& Quy tắc chọn phần tử} \\
\hline
nth-child(3) & Chọn chính xác phần tử ở vị trí thứ 3. \\
\hline
nth-child(odd) & Chọn các phần tử có vị trí lẻ (1, 3, 5, 7, 9...). \\
\hline
nth-child(even) & Chọn các phần tử có vị trí chẵn (2, 4, 6, 8, 10...). \\
\hline
nth-child(3n) & Chọn các phần tử là bội số của 3 (vị trí 3, 6, 9, 12...). \\
\hline
nth-child(3n + 1) & Chọn các phần tử ở vị trí 1, 4, 7, 10, 13... (bắt đầu từ 1, nhảy 3 bước). \\
\hline
nth-child(-n + 3) & Chọn 3 phần tử đầu tiên (vị trí 1, 2, 3). \\
\hline
nth-child(n + 4) & Chọn tất cả phần tử từ vị trí thứ 4 trở đi (4, 5, 6, 7...). \\
\hline
\end{tabularx}
\end{prettyTableBox}

\begin{conceptBox}[Tóm Lại Về :nth-child(n) vs :nth-of-type(n)]
\begin{itemize}[leftmargin=1.5em, itemsep=2pt, parsep=0pt]
    \item \codeinline{:nth-child(n)} đếm tất cả phần tử con theo thứ tự tổng thể tuyệt đối.
    \item \codeinline{:nth-of-type(n)} đếm theo loại thẻ riêng biệt trong nhóm phần tử cùng tên thẻ.
    \item Có thể dùng số nguyên cụ thể, từ khóa (\codeinline{odd}, \codeinline{even}), hoặc biểu thức số học (\codeinline{3n + 1}).
    \item Rất linh hoạt và mạnh mẽ, cực kỳ hữu ích khi bạn muốn tạo kiểu theo quy luật mà không cần viết từng class riêng cho từng phần tử.
\end{itemize}
\end{conceptBox}

\paragraph{f. Chuyên Đề 4: Bộ Chọn Đếm Ngược Từ Đáy Lên (:nth-last-child() vs :nth-last-of-type())}

Trong thiết kế giao diện động, việc định kiểu cho các phần tử đứng ở cuối hoặc gần cuối danh sách (ví dụ: nút cuối cùng, bài viết kế cuối, dòng tổng cộng ở đáy bảng) là nhu cầu rất phổ biến. CSS cung cấp hai bộ chọn đếm ngược từ đáy lên: \codeinline{:nth-last-child()} và \codeinline{:nth-last-of-type()}. Phần này sẽ phân tích chi tiết và tách biệt hai bộ chọn này.

\begin{conceptBox}[Tổng Quan Nhanh: 2 Bộ Chọn Đếm Ngược Trong CSS]
\begin{itemize}[leftmargin=1.5em, itemsep=3pt, parsep=0pt]
    \item \codeinline{:nth-last-child(an + b)}: Đếm thứ tự \textbf{từ dưới lên trên} (bắt đầu từ phần tử con cuối cùng), đếm \textbf{tất cả các phần tử con} không phân biệt loại thẻ.
    \item \codeinline{:nth-last-of-type(an + b)}: Đếm thứ tự \textbf{từ dưới lên trên}, nhưng \textbf{chỉ đếm các phần tử cùng loại thẻ} (cùng tag name).
\end{itemize}
\end{conceptBox}

%% ===================================================================
%% PHẦN F1: :nth-last-child() -- PHÂN TÍCH CHUYÊN SÂU
%% ===================================================================

\subparagraph{f1. Phân Tích Chuyên Sâu :nth-last-child() (Đếm Ngược Toàn Bộ Phần Tử Con)}

\codeinline{:nth-last-child(an + b)} là pseudo-class selector chọn phần tử dựa trên vị trí thứ tự của nó khi \textbf{đếm từ cuối danh sách đếm ngược lên đầu} (từ phần tử con cuối cùng = 1).

\begin{conceptBox}[Định Nghĩa Chi Tiết :nth-last-child()]
\codeinline{:nth-last-child()} hoạt động hoàn toàn giống với \codeinline{:nth-child()}, ngoại trừ chiều đếm bị đảo ngược:
\begin{itemize}[leftmargin=1.5em, itemsep=2pt, parsep=0pt]
    \item Phần tử con cuối cùng trong cha có chỉ số $index = 1$.
    \item Phần tử con kế cuối có chỉ số $index = 2$.
    \item Trình duyệt đếm \textbf{tất cả phần tử con} (element nodes) mà không phân biệt loại thẻ (\codeinline{<div>}, \codeinline{<p>}, \codeinline{<span>}, \codeinline{<li>}...).
    \item Biểu thức công thức $an + b$, từ khóa \codeinline{odd}, \codeinline{even}, hoặc số nguyên $n$ đều áp dụng chính xác theo chiều ngược từ dưới lên.
\end{itemize}
\end{conceptBox}

\begin{codeWindow}[Cú Pháp Cơ Bản Của :nth-last-child()]
\begin{lstlisting}[language=CSS3]
/* Select element at specific position from bottom */
selector:nth-last-child(n) {
    /* CSS properties apply here */
}

/* Select using formula from bottom */
selector:nth-last-child(an + b) {
    /* CSS properties apply here */
}
\end{lstlisting}
\end{codeWindow}

\begin{codeWindow}[Ví Dụ 1: Chọn Phần Tử Kế Cuối Bằng li:nth-last-child(2)]
\begin{lstlisting}[language=HTML5]
<ul>
  <li>Item 1</li>
  <li>Item 2</li>
  <li>Item 3</li>
  <li>Item 4</li> <!-- Selected: 2nd from bottom -->
  <li>Item 5</li> <!-- Position 1 from bottom -->
</ul>

<style>
  /* Select 2nd element from bottom */
  li:nth-last-child(2) {
    color: red;
    font-weight: bold;
  }
</style>
\end{lstlisting}
\end{codeWindow}

\begin{browserWindow}[Kết Quả Ví Dụ 1: li:nth-last-child(2)]
\sffamily\small
\begin{itemize}[leftmargin=1.5em, itemsep=2pt, parsep=0pt]
    \item Item 1 {\scriptsize\textcolor{gray!60}{(Vị trí 5 từ dưới)}}
    \item Item 2 {\scriptsize\textcolor{gray!60}{(Vị trí 4 từ dưới)}}
    \item Item 3 {\scriptsize\textcolor{gray!60}{(Vị trí 3 từ dưới)}}
    \item \textcolor{red}{\textbf{Item 4 (Chữ đỏ, in đậm)}} {\scriptsize\textcolor{gray!60}{($\leftarrow$ li:nth-last-child(2) khớp -- vị trí thứ 2 từ dưới lên)}}
    \item Item 5 {\scriptsize\textcolor{gray!60}{(Vị trí 1 từ dưới -- phần tử cuối cùng)}}
\end{itemize}
\end{browserWindow}

\begin{codeWindow}[Ví Dụ 2: Dùng Công Thức an + b Đếm Ngược Từ Đáy]
\begin{lstlisting}[language=HTML5]
<ul>
  <li>Item 1</li> <!-- Match n=2 (pos 5) -->
  <li>Item 2</li>
  <li>Item 3</li> <!-- Match n=1 (pos 3) -->
  <li>Item 4</li>
  <li>Item 5</li> <!-- Match n=0 (pos 1) -->
</ul>

<style>
  /* Select odd positions from bottom (1, 3, 5...) */
  li:nth-last-child(2n + 1) {
    background-color: #f0f0f0;
  }
</style>
\end{lstlisting}
\end{codeWindow}

\begin{browserWindow}[Kết Quả Ví Dụ 2: li:nth-last-child(2n + 1)]
\sffamily\small
\begin{itemize}[leftmargin=1.5em, itemsep=3pt, parsep=0pt]
    \item \colorbox{gray!20}{\textbf{Item 1}} {\scriptsize\textcolor{gray!60}{($\leftarrow$ Nền xám \#f0f0f0 -- Vị trí 5 từ dưới lên)}}
    \item Item 2 {\scriptsize\textcolor{gray!60}{(Vị trí 4 từ dưới -- không đổi)}}
    \item \colorbox{gray!20}{\textbf{Item 3}} {\scriptsize\textcolor{gray!60}{($\leftarrow$ Nền xám \#f0f0f0 -- Vị trí 3 từ dưới lên)}}
    \item Item 4 {\scriptsize\textcolor{gray!60}{(Vị trí 2 từ dưới -- không đổi)}}
    \item \colorbox{gray!20}{\textbf{Item 5}} {\scriptsize\textcolor{gray!60}{($\leftarrow$ Nền xám \#f0f0f0 -- Vị trí 1 từ dưới lên)}}
\end{itemize}
\end{browserWindow}

\begin{codeWindow}[Ví Dụ 3: DOM Hỗn Hợp -- Lưu Ý Về Vị Trí Tuyệt Đối Đếm Ngược]
\begin{lstlisting}[language=HTML5]
<div class="container">
  <h2>Title</h2>      <!-- Pos 4 from bottom -->
  <p>Paragraph 1</p>  <!-- Pos 3 from bottom -->
  <p>Paragraph 2</p>  <!-- Pos 2 from bottom -->
  <span>Footer</span> <!-- Pos 1 from bottom -->
</div>

<style>
  /* FAILS to select Paragraph 2:
     Because at pos 1 from bottom is <span>, pos 2 is <p>Paragraph 2</p>,
     so p:nth-last-child(1) looks for <p> at position 1 from bottom (which is <span>) */
  p:nth-last-child(1) { color: red; }
</style>
\end{lstlisting}
\end{codeWindow}

\begin{browserWindow}[Kết Quả Ví Dụ 3: DOM Hỗn Hợp Đếm Ngược]
\sffamily\small
Title {\scriptsize\textcolor{gray!60}{(Thẻ <h2> -- Vị trí 4 từ dưới)}} \par\vspace{2pt}
Paragraph 1 {\scriptsize\textcolor{gray!60}{(Thẻ <p> -- Vị trí 3 từ dưới)}} \par\vspace{2pt}
Paragraph 2 {\scriptsize\textcolor{gray!60}{(Thẻ <p> -- Vị trí 2 từ dưới, KHÔNG đỏ)}} \par\vspace{2pt}
Footer {\scriptsize\textcolor{gray!60}{(Thẻ <span> -- Vị trí 1 từ dưới)}} \par\vspace{4pt}
{\scriptsize\textcolor{gray!60}{\textit{Giải thích: p:nth-last-child(1) thất bại vì phần tử vị trí 1 từ dưới lên là <span>, không phải <p>.}}}
\end{browserWindow}

\begin{warnBox}[Các Lưu Ý Quan Trọng Khi Sử Dụng :nth-last-child()]
\textbf{1. Đếm tất cả loại phần tử con:} \codeinline{:nth-last-child()} đếm mọi thẻ con trong container khi đếm ngược từ dưới lên.
\par\vspace{4pt}
\textbf{2. Dễ nhầm lẫn vị trí trong DOM hỗn hợp:} Nếu thẻ con cuối cùng là \codeinline{<span>}, thì \codeinline{p:nth-last-child(1)} sẽ không có tác dụng.
\par\vspace{4pt}
\textbf{3. Giải pháp khi có nhiều thẻ khác loại:} Dùng \codeinline{:nth-last-of-type()} nếu chỉ muốn đếm ngược trên một loại thẻ cụ thể.
\end{warnBox}

\begin{tipBox}[Mẹo Ghi Nhớ :nth-last-child()]
\textbf{Công thức tương đương:} \codeinline{:nth-last-child(1)} $\equiv$ \codeinline{:last-child} (phần tử con cuối cùng tuyệt đối).
\end{tipBox}

\begin{conceptBox}[Ứng Dụng Thực Tế Của :nth-last-child()]
\begin{itemize}[leftmargin=1.5em, itemsep=3pt, parsep=0pt]
    \item \textbf{Ẩn đường phân cách (border) của N phần tử cuối}: Ví dụ xóa viền dưới của 2 item cuối trong danh sách.
    \item \textbf{Định dạng bảng dữ liệu}: Làm nổi bật dòng tổng cộng hoặc 2 dòng cuối bảng.
    \item \textbf{Timeline \& Nhật ký}: Style khác biệt cho các bài viết cũ nhất (ở đáy timeline).
\end{itemize}
\end{conceptBox}

%% ===================================================================
%% PHẦN F2: :nth-last-of-type() -- PHÂN TÍCH CHUYÊN SÂU
%% ===================================================================

\subparagraph{f2. Phân Tích Chuyên Sâu :nth-last-of-type() (Đếm Ngược Theo Loại Thẻ)}

\codeinline{:nth-last-of-type(an + b)} là pseudo-class selector chọn phần tử dựa trên vị trí của nó khi \textbf{chỉ lọc đếm các phần tử cùng loại thẻ (tag name) từ dưới lên trên}.

\begin{conceptBox}[Định Nghĩa Chi Tiết :nth-last-of-type()]
Khác với \codeinline{:nth-last-child()}, \codeinline{:nth-last-of-type()} \textbf{bỏ qua các thẻ khác loại} và chỉ đánh số thứ tự đếm ngược cho nhóm thẻ cùng loại:
\begin{itemize}[leftmargin=1.5em, itemsep=2pt, parsep=0pt]
    \item Trình duyệt lọc tất cả các thẻ cùng tag name (ví dụ tất cả các thẻ \codeinline{<p>}).
    \item Thẻ \codeinline{<p>} xuất hiện sau cùng trong DOM có chỉ số $index = 1$ (theo loại thẻ \codeinline{<p>}).
    \item Thẻ \codeinline{<p>} xuất hiện trước đó có chỉ số $index = 2$ (theo loại thẻ \codeinline{<p>}).
    \item Thẻ \codeinline{<h2>}, \codeinline{<span>}, \codeinline{<div>}... chen giữa hoàn toàn \textbf{không làm gián đoạn} danh sách đếm.
\end{itemize}
\end{conceptBox}

\begin{codeWindow}[Cú Pháp Cơ Bản Của :nth-last-of-type()]
\begin{lstlisting}[language=CSS3]
/* Select specific tag type counting from bottom */
selector:nth-last-of-type(n) {
    /* CSS properties apply here */
}

/* Formula for specific tag type from bottom */
selector:nth-last-of-type(an + b) {
    /* CSS properties apply here */
}
\end{lstlisting}
\end{codeWindow}

\begin{codeWindow}[Ví Dụ 1: Chọn Thẻ <p> Cuối Cùng Trong DOM Hỗn Hợp]
\begin{lstlisting}[language=HTML5]
<div class="container">
  <h2>Header</h2>
  <p>Paragraph 1</p> <!-- Pos 2 of <p> from bottom -->
  <p>Paragraph 2</p> <!-- Pos 1 of <p> from bottom -->
  <span>Footer text</span>
</div>

<style>
  /* Select last <p> tag from bottom, regardless of <span> below it */
  p:nth-last-of-type(1) {
    color: green;
    font-weight: bold;
  }
</style>
\end{lstlisting}
\end{codeWindow}

\begin{browserWindow}[Kết Quả Ví Dụ 1: p:nth-last-of-type(1)]
\sffamily\small
Header {\scriptsize\textcolor{gray!60}{(Thẻ <h2>)}} \par\vspace{2pt}
Paragraph 1 {\scriptsize\textcolor{gray!60}{(Thẻ <p> thứ 2 đếm từ dưới lên trong các thẻ <p>)}} \par\vspace{2pt}
\textcolor{green!70!black}{\textbf{Paragraph 2}} {\scriptsize\textcolor{gray!60}{($\leftarrow$ p:nth-last-of-type(1) khớp -- Thẻ <p> cuối cùng, dù có <span> đứng dưới)}} \par\vspace{2pt}
Footer text {\scriptsize\textcolor{gray!60}{(Thẻ <span> -- không bị đếm vào nhóm <p>)}}
\end{browserWindow}

\begin{codeWindow}[Ví Dụ 2: Đếm Ngược Xen Kẽ Cho Nhóm Thẻ <p>]
\begin{lstlisting}[language=HTML5]
<div>
  <h2>Title</h2>
  <p>Paragraph A</p> <!-- Pos 3 of <p> -> Matches odd (1st from top of <p>) -->
  <span>Divider</span>
  <p>Paragraph B</p> <!-- Pos 2 of <p> -> Even -->
  <p>Paragraph C</p> <!-- Pos 1 of <p> -> Matches odd -->
</div>

<style>
  /* Select odd <p> tags counting from bottom */
  p:nth-last-of-type(odd) {
    background-color: lightyellow;
  }
</style>
\end{lstlisting}
\end{codeWindow}

\begin{browserWindow}[Kết Quả Ví Dụ 2: p:nth-last-of-type(odd)]
\sffamily\small
Title {\scriptsize\textcolor{gray!60}{(Thẻ <h2>)}} \par\vspace{2pt}
\colorbox{yellow!20}{\textbf{Paragraph A}} {\scriptsize\textcolor{gray!60}{($\leftarrow$ Vị trí 3 đếm ngược trong các thẻ <p> -- Khớp odd)}} \par\vspace{2pt}
Divider {\scriptsize\textcolor{gray!60}{(Thẻ <span> -- Không bị đếm)}} \par\vspace{2pt}
Paragraph B {\scriptsize\textcolor{gray!60}{(Vị trí 2 đếm ngược trong các thẻ <p> -- Even)}} \par\vspace{2pt}
\colorbox{yellow!20}{\textbf{Paragraph C}} {\scriptsize\textcolor{gray!60}{($\leftarrow$ Vị trí 1 đếm ngược trong các thẻ <p> -- Khớp odd)}}
\end{browserWindow}

\begin{warnBox}[Các Lưu Ý Quan Trọng Khi Sử Dụng :nth-last-of-type()]
\textbf{1. Đếm độc lập cho từng loại thẻ:} Nếu trong cùng container có thẻ \codeinline{<h2>}, \codeinline{<p>}, \codeinline{<img>}, trình duyệt sẽ tạo 3 danh sách đếm ngược độc lập cho từng loại thẻ.
\par\vspace{4pt}
\textbf{2. Linh hoạt hơn :nth-last-child():} Rất thích hợp cho mã HTML thực tế nơi các thẻ tiêu đề, đoạn văn, hình ảnh thường nằm xen kẽ nhau.
\end{warnBox}

\begin{tipBox}[Mẹo Ghi Nhớ :nth-last-of-type()]
\textbf{Công thức tương đương:} \codeinline{:nth-last-of-type(1)} $\equiv$ \codeinline{:last-of-type} (phần tử cuối cùng thuộc loại thẻ đó).
\end{tipBox}

\begin{conceptBox}[Ứng Dụng Thực Tế Của :nth-last-of-type()]
\begin{itemize}[leftmargin=1.5em, itemsep=3pt, parsep=0pt]
    \item \textbf{Đoạn văn cuối cùng trong bài viết}: Thêm chữ ký, nguồn tác giả, hoặc khoảng cách lề dưới cho thẻ \codeinline{<p>} cuối cùng.
    \item \textbf{Hình ảnh cuối bài}: Tự động thêm caption hoặc bo góc cho ảnh \codeinline{<img>} cuối cùng.
    \item \textbf{Form ô nhập liệu cuối}: Đổi style cho nút bấm hoặc input cuối cùng trong Form.
\end{itemize}
\end{conceptBox}

%% ===================================================================
%% PHẦN F3: BẢNG SO SÁNH TỔNG HỢP :nth-last-child() vs :nth-last-of-type()
%% ===================================================================

\subparagraph{f3. Bảng So Sánh Chi Tiết :nth-last-child() vs :nth-last-of-type()}
\mbox{}\par\vspace{2pt}

\begin{codeWindow}[Ví Dụ Trực Quan So Sánh :nth-last-child(1) vs :nth-last-of-type(1)]
\begin{lstlisting}[language=HTML5]
<div class="demo">
  <h2>Article Header</h2>
  <p>First paragraph</p>
  <p>Second paragraph</p>
  <span>Footnote reference</span>
</div>

<style>
  /* FAILS: Last child is <span>, not <p> */
  .demo p:nth-last-child(1) { color: red; }

  /* WORKS: Selects last <p> (Second paragraph), ignoring <span> */
  .demo p:nth-last-of-type(1) { color: green; font-weight: bold; }
</style>
\end{lstlisting}
\end{codeWindow}

\begin{browserWindow}[Kết Quả So Sánh Trực Quan trong Trình Duyệt]
\sffamily\small
Article Header {\scriptsize\textcolor{gray!60}{(Thẻ <h2> -- Vị trí 4 từ dưới)}} \par\vspace{2pt}
First paragraph {\scriptsize\textcolor{gray!60}{(Thẻ <p> -- Vị trí 3 từ dưới)}} \par\vspace{2pt}
\textcolor{green!70!black}{\textbf{Second paragraph}} {\scriptsize\textcolor{gray!60}{($\leftarrow$ p:nth-last-of-type(1) KHỚP -- Thẻ <p> cuối cùng)}} \par\vspace{2pt}
Footnote reference {\scriptsize\textcolor{gray!60}{(Thẻ <span> -- Vị trí 1 từ dưới tuyệt đối)}} \par\vspace{4pt}
{\scriptsize\textcolor{gray!60}{\textit{Lưu ý: p:nth-last-child(1) KHÔNG khớp vì vị trí 1 từ dưới lên là <span>, không phải <p>.}}}
\end{browserWindow}

\begin{prettyTableBox}[Bảng So Sánh Toàn Diện :nth-last-child() vs :nth-last-of-type()]
\rowcolors{2}{white}{primaryBlue!4}
\sffamily\small
\begin{tabularx}{\linewidth}{|>{\raggedright\arraybackslash\bfseries\color{primaryBlue}}m{3.2cm}|Z|Z|}
\hline
\rowcolor{primaryBlue} \thd{\sffamily Tiêu chí so sánh} & \thd{\sffamily :nth-last-child(an + b)} & \thd{\sffamily :nth-last-of-type(an + b)} \\
\hline
Hướng đếm & Từ dưới lên trên (từ đáy lên) & Từ dưới lên trên (từ đáy lên) \\
\hline
Cơ chế đếm & Đếm \textbf{TẤT CẢ} phần tử con (không phân biệt loại thẻ) & Chỉ đếm phần tử con \textbf{CÙNG loại thẻ} (cùng tag name) \\
\hline
Ảnh hưởng bởi thẻ khác loại? & \textcolor{red}{\textbf{Có}} -- Thẻ khác loại ở đáy sẽ làm thay đổi chỉ số vị trí & \textcolor{green!70!black}{\textbf{Không}} -- Thẻ khác loại bị bỏ qua hoàn toàn \\
\hline
Tương đương logic với $n=1$ & \codeinline{:last-child} & \codeinline{:last-of-type} \\
\hline
Đối ứng với đếm từ trên xuống & \codeinline{:nth-child(an + b)} & \codeinline{:nth-of-type(an + b)} \\
\hline
Mức độ linh hoạt & Nghiêm ngặt (yêu cầu đúng vị trí tuyệt đối từ đáy) & Linh hoạt (phù hợp với HTML lộn xộn nhiều loại thẻ) \\
\hline
Khi nào dùng? & Khi cấu trúc danh sách đồng nhất (ví dụ toàn \codeinline{<li>}) & Khi cấu trúc danh sách hỗn hợp nhiều loại thẻ (\codeinline{<h2>}, \codeinline{<p>}, \codeinline{<span>}) \\
\hline
\end{tabularx}
\end{prettyTableBox}

\begin{prettyTableBox}[Bảng Tóm Tắt Quy Tắc Sử Dụng :nth-last-child() \& :nth-last-of-type()]
\rowcolors{2}{white}{primaryBlue!4}
\sffamily\small
\begin{tabularx}{\linewidth}{|Z|>{\centering\arraybackslash\bfseries}m{2.0cm}|>{\raggedright\arraybackslash}m{4.5cm}|}
\hline
\rowcolor{primaryBlue} \thd{\sffamily Quy tắc} & \thd{\sffamily Đúng / Sai} & \thd{\sffamily Ghi chú} \\
\hline
:nth-last-child(1) luôn chọn phần tử cuối cùng tuyệt đối & \textcolor{green!70!black}{Đúng} & Tương đương :last-child \\
\hline
:nth-last-child() không đếm text node hay khoảng trắng & \textcolor{green!70!black}{Đúng} & Chỉ đếm element nodes HTML \\
\hline
Có thẻ khác loại ở đáy làm lệch đếm :nth-last-child() & \textcolor{green!70!black}{Đúng} & Vì đếm tất cả phần tử con \\
\hline
:nth-last-of-type() chỉ đếm thẻ cùng loại từ dưới lên & \textcolor{green!70!black}{Đúng} & Bỏ qua thẻ khác loại \\
\hline
Muốn chọn thẻ p cuối cùng trong container hỗn hợp dùng :nth-last-of-type(1) & \textcolor{green!70!black}{Đúng} & Chuẩn xác và an toàn nhất \\
\hline
\end{tabularx}
\end{prettyTableBox}

\begin{tipBox}[Tổng Kết Nhanh Mẹo Phân Biệt 4 Bộ Chọn Vị Trí]
\begin{itemize}[leftmargin=1.5em, itemsep=3pt, parsep=0pt]
    \item \codeinline{:nth-child(n)}: Đếm từ trên xuống $\rightarrow$ Tất cả loại thẻ.
    \item \codeinline{:nth-of-type(n)}: Đếm từ trên xuống $\rightarrow$ Chỉ thẻ cùng loại.
    \item \codeinline{:nth-last-child(n)}: Đếm từ dưới lên $\rightarrow$ Tất cả loại thẻ.
    \item \codeinline{:nth-last-of-type(n)}: Đếm từ dưới lên $\rightarrow$ Chỉ thẻ cùng loại.
\end{itemize}
\end{tipBox}

\paragraph{g. Chuyên Đề 5: Bộ Chọn Con Duy Nhất (:only-child vs :only-of-type)}

Trong thiết kế giao diện Web động, có những trường hợp phần tử cha chỉ chứa duy nhất một phần tử con (ví dụ: một thông báo lỗi duy nhất trong hộp thoại, một nút bấm đơn lẻ trong card, hoặc bài viết chỉ chứa đúng một đoạn văn hay một hình ảnh). CSS3 cung cấp hai bộ chọn giả đặc biệt cho các tình huống này: \codeinline{:only-child} và \codeinline{:only-of-type}. Phần này sẽ phân tích chuyên sâu từng bộ chọn và bảng so sánh chi tiết.

\begin{conceptBox}[Tổng Quan Nhanh: 2 Cơ Chế Kiểm Tra Tính Duy Nhất]
\begin{itemize}[leftmargin=1.5em, itemsep=3pt, parsep=0pt]
    \item \codeinline{:only-child}: Phần tử là \textbf{phần tử con duy nhất tuyệt đối} của phần tử cha (không có bất kỳ anh chị em nào khác).
    \item \codeinline{:only-of-type}: Phần tử là \textbf{duy nhất trong nhóm cùng loại thẻ} (cùng tag name) trong phần tử cha (có thể có các anh chị em khác loại thẻ).
\end{itemize}
\end{conceptBox}

%% ===================================================================
%% PHẦN G1: :only-child -- PHÂN TÍCH CHUYÊN SÂU
%% ===================================================================

\subparagraph{g1. Phân Tích Chuyên Sâu :only-child (Con Duy Nhất Tuyệt Đối)}

\codeinline{:only-child} là một pseudo-class selector dùng để chọn một phần tử khi nó là \textbf{phần tử con duy nhất tuyệt đối} của phần tử cha.

\begin{conceptBox}[Định Nghĩa Chi Tiết \& Cơ Chế Hoạt Động Của :only-child]
\codeinline{:only-child} khớp với một phần tử khi và chỉ khi hai điều kiện sau thỏa mãn đồng thời:
\begin{itemize}[leftmargin=1.5em, itemsep=2pt, parsep=0pt]
    \item Phần tử là con duy nhất của phần tử cha (tức là cha chỉ có đúng 1 phần tử con HTML, không có anh chị em nào khác).
    \item Phần tử đó khớp với bộ chọn \codeinline{selector} phía trước.
    \item \textbf{Lưu ý quan trọng:} Nếu trong phần tử cha xuất hiện thêm bất kỳ phần tử con nào khác (dù khác loại thẻ như \codeinline{<span>}, \codeinline{<div>}, \codeinline{<i>}...), điều kiện \codeinline{:only-child} sẽ ngay lập tức bị phá vỡ $\rightarrow$ \textcolor{red}{\textbf{Không được chọn!}}
\end{itemize}
\par\vspace{4pt}
\textbf{Tương đương logic:} \codeinline{selector:only-child} $\equiv$ \codeinline{selector:first-child:last-child} (vừa là con đầu tiên vừa là con cuối cùng $\rightarrow$ phải là con duy nhất).
\end{conceptBox}

\begin{codeWindow}[Cú Pháp Cơ Bản Của :only-child]
\begin{lstlisting}[language=CSS3]
selector:only-child {
    /* CSS properties apply here */
}

/* Equivalent to: */
selector:first-child:last-child {
    /* Same result */
}
\end{lstlisting}
\end{codeWindow}

\begin{codeWindow}[Ví Dụ 1: Đoạn Văn Là Con Duy Nhất Trong <div>]
\begin{lstlisting}[language=HTML5]
<div>
  <p>Day la doan duy nhat</p> <!-- Single child inside <div> -->
</div>

<style>
  p:only-child {
    color: red;
    font-weight: bold;
  }
</style>
\end{lstlisting}
\end{codeWindow}

\begin{browserWindow}[Kết Quả Hiển Thị Ví Dụ 1: p:only-child Thành Công]
\sffamily\small
\textcolor{red}{\textbf{Day la doan duy nhat}} {\scriptsize\textcolor{gray!60}{($\leftarrow$ p:only-child KHỚP -- Được tô đỏ, in đậm vì là phần tử con duy nhất trong <div>)}}
\end{browserWindow}

\begin{codeWindow}[Ví Dụ 2: Trường Hợp THẤT BẠI Của :only-child Khi Có 2 Phần Tử]
\begin{lstlisting}[language=HTML5]
<div>
  <p>Doan 1</p> <!-- Siblings present -->
  <p>Doan 2</p>
</div>

<style>
  /* FAILS: <div> has 2 children, not 1 */
  p:only-child {
    color: red;
  }
</style>
\end{lstlisting}
\end{codeWindow}

\begin{browserWindow}[Kết Quả Hiển Thị Ví Dụ 2: :only-child Thất Bại]
\sffamily\small
Doan 1 {\scriptsize\textcolor{gray!60}{(Không tô đỏ -- Có anh em cùng cấp)}} \par\vspace{2pt}
Doan 2 {\scriptsize\textcolor{gray!60}{(Không tô đỏ -- Có anh em cùng cấp)}} \par\vspace{4pt}
{\scriptsize\textcolor{gray!60}{\textit{Giải thích: Dù viết p:only-child, cả hai <p> đều không bị áp dụng CSS vì chúng không phải là con duy nhất (có 2 phần tử con trong <div>).}}}
\end{browserWindow}

\begin{warnBox}[Các Lưu Ý Quan Trọng Khi Sử Dụng :only-child]
\textbf{1. Đếm tất cả các phần tử con HTML:} \codeinline{:only-child} đếm toàn bộ phần tử con element nodes (không phân biệt tag name).
\par\vspace{4pt}
\textbf{2. Khoảng trắng / Text node không bị đếm:} Các node văn bản thuần hoặc xuống dòng trong HTML không tính là phần tử con HTML.
\par\vspace{4pt}
\textbf{3. Bị phá vỡ bởi thẻ khác loại:} Nếu container chứa \codeinline{<p>Paragraph</p>} và \codeinline{<span>Note</span>}, \codeinline{p:only-child} sẽ thất bại.
\end{warnBox}

\begin{tipBox}[Mẹo Ghi Nhớ :only-child]
\textbf{Hình dung:} \codeinline{:only-child} giống như "con một" trong gia đình. Gia đình chỉ được phép có đúng 1 đứa con duy nhất.
\end{tipBox}

\begin{conceptBox}[Ứng Dụng Thực Tế Của :only-child]
\begin{itemize}[leftmargin=1.5em, itemsep=3pt, parsep=0pt]
    \item \textbf{Căn giữa nội dung đơn lẻ:} Khi một modal hoặc card chỉ có đúng 1 nút bấm hay 1 dòng thông báo, tự động căn giữa và tăng font-size.
    \item \textbf{Ẩn icon/separator:} Xóa đường gạch ngang hoặc icon phân cách khi danh sách chỉ có 1 item duy nhất.
    \item \textbf{Tự động căn chỉnh lề (Margin Reset):} Xóa khoảng cách lề thừa khi container chỉ có 1 phần tử.
\end{itemize}
\end{conceptBox}

%% ===================================================================
%% PHẦN G2: :only-of-type -- PHÂN TÍCH CHUYÊN SÂU
%% ===================================================================

\paragraph{g2. Phân Tích Chuyên Sâu :only-of-type (Duy Nhất Theo Loại Thẻ)}
\mbox{}\par\vspace{2pt}

\codeinline{:only-of-type} là pseudo-class selector dùng để chọn phần tử con duy nhất thuộc một loại (tag name) trong phần tử cha. Nếu trong phần tử cha chỉ có đúng 1 phần tử loại đó, thì nó sẽ được chọn.

\begin{conceptBox}[Định Nghĩa Chi Tiết \& Cơ Chế Hoạt Động Của :only-of-type]
\codeinline{:only-of-type} lọc và hoạt động theo các bước:
\begin{itemize}[leftmargin=1.5em, itemsep=2pt, parsep=0pt]
    \item Trình duyệt lọc tất cả các phần tử con cùng tag name (ví dụ tất cả các thẻ \codeinline{<p>}).
    \item Nếu nhóm phần tử cùng loại thẻ đó có \textbf{đúng 1 phần tử duy nhất}, phần tử đó sẽ được chọn.
    \item Các phần tử con khác loại thẻ (\codeinline{<h2>}, \codeinline{<span>}, \codeinline{<div>}...) nằm trong container hoàn toàn \textbf{bị bỏ qua} và không ảnh hưởng đến điều kiện duy nhất.
\end{itemize}
\par\vspace{4pt}
\textbf{Tương đương logic:} \codeinline{selector:only-of-type} $\equiv$ \codeinline{selector:first-of-type:last-of-type} (vừa là thẻ đầu vừa là thẻ cuối của loại đó $\rightarrow$ phải là thẻ duy nhất của loại đó).
\end{conceptBox}

\begin{codeWindow}[Cú Pháp Cơ Bản Của :only-of-type]
\begin{lstlisting}[language=CSS3]
selector:only-of-type {
  /* style apply */
}
\end{lstlisting}
\end{codeWindow}

\begin{codeWindow}[Ví Dụ 1: Thẻ <p> Duy Nhất Trong Container Có Thẻ <h2> Chen Giữa]
\begin{lstlisting}[language=HTML5]
<div class="container">
  <h2>Tieu de bai viet</h2>
  <p>Day la doan van duy nhat</p> <!-- Only <p> tag in .container -->
</div>

<style>
  p:only-of-type {
    color: red;
    font-weight: bold;
  }
</style>
\end{lstlisting}
\end{codeWindow}

\begin{browserWindow}[Kết Quả Hiển Thị Ví Dụ 1: p:only-of-type Thành Công]
\sffamily\small
Tieu de bai viet {\scriptsize\textcolor{gray!60}{(Thẻ <h2> -- Bị bỏ qua)}} \par\vspace{2pt}
\textcolor{red}{\textbf{Day la doan van duy nhat}} {\scriptsize\textcolor{gray!60}{($\leftarrow$ p:only-of-type KHỚP -- Trong container chỉ có 1 phần tử <p>)}}
\end{browserWindow}

\begin{codeWindow}[Ví Dụ 2: Trường Hợp THẤT BẠI Của :only-of-type Khi Có 2 Thẻ <p>]
\begin{lstlisting}[language=HTML5]
<div class="container">
  <h2>Tieu de bai viet</h2>
  <p>Doan 1</p> <!-- 2 <p> tags present -->
  <p>Doan 2</p>
</div>

<style>
  /* FAILS: container has 2 <p> tags */
  p:only-of-type {
    color: red;
  }
</style>
\end{lstlisting}
\end{codeWindow}

\begin{browserWindow}[Kết Quả Hiển Thị Ví Dụ 2: :only-of-type Thất Bại]
\sffamily\small
Tieu de bai viet {\scriptsize\textcolor{gray!60}{(Thẻ <h2>)}} \par\vspace{2pt}
Doan 1 {\scriptsize\textcolor{gray!60}{(Không chọn -- Có 2 thẻ <p>)}} \par\vspace{2pt}
Doan 2 {\scriptsize\textcolor{gray!60}{(Không chọn -- Có 2 thẻ <p>)}}
\end{browserWindow}

\begin{warnBox}[Các Lưu Ý Quan Trọng Khi Sử Dụng :only-of-type]
\textbf{1. Chỉ đếm các phần tử CÙNG loại thẻ:} Trình duyệt lọc ra các phần tử cùng tag name, các thẻ khác loại hoàn toàn vô hình trong phép đếm này.
\par\vspace{4pt}
\textbf{2. Linh hoạt hơn :only-child trong thực tế:} Một container thường chứa nhiều loại thẻ khác nhau (\codeinline{<h2>}, \codeinline{<p>}, \codeinline{<span>}...), nên \codeinline{:only-of-type} dễ khớp và an toàn hơn.
\end{warnBox}

\begin{prettyTableBox}[Bảng So Sánh Chi Tiết: :only-child vs :only-of-type]
\rowcolors{2}{white}{primaryBlue!4}
\sffamily\small
\begin{tabularx}{\linewidth}{|>{\raggedright\arraybackslash\bfseries\color{primaryBlue}}m{3.5cm}|Z|}
\hline
\rowcolor{primaryBlue} \thd{\sffamily Selector} & \thd{\sffamily Ý nghĩa cốt lõi} \\
\hline
:only-child & Phần tử là phần tử con duy nhất tuyệt đối của cha (đếm mọi loại thẻ HTML). \\
\hline
:only-of-type & Phần tử là duy nhất trong nhóm cùng loại (cùng tag name) trong cha (bỏ qua thẻ khác). \\
\hline
\end{tabularx}
\end{prettyTableBox}

%% ===================================================================
%% PHẦN H: :empty -- PHÂN TÍCH CHUYÊN SÂU BỘ CHỌN PHẦN TỬ RỖNG
%% ===================================================================

\paragraph{h. Phân Tích Chuyên Sâu Bộ Chọn :empty (Chọn Phần Tử Rỗng 100\%)}
\mbox{}\par\vspace{2pt}

Bộ chọn lớp giả \codeinline{:empty} trong CSS được dùng để chọn các phần tử \textbf{không có bất kỳ nội dung con nào cả}, kể cả văn bản, thẻ HTML, khoảng trắng (whitespace) hay ghi chú (comment HTML).

\begin{conceptBox}[Định Nghĩa \& Cú Pháp Chuẩn Của :empty]
\textbf{Cú pháp khai báo:}
\begin{codeWindow}[Cú Pháp Định Kiểu :empty]
\begin{lstlisting}[language=CSS3]
selector:empty {
  /* Thuoc tinh CSS ap dung cho the rong */
}
\end{lstlisting}
\end{codeWindow}
\par\vspace{4pt}
\textbf{Ý nghĩa cốt lõi:} \codeinline{:empty} chỉ chọn các phần tử không chứa bất kỳ thành phần nào bên trong --- \textbf{không có text, không có thẻ con, không có khoảng trắng, không có ghi chú comment}.
\end{conceptBox}

\begin{codeWindow}[Ví Dụ Thử Nghiệm 4 Phần Tử Box Với :empty]
\begin{lstlisting}[language=HTML5]
<div class="box box1"></div>
<div class="box box2">   </div>
<div class="box box3"><span></span></div>
<div class="box box4">Hello</div>

<style>
  /* Chi selector:empty moi co vien do va nen hong */
  div.box:empty {
    border: 2px solid red;
    width: 30px;
    height: 30px;
    background-color: #FFE6E6;
  }
</style>
\end{lstlisting}
\end{codeWindow}

\begin{browserWindow}[Kết Quả Trực Quan Hiển Thị Trên Trình Duyệt Cho 4 Thẻ Box]
\sffamily\small
\begin{center}
\setlength{\tabcolsep}{8pt}
\begin{tabular}{c c c c}
\textbf{.box1} & \textbf{.box2} & \textbf{.box3} & \textbf{.box4} \\[2pt]
{\scriptsize\codeinline{<div></div>}} & {\scriptsize\codeinline{<div>   </div>}} & {\scriptsize\codeinline{<div><span></span></div>}} & {\scriptsize\codeinline{<div>Hello</div>}} \\[6pt]
% Box 1: Matches :empty (Red 2px border, pink background)
\fcolorbox{red}{pink!35}{\parbox[c][28pt][c]{28pt}{\centering\tiny\bfseries\textcolor{red!90!black}{:empty}}} &
% Box 2: Fails (Gray border)
\fcolorbox{gray!40}{gray!5}{\parbox[c][28pt][c]{28pt}{\centering\tiny\textcolor{gray!80}{(space)}}} &
% Box 3: Fails (Gray border)
\fcolorbox{gray!40}{gray!5}{\parbox[c][28pt][c]{28pt}{\centering\tiny\textcolor{gray!80}{<span>}}} &
% Box 4: Fails (Gray border)
\fcolorbox{gray!40}{gray!5}{\parbox[c][28pt][c]{28pt}{\centering\small\textcolor{darkText}{Hello}}} \\[6pt]
\textcolor{green!60!black}{\textbf{[Có] KHỚP :empty}} & \textcolor{red!80!black}{\textbf{[Không] Bị loại}} & \textcolor{red!80!black}{\textbf{[Không] Bị loại}} & \textcolor{red!80!black}{\textbf{[Không] Bị loại}} \\[2pt]
{\scriptsize\textcolor{green!60!black}{(Khung đỏ 30px, nền hồng)}} & {\scriptsize\textcolor{gray!70}{(Chứa khoảng trắng)}} & {\scriptsize\textcolor{gray!70}{(Chứa thẻ <span>)}} & {\scriptsize\textcolor{gray!70}{(Chứa chữ "Hello")}}
\end{tabular}
\end{center}
\end{browserWindow}

\begin{prettyTableBox}[Bảng Kết Quả Kiểm Thử Thực Tế Của 4 Thẻ Box]
\rowcolors{2}{white}{primaryBlue!4}
\sffamily\small
\begin{tabularx}{\linewidth}{|>{\raggedright\arraybackslash \ttfamily\color{accentOrange}\bfseries\raggedright\arraybackslash}p{2.2cm}|>{\centering\arraybackslash}p{3.2cm}|>{\raggedright\arraybackslash}X|}
\hline
\rowcolor{primaryBlue} \thd{\sffamily Phần tử} & \thd{\sffamily Khớp :empty?} & \thd{\sffamily Giải thích nguyên nhân bản chất trong DOM} \\
\hline
.box1 & \textcolor{green!60!black}{\textbf{[Có] PASS}} & Không chứa bất kỳ nút con nào trong DOM (0 Child Nodes). \\
\hline
.box2 & \textcolor{red}{\textbf{[Không] FAIL}} & Chứa ký tự khoảng trắng (Text Node khoảng trắng). \\
\hline
.box3 & \textcolor{red}{\textbf{[Không] FAIL}} & Chứa phần tử \codeinline{<span>} bên trong (Element Node). \\
\hline
.box4 & \textcolor{red}{\textbf{[Không] FAIL}} & Chứa nội dung văn bản chữ \codeinline{"Hello"} (Text Node). \\
\hline
\end{tabularx}
\end{prettyTableBox}

\begin{conceptBox}[Cơ Chế Hoạt Động Theo Từng Bước Của :empty Trong DOM]
Bộ chọn \codeinline{:empty} hoạt động dựa trên cơ chế kiểm tra sự tồn tại của các nút con (\term{Child Nodes}) trong mô hình cây DOM (Document Object Model):
\begin{enumerate}[leftmargin=1.5em, itemsep=3pt, parsep=0pt]
    \item \textbf{Bước 1}: Trình duyệt duyệt qua danh sách các phần tử HTML trên trang.
    \item \textbf{Bước 2}: Với mỗi phần tử, trình duyệt kiểm tra các nút con (\textit{child nodes}) trực thuộc nó trong cây DOM (bao gồm Text node, Element node, Comment node).
    \item \textbf{Bước 3}: Nếu \textbf{KHÔNG} có bất kỳ nút con nào ($\rightarrow$ 0 child nodes), phần tử được chọn bởi \codeinline{:empty}.
    \item \textbf{Bước 4}: Nếu \textbf{CÓ} bất kỳ nút con nào (dù chỉ là 1 khoảng trắng, 1 câu ghi chú comment hay 1 thẻ rỗng), phần tử bị loại khỏi \codeinline{:empty}.
\end{enumerate}
\par\vspace{4pt}
\textbf{Minh họa cơ chế chọn qua các dạng cấu trúc DOM:}
\begin{itemize}[leftmargin=1.5em, itemsep=2pt, parsep=0pt]
    \item \codeinline{<div></div>} $\rightarrow$ \textcolor{green!60!black}{\textbf{[ĐƯỢC CHỌN]}} Không có bất kỳ nút con nào trong DOM.
    \item \codeinline{<div> </div>} $\rightarrow$ \textcolor{red}{\textbf{[KHÔNG CHỌN]}} Có nút văn bản (Text Node chứa khoảng trắng).
    \item \codeinline{<div><span></span></div>} $\rightarrow$ \textcolor{red}{\textbf{[KHÔNG CHỌN]}} Có nút phần tử con (Element Node \codeinline{<span>}).
    \item \codeinline{<div><!-- comment --></div>} $\rightarrow$ \textcolor{red}{\textbf{[KHÔNG CHỌN]}} Có nút ghi chú (Comment Node).
\end{itemize}
\end{conceptBox}

\begin{warnBox}[5 Quy Tắc Vàng \& Lưu Ý Quan Trọng Khi Sử Dụng :empty]
\begin{enumerate}[leftmargin=1.5em, itemsep=4pt, parsep=0pt]
    \item \textbf{:empty chỉ kiểm tra nút con trong DOM:} Không chứa văn bản, không chứa thẻ con, không chứa ghi chú comment HTML.
    \item \textbf{Khoảng trắng và xuống dòng khiến phần tử KHÔNG rỗng:} Ví dụ \codeinline{<div>\textbackslash{}n</div>} hoặc \codeinline{<div>   </div>} không được coi là rỗng vì ký tự khoảng trắng/xuống dòng là một Text Node.
    \item \textbf{Phần tử phải thật sự "trống" 100\%:} Chỉ duy nhất cú pháp viết sát nhau không kẽ hở \codeinline{<div></div>} mới thỏa mãn.
    \item \textbf{:empty KHÔNG kiểm tra thuộc tính CSS hay kích thước hiển thị:} Dù phần tử có thuộc tính \codeinline{display: none}, \codeinline{visibility: hidden}, hay chiều cao bằng 0, nếu bên trong vẫn chứa nút con DOM thì vẫn \textbf{KHÔNG được coi là :empty}. Mặc dù nhìn bằng mắt thường có vẻ giống nhau, nhưng \codeinline{:empty} chỉ quan tâm tới cây DOM chứ không quan tâm đến giao diện hiển thị!
    \item \textbf{:empty giúp dọn rác và tối ưu giao diện HTML:} Tự động ẩn các đoạn văn rỗng hoặc các vùng chứa rỗng mà không cần JavaScript can thiệp.
\end{enumerate}
\end{warnBox}

\begin{tipBox}[Mẹo Kiểm Tra Chuẩn Xác Để Phần Tử Đạt :empty]
Muốn chắc chắn một phần tử được chọn bởi \codeinline{:empty}, hãy đảm bảo 3 quy tắc:
\begin{itemize}[leftmargin=1.5em, itemsep=2pt, parsep=0pt]
    \item \textbf{KHÔNG} viết bất kỳ ký tự nào giữa thẻ mở và thẻ đóng.
    \item \textbf{KHÔNG} để khoảng trắng hay xuống dòng giữa cặp thẻ.
    \item \textbf{KHÔNG} chứa comment hay bất kỳ thẻ con nào.
\end{itemize}
\end{tipBox}

\begin{prettyTableBox}[Tóm Tắt Khả Năng Khớp Selector :empty Trong Các Trường Hợp DOM]
\rowcolors{2}{white}{primaryBlue!4}
\sffamily\small
\begin{tabularx}{\linewidth}{|>{\raggedright\arraybackslash\bfseries\color{primaryBlue}}p{5.5cm}|>{\centering\arraybackslash}p{3.2cm}|>{\raggedright\arraybackslash}X|}
\hline
\rowcolor{primaryBlue} \thd{\sffamily Cấu trúc nội dung bên trong phần tử} & \thd{\sffamily Khớp :empty?} & \thd{\sffamily Ghi chú bản chất trong cây DOM} \\
\hline
Không có gì (Sát nhau \codeinline{<div></div>}) & \textcolor{green!60!black}{\textbf{[Có] PASS}} & Thỏa mãn điều kiện rỗng 100\% (0 Child Nodes). \\
\hline
Khoảng trắng / Xuống dòng (\codeinline{<div> </div>}) & \textcolor{red}{\textbf{[Không] FAIL}} & Chứa nút văn bản khoảng trắng (Text Node). \\
\hline
Comment HTML (\codeinline{<div><!-- comment --></div>}) & \textcolor{red}{\textbf{[Không] FAIL}} & Chứa nút ghi chú (Comment Node trong DOM). \\
\hline
Thẻ con rỗng (\codeinline{<div><span></span></div>}) & \textcolor{red}{\textbf{[Không] FAIL}} & Dù thẻ con rỗng vẫn tính là 1 Element Node. \\
\hline
Nội dung văn bản chữ (\codeinline{<div>Hello</div>}) & \textcolor{red}{\textbf{[Không] FAIL}} & Chứa nút nội dung văn bản (Text Node). \\
\hline
\end{tabularx}
\end{prettyTableBox}

\begin{conceptBox}[3 Kịch Bản Ứng Dụng Thực Tế Phổ Biến Của :empty]
Bộ chọn \codeinline{:empty} cực kỳ hữu ích trong lập trình Frontend thực tế:
\begin{enumerate}[leftmargin=1.5em, itemsep=4pt, parsep=0pt]
    \item \textbf{Ẩn các phần tử rỗng tự động (Dọn dẹp layout):} Bạn có thể ẩn các thẻ đoạn văn \codeinline{<p>} hoặc vùng chứa không có nội dung để giao diện gọn gàng hơn.
\begin{codeWindow}[Ẩn Thẻ Đoạn Văn Rỗng Trong CSS]
\begin{lstlisting}[language=CSS3]
p:empty {
  display: none;
}
\end{lstlisting}
\end{codeWindow}

    \item \textbf{Tạo placeholder hoặc thông báo cảnh báo:} Hiển thị khung viền hoặc tô nền cảnh báo cho các vùng chưa có dữ liệu (ví dụ: khung tải dữ liệu hoặc khung chứa ảnh bị lỗi).
\begin{codeWindow}[Tạo Viền \& Nền Hồng Cảnh Báo Cho Div Rỗng]
\begin{lstlisting}[language=CSS3]
div:empty {
  border: 1px dashed red;
  background-color: pink;
}
\end{lstlisting}
\end{codeWindow}

    \item \textbf{Kiểm soát hiển thị động với JavaScript \& Pseudo-element:} Khi nội dung được thêm hoặc xóa bằng JS, \codeinline{:empty} giúp tự động cập nhật giao diện mà không cần thêm logic CSS phức tạp.
\begin{codeWindow}[Tự Động Chèn Chữ "Rỗng" Cho Phần Tử Không Có Dữ Liệu]
\begin{lstlisting}[language=CSS3]
.data-box:empty::before {
  content: "Rong (Chua co du lieu)";
  color: gray;
  font-style: italic;
}
\end{lstlisting}
\end{codeWindow}
\end{enumerate}
\end{conceptBox}

%% ===================================================================
%% MỞ RỘNG KIẾN THỨC NÂNG CAO VỀ :empty (SENIOR FRONTEND MASTERCLASS)
%% ===================================================================

\paragraph{Mở Rộng Kiến Thức \& Kinh Nghiệm Thực Tế Chuyên Sâu Về Selector :empty}
\mbox{}\par\vspace{2pt}

\begin{conceptBox}[1. Bẫy Định Dạng Trong React/JSX/Template Engine Với :empty]
Trong lập trình Frontend hiện đại (React JSX, Vue, EJS, Blade Template), các kỹ sư rất dễ mắc lỗi làm cho \codeinline{:empty} \textbf{thất bại im lặng} do khoảng trắng sinh ra từ code formatting:
\begin{codeWindow}[Lỗi Phổ Biến Tạo Khoảng Trắng Im Lặng Trong React/JSX]
\begin{lstlisting}[language=HTML5]
<!-- Trong JSX: Viet xuong dong hoac co whitespace giua {} -->
<div class="user-badge">
  {user.title} 
</div>

<!-- Ket qua HTML sinh ra DOM Tree: -->
<div class="user-badge"> </div> 
<!-- -> CO NHO KHOANG TRANG -> THAT BAI :empty! -->
\end{lstlisting}
\end{codeWindow}
\par\vspace{4pt}
\textbf{Giải pháp khắc phục Senior:} 
\begin{itemize}[leftmargin=1.5em, itemsep=2pt, parsep=0pt]
    \item Render có điều kiện: \codeinline{\{user.title ? <div class="user-badge">\{user.title\}</div> : null\}}
    \item Đảm bảo thẻ tự đóng hoặc sát nhau 100\% khi chuỗi rỗng: \codeinline{<div class="user-badge">\{user.title\}</div>}
\end{itemize}
\end{conceptBox}

\begin{tipBox}[2. Kỹ Thuật CSS Skeleton Loading Tự Động Với :empty]
Trong các ứng dụng Web gọi API bất đồng bộ, ta có thể kết hợp \codeinline{:empty} với \codeinline{@keyframes} để tạo hiệu ứng \textbf{Skeleton Loading} (mô phỏng khung xương tải dữ liệu) cực kỳ chuyên nghiệp mà không cần dùng đến thẻ HTML phụ:
\begin{codeWindow}[Tạo Hiệu Ứng Skeleton Loading Bằng CSS :empty]
\begin{lstlisting}[language=CSS3]
/* Khi container chua co du lieu tu API nhan ve (dang empty) */
.profile-name:empty {
  width: 150px;
  height: 20px;
  background: linear-gradient(90deg, #eee 25%, #e0e0e0 50%, #eee 75%);
  background-size: 200% 100%;
  animation: skeleton-shimmer 1.5s infinite;
  border-radius: 4px;
}

@keyframes skeleton-shimmer {
  0% { background-position: 200% 0; }
  100% { background-position: -200% 0; }
}
\end{lstlisting}
\end{codeWindow}
\textbf{Ưu điểm vượt trội:} Khi API tải xong và chèn chữ vào \codeinline{.profile-name}, thẻ lập tức mất tính chất \codeinline{:empty}, hiệu ứng Skeleton tự động biến mất và hiển thị văn bản thật mượt mà!
\end{tipBox}

\begin{warnBox}[3. Bộ Câu Hỏi Phỏng Vấn Senior Frontend Về Selector :empty]
\begin{enumerate}[leftmargin=1.5em, itemsep=5pt, parsep=0pt]
    \item \textbf{Câu hỏi 1: Thẻ tự đóng Void Elements (<input>, <img>, <br>) có được chọn bởi :empty không?} \par
    $\rightarrow$ \textbf{Trả lời:} \textbf{CÓ!} Vì các thẻ như \codeinline{<input>}, \codeinline{<img>}, \codeinline{<hr>}, \codeinline{<br>} là các phần tử rỗng (\term{Void Elements}), không có bất kỳ nút con nào trong cây DOM (0 Child Nodes). Do đó, \codeinline{input:empty} hay \codeinline{img:empty} luôn thỏa mãn điều kiện \codeinline{:empty}!

    \item \textbf{Câu hỏi 2: Thẻ có phần tử giả ::before hoặc ::after thì có khớp :empty không?} \par
    $\rightarrow$ \textbf{Trả lời:} \textbf{CÓ!} Bộ chọn phần tử giả \codeinline{::before} và \codeinline{::after} được sinh ra bởi CSS trên cây hiển thị (\term{Render Tree}), chúng \textbf{KHÔNG hề thêm nút con (\term{Child Node}) vào cây DOM}. Do đó, một thẻ \codeinline{<div class="box"></div>} có định kiểu \codeinline{.box::before \{ content: "Hi"; \}} thì thẻ đó \textbf{vẫn khớp \codeinline{div:empty} 100\%}!

    \item \textbf{Câu hỏi 3: Bộ chọn modern :has() và :blank khác gì với :empty?} \par
    $\rightarrow$ \textbf{Trả lời:} Bộ chọn \codeinline{:blank} (CSS Selectors Level 4) được thiết kế để chọn cả phần tử chứa khoảng trắng whitespace. Ngoài ra, ta có thể dùng bộ chọn modern \codeinline{:not(:has(*))} để chọn phần tử không chứa bất kỳ thẻ HTML con nào dù cho có chứa khoảng trắng.
\end{enumerate}
\end{warnBox}

\paragraph{j. Bảng Master Cheatsheet Tra Cứu Toàn Diện 11 Pseudo-Classes Cấu Trúc}

\begin{prettyTableBox}[Bảng Tra Cứu Toàn Diện Tất Cả 11 Pseudo-Classes Vị Trí Cấu Trúc và Phủ Định]
\rowcolors{2}{white}{primaryBlue!4}
\sffamily\small
\begin{tabularx}{\linewidth}{|>{\raggedright\arraybackslash \ttfamily\color{accentOrange}\bfseries}m{4.2cm}|Z|}
\hline
\rowcolor{primaryBlue} \thd{\sffamily Selector} & \thd{\sffamily Nguyên lý \& Phạm vi hoạt động} \\
\hline
:first-child & Chọn phần tử nếu nó là con đầu tiên tuyệt đối của cha. \\
\hline
:last-child & Chọn phần tử nếu nó là con cuối cùng tuyệt đối của cha. \\
\hline
:nth-child(n) & Chọn phần tử đứng ở vị trí thứ $n$ (từ trên xuống). \\
\hline
:nth-last-child(n) & Chọn phần tử đứng ở vị trí thứ $n$ (từ dưới đếm lên). \\
\hline
:only-child & Chọn phần tử nếu nó là con duy nhất trong cha. \\
\hline
:first-of-type & Chọn phần tử cùng loại thẻ đầu tiên trong cha. \\
\hline
:last-of-type & Chọn phần tử cùng loại thẻ cuối cùng trong cha. \\
\hline
:nth-of-type(n) & Chọn phần tử cùng loại thẻ thứ $n$ (từ trên xuống). \\
\hline
:nth-last-of-type(n) & Chọn phần tử cùng loại thẻ thứ $n$ (từ dưới đếm lên). \\
\hline
:only-of-type & Chọn phần tử nếu nó là thẻ duy nhất thuộc loại đó trong cha. \\
\hline
:not(selector) & Chọn tất cả các phần tử loại trừ phần tử khớp với selector bên trong. \\
\hline
\end{tabularx}
\end{prettyTableBox}

\paragraph{k. Đọc Thêm: Chuyên Đề Nâng Cao, Nguyên Lý Trình Duyệt \& Kinh Nghiệm Thực Chiến}

\begin{conceptBox}[1. Cú Pháp Modern CSS Selectors Level 4: :nth-child(An+B of Selector)]
Trình duyệt hiện đại (Chrome 111+, Firefox 113+, Safari 16.4+) hỗ trợ cú pháp mở rộng mạnh mẽ.
\par\vspace{4pt}
\textbf{Sự khác nhau}: Cú pháp cũ \codeinline{.active:nth-child(2)} chỉ kiểm tra xem phần tử thứ 2 có class \codeinline{.active} hay không. Cú pháp mới \codeinline{li:nth-child(2 of .active)} sẽ \textbf{lọc tất cả thẻ li có class .active trước}, sau đó chọn thẻ thứ 2 trong tập hợp đó!
\end{conceptBox}

\begin{codeWindow}[Ví Dụ Cú Pháp Modern CSS4 :nth-child(An+B of Selector)]
\begin{lstlisting}[language=CSS3]
/* Select 2nd element matching .active class specifically */
li:nth-child(2 of .active) { color: red; font-weight: bold; }
\end{lstlisting}
\end{codeWindow}

\begin{warnBox}[2. Nguyên Lý Trình Duyệt \& Tối Ưu Hiệu Năng Rendering (Reflow / Repaint)]
\begin{itemize}[leftmargin=1.5em, itemsep=3pt, parsep=0pt]
    \item \textbf{Cơ chế tính toán DOM Tree}: Khi trình duyệt duyệt cây DOM để match các bộ chọn \codeinline{:nth-child(an+b)}, trình duyệt phải đếm vị tríindex của từng phần tử con.
    \item \textbf{Cảnh báo Performance}: Trong các danh sách cực lớn (hàng nghìn dòng DOM), việc dùng các công thức phức tạp hoặc lồng ghép nhiều pseudo-classes cấu trúc có thể làm tăng thời gian tính toán Layout (\term{Reflow}) khi người dùng cuộn trang hoặc chèn thêm dữ liệu động.
    \item \textbf{Best Practice}: Với danh sách rất dài, nên kết hợp với kỹ thuật \term{Virtual Scrolling} hoặc dùng class cố định được render từ phía Server/JavaScript.
\end{itemize}
\end{warnBox}

\begin{tipBox}[3. Bộ Câu Hỏi Phỏng Vấn Tuyển Dụng Frontend Thường Gặp (Interview Q\&A)]
\begin{enumerate}[leftmargin=1.5em, itemsep=4pt, parsep=0pt]
    \item \textbf{Q1: Phân biệt sự khác nhau giữa \codeinline{:nth-child(2)} và \codeinline{:nth-of-type(2)}?} \\
    $\rightarrow$ \textbf{Trả lời ghi điểm}: \codeinline{:nth-child(2)} đếm vị trí thứ 2 tuyệt đối trong danh sách con của cha; nếu phần tử thứ 2 không đúng loại thẻ chỉ định thì rule bị hủy. Trong khi đó, \codeinline{:nth-of-type(2)} chỉ đếm trong tập hợp các phần tử cùng kiểu thẻ.
    \item \textbf{Q2: Tại sao bộ chọn \codeinline{:empty} lại không hoạt động khi thẻ HTML có xuống dòng hoặc khoảng trắng?} \\
    $\rightarrow$ \textbf{Trả lời ghi điểm}: Vì trong mô hình DOM HTML, khoảng trắng (whitespace) và dấu xuống dòng (newline) được trình duyệt tính là một Text Node (\codeinline{\#text}). Do đó thẻ có khoảng trắng không còn rỗng hoàn toàn nữa. Muốn chọn thẻ rỗng kể cả có khoảng trắng trong CSS hiện đại, ta dùng bộ chọn \codeinline{:blank}.
\end{enumerate}
\end{tipBox}

\begin{conceptBox}[4. Lưu Ý Về Khả Năng Truy Cập (Accessibility - a11y)]
Khi thay đổi thứ tự hiển thị trực quan hoặc màu sắc của các phần tử bằng \codeinline{:first-child} hay \codeinline{:last-child}, hãy đảm bảo thứ tự đọc của công cụ đọc màn hình (\term{Screen Reader}) và phím Tab (\term{Keyboard Navigation}) vẫn tuân theo đúng thứ tự ngữ nghĩa trong cây DOM HTML.
\end{conceptBox}

% ==============================================================================
% MỤC 8: BỘ CHỌN PHẦN TỬ GIẢ (PSEUDO-ELEMENT SELECTORS MASTERCLASS)
% ==============================================================================
\subsection{8. Chuyên Đề Phân Tích Chi Tiết \& Thực Nghiệm Về Bộ Chọn Phần Tử Giả (Pseudo-element Selectors Masterclass)}

Pseudo-element selector (bộ chọn phần tử giả) là các bộ chọn trong CSS cho phép bạn định kiểu một phần cụ thể của phần tử, mà không cần phải thêm thẻ HTML mới. Nó giống như tạo ra một phần tử ảo trong cấu trúc trang web!

Dưới đây là chi tiết về từng loại pseudo-element phổ biến và cách hoạt động của chúng:

\begin{prettyTableBox}[Bảng Tra Cứu Tóm Tắt 6 Pseudo-Element Selectors Phổ Biến]
\rowcolors{2}{white}{primaryBlue!4}
\sffamily\small
\begin{tabularx}{\linewidth}{|>{\raggedright\arraybackslash \ttfamily\color{accentOrange}\bfseries}m{2.8cm}|>{\bfseries\color{primaryBlue}}m{3.5cm}|>{\ttfamily\small}m{4.2cm}|Z|}
\hline
\rowcolor{primaryBlue} \thd{\sffamily Pseudo-element} & \thd{\sffamily Ý nghĩa} & \thd{\sffamily Ví dụ cú pháp} & \thd{\sffamily Mô tả chi tiết} \\
\hline
::before & Thêm nội dung trước nội dung của phần tử. & p::before \{ content: "[ICON] "; \} & Chèn icon hoặc văn bản vào đầu phần tử. \\
\hline
::after & Thêm nội dung sau nội dung của phần tử. & p::after \{ content: " *"; \} & Chèn icon hoặc văn bản vào cuối phần tử. \\
\hline
::first-letter & Chọn ký tự đầu tiên của nội dung phần tử. & p::first-letter \{ font-size: 2em; \} & Thường dùng để làm chữ cái đầu lớn hơn. \\
\hline
::first-line & Chọn dòng đầu tiên của nội dung phần tử. & p::first-line \{ color: blue; \} & Định dạng dòng đầu tiên (thường gặp ở bài viết). \\
\hline
::selection & Chọn phần nội dung được bôi đen (highlight). & p::selection \{ background: yellow; \} & Thay đổi màu nền khi người dùng chọn văn bản. \\
\hline
::marker & Định dạng ký hiệu của danh sách (bullet, số). & li::marker \{ color: red; \} & Thay đổi màu hoặc kiểu ký hiệu danh sách. \\
\hline
\end{tabularx}
\end{prettyTableBox}

\begin{warnBox}[Lưu Ý Quan Trọng Về Pseudo-Elements]
\begin{itemize}[leftmargin=1.5em, itemsep=4pt, parsep=0pt]
    \item \textbf{Cú pháp chuẩn:} Sử dụng hai dấu hai chấm \codeinline{::} để viết pseudo-element. (Đây là chuẩn CSS3 mới nhằm phân biệt với pseudo-class dùng 1 dấu chấm hai chấm \codeinline{:}, nhưng viết \codeinline{:} vẫn hoạt động trong một số trình duyệt cũ vì mục đích tương thích ngược).
    \item \textbf{Không tác động trực tiếp đến DOM:} Các phần tử giả này \textbf{không thật sự xuất hiện trong cây DOM} (Document Object Model), mà chỉ tồn tại trên giao diện hiển thị do trình duyệt dựng nên.
\end{itemize}
\end{warnBox}

\begin{conceptBox}[Ví Dụ Thực Tế Minh Họa Pseudo-Elements]
Dưới đây là ví dụ kết hợp bộ chọn \codeinline{::before} và \codeinline{::after} để chèn ký tự trang trí ở hai đầu văn bản:
\end{conceptBox}

\begin{codeWindow}[Mã Nguồn HTML \& CSS Minh Họa Chèn Ký Tự Trang Trí]
\begin{lstlisting}[language=HTML5]
<div class="box">Chao ban!</div>

<style>
  .box::before {
    content: "[ICON] ";
    color: gold;
  }
  .box::after {
    content: " [ICON]";
    color: red;
  }
</style>
\end{lstlisting}
\end{codeWindow}

\begin{browserWindow}[Kết Quả Hiển Thị Trên Trình Duyệt]
\sffamily\large\bfseries
\textbf{Kết quả:} [ICON] Chào bạn! [ICON]
\end{browserWindow}

\begin{conceptBox}[Ứng Dụng Thực Tế Của Pseudo-Elements]
\begin{itemize}[leftmargin=1.5em, itemsep=3pt, parsep=0pt]
    \item \textbf{Tạo hiệu ứng trang trí:} Thêm icon, ký hiệu, hoặc hình ảnh nhỏ vào giao diện.
    \item \textbf{Trang trí văn bản:} Làm chữ cái đầu lớn (drop cap) như trên các bài báo in hoặc sách cổ.
    \item \textbf{Cải thiện UX/UI:} Highlight nội dung tùy chỉnh khi người dùng thao tác bôi đen văn bản.
    \item \textbf{Tối ưu mã HTML:} Giảm số lượng thẻ HTML thừa (như thẻ \codeinline{<span>} rỗng) không cần thiết, giúp mã nguồn HTML sạch và chuẩn ngữ nghĩa.
\end{itemize}
\end{conceptBox}

\paragraph{Phân tích chi tiết từng pseudo-element selector để bạn hiểu rõ và áp dụng dễ dàng:}

%% -------------------------------------------------------------------
%% 8.1. ::BEFORE
%% -------------------------------------------------------------------
\subsubsection{8.1. Bộ Chọn ::before (Chèn Nội Dung Vào Trước Phần Tử)}

\begin{codeWindow}[Cú Pháp Cơ Bản Của ::before]
\begin{lstlisting}[language=CSS3]
selector::before {
  content: "Noi dung them vao";
  /* Cac thuoc tinh CSS khac */
}
\end{lstlisting}
\end{codeWindow}

\begin{conceptBox}[Giải Thích Chi Tiết Cấu Trúc ::before]
Chèn nội dung vào \textbf{TRƯỚC} nội dung bên trong của phần tử, nhưng không làm thay đổi nội dung thực tế trong cấu trúc mã HTML.
\begin{itemize}[leftmargin=1.5em, itemsep=3pt, parsep=0pt]
    \item \codeinline{selector}: Phần tử mà bạn muốn thêm nội dung phía trước.
    \item \codeinline{::before}: Bộ chọn phần tử giả thêm nội dung ảo vào trước phần tử được chọn.
    \item \codeinline{content}: \textbf{Bắt buộc phải có}, xác định nội dung sẽ hiển thị.
\end{itemize}
\par\vspace{4pt}
\textbf{Lưu ý quan trọng:} Nếu không có thuộc tính \codeinline{content}, thì \codeinline{::before} sẽ \textbf{không hiển thị bất kỳ thứ gì trên giao diện}!
\end{conceptBox}

\begin{codeWindow}[Ví Dụ Minh Họa Sử Dụng ::before]
\begin{lstlisting}[language=HTML5]
<!-- VD 1: Khi khong dung ::before -->
<h2>Tieu de bai viet</h2>

<!-- Khi dung ::before ap dung vao CSS -->
<h2 class="title">Tieu de bai viet</h2>

<style>
  .title::before {
    content: "[ICON] ";
    color: orange;
    font-size: 24px;
  }
</style>
\end{lstlisting}
\end{codeWindow}

\begin{browserWindow}[Kết Quả Hiển Thị ::before]
\sffamily\Large\bfseries
[ICON] Tiêu đề bài viết
\end{browserWindow}

\begin{conceptBox}[VD 2 \& Ứng Dụng Thực Tế Của ::before]
\textbf{Ứng dụng:} Thêm biểu tượng, dấu hiệu nhấn mạnh, hoặc trang trí tiêu đề mà không cần thay đổi cấu trúc mã HTML gốc.
\end{conceptBox}

%% -------------------------------------------------------------------
%% 8.2. ::AFTER
%% -------------------------------------------------------------------
\subsubsection{8.2. Bộ Chọn ::after (Chèn Nội Dung Vào Sau Phần Tử)}

\begin{codeWindow}[Cú Pháp Cơ Bản Của ::after]
\begin{lstlisting}[language=CSS3]
selector::after {
  content: "Noi dung them vao";
  /* Cac thuoc tinh CSS khac */
}
\end{lstlisting}
\end{codeWindow}

\begin{conceptBox}[Giải Thích Chi Tiết Cấu Trúc ::after]
Chèn nội dung vào \textbf{SAU} nội dung bên trong của phần tử.
\begin{itemize}[leftmargin=1.5em, itemsep=3pt, parsep=0pt]
    \item \codeinline{selector}: Phần tử mà bạn muốn thêm nội dung sau nó.
    \item \codeinline{::after}: Thêm phần tử giả ngay sau nội dung bên trong phần tử.
    \item \codeinline{content}: \textbf{Bắt buộc phải có}, nội dung muốn chèn thêm. Nếu không có \codeinline{content}, \codeinline{::after} không hiển thị gì cả!
\end{itemize}
\end{conceptBox}

\begin{codeWindow}[Ví Dụ Minh Họa Sử Dụng ::after]
\begin{lstlisting}[language=HTML5]
/* VD khai bao CSS dang sau */
.after::after {
  content: "Noi dung dang sau";
}

<!-- Ap dung thuc te cho doan luu y -->
<p class="note">Luu y quan trong</p>

<style>
  .note::after {
    content: " [WARNING]";
    color: red;
    font-weight: bold;
  }
</style>
\end{lstlisting}
\end{codeWindow}

\begin{browserWindow}[Kết Quả Hiển Thị ::after]
\sffamily\large\bfseries
Lưu ý quan trọng [WARNING]
\end{browserWindow}

\begin{conceptBox}[Ứng Dụng Thực Tế Của ::after]
\textbf{Ứng dụng:} Thêm icon cảnh báo, tạo dấu kết thúc câu, tạo hiệu ứng đường kẻ gạch chân trang trí, hoặc tạo clearfix dọn dẹp float trong CSS.
\end{conceptBox}

%% -------------------------------------------------------------------
%% 8.3. ::FIRST-LETTER
%% -------------------------------------------------------------------
\subsubsection{8.3. Bộ Chọn ::first-letter (Định Dạng Ký Tự Đầu Tiên)}

\begin{conceptBox}[Khái Niệm \& Cấu Trúc Căn Bản Của ::first-letter]
\codeinline{::first-letter} là một pseudo-element selector trong CSS, dùng để định dạng \textbf{chữ cái đầu tiên} của phần tử văn bản dạng khối (block element).
\par\vspace{4pt}
Thường dùng để tạo \term{Drop Cap} (chữ cái lớn ở đầu đoạn văn như trong sách, báo in chuyên nghiệp).
\end{conceptBox}

\begin{codeWindow}[Ví Dụ Đơn Giản Định Kiểu Chữ Cái Đầu]
\begin{lstlisting}[language=HTML5]
<p class="text">Hoc CSS rat thu vi va huu ich.</p>

<style>
  .text::first-letter {
    font-size: 36px;
    color: blue;
    font-weight: bold;
  }
</style>
\end{lstlisting}
\end{codeWindow}

\begin{browserWindow}[Kết Quả Hiển Thị]
\sffamily\small
\textbf{\Large\textcolor{blue}{H}}ọc CSS rất thú vị và hữu ích. \par\vspace{2pt}
{\scriptsize\textcolor{gray}{\textit{Kết quả: Chữ cái đầu "H" to hơn, màu xanh.}}}
\end{browserWindow}

\begin{codeWindow}[Cấu Trúc Cơ Bản Của ::first-letter]
\begin{lstlisting}[language=CSS3]
selector::first-letter {
  /* Cac thuoc tinh CSS ap dung cho chu cai dau tien */
}
\end{lstlisting}
\end{codeWindow}

\begin{conceptBox}[Giải Thích Chi Tiết ::first-letter]
\begin{itemize}[leftmargin=1.5em, itemsep=3pt, parsep=0pt]
    \item \codeinline{selector}: Chọn phần tử cần định dạng.
    \item \codeinline{::first-letter}: Chọn chữ cái đầu tiên trong phần tử đó để áp dụng style.
\end{itemize}
\end{conceptBox}

\begin{codeWindow}[Ví Dụ Minh Họa Chi Tiết Với Ký Tự Đầu]
\begin{lstlisting}[language=HTML5]
<p class="intro">Chao mung ban den voi the gioi CSS!</p>

<style>
  .intro::first-letter {
    font-size: 40px;
    color: red;
    font-weight: bold;
    margin-right: 5px;
  }
</style>
\end{lstlisting}
\end{codeWindow}

\begin{browserWindow}[Kết Quả Hiển Thị Trực Quan]
\sffamily\small
\textbf{\Huge\textcolor{red}{C}}hào mừng bạn đến với thế giới CSS!
\end{browserWindow}

\begin{conceptBox}[Phân Tích Các Thuộc Tính CSS Ảnh Hưởng Đến Ký Tự Đầu]
Các thuộc tính áp dụng trong CSS sẽ ảnh hưởng trực tiếp đến chữ cái đầu tiên:
\begin{itemize}[leftmargin=1.5em, itemsep=3pt, parsep=0pt]
    \item \codeinline{font-size: 40px} $\rightarrow$ Làm chữ cái đầu to lên.
    \item \codeinline{color: red} $\rightarrow$ Đổi màu chữ thành đỏ.
    \item \codeinline{font-weight: bold} $\rightarrow$ Chữ in đậm.
    \item \codeinline{margin-right: 5px} $\rightarrow$ Tạo khoảng cách với ký tự tiếp theo.
\end{itemize}
\end{conceptBox}

\begin{warnBox}[Thuộc Tính CSS Phổ Biến Hỗ Trợ \& Không Hỗ Trợ Cho ::first-letter]
Không phải thuộc tính CSS nào cũng áp dụng được cho \codeinline{::first-letter}! Đây là danh sách phân loại chính xác:
\par\vspace{4pt}
\textbf{Các thuộc tính hợp lệ được hỗ trợ:}
\begin{itemize}[leftmargin=1.5em, itemsep=2pt, parsep=0pt]
    \item \textbf{Kiểu chữ:} \codeinline{font}, \codeinline{font-size}, \codeinline{font-weight}, \codeinline{font-style}, \codeinline{font-variant}
    \item \textbf{Màu sắc và nền:} \codeinline{color}, \codeinline{background}, \codeinline{background-clip}, \codeinline{background-origin}
    \item \textbf{Khoảng cách và viền:} \codeinline{margin}, \codeinline{padding}, \codeinline{border}, \codeinline{border-radius}
    \item \textbf{Text style:} \codeinline{text-decoration}, \codeinline{text-transform}, \codeinline{letter-spacing}, \codeinline{line-height}
\end{itemize}
\par\vspace{4pt}
\textbf{Không hỗ trợ (hoặc bị hạn chế trong hầu hết ngữ cảnh tiêu chuẩn):} \codeinline{box-shadow}, \codeinline{width}, \codeinline{height}, \codeinline{float}, \codeinline{position}, \codeinline{transform}...
\end{warnBox}

\begin{codeWindow}[Ứng Dụng Thực Tế 1: Tạo Hiệu Ứng Chữ Cái Đầu Lớn Như Sách Cổ (Classic Drop Cap)]
\begin{lstlisting}[language=CSS3]
p::first-letter {
  font-size: 60px;
  color: #8B0000;
  float: left;
  margin-right: 10px;
  font-family: "Georgia", serif;
}
\end{lstlisting}
\end{codeWindow}

\begin{codeWindow}[Ứng Dụng Thực Tế 2: Trang Trí Tiêu Đề Bắt Mắt]
\begin{lstlisting}[language=CSS3]
h2::first-letter {
  font-size: 50px;
  text-transform: uppercase;
  color: #FF4500;
}
\end{lstlisting}
\end{codeWindow}

\begin{codeWindow}[Ứng Dụng Thực Tế 3: Tạo Điểm Nhấn Nền Vàng Cho Đoạn Văn]
\begin{lstlisting}[language=CSS3]
article::first-letter {
  font-size: 30px;
  background-color: #FFD700;
  padding: 5px;
  border-radius: 3px;
}
\end{lstlisting}
\end{codeWindow}

\begin{conceptBox}[Kết Luận Về ::first-letter]
\begin{itemize}[leftmargin=1.5em, itemsep=3pt, parsep=0pt]
    \item \codeinline{::first-letter} cực kỳ hữu ích để làm nổi bật chữ cái đầu tiên.
    \item Thường dùng trong bài viết, sách, hoặc blog để tạo điểm nhấn nghệ thuật.
    \item Kết hợp với các thuộc tính như \codeinline{float}, \codeinline{margin} có thể tạo ra hiệu ứng chữ cái lớn (drop cap) như trong sách in truyền thống!
\end{itemize}
\end{conceptBox}

%% -------------------------------------------------------------------
%% 8.4. ::FIRST-LINE
%% -------------------------------------------------------------------
\subsubsection{8.4. Bộ Chọn ::first-line (Định Dạng Dòng Đầu Tiên Của Phần Tử Khối)}

\begin{conceptBox}[Bản Chất \& Cơ Chế Hoạt Động Của ::first-line]
Bộ chọn \codeinline{::first-line} nhắm đến **dòng đầu tiên** của phần tử văn bản dạng khối (block element).
\par\vspace{4pt}
\textbf{Cấu trúc khai báo tổng quát:}
\begin{codeWindow}[Cấu Trúc Khai Báo ::first-line]
\begin{lstlisting}[language=CSS3]
selector::first-line {
  color: blue;
  font-weight: bold;
  font-size: 18px;
}
\end{lstlisting}
\end{codeWindow}
\end{conceptBox}

\begin{warnBox}[Phân Loại Thuộc Tính Hỗ Trợ \& Không Áp Dụng Được Cho ::first-line]
Do yếu tố tối ưu thuật toán dàn trang (Reflow / Repaint), trình duyệt **chỉ cho phép** áp dụng các thuộc tính liên quan đến văn bản và màu nền cho \codeinline{::first-line}:
\par\vspace{4pt}
\textbf{Các thuộc tính được phép hỗ trợ:}
\begin{itemize}[leftmargin=1.5em, itemsep=2pt, parsep=0pt]
    \item \textbf{Màu sắc \& Kiểu chữ:} \codeinline{color}, \codeinline{font}, \codeinline{font-size}, \codeinline{font-weight}, \codeinline{font-family}, \codeinline{font-style}, \codeinline{font-variant}.
    \item \textbf{Định dạng văn bản:} \codeinline{letter-spacing}, \codeinline{word-spacing}, \codeinline{text-transform}, \codeinline{text-decoration}, \codeinline{line-height}, \codeinline{clear}.
    \item \textbf{Màu nền:} \codeinline{background}, \codeinline{background-color}.
\end{itemize}
\par\vspace{4pt}
\textbf{Các thuộc tính KHÔNG áp dụng được:} \codeinline{margin}, \codeinline{padding}, \codeinline{border}, \codeinline{width}, \codeinline{height}, \codeinline{position}, \codeinline{float}...
\end{warnBox}

\begin{codeWindow}[Ví Dụ Minh Họa Định Kiểu Dòng Đầu Tiên]
\begin{lstlisting}[language=HTML5]
<p class="text">
  Day la noi dung vi du de minh hoa cach hoat dong cua pseudo-element ::first-line.
  Dong dau tien se co mau xanh va in dam.
</p>

<style>
  .text::first-line {
    color: blue;
    font-weight: bold;
    font-size: 18px;
  }
</style>
\end{lstlisting}
\end{codeWindow}

\begin{browserWindow}[Kết Quả Hiển Thị ::first-line Trực Quan]
\sffamily\small
\textbf{\textcolor{blue}{\Large ĐÂY LÀ NỘI DUNG VÍ DỤ ĐỂ MINH HỌA CÁCH HOẠT ĐỘNG CỦA PSEUDO-ELEMENT ::FIRST-LINE.}} \par\vspace{2pt}
Dòng đầu tiên sẽ có màu xanh, in đậm và phóng to chữ, nhưng các dòng sau không bị ảnh hưởng.
\end{browserWindow}

\begin{tipBox}[2 Lưu Ý Quan Trọng Khi Sử Dụng ::first-line]
\begin{enumerate}[leftmargin=1.5em, itemsep=3pt, parsep=0pt]
    \item \textbf{Ảnh hưởng bởi Responsive Viewport:} Số lượng từ nằm trong "dòng đầu tiên" phụ thuộc hoàn toàn vào chiều rộng màn hình. Khi co giãn cửa sổ trình duyệt, dòng đầu tiên sẽ tự động thay đổi các từ thuộc về nó.
    \item \textbf{Không hoạt động trên phần tử inline:} Chỉ hoạt động trên các phần tử dạng khối (\codeinline{display: block}, \codeinline{inline-block}, \codeinline{list-item}) như \codeinline{<p>}, \codeinline{<div>}, \codeinline{<article>}. KHÔNG hoạt động trên thẻ inline như \codeinline{<span>}, \codeinline{<a>}.
\end{enumerate}
\end{tipBox}

\begin{conceptBox}[Ứng Dụng Thực Tế Của ::first-line]
\begin{itemize}[leftmargin=1.5em, itemsep=3pt, parsep=0pt]
    \item \textbf{Tạo điểm nhấn cho đoạn văn bản:} Làm nổi bật câu mở đầu trong bài viết, báo chí hoặc đoạn nội dung dài.
    \item \textbf{Thiết kế tiêu đề sáng tạo:} Làm dòng đầu tiên ấn tượng hơn các dòng theo sau.
\end{itemize}
\end{conceptBox}

%% -------------------------------------------------------------------
%% 8.5. ::SELECTION
%% -------------------------------------------------------------------
\subsubsection{8.5. Bộ Chọn ::selection (Tùy Chỉnh Văn Bản Khi Người Dùng Bôi Đen)}

\begin{conceptBox}[Định Nghĩa \& Cấu Trúc Căn Bản Của ::selection]
\codeinline{::selection} là một pseudo-element selector dùng để định dạng (\term{styling}) phần nội dung mà người dùng bôi đen (\term{highlight}) trên trang web bằng chuột hoặc thao tác vuốt cảm ứng!
\par\vspace{4pt}
\textbf{Cú pháp chuẩn:}
\begin{codeWindow}[Cú Pháp Định Kiểu ::selection]
\begin{lstlisting}[language=CSS3]
::selection {
  background: #ffeb3b; /* Mau nen vang sang */
  color: #d50000;      /* Mau chu do dam */
}
\end{lstlisting}
\end{codeWindow}
\end{conceptBox}

\begin{codeWindow}[Ví Dụ Minh Họa Nâng Cao Với ::selection]
\begin{lstlisting}[language=HTML5]
<p class="custom-select">Hay boi den doan van ban nay de xem hieu ung cua ::selection!</p>

<style>
  .custom-select::selection {
    background: #ffeb3b;
    color: #d50000;
    text-shadow: 1px 1px 2px #000;
  }
</style>
\end{lstlisting}
\end{codeWindow}

\begin{browserWindow}[Kết Quả Hiển Thị Khi Bôi Đen Nội Dung]
\sffamily\small
Khi bạn bôi đen đoạn văn bản, nó sẽ hiển thị: \colorbox{yellow}{\textcolor{red!80!black}{\textbf{Nền vàng sáng, chữ đỏ đậm, bóng chữ}}}.
\end{browserWindow}

\begin{prettyTableBox}[Bảng Phân Loại Thuộc Tính Hỗ Trợ Cho ::selection]
\rowcolors{2}{white}{primaryBlue!4}
\sffamily\small
\begin{tabularx}{\linewidth}{|>{\raggedright\arraybackslash \bfseries\color{primaryBlue}}m{4.5cm}|Z|}
\hline
\rowcolor{primaryBlue} \thd{\sffamily Nhóm thuộc tính} & \thd{\sffamily Danh sách thuộc tính được hỗ trợ} \\
\hline
Được hỗ trợ \textcolor{green!60!black}{\textbf{[Có]}} & \codeinline{color} (Màu chữ), \codeinline{background} / \codeinline{background-color} (Màu nền), \codeinline{text-shadow} (Đổ bóng chữ), \codeinline{background-clip} (Kiểm soát vùng nền). \\
\hline
Không hỗ trợ \textcolor{red}{\textbf{[Không]}} & \codeinline{border}, \codeinline{padding}, \codeinline{margin}, \codeinline{box-shadow}, \codeinline{font-size}... \\
\hline
\end{tabularx}
\end{prettyTableBox}

\begin{conceptBox}[3 Ứng Dụng Thực Tế Phổ Biến Của ::selection]
\begin{enumerate}[leftmargin=1.5em, itemsep=3pt, parsep=0pt]
    \item \textbf{Tạo trải nghiệm người dùng tốt hơn (UX):} Giúp nội dung bôi đen dễ nhìn, nổi bật và dịu mắt.
    \item \textbf{Đồng bộ nhận diện thương hiệu:} Dùng màu sắc chủ đạo của thương hiệu khi người dùng bôi đen văn bản.
    \item \textbf{Mẹo chống sao chép nội dung (tạm thời):} Đổi màu chữ thành trắng trên màu nền trắng khi bôi đen để gây khó khăn nhẹ cho việc xem nội dung bị copy.
\end{enumerate}
\end{conceptBox}

%% -------------------------------------------------------------------
%% 8.6. ::MARKER
%% -------------------------------------------------------------------
\subsubsection{8.6. Bộ Chọn ::marker (Tùy Chỉnh Ký Hiệu Đầu Dòng Danh Sách <li>)}

\begin{conceptBox}[Định Nghĩa \& Đối Tượng Nhắm Tới Của ::marker]
\codeinline{::marker} chọn phần ký hiệu đánh dấu của danh sách, bao gồm:
\begin{itemize}[leftmargin=1.5em, itemsep=2pt, parsep=0pt]
    \item Dấu chấm tròn ($\bullet$), ô vuông trong danh sách không thứ tự (\codeinline{<ul>}).
    \item Số thứ tự (1, 2, 3...) trong danh sách có thứ tự (\codeinline{<ol>}).
    \item Ký hiệu tùy chỉnh tạo bởi \codeinline{list-style-type} hoặc thuộc tính \codeinline{content}.
\end{itemize}
\end{conceptBox}

\begin{prettyTableBox}[Phân Tích 3 Trường Hợp Cú Pháp Bộ Chọn ::marker]
\rowcolors{2}{white}{primaryBlue!4}
\sffamily\small
\begin{tabularx}{\linewidth}{|>{\raggedright\arraybackslash \ttfamily\color{accentOrange}\bfseries}m{3.2cm}|Z|}
\hline
\rowcolor{primaryBlue} \thd{\sffamily Cú pháp bộ chọn} & \thd{\sffamily Cách đọc \& Ý nghĩa phạm vi chọn} \\
\hline
li::marker & **Trường hợp 1:** Chọn phần marker của tất cả thẻ \codeinline{<li>} trong mọi danh sách trên trang web. \\
\hline
ul li::marker & **Trường hợp 2:** Chỉ chọn phần marker (dấu chấm $\bullet$) của thẻ \codeinline{<li>} nằm bên trong danh sách không có thứ tự (\codeinline{<ul>}). \\
\hline
ol li::marker & **Trường hợp 3:** Chỉ chọn phần marker (số thứ tự 1, 2, 3) của thẻ \codeinline{<li>} nằm bên trong danh sách có thứ tự (\codeinline{<ol>}). \\
\hline
\end{tabularx}
\end{prettyTableBox}

\paragraph{3 Ví dụ minh họa thực tế tùy chỉnh ::marker:}

\begin{codeWindow}[Ví Dụ 1: Tùy Chỉnh Màu Sắc Và Kích Thước Marker Mặc Định]
\begin{lstlisting}[language=HTML5]
<ul>
  <li>Muc 1</li>
  <li>Muc 2</li>
</ul>

<style>
  li::marker {
    color: blue;
    font-size: 24px;
  }
</style>
\end{lstlisting}
\end{codeWindow}

\begin{codeWindow}[Ví Dụ 2: Thay Thế Marker Bằng Icon Hoặc Biểu Tượng Mũi Tên]
\begin{lstlisting}[language=HTML5]
<ul class="icon-list">
  <li>Hoc HTML5</li>
  <li>Hoc CSS3</li>
</ul>

<style>
  .icon-list li::marker {
    content: "-> ";
    color: orange;
    font-size: 18px;
  }
</style>
\end{lstlisting}
\end{codeWindow}

\begin{codeWindow}[Ví Dụ 3: Kết Hợp Danh Sách Có Thứ Tự Với CSS Counter (Bước 1, Bước 2...)]
\begin{lstlisting}[language=HTML5]
<ol class="step-list">
  <li>Phan dau tien</li>
  <li>Phan thu hai</li>
  <li>Phan thu ba</li>
</ol>

<style>
  .step-list li::marker {
    content: "Buoc " counter(list-item) ": ";
    color: purple;
    font-weight: bold;
  }
</style>
\end{lstlisting}
\end{codeWindow}

\begin{browserWindow}[Kết Quả Hiển Thị Ví Dụ 3: Bước 1, Bước 2,...]
\sffamily\small
\textcolor{purple}{\textbf{Bước 1:}} Phần đầu tiên \par
\textcolor{purple}{\textbf{Bước 2:}} Phần thứ hai \par
\textcolor{purple}{\textbf{Bước 3:}} Phần thứ ba
\end{browserWindow}

\begin{prettyTableBox}[Phân Loại Thuộc Tính Hỗ Trợ \& Không Hỗ Trợ Cho ::marker]
\rowcolors{2}{white}{primaryBlue!4}
\sffamily\small
\begin{tabularx}{\linewidth}{|>{\raggedright\arraybackslash \bfseries\color{primaryBlue}}m{4.5cm}|Z|}
\hline
\rowcolor{primaryBlue} \thd{\sffamily Loại thuộc tính} & \thd{\sffamily Danh sách chi tiết} \\
\hline
Được hỗ trợ \textcolor{green!60!black}{\textbf{[Có]}} & \codeinline{color}, \codeinline{font}, \codeinline{font-size}, \codeinline{font-weight}, \codeinline{font-family}, \codeinline{letter-spacing}, \codeinline{text-transform}, \codeinline{content} (thay thế marker). \\
\hline
Không hỗ trợ \textcolor{red}{\textbf{[Không]}} & \codeinline{margin}, \codeinline{padding}, \codeinline{background}, \codeinline{border}, \codeinline{width}, \codeinline{height}... \\
\hline
\end{tabularx}
\end{prettyTableBox}

\begin{warnBox}[3 Lưu Ý Vàng Khi Sử Dụng ::marker]
\begin{enumerate}[leftmargin=1.5em, itemsep=3pt, parsep=0pt]
    \item \textbf{Phải là phần tử danh sách:} \codeinline{::marker} chỉ hoạt động với các phần tử có thuộc tính \codeinline{display: list-item}.
    \item \textbf{Ảnh hưởng toàn bộ danh sách:} Quy tắc \codeinline{li::marker} sẽ áp dụng cho tất cả các mục. Nếu muốn marker riêng cho từng mục, bạn cần thêm class riêng cho từng thẻ \codeinline{<li>}.
    \item \textbf{Hỗ trợ trình duyệt hiện đại:} Tất cả các trình duyệt hiện đại (Chrome, Firefox, Edge, Safari) đều đã hỗ trợ hoàn hảo \codeinline{::marker}.
\end{enumerate}
\end{warnBox}

\begin{conceptBox}[Ứng Dụng Thực Tế Của ::marker]
\begin{itemize}[leftmargin=1.5em, itemsep=2pt, parsep=0pt]
    \item \textbf{Tạo danh sách đẹp mắt:} Dùng icon, biểu tượng mũi tên hoặc ký hiệu đặc biệt làm dấu đầu dòng.
    \item \textbf{Hướng dẫn từng bước:} Thay số thứ tự thành "Bước 1:", "Bước 2:", v.v.
    \item \textbf{Điểm nhấn thiết kế:} Làm nổi bật danh sách sản phẩm/tính năng mà không cần thêm thẻ HTML phức tạp.
\end{itemize}
\end{conceptBox}

%% -------------------------------------------------------------------
%% 8.7. BẢNG MASTER SUMMARY VÀ TỔNG KẾT
%% -------------------------------------------------------------------
\subsubsection{8.7. Bảng Tổng Hợp Cheatsheet \& Nguyên Tắc Thiết Kế Thực Tế}

\begin{prettyTableBox}[Bảng Cheatsheet Tra Cứu Toàn Diện 6 Pseudo-Element Selectors]
\rowcolors{2}{white}{primaryBlue!4}
\sffamily\small
\begin{tabularx}{\linewidth}{|>{\raggedright\arraybackslash \ttfamily\color{accentOrange}\bfseries}m{2.5cm}|>{\bfseries\color{primaryBlue}}m{3.8cm}|>{\ttfamily\small}m{3.6cm}|Z|}
\hline
\rowcolor{primaryBlue} \thd{\sffamily Pseudo-element} & \thd{\sffamily Chức năng chính} & \thd{\sffamily Thuộc tính hỗ trợ} & \thd{\sffamily Ứng dụng thực tế tiêu biểu} \\
\hline
::before & Chèn nội dung trước phần tử & content, color, font, layout & Thêm icon, dấu hiệu, nhãn trang trí trước tiêu đề hoặc nút bấm. \\
\hline
::after & Chèn nội dung sau phần tử & content, color, font, layout & Thêm chú thích, hiệu ứng kết thúc, hoặc ký hiệu cảnh báo, clearfix. \\
\hline
::first-letter & Chọn chữ cái đầu tiên & font-size, color, float, margin & Tạo drop cap (chữ đầu lớn) nghệ thuật như trong sách, báo in. \\
\hline
::first-line & Chọn dòng đầu tiên & font, color, text-transform & Làm nổi bật dòng mở đầu bài viết blog, báo chí. \\
\hline
::selection & Chọn nội dung được bôi đen & background, color, text-shadow & Tạo hiệu ứng đặc biệt khi người dùng chọn văn bản. \\
\hline
::marker & Chọn ký hiệu danh sách & color, content, font-size & Tùy chỉnh bullet, số thứ tự trong danh sách \codeinline{<li>}. \\
\hline
\end{tabularx}
\end{prettyTableBox}


\begin{conceptBox}[Tổng Kết \& Ứng Dụng Thực Tế Nâng Cao]
\begin{itemize}[leftmargin=1.5em, itemsep=4pt, parsep=0pt]
    \item \textbf{Tối ưu HTML:} Giảm số lượng thẻ HTML không cần thiết, giữ mã nguồn sạch sẽ và chuẩn W3C semantic.
    \item \textbf{Nâng cao UX/UI:} Tạo hiệu ứng tương tác thị giác đẹp mắt, tăng trải nghiệm người dùng khi đọc và tương tác.
    \item \textbf{Dễ bảo trì mã nguồn:} Chỉ cần thay đổi CSS trong một file duy nhất mà không cần sửa cấu trúc mã HTML hàng loạt.
\end{itemize}
\par\vspace{6pt}
\textbf{Ví dụ thực tiễn:} Khi làm menu điều hướng hoặc card sản phẩm, bạn có thể dùng \codeinline{::before} để thêm biểu tượng icon nhỏ, hoặc \codeinline{::selection} để tạo trải nghiệm đọc thương hiệu thú vị hơn.
\end{conceptBox}

%% -------------------------------------------------------------------
%% 8.8. ĐỌC THÊM: NÂNG CAO KIẾN THỨC VÀ KINH NGHIỆM THỰC CHIẾN
%% -------------------------------------------------------------------
\subsubsection{8.8. Đọc Thêm: Chuyên Đề Nâng Cao Về Pseudo-Elements (Mở Rộng Kiến Thức Thực Chiến)}

\begin{conceptBox}[1. Các Giá Trị Đa Dạng Của Thuộc Tính content Trong ::before / ::after]
Thuộc tính \codeinline{content} trong \codeinline{::before} và \codeinline{::after} không chỉ nhận chuỗi văn bản đơn thuần, mà còn hỗ trợ nhiều kiểu giá trị cực kỳ mạnh mẽ:
\begin{itemize}[leftmargin=1.5em, itemsep=3pt, parsep=0pt]
    \item \codeinline{content: "Chuoi van ban";}: Chèn văn bản trực tiếp.
    \item \codeinline{content: "";}: Chuỗi rỗng (Thường kết hợp với \codeinline{display: block/inline-block}, \codeinline{width}, \codeinline{height}, \codeinline{background} để vẽ các hình khối trang trí, đường kẻ, hoặc overlay).
    \item \codeinline{content: attr(attribute-name);}: Lấy giá trị thuộc tính HTML động. Rất phổ biến để tạo Tooltip hoặc hiển thị đường dẫn URL khi in trang web.
    \item \codeinline{content: url(image.png);}: Chèn hình ảnh trực tiếp.
    \item \codeinline{content: counter(name);}: Tạo bộ đếm số tự động trong CSS (CSS Counters).
\end{itemize}
\end{conceptBox}

\begin{codeWindow}[Ví Dụ Ứng Dụng content: attr(data-tooltip) Để Tạo Tooltip Thuần CSS]
\begin{lstlisting}[language=HTML5]
<button class="btn" data-tooltip="Nhan vao day de luu du lieu">Luu bai viet</button>

<style>
  .btn { position: relative; }
  .btn:hover::after {
    content: attr(data-tooltip);
    position: absolute;
    bottom: 100%;
    left: 50%;
    transform: translateX(-50%);
    background: #333;
    color: #fff;
    padding: 4px 8px;
    border-radius: 4px;
    white-space: nowrap;
    font-size: 12px;
  }
</style>
\end{lstlisting}
\end{codeWindow}

\begin{browserWindow}[Kết Quả: Tooltip Động Thuần CSS Với content: attr()]
\begin{tcolorbox}[arc=4pt, colback=white, colframe=primaryBlue!30, boxrule=0.8pt, top=10pt, bottom=10pt, left=12pt, right=12pt]
  \begin{center}
    \begin{tcolorbox}[colback=gray!90!black, colframe=gray!90!black, arc=3pt, boxrule=0pt, top=3pt, bottom=3pt, left=6pt, right=6pt, width=5.5cm, halign=center]
      {\sffamily\small\color{white} Nhấn vào đây để lưu dữ liệu}
    \end{tcolorbox}\par\vspace{-2pt}
    \begin{tcolorbox}[colback=primaryBlue, colframe=primaryBlue, arc=4pt, boxrule=0.8pt, top=6pt, bottom=6pt, left=12pt, right=12pt, width=4.5cm, halign=center]
      {\sffamily\bfseries\color{white} Lưu bài viết}
    \end{tcolorbox}
  \end{center}\par\vspace{2pt}
  {\sffamily\small\color{gray!70} → Khi di chuột vào nút, \codeinline{::after} lấy thuộc tính \codeinline{data-tooltip} hiển thị làm khung chú thích phía trên.}
\end{tcolorbox}
\end{browserWindow}

\begin{prettyTableBox}[2. Bảng So Sánh Chi Tiết: Pseudo-class (:) vs Pseudo-element (::)]
\rowcolors{2}{white}{primaryBlue!4}
\sffamily\small
\begin{tabularx}{\linewidth}{|>{\raggedright\arraybackslash \bfseries\color{primaryBlue}}m{3.5cm}|Z|Z|}
\hline
\rowcolor{primaryBlue} \thd{\sffamily Tiêu chí} & \thd{\sffamily Pseudo-class (:)} & \thd{\sffamily Pseudo-element (::)} \\
\hline
Cú pháp tiêu chuẩn CSS3 & Sử dụng 1 dấu hai chấm (\codeinline{:hover}) & Sử dụng 2 dấu hai chấm (\codeinline{::before}) \\
\hline
Bản chất tác động & Chọn phần tử dựa trên \textbf{trạng thái tương tác} hoặc \textbf{vị trí vị thế DOM}. & Định kiểu cho một \textbf{phần nội dung cụ thể} hoặc tạo \textbf{phần tử ảo}. \\
\hline
Sự tồn tại trong DOM & Nhắm trực tiếp vào thẻ HTML thực tế đã có trong cây DOM. & Tạo phần tử ảo hiển thị trên màn hình nhưng KHÔNG tồn tại trong cây DOM. \\
\hline
Ví dụ tiêu biểu & \codeinline{:hover}, \codeinline{:focus}, \codeinline{:nth-child(2)} & \codeinline{::before}, \codeinline{::after}, \codeinline{::first-letter} \\
\hline
\end{tabularx}
\end{prettyTableBox}

\begin{warnBox}[3. Lưu Ý Quan Trọng Về Accessibility (Khả Năng Truy Cập - a11y) \& Hiệu Năng]
\begin{itemize}[leftmargin=1.5em, itemsep=4pt, parsep=0pt]
    \item \textbf{Công cụ đọc màn hình (Screen Readers):} Một số công cụ đọc màn hình có thể đọc nội dung trong thuộc tính \codeinline{content:"..."} của \codeinline{::before} và \codeinline{::after}. Tránh chèn các ký tự trang trí không có nghĩa ngữ văn bản (ví dụ: \codeinline{content: ">>> ";}) vì có thể gây phiền cho người khiếm thị. Nếu nội dung chỉ mang tính trang trí, hãy cân nhắc áp dụng thuộc tính \codeinline{aria-hidden="true"}.
    \item \textbf{Modern CSS Syntax for Alt Text:} Trong CSS Modern, bạn có thể truyền văn bản thay thế alt trực tiếp trong \codeinline{content}: \codeinline{content: "STAR" / "Bieu tuong ngoi sao";}.
    \item \textbf{Hiệu năng Rendering:} Các pseudo-elements được xử lý rất nhanh bởi GPU của trình duyệt. Tuy nhiên khi kết hợp animation biến đổi (\codeinline{transform}, \codeinline{opacity}) trên \codeinline{::before}/\codeinline{::after}, hãy luôn ưu tiên sử dụng thuộc tính \codeinline{will-change} để tránh hiện tượng giật khung hình (\term{Jank}).
\end{itemize}
\end{warnBox}

\begin{tipBox}[4. Bộ Câu Hỏi Phỏng Vấn Tuyển Dụng Frontend Chuyên Sâu (Interview Q\&A)]
\begin{enumerate}[leftmargin=1.5em, itemsep=4pt, parsep=0pt]
    \item \textbf{Q1: Phân biệt sự khác biệt cốt lõi giữa 1 dấu hai chấm (:) và 2 dấu hai chấm (::) trong CSS?} \\
    $\rightarrow$ \textbf{Trả lời ghi điểm}: Dấu single colon (\codeinline{:}) dùng cho Pseudo-classes đại diện cho trạng thái tương tác của phần tử có sẵn trong DOM (như \codeinline{:hover}, \codeinline{:focus}, \codeinline{:nth-child}). Dấu double colon (\codeinline{::}) dùng cho Pseudo-elements đại diện cho phần tử ảo không có trong DOM hoặc một phần nội dung cụ thể (như \codeinline{::before}, \codeinline{::after}, \codeinline{::first-letter}).
    \item \textbf{Q2: Tại sao không thể sử dụng ::before hay ::after trên các thẻ <img>, <input>, <br>?} \\
    $\rightarrow$ \textbf{Trả lời ghi điểm}: Vì các thẻ như \codeinline{<img>}, \codeinline{<input>}, \codeinline{<br>} là các phần tử thay thế (\term{Replaced Elements} / \term{Void Elements}). Chúng không chứa phần nội dung con bên trong (no child content), do đó trình duyệt không có vị trí con đầu tiên (first child) hay con cuối cùng (last child) để chèn phần tử giả \codeinline{::before} hay \codeinline{::after}.
    \item \textbf{Q3: Điều gì xảy ra nếu bạn quên thuộc tính content khi khai báo ::before?} \\
    $\rightarrow$ \textbf{Trả lời ghi điểm}: Thuộc tính \codeinline{content} là điều kiện tiên quyết bắt buộc. Nếu không có \codeinline{content} (hoặc nếu set \codeinline{content: none}), trình duyệt sẽ bỏ qua và không khởi tạo phần tử giả đó trên cây Render Tree.
    \item \textbf{Q4: Kỹ thuật Micro Clearfix giải quyết vấn đề gì bằng ::after?} \\
    $\rightarrow$ \textbf{Trả lời ghi điểm}: Giải quyết hiện tượng sụp đổ chiều cao (\term{Height Collapse}) của thẻ cha khi các thẻ con dùng \codeinline{float}. Bằng cách chèn \codeinline{::after} có \codeinline{content: ""; display: table; clear: both;}, thẻ cha tự động bao bọc trọn vẹn toàn bộ chiều cao của các thẻ con float mà không cần chèn thẻ HTML rỗng.
    \item \textbf{Q5: Làm thế nào để đặt một phần tử giả ::before nằm phía sau background của thẻ cha?} \\
    $\rightarrow$ \textbf{Trả lời ghi điểm}: Thiết lập \codeinline{position: relative} và \codeinline{z-index: 1} (hoặc \codeinline{isolation: isolate}) trên thẻ cha để tạo một Stacking Context độc lập, sau đó thiết lập \codeinline{position: absolute} và \codeinline{z-index: -1} trên \codeinline{::before}.
\end{enumerate}
\end{tipBox}

%% -------------------------------------------------------------------
%% 8.9. XỬ LÝ XUNG ĐỘT KHI SỬ DỤNG PSEUDO-ELEMENT SELECTORS
%% -------------------------------------------------------------------
\subsubsection{8.9. Phân Tích \& Xử Lý Xung Đột Khi Sử Dụng Pseudo-elements (CSS Conflicts Masterclass)}

Khi phát triển giao diện Web phức tạp, việc khai báo nhiều quy tắc CSS cho cùng một bộ chọn phần tử giả có thể gây ra xung đột thị giác. Dưới đây là phân tích chi tiết 4 dạng xung đột phổ biến nhất cùng giải pháp xử lý triệt để:

\begin{conceptBox}[1. Xung Đột Khi Khai Báo Nhiều Lần Cho Cùng Một Pseudo-element]
\textbf{Nguyên lý duy nhất (Single Instance Rule):} Mỗi phần tử HTML chỉ sở hữu duy nhất **một instance (thể hiện ảo)** của từng loại phần tử giả (\codeinline{::before} hoặc \codeinline{::after}).
\par\vspace{4pt}
\textbf{Cơ chế xung đột (Cascade Override):} Nếu bạn khai báo nhiều khối CSS cùng nhắm vào \codeinline{.class::before}, trình duyệt sẽ áp dụng quy tắc Cascading mặc định: **Quy tắc xuất hiện sau cùng trong mã nguồn sẽ ghi đè hoàn toàn quy tắc trước đó**!
\end{conceptBox}

\begin{codeWindow}[Ví Dụ Xung Đột Khai Báo Nhiều Khối ::before]
\begin{lstlisting}[language=HTML5]
<p class="text">Chao ban!</p>

<style>
  /* Quy tac 1: Bi ghi de */
  .text::before {
    content: "Xin chao! ";
    color: red;
  }
  /* Quy tac 2: Chien thang (viet sau) */
  .text::before {
    content: "Hello! ";
    color: blue;
  }
</style>
\end{lstlisting}
\end{codeWindow}

\begin{browserWindow}[Kết Quả Hiển Thị: Quy Tắc Viết Sau Ghi Đè Quy Tắc Trước]
\sffamily\small
\textbf{Kết quả:} \textcolor{blue}{\textbf{Hello! }}Chào bạn! \par\vspace{2pt}
{\scriptsize\textcolor{gray!70}{\textit{Giải thích: Quy tắc thứ hai (.text::before có content: "Hello! ") được viết sau nên ghi đè quy tắc thứ nhất.}}}
\end{browserWindow}

\begin{tipBox}[Giải Pháp Xử Lý Xung Đột Khai Báo Trùng]
\textbf{Giải pháp:} Gộp tất cả thuộc tính định kiểu vào một khối CSS duy nhất hoặc sử dụng các selector có độ đặc hiệu (Specificity) cao hơn để phân biệt ngữ cảnh rõ ràng:
\begin{codeWindow}[Mã Nguồn Giải Pháp Gộp Style Chuẩn]
\begin{lstlisting}[language=CSS3]
.text::before {
  content: "Xin chao! ";
  color: red;
  font-weight: bold;
}
\end{lstlisting}
\end{codeWindow}
\end{tipBox}

\begin{conceptBox}[2. Xung Đột Khi Thẻ Cha Mang Thuộc Tính display: none]
\textbf{Hiện tượng:} \codeinline{::before} và \codeinline{::after} tạo ra nội dung ảo nằm **bên trong** phần tử cha trên Render Tree.
\par\vspace{4pt}
\textbf{Hệ quả:} Nếu thẻ cha bị ẩn bằng thuộc tính \codeinline{display: none}, toàn bộ phần tử giả con bên trong cũng sẽ **biến mất 100\%** khỏi màn hình hiển thị, bất kể bạn có cố gắng đặt \codeinline{display: block} hay \codeinline{visibility: visible} riêng cho phần tử giả!
\end{conceptBox}

\begin{codeWindow}[Ví Dụ Thẻ Cha display: none Làm Biến Mất Pseudo-element]
\begin{lstlisting}[language=HTML5]
<p class="hidden-box">Noi dung bai viet</p>

<style>
  p.hidden-box {
    display: none; /* An toan bo the cha */
  }
  p.hidden-box::before {
    content: "Khong thay toi dau!";
    color: red;
  }
</style>
\end{lstlisting}
\end{codeWindow}

\begin{browserWindow}[Kết Quả Hiển Thị: Không Có Gì Hiển Thị]
\sffamily\small
\textit{(Không hiển thị bất kỳ thành phần nào trên màn hình vì thẻ cha p bị ẩn hoàn toàn bởi display: none)}
\end{browserWindow}

\begin{tipBox}[Giải Pháp Khi Cần Ẩn Nội Dung Thật Nhưng Giữ Phần Tử Giả]
\textbf{Giải pháp:} Nếu muốn ẩn nội dung thật của thẻ cha mà vẫn duy trì phần tử giả hiển thị, bạn có thể thiết lập \codeinline{font-size: 0} trên thẻ cha và đặt lại \codeinline{font-size} chuẩn trên phần tử giả, hoặc dùng thuộc tính \codeinline{visibility: hidden} kết hợp \codeinline{visibility: visible} trên phần tử giả!
\end{tipBox}

\begin{conceptBox}[3. Xung Đột Tương Tác Giữa ::first-letter / ::first-line Với ::before / ::after]
\textbf{Nguyên lý tương tác:} Bộ chọn \codeinline{::first-letter} và \codeinline{::first-line} được thiết kế để định dạng **văn bản thực tế** trong thẻ HTML. Chúng **KHÔNG tự động tác động** lên ký tự hoặc nội dung ảo được chèn từ thuộc tính \codeinline{content} của \codeinline{::before}!
\end{conceptBox}

\begin{codeWindow}[Ví Dụ ::first-letter Không Ảnh Hưởng Đến Nội Dung Ảo Từ ::before]
\begin{lstlisting}[language=HTML5]
<p class="intro">Chao mung ban!</p>

<style>
  p.intro::before {
    content: "STAR "; /* Noi dung ao */
  }
  p.intro::first-letter {
    color: red;
    font-size: 30px;
  }
</style>
\end{lstlisting}
\end{codeWindow}

\begin{browserWindow}[Kết Quả Hiển Thị: ::first-letter Định Kiểu Cho Chữ C]
\sffamily\small
STAR \textbf{\Huge\textcolor{red}{C}}hào mừng bạn! \par\vspace{2pt}
{\scriptsize\textcolor{gray!70}{\textit{Giải thích: ::first-letter tác động vào chữ cái đầu "C" của nội dung thật, không phải chữ "S" trong nội dung giả "STAR".}}}
\end{browserWindow}

\begin{tipBox}[Giải Pháp Định Kiểu Riêng Cho Nội Dung Giả]
\textbf{Giải pháp:} Nếu muốn định dạng kiểu chữ, màu sắc hay kích thước cho nội dung giả từ \codeinline{::before}, hãy trực tiếp khai báo các thuộc tính đó bên trong khối \codeinline{::before}:
\begin{codeWindow}[Khai Báo Định Kiểu Trực Tiếp Trong ::before]
\begin{lstlisting}[language=CSS3]
p.intro::before {
  content: "STAR ";
  color: green;
  font-size: 30px;
  font-weight: bold;
}
\end{lstlisting}
\end{codeWindow}
\end{tipBox}

\begin{conceptBox}[4. Xung Đột Thứ Tự CSS Về Độ Ưu Tiên (Specificity Conflict)]
Khi hai bộ chọn cùng nhắm tới phần tử giả của một thẻ (ví dụ: \codeinline{h1::before} và \codeinline{h1.special::before}), trình duyệt sẽ tính toán điểm **CSS Specificity** để quyết định quy tắc nào chiến thắng.
\end{conceptBox}

\begin{codeWindow}[Ví Dụ Bộ Chọn Cụ Thể Hơn Chiến Thắng]
\begin{lstlisting}[language=HTML5]
<h1 class="special">Tieu de trang web</h1>

<style>
  /* Specificity: (0,0,0,2) - Class va Element */
  h1.special::before {
    content: "[SPECIAL] ";
    color: blue;
  }

  /* Specificity: (0,0,0,1) - Element */
  h1::before {
    content: "[DEFAULT] ";
    color: orange;
  }
</style>
\end{lstlisting}
\end{codeWindow}

\begin{browserWindow}[Kết Quả Hiển Thị: Specificity Cao Hơn Chiến Thắng]
\sffamily\small
\textbf{\Large\textcolor{blue}{[SPECIAL] }} Tiêu đề trang web \par\vspace{2pt}
{\scriptsize\textcolor{gray!70}{\textit{Giải thích: Thẻ h1 có class "special" nên bộ chọn h1.special::before (điểm cao hơn) chiến thắng.}}}
\end{browserWindow}

\begin{prettyTableBox}[Bảng Tra Cứu Toàn Diện Các Dạng Xung Đột, Nguyên Nhân \& Giải Pháp Xử Lý]
\rowcolors{2}{white}{primaryBlue!4}
\sffamily\small
\begin{tabularx}{\linewidth}{|>{\raggedright\arraybackslash \ttfamily\color{accentOrange}\bfseries}m{2.5cm}|>{\bfseries\color{primaryBlue}}m{3.5cm}|Z|Z|}
\hline
\rowcolor{primaryBlue} \thd{\sffamily Pseudo-element} & \thd{\sffamily Bản chất đối tượng} & \thd{\sffamily Nguyên nhân xung đột thường gặp} & \thd{\sffamily Giải pháp xử lý chuẩn mực} \\
\hline
::before / ::after & Nội dung ảo chèn vào đầu/cuối phần tử & Trùng lặp khai báo nhiều khối CSS; bị ẩn khi thẻ cha dùng \codeinline{display: none}. & Dùng selector cụ thể hơn để phân tách; điều chỉnh logic hiển thị bằng \codeinline{visibility}. \\
\hline
::first-letter & Chữ cái đầu tiên của văn bản thật & Không tự động tác động lên nội dung ảo của \codeinline{::before}. & Khai báo định kiểu font/color trực tiếp trong khối \codeinline{::before}. \\
\hline
::first-line & Dòng đầu tiên của văn bản khối & Không áp dụng được các thuộc tính layout như \codeinline{margin}, \codeinline{padding}, \codeinline{width}. & Chỉ dùng các thuộc tính văn bản/màu sắc hợp lệ (\codeinline{color}, \codeinline{font}, \codeinline{line-height}). \\
\hline
Thứ tự CSS (Cascade) & Quy tắc ghi nguồn từ trên xuống dưới & Khai báo sau ghi đè khai báo trước nếu cùng Specificity. & Kiểm tra thứ tự nạp file CSS, gộp thuộc tính trùng lặp. \\
\hline
Độ ưu tiên (Specificity) & Điểm số bộ chọn trình duyệt tính toán & Class mạnh hơn Element, ID mạnh hơn Class. & Nắm vững công thức Specificity; hạn chế tối đa việc lạm dụng \codeinline{!important}. \\
\hline
\end{tabularx}
\end{prettyTableBox}

\begin{tipBox}[3 Lời Khuyên Thực Chiến Tối Ưu Quản Lý Pseudo-elements]
\begin{enumerate}[leftmargin=1.5em, itemsep=3pt, parsep=0pt]
    \item \textbf{Viết CSS rõ ràng, mô-đun hóa:} Đừng chồng chéo quá nhiều phần tử giả không cần thiết trên cùng một thẻ HTML.
    \item \textbf{Kiểm tra điểm Specificity:} Khi thấy phần tử giả không hiển thị đúng màu hoặc content, hãy kiểm tra xem có selector nào mạnh hơn đang ghi đè hay không.
    \item \textbf{Thành thạo công cụ Chrome DevTools:} Mở tab *Elements* trên trình duyệt, click mở rộng thẻ HTML target để soi trực tiếp phần tử ảo \codeinline{::before} / \codeinline{::after} và kiểm tra các dòng CSS đang bị gạch ngang (ghi đè).
\end{enumerate}
\end{tipBox}




% ==============================================================================
% MỤC 9: THỨ TỰ ƯU TIÊN TRONG CSS (CSS SPECIFICITY MASTERCLASS)
% ==============================================================================
\section{Thứ Tự Ưu Tiên Trong CSS (CSS Specificity Masterclass)}


Trong thiết kế giao diện Web thực tế, một phần tử HTML thường nằm trong phạm vi ảnh hưởng của nhiều quy tắc CSS khác nhau (ví dụ: vừa có style từ thẻ \codeinline{p}, vừa có class \codeinline{.text}, vừa có ID \codeinline{\#main}, lại vừa có thuộc tính \codeinline{style="..."} trực tiếp). Khi xảy ra xung đột giữa các quy tắc này, trình duyệt sẽ dùng cơ chế \term{CSS Specificity} (Thứ tự ưu tiên / Tính đặc hiệu) để tính toán điểm số và quyết định quy tắc nào sẽ chiến thắng.

\begin{conceptBox}[Khái Niệm Specificity (Tính Đặc Hiệu) Là Gì?]
\textbf{CSS Specificity (Tính đặc hiệu)} là thuật toán và hệ thống quy tắc mà trình duyệt sử dụng để quyết định quy tắc CSS nào sẽ được áp dụng cho một phần tử HTML khi có nhiều quy tắc cùng nhắm vào phần tử đó.
\par\vspace{4pt}
\textbf{Vì sao cần Specificity?} Khi trang web phát triển lớn, một phần tử (ví dụ: nút bấm \codeinline{<button id="submit" class="btn">}) thường khớp với nhiều quy tắc CSS khác nhau. Specificity giúp trình duyệt tính toán điểm số công bằng để chọn ra quy tắc mạnh nhất dựa trên độ cụ thể của bộ chọn, \textbf{không phụ thuộc vào thứ tự viết code} (trừ khi hai bộ chọn có điểm số bằng nhau).
\end{conceptBox}

%% ===================================================================
%% PHẦN 8.1: BẢNG XẾP HẠNG ĐỘ ƯU TIÊN SPECIFICITY
%% ===================================================================
\subsection{1. Bảng Xếp Hạng Chi Tiết Độ Ưu Tiên Specificity (6 Bậc Giá Trị)}

\begin{prettyTableBox}[Bảng Xếp Hạng Chi Tiết Độ Ưu Tiên Trong CSS]
\rowcolors{2}{white}{primaryBlue!4}
\sffamily\small
\begin{tabularx}{\linewidth}{|c|>{\raggedright\arraybackslash \bfseries\color{primaryBlue}}m{4.2cm}|>{\ttfamily\color{darkText}}m{4.5cm}|c|Z|}
\hline
\rowcolor{primaryBlue} \thd{\sffamily Cấp độ} & \thd{\sffamily Loại bộ chọn (Selector)} & \thd{\sffamily Cú pháp minh họa} & \thd{\sffamily Specificity} & \thd{\sffamily Đặc tính \& Độ mạnh} \\
\hline
(1) & !important trong CSS & \codeinline{color: red !important;} & — & \textcolor{red}{\textbf{[Cao nhất]}} -- Ghi đè tất cả (trừ inline style có !important) \\
\hline
(2) & Style nội tuyến (Inline style) & \codeinline{<div style="color:red;">} & (1,0,0,0) & \textcolor{orange!90!black}{\textbf{Nội tuyến}} -- Rất mạnh, ghi đè mọi ID và Class \\
\hline
(3) & Bộ chọn ID (\#id) & \codeinline{\#header \{ color: blue; \}} & (0,1,0,0) & \textcolor{primaryBlue}{\textbf{ID Selector}} -- Rất mạnh, nhưng hạn chế lạm dụng \\
\hline
(4) & Class, Attribute, Pseudo-class & \codeinline{.btn}, \codeinline{[type="text"]}, \codeinline{:hover} & (0,0,1,0) & \textcolor{green!60!black}{\textbf{Chung nhóm}} -- Linh hoạt và được sử dụng phổ biến nhất \\
\hline
(5) & Thẻ HTML \& Pseudo-element & \codeinline{div}, \codeinline{h1}, \codeinline{p}, \codeinline{::before} & (0,0,0,1) & \textcolor{gray}{\textbf{Thẻ Element}} -- Mức độ yếu nhất trong các bộ chọn \\
\hline
(6) & Bộ chọn chung \& Kế thừa & \codeinline{* \{ margin: 0; \}}, \codeinline{body \{ color \}} & — & \textcolor{gray!70}{\textbf{Mặc định}} -- Yếu nhất, dễ dàng bị bất kỳ bộ chọn nào ghi đè \\
\hline
\end{tabularx}
\end{prettyTableBox}

%% ===================================================================
%% PHẦN 8.2: QUY TRÌNH 4 BƯỚC DUYỆT DOM TÍNH SPECIFICITY
%% ===================================================================
\subsection{2. Quy Trình 4 Bước Duyệt DOM \& Công Thức (A, B, C, D) / (a, b, c, d)}

Trình duyệt thực hiện quá trình tính toán và áp dụng kiểu dáng theo 4 bước tiêu chuẩn:

\begin{conceptBox}[4 Bước Xử Lý Thứ Tự Ưu Tiên Của Trình Duyệt]
\textbf{Bước 1: Thu thập tất cả các quy tắc CSS áp dụng cho phần tử.} \par
Ví dụ: Với thẻ \codeinline{<p id="main" class="text">}, trình duyệt gom tất cả các quy tắc khớp với phần tử này như \codeinline{p}, \codeinline{.text}, \codeinline{\#main}, \codeinline{p.text}, \codeinline{\#main.text}...
\par\vspace{4pt}
\textbf{Bước 2: Tính điểm specificity cho từng bộ chọn theo dạng 4 nhóm (A, B, C, D) hoặc (a, b, c, d).}
\begin{itemize}[leftmargin=1.5em, itemsep=2pt, parsep=0pt]
    \item \textbf{Nhóm A / a (1,0,0,0):} Style nội tuyến (Inline styles) đặt trực tiếp trong thuộc tính \codeinline{style="..."}.
    \item \textbf{Nhóm B / b (0,1,0,0):} Số lượng ID selectors (\codeinline{\#id}).
    \item \textbf{Nhóm C / c (0,0,1,0):} Số lượng Class (\codeinline{.text}), Attribute (\codeinline{[type="text"]}), Pseudo-class (\codeinline{:hover}, \codeinline{:focus}, \codeinline{:not()}).
    \item \textbf{Nhóm D / d (0,0,0,1):} Số lượng Element (thẻ HTML như \codeinline{div}, \codeinline{p}) và Pseudo-element (\codeinline{::before}, \codeinline{::after}).
\end{itemize}
\par\vspace{4pt}
\textbf{Bước 3: So sánh specificity theo thứ tự ưu tiên $A > B > C > D$ (hoặc $a > b > c > d$).} \par
Trình duyệt so sánh từng chỉ số từ trái sang phải (giống như so sánh số nguyên nhiều chữ số):
\begin{itemize}[leftmargin=1.5em, itemsep=2pt, parsep=0pt]
    \item \codeinline{p} $\rightarrow$ Specificity = \codeinline{(0,0,0,1)}
    \item \codeinline{.text} $\rightarrow$ Specificity = \codeinline{(0,0,1,0)}
    \item \codeinline{\#main} $\rightarrow$ Specificity = \codeinline{(0,1,0,0)}
    \item \codeinline{style="..."} $\rightarrow$ Specificity = \codeinline{(1,0,0,0)}
    \item \codeinline{p.text:hover} $\rightarrow$ Specificity = \codeinline{(0,0,2,1)}
\end{itemize}
\par\vspace{4pt}
\textbf{Bước 4: Áp dụng quy tắc có specificity cao nhất.}
\begin{itemize}[leftmargin=1.5em, itemsep=2pt, parsep=0pt]
    \item Nếu nhiều quy tắc có \textbf{cùng điểm specificity}, quy tắc nào \textbf{viết sau trong file CSS sẽ thắng} (theo thứ tự xuất hiện \term{Source Order}).
    \item Nếu có từ khóa \codeinline{!important}, trình duyệt sẽ cấp quyền ưu tiên tuyệt đối để đè qua tất cả chỉ số specificity thông thường.
\end{itemize}
\end{conceptBox}

\begin{prettyTableBox}[Bảng Định Nghĩa 4 Thành Phần Của Công Thức (A, B, C, D) / (a, b, c, d)]
\rowcolors{2}{white}{primaryBlue!4}
\sffamily\small
\begin{tabularx}{\linewidth}{|c|>{\raggedright\arraybackslash \bfseries\color{primaryBlue}}m{4.8cm}|c|Z|}
\hline
\rowcolor{primaryBlue} \thd{\sffamily Ký hiệu} & \thd{\sffamily Thành phần quy đổi} & \thd{\sffamily Giá trị mẫu} & \thd{\sffamily Quy tắc cộng điểm} \\
\hline
A / a & Style nội tuyến (Inline style) & (1,0,0,0) & \textbf{+1} nếu phần tử được viết trực tiếp trong thuộc tính \codeinline{style="..."} \\
\hline
B / b & Số lượng ID selectors (\codeinline{\#id}) & (0,1,0,0) & \textbf{+1} cho mỗi ID xuất hiện trong bộ chọn \\
\hline
C / c & Class, Attribute, Pseudo-classes & (0,0,1,0) & \textbf{+1} cho mỗi \codeinline{.class}, \codeinline{[attr]}, \codeinline{:hover}, \codeinline{:nth-child()} \\
\hline
D / d & Thẻ HTML (Elements) \& Pseudo-elements & (0,0,0,1) & \textbf{+1} cho mỗi thẻ HTML (\codeinline{div}, \codeinline{p}) và \codeinline{::before}, \codeinline{::after} \\
\hline
\end{tabularx}
\end{prettyTableBox}

\begin{noteBox}[Lưu Ý Quan Trọng Khi Tính Điểm Specificity]
\begin{itemize}[leftmargin=1.5em, itemsep=3pt, parsep=0pt]
    \item \textbf{Cách so sánh}: Trình duyệt so sánh từ trái sang phải: $A$ lớn hơn thì thắng; nếu $A$ bằng nhau thì so sánh $B$; nếu $B$ bằng nhau thì so sánh $C$; nếu $C$ bằng nhau thì so sánh $D$.
    \item \textbf{Không quy đổi thập phân}: Điểm $(0,0,1,0)$ đại diện cho 1 Class luôn luôn mạnh hơn bất kỳ số lượng thẻ HTML nào (ví dụ $(0,0,0,11)$ gồm 11 thẻ HTML vẫn thua 1 Class).
    \item \textbf{Universal Selector (\codeinline{*}) và Combinators (\codeinline{>}, \codeinline{+}, \codeinline{\~{}})}: Có điểm số là $(0,0,0,0)$, không cộng thêm điểm specificity.
    \item \textbf{File Internal (\codeinline{<style>}) vs External (\codeinline{.css})}: Vị trí khai báo CSS không làm thay đổi specificity, chỉ bị ảnh hưởng bởi specificity và thứ tự xuất hiện.
\end{itemize}
\end{noteBox}

%% ===================================================================
%% PHẦN 8.3: BẢNG SO SÁNH ĐỘ ƯU TIÊN CÁC BỘ CHỌN THƯỜNG GẶP
%% ===================================================================
\subsection{3. Bảng Tra Cứu Specificity Của Các Bộ Chọn Thường Gặp (15 Ví Dụ)}

\begin{prettyTableBox}[Bảng Tra Cứu Specificity Của Các Bộ Chọn CSS Phổ Biến (Tự Động Sắp Xếp Từ Thấp Đến Cao)]
\rowcolors{2}{white}{primaryBlue!4}
\sffamily\small
\begin{tabularx}{\linewidth}{|>{\raggedright\arraybackslash \ttfamily\color{accentOrange}\bfseries}m{6.2cm}|c|Z|}
\hline
\rowcolor{primaryBlue} \thd{\sffamily Bộ chọn CSS (Selector)} & \thd{\sffamily Specificity (A,B,C,D)} & \thd{\sffamily Phân tích chi tiết thành phần} \\
\hline
h1 & (0,0,0,1) & 1 thẻ HTML (\codeinline{h1}) \\
\hline
div & (0,0,0,1) & 1 thẻ HTML (\codeinline{div}) \\
\hline
.box / .btn & (0,0,1,0) & 1 class (\codeinline{.box} hoặc \codeinline{.btn}) \\
\hline
\#header / \#main & (0,1,0,0) & 1 ID (\codeinline{\#header} hoặc \codeinline{\#main}) \\
\hline
div p & (0,0,0,2) & 2 thẻ HTML (\codeinline{div} + \codeinline{p}) \\
\hline
p::first-line & (0,0,0,2) & 1 thẻ HTML (\codeinline{p}) + 1 pseudo-element (\codeinline{::first-line}) \\
\hline
a:hover & (0,0,1,1) & 1 thẻ HTML (\codeinline{a}) + 1 pseudo-class (\codeinline{:hover}) \\
\hline
input[type="text"] & (0,0,1,1) & 1 thẻ HTML (\codeinline{input}) + 1 attribute (\codeinline{[type]}) \\
\hline
div:not(.active) & (0,0,1,1) & 1 thẻ HTML (\codeinline{div}) + 1 class bên trong \codeinline{:not()} (\codeinline{.active}) \\
\hline
.nav ul li a & (0,0,1,3) & 1 class (\codeinline{.nav}) + 3 thẻ HTML (\codeinline{ul}, \codeinline{li}, \codeinline{a}) \\
\hline
\#main .button / div\#main .btn:hover & (0,1,2,1) & 1 ID (\codeinline{\#main}) + 2 class (\codeinline{.btn}, \codeinline{:hover}) + 1 thẻ (\codeinline{div}) \\
\hline
a[href][data-id].active:hover & (0,0,4,1) & 1 thẻ (\codeinline{a}) + 2 attributes (\codeinline{[href]}, \codeinline{[data-id]}) + 2 class (\codeinline{.active}, \codeinline{:hover}) \\
\hline
section\#hero .content p & (0,1,1,2) & 1 ID (\codeinline{\#hero}) + 1 class (\codeinline{.content}) + 2 thẻ (\codeinline{section}, \codeinline{p}) \\
\hline
html body .wrapper \#main-content p & (0,1,1,3) & 1 ID (\codeinline{\#main-content}) + 1 class (\codeinline{.wrapper}) + 3 thẻ (\codeinline{html}, \codeinline{body}, \codeinline{p}) \\
\hline
style="color:red" (inline style) & (1,0,0,0) & 1 thuộc tính nội tuyến \codeinline{style="..."} \\
\hline
\end{tabularx}
\end{prettyTableBox}

%% ===================================================================
%% PHẦN 8.4: NHỮNG QUY TẮC ĐẶC BIỆT TRONG SPECIFICITY
%% ===================================================================
\subsection{4. Những Quy Tắc Đặc Biệt Trong Specificity}

\paragraph{4.1. Từ Khóa !important Không Được Tính Vào Specificity Nhưng Ghi Đè Tất Cả}

\begin{conceptBox}[Tại Sao !important Không Tính Vào Specificity?]
\codeinline{!important} là một quy tắc ưu tiên tuyệt đối trong thuật toán Cascade của trình duyệt, nhưng nó \textbf{không được cộng điểm vào bộ số Specificity $(A,B,C,D)$}. Điều này có nghĩa là, dù một quy tắc có specificity thấp, nếu có \codeinline{!important}, nó vẫn được áp dụng thay vì các quy tắc có specificity cao hơn nhưng không có \codeinline{!important}.
\end{conceptBox}

\begin{codeWindow}[Ví Dụ 1: Quy Tắc Specificity Cao Hơn Nhưng Không Có !important]
\begin{lstlisting}[language=CSS3]
/* Specificity: (0,1,0,1) - Khong co !important */
#content p {
  color: red;
}

/* Specificity: (0,0,0,1) - Thap hon nhung CO !important -> WINS! */
p {
  color: blue !important;
}
\end{lstlisting}
\end{codeWindow}

\begin{browserWindow}[Kết Quả Ví Dụ 1: Văn Bản <p> Có Màu Xanh (Blue)]
\sffamily\small
\textcolor{blue}{\textbf{Đoạn văn hiển thị màu blue}} \par\vspace{4pt}
{\scriptsize\textcolor{gray!60}{\textit{Giải thích: \#content p có specificity (0,1,0,1), p có specificity (0,0,0,1) thấp hơn. Nhưng nhờ từ khóa !important, quy tắc p \{ color: blue !important; \} ghi đè tất cả và màu xanh (blue) chiến thắng.}}}
\end{browserWindow}

\begin{codeWindow}[Ví Dụ 2: Nếu Cả Hai Quy Tắc Đều Có !important]
\begin{lstlisting}[language=CSS3]
/* Specificity: (0,1,0,1) + !important -> WINS! */
#content p {
  color: red !important;
}

/* Specificity: (0,0,0,1) + !important -> LOSES! */
p {
  color: blue !important;
}
\end{lstlisting}
\end{codeWindow}

\begin{browserWindow}[Kết Quả Ví Dụ 2: Văn Bản <p> Có Màu Đỏ (Red)]
\sffamily\small
\textcolor{red}{\textbf{Đoạn văn hiển thị màu red}} \par\vspace{4pt}
{\scriptsize\textcolor{gray!60}{\textit{Giải thích: Khi cả hai quy tắc đều có !important, quyền ưu tiên tuyệt đối triệt tiêu lẫn nhau. Trình duyệt sẽ quay lại so sánh specificity: \#content p có (0,1,0,1) cao hơn p (0,0,0,1), nên màu đỏ (red) chiến thắng.}}}
\end{browserWindow}

\paragraph{4.2. Universal Selector (*) Và Combinators (>, +, \~) Không Ảnh Hưởng Đến Specificity}

\begin{conceptBox}[Tại Sao *, >, +, \~ Khóa Ảnh Hưởng Đến Specificity?]
\begin{itemize}[leftmargin=1.5em, itemsep=3pt, parsep=0pt]
    \item \textbf{Universal selector (\codeinline{*}):} Là bộ chọn chung áp dụng cho mọi phần tử, nhưng không có giá trị specificity $(0,0,0,0)$ vì nó không giúp phân biệt một phần tử cụ thể.
    \item \textbf{Combinators (\codeinline{>}, \codeinline{+}, \codeinline{\~{}}):} Chỉ định quan hệ giữa các phần tử (con trực tiếp, anh em liền kề, anh em chung cha), nhưng không làm tăng specificity.
\end{itemize}
\end{conceptBox}

\begin{codeWindow}[Ví Dụ 3: Universal Selector (*)]
\begin{lstlisting}[language=CSS3]
/* Specificity: (0,0,0,0) */
* {
  color: green;
}

/* Specificity: (0,0,0,1) - WINS! */
p {
  color: red;
}
\end{lstlisting}
\end{codeWindow}

\begin{browserWindow}[Kết Quả Ví Dụ 3: Thẻ <p> Có Màu Đỏ (Red)]
\sffamily\small
\textcolor{red}{\textbf{Đoạn văn hiển thị màu red}} \par\vspace{4pt}
{\scriptsize\textcolor{gray!60}{\textit{Giải thích: * \{ color: green; \} có specificity (0,0,0,0), p \{ color: red; \} có specificity (0,0,0,1) cao hơn nên màu đỏ chiến thắng.}}}
\end{browserWindow}

\begin{codeWindow}[Ví Dụ 4: Combinator (>) Không Làm Tăng Specificity]
\begin{lstlisting}[language=CSS3]
/* Specificity: (0,0,0,2) - 2 elements (div, p), combinator > co diem = 0 */
div > p {
  color: blue;
}

/* Specificity: (0,0,0,1) - 1 element */
p {
  color: red;
}
\end{lstlisting}
\end{codeWindow}

\begin{browserWindow}[Kết Quả Ví Dụ 4: div > p (0,0,0,2) Thắng p (0,0,0,1)]
\sffamily\small
\textcolor{blue}{\textbf{Đoạn văn trong div có màu blue}} \par\vspace{4pt}
{\scriptsize\textcolor{gray!60}{\textit{Giải thích: div > p gồm 2 thẻ element (div và p) nên có specificity (0,0,0,2) cao hơn p (0,0,0,1). Bản thân dấu > có điểm = 0.}}}
\end{browserWindow}

\paragraph{4.3. CSS Internal (<style>) Và External (.css File) Không Ảnh Hưởng Độ Ưu Tiên}

\begin{conceptBox}[Tại Sao Cách Viết CSS (Internal Hay External) Không Ảnh Hưởng Độ Ưu Tiên?]
Tất cả các quy tắc CSS (dù nằm trong thẻ \codeinline{<style>} hay file \codeinline{.css} bên ngoài) đều có cùng mức độ ưu tiên dựa trên specificity và vị trí xuất hiện trong tài liệu HTML. Nếu hai quy tắc có cùng specificity, quy tắc nào xuất hiện sau sẽ ghi đè quy tắc trước.
\end{conceptBox}

\begin{codeWindow}[Ví Dụ 5: Internal CSS vs External CSS]
\begin{lstlisting}[language=HTML5]
<!-- File index.html -->
<head>
  <!-- Internal CSS viet TRUOC -->
  <style>
    p { color: red; }
  </style>
  
  <!-- External CSS duoc link SAU -->
  <link rel="stylesheet" href="style.css">
</head>
<body>
  <p>Doan van nay co mau gi?</p>
</body>
\end{lstlisting}
\end{codeWindow}

\begin{browserWindow}[Kết Quả Ví Dụ 5: Thẻ <p> Có Màu Xanh (Blue) Vì style.css Được Gọi Sau]
\sffamily\small
\textcolor{blue}{\textbf{Đoạn văn hiển thị màu blue}} \par\vspace{4pt}
{\scriptsize\textcolor{gray!60}{\textit{Giải thích: Cả hai quy tắc p đều có specificity (0,0,0,1). Vì style.css được chèn sau thẻ <style>, quy tắc p \{ color: blue; \} trong style.css sẽ ghi đè quy tắc trước đó.}}}
\end{browserWindow}

\begin{prettyTableBox}[Bảng Tóm Tắt Quy Tắc Đặc Biệt Trong Specificity]
\rowcolors{2}{white}{primaryBlue!4}
\sffamily\small
\begin{tabularx}{\linewidth}{|>{\raggedright\arraybackslash \bfseries\color{primaryBlue}}m{5.5cm}|Z|}
\hline
\rowcolor{primaryBlue} \thd{\sffamily Quy tắc} & \thd{\sffamily Giải thích nguyên lý} \\
\hline
\codeinline{!important} không tính vào specificity & Nhưng nó ghi đè tất cả, trừ khi một quy tắc khác cũng có \codeinline{!important} và có specificity cao hơn. \\
\hline
Universal selector (\codeinline{*}) không ảnh hưởng & Vì nó không chọn phần tử cụ thể nào, điểm specificity bằng $(0,0,0,0)$. \\
\hline
Combinators (\codeinline{>}, \codeinline{+}, \codeinline{\~{}}) không ảnh hưởng & Vì chúng chỉ định quan hệ giữa phần tử nhưng không làm tăng điểm specificity. \\
\hline
Internal và External CSS không ảnh hưởng & Độ ưu tiên được quyết định strictly bởi specificity và thứ tự xuất hiện trong tài liệu DOM. \\
\hline
\end{tabularx}
\end{prettyTableBox}

%% ===================================================================
%% PHẦN 8.5: KHI NÀO QUY TẮC "VIẾT SAU ĐÈ VIẾT TRƯỚC" ĐÚNG / SAI?
%% ===================================================================
\subsection{5. Phân Tích Chi Tiết: Khi Nào Quy Tắc "Viết Sau Đè Viết Trước" Đúng Hoặc Sai?}

\paragraph{5.1. Khi Nào Quy Tắc Này ĐÚNG?}
\begin{conceptBox}[Điều Kiện Để Quy Tắc Source Order Áp Dụng]
Nếu hai quy tắc CSS có \textbf{cùng điểm specificity} (và không có \codeinline{!important}), quy tắc nào xuất hiện sau trong mã nguồn CSS sẽ được áp dụng.
\end{conceptBox}

\begin{codeWindow}[Ví Dụ 1: Hai Quy Tắc Có Cùng Specificity (0,0,0,1)]
\begin{lstlisting}[language=CSS3]
/* Viet truoc */
p { color: red; }

/* Viet SAU -> WINS! */
p { color: blue; }
\end{lstlisting}
\end{codeWindow}

\begin{browserWindow}[Kết Quả: Thẻ <p> Có Màu Xanh (Blue)]
\sffamily\small
\textcolor{blue}{\textbf{Đoạn văn có màu blue}} \par\vspace{4pt}
{\scriptsize\textcolor{gray!60}{\textit{Giải thích: Cả hai quy tắc p đều có specificity (0,0,0,1). Quy tắc p \{ color: blue; \} xuất hiện sau nên ghi đè lên quy tắc trước.}}}
\end{browserWindow}

\paragraph{5.2. Khi Nào Quy Tắc Này KHÔNG ĐÚNG Nữa? (3 Trường Hợp)}

Có 3 trường hợp khiến quy tắc "viết sau đè viết trước" hoàn toàn không còn đúng:

\begin{warnBox}[Trường Hợp 1: Khi Một Quy Tắc Có !important]
\textbf{\codeinline{!important} luôn ghi đè tất cả}, trừ khi quy tắc khác cũng có \codeinline{!important} và có specificity cao hơn. Dù quy tắc có \codeinline{!important} viết trước, nó vẫn thắng quy tắc viết sau!
\end{warnBox}

\begin{codeWindow}[Ví Dụ 2: !important Ghi Đè Quy Tắc Viết Sau]
\begin{lstlisting}[language=CSS3]
/* Viet TRUOC nhung co !important -> WINS! */
p { color: red !important; }

/* Viet SAU nhung khong co !important -> LOSES! */
p { color: blue; }
\end{lstlisting}
\end{codeWindow}

\begin{browserWindow}[Kết Quả: Thẻ <p> Có Màu Đỏ (Red) Dù Quy Tắc Viết Sau Đặt Màu Blue]
\sffamily\small
\textcolor{red}{\textbf{Đoạn văn có màu red}} \par\vspace{4pt}
{\scriptsize\textcolor{gray!60}{\textit{Giải thích: Quy tắc p \{ color: red !important; \} viết trước nhưng có !important nên chiến thắng hoàn toàn quy tắc p \{ color: blue; \} viết sau.}}}
\end{browserWindow}

\begin{warnBox}[Trường Hợp 2: Khi Một Quy Tắc Có Specificity Cao Hơn]
Nếu một quy tắc có specificity cao hơn, nó sẽ ghi đè lên quy tắc có specificity thấp hơn, \textbf{dù xuất hiện trước hay sau}!
\end{warnBox}

\begin{codeWindow}[Ví Dụ 3: Specificity Cao Hơn Ghi Đè Quy Tắc Viết Sau]
\begin{lstlisting}[language=CSS3]
/* Specificity: (0,0,0,1) - Viet TRUOC */
p { color: red; }

/* Specificity: (0,1,0,1) - (*@ Viết SAU nhưng specificity CAO HƠN -> WINS! @*) */
#content p { color: blue; }
\end{lstlisting}
\end{codeWindow}

\begin{browserWindow}[Kết Quả: \#content p (0,1,0,1) Thắng p (0,0,0,1)]
\sffamily\small
\textcolor{blue}{\textbf{Đoạn văn trong \#content có màu blue}} \par\vspace{4pt}
{\scriptsize\textcolor{gray!60}{\textit{Giải thích: \#content p có specificity (0,1,0,1) cao hơn p (0,0,0,1). Dù bạn có đảo vị trí viết p sau \#content p thì \#content p vẫn luôn chiến thắng!}}}
\end{browserWindow}

\begin{warnBox}[Trường Hợp 3: Khi CSS Nội Tuyến (Inline Styles) Được Sử Dụng]
Inline styles (\codeinline{style="..."}) có specificity cao nhất $(1,0,0,0)$, trừ khi bị \codeinline{!important} ghi đè. Nó luôn ghi đè mọi quy tắc trong file CSS bên ngoài bất kể thứ tự.
\end{warnBox}

\begin{codeWindow}[Ví Dụ 4: Inline Styles Ghi Đè CSS Bên Ngoài]
\begin{lstlisting}[language=HTML5]
<!-- (*@ Inline Style có Specificity = (1,0,0,0) @*) -->
<p style="color: green;">Doan van nay co mau gi?</p>

<style>
  /* Specificity: (0,0,0,1) - Viet SAU nhung thua Inline Style */
  p { color: red; }
</style>
\end{lstlisting}
\end{codeWindow}

\begin{browserWindow}[Kết Quả: Thẻ <p> Có Màu Xanh Lá (Green) Nhờ Inline Style]
\sffamily\small
\textcolor{green!60!black}{\textbf{Đoạn văn hiển thị màu green}} \par\vspace{4pt}
{\scriptsize\textcolor{gray!60}{\textit{Giải thích: style="color: green;" có specificity (1,0,0,0) cao hơn p \{ color: red; \} (0,0,0,1), nên đoạn văn có màu green dù CSS bên ngoài đặt màu red.}}}
\end{browserWindow}

\begin{prettyTableBox}[Bảng Tổng Kết: Khi Nào Quy Tắc "Cái Nào Viết Sau Sẽ Áp Dụng" Đúng Hoặc Sai?]
\rowcolors{2}{white}{primaryBlue!4}
\sffamily\small
\begin{tabularx}{\linewidth}{|>{\raggedright\arraybackslash \bfseries\color{primaryBlue}}Z|c|>{\bfseries}Z|}
\hline
\rowcolor{primaryBlue} \thd{\sffamily Trường hợp xét duyệt} & \thd{\sffamily Quy tắc sau được áp dụng?} & \thd{\sffamily Phân tích nguyên do} \\
\hline
Hai quy tắc có cùng specificity, không có \codeinline{!important} & \textcolor{green!60!black}{\textbf{ĐÚNG}} & Áp dụng nguyên tắc Source Order chuẩn. \\
\hline
Một quy tắc có \codeinline{!important}, quy tắc kia không có & \textcolor{red}{\textbf{SAI}} & Quy tắc có \codeinline{!important} luôn thắng. \\
\hline
Một quy tắc có specificity cao hơn & \textcolor{red}{\textbf{SAI}} & Quy tắc có specificity cao hơn luôn thắng. \\
\hline
Một quy tắc là Inline styles (\codeinline{style="..."}) & \textcolor{red}{\textbf{SAI}} & Inline styles có điểm $(1,0,0,0)$ cao hơn. \\
\hline
\end{tabularx}
\end{prettyTableBox}

%% ===================================================================
%% PHẦN 8.6: BÀI TẬP VÀ VÍ DỤ PHÂN TÍCH CHUYÊN SÂU
%% ===================================================================
\subsection{6. Bài Tập \& Ví Dụ Phân Tích Chuyên Sâu Độ Ưu Tiên}

\paragraph{Ví Dụ Chuyên Sâu 1: Phân Tích Độ Ưu Tiên Danh Sách <li> Trong <ul>}

Xét 2 bộ chọn đối lập nhắm vào danh sách:
\begin{codeWindow}[Ví Dụ Danh Sách: li vs ul li]
\begin{lstlisting}[language=CSS3]
/* Specificity: (0,0,0,1) - 1 element */
li { color: blue; }

/* Specificity: (0,0,0,2) - 2 elements (ul + li) -> WINS! */
ul li { color: red; }
\end{lstlisting}
\end{codeWindow}

\begin{browserWindow}[Kết Quả: Tất Cả Thẻ <li> Đều Có Màu Đỏ (Red)]
\sffamily\small
\textcolor{red}{\textbf{1. Mục danh sách thứ nhất (Red)}} \par
\textcolor{red}{\textbf{2. Mục danh sách thứ hai (Red)}} \par\vspace{4pt}
{\scriptsize\textcolor{gray!60}{\textit{Giải thích: ul li có specificity (0,0,0,2) cao hơn li (0,0,0,1). Do đó tất cả thẻ <li> bên trong <ul> sẽ hiển thị màu đỏ.}}}
\end{browserWindow}

\begin{conceptBox}[Làm Thế Nào Để Đổi Sang Màu Xanh (Blue)? (3 Cách Xử Lý)]
\begin{enumerate}[leftmargin=1.5em, itemsep=3pt, parsep=0pt]
    \item \textbf{Cách 1: Đổi thứ tự CSS?} Đảo \codeinline{ul li \{ color: red; \}} lên trước và \codeinline{li \{ color: blue; \}} xuống sau? \textbf{Vẫn có màu đỏ!} Vì \codeinline{ul li} $(0,0,0,2) > li (0,0,0,1)$, specificity cao hơn luôn thắng bất kể vị trí!
    \item \textbf{Cách 2: Tăng specificity cho bộ chọn muốn đổi?} Viết \codeinline{ul > li \{ color: blue; \}}? Specificity của \codeinline{ul > li} vẫn là $(0,0,0,2)$ bằng với \codeinline{ul li}. Nếu viết \codeinline{ul > li} sau \codeinline{ul li}, màu xanh sẽ thắng nhờ Source Order.
    \item \textbf{Cách 3: Sử dụng !important?} Đặt \codeinline{li \{ color: blue !important; \}} sẽ ngay lập tức ghi đè màu đỏ thành màu xanh!
\end{enumerate}
\end{conceptBox}

\paragraph{Ví Dụ Chuyên Sâu 2: Chuỗi 10 Cấp Độ Leo Thang Specificity Tăng Dần}

Để hiểu thấu đáo bản chất thuật toán so sánh Specificity của trình duyệt, chúng ta cùng phân tích một ví dụ kinh điển với 10 quy tắc CSS cùng nhắm vào thẻ \codeinline{<h2>}. Điểm Specificity sẽ leo thang liên tục từ 1 điểm đến 213 điểm qua từng cấp độ.

\begin{codeWindow}[1. Cấu Trúc HTML Mẫu Thử Nghiệm]
\begin{lstlisting}[language=HTML5]
<div id="main">
  <div id="head" class="head">
    <h2 class="title">Specificity Test</h2>
  </div>
</div>
\end{lstlisting}
\end{codeWindow}

\begin{codeWindow}[2. Chuỗi 10 Quy Tắc CSS Leo Thang Specificity Tăng Dần]
\begin{lstlisting}[language=CSS3]
/* Cap 1: The don le (0,0,0,1) -> ~1 pt */
h2 { color: red; }

/* Cap 2: Class + The (0,0,1,1) -> ~11 pts */
h2.title { color: blue; }

/* Cap 3: Them div cha (0,0,1,2) -> ~12 pts */
div h2.title { color: orange; }

/* Cap 4: Them class .head (0,0,2,1) -> ~21 pts */
.head h2.title { color: green; }

/* Cap 5: Class cha + Tag div (0,0,2,2) -> ~22 pts */
div.head h2.title { color: gray; }

/* Cap 6: Them Pseudo-class :last-child (0,0,3,2) -> ~32 pts */
div.head h2.title:last-child { color: yellow; }

/* Cap 7: Them 1 ID #main (0,1,2,2) -> ~122 pts */
#main div.head h2.title { color: pink; }

/* Cap 8: Them ID thu 2 #head (0,2,1,2) -> ~212 pts */
#main div#head h2.title { color: purple; }

/* Cap 9: Pseudo-element chu dau (0,2,1,3) -> ~213 pts */
#main div#head h2.title::first-letter { color: red; }

/* Cap 10: Pseudo-element it class (0,2,0,3) -> ~203 pts */
#main div#head h2::first-letter { color: yellow; }
\end{lstlisting}
\end{codeWindow}

\begin{prettyTableBox}[3. Bảng Kết Quả Phân Tích 10 Cấp Độ Leo Thang Specificity]
\rowcolors{2}{white}{primaryBlue!4}
\sffamily\small
\begin{tabularx}{\linewidth}{|c|>{\raggedright\arraybackslash \ttfamily\color{accentOrange}\bfseries}m{5.5cm}|c|Z|}
\hline
\rowcolor{primaryBlue} \thd{\sffamily Cấp} & \thd{\sffamily Selector CSS} & \thd{\sffamily Specificity} & \thd{\sffamily Kết quả \& Phân tích lý do} \\
\hline
1 & h2 & (0,0,0,1) & \textcolor{red!85!black}{\textbf{Red}} -- Bị Cấp 2 đè vì Cấp 2 có thêm Class \codeinline{.title}. \\
\hline
2 & h2.title & (0,0,1,1) & \textcolor{blue!85!black}{\textbf{Blue}} -- Bị Cấp 3 đè vì Cấp 3 có thêm thẻ \codeinline{div} cha. \\
\hline
3 & div h2.title & (0,0,1,2) & \textcolor{orange!90!black}{\textbf{Orange}} -- Bị Cấp 4 đè vì Cấp 4 có 2 Class: $C{=}2 > C{=}1$. \\
\hline
4 & .head h2.title & (0,0,2,1) & \textcolor{green!60!black}{\textbf{Green}} -- Bị Cấp 5 đè vì Cấp 5 có thêm thẻ \codeinline{div}. \\
\hline
5 & div.head h2.title & (0,0,2,2) & \textcolor{gray!80!black}{\textbf{Gray}} -- Bị Cấp 6 đè vì \codeinline{:last-child} cộng thêm 1C. \\
\hline
6 & div.head h2.title:last-child & (0,0,3,2) & \textcolor{yellow!75!black}{\textbf{Yellow}} -- Bị Cấp 7 đè bẹp vì ID \codeinline{\#main}: $B{=}1 > B{=}0$. \\
\hline
7 & \#main div.head h2.title & (0,1,2,2) & \textcolor{magenta!90!black}{\textbf{Pink}} -- Bị Cấp 8 đè vì thêm ID thứ 2 \codeinline{\#head}: $B{=}2$. \\
\hline
8 & \#main div\#head h2.title & (0,2,1,2) & \textcolor{purple!90!black}{\textbf{Purple [VÔ ĐỊCH H2]}} -- Điểm cao nhất nhắm vào thân \codeinline{h2}! \\
\hline
9 & \#main div\#head h2.title::first-letter & (0,2,1,3) & \textcolor{red!90!black}{\textbf{Red [VÔ ĐỊCH CHỮ 'S']}} -- Nhắm vào \codeinline{::first-letter}, thắng Cấp 10. \\
\hline
10 & \#main div\#head h2::first-letter & (0,2,0,3) & \textcolor{yellow!75!black}{\textbf{Yellow}} -- Thua Cấp 9 vì thiếu Class \codeinline{.title}: $C{=}0 < C{=}1$. \\
\hline
\end{tabularx}
\end{prettyTableBox}

\begin{conceptBox}[4. Giải Đáp 3 Thắc Mắc Cốt Lõi Về Thuật Toán So Sánh]
\begin{enumerate}[leftmargin=1.5em, itemsep=4pt, parsep=0pt]
    \item \textbf{Tại sao Cấp 4 (0,0,2,1) màu Green lại thắng Cấp 3 (0,0,1,2) màu Orange?} \par
    Khi so sánh hai điểm số $(0,0,2,1)$ và $(0,0,1,2)$, trình duyệt so sánh từ trái sang phải: Nhóm A bằng nhau (0=0), Nhóm B bằng nhau (0=0), đến Nhóm C: Cấp 4 có 2 Class $(C=2)$, Cấp 3 chỉ có 1 Class $(C=1)$. Vì $2 > 1$, Cấp 4 chiến thắng \textbf{mặc dù Cấp 3 có nhiều thẻ HTML hơn}!
    
    \item \textbf{Tại sao Cấp 7 (0,1,2,2) màu Pink lại đánh bại Cấp 6 (0,0,3,2) màu Yellow?} \par
    Cấp 6 dù gom tới 3 Class/Pseudo-class (điểm $C=3$), nhưng Cấp 7 xuất hiện 1 bộ chọn ID \codeinline{\#main} (điểm $B=1$). Vì nhóm ID ($B$) nằm ở vị trí ưu tiên cao hơn nhóm Class ($C$), nên $B=1$ lập tức đè bẹp $B=0$ bất kể Cấp 6 có bao nhiêu Class đi nữa!
    
    \item \textbf{Tại sao chữ 'S' lại có màu Red còn toàn bộ chữ 'pecificity Test' lại màu Purple?} \par
    \begin{itemize}[leftmargin=1.2em, itemsep=2pt]
        \item Cấp 8 nhắm vào toàn bộ thẻ \codeinline{<h2>} và đạt điểm Specificity $(0,2,1,2)$ cao nhất $\rightarrow$ Tô màu \textbf{\textcolor{purple}{Purple}} cho thân thẻ.
        \item Cấp 9 dùng pseudo-element \codeinline{::first-letter} nhắm riêng vào chữ cái đầu tiên \codeinline{'S'} với điểm $(0,2,1,3)$ thắng Cấp 10 $(0,2,0,3)$ $\rightarrow$ Tô màu \textbf{\textcolor{red}{Red}} riêng cho chữ 'S'.
    \end{itemize}
\end{enumerate}
\end{conceptBox}

\begin{browserWindow}[5. Giao Diện Mô Phỏng Kết Quả Trình Duyệt Thực Tế]
\centering\vspace{8pt}
{\sffamily\Huge\bfseries \textcolor{red}{S}\textcolor{purple}{pecificity Test}}
\par\vspace{10pt}
\sffamily\small
\mbox{\textcolor{red}{\textbf{Chữ 'S' màu Đỏ (Red) [Quy tắc 9 thắng]}}} \qquad \mbox{\textcolor{purple}{\textbf{Chữ 'pecificity Test' màu Tím (Purple) [Quy tắc 8 thắng]}}}
\par\vspace{6pt}
\end{browserWindow}

%% ===================================================================
%% PHẦN 8.7: GIẢI THÍCH CHI TIẾT :not(), :is(), :where()
%% ===================================================================
\subsection{7. Pseudo-class :not(), :is() Và :where() Trong CSS Modern}

\paragraph{7.1. Pseudo-class :not() -- Bản Thân Không Tính Điểm}

\begin{conceptBox}[Quy Tắc Đặc Biệt Của :not()]
\begin{itemize}[leftmargin=1.5em, itemsep=3pt, parsep=0pt]
    \item \textbf{Bản thân \codeinline{:not()} KHÔNG được tính điểm specificity.} Nó là một khung vô hình.
    \item \textbf{Chỉ có nội dung bên trong ngoặc đơn mới được tính.} Ví dụ: \codeinline{:not(.active)} cộng 1 class (\codeinline{.active}) vào nhóm C.
    \item \textbf{Kết quả}: \codeinline{div:not(.active)} có specificity = \codeinline{(0,0,1,1)} -- gồm 1 thẻ \codeinline{div} (D) + 1 class \codeinline{.active} (C).
\end{itemize}
\end{conceptBox}

\paragraph{7.2. Pseudo-class :is() -- Lấy Specificity Của Đối Số Cao Nhất}

\begin{conceptBox}[Quy Tắc Của :is()]
\codeinline{:is()} gom nhóm selector và lấy specificity của \textbf{ đối số có điểm CAO NHẤT} trong danh sách. Ví dụ: \codeinline{:is(.btn, \#submit)} lấy specificity của \codeinline{\#submit} $(0,1,0,0)$.
\end{conceptBox}

\paragraph{7.3. Pseudo-class :where() -- Specificity Luôn Bằng 0}

\begin{conceptBox}[Quy Tắc Của :where() -- Zero Specificity]
\codeinline{:where()} luôn có \textbf{specificity = (0,0,0,0)}, bất kể nội dung bên trong là gì. Rất hữu ích để viết CSS mặc định trong thư viện dễ ghi đè.
\end{conceptBox}

\begin{prettyTableBox}[Bảng So Sánh Nhanh :not() vs :is() vs :where()]
\rowcolors{2}{white}{primaryBlue!4}
\sffamily\small
\begin{tabularx}{\linewidth}{|>{\raggedright\arraybackslash \bfseries\color{primaryBlue}\ttfamily}m{3.2cm}|Z|Z|Z|}
\hline
\rowcolor{primaryBlue} \thd{\sffamily Pseudo-class} & \thd{\sffamily Cách tính Specificity} & \thd{\sffamily Ví dụ} & \thd{\sffamily Specificity kết quả} \\
\hline
:not(.active) & Lấy specificity của đối số bên trong & \codeinline{div:not(.active)} & (0,0,1,1) \\
\hline
:is(.btn, \#id) & Lấy specificity của đối số \textbf{cao nhất} & \codeinline{:is(.btn, \#submit)} & (0,1,0,0) \\
\hline
:where(.btn, \#id) & \textbf{Luôn = 0}, bất kể nội dung & \codeinline{:where(\#main, .box)} & (0,0,0,0) \\
\hline
\end{tabularx}
\end{prettyTableBox}

%% ===================================================================
%% PHẦN 8.8: SƠ ĐỒ TÍNH ĐIỂM TỪNG BƯỚC CHO SELECTOR PHỨC TẠP
%% ===================================================================
\subsection{8. Sơ Đồ Tính Điểm Specificity Từng Bước (Step-by-Step)}

\paragraph{Ví dụ minh họa: Tính specificity cho \codeinline{div\#main .btn:hover}}

\begin{prettyTableBox}[Sơ Đồ Tính Điểm Từng Bước Cho Selector: div\#main .btn:hover]
\rowcolors{2}{white}{primaryBlue!4}
\sffamily\small
\begin{tabularx}{\linewidth}{|c|>{\raggedright\arraybackslash \ttfamily\color{accentOrange}\bfseries}m{3.5cm}|>{\bfseries\color{primaryBlue}}m{5cm}|c|Z|}
\hline
\rowcolor{primaryBlue} \thd{\sffamily Bước} & \thd{\sffamily Thành phần} & \thd{\sffamily Phân loại} & \thd{\sffamily Nhóm} & \thd{\sffamily Cộng điểm} \\
\hline
1 & div & Thẻ HTML (Element) & D & +1 vào D \\
\hline
2 & \#main & ID Selector & B & +1 vào B \\
\hline
3 & .btn & Class Selector & C & +1 vào C \\
\hline
4 & :hover & Pseudo-class & C & +1 vào C \\
\hline
\multicolumn{3}{|r|}{\textbf{Tổng cộng:}} & \multicolumn{2}{l|}{\textbf{A=0, B=1, C=2, D=1}} \\
\hline
\multicolumn{3}{|r|}{\textbf{Specificity:}} & \multicolumn{2}{l|}{\codeinline{(0, 1, 2, 1)}} \\
\hline
\end{tabularx}
\end{prettyTableBox}

%% ===================================================================
%% PHẦN 8.9: KINH NGHIỆM THỰC TẾ & BEST PRACTICES
%% ===================================================================
\subsection{9. Kinh Nghiệm Thực Tế \& Tối Ưu Kiến Trúc CSS}

\begin{prettyTableBox}[Bảng Quy Tắc "Nên \& Không Nên" Trong Thực Tế Lập Trình Frontend]
\rowcolors{2}{white}{primaryBlue!4}
\sffamily\small
\begin{tabularx}{\linewidth}{|>{\raggedright\arraybackslash \bfseries\color{red!80!black}}Z|>{\bfseries\color{green!60!black}}Z|}
\hline
\rowcolor{primaryBlue} \thd{\sffamily [KHÔNG NÊN] làm trong CSS} & \thd{\sffamily [NÊN] thực hiện (Best Practices)} \\
\hline
Dùng ID làm selector chính trong CSS (\codeinline{\#header \{ background: blue; \}}) làm code bị cứng, khó ghi đè. & Sử dụng Class cho tất cả các bộ chọn (\codeinline{.header \{ background: blue; \}}) để linh hoạt và tái sử dụng. \\
\hline
Lạm dụng \codeinline{!important} để giải quyết xung đột CSS tạm thời, gây rác mã nguồn. & Chỉ dùng \codeinline{!important} cho các utility class đặc biệt (như \codeinline{.d-none}) hoặc override thư viện ngoài. \\
\hline
Lồng selector quá sâu: \codeinline{body div.container section.main article.post p} gây tăng điểm specificity không cần thiết. & Viết selector ngắn gọn, phẳng (\term{Flat Selectors}): \codeinline{.post p} hoặc áp dụng phương pháp luận đặt tên BEM (\codeinline{.post\_\_text}). \\
\hline
\end{tabularx}
\end{prettyTableBox}

\begin{prettyTableBox}[Bảng Mẹo Nhớ Nhanh Sức Mạnh Của Các Loại Selector (Memory Tips)]
\rowcolors{2}{white}{primaryBlue!4}
\sffamily\small
\begin{tabularx}{\linewidth}{|>{\bfseries\raggedright}m{4.5cm}|Z|}
\hline
\rowcolor{primaryBlue} \thd{\sffamily Mức độ Selector} & \thd{\sffamily Mẹo ghi nhớ dễ thuộc lòng} \\
\hline
\textcolor{red!85!black}{!important} & \textbf{Nói là làm} -- Luôn luôn chiến thắng mọi bộ chọn khác. \\
\hline
\textcolor{orange!85!black}{Inline Style \scriptsize(1,0,0,0)} & \textbf{Gần nhất phần tử} -- Viết trực tiếp trong thuộc tính \codeinline{style="..."} trên thẻ HTML. \\
\hline
\textcolor{primaryBlue}{ID Selector \scriptsize(0,1,0,0)} & \textbf{Định danh duy nhất} -- Rất mạnh mẽ, định danh duy nhất cho phần tử. \\
\hline
\textcolor{green!60!black}{Class / Attr / Pseudo-class \scriptsize(0,0,1,0)} & \textbf{Chung nhóm} -- Ổn định, linh hoạt, được sử dụng phổ biến nhất trong thực tế. \\
\hline
\textcolor{gray!80!black}{Element / Pseudo-element \scriptsize(0,0,0,1)} & \textbf{Thẻ HTML} -- Mức độ yếu nhất trong các loại bộ chọn trực tiếp. \\
\hline
\textcolor{gray!60}{Inheritance / Kế thừa \scriptsize(0,0,0,0)} & \textbf{Mặc định} -- Dễ bị bất kỳ bộ chọn trực tiếp nào ghi đè. \\
\hline
\end{tabularx}
\end{prettyTableBox}

\begin{conceptBox}[Tổng Kết Quy Tắc Vàng Về Thứ Tự Ưu Tiên CSS]
\sffamily
\textbf{Công thức tổng quát:} \quad \codeinline{Specificity = (A, B, C, D)} \quad \text{hoặc} \quad \codeinline{Specificity = (a, b, c, d)}
\par\vspace{8pt}
\textbf{Chuỗi thứ tự ưu tiên tuyệt đối từ cao xuống thấp (1 $\rightarrow$ 6):}
\par\vspace{6pt}
\begin{enumerate}[leftmargin=1.8em, itemsep=4pt, parsep=0pt]
    \item \textcolor{red!85!black}{\textbf{!important trong CSS}} -- Cấp quyền ưu tiên tuyệt đối (ghi đè mọi bộ chọn).
    \item \textcolor{orange!85!black}{\textbf{Style nội tuyến (Inline style) \codeinline{(1,0,0,0)}}} -- Đặt trực tiếp trong thuộc tính \codeinline{style="..."}.
    \item \textcolor{primaryBlue}{\textbf{Bộ chọn ID (\#id) \codeinline{(0,1,0,0)}}} -- Mức ưu tiên cao cho định danh duy nhất.
    \item \textcolor{green!60!black}{\textbf{Class, Attribute, Pseudo-class \codeinline{(0,0,1,0)}}} -- Nhóm bộ chọn phổ biến (\codeinline{.class}, \codeinline{[attr]}, \codeinline{:hover}).
    \item \textcolor{gray!80!black}{\textbf{Thẻ HTML \& Pseudo-element \codeinline{(0,0,0,1)}}} -- Nhóm thẻ phần tử (\codeinline{div}, \codeinline{p}, \codeinline{::before}).
    \item \textcolor{gray!50!black}{\textbf{Kế thừa \& Universal Selector \codeinline{(0,0,0,0)}}} -- Kế thừa mặc định và bộ chọn chung (\codeinline{*}) -- \textit{Yếu nhất}.
\end{enumerate}
\end{conceptBox}


%% ==============================================================================
%% MỤC 10: THUỘC TÍNH OPACITY (ĐỘ MỜ & ĐỘ TRONG SUỐT TRONG CSS)
%% ==============================================================================
\section{Thuộc Tính Opacity (Độ Mờ \& Độ Trong Suốt Trong CSS)}
\label{sec:css_opacity_masterclass}

Thuộc tính \codeinline{opacity} trong CSS cho phép lập trình viên kiểm soát mức độ hiển thị, độ mờ và độ trong suốt của bất kỳ phần tử HTML nào trên giao diện web.

\paragraph{1. Opacity Là Gì?}

\begin{conceptBox}[Khái Niệm \& Giá Trị Của Thuộc Tính Opacity]
\textbf{Opacity} xác định độ trong suốt của một phần tử HTML. Nó quyết định xem ánh sáng/màu nền phía sau phần tử có thể nhìn xuyên qua được hay không.
\par\vspace{4pt}
\textbf{Dải giá trị chấp nhận từ 0 đến 1:}
\begin{itemize}[leftmargin=1.5em, itemsep=3pt, parsep=0pt]
    \item \textbf{0} $\rightarrow$ Hoàn toàn trong suốt (phần tử vô hình nhưng vẫn chiếm vị trí không gian trên luồng tài liệu).
    \item \textbf{1} $\rightarrow$ Không trong suốt (hiển thị rõ ràng $100\%$ như bình thường).
    \item \textbf{Giữa 0 và 1 (ví dụ 0.1 -- 0.9)} $\rightarrow$ Mờ dần theo tỷ lệ phần trăm (ví dụ: \codeinline{0.5} mờ $50\%$, \codeinline{0.2} mờ $80\%$).
\end{itemize}
\end{conceptBox}

\paragraph{2. Cú Pháp Khai Báo CSS:}

\begin{codeWindow}[Cú Pháp Định Mức Độ Mờ Opacity]
\begin{lstlisting}[language=CSS3]
selector {
    opacity: value; /* Gia tri so tu 0.0 den 1.0 */
}
\end{lstlisting}
\end{codeWindow}

\begin{itemize}[leftmargin=1.5em, itemsep=2pt, parsep=0pt]
    \item \textbf{0}: Hoàn toàn trong suốt (ẩn hoàn toàn).
    \item \textbf{1}: Hoàn toàn không trong suốt (hiển thị rõ ràng).
    \item \textbf{Giữa 0 và 1}: Mờ dần theo mức độ (ví dụ: \codeinline{0.5} là mờ $50\%$).
\end{itemize}

\paragraph{3. Ví Dụ Minh Họa \& Kết Quả Trực Quan Trên Trình Duyệt:}

\begin{codeWindow}[Khai Báo 3 Mức Độ Mờ Cho Các Khối Div]
\begin{lstlisting}[language=HTML5]
<div class="box full-opacity">Opacity: 1</div>
<div class="box semi-opacity">Opacity: 0.5</div>
<div class="box low-opacity">Opacity: 0.2</div>

<style>
.box {
    width: 200px;
    height: 100px;
    margin: 10px;
    padding: 10px;
    color: white;
    text-align: center;
    line-height: 100px;
    font-size: 20px;
}
.full-opacity {
    background: #ff0000;
    opacity: 1; /* Khong mo */
}
.semi-opacity {
    background: #007bff;
    opacity: 0.5; /* Mo 50% */
}
.low-opacity {
    background: #28a745;
    opacity: 0.2; /* Mo 80% */
}
</style>
\end{lstlisting}
\end{codeWindow}

\begin{browserWindow}[Kết Quả Trực Quan 3 Khối Box Trên Nền Ca-rô Trong Suốt Mờ Trình Duyệt]
\sffamily\small
\begin{center}
\begin{tikzpicture}[scale=0.95, transform shape]
  % 1. Background Grid (Nen ca-ro mo phong Chrome/Photoshop transparency grid)
  \begin{scope}
    \clip (-0.2,-0.2) rectangle (11.5, 2.2);
    \foreach \x in {0,0.4,...,12} {
      \foreach \y in {0,0.4,...,2.4} {
        \pgfmathparse{int(mod(\x*2.5 + \y*2.5, 2))}
        \ifnum\pgfmathresult=0
          \fill[gray!20] (\x,\y) rectangle (\x+0.4,\y+0.4);
        \else
          \fill[white] (\x,\y) rectangle (\x+0.4,\y+0.4);
        \fi
      }
    }
  \end{scope}
  \draw[gray!35, rounded corners=4pt] (-0.2,-0.2) rectangle (11.5, 2.2);

  % Box 1: Opacity 1.0 (Solid Red - #ff0000)
  \node[draw=red!90!black, fill=red, fill opacity=1.0, text=white, font=\bfseries\sffamily, minimum width=2.9cm, minimum height=1.4cm, rounded corners=3pt] at (1.7, 1.0) {Opacity 1.0};
  \node[anchor=north, font=\scriptsize\sffamily, text=green!60!black] at (1.7, -0.3) {\textbf{Box 1: Đỏ đặc 100\%}};
  \node[anchor=south, font=\tiny\ttfamily\bfseries, text=accentOrange] at (1.7, 2.3) {opacity: 1;};

  % Box 2: Opacity 0.5 (Semi-transparent Blue - #007bff)
  \node[draw=blue!90!black, fill=blue!80!cyan, fill opacity=0.5, text=white, text opacity=1, font=\bfseries\sffamily, minimum width=2.9cm, minimum height=1.4cm, rounded corners=3pt] at (5.6, 1.0) {Opacity 0.5};
  \node[anchor=north, font=\scriptsize\sffamily, text=primaryBlue] at (5.6, -0.3) {\textbf{Box 2: Xanh mờ 50\%}};
  \node[anchor=south, font=\tiny\ttfamily\bfseries, text=accentOrange] at (5.6, 2.3) {opacity: 0.5;};

  % Box 3: Opacity 0.2 (Low Opacity Green - #28a745)
  \node[draw=green!70!black, fill=green!70!black, fill opacity=0.2, text=green!30!black, text opacity=1, font=\bfseries\sffamily, minimum width=2.9cm, minimum height=1.4cm, rounded corners=3pt] at (9.5, 1.0) {Opacity 0.2};
  \node[anchor=north, font=\scriptsize\sffamily, text=red!80!black] at (9.5, -0.3) {\textbf{Box 3: Xanh mờ 80\%}};
  \node[anchor=south, font=\tiny\ttfamily\bfseries, text=accentOrange] at (9.5, 2.3) {opacity: 0.2;};
\end{tikzpicture}
\end{center}
\end{browserWindow}

\paragraph{4. Lưu Ý Quan Trọng Khi Sử Dụng Opacity:}

\begin{warnBox}[3 Cạm Bẫy Quan Trọng Cần Ghi Nhớ Với Opacity]
\begin{enumerate}[leftmargin=1.5em, itemsep=4pt, parsep=0pt]
    \item \textbf{Opacity làm mờ CẢ phần tử lẫn tất cả nội dung con bên trong:} Khi bạn áp dụng \codeinline{opacity} cho một thẻ cha, toàn bộ chữ, hình ảnh, nút bấm con bên trong thẻ cha đó \textbf{đều bị mờ theo}.
\begin{codeWindow}[Lỗi Tự Động Mờ Nội Dung Con]
\begin{lstlisting}[language=CSS3]
.parent {
    background: purple;
    opacity: 0.5; /* Chu va nut bam con ben trong cung bi mo 50%! */
}
\end{lstlisting}
\end{codeWindow}

    \item \textbf{Phần tử vẫn chiếm không gian luồng tài liệu dù \codeinline{opacity: 0}:} Phần tử không bị xóa khỏi bố cục trang mà chỉ bị biến thành vô hình. (Nếu muốn ẩn hoàn toàn và xóa khỏi luồng bố cục, hãy dùng \codeinline{display: none;}).

    \item \textbf{Không thể click tương tác khi phần tử vô hình (\codeinline{opacity: 0}):} Mặc dù phần tử vẫn tồn tại nhưng người dùng không thể nhìn thấy để thao tác (tương tự như bị vô hình).
\end{enumerate}
\end{warnBox}

\paragraph{5. Thay Thế Bằng RGBA Để Chỉ Làm Mờ Nền (Giữ Chữ Rõ Nét):}

Nếu bạn muốn \textbf{chỉ làm mờ màu nền} nhưng vẫn giữ nội dung văn bản và phần tử con hiển thị rõ nét $100\%$, hãy sử dụng giá trị màu \codeinline{rgba()} thay vì thuộc tính \codeinline{opacity}:

\begin{codeWindow}[Sử Dụng RGBA Để Chỉ Mờ Màu Nền]
\begin{lstlisting}[language=CSS3]
.box {
    background: rgba(255, 0, 0, 0.5); /* Nen do mo 50%, van ban ben trong van ro 100% */
}
\end{lstlisting}
\end{codeWindow}

\begin{conceptBox}[Cấu Trúc Giá Trị Màu RGBA (Red, Green, Blue, Alpha)]
\textbf{RGBA (Red, Green, Blue, Alpha):}
\begin{itemize}[leftmargin=1.5em, itemsep=2pt, parsep=0pt]
    \item \textbf{R}: Giá trị kênh Đỏ (Red: $0 \rightarrow 255$).
    \item \textbf{G}: Giá trị kênh Xanh lá (Green: $0 \rightarrow 255$).
    \item \textbf{B}: Giá trị kênh Xanh dương (Blue: $0 \rightarrow 255$).
    \item \textbf{A}: Kênh Alpha độ trong suốt ($0 \rightarrow 1$).
\end{itemize}
\par\vspace{4pt}
\textbf{Ưu điểm vượt trội của RGBA:} Chỉ làm mờ riêng màu nền, nội dung văn bản và các phần tử con bên trong \textbf{hoàn toàn không bị ảnh hưởng}!
\end{conceptBox}

\paragraph{6. Các Ứng Dụng Thực Tế Phổ Biến Của Opacity:}

\begin{conceptBox}[4 Kịch Bản Ứng Dụng Opacity Trong Dự Án Web Thực Tế]
\begin{enumerate}[leftmargin=1.5em, itemsep=5pt, parsep=0pt]
    \item \textbf{Hiệu ứng di chuột (Hover Effect):} Làm mờ nhẹ thẻ hoặc ảnh khi di chuột qua để báo hiệu khả năng nhấp chuột:
\begin{codeWindow}[Làm Mờ Khi Di Chuột Di Vào Card]
\begin{lstlisting}[language=CSS3]
.card:hover {
    opacity: 0.7;
}
\end{lstlisting}
\end{codeWindow}

    \item \textbf{Màn hình phủ tải trang (Loading Screen Overlay):} Làm mờ tối toàn bộ trang web khi hiển thị cửa sổ tải dữ liệu:
\begin{codeWindow}[Tạo Màn Phủ Tối Overlay]
\begin{lstlisting}[language=CSS3]
.overlay {
    position: fixed;
    top: 0;
    left: 0;
    width: 100%;
    height: 100%;
    background: rgba(0, 0, 0, 0.5);
}
\end{lstlisting}
\end{codeWindow}

    \item \textbf{Hiệu ứng ảnh mờ dần / Hiện dần (CSS Fade Animation):}
\begin{codeWindow}[Tạo Animation Hiện Dần Vừa Mờ Đến Rõ]
\begin{lstlisting}[language=CSS3]
@keyframes fade {
    0% { opacity: 0; }
    100% { opacity: 1; }
}
\end{lstlisting}
\end{codeWindow}

    \item \textbf{Làm nổi bật \& Vô hiệu hóa nút bấm (Disabled State):} Kết hợp \codeinline{opacity} với \codeinline{pointer-events: none} để vừa làm mờ vừa khóa tương tác nhấp chuột:
\begin{codeWindow}[Vô Hiệu Hóa Thao Tác Click Vẫn Làm Mờ Nut]
\begin{lstlisting}[language=CSS3]
.disabled {
    opacity: 0.5;
    pointer-events: none; /* Vo hieu hoa moi tuong tac click */
}
\end{lstlisting}
\end{codeWindow}
\end{enumerate}
\end{conceptBox}

\paragraph{7. Bảng Tổng Kết Trực Quan Kiến Thức Opacity:}

\begin{prettyTableBox}[Bảng Tổng Kết Tra Cứu Toàn Diện Về Thuộc Tính Opacity \& Giải Pháp RGBA]
\rowcolors{2}{white}{primaryBlue!4}
\sffamily\small
\begin{tabularx}{\linewidth}{|>{\raggedright\arraybackslash\ttfamily\color{accentOrange}\bfseries}p{3.2cm}|>{\raggedright\arraybackslash}X|>{\raggedright\arraybackslash\ttfamily\color{primaryBlue}}p{4.2cm}|}
\hline
\rowcolor{primaryBlue} \thd{\sffamily Thuộc tính / Giá trị} & \thd{\sffamily Ý nghĩa \& Hành vi hiển thị} & \thd{\sffamily Cú pháp ví dụ minh họa} \\
\hline
opacity: 1 & Phần tử hoàn toàn hiển thị (không mờ). & opacity: 1; \\
\hline
opacity: 0.5 & Phần tử và mọi thẻ con mờ $50\%$. & opacity: 0.5; \\
\hline
opacity: 0 & Phần tử hoàn toàn trong suốt (vẫn chiếm bố cục). & opacity: 0; \\
\hline
pointer-events & Kết hợp với opacity để vô hiệu hóa click chuột. & pointer-events: none; \\
\hline
rgba(...) & Chỉ làm mờ màu nền, giữ nguyên chữ con $100\%$. & background: rgba(0, 0, 0, 0.5); \\
\hline
\end{tabularx}
\end{prettyTableBox}

\begin{conceptBox}[Tóm Lại Cốt Lõi Kiến Thức]
\begin{itemize}[leftmargin=1.5em, itemsep=3pt, parsep=0pt]
    \item \codeinline{opacity} là công cụ cực kỳ mạnh mẽ để tạo hiệu ứng mờ, lớp phủ overlay, và làm nổi bật nội dung giao diện.
    \item Thuộc tính này tác động đến \textbf{cả phần tử cha lẫn toàn bộ phần tử con}. Nếu chỉ muốn mờ màu nền, hãy sử dụng \codeinline{rgba()} hoặc kết hợp \codeinline{pointer-events} để kiểm soát trải nghiệm người dùng tốt nhất.
\end{itemize}
\end{conceptBox}

%% ==============================================================================
%% MỞ RỘNG KIẾN THỨC NÂNG CAO VỀ OPACITY (SENIOR FRONTEND MASTERCLASS)
%% ==============================================================================
\paragraph{Mở Rộng Kiến Thức \& Kinh Nghiệm Thực Tế Chuyên Sâu Về Opacity}
\mbox{}\par\vspace{2pt}

\begin{conceptBox}[1. Nguyên Lý Tạo Stacking Context Trực Tiếp Từ Opacity < 1]
Trong cơ chế render của trình duyệt, khi một phần tử có thuộc tính \codeinline{opacity < 1}, trình duyệt tự động tạo ra một \textbf{Ngữ cảnh xếp chồng mới (Stacking Context)} trên phần tử đó!
\par\vspace{4pt}
\textbf{Hệ quả quan trọng:} Thẻ con chứa \codeinline{z-index: 9999} bên trong phần tử có \codeinline{opacity: 0.99} sẽ \textbf{KHÔNG THỂ} ngoi lên trên các phần tử bên ngoài có z-index cao hơn thẻ cha! Đây là nguyên nhân khiến nhiều lập trình viên Junior loay hoay không hiểu tại sao z-index của thẻ con bị vô hiệu hóa.
\end{conceptBox}

\begin{prettyTableBox}[So Sánh 3 Phương Pháp Ẩn Phần Tử HTML Trong CSS (opacity vs visibility vs display)]
\rowcolors{2}{white}{primaryBlue!4}
\sffamily\small
\begin{tabularx}{\linewidth}{|>{\raggedright\arraybackslash\bfseries\color{primaryBlue}}p{3.6cm}|>{\raggedright\arraybackslash}X|>{\raggedright\arraybackslash}X|>{\raggedright\arraybackslash}X|}
\hline
\rowcolor{primaryBlue} \thd{\sffamily Tiêu chí so sánh} & \thd{\sffamily opacity: 0;} & \thd{\sffamily visibility: hidden;} & \thd{\sffamily display: none;} \\
\hline
Chiếm không gian layout? & \textcolor{green!60!black}{\textbf{CÓ}} (Vẫn chiếm diện tích). & \textcolor{green!60!black}{\textbf{CÓ}} (Vẫn chiếm diện tích). & \textcolor{red}{\textbf{KHÔNG}} (Xóa khỏi luồng layout). \\
\hline
Click / Sự kiện DOM? & Vẫn nhận click ngoại trừ khi dùng \codeinline{pointer-events: none}. & \textcolor{red}{\textbf{KHÔNG}} nhận sự kiện mouse/click. & \textcolor{red}{\textbf{KHÔNG}} nhận bất kỳ sự kiện nào. \\
\hline
Hỗ trợ CSS Transition / Animation? & \textcolor{green!60!black}{\textbf{CÓ}} (Mượt mà từ $0 \rightarrow 1$). & Hỗ trợ hạn chế (chỉ bật/tắt ở cuối). & \textcolor{red}{\textbf{KHÔNG}} (Biến mất tức thì, không transition). \\
\hline
Tăng tốc phần cứng GPU? & \textcolor{green!60!black}{\textbf{CÓ}} (Dùng GPU Composite Layer). & \textcolor{red}{\textbf{KHÔNG}}. & \textcolor{red}{\textbf{KHÔNG}}. \\
\hline
\end{tabularx}
\end{prettyTableBox}

\begin{tipBox}[2. Tối Ưu Hiệu Năng Animation Bằng Tăng Tốc Phần Cứng GPU]
\codeinline{opacity} cùng với \codeinline{transform} là hai thuộc tính CSS duy nhất được xử lý trực tiếp trên tầng \textbf{Composite Layer của GPU (Graphics Processing Unit)}. Khi thay đổi \codeinline{opacity}, trình duyệt \textbf{KHÔNG cần phải tính toán lại Layout (Reflow) hay vẽ lại pixel (Repaint)}, giúp hiệu ứng đạt mượt mà chuẩn \textbf{60 FPS} ngay cả trên thiết bị di động cấu hình yếu.
\end{tipBox}

\begin{warnBox}[3. Bộ Câu Hỏi Phỏng Vấn Senior Frontend Về Opacity]
\begin{enumerate}[leftmargin=1.5em, itemsep=5pt, parsep=0pt]
    \item \textbf{Câu hỏi 1: Làm thế nào để thẻ cha mờ nhưng thẻ con bên trong vẫn giữ nguyên opacity 100\%?} \par
    $\rightarrow$ \textbf{Trả lời:} Không thể đặt \codeinline{opacity} trên thẻ cha rồi chỉnh thẻ con \codeinline{opacity: 1} vì opacity mang tính chất nhân dồn ($0.5 \times 1.0 = 0.5$). Giải pháp là dùng màu \codeinline{rgba()} cho nền thẻ cha, hoặc tách thẻ con ra thành phần tử đồng cấp rồi xếp chồng bằng CSS \codeinline{position: absolute}.

    \item \textbf{Câu hỏi 2: Tại sao khi đổi opacity từ 1 về 0 lại không click được nếu dùng pointer-events: none?} \par
    $\rightarrow$ \textbf{Trả lời:} Thuộc tính \codeinline{pointer-events: none} giúp phần tử hoàn toàn bỏ qua mọi sự kiện chuột, giúp các sự kiện click xuyên qua phần tử vô hình để tác động tới phần tử bên dưới.

    \item \textbf{Câu hỏi 3: Sự khác biệt bản chất giữa opacity: 0 và display: none trong SEO và Accessibility (SEO \& Màn hình đọc Screen Reader)?} \par
    $\rightarrow$ \textbf{Trả lời:} \codeinline{display: none} ẩn hoàn toàn với cả người dùng lẫn Màn hình đọc (\term{Screen Reader}). Trong khi \codeinline{opacity: 0} vẫn được trình đọc màn hình đọc dữ liệu văn bản bên trong cho người khiếm thị.
\end{enumerate}
\end{warnBox}


% ==============================================================================
% MỤC 11: THUỘC TÍNH CURSOR TRONG CSS (CSS CURSOR MASTERCLASS)
% ==============================================================================
\section{Thuộc Tính Cursor Trong CSS (CSS Cursor Masterclass)}
\label{sec:css_cursor_masterclass}

Thuộc tính \codeinline{cursor} trong CSS cho phép bạn thay đổi kiểu con trỏ chuột khi người dùng di chuyển chuột lên phần tử HTML. Nó giúp cải thiện trải nghiệm người dùng bằng cách cho biết hành vi có thể xảy ra (ví dụ: click được, không cho phép, đang chờ...).

\paragraph{1. Cú Pháp Cơ Bản Của Thuộc Tính Cursor:}

\begin{conceptBox}[Khái Niệm Thuộc Tính Cursor Trong CSS]
\textbf{Thuộc tính \codeinline{cursor}} trong CSS quy định hình dáng hiển thị của con trỏ chuột khi người dùng di chuyển chuột lên phần tử HTML. Nó giúp cải thiện trải nghiệm người dùng bằng cách cho biết hành vi có thể xảy ra (ví dụ: click được, không cho phép, đang chờ...).
\end{conceptBox}

\begin{codeWindow}[Cú Pháp Cơ Bản Của Thuộc Tính Cursor]
\begin{lstlisting}[language=CSS3]
selector {
  cursor: gia_tri_cursor;
}
\end{lstlisting}
\end{codeWindow}

\paragraph{2. Các Giá Trị Phổ Biến Của Cursor:}

\begin{prettyTableBox}[Các Giá Trị Phổ Biến Của Thuộc Tính Cursor Trong CSS]
\rowcolors{2}{white}{primaryBlue!4}
\sffamily\small
\begin{tabularx}{\linewidth}{|>{\raggedright\arraybackslash \ttfamily\color{primaryBlue}\bfseries}K{3.2cm}|Z|>{\sffamily\small}Z|}
\hline
\rowcolor{primaryBlue} \thd{Giá Trị Cursor} & \thd{Ý Nghĩa / Biểu Tượng Con Trỏ Chuột} & \thd{Ví Dụ Tình Huống Áp Dụng} \\
\hline
default & Mũi tên bình thường & Mặc định trên văn bản tĩnh \\
\hline
pointer & Bàn tay chỉ & Dùng cho nút, liên kết, phần tử có thể click \\
\hline
text & Hình chữ I (\term{I-beam}) & Dùng cho vùng nhập liệu như \codeinline{<input>}, \codeinline{<textarea>} \\
\hline
move & Mũi tên 4 chiều & Dùng cho phần tử có thể kéo (\term{drag}) \\
\hline
not-allowed & Vòng tròn gạch chéo & Khi chức năng bị vô hiệu hóa \\
\hline
wait & Đồng hồ cát & Khi đang xử lý, chờ phản hồi \\
\hline
help & Dấu chấm hỏi & Dùng cho tooltip hoặc hỗ trợ \\
\hline
crosshair & Dấu cộng (+) nhỏ & Dùng trong vẽ, bản đồ, hoặc chọn tọa độ \\
\hline
grab / grabbing & Bàn tay nắm hờ / đang nắm & Khi kéo một phần tử (\term{drag \& drop}) \\
\hline
none & Không hiển thị con trỏ chuột & Ẩn con trỏ (cẩn thận khi dùng) \\
\hline
auto & Trình duyệt tự chọn phù hợp & Giá trị mặc định khi không khai báo cụ thể \\
\hline
url(image.png) & Sử dụng ảnh tùy chỉnh làm con trỏ & Game, ứng dụng đặc biệt (phải nhỏ gọn và định dạng \codeinline{.cur}, \codeinline{.png}) \\
\hline
\end{tabularx}
\end{prettyTableBox}

\paragraph{3. Các Ví Dụ Thực Tế Về Khai Báo Cursor:}

\begin{codeWindow}[4 Tình Huống Sử Dụng Cursor Thực Tế]
\begin{lstlisting}[language=HTML5]
<!-- 1. Nut co the nhan: -->
<button style="cursor: pointer;">Click me</button>

<!-- 2. Vung dang xu ly: -->
<style>
.loading-area {
  cursor: wait;
}

/* 3. Keo tha: */
.draggable {
  cursor: grab;
}
.draggable:active {
  cursor: grabbing;
}

/* 4. An con tro chuot: */
.hide-cursor {
  cursor: none;
}
</style>
\end{lstlisting}
\end{codeWindow}

\begin{browserWindow}[Mô Phỏng Trực Quan Kết Quả Hiển Thị 4 Tình Huống Cursor]
\sffamily\small
\begin{tcolorbox}[enhanced, colback=white, colframe=primaryBlue!25, arc=6pt, boxrule=0.8pt, top=10pt, bottom=10pt, left=12pt, right=12pt]
  \begin{tabularx}{\linewidth}{Y Y Y Y}
    \tcbox[on line, arc=4pt, colback=primaryBlue, colframe=primaryBlue, boxsep=3pt, left=8pt, right=8pt, top=4pt, bottom=4pt]{\color{white}\bfseries Click me} &
    \tcbox[on line, arc=4pt, colback=gray!15, colframe=gray!40, boxsep=3pt, left=8pt, right=8pt, top=4pt, bottom=4pt]{\color{gray!80!black}\bfseries Đang xử lý...} &
    \tcbox[on line, arc=4pt, colback=orange!10!white, colframe=orange!60!black, boxsep=3pt, left=8pt, right=8pt, top=4pt, bottom=4pt]{\color{orange!90!black}\bfseries Kéo thả tôi} &
    \tcbox[on line, arc=4pt, colback=red!10, colframe=red!40, boxsep=3pt, left=8pt, right=8pt, top=4pt, bottom=4pt]{\color{red!80!black}\bfseries Ẩn con trỏ} \\
    \vspace{2pt}\textcolor{primaryBlue}{\small (cursor: pointer)} &
    \vspace{2pt}\textcolor{gray!80!black}{\small (cursor: wait)} &
    \vspace{2pt}\textcolor{orange!90!black}{\small (cursor: grab)} &
    \vspace{2pt}\textcolor{red!80!black}{\small (cursor: none)} \\
  \end{tabularx}
\end{tcolorbox}
\end{browserWindow}

\paragraph{4. Demo Tổ Hợp Các Kiểu Con Trỏ Chuột:}

\begin{codeWindow}[Mã Demo Tổ Hợp Các Kiểu Con Trỏ]
\begin{lstlisting}[language=HTML5]
<div style="cursor: pointer;">Click duoc</div>
<div style="cursor: not-allowed;">Khong cho phep</div>
<div style="cursor: help;">Can tro giup</div>
<div style="cursor: text;">Vung nhap lieu</div>
<div style="cursor: crosshair;">Ve toa do</div>
\end{lstlisting}
\end{codeWindow}

\begin{browserWindow}[Mô Phỏng Trực Quan Demo Tổ Hợp Các Kiểu Con Trỏ Chuột]
\sffamily\small
\begin{tcolorbox}[enhanced, colback=white, colframe=primaryBlue!25, arc=6pt, boxrule=0.8pt, top=10pt, bottom=10pt, left=12pt, right=12pt]
  \begin{tabularx}{\linewidth}{|X|c|c|}
    \hline
    \rowcolor{primaryBlue!10} \textbf{Phần Tử \& Văn Bản Hiển Thị} & \textbf{Biểu Tượng Con Trỏ} & \textbf{Kiểu Khai Báo CSS} \\
    \hline
    \textbf{Click được} & Bàn tay chỉ & \codeinline{cursor: pointer;} \\
    \hline
    \textbf{Không cho phép} & Vòng tròn gạch chéo & \codeinline{cursor: not-allowed;} \\
    \hline
    \textbf{Cần trợ giúp} & Dấu chấm hỏi & \codeinline{cursor: help;} \\
    \hline
    \textbf{Vùng nhập liệu} & Con trỏ chữ I & \codeinline{cursor: text;} \\
    \hline
    \textbf{Vẽ tọa độ} & Dấu cộng (+) nhỏ & \codeinline{cursor: crosshair;} \\
    \hline
  \end{tabularx}
\end{tcolorbox}
\end{browserWindow}

\begin{warnBox}[Lưu Ý Khi Sử Dụng Thuộc Tính Cursor]
\begin{itemize}[leftmargin=1.5em, itemsep=3pt, parsep=0pt]
    \item \textbf{Phản ánh tình trạng logic}: Nên dùng \codeinline{cursor} để phản ánh tình trạng logic: nếu phần tử không thể nhấn, không nên dùng \codeinline{pointer}.
    \item \textbf{Hỗ trợ trình duyệt}: Một số giá trị \codeinline{cursor} như \codeinline{grab}, \codeinline{grabbing}, \codeinline{not-allowed} cần trình duyệt hiện đại (Edge, Chrome, Firefox, Safari đều hỗ trợ).
\end{itemize}
\end{warnBox}

\paragraph{5. Chuyên Sâu: Thuộc Tính cursor: pointer; Là Gì?}

\begin{conceptBox}[Khái Niệm \& Cách Hoạt Động Của cursor: pointer;]
\textbf{\codeinline{cursor: pointer;}} là một thuộc tính CSS dùng để thay đổi con trỏ chuột thành biểu tượng bàn tay chỉ khi người dùng đưa chuột vào một phần tử -- biểu thị rằng phần tử đó có thể nhấp được (\term{clickable}). Khi bạn áp dụng cho một phần tử (ví dụ: nút, liên kết, ảnh...), khi rê chuột vào phần tử đó, con trỏ chuột sẽ biến thành bàn tay -- giống như khi rê vào thẻ \codeinline{<a>}.
\end{conceptBox}

\begin{codeWindow}[Cú Pháp Khai Báo cursor: pointer;]
\begin{lstlisting}[language=CSS3]
.element {
  cursor: pointer;
}
\end{lstlisting}
\end{codeWindow}

\paragraph{Các Tình Huống Nên Dùng \codeinline{cursor: pointer;}:}
\begin{itemize}[leftmargin=1.5em, itemsep=3pt, parsep=0pt]
    \item \textbf{Nút tùy chỉnh}: Các thẻ không phải \codeinline{<button>} hoặc \codeinline{<a>} nhưng dùng như nút (ví dụ: \codeinline{<div class="btn">}, \codeinline{<span class="clickable">}).
    \item \textbf{Ảnh có thể nhấn}: \codeinline{<img onclick="..." />}
    \item \textbf{Các vùng có sự kiện onclick}: \codeinline{<div onclick="...">}, \codeinline{<li>}, \codeinline{<section>} gắn \codeinline{onclick}.
\end{itemize}

\paragraph{Ví Dụ Thực Tế Với Nút Tùy Chỉnh:}

\begin{codeWindow}[Mã Nguồn Nút Tùy Chỉnh Sử Dụng cursor: pointer;]
\begin{lstlisting}[language=HTML5]
<style>
  .custom-button {
    background: #007bff;
    color: white;
    padding: 10px 16px;
    border-radius: 6px;
    display: inline-block;
    cursor: pointer;
    user-select: none;
  }
  .custom-button:hover {
    background: #0056b3;
  }
</style>
<div class="custom-button" onclick="alert('Ban da nhan!')">
  Nhan vao toi!
</div>
\end{lstlisting}
\end{codeWindow}

\begin{browserWindow}[Mô Phỏng Trực Quan Giao Diện Nút Tùy Chỉnh (.custom-button)]
\sffamily\small
\centering
\begin{tcolorbox}[enhanced, colback=white, colframe=primaryBlue!25, arc=8pt, boxrule=0.8pt, width=0.85\linewidth, top=12pt, bottom=12pt, left=16pt, right=16pt]
  \centering
  \tcbox[on line, arc=6pt, colback=primaryBlue, colframe=primaryBlue, boxsep=4pt, left=14pt, right=14pt, top=6pt, bottom=6pt, drop shadow=black!10!white]{%
    \color{white}\bfseries\sffamily Nhấn vào tôi!%
  }
  \par\vspace{8pt}
  \textcolor{secondaryTeal}{\small ($\rightarrow$ Khi di chuột vào nút trên: Con trỏ biến thành hình bàn tay, người dùng nhận biết ngay có thể click).}
\end{tcolorbox}
\end{browserWindow}

\begin{noteBox}[Trải Nghiệm Người Dùng \& Lưu Ý Quan Trọng]
Khi rê chuột vào \codeinline{div}, con trỏ sẽ biến thành bàn tay, tạo trải nghiệm giống nút bấm.
\par\vspace{4pt}
\textbf{Lưu ý quan trọng:}
\begin{itemize}[leftmargin=1.5em, itemsep=3pt, parsep=0pt]
    \item \codeinline{cursor: pointer} chỉ thay đổi biểu tượng con trỏ, \textbf{không làm phần tử có hành động click} nếu không có JS hoặc không phải \codeinline{<a>}, \codeinline{<button>}.
    \item Nên kết hợp \codeinline{cursor: pointer;} với xử lý sự kiện click để người dùng không bị nhầm.
\end{itemize}
\end{noteBox}

\subsection{Mục "Đọc Thêm": Mở Rộng Kiến Thức Về Thuộc Tính Cursor}

\begin{conceptBox}[Đọc Thêm: Các Kỹ Thuật Chuyên Sâu Về Cursor Trong Thiết Kế Web]
Mục này mở rộng các kiến thức chuyên sâu về thuộc tính \codeinline{cursor}, giúp lập trình viên tối ưu trải nghiệm người dùng (UX), tăng khả năng truy cập (Accessibility) và hiểu rõ nguyên lý hoạt động của trình duyệt.
\end{conceptBox}

\paragraph{1. Tùy Chỉnh Con Trỏ Bằng Hình Ảnh (\codeinline{cursor: url(...)}) Nâng Cao:}
Trình duyệt cho phép sử dụng hình ảnh tùy chỉnh làm con trỏ chuột bằng hàm \codeinline{url()}:

\begin{codeWindow}[Cú Pháp Tùy Chỉnh Cursor Với Hình Ảnh \& Hotspot]
\begin{lstlisting}[language=CSS3]
/* Cu phap chuan bat buoc co Hotspot coordinate va Fallback: */
.custom-cursor {
    cursor: url('custom-cursor.png') 10 10, auto;
}
\end{lstlisting}
\end{codeWindow}

\begin{itemize}[leftmargin=1.5em, itemsep=3pt, parsep=0pt]
    \item \textbf{Hotspot Coordinate ($x, y$)}: Hai tham số số thực nằm sau đường dẫn ảnh (ví dụ: \codeinline{10 10}) xác định điểm tác động nhấp chuột chính xác trên hình ảnh con trỏ. Nếu không ghi, điểm mặc định là $(0, 0)$ (góc trên bên trái).
    \item \textbf{Giá trị dự phòng (\term{Fallback Value})}: Bắt buộc phải có một con trỏ chuẩn (như \codeinline{auto}, \codeinline{pointer}) đứng sau danh sách \codeinline{url()}. Nếu file ảnh không tải được, trình duyệt sẽ dùng giá trị dự phòng này.
    \item \textbf{Kích thước \& Định dạng ảnh tối ưu}: Nên sử dụng định dạng \codeinline{.cur}, \codeinline{.png} hoặc \codeinline{.svg} có độ phân giải tối ưu $32 \times 32\text{px}$ (hoặc tối đa $64 \times 64\text{px}$). Hầu hết trình duyệt sẽ từ chối hiển thị các file ảnh con trỏ có kích thước quá lớn vì lý do hiệu năng và UX.
\end{itemize}

\paragraph{2. Tương Tác Trên Thiết Bị Cảm Ứng (Touch Devices):}
Trên thiết bị di động (Smartphone/Tablet):
\begin{itemize}[leftmargin=1.5em, itemsep=3pt, parsep=0pt]
    \item Không có khái niệm con trỏ chuột vật lý di chuyển (\term{hover state}), do đó thuộc tính \codeinline{cursor} bị trình duyệt di động hoàn toàn bỏ qua.
    \item Không nên xây dựng luồng tương tác quan trọng chỉ dựa vào trạng thái \codeinline{:hover} hoặc sự thay đổi con trỏ chuột.
    \item Có thể dùng CSS Media Query \codeinline{@media (hover: hover) and (pointer: fine)} để giới hạn các kiểu con trỏ nâng cao chỉ áp dụng cho thiết bị dùng chuột/trackpad.
\end{itemize}

\paragraph{3. Khả Năng Truy Cập (Accessibility - WCAG) \& Best Practices:}
\begin{itemize}[leftmargin=1.5em, itemsep=3pt, parsep=0pt]
    \item \textbf{Kết hợp với \codeinline{user-select: none}}: Khi tạo nút tùy chỉnh bằng thẻ \codeinline{<div>} hay \codeinline{<span>}, hãy luôn đặt thuộc tính \codeinline{user-select: none} để ngăn văn bản bên trong bị bôi đen vô ý khi người dùng nhấp đúp liên tục.
    \item \textbf{Bổ sung thuộc tính ARIA}: Khi sử dụng \codeinline{cursor: not-allowed} cho phần tử bị vô hiệu hóa, nên kết hợp thêm thuộc tính HTML \codeinline{disabled} hoặc \codeinline{aria-disabled="true"} để công cụ màn hình đọc (\term{Screen Reader}) truyền đạt chính xác cho người dùng khiếm thị.
    \item \textbf{Tránh tạo hiểu nhầm}: Tuyệt đối không dùng \codeinline{cursor: pointer} cho văn bản tĩnh không có khả năng nhấp, hoặc \codeinline{cursor: text} cho các phần tử nút bấm.
\end{itemize}

\paragraph{4. Tip Phỏng Vấn Senior Frontend:}
\begin{tipBox}[Câu Hỏi Phỏng Vấn: Dùng div Với cursor: pointer Hay Dùng Thẻ button Bản Địa?]
\textbf{Câu hỏi:} Có nên dùng \codeinline{<div class="btn" style="cursor: pointer;">} thay thế cho thẻ \codeinline{<button>} bản địa không?
\par\vspace{4pt}
\textbf{Trả lời chuẩn phỏng vấn:} \textbf{Không nên.} Mặc dù \codeinline{cursor: pointer} giúp thẻ \codeinline{<div>} hiển thị con trỏ bàn tay giống nút bấm, thẻ \codeinline{<div>} thiếu các tính năng Accessibility \& Keyboard Navigation bản địa như: không kích hoạt được bằng phím \codeinline{Enter}/\codeinline{Space}, không tự động nhận lấy tiêu điểm (\term{focus state}), không submit được form và không thông báo đúng loại phần tử cho trình đọc màn hình (\term{Screen Readers}). Hãy luôn ưu tiên dùng thẻ \codeinline{<button>} cho mọi hành động tương tác nút bấm.
\end{tipBox}


% ==============================================================================
% CHƯƠNG 7: BỐ CỤC MODERN CSS: FLEXBOX & CSS GRID
% ==============================================================================
\chapter{Bố Cục Modern CSS: Flexbox \& CSS Grid}
\label{chap:flexbox_grid}

\section{Bố Cục Một Chiều Với CSS Flexbox}
Các thuộc tính Container: \codeinline{display: flex}, \codeinline{flex-direction}, \codeinline{justify-content}, \codeinline{align-items}, \codeinline{gap}.
Thuộc tính Flex Items: \codeinline{flex-grow}, \codeinline{flex-shrink}, \codeinline{flex-basis}.

\section{Bố Cục Hai Chiều Với CSS Grid}
Khai báo lưới \codeinline{grid-template-columns}, \codeinline{grid-template-rows}, kỹ thuật \codeinline{repeat(auto-fit, minmax(250px, 1fr))} giúp responsive tự động không cần Media Query.


% ==============================================================================
% CHƯƠNG 8: RESPONSIVE WEB DESIGN NÂNG CAO
% ==============================================================================
\chapter{Responsive Web Design Nâng Cao}
\label{chap:responsive_design}

\section{Tư Duy Mobile-First}
Xây dựng giao diện từ màn hình nhỏ nhất (Mobile), mở rộng dần sang Tablet và Desktop để tối ưu hiệu năng.

\section{Fluid Typography \& Container Queries}
Sử dụng hàm CSS \codeinline{clamp(1rem, 2.5vw, 2rem)} cho cỡ chữ linh hoạt và tính năng hiện đại \term{Container Queries}.


% ==============================================================================
% CHƯƠNG 9: CSS VARIABLES & BEM ARCHITECTURE
% ==============================================================================
\chapter{CSS Variables \& BEM Architecture}
\label{chap:css_variables_bem}

\section{CSS Custom Properties (CSS Variables)}
Khai báo biến toán cục trong \codeinline{:root} để làm Dark Mode và Design System.

\section{Phương Pháp Đặt Tên BEM (Block - Element - Modifier)}
Tổ chức CSS sạch sẽ, tránh xung đột tên lớp bằng BEM (ví dụ: \codeinline{.card\_\_title--active}).


% ==============================================================================
% CHƯƠNG 10: CSS ANIMATIONS & TRANSITIONS
% ==============================================================================
\chapter{CSS Animations \& Transitions}
\label{chap:animations_transitions}

\section{Hiệu Ứng Chuyển Động Với CSS Transitions}
Thuộc tính \codeinline{transition}, \codeinline{transition-timing-function} (\codeinline{ease-in-out}, \codeinline{cubic-bezier}).

\section{Keyframe Animations}
Viết animation phức tạp bằng \codeinline{@keyframes} và biến đổi không gian 2D/3D \codeinline{transform}.


\end{document}
